\documentclass[a4paper,11pt]{report}

% Moderne deutsche Sprachunterstuetzung
\usepackage[ngerman]{babel}
\usepackage[T1]{fontenc}
\usepackage[utf8]{inputenc}
\usepackage{lmodern}

% Farben und Mathematik
\usepackage[usenames,dvipsnames]{xcolor}
\usepackage{amssymb}
\usepackage{amsmath}
\usepackage{mathtools}
\usepackage{wasysym}

% Grafiken und Diagramme
\usepackage{graphicx}
\usepackage{tikz}
\usepackage{pgfplots}
\usetikzlibrary{patterns}
\pgfplotsset{compat=1.18}

% Seitenlayout
\usepackage{geometry}
\geometry{
	left=1cm,
	right=2cm,
	top=1cm,
	bottom=2cm,
	bindingoffset=2mm
}

% Links am Ende des Vorspanns
\usepackage[
colorlinks=true,
urlcolor=blue,
linkcolor=black
]{hyperref}
\setlength{\parindent}{0pt}
\begin{document}

\setcounter{secnumdepth}{3}
\setcounter{tocdepth}{3}
	

	
%	\maketitle
	\setcounter{secnumdepth}{3}
	\setcounter{tocdepth}{3}
	
	
	
	%	\maketitle
	\clearpage
	\thispagestyle{empty}
	
	\vspace*{\fill}
	
\begin{center}
	
	{\fontsize{24}{29}\selectfont
		\textbf{Analysis I}}\\[0.8em]
	
	{\fontsize{16}{20}\selectfont
		Grundlagen, Methoden und Anwendungen}\\[2em]
	
	{\fontsize{14}{18}\selectfont
		Dipl.-Ing.\,(FH) Christian Weilharter}\\[2em]
	
	\textcopyright~2025, Christian Weilharter, Traunstein\\[2em]
		
		%\begin{flushleft}
		\begin{tabular}{@{}l l}
			
			
			%	\textbf{ISBN: (Print)} & 978-3-912302-00-4 \\[0.5em]
			\textbf{ISBN: (E-Book)} & 978-3-912302-20-2 \\[0.5em]
			\textbf{Auflage:} & 1.~Auflage 2027 \\[0.5em]
			\textbf{Satz:} & \LaTeX \\[0.5em]
			\textbf{Verlag:} & Christian Weilharter \\[0.5em]
			%	\textbf{Druck:} & Amazon KDP (Print-on-Demand) \\[0.5em]
			
			\textbf{E-Book-Ausgabe:} & Zenodo\\[0.5em]
			%	&Amazon Kindle Direct Publishing \\[0.5em]
			\textbf{Kontakt:} & \href{mailto:info@mathandphysics.de}{info@mathandphysics.de}\\[0.5em]
			\textbf{Web:} & \href{https://www.mathandphysics.de}{www.mathandphysics.de}\\
			
			
		\end{tabular}
		%	\end{flushleft}
		\end{center}
		
		
		\vspace*{\fill}
		
		\clearpage
	
	
	
	% Vorwort
	\chapter*{Vorwort}
	\addcontentsline{toc}{chapter}{Vorwort}

	
	Dieses Skript entstand ursprünglich im Jahr 2020 als persönliche Ausarbeitung zu den damals verfügbaren Lernvideos von OnlineTutorium.com im Bereich Analysis I. Ziel war es, die dort behandelten mathematischen Inhalte in kompakter Form zusammenzufassen und durch eigene Rechnungen und Notizen nachvollziehbar zu machen.
	
	Die ursprünglichen Videos dienten dabei als didaktische Grundlage und Orientierung für die Reihenfolge einzelner Themen. Das vorliegende Skript wurde seitdem überarbeitet, ergänzt und in eine eigenständige Fassung überführt. Erklärungen, mathematische Darstellungen und Beispiele wurden teilweise neu formuliert, erweitert oder ersetzt.
	
	Im Mittelpunkt steht nicht eine möglichst umfassende Darstellung der Analysis, sondern ein verständlicher und rechnerisch orientierter Zugang zu ihren grundlegenden Methoden. Definitionen und Sätze werden daher möglichst direkt mit Herleitungen und Beispielen verbunden.
	
	Das Skript erhebt keinen Anspruch, ein vollständiges Lehrbuch zu ersetzen. Es soll vielmehr eine kompakte Begleitung zum Studium und eine Möglichkeit zur Wiederholung zentraler Inhalte der Analysis I bieten.
	
	Mein Dank gilt den damaligen Autoren und Dozenten von OnlineTutorium.com für die verständliche Vermittlung der mathematischen Grundlagen und die Anregung zur ursprünglichen Erstellung dieses Skripts.
	
	Christian Weilharter
	November 2026

	
	
	
	\clearpage
	
	% Inhaltsverzeichnis
	
	
	\tableofcontents
	\hyphenation{exis-tiert}
	
	% Die 15 Originaldateien muessen neben dieser Hauptdatei liegen.
	
\chapter{Mengen \& Abbildungen}
\section{Einführung in die Analysis I}

Die Analysis ist ein grundlegendes Teilgebiet
der Mathematik. Sie beschäftigt sich insbesondere
mit der Untersuchung von Funktionen, ihren
Eigenschaften und ihrem Verhalten.

Dieses Skript behandelt die wesentlichen
Grundlagen der Analysis I und führt schrittweise
in die wichtigsten mathematischen Methoden ein.

Ausgehend von Mengen, Abbildungen und der
vollständigen Induktion werden zunächst die
komplexen Zahlen sowie Folgen und deren
Konvergenz untersucht.

Darauf aufbauend folgen Grenzwerte und
Stetigkeit von Funktionen, elementare Funktionen,
Differentiation und höhere Ableitungen.

Weitere Schwerpunkte bilden die Integration,
die Untersuchung unendlicher Reihen und
Potenzreihen sowie die Taylorentwicklung.

Ergänzend wird die Fourieranalyse als
weiterführende Methode zur Darstellung
periodischer Funktionen behandelt.

Im Mittelpunkt stehen nachvollziehbare
Herleitungen, anschauliche Beispiele und
Übungsaufgaben. Den Abschluss bildet ein
Kapitel mit Aufgaben zur Wiederholung
und Vertiefung der behandelten Inhalte.
\section{Mengen }

Georg Cantor beschrieb 1895 eine Menge sinngemäß als Zusammenfassung
bestimmter, wohlunterschiedener Objekte zu einem Ganzen.

\textbf{Beispiele:}

\begin{align*}
 S &:= \{\text{Berlin},\text{ Hamburg},\text{ Köln},\text{ Hannover}\},\\
 \sqrt{2} &\in \mathbb{R},\qquad \{\sqrt{2}\}\subseteq\mathbb{R},\\
 \emptyset &:= \{\}\subseteq\mathbb{R},\\
 \infty &\notin\mathbb{R},\\
 \{1,2,7,\text{Berlin}\} &\nsubseteq\mathbb{R}.
\end{align*}

Die Menge $\mathbb{R}$ umfasst die rationalen und die irrationalen
Zahlen. Beispielsweise gehören $\frac{1}{2}$, $\sqrt{2}$ und $\pi$
zu den reellen Zahlen.

\textbf{Wichtige Zahlenmengen:}

\begin{align*}
 \mathbb{N} &:= \{0,1,2,3,\ldots\},\\
 \mathbb{Z} &:= \{\ldots,-2,-1,0,1,2,\ldots\},\\
 \mathbb{Q} &:= \left\{\frac{p}{q}\mid p\in\mathbb{Z},\;
                       q\in\mathbb{N}\setminus\{0\}\right\},\\
 \mathbb{C} &:= \{a+\mathrm{i}b\mid a,b\in\mathbb{R},\;\mathrm{i}^2=-1\}.
\end{align*}

\textbf{Hinweis:} Die natürlichen Zahlen werden nicht überall gleich
definiert. In diesem Skript gilt $0\in\mathbb{N}$. Alternativ wird
häufig $\mathbb{N}=\{1,2,3,\ldots\}$ und
$\mathbb{N}_0=\{0,1,2,\ldots\}$ verwendet.

Es gelten die Inklusionen
\[
 \mathbb{N}\subseteq\mathbb{Z}\subseteq\mathbb{Q}
 \subseteq\mathbb{R}\subseteq\mathbb{C}.
\]

\textbf{Mathematische Schreibweisen:}

\begin{align*}
 x\in A &\quad\text{bedeutet: $x$ ist Element von $A$,}\\
 x\notin A &\quad\text{bedeutet: $x$ ist kein Element von $A$,}\\
 A\subseteq B &\quad\text{bedeutet: $A$ ist Teilmenge von $B$,}\\
 \emptyset &\quad\text{bezeichnet die leere Menge.}
\end{align*}
In einer Mengenbeschreibung bedeutet $\mid$ so viel wie
„mit der Eigenschaft“.

\subsection{Intervalle}

\textbf{Intervalle:}\\
Seien $a,b\in\mathbb{R}$ mit $a<b$.\\
$\blacktriangleright$ Abgeschlossenes Intervall\\  
\begin{quote}
	$[a,b]:=\{x\in\mathbb{R}\;|\;a\leq x \leq b\}$\vspace*{4mm}\\
	\begin{tikzpicture}
		\node at (2,1) [rotate=0,color=red] {$[$};
		\node at (2,0.5) [rotate=0, color=red] {$a$};
		\draw[->] (0,1) -- (6,1) node [below]{$\mathbb{R}$};
		\draw  (2,1) -- (4,1)[ color=red] ;
		\node at (4,1) [rotate=0,color=red] {$]$};
		\node at (4,0.5) [rotate=0, color=red] {$b$};
	\end{tikzpicture}\\
	z.B.: $[1,3]$\\
	\begin{tikzpicture}
		\node at (2,1) [rotate=0,color=red] {$[$};
		\node at (2,0.5) [rotate=0, color=red] {$1$};
		\draw[->] (0,1) -- (6,1) node [below]{$\mathbb{R}$};
		\draw  (2,1) -- (4,1)[ color=red] ;
		\node at (4,1) [rotate=0,color=red] {$]$};
		\node at (4,0.5) [rotate=0, color=red] {$3$};
	\end{tikzpicture}\\
\end{quote}
$\blacktriangleright$ Rechts halboffenes Intervall\\ 
\begin{quote}
	$[a,b[:=\{x\in\mathbb{R}\;|\;a\leq x < b\}$\vspace*{4mm}\\
	\begin{tikzpicture}
		\node at (2,1) [rotate=0,color=red] {$[$};
		\node at (2,0.5) [rotate=0, color=red] {$a$};
		\draw[->] (0,1) -- (6,1) node [below]{$\mathbb{R}$};
		\draw  (2,1) -- (4,1)[ color=red] ;
		\node at (4,1) [rotate=0,color=black] {$[$};
		\node at (4,0.5) [rotate=0, color=black] {$b$};
	\end{tikzpicture}\\
	$b\notin[a,b[$\\
	(analog: das links halboffene Intervall: $\quad  ]a,b]$) \\
\end{quote}
$\blacktriangleright$ Offenes Intervall\\ 
\begin{quote}
	$]a,b[:=\{x\in\mathbb{R}\;|\;a<x < b\}$\vspace*{4mm}\\
	\begin{tikzpicture}
		\node at (2,1) [rotate=0,color=black] {$]$};
		\node at (2,0.5) [rotate=0, color=black] {$a$};
		\draw[->] (0,1) -- (6,1) node [below]{$\mathbb{R}$};
		\draw  (2,1) -- (4,1)[ color=red] ;
		\node at (4,1) [rotate=0,color=black] {$[$};
		\node at (4,0.5) [rotate=0, color=black] {$b$};
	\end{tikzpicture}\ \\
\end{quote}
$\blacktriangleright$ Spezielle Teilintervalle\\
\begin{quote}
	Geometrisch gedeutet ist dies eine Teilgerade der reellen Zahlen.\vspace*{4mm}\\ $]-\infty,a[:=\{x\in\mathbb{R}\;|\;x<a\}$\vspace*{4mm}\\
	\begin{tikzpicture}
		
		
		\draw[<->] (0,1) -- (6,1) node [below]{$\mathbb{R}$};
		\draw  (0,1) -- (4,1)[ color=red] ;
		\node at (4,1) [rotate=0,color=black] {$[$};
		\node at (4,0.5) [rotate=0, color=black] {$a$};
	\end{tikzpicture}\ \\
\end{quote}
Das offene Intervall $]a,b[$ wird auch als $(a,b)$ geschrieben.\\
Das rechts halboffene Intervall $[a,b[$ wird auch als $[a,b)$ geschrieben.\\
Das links halboffene Intervall $]a,b]$ wird auch als $(a,b]$ geschrieben.\\
\subsection{Durchschnitt, Vereinigung und Differenzmenge}

\textbf{Durchschnitt, Vereinigung und Differenzmenge}\\
$\blacktriangleright$ Durchschnitt\\
\begin{quote}
	
	\tikzset{
		schraffiert/.style={pattern=horizontal lines,pattern color=#1},
		schraffiert/.default=black
	}
	\begin{tikzpicture}
		\def\kreisUeins{(0.5,2) circle [radius=1.5]}
		\def\kreisUzwei{(1.5,2) circle [radius=1.5]}
		\begin{scope} % lokal halten der Wirkung von \clip
			\clip\kreisUeins;
			\draw[schraffiert=green]\kreisUzwei;
		\end{scope}
		\draw\kreisUeins;
		\draw\kreisUzwei;
		\node[anchor=east] at (-0.3,2) {$A$};
		\node[anchor=west] at (2.2,2) {$B$};
	\end{tikzpicture}\\
	
	$\qquad \quad A\cap B$\\
	Beispiel:\\
	Bilde: $[1,2] \cap [2,3]$\\
	\begin{tikzpicture}
		
		
		\draw[->] (0,1) -- (6,1) node [below]{$\mathbb{R}$};
		
		\node at (1,1) [rotate=0,color=black] {$[$};
		\node at (1,0.5) [rotate=0, color=black] {$1$};
		\node at (2,1) [rotate=0,color=black] {$[$};
		\node at (2,0.5) [rotate=0, color=black] {$2$};
		
		
	\end{tikzpicture}\ \\
	\begin{tikzpicture}
		
		
		\draw[->] (0,1) -- (6,1) node [below]{$\mathbb{R}$};
		
		\node at (2,1) [rotate=0,color=black] {$[$};
		\node at (2,0.5) [rotate=0, color=black] {$2$};
		\node at (3,1) [rotate=0,color=black] {$]$};
		\node at (3,0.5) [rotate=0, color=black] {$3$};
		
		
	\end{tikzpicture}\ \\
	$[1,2] \cap [2,3]=\{2\}$\\
\end{quote}
$\blacktriangleright$ Vereinigung\\
\begin{quote}
	
	\tikzset{
		schraffiert/.style={pattern=horizontal lines,pattern color=#1},
		schraffiert/.default=black
	}
	
	\begin{tikzpicture}
		\def\kreisUeins{(0.5,2) circle [radius=1.5]}
		\def\kreisUzwei{(1.5,2) circle [radius=1.5]}
		\begin{scope} % lokal halten der Wirkung von \clip
			%\clip\kreisUeins;
			\draw[schraffiert=green]\kreisUzwei;
			\draw[schraffiert=green]\kreisUeins;
		\end{scope}
		\draw\kreisUeins;
		\draw\kreisUzwei;
		\node[anchor=east] at (-0.3,2) {$A$};
		\node[anchor=west] at (2.2,2) {$B$};
	\end{tikzpicture}\\
	Beispiel:\\
	Bilde: $[1,2]\cup [2,3] $\\
	\begin{tikzpicture}
		
		
		\draw[->] (0,1) -- (6,1) node [below]{$\mathbb{R}$};
		
		\node at (1,1) [rotate=0,color=black] {$[$};
		\node at (1,0.5) [rotate=0, color=black] {$1$};
		\node at (2,1) [rotate=0,color=black] {$]$};
		\node at (2,0.5) [rotate=0, color=black] {$2$};
		
		
	\end{tikzpicture}\ \\
	\begin{tikzpicture}
		
		
		\draw[->] (0,1) -- (6,1) node [below]{$\mathbb{R}$};
		
		\node at (2,1) [rotate=0,color=black] {$[$};
		\node at (2,0.5) [rotate=0, color=black] {$2$};
		\node at (3,1) [rotate=0,color=black] {$]$};
		\node at (3,0.5) [rotate=0, color=black] {$3$};
		
		
	\end{tikzpicture}\ \\
	$[1,2]\cup [2,3]=[1,3]$\\
\end{quote}
\pagebreak
$\blacktriangleright$ Differenzmenge\\
\begin{quote}
	
	\tikzset{
		schraffiert/.style={pattern=horizontal lines,pattern color=#1},
		schraffiert/.default=black
	}
	\begin{tikzpicture}
		\def\kreisUeins{(0.5,2) circle [radius=1.5]}
		\def\kreisUzwei{(1.5,2) circle [radius=1.5]}
		\begin{scope} % lokal halten der Wirkung von \clip
			\clip\kreisUeins;
			\draw[schraffiert=green]\kreisUzwei;
			\draw[schraffiert=green]\kreisUeins;
		\end{scope}
		\draw\kreisUeins;
		\draw\kreisUzwei;
		\fill[white] (1.5,2) circle (1.5);
		\draw (0.5,2) circle (1.5);
		\node[anchor=east] at (-0.3,2) {$A$};
		\node[anchor=west] at (2.2,2) {$B$};
	\end{tikzpicture}\\
	Beispiel:\\
	$[1,3] \; \backslash \;] \frac{3}{2},\frac{5}{2} ]$	\\
	\begin{tikzpicture}
		
		
		\draw[->] (0,1) -- (6,1) node [below]{$\mathbb{R}$};
		
		\node at (1,1) [rotate=0,color=black] {$[$};
		\node at (1,0.5) [rotate=0, color=black] {$1$};
		\node at (3,1) [rotate=0,color=black] {$]$};
		\node at (3,0.5) [rotate=0, color=black] {$3$};
		
		
	\end{tikzpicture}\ \\
	\begin{tikzpicture}
		
		
		\draw[->] (0,1) -- (6,1) node [below]{$\mathbb{R}$};
		
		\node at (1.5,1) [rotate=0,color=black] {$]$};
		\node at (1.5,0.5) [rotate=0, color=black] {$\frac{3}{2}$};
		\node at (2.5,1) [rotate=0,color=black] {$]$};
		\node at (2.5,0.5) [rotate=0, color=black] {$\frac{5}{2}$};
		
		
	\end{tikzpicture}\ \\
	$[1,3] \;\backslash \;] \frac{3}{2},\frac{5}{2} ]=[1,\frac{3}{2}] \cup ]\frac{5}{2}, 3]$	\\
	
	\begin{tikzpicture}
		
		
		\draw[->] (0,1) -- (6,1) node [below]{$\mathbb{R}$};
		
		\node at (1,1) [rotate=0,color=black] {$[$};
		\node at (1,0.5) [rotate=0, color=black] {$1$};
		\node at (1.5,1) [rotate=0,color=black] {$]$};
		\node at (1.5,0.5) [rotate=0, color=black] {$\frac{3}{2}$};
		\node at (2.5,1) [rotate=0,color=black] {$]$};
		\node at (2.5,0.5) [rotate=0, color=black] {$\frac{5}{2}$};
		\node at (3,1) [rotate=0,color=black] {$]$};
		\node at (3,0.5) [rotate=0, color=black] {$3$};
	\end{tikzpicture}\ \\	
\end{quote}
\subsection{Aufgabe zu Mengen 1}

\textbf{Aufgabe:}

Bestimme die folgende Menge $\mathbb{M}$:
\[
 \mathbb{M}=\{x\in\mathbb{Z}\mid 3\text{ teilt }x\}
 \cap\{x\in\mathbb{Z}\mid 4\text{ teilt }x\}.
\]

\textbf{Lösung:}

Die durch 3 bzw. 4 teilbaren ganzen Zahlen bilden die Mengen
\[
 \{3n\mid n\in\mathbb{Z}\}
 \quad\text{bzw.}\quad
 \{4n\mid n\in\mathbb{Z}\}.
\]
Da 3 und 4 teilerfremd sind, beträgt ihr kleinstes gemeinsames
Vielfaches $\operatorname{kgV}(3,4)=12$. Genau die Vielfachen von 12
sind daher durch beide Zahlen teilbar. Folglich gilt:
\[
 \boxed{\mathbb{M}=\{12n\mid n\in\mathbb{Z}\}.}
\]

\subsection{Aufgabe zu Mengen 2}

\textbf{Aufgabe:}

Bestimme die folgende Menge $\mathbb{M}$:

\begin{align*}
	\mathbb{M}
	= \bigcup_{m\in\mathbb{N}\setminus\{0\}}
	\{x\in\mathbb{Q}\mid m\cdot x\in\mathbb{Z}\}
\end{align*}

\textbf{Lösung:}

Betrachten wir zunächst den Fall $m=2$:

\begin{align*}
	M_2=\{x\in\mathbb{Q}\mid 2x\in\mathbb{Z}\}
\end{align*}

Welche rationalen Zahlen ergeben mit 2 multipliziert
eine ganze Zahl?

Beispiele:

\begin{align*}
	\frac{1}{2}\cdot 2 &= 1\\
	-\frac{1}{2}\cdot 2 &= -1\\
	\frac{7}{2}\cdot 2 &= 7
\end{align*}

Damit erhalten wir:

\begin{align*}
	M_2
	= \left\{\frac{n}{2}\mid n\in\mathbb{Z}\right\}
\end{align*}

Für ein beliebiges festes
$m\in\mathbb{N}\setminus\{0\}$ gilt entsprechend:

\begin{align*}
	M_m
	&= \{x\in\mathbb{Q}\mid mx\in\mathbb{Z}\}\\
	&= \left\{\frac{n}{m}\mid n\in\mathbb{Z}\right\}
\end{align*}

Vereinigt man diese Mengen über alle positiven
natürlichen Zahlen $m$, so erhält man sämtliche
rationalen Zahlen.

Denn jede rationale Zahl lässt sich als Bruch
$\frac{n}{m}$ mit $n\in\mathbb{Z}$ und
$m\in\mathbb{N}\setminus\{0\}$ darstellen.

Somit gilt:

\begin{align*}
	\boxed{
		\mathbb{M}
		= \bigcup_{m=1}^{\infty} M_m
		= \bigcup_{m=1}^{\infty}
		\left\{\frac{n}{m}\mid n\in\mathbb{Z}\right\}
		= \mathbb{Q}
	}
\end{align*}



\subsection{Aufgabe zu Mengen 3}

\textbf{Aufgabe:}

Bestimme die folgende Menge:
\[
 \mathbb{M}=\bigcap_{m\in\mathbb{R}}
 \{x\in\mathbb{R}\mid(x-m)(x-2)(x-3)=0\}.
\]

\textbf{Lösung:}

Für $m=1$ gilt:
\[
 \{x\in\mathbb{R}\mid(x-1)(x-2)(x-3)=0\}=\{1,2,3\}.
\]
Für $m=\sqrt{2}$ gilt:
\[
 \{x\in\mathbb{R}\mid(x-\sqrt{2})(x-2)(x-3)=0\}
 =\{\sqrt{2},2,3\}.
\]
Für ein beliebiges $m\in\mathbb{R}$ ist die jeweilige Menge
\[
 M_m=\{x\in\mathbb{R}\mid(x-m)(x-2)(x-3)=0\}
 =\{m,2,3\}.
\]
Dabei gilt auch für $m=2$ oder $m=3$ die Mengenschreibweise
$M_m=\{m,2,3\}$; doppelte Elemente werden in Mengen nicht mehrfach
gezählt.

Die Werte $x=2$ und $x=3$ erfüllen die Gleichung für jedes
$m\in\mathbb{R}$, weil jeweils einer der beiden festen Faktoren
verschwindet. Für ein $x\notin\{2,3\}$ lässt sich dagegen ein
$m\neq x$ wählen. Dann ist keiner der drei Faktoren null.
Somit ist kein weiterer Wert in allen Mengen enthalten.

Es folgt:
\[
 \boxed{\mathbb{M}=\bigcap_{m\in\mathbb{R}}M_m=\{2,3\}.}
\]


\section{Abbildungen/Funktionen}

\textbf{Definition (Abbildung oder Funktion):}

Seien $A$ und $B$ zwei Mengen. Eine Abbildung
$f:A\rightarrow B$ ordnet jedem $x\in A$ genau
ein Element $f(x)\in B$ zu:
\[
f:A\rightarrow B,\qquad x\mapsto f(x).
\]
$A$ heißt Definitionsbereich und $B$ Zielmenge.
Das tatsächliche Bild ist eine Teilmenge der Zielmenge.
Für $A\subseteq\mathbb R$ und $B\subseteq\mathbb R$
liegt insbesondere eine reellwertige Funktion
einer reellen Variablen vor.

Beispiel:\\
\textbullet $\qquad f_1:\mathbb{R}\rightarrow \mathbb{R}\text{ ,  }  x\rightarrow 2x $\vspace*{2mm}\\
\begin{center}
	\begin{tikzpicture}
		
		% Gitternetzlinien
		\draw[xstep=1cm, ystep=1cm, very thin, color=lightgray] (-4,4) grid (4,-2);
		% x-Achse
		\draw[->] (-4,0) -- (4,0) node[below,xshift=-1em] {$\mathbb{R}$}; 
		%Zahlen auf x-Achse
		\foreach \x in {-3,...,-1}
		\draw[shift={(\x,0)},color=black] (0pt,2pt) -- (0pt,-2pt) node[below] {\footnotesize $\x$};
		\foreach \x in {1,...,3}
		\draw[shift={(\x,0)},color=black] (0pt,2pt) -- (0pt,-2pt) node[below] {\footnotesize $\x$};
		% y-Achse 
		\draw[->] (0,-2) -- (0,4) node[left,yshift=-1em] {$\mathbb{R}$};%node[above left]
		%Zahlen auf y-Achse
		\foreach \y in {-1,...,-1}
		\draw[shift={(0,\y)},color=black] (2pt,0pt) -- (-2pt,0pt) node[left] {\footnotesize $\y$};
		\foreach \y in {1,...,3}
		\draw[shift={(0,\y)},color=black] (2pt,0pt) -- (-2pt,0pt) node[left] {\footnotesize $\y$};
		%Ursprung
		\draw[color=black] (0pt,-10pt) node[right] {\footnotesize $0$};
		
		
		
		\draw[thick,color=black] plot[samples=200, domain=-1:2] 
		(\x,{2*\x}) node[below, xshift=2em]{$f_1(x)=2x$};
		
		
		
	\end{tikzpicture}\vspace*{8mm}\\
\end{center}
\textbullet $\qquad f_2:\mathbb{R}\rightarrow \mathbb{R}\text{ ,  }  x\rightarrow x^2 $\vspace*{2mm}\\
\begin{center}
	\begin{tikzpicture}
		
		% Gitternetzlinien
		\draw[xstep=1cm, ystep=1cm, very thin, color=lightgray] (-4,4) grid (4,-2);
		% x-Achse
		\draw[->] (-4,0) -- (4,0) node[below,xshift=-1em] {$\mathbb{R}$}; 
		%Zahlen auf x-Achse
		\foreach \x in {-3,...,-1}
		\draw[shift={(\x,0)},color=black] (0pt,2pt) -- (0pt,-2pt) node[below] {\footnotesize $\x$};
		\foreach \x in {1,...,3}
		\draw[shift={(\x,0)},color=black] (0pt,2pt) -- (0pt,-2pt) node[below] {\footnotesize $\x$};
		% y-Achse 
		\draw[->] (0,-2) -- (0,4) node[left,yshift=-1em] {$\mathbb{R}$};%node[above left]
		%Zahlen auf y-Achse
		\foreach \y in {-1,...,-1}
		\draw[shift={(0,\y)},color=black] (2pt,0pt) -- (-2pt,0pt) node[left] {\footnotesize $\y$};
		\foreach \y in {1,...,3}
		\draw[shift={(0,\y)},color=black] (2pt,0pt) -- (-2pt,0pt) node[left] {\footnotesize $\y$};
		%Ursprung
		\draw[color=black] (0pt,-10pt) node[right] {\footnotesize $0$};
		
		
		
		\draw[thick,color=black] plot[samples=200, domain=-2:2] 
		(\x,{\x*\x}) node[below, xshift=2.2em]{$f_2(x)=x^2$};
		
		
		
	\end{tikzpicture}\\
\end{center}
\newpage
\noindent
\textbullet  \;\;  Dies ist keine Abbildung.\\
\begin{center}
	
	\begin{tikzpicture}
		
		% Gitternetzlinien
		\draw[xstep=1cm, ystep=1cm, very thin, color=lightgray] (-4,4) grid (4,-2);
		% x-Achse
		\draw[->] (-4,0) -- (4,0) node[below,xshift=-1em] {$\mathbb{R}$}; 
		%Zahlen auf x-Achse
		\foreach \x in {-3,...,-1}
		\draw[shift={(\x,0)},color=black] (0pt,2pt) -- (0pt,-2pt) node[below] {\footnotesize $\x$};
		\foreach \x in {1,...,2}
		\draw[shift={(\x,0)},color=black] (0pt,2pt) -- (0pt,-2pt) node[below] {\footnotesize $\x$};
		% y-Achse 
		\draw[->] (0,-2) -- (0,4) node[left,yshift=-1em] {$\mathbb{R}$};%node[above left]
		%Zahlen auf y-Achse
		\foreach \y in {-1,...,-1}
		\draw[shift={(0,\y)},color=black] (2pt,0pt) -- (-2pt,0pt) node[left] {\footnotesize $\y$};
		\foreach \y in {1,...,3}
		\draw[shift={(0,\y)},color=black] (2pt,0pt) -- (-2pt,0pt) node[left] {\footnotesize $\y$};
		%Ursprung
		\draw[color=black] (0pt,-10pt) node[right] {\footnotesize $0$};
		
		
		
		\draw[thick,color=black] plot[samples=200, domain=0:4] 
		(\x,{sqrt(\x)});
		
		\draw[thick,color=black] plot[samples=200, domain=0:4] 
		(\x,{-sqrt(\x)});
		\draw[color=blue] (-0.7,-1.73) node[right] {\footnotesize $y_2$};
		\draw[color=blue] (-0.7,1.73) node[right] {\footnotesize $y_1$};
		\draw[color=blue,dashed] (3,1.73) -- (3,0);
		\draw[color=blue,dashed] (3,-0.5) -- (3,-1.73);
		\draw[color=blue,dashed] (0,1.73) -- (3,1.73) ;
		\draw[color=blue,dashed] (0,-1.73) -- (3,-1.73);
		\draw[color=blue] (2.75,-10pt) node[right] {\footnotesize $x_0$};
	\end{tikzpicture}\\
\end{center}
Dies ist ein Widerspruch zur Definition einer Abbildung. Den Wert $x_0$ werden zwei Elemente und zwar $y_1$ und $y_2$ zugewiesen.\\
\subsection{Einschränkungen von Abbildungen}

\textbf{Definition (Einschränkung):}

Sei $f:D_f\rightarrow W_f$ eine Abbildung und
$C\subseteq D_f$. Die Einschränkung von $f$ auf $C$
ist die Abbildung
\[
f|_C:C\rightarrow W_f,\qquad
f|_C(x)=f(x)\quad\text{für alle }x\in C.
\]
Die Funktionsvorschrift bleibt gleich, nur der
Definitionsbereich wird eingeschränkt.

Beispiele:\\
\textbullet$\qquad f_2:\mathbb{R}\longmapsto [0,\infty[$ mit $x\longmapsto x^2$\\
\begin{center}
	\begin{tikzpicture}
		
		% Gitternetzlinien
		\draw[xstep=1cm, ystep=1cm, very thin, color=lightgray] (-4,4) grid (4,-2);
		% x-Achse
		\draw[->] (-4,0) -- (4,0) node[below,xshift=-1em] {$x$}; 
		%Zahlen auf x-Achse
		\foreach \x in {-3,...,-1}
		\draw[shift={(\x,0)},color=black] (0pt,2pt) -- (0pt,-2pt) node[below] {\footnotesize $\x$};
		\foreach \x in {1,...,3}
		\draw[shift={(\x,0)},color=black] (0pt,2pt) -- (0pt,-2pt) node[below] {\footnotesize $\x$};
		% y-Achse 
		\draw[->] (0,-2) -- (0,4) node[left,yshift=-1em] {$f(x)$};%node[above left]
		%Zahlen auf y-Achse
		\foreach \y in {-1,...,-1}
		\draw[shift={(0,\y)},color=black] (2pt,0pt) -- (-2pt,0pt) node[left] {\footnotesize $\y$};
		\foreach \y in {1,...,3}
		\draw[shift={(0,\y)},color=black] (2pt,0pt) -- (-2pt,0pt) node[left] {\footnotesize $\y$};
		%Ursprung
		\draw[color=black] (0pt,-10pt) node[right] {\footnotesize $0$};
		
		
		
		\draw[thick,color=black] plot[samples=200, domain=-2:2] 
		(\x,{\x*\x}) node[below, xshift=1em]{$f_2(x)$};
		
		
		
	\end{tikzpicture}\\
\end{center}
\pagebreak
\textbullet $\qquad f_2|_{[0,\infty[}:[0,\infty[\; \longmapsto [0,\infty [\;$ mit $ x \longmapsto x^2 $\\
\begin{center}
	\begin{tikzpicture}
		% Gitternetzlinien
		\draw[xstep=1cm, ystep=1cm, very thin, color=lightgray] (-4,4) grid (4,-2);
		% x-Achse
		\draw[->] (-4,0) -- (4,0) node[below,xshift=-1em] {$x$}; 
		%Zahlen auf x-Achse
		\foreach \x in {-3,...,-1}
		\draw[shift={(\x,0)},color=black] (0pt,2pt) -- (0pt,-2pt) node[below] {\footnotesize $\x$};
		\foreach \x in {1,...,3}
		\draw[shift={(\x,0)},color=black] (0pt,2pt) -- (0pt,-2pt) node[below] {\footnotesize $\x$};
		% y-Achse 
		\draw[->] (0,-2) -- (0,4) node[left,yshift=-1em] {$f(x)$};%node[above left]
		%Zahlen auf y-Achse
		\foreach \y in {-1,...,-1}
		\draw[shift={(0,\y)},color=black] (2pt,0pt) -- (-2pt,0pt) node[left] {\footnotesize $\y$};
		\foreach \y in {1,...,3}
		\draw[shift={(0,\y)},color=black] (2pt,0pt) -- (-2pt,0pt) node[left] {\footnotesize $\y$};
		%Ursprung
		\draw[color=black] (0pt,-10pt) node[right] {\footnotesize $0$};
		\draw[thick,color=black] plot[samples=200, domain=0:2] 
		(\x,{\x*\x}) node[below, xshift=2.5em]{$f_2|_{[0,\infty[}(x)$};
		
		
		
	\end{tikzpicture}\\
\end{center}
Dies ist eine Funktion mit der Einschränkung von $f_2$ auf die Menge:$[0,\infty[$\vspace*{2mm}\\
\textbullet $\qquad f_1:\mathbb{R}\longmapsto\mathbb{R}\;$ mit $x\longmapsto 2x $\\

\begin{center}
	\begin{tikzpicture}
		% Gitternetzlinien
		\draw[xstep=1cm, ystep=1cm, very thin, color=lightgray] (-4,4) grid (4,-2);
		% x-Achse
		\draw[->] (-4,0) -- (4,0) node[below,xshift=-1em] {$x$}; 
		%Zahlen auf x-Achse
		\foreach \x in {-3,...,-1}
		\draw[shift={(\x,0)},color=black] (0pt,2pt) -- (0pt,-2pt) node[below] {\footnotesize $\x$};
		\foreach \x in {1,...,3}
		\draw[shift={(\x,0)},color=black] (0pt,2pt) -- (0pt,-2pt) node[below] {\footnotesize $\x$};
		% y-Achse 
		\draw[->] (0,-2) -- (0,4) node[left,yshift=-1em] {$f(x)$};%node[above left]
		%Zahlen auf y-Achse
		\foreach \y in {-1,...,-1}
		\draw[shift={(0,\y)},color=black] (2pt,0pt) -- (-2pt,0pt) node[left] {\footnotesize $\y$};
		\foreach \y in {1,...,3}
		\draw[shift={(0,\y)},color=black] (2pt,0pt) -- (-2pt,0pt) node[left] {\footnotesize $\y$};
		%Ursprung
		\draw[color=black] (0pt,-10pt) node[right] {\footnotesize $0$};
		\draw[thick,color=black] plot[samples=200, domain=-1:2] 
		(\x,{\x*2}) node[below, xshift=2.5em]{$f_1(x)$};
		
		
		
	\end{tikzpicture}\\
\end{center}
\textbullet $\qquad f_1|_{\mathbb{Z}}:\mathbb{Z}\longmapsto\mathbb{Z}\;$ mit $x\longmapsto 2x $$\qquad$ Bemerkung: ($\longmapsto\mathbb{Z}$, weil $2\cdot x$ immer $\in\mathbb{Z}$ ist)\\

\begin{center}
	\begin{tikzpicture}
		% Gitternetzlinien
		\draw[xstep=1cm, ystep=1cm, very thin, color=lightgray] (-4,4) grid (4,-2);
		% x-Achse
		\draw[->] (-4,0) -- (4,0) node[below,xshift=-1em] {$x$}; 
		%Zahlen auf x-Achse
		\foreach \x in {-3,...,-1}
		\draw[shift={(\x,0)},color=black] (0pt,2pt) -- (0pt,-2pt) node[below] {\footnotesize $\x$};
		\foreach \x in {1,...,3}
		\draw[shift={(\x,0)},color=black] (0pt,2pt) -- (0pt,-2pt) node[below] {\footnotesize $\x$};
		% y-Achse 
		\draw[->] (0,-2) -- (0,4) node[left,yshift=-1em] {$f(x)$};%node[above left]
		%Zahlen auf y-Achse
		\foreach \y in {-1,...,-1}
		\draw[shift={(0,\y)},color=black] (2pt,0pt) -- (-2pt,0pt) node[left] {\footnotesize $\y$};
		\foreach \y in {1,...,3}
		\draw[shift={(0,\y)},color=black] (2pt,0pt) -- (-2pt,0pt) node[left] {\footnotesize $\y$};
		%Ursprung
		\draw[color=black] (0pt,-10pt) node[right] {\footnotesize $0$};
		
		
		\draw[color=blue] (1,2) node {\textbullet};
		\draw[color=blue] (2,4) node {\textbullet};
		\draw[color=blue] (0,0) node {\textbullet};
		\draw[color=blue] (-1,-2) node {\textbullet};
		\draw[color=blue] (2,3) node {$f_1|_{\mathbb{Z}}$};
		
	\end{tikzpicture}\\
\end{center}
\pagebreak
\subsection{Bild einer Abbildung}

\textbf{Definition (Bild):}\\
Sei $f:\mathbb{D}_f\longmapsto\mathbb{W}_f$ eine Abbildung.\\
Das Bild von $f$ ist die Menge:\\
\begin{align*}
	\mathrm{Bild}(f)=f(\mathbb{D}_f)=\{f(x)\mid x\in\mathbb{D}_f\}\subseteq\mathbb{W}_f
\end{align*}
Beispiele:\\
\textbullet$\qquad f_2:\mathbb{R}\longmapsto [0,\infty[ \qquad$ mit $x\longmapsto x^2$\\
\begin{center}
	\begin{tikzpicture}
		
		% Gitternetzlinien
		\draw[xstep=1cm, ystep=1cm, very thin, color=lightgray] (-4,4) grid (4,-2);
		% x-Achse
		\draw[->] (-4,0) -- (4,0) node[below,xshift=-1em] {$x$}; 
		%Zahlen auf x-Achse
		\foreach \x in {-3,...,-1}
		\draw[shift={(\x,0)},color=black] (0pt,2pt) -- (0pt,-2pt) node[below] {\footnotesize $\x$};
		\foreach \x in {1,...,3}
		\draw[shift={(\x,0)},color=black] (0pt,2pt) -- (0pt,-2pt) node[below] {\footnotesize $\x$};
		% y-Achse 
		\draw[->] (0,-2) -- (0,4) node[left,yshift=-1em] {$f(x)$};%node[above left]
		%Zahlen auf y-Achse
		\foreach \y in {-1,...,-1}
		\draw[shift={(0,\y)},color=black] (2pt,0pt) -- (-2pt,0pt) node[left] {\footnotesize $\y$};
		\foreach \y in {1,...,3}
		\draw[shift={(0,\y)},color=black] (2pt,0pt) -- (-2pt,0pt) node[left] {\footnotesize $\y$};
		%Ursprung
		\draw[color=black] (0pt,-10pt) node[right] {\footnotesize $0$};
		
		
		
		\draw[thick,color=black] plot[samples=200, domain=-2:2] 
		(\x,{\x*\x}) node[below, xshift=1em]{$f_2(x)$};
		\draw[thick,color=red] (0,3.7) node[rotate=-90] {$]\infty$};
		\draw[thick,color=red] (0,0) node[rotate=-90] {$]$};
		\draw[thick,color=red] (0,0) -- (0,3.3); 
	\end{tikzpicture}\\
\end{center}
\begin{align*}
	\mathrm{Bild}_{(f_{2})}=\textcolor{red}{[0,\infty[}=\mathbb{W}_f
\end{align*}
\textbullet\quad Bestimme das Bild von $f_2|_{[2,4]}$ für $f_2(x)=x^2$.
\begin{center}
	\begin{tikzpicture}[x=1cm,y=0.25cm]
\draw[step=1, very thin, color=lightgray] (0,0) grid (5,18);
\draw[->] (0,0)--(5,0) node[right] {$x$};
\draw[->] (0,0)--(0,18) node[above] {$f_2(x)$};
\foreach \x in {1,2,3,4}
  \draw (\x,0.2)--(\x,-0.2) node[below] {\small $\x$};
\foreach \y in {4,8,12,16}
  \draw (0.08,\y)--(-0.08,\y) node[left] {\small $\y$};
\draw[very thick,blue,domain=2:4,samples=100]
  plot (\x,{\x*\x});
\fill[blue] (2,4) circle[radius=2.5pt];
\fill[blue] (4,16) circle[radius=2.5pt];
\draw[red,very thick] (0,4)--(0,16);
\end{tikzpicture}\\
\end{center}
\begin{align*}
	\operatorname{Bild}(f_2|_{[2,4]})=[4,16]
\end{align*}
\newpage
\noindent
\textbullet\quad Bestimme das Bild von $f_2|_{\{5\}}$ für $f_2(x)=x^2$.
\begin{center}
	\begin{tikzpicture}[x=0.9cm,y=0.15cm]
\draw[xstep=1,ystep=5, very thin, color=lightgray] (0,0) grid (6,30);
\draw[->] (0,0)--(6,0) node[right] {$x$};
\draw[->] (0,0)--(0,30) node[above] {$f_2(x)$};
\draw (5,0.35)--(5,-0.35) node[below] {$5$};
\draw (0.1,25)--(-0.1,25) node[left] {$25$};
\draw[dashed] (5,0)--(5,25)--(0,25);
\fill[blue] (5,25) circle[radius=3pt];
\end{tikzpicture}\\
\end{center}
\begin{align*}
	\operatorname{Bild}(f_2|_{\{5\}})=\{25\}
\end{align*}
\subsection{Urbild von Abbildungen}

\textbf{Definition (Urbild):}

Für eine Teilmenge $V\subseteq W_f$ ist das Urbild
von $V$ unter $f$ definiert durch
\[
f^{-1}(V)=\{x\in D_f\mid f(x)\in V\}\subseteq D_f.
\]
Das Urbild einer Menge ist auch dann definiert,
wenn $f$ keine Umkehrfunktion besitzt.

\begin{center}
	
	
	\begin{tikzpicture} 
		\draw[color=black] (5,0) ellipse (60pt and 30pt); 
		\draw[color=red] (0,0) ellipse (60pt and 30pt);
		\draw[color=blue] (0.5,0) ellipse (25pt and 15pt);
		\draw[color=green] (5.5,0) ellipse (25pt and 15pt);
		
		\draw[color=black] (4,-.5) node { $\mathbb{W}_f$};
		\draw[color=red] (-1.3,-.5) node { $\mathbb{D}_f$};
		\draw[color=blue] (0.35,0) node { $f^{-1}(\mathbb{V})$};
		\draw[color=green] (5.7,0) node { $\mathbb{V}$};
		
		\draw[->,orange] (1, 0) .. controls(2, 1) and(4,1) .. node[above,sloped] {f}(5, 0);
	\end{tikzpicture}
	
\end{center}
Beispiel:\\
$f_2 :\mathbb{R}\longmapsto[0,\infty[$ mit $x\longmapsto x^2$\\

\textbullet\;\; Bestimme das Urbild von: $f_2^{-1}([0,\infty[)$



\begin{center}
	\begin{tikzpicture}
		
		% Gitternetzlinien
		\draw[xstep=1cm, ystep=1cm, very thin, color=lightgray] (-4,4) grid (4,-2);
		% x-Achse
		\draw[->] (-4,0) -- (4,0) node[below,xshift=-1em] {$x$}; 
		%Zahlen auf x-Achse
		\foreach \x in {-3,...,-1}
		\draw[shift={(\x,0)},color=black] (0pt,2pt) -- (0pt,-2pt) node[below] {\footnotesize $\x$};
		\foreach \x in {1,...,3}
		\draw[shift={(\x,0)},color=black] (0pt,2pt) -- (0pt,-2pt) node[below] {\footnotesize $\x$};
		% y-Achse 
		\draw[->] (0,-2) -- (0,4) node[left,yshift=-1em] {$f(x)$};%node[above left]
		%Zahlen auf y-Achse
		\foreach \y in {-1,...,-1}
		\draw[shift={(0,\y)},color=black] (2pt,0pt) -- (-2pt,0pt) node[left] {\footnotesize $\y$};
		\foreach \y in {1,...,3}
		\draw[shift={(0,\y)},color=black] (2pt,0pt) -- (-2pt,0pt) node[left] {\footnotesize $\y$};
		%Ursprung
		\draw[color=black] (0pt,-10pt) node[right] {\footnotesize $0$};
		
		
		
		\draw[thick,color=black] plot[samples=200, domain=-2:2] 
		(\x,{\x*\x}) node[below, xshift=1em]{$f_2(x)$};
		\draw[thick,color=red] (0,3.7) node[rotate=-90] {$]\infty$};
		\draw[thick,color=red] (0,0) node[rotate=-90] {$]$};
		\draw[thick,color=red] (0,0) -- (0,3.3); 
		\draw[color=blue,thick] (-4,0) -- (4,0) node {$[$};
		\draw[color=blue,thick] (-4,0) node {$]$}; 
	\end{tikzpicture}\\
	
\end{center}
\begin{align*}
	f_2^{-1}(\textcolor{red}{[0,\infty[})=\textcolor{blue}{\mathbb{R}}
\end{align*}

\textbullet\;\; Bestimme das Urbild von: $f_2^{-1}([2,4])$
\begin{center}
	\begin{tikzpicture}
		
		% Gitternetzlinien
		\draw[xstep=1cm, ystep=1cm, very thin, color=lightgray] (-4,5) grid (4,-1);
		% x-Achse
		\draw[->] (-4,0) -- (4,0) node[below,xshift=-1em] {$x$}; 
		%Zahlen auf x-Achse
		\foreach \x in {-3,...,-1}
		\draw[shift={(\x,0)},color=black] (0pt,2pt) -- (0pt,-2pt);
		\foreach \x in {1,...,3}
		\draw[shift={(\x,0)},color=black] (0pt,2pt) -- (0pt,-2pt);
		\draw[color=blue] (-1.8,-15pt) node[right] {\footnotesize $- \sqrt[]{2}$};
		\draw[color=blue] (-2.5,-15pt) node[right] {\footnotesize $-2$};
		\draw[color=blue] (1.2,-15pt) node[right] {\footnotesize $ \sqrt[]{2}$};
		\draw[color=blue] (1.8,-15pt) node[right] {\footnotesize $2$};
		% y-Achse 
		\draw[->] (0,-1) -- (0,5) node[left,yshift=-1em] {$f(x)$};%node[above left]
		%Zahlen auf y-Achse
		%\foreach \y in {0,...,-1}
		%\draw[shift={(0,\y)},color=black] (2pt,0pt) -- (-2pt,0pt) node[left] {\footnotesize $\y$};
		\foreach \y in {1,...,4}
		\draw[shift={(0,\y)},color=black] (2pt,0pt) -- (-2pt,0pt) node[left] {\footnotesize $\y$};
		%Ursprung
		\draw[color=black] (0pt,-10pt) node[right] {\footnotesize $0$};
		
		
		
		\draw[thick,color=black] plot[samples=200, domain=-2.23:2.23] 
		(\x,{\x*\x}) node[below, xshift=1em]{$f_2(x)$};
		\draw[color=green,dashed,thick] (-2,4) -- (2,4);
		\draw[color=green,dashed,thick] (-2,0) -- (-2,4);
		\draw[color=green,dashed,thick] (2,0) -- (2,4);
		\draw[color=red,thick] (0,2)node[rotate=90] {[} -- (0,4);
		\draw[color=green,dashed,thick] (-1.41,2) -- (1.41,2);
		\draw[color=green,dashed,thick] (-1.41,0) -- (-1.41,2);
		\draw[color=green,dashed,thick] (1.41,0) -- (1.41,2);
		\draw[color=blue,thick] (-2,0)node{[} -- (-1.41,0)node{]};
		\draw[color=blue,thick] (1.41,0)node{[} -- (2,0)node{]};
		\draw[color=red,thick] (0,2)node[rotate=90] {[} -- (0,4)node[rotate=90] {]};
		
	\end{tikzpicture}\\
\end{center}
\begin{align*}
	f_2^{-1}(\textcolor{red}{[2,4]})=\textcolor{blue}{[-2,-\sqrt{2}]\cup[\sqrt{2},2]}
\end{align*}
\textbullet\;\; Bestimme das Urbild von: $f_2^{-1}(\{9\})$
\begin{center}
	\begin{tikzpicture}[x=0.85cm,y=0.38cm]
\draw[xstep=1,ystep=1, very thin, color=lightgray] (-4,0) grid (4,11);
\draw[->] (-4,0)--(4,0) node[right] {$x$};
\draw[->] (0,-0.5)--(0,11) node[above] {$f_2(x)$};
\draw[thick,domain=-3.3:3.3,samples=100] plot (\x,{\x*\x});
\draw[red,dashed] (-3,9)--(3,9);
\draw[green!60!black,dashed] (-3,9)--(-3,0) (3,9)--(3,0);
\fill[blue] (-3,0) circle[radius=2.8pt] (3,0) circle[radius=2.8pt];
\node[below] at (-3,0) {$-3$};
\node[below] at (3,0) {$3$};
\node[left] at (0,9) {$9$};
\end{tikzpicture}\\
\end{center}
\begin{align*}
	f_2^{-1}(\textcolor{red}{\{9\}})=\textcolor{blue}{\{-3,3\}}
\end{align*}
\subsection{Aufgabe zum Thema Bild und Urbild}

Sei $f:\mathbb{R}\longmapsto \mathbb{R}$ mit $f(x)=2x^2-2$ gegeben.\vspace*{2mm}\\
\begin{tabular}[t]{lll}% l=linksbündig r=rechtsbündig
	(i) & Bestimme $\mathrm{Bild}(f)$.  \\
	(ii) & Berechne $f^{-1}([6,30])$. \\
	(iii)&Berechne $f^{-1}(\{\frac{23}{7}\})$.\\	
\end{tabular} \vspace*{2mm}\\
zu (i):
\begin{quote}
	$\mathrm{Bild}(f)=?$\\
	Im positiven Bereich, wenn man große x-Werte in $f(x)$ einsetzt laufen die Werte in das Unendliche $\Rightarrow\infty [ $ .\\
	Die Funktion kann nicht kleiner als -2 werden $\Rightarrow [-2$ .\\
	\begin{align*}
		\mathrm{Bild}(f)=[-2,\infty[
	\end{align*}
\end{quote}
\begin{center}
	\begin{tikzpicture}
		
		% Gitternetzlinien
		\draw[xstep=1cm, ystep=1cm, very thin, color=lightgray] (-4,4) grid (4,-3);
		% x-Achse
		\draw[->] (-4,0) -- (4,0) node[below,xshift=-1em] {$x$}; 
		%Zahlen auf x-Achse
		\foreach \x in {-3,...,-1}
		\draw[shift={(\x,0)},color=black] (0pt,2pt) -- (0pt,-2pt) node[below] {\footnotesize $\x$};
		\foreach \x in {1,...,3}
		\draw[shift={(\x,0)},color=black] (0pt,2pt) -- (0pt,-2pt) node[below] {\footnotesize $\x$};
		% y-Achse 
		\draw[->] (0,-3) -- (0,4) node[left,yshift=-1em] {$y$};%node[above left]
		%Zahlen auf y-Achse
		\foreach \y in {-1,...,-1}
		\draw[shift={(0,\y)},color=black] (2pt,0pt) -- (-2pt,0pt) node[left] {\footnotesize $\y$};
		\foreach \y in {1,...,3}
		\draw[shift={(0,\y)},color=black] (2pt,0pt) -- (-2pt,0pt) node[left] {\footnotesize $\y$};
		%Ursprung
		\draw[color=black] (0pt,-10pt) node[right] {\footnotesize $0$};
		
		
		
		\draw[thick,color=black] plot[samples=200, domain=-1.7:1.7] 
		(\x,{2*\x*\x-2}) node[below, xshift=1em]{$f(x)$};
		\draw[thick,color=red] (0,3.7) node[rotate=-90] {$]\infty$};
		\draw[thick,color=red] (0,0) node[rotate=-90] {$]$};
		\draw[thick,color=red] (0,-2) node[rotate=-90] {$]$}-- (0,3.3); 
		\draw[color=blue,thick] (-4,0) -- (4,0) node {$[$};
		\draw[color=blue,thick] (-4,0) node {$]$}; 
	\end{tikzpicture}\\
	
	
\end{center}
zu (ii):
\begin{quote}
	Bestimme Urbild von: $f^{-1}({[6,30]})$\vspace*{2mm}\\
	Berechnung der Intervallgrenzen:\\
	$\blacktriangleright \quad f(x)=6\Leftrightarrow 2x^2-2=6\Leftrightarrow 2x^2=8\Leftrightarrow x^2=4 \Rightarrow x_{1,2}=\pm 2 $\\
	$\blacktriangleright \quad f(x)=30\Leftrightarrow2x^2-2=30\Leftrightarrow x^2=16\Rightarrow x_{1,2}=\pm 4$\\
	\begin{align*}
		f^{-1}([6,30])=[-4,-2]\cup [2,4]
	\end{align*}
\end{quote}
\begin{center}
	\begin{tikzpicture}[x=0.88cm,y=0.15cm]
\draw[xstep=1,ystep=4, very thin, color=lightgray] (-4,-3) grid (4,34);
\draw[->] (-4.7,0)--(4.7,0) node[right] {$x$};
\draw[->] (0,-3)--(0,34) node[above] {$f(x)$};
\draw[thick,domain=-4.15:4.15,samples=120]
  plot (\x,{2*\x*\x-2});
\draw[red,dashed] (-4.3,6)--(4.3,6);
\draw[red,dashed] (-4.3,30)--(4.3,30);
\node[left] at (0,6) {$6$};
\node[left] at (0,30) {$30$};
\draw[blue,ultra thick] (-4,0)--(-2,0) (2,0)--(4,0);
\foreach \x in {-4,-2,2,4}{
 \fill[blue] (\x,0) circle[radius=2.7pt];
 \node[below] at (\x,0) {\small $\x$};
}
\end{tikzpicture}\\
\end{center}
zu (iii):\\
\begin{quote}
	Bestimme Urbild von: $f^{-1}(\{\frac{23}{7}\})$\vspace*{2mm}\\
	Berechnung des Urbildes:\\
	$\blacktriangleright\quad f(x)=\frac{23}{7}	\Leftrightarrow 2x^2-2=\frac{23}{7}\Leftrightarrow 2x^2=\frac{23}{7}+2=\frac{37}{7}\Leftrightarrow x^2=\frac{37}{14}\Rightarrow x_{1,2}=\pm \sqrt{\smash[b]{\scriptscriptstyle\frac{37}{14}}}$\\
	\begin{align*}
		f^{-1}(\{\tfrac{23}{7}\})=\{ -\sqrt{\scriptscriptstyle\frac{37}{14}},\sqrt{\scriptscriptstyle\frac{37}{14}}\}
	\end{align*}
\end{quote}
\begin{center}
	\begin{tikzpicture}[x=1.3cm,y=0.75cm]
\draw[xstep=0.5,ystep=1, very thin, color=lightgray] (-2,-2) grid (2,6);
\draw[->] (-2.4,0)--(2.4,0) node[right] {$x$};
\draw[->] (0,-2.5)--(0,6) node[above] {$f(x)$};
\draw[thick,domain=-2:2,samples=120]
  plot (\x,{2*\x*\x-2});
\draw[red,dashed] (-2,3.285714)--(2,3.285714);
\node[left] at (0,3.285714) {$\frac{23}{7}$};
\draw[green!60!black,dashed]
  (-1.625687,3.285714)--(-1.625687,0)
  (1.625687,3.285714)--(1.625687,0);
\fill[blue] (-1.625687,0) circle[radius=2.5pt];
\fill[blue] (1.625687,0) circle[radius=2.5pt];
\node[below,align=center] at (-1.625687,0)
  {\scriptsize $-\sqrt{\frac{37}{14}}$};
\node[below,align=center] at (1.625687,0)
  {\scriptsize $\sqrt{\frac{37}{14}}$};
\end{tikzpicture}\\
\end{center}

\section{Injektivität, Surjektivität und Bijektivität von Abbildungen}

\subsection{Injektivität von Abbildungen}

\textbf{Definition (injektiv):}\\
Eine Abb. $f:\mathbb{D}_f\longmapsto\mathbb{W}_f$ heißt injektiv, wenn $f(x_1)=f(x_2)$ nur für $x_1=x_2$ gilt.\\
Beispiel:\\
$f_2:\mathbb{R}\longmapsto[0,\infty[$ mit $x\longmapsto x^2$\\
\textbullet \;\;Zeige, dass $f_2$ nicht injektiv ist.\\
\begin{quote}
	$f_2:\mathbb{R}\longmapsto [0,\infty[$ mit $x\longmapsto x^2$\\
	$ f_2(-1)=(-1)^2=1$\\
	$f_2(1)=(1)^2=1$\\
	es muß gelten:\\
	$f_2(x_1)=f(x_2)$ nur für $x_1=x_2$\\
	Widerspruchsbeweis:\\
	Für $x_1=-1$ gilt $f(x_1)=1$ und für $x_2=1$ gilt $f_2(x_2)=1\Rightarrow x_1\neq x_2\wedge f(x_1)=f(x_2).\\\text{Widerspruch zur Definition} \Rightarrow f_2(x) \text{ ist nicht injektiv.}_{\Square} $ \\ 
	
\end{quote}

\begin{center}
	\begin{tikzpicture}
		
		% Gitternetzlinien
		\draw[xstep=1cm, ystep=1cm, very thin, color=lightgray] (-4,4) grid (4,-2);
		% x-Achse
		\draw[->] (-4,0) -- (4,0) node[below,xshift=-1em] {$x$}; 
		%Zahlen auf x-Achse
		\foreach \x in {-3,...,-1}
		\draw[shift={(\x,0)},color=black] (0pt,2pt) -- (0pt,-2pt) node[below] {\footnotesize $\x$};
		\foreach \x in {1,...,3}
		\draw[shift={(\x,0)},color=black] (0pt,2pt) -- (0pt,-2pt) node[below] {\footnotesize $\x$};
		% y-Achse 
		\draw[->] (0,-2) -- (0,4) node[left,yshift=-1em] {$f(x)$};%node[above left]
		%Zahlen auf y-Achse
		\foreach \y in {-1,...,-1}
		\draw[shift={(0,\y)},color=black] (2pt,0pt) -- (-2pt,0pt) node[left] {\footnotesize $\y$};
		\foreach \y in {1,...,3}
		\draw[shift={(0,\y)},color=black] (2pt,0pt) -- (-2pt,0pt) node[left] {\footnotesize $\y$};
		%Ursprung
		\draw[color=black] (0pt,-10pt) node[right] {\footnotesize $0$};
		
		
		
		\draw[thick,color=black] plot[samples=200, domain=-2:2] 
		(\x,{\x*\x}) node[below, xshift=1em]{$f_2(x)$};
		
		\draw[color=blue,thick,dashed] (-1,1) -- (1,1);
		\draw[->,color=blue,thick,dashed] (-1,0) -- (-1,1);
		\draw[->,color=blue,thick,dashed] (1,0) -- (1,1);
		
	\end{tikzpicture}\\
\end{center}
\newpage
\noindent
\textbullet \;\; $f_{2_{|[0,\infty[}}$  ist injektiv .\\
\begin{center}
	\begin{tikzpicture}
		
		% Gitternetzlinien
		\draw[xstep=1cm, ystep=1cm, very thin, color=lightgray] (-4,4) grid (4,-2);
		% x-Achse
		\draw[->] (-4,0) -- (4,0) node[below,xshift=-1em] {$x$}; 
		%Zahlen auf x-Achse
		\foreach \x in {-3,...,-1}
		\draw[shift={(\x,0)},color=black] (0pt,2pt) -- (0pt,-2pt) node[below] {\footnotesize $\x$};
		\foreach \x in {1,...,3}
		\draw[shift={(\x,0)},color=black] (0pt,2pt) -- (0pt,-2pt) node[below] {\footnotesize $\x$};
		% y-Achse 
		\draw[->] (0,-2) -- (0,4) node[left,yshift=-1em] {$f(x)$};%node[above left]
		%Zahlen auf y-Achse
		\foreach \y in {-1,...,-1}
		\draw[shift={(0,\y)},color=black] (2pt,0pt) -- (-2pt,0pt) node[left] {\footnotesize $\y$};
		\foreach \y in {1,...,3}
		\draw[shift={(0,\y)},color=black] (2pt,0pt) -- (-2pt,0pt) node[left] {\footnotesize $\y$};
		%Ursprung
		\draw[color=black] (0pt,-10pt) node[right] {\footnotesize $0$};
		
		
		
		\draw[thick,color=black] plot[samples=200, domain=0:2] 
		(\x,{\x*\x}) node[below, xshift=2em]{$f_{2_{|[0,\infty[}}(x)$};
		
		
		
	\end{tikzpicture}\\
\end{center}
\begin{quote}
	Merke: Die Injektivität hängt stark vom $ \mathbb{D}_f$ ab.
\end{quote}
\newpage
\noindent

Beispiel: $f(x)=2$\\
\begin{center}
	\begin{tikzpicture}
		
		% Gitternetzlinien
		\draw[xstep=1cm, ystep=1cm, very thin, color=lightgray] (-4,4) grid (4,-2);
		% x-Achse
		\draw[->] (-4,0) -- (4,0) node[below,xshift=-1em] {$x$}; 
		%Zahlen auf x-Achse
		\foreach \x in {-3,...,-1}
		\draw[shift={(\x,0)},color=black] (0pt,2pt) -- (0pt,-2pt) node[below] {\footnotesize $\x$};
		\foreach \x in {1,...,3}
		\draw[shift={(\x,0)},color=black] (0pt,2pt) -- (0pt,-2pt) node[below] {\footnotesize $\x$};
		% y-Achse 
		\draw[->] (0,-2) -- (0,4) node[left,yshift=-1em] {$f(x)$};%node[above left]
		%Zahlen auf y-Achse
		\foreach \y in {-1,...,-1}
		\draw[shift={(0,\y)},color=black] (2pt,0pt) -- (-2pt,0pt) node[left] {\footnotesize $\y$};
		\foreach \y in {1,...,3}
		\draw[shift={(0,\y)},color=black] (2pt,0pt) -- (-2pt,0pt) node[left] {\footnotesize $\y$};
		%Ursprung
		\draw[color=black] (0pt,-10pt) node[right] {\footnotesize $0$};
		
		
		
		\draw[thick,color=black] plot[samples=200, domain=-4:4] 
		(\x,{2}) node[below, xshift=-2em]{$f(x)=2$};
		
		
		
	\end{tikzpicture}\\
\end{center}
Diese Funktion ist nicht injektiv. Jedem Wert der $ x$-Achse wird ein Wert an der $y$-Achse zugewiesen.\vspace*{2mm}
\\
Beispiel: $f(x)=3\sin(x)+1$\\
\begin{center}
	\begin{tikzpicture}
		
		% Gitternetzlinien
		\draw[xstep=1cm, ystep=1cm, very thin, color=lightgray] (-4,4) grid (4,-2);
		% x-Achse
		\draw[->] (-4,0) -- (4,0) node[below,xshift=-1em] {$x$}; 
		%Zahlen auf x-Achse
		\foreach \x in {-3,...,-1}
		\draw[shift={(\x,0)},color=black] (0pt,2pt) -- (0pt,-2pt) node[below] {\footnotesize $\x$};
		\foreach \x in {1,...,3}
		\draw[shift={(\x,0)},color=black] (0pt,2pt) -- (0pt,-2pt) node[below] {\footnotesize $\x$};
		% y-Achse 
		\draw[->] (0,-2) -- (0,4) node[left,yshift=-1em] {$f(x)$};%node[above left]
		%Zahlen auf y-Achse
		\foreach \y in {-1,...,-1}
		\draw[shift={(0,\y)},color=black] (2pt,0pt) -- (-2pt,0pt) node[left] {\footnotesize $\y$};
		\foreach \y in {1,...,3}
		\draw[shift={(0,\y)},color=black] (2pt,0pt) -- (-2pt,0pt) node[left] {\footnotesize $\y$};
		%Ursprung
		\draw[color=black] (0pt,-10pt) node[right] {\footnotesize $0$};
		
		
		
		\draw[thick,color=black] plot[samples=200, domain=-4:4] 
		(\x,{3*sin(\x r)+1}) node[below, xshift=-2em]{$f(x)$};
		
		\draw[thick,color=blue] plot[samples=200, domain=-4:-3] 
		(\x,{3*sin(\x r)+1}) node[below, xshift=3em,yshift=2em]{$f_{|[-4,-3]}(x)$};
		\draw[thick,color=red] plot[samples=200, domain=-1:1] 
		(\x,{3*sin(\x r )+1}) node[below, xshift=1.5em,yshift=-1em]{$f_{|[-1,1]}(x)$};
		
	\end{tikzpicture}\\
\end{center}
Die Funktion $f(x)$ ist nicht injektiv. Durch Einschränkung der  Funktion $f(x)$ kann ein Teil injektiv werden. Die Funktion $f_{|[-4,-3]}(x)$ und $f_{|[-1,1]}(x)$ sind injektiv.\\
\subsection{Surjektivität und Bijektivität von Abbildungen}

\textbf{Definition (surjektiv):}\\
Eine Abbildung
$f:\mathbb{D}_f\rightarrow\mathbb{W}_f$
heißt surjektiv, wenn zu jedem Element
$y$ der Zielmenge mindestens ein Element
$x$ des Definitionsbereichs existiert,
sodass $f(x)=y$ gilt.

Mathematisch:

\begin{align*}
	\forall y\in\mathbb{W}_f\;
	\exists x\in\mathbb{D}_f:
	\quad f(x)=y
\end{align*}




Beispiel:\\
\textbullet $\;\;f_2:\mathbb{R}\longmapsto\mathbb{R}$ mit $x\longmapsto x^2$\\
\begin{quote}
	Für $y=-1$ existiert kein $x$ mit $x^2=y=-1$\\
	$f_2$ ist nicht surjektiv.\vspace*{2mm}\\
\end{quote}
\pagebreak
\textbullet$\;\;\widetilde{f_2}: \mathbb{R}\longmapsto [0,\infty,[$ mit $x\longmapsto x^2$\\
\begin{quote}
	Man sucht sich einen beliebigen Wert aus der Bildmenge ($ y=.. $) und zeige, dass gilt: $ \widetilde{f_2}(x)=y $. Da $ y=x^2 $ wählt man $ x=\sqrt{y} $ und setzt dies in $ \widetilde{f_2}(x) $ ein.
	
	Sei $y\in [0,\infty,[$ und wähle $ x=\sqrt{y} $\\
	$\widetilde{f_2}(x)=\widetilde{f_2}(\sqrt{y})=(\sqrt{y})^2=y \mapsto$ surjektiv\\
	
\end{quote}

\begin{center}
	\begin{tikzpicture}
		
		% Gitternetzlinien
		\draw[xstep=1cm, ystep=1cm, very thin, color=lightgray] (-4,4) grid (4,-2);
		% x-Achse
		\draw[->] (-4,0) -- (4,0) node[below,xshift=-1em] {$x$}; 
		%Zahlen auf x-Achse
		\foreach \x in {-3,...,-1}
		\draw[shift={(\x,0)},color=black] (0pt,2pt) -- (0pt,-2pt) node[below] {\footnotesize $\x$};
		\foreach \x in {1,...,3}
		\draw[shift={(\x,0)},color=black] (0pt,2pt) -- (0pt,-2pt) node[below] {\footnotesize $\x$};
		% y-Achse 
		\draw[->] (0,-2) -- (0,4) node[left,yshift=-1em] {$f(x)$};%node[above left]
		%Zahlen auf y-Achse
		\foreach \y in {-1,...,-1}
		\draw[shift={(0,\y)},color=black] (2pt,0pt) -- (-2pt,0pt) node[left] {\footnotesize $\y$};
		\foreach \y in {1,...,3}
		\draw[shift={(0,\y)},color=black] (2pt,0pt) -- (-2pt,0pt) node[left] {\footnotesize $\y$};
		%Ursprung
		\draw[color=black] (0pt,-10pt) node[right] {\footnotesize $0$};
		
		
		
		\draw[thick,color=black] plot[samples=200, domain=-2:2] 
		(\x,{\x*\x}) node[below, xshift=1em]{$f_2(x)$};
		
		\draw[color=blue,thick,dashed] (-1,1) -- (1,1) node[xshift=2em] {$\widetilde{f}_2$};
		\draw[color=blue,thick,dashed] (-1,0) -- (-1,1);
		\draw[color=blue,thick,dashed] (1,0) -- (1,1);
		\draw[color=blue,thick,dashed] (0.05,0) -- (0.05,4);
		\draw[color=blue,thick,dashed] (0,3) node[right] {$\mathbb{W}_{\widetilde{f}_2}$};
		\draw[color=red,thick,dashed] (-0.05,-2) -- (-0.05,4);
		\draw[color=red,thick,dashed] (-1,-1.7) node[right] {$\mathbb{W}_{f_2}$};
	\end{tikzpicture}\\
\end{center}
\begin{quote}
	Nicht für jedes $ \textcolor{red}{y\in \mathbb{W}_{f_2}} $ existiert ein $x\in\mathbb{D}_{f_2}$, deswegen ist die Funktion nicht surjektiv.(z.B: $y=-1$)\\
	Für jedes $\textcolor{blue}{y\in\mathbb{W}_{\widetilde{f}_2}}$ existiert ein $x\in\mathbb{D}_{\widetilde{f}_2}$, deswegen ist die Funktion surjektiv.\\	
\end{quote}
Merke: Surjektivität hängt stark vom $\mathbb{W}_f$ ab.\vspace*{4mm}\\
\textbf{Definition (bijektiv):}\\
Eine Funktion $f:\mathbb{D}_f\longmapsto\mathbb{W}_f$ heißt bijektiv, wenn sie injektiv und surjektiv zugleich ist!\vspace*{2mm}\\
Merke:\\
Es gibt Funktionen, die weder injektiv noch surjektiv sind.\\
Beispiel: $f_2:\mathbb{R}\longmapsto\mathbb{R}$ mit $x\longmapsto x^2$ ist weder injektiv noch surjektiv.\\
\subsection{Aufgaben zur Injektivität und Surjektivität Teil I}

Überprüfe die folgende Funktion auf Injektivität und Surjektivität:
\begin{align*}
	f:\mathbb{R}\longmapsto\mathbb{R} \text{ mit } f(x)=(x-4)^3
\end{align*}
\textbf{Zur Injektivität}\\
Falls man weis, dass die Funktion nicht injektiv ist kann man ein Gegenbeispiel konstruieren.\\
Beweis der Injektivität:\\
Betrachte:\\ $f(x_1)=f(x_2)\Leftrightarrow (x_1-4)^3=(x_2-4)^3$\\
z.z: $x_1=x_2$\\
\underline{1.Fall:} $(x_1-4)\neq 0$\\
dann erhalten wir:\\
\begin{align*}
	\bigg(\frac{(x_2-4)}{(x_1-4)}\bigg)^3=1\Leftrightarrow\frac{(x_2-4)}{(x_1-4)}=1\Leftrightarrow x_2-4=x_1-4 \Leftrightarrow x_2=x_1
\end{align*} 
\underline{2.Fall:}$(x_1-4)=0\Leftrightarrow x_1=4$
\begin{align*}
	0=(x_1-4)^3=(x_2-4)^3=0\Leftrightarrow x_2=4\Leftrightarrow x_1=x_2
\end{align*}
\begin{align*}
	\text{also }x_1=x_{2_{\Square}}\text{    d.h. f ist injektiv}
\end{align*}
\subsection{Aufgaben zur Injektivität und Surjektivität Teil II}

\textbf{Zur Surjektivität}\\
Falls man weis, dass die Funktion nicht surjektiv ist kann man ein Gegenbeispiel konstruieren.\vspace*{2mm}\\
Beweis der Surjektivität:\\
Betrachte:\\ $f:\mathbb{R}\longmapsto\mathbb{R} \text{ mit } f(x)=(x-4)^3$\\
Sei $y\in\mathbb{R}$ beliebig. Es existiert ein $x\in\mathbb{R}$ mit $f(x)$ also mit $(x-4)^3=y$\vspace*{2mm}\\
Die Gleichung auf $x$ auflösen:\\
(Fallunterscheidung wegen negativer Wurzel)\\
\underline{1.Fall:} $y\in [0,\infty[$\\
\begin{quote}
	Betrachte:\\	
	$(x-4)^3=y\Leftrightarrow (x-4)=\sqrt[3]{y}\Leftrightarrow x=\sqrt[3]{y}+4\\$
\end{quote}
\begin{quote}
	Probe:\\
	$ f(x)=f(\sqrt[3]{y}+4)=((\sqrt[3]{y}+4)-4)^3=y$
\end{quote}
\underline{2.Fall:} $y\in \;]-\infty,0[$
\begin{quote}
	Betrachte:\\
	$(x-4)^3=y\Leftrightarrow-(-x+4)^3=y$\\
	Substituiere $y=-\widetilde{y}$\\
	$\Leftrightarrow-(-x+4)^3=-\widetilde{y}\Leftrightarrow(-x+4)^3=\widetilde{y}\Leftrightarrow -x+4=\sqrt[3]{\widetilde{y}}\Leftrightarrow x=-\sqrt[3]{\widetilde{y}}+4$
\end{quote}
\begin{quote}
	Probe:\\
	$f(x)=f(-\sqrt[3]{\widetilde{y}}+4)=((-\sqrt[3]{\widetilde{y}}+4)-4)^3=\big(-\sqrt[3]{\widetilde{y}\widetilde{y}}\big)^3=-\widetilde{y}=y_{_{\Square}}$
\end{quote}
Daraus folgt die Funktion ist surjektiv, und da die Funktion auch injektiv ist, ist die Funktion bijektiv.\\
\section{Komposition von Abbildungen}

\textbf{Definition (Komposition):}

Seien $f:D_f\rightarrow W_f$ und $g:D_g\rightarrow W_g$
zwei Abbildungen. Falls das Bild von $f$ im Definitionsbereich
von $g$ enthalten ist, also
\[
\operatorname{Bild}(f)\subseteq D_g,
\]
ist die Komposition $g\circ f$ auf ganz $D_f$ definiert durch
\[
g\circ f:D_f\rightarrow W_g,
\qquad (g\circ f)(x)=g(f(x)).
\]

Dabei wird zuerst $f$ und anschließend $g$ angewendet.
Es genügt insbesondere, wenn $W_f\subseteq D_g$ gilt.
Die schwächere Bedingung $\operatorname{Bild}(f)\subseteq D_g$
ist jedoch entscheidend, denn die Zielmenge $W_f$ kann
größer als das tatsächliche Bild sein.

% Die Abbildung verdeutlicht die Hintereinanderausführung.
\begin{center}
	\begin{tikzpicture} 
		\draw[color=black] (5,0) ellipse (60pt and 30pt); 
		\draw[color=red] (0,0) ellipse (60pt and 30pt);
		\draw[color=blue] (10,0) ellipse (60pt and 30pt);
		%\draw[color=green] (5.5,0) ellipse (25pt and 15pt);
		
		\draw[color=black] (4.2,-.5) node { $\mathrm{Bild}(f)\subseteq\mathbb{D}_g$};
		\draw[color=red] (-1.3,-.5) node { $\mathbb{D}_f$};
		\draw[color=blue] (9,-0.5) node { $\mathbb{W}_g$};
		%\draw[color=green] (5.7,0) node { $\mathbb{V}$};
		\draw[color=blue,thick](1,0) node {$\bullet$};
		\draw[color=blue,thick](5,0) node {$\bullet$};
		\draw[color=blue,thick](9,0) node {$\bullet$};
		\draw[->,orange] (1, 0) .. controls(2, 1) and(4,1) .. node[above,sloped] {$f$}(5, 0);
		\draw[->,orange] (5, 0) .. controls(6, 1) and(8,1) .. node[above,sloped] {$g$}(9, 0);
		
		\draw[->,orange] (1, 0) .. controls(3, -2.5) and(7,-2.5) .. node[above,sloped] {$g\circ f$}(9, 0);
	\end{tikzpicture}
	
\end{center}

Ist $\operatorname{Bild}(f)$ nicht vollständig in $D_g$ enthalten,
kann $g\circ f$ dennoch auf der Teilmenge
\[
D_{g\circ f}=\{x\in D_f\mid f(x)\in D_g\}
\]
definiert werden. Diese Teilmenge kann auch leer sein.

\textbf{Beispiel:}

Seien
\[
\begin{aligned}
f:]0,1]&\rightarrow[1,\infty[,
& f(x)&=\frac1x,\\
g:[0,\infty[&\rightarrow[0,\infty[,
& g(x)&=\sqrt{x}.
\end{aligned}
\]

Untersuche die Kompositionen (i) $g\circ f$ und (ii) $f\circ g$.

\textbf{(i) Komposition $g\circ f$:}

Es gilt
\[
\operatorname{Bild}(f)=[1,\infty[\ \subseteq\ D_g=[0,\infty[.
\]
Damit ist $g\circ f$ auf dem gesamten Definitionsbereich
von $f$ definiert:
\[
\boxed{
(g\circ f)(x)=g\!\left(\frac1x\right)=\sqrt{\frac1x},
\qquad x\in\,]0,1].
}
\]
Also gilt
\[
g\circ f:\,]0,1]\rightarrow[0,\infty[.
\]
Tatsächlich ist das Bild von $g\circ f$ die Menge $[1,\infty[$.

\textbf{(ii) Komposition $f\circ g$:}

Hier gilt
\[
\operatorname{Bild}(g)=[0,\infty[\ \nsubseteq\ D_f=\,]0,1].
\]
Daher ist $f\circ g$ \emph{nicht auf ganz} $D_g=[0,\infty[$
definiert. Dennoch kann die Komposition auf einer geeigneten
Teilmenge gebildet werden. Aus
\[
\sqrt{x}\in\,]0,1]
\quad\Longleftrightarrow\quad
0<x\leq1
\]
folgt der zulässige Definitionsbereich $D_{f\circ g}=\,]0,1]$.
Somit gilt
\[
\boxed{
(f\circ g)(x)=f(\sqrt{x})=\frac1{\sqrt{x}},
\qquad x\in\,]0,1].
}
\]
Das Bild der eingeschränkten Komposition ist $[1,\infty[$.

\subsection{Aufgabe zur Komposition von Abbildungen}

\textbf{Aufgabe:}

Betrachte die Abbildungen
\[
\begin{aligned}
f:\mathbb R&\rightarrow\mathbb R,
&f(x)&=1+x^4,\\
h:\mathbb R&\rightarrow\mathbb R,
&h(x)&=1+x^3.
\end{aligned}
\]
Untersuche:

(i) Existiert eine Funktion $g$, sodass $h=g\circ f$
für alle $x\in\mathbb R$ gilt?

(ii) Existiert eine reellwertige Funktion $g$, sodass
$h=f\circ g$ für alle $x\in\mathbb R$ gilt?

\textbf{Lösung zu (i):}

Angenommen, eine Funktion $g$ mit $h=g\circ f$
existiert auf ganz $\mathbb R$. Dann müsste gelten:
\[
g(1+x^4)=1+x^3
\qquad\text{für alle }x\in\mathbb R.
\]
Da $f(1)=f(-1)=2$, folgt einerseits
\[
g(2)=h(1)=2,
\]
andererseits aber
\[
g(2)=h(-1)=0.
\]
Dies ist ein Widerspruch. Daher existiert \emph{keine}
solche Funktion $g$ auf ganz $\mathbb R$.

\textbf{Warum funktioniert der zunächst naheliegende Ansatz nicht?}

Für $t\geq1$ könnte man setzen:
\[
g(t)=1+\bigl(\sqrt[4]{t-1}\bigr)^3.
\]
Dann erhält man jedoch
\[
(g\circ f)(x)=1+\bigl(\sqrt[4]{x^4}\bigr)^3
=1+|x|^3.
\]
Dies stimmt nur für $x\geq0$ mit $h(x)=1+x^3$ überein.
Der Ansatz funktioniert also, wenn $f$ und $h$ auf
$[0,\infty[$ eingeschränkt werden; dann kann man
$g:[1,\infty[\rightarrow[1,\infty[$ wie oben definieren.

\textbf{Lösung zu (ii):}

Angenommen, es existiert eine reellwertige Funktion $g$
mit $h=f\circ g$ auf ganz $\mathbb R$. Dann müsste gelten:
\[
1+x^3=f(g(x))=1+(g(x))^4.
\]
Für jedes reelle $g(x)$ gilt $(g(x))^4\geq0$.
Somit wäre $h(x)\geq1$ für alle $x\in\mathbb R$.
Für $x=-1$ gilt jedoch $h(-1)=0$. Dies ist ein Widerspruch.
Daher existiert auch hier \emph{keine} solche Funktion $g$
auf ganz $\mathbb R$.

Für $x\geq0$ ist hingegen
\[
g(x)=x^{3/4}
\]
reell definiert, und es gilt
\[
(f\circ g)(x)=1+\bigl(x^{3/4}\bigr)^4
=1+x^3=h(x).
\]
Auf dem eingeschränkten Definitionsbereich $[0,\infty[$
ist eine solche Darstellung also möglich.

\textbf{Merke:}

Bei Kompositionen sind die Definitionsbereiche entscheidend.
Außerdem ist bei geraden Wurzeln zu beachten:
\[
\sqrt[4]{x^4}=|x|.
\]

\section{Umkehrabbildung}

\textbf{Definition (Umkehrabbildung):}

Seien $D_f$ und $W_f$ Mengen und
$f:D_f\rightarrow W_f$ eine bijektive Abbildung.
Dann gibt es zu jedem $y\in W_f$ genau ein $x\in D_f$
mit $f(x)=y$.

Die Umkehrabbildung ist definiert durch
\[
f^{-1}:W_f\rightarrow D_f,\qquad
y\mapsto x \quad\text{mit } f(x)=y.
\]
Sie wird auch Umkehrfunktion genannt.

Die folgende Zeichnung verdeutlicht die Zuordnung
$x\mapsto f(x)$ und die Rückzuordnung
$y\mapsto f^{-1}(y)$.

\begin{center}
\begin{tikzpicture}
		
		% Gitternetzlinien
		\draw[xstep=1cm, ystep=1cm, very thin, color=lightgray] (-4,4) grid (4,-2);
		% x-Achse
		\draw[->] (-4,0) -- (4,0) node[below,xshift=-1em] {$x$}; 
		%Zahlen auf x-Achse
		\foreach \x in {-3,...,-1}
		\draw[shift={(\x,0)},color=black] (0pt,2pt) -- (0pt,-2pt) node[below] {\footnotesize $\x$};
		\foreach \x in {1,...,3}
		\draw[shift={(\x,0)},color=black] (0pt,2pt) -- (0pt,-2pt) node[below] {\footnotesize $\x$};
		% y-Achse 
		\draw[->] (0,-2) -- (0,4) node[left,yshift=-1em] {$f(x)$};%node[above left]
		%Zahlen auf y-Achse
		\foreach \y in {-1,...,-1}
		\draw[shift={(0,\y)},color=black] (2pt,0pt) -- (-2pt,0pt) node[left] {\footnotesize $\y$};
		\foreach \y in {1,...,3}
		\draw[shift={(0,\y)},color=black] (2pt,0pt) -- (-2pt,0pt) node[left] {\footnotesize $\y$};
		%Ursprung
		\draw[color=black] (0pt,-10pt) node[right] {\footnotesize $0$};
		
		
		
		\draw[thick,color=black] plot[samples=200, domain=0:2] 
		(\x,{\x*\x}) node[below, xshift=2em]{$f(x)$};
		
		\draw[->,color=blue,thick,dashed] (1,0) -- (1,1) ;
		\draw[->,color=blue,thick,dashed] (1,1) --  (0,1) ;
		
		\draw[color=blue,thick,dashed] (1,-0.7) node {$x_1$}  ;
		\draw[color=blue,thick,dashed] (-0.8,1) node {$f(x_1)$}  ;
		
		\draw[->,color=red,thick,dashed] (0,2.25) -- (1.5,2.25) ;
		\draw[->,color=red,thick,dashed] (1.5,2.25) --  (1.5,0) ;
		
		\draw[color=red,thick,dashed] (1.5,-0.7) node {$x_2$}  ;
		\draw[color=red,thick,dashed] (-0.8,2.25) node {$y$}  ;
		
	\end{tikzpicture}
\end{center}

\textbf{Beispiel:}

Wir betrachten die quadratische Funktion mit
eingeschränktem Definitionsbereich:
\begin{align*}
f&:[0,\infty[\ \rightarrow [0,\infty[,
& x&\mapsto x^2,\\
f^{-1}&:[0,\infty[\ \rightarrow [0,\infty[,
& x&\mapsto \sqrt{x}.
\end{align*}

\begin{center}
\begin{tikzpicture}
		% Gitternetzlinien
		\draw[xstep=1cm, ystep=1cm, very thin, color=lightgray] (-4,4) grid (4,-2);
		% x-Achse
		\draw[->] (-4,0) -- (4,0) node[below,xshift=-1em] {$x$}; 
		%Zahlen auf x-Achse
		\foreach \x in {-3,...,-1}
		\draw[shift={(\x,0)},color=black] (0pt,2pt) -- (0pt,-2pt) node[below] {\footnotesize $\x$};
		\foreach \x in {1,...,3}
		\draw[shift={(\x,0)},color=black] (0pt,2pt) -- (0pt,-2pt) node[below] {\footnotesize $\x$};
		% y-Achse 
		\draw[->] (0,-2) -- (0,4) node[left,yshift=-1em] {$f(x)$};%node[above left]
		%Zahlen auf y-Achse
		\foreach \y in {-1,...,-1}
		\draw[shift={(0,\y)},color=black] (2pt,0pt) -- (-2pt,0pt) node[left] {\footnotesize $\y$};
		\foreach \y in {1,...,3}
		\draw[shift={(0,\y)},color=black] (2pt,0pt) -- (-2pt,0pt) node[left] {\footnotesize $\y$};
		%Ursprung
		\draw[color=black] (0pt,-10pt) node[right] {\footnotesize $0$};
		\draw[thick,color=black] plot[samples=200, domain=-2:4] 
		(\x,{\x}) node[below, xshift=-1.5em,yshift=-1.7em]{$g(x)$};
		
		\draw[thick,color=red] plot[samples=200, domain=0:2] 
		(\x,{\x*\x}) node[below, xshift=1.5em]{$f(x)$};
		\draw[thick,color=blue] plot[samples=200, domain=0:4] 
		(\x,{sqrt(\x)}) node[below, xshift=-1.4em,yshift=-.6em]{$f^{-1}(x)$};
	\end{tikzpicture}
\end{center}

\textbf{Eigenschaften:}

Für eine bijektive Funktion gelten die Identitäten
\begin{align*}
(f^{-1}\circ f)(x)&=x
&&\text{für alle }x\in D_f,\\
(f\circ f^{-1})(y)&=y
&&\text{für alle }y\in W_f.
\end{align*}

Die Graphen von $f$ und $f^{-1}$ sind bei reellen
Funktionen Spiegelbilder an der Geraden $y=x$,
sofern beide Graphen in demselben Koordinatensystem
dargestellt werden.

Für das Beispiel $f(x)=x^2$ mit $x\geq 0$ folgt:
\begin{align*}
(f^{-1}\circ f)(x)
&=\sqrt{x^2}=|x|=x
&&\text{für }x\geq 0,\\
(f\circ f^{-1})(y)
&=(\sqrt{y})^2=y
&&\text{für }y\geq 0.
\end{align*}

\textbf{Wichtig:}
Die Umkehrfunktion $f^{-1}$ ist nicht mit der
Kehrwertfunktion $1/f(x)$ zu verwechseln.
Außerdem ist das Urbild $f^{-1}(V)$ einer Menge $V$
auch ohne Umkehrfunktion definiert; die Notation
wird aus dem Zusammenhang unterschieden.

\subsection{Aufgabe zur Umkehrabbildung Teil I}

\textbf{Aufgabe:}

Gegeben ist $f(x)=x^2-2x$.

\begin{enumerate}
\item[(i)] Bestimme möglichst große Intervalle,
auf denen $f$ injektiv und somit bei geeigneter
Zielmenge umkehrbar ist.
\item[(ii)] Bestimme die Umkehrfunktion für
die Einschränkung auf $]-\infty,1]$.
\end{enumerate}

\textbf{Lösung zu (i):}

Durch quadratische Ergänzung ergibt sich:
\[
f(x)=x^2-2x=(x-1)^2-1.
\]
Die Parabel hat ihren Scheitelpunkt bei $(1,-1)$
und die Nullstellen $x=0$ und $x=2$.

\begin{center}
\begin{tikzpicture}
			
			% Gitternetzlinien
			\draw[xstep=1cm, ystep=1cm, very thin, color=lightgray] (-4,4) grid (4,-2);
			% x-Achse
			\draw[->] (-4,0) -- (4,0) node[below,xshift=-1em] {$x$}; 
			%Zahlen auf x-Achse
			\foreach \x in {-3,...,-1}
			\draw[shift={(\x,0)},color=black] (0pt,2pt) -- (0pt,-2pt) node[below] {\footnotesize $\x$};
			\foreach \x in {1,...,3}
			\draw[shift={(\x,0)},color=black] (0pt,2pt) -- (0pt,-2pt) node[below] {\footnotesize $\x$};
			% y-Achse 
			\draw[->] (0,-2) -- (0,4) node[left,yshift=-1em] {$f(x)$};%node[above left]
			%Zahlen auf y-Achse
			\foreach \y in {-1,...,-1}
			\draw[shift={(0,\y)},color=black] (2pt,0pt) -- (-2pt,0pt) node[left] {\footnotesize $\y$};
			\foreach \y in {1,...,3}
			\draw[shift={(0,\y)},color=black] (2pt,0pt) -- (-2pt,0pt) node[left] {\footnotesize $\y$};
			%Ursprung
			\draw[color=black] (0pt,-10pt) node[right] {\footnotesize $0$};
			
			
			
			\draw[thick,color=red] plot[samples=200, domain=-1.2:1] 
			(\x,{\x*\x-2*\x}) node[below, xshift=1em]{$f(x)$};
			
			\draw[thick,color=blue] plot[samples=200, domain=1:3.2] 
			(\x,{\x*\x-2*\x});
			
		\end{tikzpicture}
\end{center}

Auf $]-\infty,1]$ ist die Funktion streng fallend,
auf $[1,\infty[$ streng steigend. Daher ist sie
auf jedem dieser beiden Intervalle injektiv.

Wählt man jeweils $[-1,\infty[$ als Zielmenge,
erhält man zwei bijektive Einschränkungen:
\begin{align*}
f_- &: ]-\infty,1]\rightarrow[-1,\infty[,
&x&\mapsto x^2-2x,\\
f_+ &: [1,\infty[\rightarrow[-1,\infty[,
&x&\mapsto x^2-2x.
\end{align*}

Beide Intervalle sind bezüglich der Injektivität
maximal durch Mengeninklusion.
Keines der beiden ist eine Obermenge des anderen.
Für Teil (ii) wählen wir den linken Ast $f_-$.

\subsection{Aufgabe zur Umkehrabbildung Teil II}

\textbf{Lösung zu (ii):}

Gesucht ist die Umkehrfunktion von
\[
f_-:]-\infty,1]\rightarrow[-1,\infty[,
\qquad f_-(x)=x^2-2x.
\]

Wir setzen $y=f_-(x)$ und lösen nach $x$ auf:
\begin{align*}
y &= x^2-2x,\\
y+1 &= x^2-2x+1=(x-1)^2,\\
x-1 &= \pm\sqrt{y+1},\\
x &= 1\pm\sqrt{y+1}.
\end{align*}

Auf dem gewählten Definitionsbereich gilt $x\leq 1$.
Daher ist der Minuszweig auszuwählen:
\[
\boxed{f_-^{-1}(y)=1-\sqrt{y+1},\qquad y\geq-1.}
\]

\textbf{Probe:}

Für $y\in[-1,\infty[$ folgt:
\begin{align*}
(f_-\circ f_-^{-1})(y)
 &=\bigl(1-\sqrt{y+1}\bigr)^2
   -2\bigl(1-\sqrt{y+1}\bigr)\\
 &=1-2\sqrt{y+1}+(y+1)-2+2\sqrt{y+1}\\
 &=y.
\end{align*}

Für $x\in]-\infty,1]$ gilt:
\begin{align*}
(f_-^{-1}\circ f_-)(x)
 &=1-\sqrt{x^2-2x+1}\\
 &=1-\sqrt{(x-1)^2}\\
 &=1-|x-1|\\
 &=1-(1-x)=x.
\end{align*}
Der letzte Schritt verwendet $x\leq 1$ und damit
$|x-1|=1-x$.

\textbf{Ergebnis:}
\[
\boxed{
f_-^{-1}:[-1,\infty[\rightarrow]-\infty,1],
\qquad y\mapsto1-\sqrt{y+1}.
}
\]

\textbf{Ergänzung:}
Für die ebenfalls maximal injektive Einschränkung
$f_+$ auf $[1,\infty[$ lautet die Umkehrfunktion
\[
f_+^{-1}(y)=1+\sqrt{y+1},
\qquad y\in[-1,\infty[.
\]










    \chapter{Algebraische Grundlagen und Vollständige Induktion}
 %2.1
\section{Summen und Produkte}


\begin{center}
	\fbox{\parbox [t] {90mm} {
			Seien $n,m \in\mathbb{Z}$ mit $m\leqq n$. Dann definieren wir:\\
			\begin{center}
				\begin{quote}
					$\sum\limits_{k=m}^{n} a_k:=a_m+a_{m+1}+\ldots + a_n$\\
					$\prod\limits_{k=m}^{n} a_k:=a_m\cdot a_{m+1}\cdot\ldots\cdot a_n$
				\end{quote}
			\end{center}
	}}	
\end{center}
$k$ ist der Laufindex\\
Summenzeichen $\sum$ ist der griechische Buchstabe Sigma.\\
Produktzeichen $\prod$ ist der griechische Buchstabe Pi.

\begin{center}
	\fbox{\parbox [t] {70mm} {
			Seien $n,m \in\mathbb{Z}$ mit $m > n$. Dann gilt:\\
			\begin{center}
				
				$\sum\limits_{k=m}^{n} a_k:=0$\\
				$\prod\limits_{k=m}^{n} a_k:=1$\\
				
			\end{center}
			
			
	}}
\end{center}
Beispiel:\\
\textbullet$\;\;14=2+3+4+5=\sum\limits_{k=1}^{5}k$\\
\textbullet$\;\;\sum\limits_{k=1}^{0}=0$ (neutrales Element der Addition)\\
\textbullet$\;\;\prod\limits_{k=7}^{-3}=1$ (neutrales Element der Multiplikation)\vspace*{4mm}\\


Definition:\vspace*{2mm}\\
\begin{center}
	\fbox{\parbox  {70mm} {
			Für $n\in\mathbb{N}$ definiert man: n-Falkultät\\
			\begin{align*}
				n!:=\prod\limits_{k=1}^{n}k=1\cdot 2\cdot 3\cdot \ldots\cdot (n-1)n
			\end{align*}
	}}
\end{center}
Beispiel:\\
$5!=1\cdot 2\cdot 3 \cdot 4 \cdot 5=120$\\
$1!=1$\\
$2!=1\cdot2=2$\\
$3!=1 \cdot 2 \cdot 3=6$\\
$0!=\prod\limits_{k=1}^{0}k=1$ da $k>0$\\
\subsection{Rechenregeln für Summen und Produkte}

\begin{center}
	\fbox{\parbox  {90mm} {
			Seien $n,m\in\mathbb{Z}$ mit $m\leq n$ und $j\in\mathbb{N}$, dann gilt:\\
			\begin{align*}
				\sum\limits_{k=m}^{n+j}a_k&=\sum\limits_{k=m}^{n}a_k + \sum\limits_{k=n+1}^{n+j}a_k\\
				\prod\limits_{k=m}^{n+j}a_k&=\prod\limits_{k=m}^{n}a_k \cdot \prod\limits_{k=n+1}^{n+j}a_k
			\end{align*}
	}}
\end{center}
Beispiel:\\

$\underbrace{\sum\limits_{k=2}^{5}k}_{2+3+4+5=14}=\sum\limits_{k=2}^{3+2}k=\underbrace{\sum\limits_{k=2}^{3}k}_{2+3=5}+\underbrace{\sum\limits_{k=3+1}^{5}k}_{4+5=9}$\\
Beispiel:\\
$n!=\prod\limits_{k=1}^{n}k$\\
sei $n\in \mathbb{N}$ mit $m\leqq n$ \\
$n=m+(n-m)$ mit $(n-m)\in \mathbb{N}$\\
dann gilt:\\
$n!=\prod\limits_{k=1}^{n}k=\prod\limits_{k=1}^{\scriptscriptstyle m+(n-m)}k=\underbrace{\prod\limits_{k=1}^{m}k}_{m!}\cdot\prod\limits_{\scriptscriptstyle k=m+1}^{\scriptscriptstyle m+(n-m)}k$\\
$\overset{\div m!}{\Leftrightarrow}\frac{n!}{m!}=\prod\limits_{\scriptscriptstyle k=m+1}^{\scriptscriptstyle m+(n-m)}k=\prod\limits_{\scriptscriptstyle k=m+1}^{n}k=(m+1)\cdot (m+2)\cdot (m+3)\ldots n$\vspace*{2mm}\\
\begin{center}
	\boxed{
		\frac{n!}{m!}=\prod\limits_{\scriptscriptstyle  k=m+1}^{n}k=(m+1)\cdot (m+2)\cdot (m+3)\cdot\ldots\cdot n
	}
\end{center}
\subsection{Aufgabe zu Summen und Produkte}

Berechne folgende Summenen und Produkte.\\
(i)$\;\;\prod\limits_{n=-3}^{3}=4^n $\\
(ii)$\;\sum\limits_{k=1}^{3}\bigg(\prod\limits_{j=1}^{k}i\bigg)  $\vspace*{2mm}\\
\newpage
\noindent
zu(i):
\begin{quote}
	$\prod\limits_{n=-3}^{3}=4^n $\\
	$ 4^n=? $\\
	$ a^b\cdot a^c=a^{b+c} $\\
	also gilt: $ n=\sum\limits_{n=-3}^{3}n=-3+-2+-1+0+1+2+3=0 $\\
	$\prod\limits_{n=-3}^{3}=4^n\Leftrightarrow\prod\limits_{n=-3}^{3}=4^{\big(\scriptscriptstyle\sum\limits_{n=-3}^{n}n\big)}\Leftrightarrow\prod\limits_{n=-3}^{3}=4^0=1$  
\end{quote}
zu (ii):
\begin{quote}
	$ \sum\limits_{k=1}^{3}\bigg(\prod\limits_{j=1}^{k}i\bigg) = \underbrace {\prod\limits_{j=1}^{1}}_{1!}j+\underbrace {\prod\limits_{j=1}^{2}}_{2!}j+\underbrace {\prod\limits_{j=1}^{3}}_{3!}j=1+1\cdot2+1\cdot2\cdot3=1+2+6=9 $
\end{quote}
\subsection{Indexverschiebung von Summen und Produkte}

$ \sum\limits_{k=2}^{3}k=2+3=(5-3) +(6-3)=\sum\limits_{k=5}^{6}(k-3)$\\
\begin{center}
	\fbox{\parbox  {70mm} {
			Seien $n,m\in\mathbb{Z}$ mit $ m<n $ und $ j\in\mathbb{N} $\\
			\begin{align*}
				\sum\limits_{k=m}^{n}=\sum\limits_{\scriptscriptstyle k=m+j}^{n+j}a_{k-j}\\
				\prod\limits_{k=m}^{n}=\prod\limits_{\scriptscriptstyle k=m+j}^{n+j}a_{k-j}
			\end{align*}
	}}\\
	
\end{center}
$ \sum\limits_{k=2}^{k}k=2+3=\sum\limits_{\scriptscriptstyle k=2+3}^{3+3}(k-3)$\vspace*{2mm}\\
Beispiel:\\
\textbullet\;\; Indexverschiebung der Summe $ \sum\limits_{k=2}^{n}3^k  $ um $ +7 $.
\begin{quote}
	$ \sum\limits_{k=2}^{n}3^k=\sum\limits_{\scriptscriptstyle k=2+7}^{n+7}3^{(k-7)}=\sum\limits_{k=9}^{n+7}3^{(k-7)} $
\end{quote}
\textbullet\;\; Indexverschiebung des Produktes $ \prod\limits_{k=5}^k(k-1)^k  $ um $ +5 $.\\
\begin{quote}
	$ \prod\limits_{k=5}^k(k-1)^k=\prod\limits_{k=10}^{n+5}(k-5)((k-5)-1)^{(k-5)}=\prod\limits_{k=10}^{n+5}(k-5)((k-6)^{(k-5)} $
\end{quote}
\subsection{Aufgabe I zur Indexverschiebung}

Berechne:
$ \sum\limits_{k=1}^{25}(1+k)^2-\sum\limits_{s=4}^{28}=s^2-4s+3  $\\
Durch Anpassen der Grenzen durch Indexverschiebung kann man die Summen leichter zusammenfassen.\\
NR.:\\
\begin{quote}
	$\sum\limits_{s=4}^{28}=s^2-4s+3  $ Indexverschiebung um -3\\
	$ \sum\limits_{k=4}^{28}=k^2-4k+3= \sum\limits_{k=1}^{25}=(k+3)^2-4(k+3)+3$\\
\end{quote}
\begin{align*}
	\sum\limits_{k=1}^{25}(1+k)^2-\sum\limits_{s=4}^{28}=s^2-4s+3&=\sum\limits_{k=1}^{25}(1+k)^2-((k+3)^2-4(k+3)+3) \\
	&=\sum\limits_{k=1}^{25}(1+k)^2-(k+3)^2+4(k+3)-3\\
	&=\sum\limits_{k=1}^{25}1+2k+k^2-k^2-6k-9+4k+9\\
	&=\sum\limits_{k=1}^{25}1=1+1+1+\ldots = 25\\
\end{align*}
\subsection{Aufgabe II zur Indexverschiebung}

Schreibe die folgende Summe in der Form:\\
(i) \;\;$ \sum\limits_{k=1}^{n}a_k $ und
(ii) $ \sum\limits_{j=6}^{n+5}b_k $
mit $ -1+3-5+7-9\pm\ldots\mp n $ und $n$ soll ungerade sein.\vspace*{6mm}\\
zu (i): $\sum\limits_{k=1}^{n}a_k  $\\
\begin{quote}
	Ein alternierendes Vorzeichen bekommt man mit der Summe:\\
	$\sum\limits_{k=1}^{n}(-1)^k=(-1)^1+(-1)^2+(-1)^3\pm\ldots\mp n=-1+1-1+1\pm\ldots\mp n$\vspace*{2mm}\\
	Ungerade Zahlen bekommt man mit der Summe:\\
	$\sum\limits_{k=1}^{n}2k+1=3+5+7+....+2k+1 $\\
	$\sum\limits_{k=1}^{n}(-1)^k (2k+1)=-3+5-7\pm\ldots\mp 2k+1$\\
	$\text{Es muß mit der Zahl 1 gestartet werden, also erzeugt man die ungerade Zahl mit }2k-1$\\
	$\sum\limits_{k=1}^{n}(-1)^k (2k-1)=-1+3-5\pm\ldots\mp 2k-1$\\
	\begin{align*}
		\sum\limits_{k=1}^{n}a_k=\sum\limits_{k=1}^{n}(-1)^k (2k-1)
	\end{align*}
\end{quote}
\newpage
\noindent
zu (ii): $ \sum\limits_{j=6}^{n+5}b_k $\\
\begin{quote}
	Mit Indexverschiebung der Summe $\sum\limits_{k=1}^{n}(-1)^k (2k-1)  $ mit 5, bekommt man:\\
	\begin{align*}
		\sum\limits_{k=1+5}^{n+5}(-1)^{k-5} (2(k-5)-1)&=(-1)^k(-1)^{-5}(2k-11)\\&=(-1)^k(-1)(2k-11)\\&=(-1)^{k-1}(2k-11)
	\end{align*}
	
\end{quote}
\begin{align*}
	\sum\limits_{k=6}^{n+5}b_k=\sum\limits_{k=6}^{n+5}(-1)^{k-1}(2k-11)
\end{align*}



%2.2
\section{Vollständige Induktion}

Die vollständige Induktion ist eine Beweismethode um Aussagen zu beweisen.\\
Beweise folgende Aussage: Für alle $ n\in\mathbb{N} $ gilt $ \sum\limits_{k=1}^{n}k=\frac{n(n+1)}{2} $.\\
Anleitung:\\
\textbf{(I.A.) Induktionsanfang}\\
Beweise die Aussage für erstes n:\\
$ n=0:\sum\limits_{k=1}^{0}k=0=\frac{0(0+1)}{2}_{\Square} $\\
(Obere Grenze der Summe ist größer als die untere Grenze, deshalb ist die Summe=0)\\
\textbf{(I.V.) Induktionsvoraussetzung}\\
für ein $ n\in \mathbb{N} $ gilt:\\
\begin{align}
	\sum\limits_{k=1}^{n}k=\frac{n(n+1)}{2}
\end{align}
\textbf{(I.S.) Induktionsschritt}\\
Zu zeigen: Dass die Aussage auch für $ n+1 $ gilt. Unter der Annahme , dass (I.V.) gilt beweise wir die Aussage:\\
\begin{align}
	\sum\limits_{k=1}^{n+1}k=\frac{(n+1)((n+1)+1)}{2}=\frac{n^2+2n+n+2}{2}
\end{align}
Betrachte:\\
\begin{align}
	\sum\limits_{k=1}^{n+1}k=\sum\limits_{k=1}^{n}k+(n+1)&\overset{(I.V.)}{=}\underbrace{\frac{n(n+1)}{2}}_{\text{(II.1)}}+(n+1)\nonumber\\
	&\;\;=\frac{n(n+1)+2(n+1)}{2}\nonumber\\
	&\;\;=\frac{n^2+2n+n+2}{2}\nonumber\\
	\sum\limits_{k=1}^{n+1}k&\;\;=\frac{n^2+2n+n+2}{2}_{\Square}
\end{align}
Man hat gezeigt:\\
Die Aussage (II.2) ist gleich der Betrachtung (II.3) mit Aussnutzung der Induktionsvoraussetzung (II.1).
\subsection{Vollständige Induktion Aufgabe 1}

Beweise: Für $ n\in\mathbb{N} $ gilt $ \sum\limits_{i=1}^{n}(2i-1)=n^2 $\\
\textbf{(I.A.)}\\
$ n=0:\sum\limits_{i=1}^{0}(2i-1)=0=0^2_{\;\Square} $\\
\textbf{(I.V.)}\\
Für ein $ n\in\mathbb{N} $ gilt $\sum\limits_{i=1}^{n}(2i-1)=n^2  $\\
\textbf{(I.S.)}\\
Z.z.: $ \sum\limits_{i=1}^{n+1}(2i-1) =(n+1)^2=n^2+2n+1$\\
Dies muß man jetzt beweisen:\\
Betrachte:\\
\begin{align*}
	\sum\limits_{i=1}^{n+1}(2i-1)=\underbrace {\sum\limits_{i=1}^{n}(2i-1)}_{I.V.}+(2(n+1)-1)\overset{(I.V.)}{=}n^2+(2(n+1)-1)=n^2+2n+1_{\;\Square}
\end{align*}
\subsection{Vollständige Induktion Aufgabe 2}

Beweise: Für alle $ n\in \mathbb{N}$ mit $ n \geq 2 $ gilt:$ \prod\limits_{k=2}^{n}\big(1-\frac{1}{k}\big)=\frac{1}{n} $\\
$ \scriptscriptstyle (k=2 \text{ und } n\geq2 \text{ : Die Zahl 2 ist Zufall!}) $\vspace*{2mm}\\
\textbf{(I.A.)}\\
$ n=2:\prod\limits_{k=2}^{2}\big(1-\frac{1}{2}\big)=\frac{1}{2}_{\;\;\Square} $\\
\textbf{(I.V.)}\\
Es gilt für ein $ n\in \mathbb{N} $ mit $ n\geq 2 $\\
\begin{align*}
	\prod\limits_{k=2}^{n}\bigg(1-\frac{1}{k}\bigg)=\frac{1}{n}
\end{align*}
\textbf{(I.S.)}\\
z.Z: $ \prod\limits_{k=2}^{n+1}\big(1-\frac{1}{k}\big)=\frac{1}{n+1} $\\
\newpage
\noindent
Betrachte:\\
\begin{align*}
	\prod\limits_{k=2}^{n+1}\bigg(1-\frac{1}{k}\bigg)=\underbrace {\prod\limits_{k=2}^{n}\bigg(1-\frac{1}{k}\bigg)}_{(I.V.)}\cdot\bigg(1-\frac{1}{n+1}\bigg) & \overset{(I.V.)}{=}\frac{1}{n}\bigg(1-\frac{1}{n+1}\bigg)\\
	&\;\;=\frac{1}{n} \bigg(\frac{(n+1)-1}{n+1}\bigg)\\
	&\;\;=\frac{1}{n}\bigg(\frac{n}{n+1}\bigg)\\
	\prod\limits_{k=2}^{n+1}\bigg(1-\frac{1}{k}\bigg)&\;\;=\frac{1}{n+1}_{\;\;\Square}\\
\end{align*}
\subsection{Vollständige Induktion Aufgabe 3}

Beweise: Für $ n\geq 1 $ gilt $ \sum\limits_{k=1}^{n} (-1)^k\cdot k^2=(-1)^n \frac{n(n+1)}{2}$\vspace*{2mm}\\
\textbf{I.A.}\\
Für $n=1$ gilt:
\begin{align*}
	\sum\limits_{k=1}^{1}(-1)^1\cdot 1^2=-1=(-1)^1\cdot\frac{1(1+1)}{2}_{\;\;\Square}\\
\end{align*}
\textbf{I.V.}\\
Es gilt für ein $ n\in \mathbb{N} $ mit $ n\geq 1 $\\
\begin{align}
	\sum\limits_{k=1}^{n} (-1)^k\cdot k^2=(-1)^n \frac{n(n+1)}{2}
\end{align}
\textbf{I.S.}\\
Zu zeigen:\\
\begin{align*}
	\sum\limits_{k=1}^{n+1}(-1)^k\cdot k^2=(-1)^{n+1}\cdot \frac{(n+1)((n+1)+1)}{2}=(-1)^{n+1}\cdot \frac{(n+1)(n+2)}{2}=(-1)^{n+1}\cdot\frac{n^2+3n+2}{2}
\end{align*}
\newpage
\noindent
Betrachte:\\
\begin{align*}
	\sum\limits_{k=1}^{n+1}(-1)^k\cdot k^2&=\underbrace {\sum\limits_{k=1}^{n} (-1)^k\cdot k^2}_{\text{(II.4) }}+ (-1)^{n+1}\cdot(n+1)^2\\
	&\overset{(I.V.)}{=}\underbrace {(-1)^n\cdot \frac{n(n+1)}{2}}_{\text{(II.4) }}+(-1)^{n+1} \cdot(n+1)^2\\
	&\;\;=\frac{(-1)^n n(n+1)+2(-1)^{n+1}(n+1)^2 }{2}\\
	&\;\;=\frac{(-1)^n (n^2+n)+2(-1)^{n+1}(n+1)^2 }{2}& |(-1)\cdot(-1) \text{ ergänzen }\\
	&\;\;=\frac{(-1)^n\cdot(-1)\cdot(-1) (n^2+n)+2(-1)^{n+1}(n+1)^2}{2}\\
	&\;\;=\frac{(-1)^{n+1}(-n^2-n)+2(-1)^{n+1}(n+1)^2}{2}\\
	&\;\;=(-1)^{n+1}\cdot\frac{(-n^2-n)+2(n+1)^2}{2}\\
	&\;\;=(-1)^{n+1}\cdot\frac{-n^2-n+2n^2+4n+2}{2}\\
	\sum\limits_{k=1}^{n+1}(-1)^k\cdot k^2&\;\;=(-1)^{n+1}\cdot\frac{n^2+3n+2}{2}_{\;\;\Square}\\
\end{align*}


%2.3
\section{Binomischer Lehrsatz und Geometrische Reihe}

\textbullet\;\;Binomischer Lehrsatz
\begin{center}
	\fbox{\parbox  {66mm} {
			\begin{center}
				Binominalkoeffizient:\\
			\end{center}
			Seien $ n,k \in\mathbb{N} $ und $ k>0 $ dann gilt:\\
			\begin{align*}
				&\text{für }n\geq k \\
				&\binom{n}{k}:=\frac{n!}{k!(n-k)!}\;\;\;\;\in \mathbb{N}\\
				&\binom{n}{n}:=1\\
				&\binom{n}{0}:=1\\
				& \text{für }n<k \\
				&\binom{n}{k}:=0
			\end{align*}
	}}
\end{center}
\newpage
\noindent
Beispiel:\\
\begin{align*}
	\binom{3}{2}&=\frac{3!}{2!(3-2)!}=\frac{1\cdot2\cdot3}{1\cdot2(1)}=3 \\
	\binom{5}{3}&=\frac{5!}{3!(5-3)!}=\frac{1\cdot2\cdot3\cdot4\cdot5}{1\cdot2\cdot3\cdot2}=10 \\
\end{align*}

\begin{center}
	\fbox{\parbox  {76mm} {
			\begin{center}
				Binomischer Satz
			\end{center}
			Seien $ a,b\in \mathbb{R} $, dann gilt für alle $ n \in \mathbb{N} $\\
			\begin{align*}
				(a+b)^n=\sum\limits_{k=0}^{n}\binom{n}{k}a^{n-k} \;\;b^k
			\end{align*}	
	}}
	\
\end{center}
Beispiel:\\
\begin{align*}
	\big(3-\sqrt{2}\big)^2&=\sum\limits_{k=0}^{2}\binom{2}{k}3^{2-k}\;\;\big(-\sqrt{2}\big)^k\\
	&=\binom{2}{0}3^{2-0}\;\;\big(-\sqrt{2}\big)^0+\binom{2}{1}3^{2-1}\;\;\big(-\sqrt{2}\big)^1+\binom{2}{2}3^{2-2}\;\;\big(-\sqrt{2}\big)^2\\
	&=1\cdot 9+2\cdot 3\big(-\sqrt{2}\big)+1\cdot 1\cdot 2\\
	&=9-6\sqrt{2}+2
\end{align*}
\textbullet\;\;Geometrische Reihe
\begin{center}
	\fbox{\parbox  {80mm} {
			\begin{center}
				Geometrische Reihe\vspace*{2mm}\\
				$\sum\limits_{k=0}^{n}x^k =     \left\{ \begin{array}{rcl}
					n+1 & \mbox{für} & x=1\\
					\frac{1-x^{n+1}}{1-x} & \mbox{für} & x\neq  1 \\
					
				\end{array}\right.$	
			\end{center}
	}}
\end{center}
Beispiel:\\
\begin{align*}
	\sum\limits_{k=0}^{3}2^k=\frac{1-2^4}{1-2}=\frac{1-16}{1-2}=15
\end{align*}
\newpage
\subsection{Binomischer Satz Aufgabe}

Zeige mit den binomischen Satz, dass $ \big(1+\sqrt{5}\big)^n+\big(1-\sqrt{5}\big)^n \in\mathbb{N} $ ist.\vspace*{4mm} \\
Für alle $ n \in \mathbb{N} $ gilt:\\
\begin{align*}
	\text{\textbullet }\;\;\big(1+\sqrt{5}\big)^n&=\sum\limits_{k=0}^{n}\binom{n}{k}1^{n-k}\big(\sqrt{5}\big)^k=\sum\limits_{k=0}^{n}\binom{n}{k}\big(\sqrt{5}\big)^k\\
	\text{\textbullet }\;\;\big(1-\sqrt{5}\big)^n&=\sum\limits_{k=0}^{n}\binom{n}{k}1^{n-k}\big(-\sqrt{5}\big)^k=\sum\limits_{k=0}^{n}\binom{n}{k}\big(-\sqrt{5}\big)^k\\
\end{align*}
Damit ist die Aussage äquivalent zu:\\
\begin{align*}
	&\sum\limits_{k=0}^{n}\binom{n}{k}\big(\sqrt{5}\big)^k+\sum\limits_{k=0}^{n}\binom{n}{k}\big(-\sqrt{5}\big)^k \in \mathbb{N}\\
	&\sum\limits_{k=0}^{n}\binom{n}{k}\bigg(\big(\sqrt{5}\big)^k+\big(-\sqrt{5}\big)^k\bigg)& |\big(-\sqrt{5}\big)^k=\underbrace {-1^k}_{Vorzeichenwechsel}\big(\sqrt{5}\big)^k\\
	&\sum\limits_{k=0}^{n}\binom{n}{k}\bigg(\big(\sqrt{5}\big)^k-\big(-1^k\sqrt{5}\big)^k\bigg)\\
	&\sum\limits_{k=0}^{n}\binom{n}{k}\bigg(\sqrt{5}\big)^k(1+(-1)^k\bigg)\\
\end{align*}
Wir zeigen, dass jeder Summant von der Summe aus $\in \mathbb{N} $ ist, dann ist auch die Summe $\in \mathbb{N} $\\
d.h. $ \binom{n}{k}\big(\sqrt{5}\big)^k(1+(-1)^k) \in \mathbb{N}$ für $ k \in \mathbb{N} $\\
$ \binom{n}{k}  \in \mathbb{N}$\\$\Rightarrow$ dann müssen noch die einzelnen Faktoren von $ \big(\sqrt{5}\big)^k\big(1+(-1)^k\big) \in \mathbb{N}$ sein.\vspace*{2mm} \\
\underline{1.Fall (k ist gerade):}\\
$ k=2p $\\
\begin{align*}
	&\big(\sqrt{5}\big)^k\Big(1+(-1)^k\Big)\\
	&\big(\sqrt{5}\big)^{2p}\Big(1+(-1)^{2p}\Big)\\
	&\Big(\big(\sqrt{5}\big)^2\Big)^p\Big(1+\big((-1)^2 \big)^p\Big)\\
	&5^p(1+1^p)\\
	&5^p(2)\in \mathbb{N}_{\;\;\Square}
\end{align*}
\newpage
\noindent
\underline{2.Fall (k ist ungerade):}\\
$ k=2p+1 $\\
\begin{align*}
	&\big(\sqrt{5}\big)^k\big(1+(-1)^k\big)\\
	&\big(\sqrt{5}\big)^{2p+1}\Big(1+\big(-1\big)^{2p+1}\Big)\\
	&\big(\sqrt{5}\big)^{2p}\big(\sqrt{5}\big)^1 \big(1+\big(\big(-1\big)^{2p}(-1)^1\big)\\
	&5^p\sqrt{5}(1+(-1))\\
	&5^p\sqrt{5}\cdot 0 \in \mathbb{N}_{\;\;\Square}
\end{align*}

%2.4
\section{Absolutbetrag und reelle Lösungen von Gleichungen}

Der (Absolut)-Betrag ist die folgende Funktion:\\
\begin{center}
	$|\cdot|:\mathbb{R}\longmapsto [o,\infty[$\;\;  mit $ |x| =     \left\{ \begin{array}{rcl}
		x & \mbox{für} & x\geq0\\
		-x & \mbox{für} & x< 0 \;\;\;\;\;\;\;\; \text{(II.5) }\\
		
	\end{array}\right.$
\end{center}
\begin{center}
	\fbox{\parbox  {60mm} {
			\begin{center}
				Mit den Eigenschaften:\\
			\end{center}
			\begin{align*}
				&|x\cdot y|=|x|\cdot|y|\\
				&\text{\tiny Dreiecksungleichung: }\\
				&|x+y|\leq |x|+|y|
			\end{align*}
	}}
\end{center}
Beispiel:\\
\textbullet\;\;$|-7|=-(-7) =7 $ \tiny{(in Funktion (II.5) einsetzen)}\normalsize \\
\textbullet \;\;$|7|=7$\\
\textbullet \;\;$ x<2:|x+2|=-(x+2) $\\

Funktion: $ f(x)=|x| $\vspace*{2mm}\\
\begin{center}
	
	\begin{tikzpicture}
		
		% Gitternetzlinien
		\draw[xstep=1cm, ystep=1cm, very thin, color=lightgray] (-4,4) grid (4,-2);
		% x-Achse
		\draw[->] (-4,0) -- (4,0) node[below,xshift=-1em] {$x$}; 
		%Zahlen auf x-Achse
		\foreach \x in {-3,...,-1}
		\draw[shift={(\x,0)},color=black] (0pt,2pt) -- (0pt,-2pt) node[below] {\footnotesize $\x$};
		\foreach \x in {1,...,3}
		\draw[shift={(\x,0)},color=black] (0pt,2pt) -- (0pt,-2pt) node[below] {\footnotesize $\x$};
		% y-Achse 
		\draw[->] (0,-2) -- (0,4) node[left,yshift=-1em] {$y$};%node[above left]
		%Zahlen auf y-Achse
		\foreach \y in {-1,...,-1}
		\draw[shift={(0,\y)},color=black] (2pt,0pt) -- (-2pt,0pt) node[left] {\footnotesize $\y$};
		\foreach \y in {1,...,3}
		\draw[shift={(0,\y)},color=black] (2pt,0pt) -- (-2pt,0pt) node[left] {\footnotesize $\y$};
		%Ursprung
		\draw[color=black] (0pt,-10pt) node[right] {\footnotesize $0$};
		
		\draw[thick,color=black] plot[samples=200, domain=0:4] 
		(\x,{\x}) node[below, xshift=-6em] {$f(x)=|x|$};
		
		\draw[thick,color=black] plot[samples=200, domain=-4:0] 
		(\x,{-\x}) node[below, xshift=1em]{};
		
		
		
	\end{tikzpicture}\vspace*{4mm}\\
\end{center}
\newpage
\noindent
Löse folgende reelle Gleichung:\\
\begin{center}
	$ |x+2|-|x-2|=1 $ mit $ x \in \mathbb{R} $
\end{center}
Es muß eine Fallunterscheidung durchgeführt werden und zwar da, wo sich die Vorzeichen ändern.\\
Die Vorzeichen ändern sich bei -2 und +2.\\
Am Zahlenstrahl hat man dann 3 Fälle.\vspace*{2mm}
\\
\begin{tikzpicture}
	%\draw[->,>=stealth] (0,0)--++(-0.5,-0.5)--++(0.5,0)--++(-0.5,-0.5);
	
	\node at (2,0) [red,rotate=0] {$]$};
	\node at (1.8,-0.6) [black,rotate=0] {$-2$};
	\draw [red](0,0) -- (2,0) ;
	\draw [blue](2,0) -- (4,0) ;
	\node at (4,0) [green,rotate=0] {$[$};
	\draw [->,green](4,0) -- (6,0) ;
	\node at (4,-0.6) [black,rotate=0] {$2$};
\end{tikzpicture}\\
1.Fall: $ \textcolor{red}{x\leq2} $\\
2.Fall: $\textcolor{blue}{-2<x<2 } $\\
3.Fall: $ \textcolor{green}{x\geq 2} $\vspace*{2mm}\\
\underline{1.Fall:}\\
Für $x\leq-2  $ gilt:\\
\begin{align*}
	-(x+2)-(-(x-2))=-x-2+x-2=-4=1_{ \;\;\tikz \draw[->,>=stealth] (0,0)--++(-0.2,-0.2)--++(0.2,0)--++(-0.2,-0.2);}
\end{align*}
Für $ x\leq-2 $ gibt es keine Lösung!\vspace*{2mm}\\
\underline{2.Fall:}\\
Für $ -2<x<2 $ gilt:\\
\begin{align*}
	x+2-(-(x-2))=x+2+x-2=2x=1
\end{align*}
d.h. $ x=\frac{1}{2} $ ist eine Lösung.\vspace*{2mm}\\
\underline{3.Fall:}\\
$ 2\leq x $\\
\begin{align*}
	x+2-(x-2)=4=1_{ \;\;\tikz \draw[->,>=stealth] (0,0)--++(-0.2,-0.2)--++(0.2,0)--++(-0.2,-0.2);}
\end{align*}
Für $ 2\leq x $ gibt es keine Lösung!\\
\\
Fazit: Einzige Lösung ist $ x=\frac{1}{2} $.\\
\subsection{Aufgabe 1 zum Thema reelle Lösungen von Gleichungen}

Bestimme alle reellen Lösungen dieser Gleichung:\\
\begin{align*}
	\big||3x|-|-5|\big|=2
\end{align*}
Die Gleichung ist äquilalent zu:\\
\begin{align*}
	\big||3x|-|-5|\big|=2\Leftrightarrow\big||3x|-5\big|=2
\end{align*}
\newpage
\noindent
\underline{1.Fall:}\\
$ x\geq 0:\big||3x|-|-5|\big|=|3x-5| $ \\
\begin{quote}
	\underline{1.1Fall:}\\
	$x\geq \frac{5}{3}:|3x-5|=3x-5=2\Leftrightarrow 3x=7\Leftrightarrow x=\frac{7}{3}$\\
	\underline{1.2Fall:}\\
	$0<x<\frac{5}{3}:|3x-5|=-(3x-5)=-3x+5=2\Leftrightarrow -3x=-3\Leftrightarrow x=1$\\
\end{quote}
\underline{2.Fall:}\\
$x<0:\big||3x|-5\big|=|-3x-5|=|-1(3x+5)|=|(-1)||3x+5|=|3x+5|=2$\\
\begin{quote}
	\underline{2.1Fall:}\\
	$x<-\frac{5}{3}:|3x+5|=-(3x+5) =-3x-5=2\Leftrightarrow-3x=7\Leftrightarrow x=-\frac{7}{3} $\\
	\underline{2.2Fall:}\\
	$0\geq x\geq -\frac{5}{3} : |3x+5|=3x+5=2\Leftrightarrow 3x=-3\Leftrightarrow x=-1 $\\
\end{quote}
Fazit: $x_1=\frac{7}{3},\; x_2=1,\; x_3=-\frac{7}{3},\; x_4=-1$ lösen die Gleichung.\\

\subsection{Aufgabe 2 zum Thema reelle Lösungen von Gleichungen}

Bestimme alle reellen Lösungen der folgenden Ungleichung:\\
\begin{align*}
	x^3+11x<6x^2+6
\end{align*}
Zuerst Gleichung umstellen:\
\begin{align*}
	\Leftrightarrow \underbrace {x^3-6x+11x-6}_{=p(x)=(x-x_1)(x-x_2)(x-x_3)}<0
\end{align*}
NR:\\
Faktorisierung des Polynoms $ p(x) $\\
Durch hinschauen erhalten wir $ x_1=1 $ eine Nullstelle.\\
D.h. $ (x-x_1)=(x_1-1) $ teilt $ p(x) $\\
Mittels Polynomdivision ergibt sich:\\
\begin{center}
	\begin{tabular}[t]{lll}% l=linksbündig r=rechtsbündig
		&\;\;\;$(x^3-6x^2+11x-6)\div (x-1)=x^2-5x+6 $ \\
		&$-\underline{(x^3-\;\;x^2)} $  \\
		&\hspace{1.1cm}$ -5x^2+11x $\\
		&\hspace{0.6cm}$-\underline{(-5x^2+\;\;5x) }$\\
		&\hspace{2.8cm}$6x-6 $\\
		&\hspace{2.43cm}$\underline{-(6x-6)} $\\
		&\hspace{3.8cm}$0$\\
	\end{tabular} \\
\end{center}
\newpage
\noindent

$ p(x)=(x^2-5x+6)(x-1) $\\
Lösung der quatratischen Gleichung $ x^2-5x+6=0$.\\
\begin{align*}
	x^2-5x+6&=0\\
	x^2-5x+\frac{25}{4}-\frac{25}{4}+\frac{24}{4}&=0\\
	\bigg(x-\frac{5}{2}\bigg)^2-\frac{1}{4}&=0\\
	\bigg(x-\frac{5}{2}\bigg)^2&=\frac{1}{4}\\
	x-\frac{5}{2}&=\pm \sqrt{\frac{1}{4}}\\
	x_{2,3}&=\frac{5}{2}\pm \frac{1}{2}\\
	x_2&=3\\
	x_3&=2
\end{align*}
$ p(x)=(x-3)(x-2)(x-1) $\\
Jetzt kann man die Ungleichung durch die Faktorisierung leicht lösen:\\
\begin{align*}
	(x-3)(x-2)(x-1)<0
\end{align*}
Durch die Nullstellen $ x_1=1, x_2=3, x_3=2 $ bekommt man vier Fälle:\\
Am Zahlenstrahl hat man dann 4 Fälle.\vspace*{2mm}
\\
\begin{tikzpicture}
	%\draw[->,>=stealth] (0,0)--++(-0.5,-0.5)--++(0.5,0)--++(-0.5,-0.5);
	
	\node at (2,0) [red,rotate=0] {$][$};
	\node at (2,-0.6) [black,rotate=0] {$1$};
	\draw [red](0,0) -- (2,0) ;
	\draw [blue](2,0) -- (4,0) ;
	\node at (4,0) [green,rotate=0] {$][$};
	\draw [green](4,0) -- (6,0) ;
	\node at (4,-0.6) [black,rotate=0] {$2$};
	\draw [->,black](6,0) -- (8,0) ;
	\node at (6,-0.6) [black,rotate=0] {$3$};
	\node at (6,0) [black,rotate=0] {$][$};
\end{tikzpicture}\\
Bei $ x_1=1, x_2=3, x_3=2 $ wird das Polynom zu Null. Also keine Lösung, da Polynom $ <0 $ werden soll.\\
Für die 4 Fälle gilt dann:\\
\begin{center}
	\( \begin{array}{|l|c|c|c|c|}
		\hline       &(x-1)&(x-2)&(x-3)&p(x)\\
		\hline 1.Fall: x<1&-&-&-&-\\
		\hline 2.Fall: 1<x<2&+&-&-&+\\
		\hline 3.Fall: 2<x<3&+&+&-&-\\
		\hline 4.Fall: x>3&+&+&+&+\\
		\hline
	\end{array} \)
\end{center}
Fazit:\\
Bei $ x<1 $ und  $ 2<x<3 $ wird das Polynom negativ. Also ist die Lösung:\\
\begin{align*}
	x\in \;]-\infty,1[\;\; \cup \;\;]2,3[
\end{align*}









%\chapter{Komplexe Zahlen}
 \setcounter{section}{3}
\setcounter{subsection}{0}
\setcounter{subsubsection}{1}
\setcounter{secnumdepth}{3}
\subsection{Grundlagen der komplexen Zahlen }
\href{https://youtu.be/k6ZBEYDOrY0}{Video 03}\\
\begin{center}
	\fbox{\parbox  {65mm} {
	\begin{center}
		Definition: komplexe Zahlen:\\
	\end{center}
	\begin{center}
	$ \mathbb{C}:={x+y\mathrm{i}}\;|\;x,y\in \mathbb{R} \text{ mit } \mathrm{i}^2=-1$
	\end{center}
	
}}
\end{center}
Mit der Menge $  \mathbb{C} $ hat man die gleichen Grundrechenarten wie in der Menge $ \mathbb{R} $\vspace*{2mm}\\
Beispiel:\\
\textbullet $ \;\;(7+6\mathrm{i})+(1+1\mathrm{i})=(7+1)+(6\mathrm{i}+1\mathrm{i})=8+7\mathrm{i} \in  \mathbb{C}$\\
\textbullet $ \;\;(1+2\mathrm{i}) (2+1\mathrm{i})=2+1\mathrm{i}+4\mathrm{i}+2\mathrm{i}^2=2+5\mathrm{i}+2(-1)=0+5\mathrm{i}=5\mathrm{i} \in  \mathbb{C}$\\
\textbullet $ \;\;7+0\mathrm{i}=7 \in \mathbb{R}\subset  \mathbb{C} $\\
\textbullet $\;\; 0+6\mathrm{i} \in  \mathbb{C} $\\
\begin{center}
	\fbox{\parbox  {65mm} {
	\begin{center}
		Definition: Realteil\\
		$ \mathrm{Re}: \mathbb{C}\longmapsto \mathbb{R} \text{ mit } \mathrm{Re}(x+\mathrm{i}y)=x $
	\end{center}
    \begin{center}
    	Definition: Imaginärteil\\
    	$ \mathrm{Im}: \mathbb{C}\longmapsto \mathbb{R} \text{ mit } \mathrm{Im}(x+\mathrm{i}y)=y $
    	
    \end{center}
}}
\end{center}
Beispiel:\\
\textbullet$ \;\;\mathrm{Re}(7+6\mathrm{i})=7 $\\
\textbullet$ \;\; \mathrm{Im}(7+6\mathrm{i})=6 $\vspace*{4mm}\\
\newpage
\noindent
Darstellung in der Gaußschen Zahlenebene:\\
\begin{center}
\begin{tikzpicture}

% Gitternetzlinien
\draw[xstep=1cm, ystep=1cm, very thin, color=lightgray] (0,7) grid (8,0);
% x-Achse
\draw[->] (0,0) -- (8,0) node[below,xshift=-1em] {$\mathrm{Im}$}; 
%Zahlen auf x-Achse
\foreach \x in {1,...,7}
\draw[shift={(\x,0)},color=black] (0pt,2pt) -- (0pt,-2pt) node[below] {\footnotesize $\x$};
%\foreach \x in {1,...,3}
%\draw[shift={(\x,0)},color=black] (0pt,2pt) -- (0pt,-2pt) node[below] {\footnotesize $\x$};
% y-Achse 
\draw[->] (0,0) -- (0,7) node[left,yshift=-4em,xshift=-1em] {$\mathrm{Re}$};%node[above left]
%Zahlen auf y-Achse
%\foreach \y in {-1,...,-1}
%\draw[shift={(0,\y)},color=black] (2pt,0pt) -- (-2pt,0pt) node[left] {\footnotesize $\y$};
\foreach \y in {1,...,6}
\draw[shift={(0,\y)},color=black] (2pt,0pt) -- (-2pt,0pt) node[left] {\footnotesize $\y$};
%Ursprung
\draw[color=black] (0pt,-10pt) node[right] {\footnotesize $0$};
\draw[->] (0,0) -- (7,6) node[yshift=1em]{\footnotesize $z=7+6\mathrm{i}$};


%\draw[thick,color=black] plot[samples=200, domain=2:7] 
%(\x,{\x-1}) node[below, xshift=1em]{$g(x)$};



\end{tikzpicture}\\
\end{center}
Interpretation: Darstellung als Punkt $ (7,6) $ oder als Vektor $ \binom{7}{6} $\vspace*{4mm}\\
Addition von komplexen Zahlen ist gleich der Addition der Vektoren $ \binom{5}{2}+\binom{2}{4} $.\\
$ (5+2\mathrm{i})+(2+4\mathrm{i})=(7+6\mathrm{i}) $\\
\begin{center}
	\begin{tikzpicture}
	
	% Gitternetzlinien
	\draw[xstep=1cm, ystep=1cm, very thin, color=lightgray] (0,7) grid (8,0);
	% x-Achse
	\draw[->] (0,0) -- (8,0) node[below,xshift=-1em] {$Re$}; 
	%Zahlen auf x-Achse
	\foreach \x in {1,...,7}
	\draw[shift={(\x,0)},color=black] (0pt,2pt) -- (0pt,-2pt) node[below] {\footnotesize $\x$};
	%\foreach \x in {1,...,3}
	%\draw[shift={(\x,0)},color=black] (0pt,2pt) -- (0pt,-2pt) node[below] {\footnotesize $\x$};
	% y-Achse 
	\draw[->] (0,0) -- (0,7) node[left,yshift=-4em,xshift=-1em] {$\mathrm{Im}$};%node[above left]
	%Zahlen auf y-Achse
	%\foreach \y in {-1,...,-1}
	%\draw[shift={(0,\y)},color=black] (2pt,0pt) -- (-2pt,0pt) node[left] {\footnotesize $\y$};
	\foreach \y in {1,...,6}
	\draw[shift={(0,\y)},color=black] (2pt,0pt) -- (-2pt,0pt) node[left] {\footnotesize $\y$};
	%Ursprung
	\draw[color=black] (0pt,-10pt) node[right] {\footnotesize $0$};
	\draw[->] (0,0) -- (7,6) node[yshift=1em]{\footnotesize $7+6i$};
	\draw[->] (0,0) -- (5,2) node[yshift=-1em]{\footnotesize $5+2i$};
	\draw[->] (0,0) -- (2,4) node[yshift=1em]{\footnotesize $2+4i$};
	\draw[->,color=blue,dashed] (5,2) -- (7,6) ;
	%\draw[thick,color=black] plot[samples=200, domain=2:7] 
	%(\x,{\x-1}) node[below, xshift=1em]{$g(x)$};
	
	
	
	\end{tikzpicture}\\
\end{center}
\begin{center}
	\fbox{\parbox  {80mm} {
	\begin{center}
		Betrag einer komplexen Zahl\vspace*{2mm}\\
		$  \mathbb{C}\longmapsto [0,\infty[ \subset \mathbb{R} \text{ mit } |x+\mathrm{i}y|=\sqrt{x^2+y^2} $
	\end{center}
}}
\end{center}
Beispiel:\\
$ |7+6\mathrm{i}|=\sqrt{7^2+6^2}=\sqrt{85} $\\
\begin{center}
	\fbox{\parbox  {80mm} {
	\begin{center}
		Kunjugation einer komplexen Zahl \vspace*{2mm}\\
		$  \mathbb{C}\longmapsto  \mathbb{C}\text{ mit }\overline{(x+\mathrm{i}y)}=(x-\mathrm{i}y) $
	\end{center}
}}
\end{center}
Beispiel:\vspace*{2mm}\\
$\overline{(3+3\mathrm{i})}=(3-3\mathrm{i})$\\
\begin{center}
\begin{tikzpicture}

% Gitternetzlinien
\draw[xstep=1cm, ystep=1cm, very thin, color=lightgray] (-4,4) grid (4,-4);
% x-Achse
\draw[->] (-4,0) -- (4,0) node[below,xshift=-1em] {$\mathrm{Re}$}; 
%Zahlen auf x-Achse
\foreach \x in {-3,...,-1}
\draw[shift={(\x,0)},color=black] (0pt,2pt) -- (0pt,-2pt) node[below] {\footnotesize $\x$};
\foreach \x in {1,...,3}
\draw[shift={(\x,0)},color=black] (0pt,2pt) -- (0pt,-2pt) node[below] {\footnotesize $\x$};
% y-Achse 
\draw[->] (0,-4) -- (0,4) node[left,yshift=-1em] {$\mathrm{Im}$};%node[above left]
%Zahlen auf y-Achse
\foreach \y in {-3,...,-1}
\draw[shift={(0,\y)},color=black] (2pt,0pt) -- (-2pt,0pt) node[left] {\footnotesize $\y$};
\foreach \y in {1,...,3}
\draw[shift={(0,\y)},color=black] (2pt,0pt) -- (-2pt,0pt) node[left] {\footnotesize $\y$};
%Ursprung
\draw[color=black] (0pt,-10pt) node[right] {\footnotesize $0$};

\draw[thick,color=black,->] plot[samples=200, domain=0:3] 
(\x,{-1*\x}) node[below, xshift=0.8em] {$\overline{(3+3\mathrm{i})}$};

\draw[thick,color=blue,->] plot[samples=200, domain=0:3] 
(\x,{\x}) node[below, xshift=1em]{$3+3\mathrm{i}$};



\end{tikzpicture}\\
\end{center}
Die Kunjugation einer komplexen Zahl entspricht der Spiegelung des Vektors  an der Re-Achse.\\
\subsubsection{Aufgabe 1 zu Grundlagen der komplexen Zahlen }
\href{https://youtu.be/xC24jiGbXjo}{Video 05}\\
Skiziere die folgende Menge in die Gaußschene Zahlenebene:\\
\begin{align*}
\mathbb{A}= \{z \in  \mathbb{C}\;|\; 1\leq|z| \leq 4\}\cap\{z\in  \mathbb{C}\;|\; \mathrm{Re}_{(z)}\geq 1\}
\end{align*}
Zuerst für $ z=x+\mathrm{i}y $, $ |z|=\sqrt{x^2+y^2} $ und für $ \mathrm{Re}_{z}=x$ einsetzen.\\
\begin{align*}
\Leftrightarrow\mathbb{A}= \{x+\mathrm{i}y \in  \mathbb{C}\;|\; 1\leq \sqrt{x^2+y^2} \leq 4\}\cap\{x+\mathrm{i}y\in  \mathbb{C}\;|\; x\geq 1\}
\end{align*}
Zur Darstellung in die Zahlenebene benötigt man die Kreisgleichung.\\
\begin{center}
\fbox{\parbox  {60mm} {
		\begin{center}
				Kreisgleichung
		\end{center}
		\begin{align*}
		(x-x_0)^2+(y-y_0)^2=r^2\\
		\scriptscriptstyle \text{Radius }=r \text{ Mittelpunkt= }\scriptscriptstyle (x_0,y_0)
		\end{align*}
	%	$\scriptscriptstyle r \overset{\wedge}{=} \text{Radius } \text{ Mittelpunkt= }\scriptscriptstyle (x_0,y_0)  $
}}
\end{center}
Dann gilt:\\
\begin{align*}
\sqrt{(x-0)^2+(y-0)^2}=4\Leftrightarrow (x-0)^2+(y-0)^2=4^2\Rightarrow \text{ Kreis mit } r=4 \text{ vom Mittelpunkt } (0,0)\\
\end{align*}
\begin{center}
\begin{tikzpicture}[scale=.7]
\filldraw [blue,opacity=0.5] (0,0) circle (4);
\filldraw [white] (0,0) circle (1);
\filldraw [green,opacity=0.5] (1,-6) rectangle (6,6);
% Gitternetzlinien
\draw[xstep=1cm, ystep=1cm, very thin, color=lightgray] (-6,6) grid (6,-6);
% x-Achse
\draw[->] (-6,0) -- (6,0) node[below,xshift=-0.6em] {$\mathrm{Re}$}; 
%Zahlen auf x-Achse
\foreach \x in {-5,...,-1}
\draw[shift={(\x,0)},color=black] (0pt,2pt) -- (0pt,-2pt) node[below] {\footnotesize $\x$};
\foreach \x in {1,...,5}
\draw[shift={(\x,0)},color=black] (0pt,2pt) -- (0pt,-2pt) node[below] {\footnotesize $\x$};
% y-Achse 
\draw[->] (0,-6) -- (0,6) node[left,yshift=-0.6em] {$\mathrm{Im}$};%node[above left]
%Zahlen auf y-Achse
\foreach \y in {-5,...,-1}
\draw[shift={(0,\y)},color=black] (2pt,0pt) -- (-2pt,0pt) node[left] {\footnotesize $\y$};
\foreach \y in {1,...,5}
\draw[shift={(0,\y)},color=black] (2pt,0pt) -- (-2pt,0pt) node[left] {\footnotesize $\y$};
%Ursprung
\draw[color=black] (0pt,-10pt) node[right] {\footnotesize $0$};

%\draw[thick,color=black] plot[samples=200, domain=2:4] 
%(\x,{2}) node[below, xshift=-5em] {$h(x)$};

%\draw[thick,color=black] plot[samples=200, domain=-1:2] 
%(\x,{\x-1}) node[below, xshift=1em]{$g(x)$};

%\filldraw [blue,opacity=0.5] (0,0) circle (4);
%\filldraw [white,opacity=0.9] (0,0) circle (1);
%\filldraw [fill=green,draw=blue] (2,-2) rectangle (4,4);
%\filldraw[fill=cyan, draw=blue] (0,0) -- (0,0) arc (90:-90:10mm) -- (0,0);
%\filldraw[fill=cyan,draw=blue ](-1,1) -- (0,1)  arc[radius = 10mm, start angle= 0, end angle= 180];
%\draw[rotate=270,color=blue] (1,1)  arc[radius = 10mm, start angle= 0, end angle= 180];

\filldraw[blue,opacity=0.2] (1,3.88)--(1,-3.88)  arc[radius = 40mm, start angle= -75.5, end angle= 75.5];
\draw[color=black] (2,1) node[right] {$\mathbb{A}$};
\draw[color=blue] (-5.5,4) node[right] {\footnotesize $1\leq\sqrt{x^2+y^2}\leq4$};
\draw[color=black] (1.5,5) node[right] {\footnotesize $x\geq1$};
\end{tikzpicture}\\
\end{center}
\subsubsection{Aufgabe 2 zu Grundlagen der komplexen Zahlen}
\href{https://youtu.be/0TRhS_AU7g0}{Video 06}\\
Skiziere die folgende Menge in die Gaußschene Zahlenebene:\\
\begin{align*}
\mathbb{A}=\{z \in \mathbb{Z}\;|\;|\mathrm{Im}(z)|\leq 2\cdot|\mathrm{Re}(z)|\}\cap\{z\in \mathbb{C}\;|\;|z|\leq 1 \}
\end{align*}
Zuerst für $ z=x+\mathrm{i}y $, $ \mathrm{Im}(z)=y $ und $ \mathrm{Re}(z)=x$ und $ |z|=\sqrt{x^2+y^2} $ ersetzen.\\
\begin{align*}
\mathbb{A}=\{x+\mathrm{i}y \in \mathbb{Z}\;|\;|y|\leq 2\cdot|x|\}\cap\{x+\mathrm{i}y\in \mathbb{C}\;|\;\sqrt{x^2+y^2}\leq 1 \}
\end{align*}
Vorgehensweise:\\
\textbullet \;\;Zeichne die Menge $|y|\leq 2\cdot|x| $im I.Quatranten:\\
\begin{quote}
Es gilt: $ y\leq2x  $\\	
Zeichne zuerst die Menge: $ y=2x $ \\
Wähle einen Punkt für die, die Gleichung $ y\leq2x $ gilt:\\Für den  Punkt$P=(1,1)  $ erhält man $1\leq2\cdot1  $, deshalb ist die Menge unterhalb der Geraden.\\

\begin{center}
\begin{tikzpicture}
\filldraw [red,opacity=0.5] (0,0)--(1,2)--(2,4)--(4,4)--(4,0)--(0,0);
% Gitternetzlinien
\draw[xstep=1cm, ystep=1cm, very thin, color=lightgray] (-4,4) grid (4,-2);
% x-Achse
\draw[->] (-4,0) -- (4,0) node[below,xshift=-1em] {$\mathrm{Re}$}; 
%Zahlen auf x-Achse
\foreach \x in {-3,...,-1}
\draw[shift={(\x,0)},color=black] (0pt,2pt) -- (0pt,-2pt) node[below] {\footnotesize $\x$};
\foreach \x in {1,...,3}
\draw[shift={(\x,0)},color=black] (0pt,2pt) -- (0pt,-2pt) node[below] {\footnotesize $\x$};
% y-Achse 
\draw[->] (0,-2) -- (0,4) node[left,yshift=-1em] {$\mathrm{Im}$};%node[above left]
%Zahlen auf y-Achse
\foreach \y in {-1,...,-1}
\draw[shift={(0,\y)},color=black] (2pt,0pt) -- (-2pt,0pt) node[left] {\footnotesize $\y$};
\foreach \y in {1,...,3}
\draw[shift={(0,\y)},color=black] (2pt,0pt) -- (-2pt,0pt) node[left] {\footnotesize $\y$};
%Ursprung
\draw[color=black] (0pt,-10pt) node[right] {\footnotesize $0$};

\draw[thick,color=black] plot[samples=200, domain=0:2] 
(\x,{2*\x}) node[below, xshift=1.5em] {$y=2x$};





\fill [black] (1,1) circle (.051cm) ;
\draw (1,1) node[xshift=-0.6em]{P};





\end{tikzpicture}\\
\end{center}
\end{quote}
\textbullet \;\;Zeichne die Menge im II.Quatranten:\\

\begin{quote}
Es gilt: $ y\leq-2x  $\\	
Zeichne zuerst die Menge: $ y=-2x $ \\
Wähle einen Punkt für die, die Gleichung $ y\leq-2x $ gilt:\\Für den  Punkt $P=(-1,1)  $ erhält man $-1\leq2\cdot1  $, deshalb ist die Menge unterhalb der Geraden.\\
\begin{center}
	\begin{tikzpicture}
	\filldraw [red,opacity=0.5] (0,0)--(-1,2)--(-2,4)--(-4,4)--(-4,0)--(0,0);
	% Gitternetzlinien
	\draw[xstep=1cm, ystep=1cm, very thin, color=lightgray] (-4,4) grid (4,-2);
	% x-Achse
	\draw[->] (-4,0) -- (4,0) node[below,xshift=-1em] {$\mathrm{Re}$}; 
	%Zahlen auf x-Achse
	\foreach \x in {-3,...,-1}
	\draw[shift={(\x,0)},color=black] (0pt,2pt) -- (0pt,-2pt) node[below] {\footnotesize $\x$};
	\foreach \x in {1,...,3}
	\draw[shift={(\x,0)},color=black] (0pt,2pt) -- (0pt,-2pt) node[below] {\footnotesize $\x$};
	% y-Achse 
	\draw[->] (0,-2) -- (0,4) node[left,yshift=-1em] {$\mathrm{Im}$};%node[above left]
	%Zahlen auf y-Achse
	\foreach \y in {-1,...,-1}
	\draw[shift={(0,\y)},color=black] (2pt,0pt) -- (-2pt,0pt) node[left] {\footnotesize $\y$};
	\foreach \y in {1,...,3}
	\draw[shift={(0,\y)},color=black] (2pt,0pt) -- (-2pt,0pt) node[left] {\footnotesize $\y$};
	%Ursprung
	\draw[color=black] (0pt,-10pt) node[right] {\footnotesize $0$};
	
	\draw[thick,color=black] plot[samples=200, domain=-2:0] 
	(\x,{-2*\x}) node[below, xshift=-6.5em, yshift=9em] {$y=-2x$};
	
	
	
	
	
	\fill [black] (-1,1) circle (.051cm) ;
	\draw (-1,1) node[xshift=-0.6em]{P};
	
	
	
	
	
	\end{tikzpicture}\\
\end{center}
\end{quote}
\textbullet \;\;Zeichne die Menge im III.Quatranten:\\
\begin{quote}
Es gilt: $ -y\leq-2x  $\\	
Zeichne zuerst die Menge: $ -y=-2x $ \\
Wähle einen Punkt für die, die Gleichung $ -y\leq-2x $ gilt:\\Für den  Punkt $P=(-1,-1)  $ erhält man $1\leq-2\cdot-1  $, deshalb ist die Menge oberhalb der Geraden.\\
\begin{center}
	\begin{tikzpicture}
	\filldraw [red,opacity=0.5] (0,0)--(-4,0)--(-4,-4)--(-2,-4)--(0,0);
	% Gitternetzlinien
	\draw[xstep=1cm, ystep=1cm, very thin, color=lightgray] (-4,-4) grid (4,4);
	% x-Achse
	\draw[->] (-4,0) -- (4,0) node[below,xshift=-1em] {$Re$}; 
	%Zahlen auf x-Achse
	\foreach \x in {-3,...,-1}
	\draw[shift={(\x,0)},color=black] (0pt,2pt) -- (0pt,-2pt) node[below] {\footnotesize $\x$};
	\foreach \x in {1,...,3}
	\draw[shift={(\x,0)},color=black] (0pt,2pt) -- (0pt,-2pt) node[below] {\footnotesize $\x$};
	% y-Achse 
	\draw[->] (0,-4) -- (0,4) node[left,yshift=-1em] {$Im$};%node[above left]
	%Zahlen auf y-Achse
	\foreach \y in {-1,...,-3}
	\draw[shift={(0,\y)},color=black] (2pt,0pt) -- (-2pt,0pt) node[left] {\footnotesize $\y$};
	\foreach \y in {1,...,3}
	\draw[shift={(0,\y)},color=black] (2pt,0pt) -- (-2pt,0pt) node[left] {\footnotesize $\y$};
	%Ursprung
	\draw[color=black] (0pt,-10pt) node[right] {\footnotesize $0$};
	
	\draw[thick,color=black] plot[samples=200, domain=-2:0] 
	(\x,{2*\x}) node[below, xshift=-7em, yshift=-8em] {$-y=-2x$};
	
	
	
	
	
	\fill [black] (-1,-1) circle (.051cm) ;
	\draw (-1,-1) node[xshift=-0.6em]{P};
	
	
	
	
	
	\end{tikzpicture}\\
\end{center}
\end{quote}
\newpage
\textbullet \;\;Zeichne die Menge im IV.Quatranten:\\

\begin{quote}
	Es gilt: $ -y\leq 2x  $\\	
	Zeichne zuerst die Menge: $ -y= 2x $ \\
	Wähle einen Punkt für die, die Gleichung $ -y\leq 2x $ gilt:\\Für den  Punkt $P=(1,-1)  $ erhält man $1\leq2\cdot1  $, deshalb ist die Menge oberhalb der Geraden.\\
	\begin{center}
		\begin{tikzpicture}
		\filldraw [red,opacity=0.5] (0,0)--(4,0)--(4,-4)--(2,-4)--(0,0);
		% Gitternetzlinien
		\draw[xstep=1cm, ystep=1cm, very thin, color=lightgray] (-4,4) grid (4,-4);
		% x-Achse
		\draw[->] (-4,0) -- (4,0) node[below,xshift=-1em] {$\mathrm{Re}$}; 
		%Zahlen auf x-Achse
		\foreach \x in {-3,...,-1}
		\draw[shift={(\x,0)},color=black] (0pt,2pt) -- (0pt,-2pt) node[below] {\footnotesize $\x$};
		\foreach \x in {1,...,3}
		\draw[shift={(\x,0)},color=black] (0pt,2pt) -- (0pt,-2pt) node[below] {\footnotesize $\x$};
		% y-Achse 
		\draw[->] (0,-4) -- (0,4) node[left,yshift=-1em] {$\mathrm{Im}$};%node[above left]
		%Zahlen auf y-Achse
		\foreach \y in {-1,...,-3}
		\draw[shift={(0,\y)},color=black] (2pt,0pt) -- (-2pt,0pt) node[left] {\footnotesize $\y$};
		\foreach \y in {1,...,3}
		\draw[shift={(0,\y)},color=black] (2pt,0pt) -- (-2pt,0pt) node[left] {\footnotesize $\y$};
		%Ursprung
		\draw[color=black] (0pt,-10pt) node[right] {\footnotesize $0$};
		
		\draw[thick,color=black] plot[samples=200, domain=0:2] 
		(\x,{-2*\x}) node[below, xshift=2.5em, yshift=2.5em] {$-y=2x$};
		
		
		
		
		
		\fill [black] (1,-1) circle (.051cm) ;
		\draw (1,-1) node[xshift=-0.6em]{P};
		
		
		
		
		
		\end{tikzpicture}\\
	\end{center}
\end{quote}
Die gesamte Menge $|y|\leq 2\cdot|x|$ ist dann:\\
\begin{center}
	\begin{tikzpicture}
	\filldraw [red,opacity=0.5] (0,0)--(4,0)--(4,-4)--(2,-4)--(0,0);
	\filldraw [red,opacity=0.5] (0,0)--(-1,2)--(-2,4)--(-4,4)--(-4,0)--(0,0);
	\filldraw [red,opacity=0.5] (0,0)--(1,2)--(2,4)--(4,4)--(4,0)--(0,0);
	\filldraw [red,opacity=0.5] (0,0)--(-4,0)--(-4,-4)--(-2,-4)--(0,0);
	% Gitternetzlinien
	\draw[xstep=1cm, ystep=1cm, very thin, color=lightgray] (-4,-4) grid (4,4);
	% x-Achse
	\draw[->] (-4,0) -- (4,0) node[below,xshift=-1em] {$\mathrm{Re}$}; 
	%Zahlen auf x-Achse
	\foreach \x in {-3,...,-1}
	\draw[shift={(\x,0)},color=black] (0pt,2pt) -- (0pt,-2pt) node[below] {\footnotesize $\x$};
	\foreach \x in {1,...,3}
	\draw[shift={(\x,0)},color=black] (0pt,2pt) -- (0pt,-2pt) node[below] {\footnotesize $\x$};
	% y-Achse 
	\draw[->] (0,-4) -- (0,4) node[left,yshift=-1em] {$\mathrm{Im}$};%node[above left]
	%Zahlen auf y-Achse
	\foreach \y in {-1,...,-3}
	\draw[shift={(0,\y)},color=black] (2pt,0pt) -- (-2pt,0pt) node[left] {\footnotesize $\y$};
	\foreach \y in {1,...,3}
	\draw[shift={(0,\y)},color=black] (2pt,0pt) -- (-2pt,0pt) node[left] {\footnotesize $\y$};
	%Ursprung
	\draw[color=black] (0pt,-10pt) node[right] {\footnotesize $0$};
	
	\draw[thick,color=black] plot[samples=200, domain=-2:0] 
	(\x,{2*\x}) node[below, xshift=-7em, yshift=-8em] {$-y=-2x$};
	
	\draw[thick,color=black] plot[samples=200, domain=-2:0] 
	(\x,{-2*\x}) node[below, xshift=-6.5em, yshift=9em] {$y=-2x$};
	
	\draw[thick,color=black] plot[samples=200, domain=0:2] 
	(\x,{2*\x}) node[below, xshift=1.5em] {$y=2x$};
	
	\draw[thick,color=black] plot[samples=200, domain=0:2] 
	(\x,{-2*\x}) node[below, xshift=2.5em, yshift=2.5em] {$-y=2x$};
	\fill [black] (-1,-1) circle (.051cm) ;
	\draw (-1,-1) node[xshift=-0.6em]{P};
	
	\fill [black] (-1,1) circle (.051cm) ;
	\draw (-1,1) node[xshift=-0.6em]{P};
	
	\fill [black] (1,1) circle (.051cm) ;
	\draw (1,1) node[xshift=-0.6em]{P};
		\fill [black] (1,-1) circle (.051cm) ;
	\draw (1,-1) node[xshift=-0.6em]{P};
	\end{tikzpicture}\\
\end{center}
\textbullet\;\;Zeichne die Menge $ \sqrt{x^2+y^2}\leq 1 $\\
\begin{quote}
	\begin{align*}
	\sqrt{x^2+y^2}\leq1\Leftrightarrow \sqrt{(x-x_0)^2+(y-y_0)^2}\leq 1\Leftrightarrow (x-x_0)^2+(y-y_0)^2=1^1
	\end{align*}
\end{quote}
Zeichne nun $ (x-x_0)^2+(y-y_0)^2=1 $ , die gesamte Menge ist dann von $ r=1 $ bis $ r=0 $.\\


\begin{center}
	\begin{tikzpicture}
	\filldraw [blue,opacity=0.5] (0,0) circle (1);
	\filldraw [red,opacity=0.5] (0,0)--(4,0)--(4,-4)--(2,-4)--(0,0);
	\filldraw [red,opacity=0.5] (0,0)--(-1,2)--(-2,4)--(-4,4)--(-4,0)--(0,0);
	\filldraw [red,opacity=0.5] (0,0)--(1,2)--(2,4)--(4,4)--(4,0)--(0,0);
	\filldraw [red,opacity=0.5] (0,0)--(-4,0)--(-4,-4)--(-2,-4)--(0,0);
	% Gitternetzlinien
	\draw[xstep=1cm, ystep=1cm, very thin, color=lightgray] (-4,-4) grid (4,4);
	% x-Achse
	\draw[->] (-4,0) -- (4,0) node[below,xshift=-1em] {$\mathrm{Re}$}; 
	%Zahlen auf x-Achse
	\foreach \x in {-3,...,-1}
	\draw[shift={(\x,0)},color=black] (0pt,2pt) -- (0pt,-2pt) node[below] {\footnotesize $\x$};
	\foreach \x in {1,...,3}
	\draw[shift={(\x,0)},color=black] (0pt,2pt) -- (0pt,-2pt) node[below] {\footnotesize $\x$};
	% y-Achse 
	\draw[->] (0,-4) -- (0,4) node[left,yshift=-1em] {$\mathrm{Im}$};%node[above left]
	%Zahlen auf y-Achse
	\foreach \y in {-1,...,-3}
	\draw[shift={(0,\y)},color=black] (2pt,0pt) -- (-2pt,0pt) node[left] {\footnotesize $\y$};
	\foreach \y in {1,...,3}
	\draw[shift={(0,\y)},color=black] (2pt,0pt) -- (-2pt,0pt) node[left] {\footnotesize $\y$};
	%Ursprung
	\draw[color=black] (0pt,-10pt) node[right] {\footnotesize $0$};
	
	\draw[thick,color=black] plot[samples=200, domain=-2:0] 
	(\x,{2*\x}) node[below, xshift=-7em, yshift=-8em] {$-y=-2x$};
	
	\draw[thick,color=black] plot[samples=200, domain=-2:0] 
	(\x,{-2*\x}) node[below, xshift=-6.5em, yshift=9em] {$y=-2x$};
	
	\draw[thick,color=black] plot[samples=200, domain=0:2] 
	(\x,{2*\x}) node[below, xshift=1.5em] {$y=2x$};
	
	\draw[thick,color=black] plot[samples=200, domain=0:2] 
	(\x,{-2*\x}) node[below, xshift=2.5em, yshift=2.5em] {$-y=2x$};
	
	
	
	
	
	
	\draw (-0.5,0.3) node[xshift=-0.5em]{$\mathbb{A}$};
	\end{tikzpicture}\\
\end{center}
\subsubsection{Aufgabe 1 zur Darstellung komplexer Zahlen in kartesischen Koordinaten}
\href{https://youtu.be/FL3zN9eSh_c}{Video 07}\\
Stelle folgende komplexe Zahl in der Form $ x+\mathrm{i}y\; |\; x,y \in \mathbb{R}$ dar:\\
\begin{align*}
\frac{8-\mathrm{i}}{5-1}
\end{align*}
Durch Erweiterung mit $ (5+1) $ und Ausnützung der 3. binomischen Formel erhält man:\\
\begin{align*}
\frac{8-\mathrm{i}}{5-\mathrm{i}}=\frac{(8-\mathrm{i})(5+\mathrm{i})}{(5-\mathrm{i})(5+\mathrm{i})}=\frac{40+8\mathrm{i}-5\mathrm{i}-\mathrm{i}^2}{5^2-\mathrm{i}^2}=\frac{40+3\mathrm{i}-(-1)}{25-(-1)}=\frac{41+3\mathrm{i}}{26}=\frac{41}{26}+\mathrm{i}\frac{3}{26} \text{ mit }x \in \mathbb{R}
\end{align*}
\subsubsection{Aufgabe 2 zur Darstellung komplexer Zahlen in kartesischen Koordinaten}
\href{https://youtu.be/fm9W_BQeuOA}{Video 08}\\
Stelle folgende komplexe Zahl in der Form $ x+\mathrm{i}y $ mit $ x,y \in \mathbb{R} $ dar.\\
\begin{align*}
\sum\limits_{n=1}^{493836} \mathrm{i}^n\;\; \;\; n\in\mathbb{N}
\end{align*}
Für die Summe gilt doch:
\begin{align*}
\sum\limits_{n=1}^{493836} \mathrm{i}^n= \mathrm{i}^1+\mathrm{i}^2+\mathrm{i}^3\ldots 493836\\
\end{align*}
\newpage
\noindent
Es gilt:\\
\begin{align*}
&\mathrm{i}^1=\mathrm{i}\\
&\mathrm{i}^2=-\mathrm{i}\\
&\mathrm{i}^4=\mathrm{i}^2\cdot \mathrm{i}^2=(-1)(-1)=1\\
&\sum\limits_{n=1}^{493836} \mathrm{i}^n= \mathrm{i}^1+\mathrm{i}^2+\mathrm{i}^3+\mathrm{i}^4+\underbrace {\mathrm{i}^5}_{=\mathrm{i}^4\cdot \mathrm{i}^1=\mathrm{i}^1}+\underbrace {\mathrm{i}^6}_{=\mathrm{i}^4\cdot \mathrm{i}^2=\mathrm{i}^2}+\underbrace {\mathrm{i}^7}_{=\mathrm{i}^4+\mathrm{i}^3=\mathrm{i}^3}+\underbrace {\mathrm{i}^8}_{\mathrm{i}^4\cdot \mathrm{i}^4=1=\mathrm{i}^4}+\underbrace {\mathrm{i}^9}_{=\mathrm{i}^4\cdot \mathrm{i}^4\cdot \mathrm{i}=\mathrm{i}^1}\ldots\\
\end{align*}
Die ersten Summen 1-4 sind also gleich mit den Summen 5-8 und den Summen 9-12 usw..\\
Es gilt dann:\\
\begin{align*}
& \sum\limits_{n=1}^{493836} \mathrm{i}^n=\sum\limits_{n=1}^{\frac{493836}{4}}(\mathrm{i}+\mathrm{i}^2+\mathrm{i}^3+\mathrm{i}^4)\\
&\hspace{1.4cm}=\frac{493836}{4}(\mathrm{i}+\mathrm{i}^2+\mathrm{i}^3+\mathrm{i}^4)=\frac{493836}{4}(\mathrm{i}+(-1)+(-1\cdot \mathrm{i})+1)=\frac{493836}{4}\cdot 0\\
&\hspace{1.4cm}=0+0\mathrm{i}\\
\end{align*}
\subsection{Polarkoordinaten komplexer Zahlen}
\subsubsection{Polarkoordinaten komplexer Zahlen - Herleitung}
\href{https://youtu.be/8uVqLVXBXL4}{Video 09}\\
Für $ z=x+\mathrm{i}y $ und $ |z|=1 $ gilt:\\
\begin{align*}
&z=\cos(\alpha)+i\cdot \sin(\alpha) \text{ für }|z|=1\\
\end{align*}
\begin{center}

\begin{tikzpicture}[scale=2.5]

% Gitternetzlinien
\draw[xstep=1cm, ystep=1cm, very thin, color=lightgray] (-2,-2) grid (2,2);
% x-Achse
\draw[->] (-2,0) -- (2,0) node[below,xshift=-1em] {$\mathrm{Re}$}; 
%Zahlen auf x-Achse
\foreach \x in {-1,...,-1}
\draw[shift={(\x,0)},color=black] (0pt,2pt) -- (0pt,-2pt) node[below] {\footnotesize $\x$};
\foreach \x in {1,...,1}
\draw[shift={(\x,0)},color=black] (0pt,2pt) -- (0pt,-2pt) node[below] {\footnotesize $\x$};
% y-Achse 
\draw[->] (0,-2) -- (0,2) node[left,yshift=-1em] {$\mathrm{Im}$};%node[above left]
%Zahlen auf y-Achse
\foreach \y in {-1,...,-1}
\draw[shift={(0,\y)},color=black] (2pt,0pt) -- (-2pt,0pt) node[left] {\footnotesize $\y$};
\foreach \y in {1,...,1}
\draw[shift={(0,\y)},color=black] (2pt,0pt) -- (-2pt,0pt) node[left] {\footnotesize $\y$};
%Ursprung
%\draw[color=black] (0pt,-10pt) node[right] {\footnotesize $0$};
	\draw[color=black] (-4pt,-4pt) node[right] {\footnotesize $0$};
%\draw[thick,color=black] plot[samples=200, domain=2:4] 
%(\x,{2}) node[below, xshift=-5em] {$h(x)$};

%\draw[thick,color=black] plot[samples=200, domain=-1:2] 
%(\x,{\x-1}) node[below, xshift=1em]{$g(x)$};

\draw [blue] (0,0) circle (1);
\filldraw[fill=green!20,draw=green!50!black] (0,0) -- (5mm,0mm) arc (0:45:5mm) -- cycle; 
\draw [blue,->] (0,0)-.(0.7,0.7)node[below,xshift= 0.5em,yshift=1em] {\footnotesize $z$};
\draw [blue,dashed] (0,.7)-.(0.7,0.7);
\draw [blue,dashed] (0.7,0)-.(0.7,0.7);
\draw [blue] (0.7,0) node[below,xshift= 0em,yshift=-0.2em] {\footnotesize $x$};
\draw [blue] (0.3,0.3) node[below,xshift= -0.5em,yshift=1.8em,rotate=45] {\footnotesize $|z|$};
\draw [blue] (0,0.7) node[below,xshift=- 0.51em,yshift=1em] {\footnotesize $y$};


\draw [red,thick] (0,0)-.(0.7,0) node[midway,xshift= 0em,yshift=-0.5em] {\footnotesize $\cos(\alpha)$};
\draw [red,thick] (0,0)-.(0,0.7) node[midway,xshift= -0.50em,yshift=-0em,rotate=90] {\footnotesize $\sin(\alpha)$};
\draw [red] (0.2,0.2) node[xshift= 0em,yshift=-0.6em,rotate=0] {\footnotesize $\alpha$};
\end{tikzpicture}\\
\end{center}
\newpage
\noindent
Für $ z \in \mathbb{C} $ beliebig gilt:\\
\begin{align*}
z=x+\mathrm{i}y&=|z|\cos(\alpha)+|z|\;i\;\sin(\alpha)\\
&=|z|(\cos(\alpha)+\mathrm{i} \; \sin(\alpha))
\end{align*}
\begin{center}

\fbox{\parbox  {70mm} {
		\begin{align*}
		z=|z|(\cos(\alpha)+\mathrm{i}\; \sin(\alpha))
		\text{  für } \alpha \in [0,2\pi[
		\end{align*}
}}
\end{center}
\begin{center}

\begin{tikzpicture}[scale=2.5]

% Gitternetzlinien
\draw[xstep=1cm, ystep=1cm, very thin, color=lightgray] (-2,-2) grid (2,2);
% x-Achse
\draw[->] (-2,0) -- (2,0) node[below,xshift=-2em] {$\mathrm{Re}$}; 
%Zahlen auf x-Achse
\foreach \x in {-1,...,-1}
\draw[shift={(\x,0)},color=black] (0pt,2pt) -- (0pt,-2pt) node[below] {\footnotesize $\x$};
\foreach \x in {1,...,1}
\draw[shift={(\x,0)},color=black] (0pt,2pt) -- (0pt,-2pt) node[below] {\footnotesize $\x$};
% y-Achse 
\draw[->] (0,-2) -- (0,2) node[left,yshift=-0.6em] {$\mathrm{Im}$};%node[above left]
%Zahlen auf y-Achse
\foreach \y in {-1,...,-1}
\draw[shift={(0,\y)},color=black] (2pt,0pt) -- (-2pt,0pt) node[left] {\footnotesize $\y$};
\foreach \y in {1,...,1}
\draw[shift={(0,\y)},color=black] (2pt,0pt) -- (-2pt,0pt) node[left] {\footnotesize $\y$};
%Ursprung
%\draw[color=black] (0pt,-10pt) node[right] {\footnotesize $0$};
	\draw[color=black] (-4pt,-4pt) node[right] {\footnotesize $0$};
%\draw[thick,color=black] plot[samples=200, domain=2:4] 
%(\x,{2}) node[below, xshift=-5em] {$h(x)$};

%\draw[thick,color=black] plot[samples=200, domain=-1:2] 
%(\x,{\x-1}) node[below, xshift=1em]{$g(x)$};

\draw [blue] (0,0) circle (1.8);
\filldraw[fill=green!20,draw=green!50!black] (0,0) -- (5mm,0mm) arc (0:45:5mm) -- cycle; 
\draw [blue,->] (0,0)-.(1.27,1.27)node[below,xshift= 0.5em,yshift=1em] {\footnotesize $z$};
\draw [blue,dashed] (0,1.27)-.(1.27,1.27);
\draw [blue,dashed] (1.27,0)-.(1.27,1.27);
\draw [blue] (1.27,0) node[below,xshift= 0em,yshift=-0.2em] {\footnotesize $x$};
\draw [blue] (0.5,0.5) node[below,xshift= -0.5em,yshift=1.8em,rotate=45] {\footnotesize $|z|$};
\draw [blue] (0,1.27) node[below,xshift=- 0.51em,yshift=1em] {\footnotesize $y$};


\draw [red,thick] (0,0)-.(1.27,0) node[midway,xshift= 0em,yshift=-0.5em] {\footnotesize $|z|\cos(\alpha)$};
\draw [red,thick] (0,0)-.(0,1.27) node[midway,xshift= -0.50em,yshift=0em,rotate=90] {\footnotesize $|z|\sin(\alpha)$};
\draw [red] (0.2,0.2) node[xshift= 0em,yshift=-0.6em,rotate=0] {\footnotesize $\alpha$};
\end{tikzpicture}\\
\end{center}
Beispiel:\\
\textbullet \;\;$ z_1=1=1+0\mathrm{i}=|1|(1+0\mathrm{i}) =1(\cos(0) +\mathrm{i}\;\sin(0))$\\
\textbullet \;\;$ z_2=\mathrm{i}=0+1\mathrm{i}\;=|1|(0+1\;\mathrm{i})=1(\cos(\frac{\pi}{2})+\mathrm{i}\;\sin(\frac{\pi}{2})) $\\
\textbullet \;\;$ z_3=1,5+1,6\;\mathrm{i}=|\sqrt{1,5^2+1,6^2}|(\underbrace {\scriptstyle\frac{1,5}{2,19}}_{=\cos(\alpha)}+\mathrm{i}\underbrace {\scriptscriptstyle\frac{1,6}{2,19}}_{=\sin(\alpha)})=2,19(\cos(0,684)+\mathrm{i}\;\sin(0,730))$\\
\begin{center}
	
	\begin{tikzpicture}[scale=2.5]
	
	% Gitternetzlinien
	\draw[xstep=1cm, ystep=1cm, very thin, color=lightgray] (-2,-2) grid (2,2);
	% x-Achse
	\draw[->] (-2,0) -- (2,0) node[below,xshift=-1em] {$\mathrm{Re}$}; 
	%Zahlen auf x-Achse
	\foreach \x in {-1,...,-1}
	\draw[shift={(\x,0)},color=black] (0pt,2pt) -- (0pt,-2pt) node[below] {\footnotesize $\x$};
	\foreach \x in {1,...,1}
	\draw[shift={(\x,0)},color=black] (0pt,2pt) -- (0pt,-2pt) node[below] {\footnotesize $\x$};
	% y-Achse 
	\draw[->] (0,-2) -- (0,2) node[left,yshift=-1em] {$\mathrm{Im}$};%node[above left]
	%Zahlen auf y-Achse
	\foreach \y in {-1,...,-1}
	\draw[shift={(0,\y)},color=black] (2pt,0pt) -- (-2pt,0pt) node[left] {\footnotesize $\y$};
	\foreach \y in {1,...,1}
	\draw[shift={(0,\y)},color=black] (2pt,0pt) -- (-2pt,0pt) node[left] {\footnotesize $\y$};
	%Ursprung
	%\draw[color=black] (0pt,-10pt) node[right] {\footnotesize $0$};
		\draw[color=black] (-4pt,-4pt) node[right] {\footnotesize $0$};
	%\draw[thick,color=black] plot[samples=200, domain=2:4] 
	%(\x,{2}) node[below, xshift=-5em] {$h(x)$};
	
	%\draw[thick,color=black] plot[samples=200, domain=-1:2] 
	%(\x,{\x-1}) node[below, xshift=1em]{$g(x)$};
	
	
	\draw [green,->,thick] (0,0)-.(0,1)node[below,xshift= 0.5em,yshift=1em] {\footnotesize $z_2$};
	\draw [blue,->,thick] (0,0)-.(1,0)node[below,xshift= 0.5em,yshift=1em] {\footnotesize $z_1$};
	
	\draw [red,->,thick] (0,0)-.(1.5,1.6)node[below,xshift= 0.5em,yshift=1em] {\footnotesize $z_3$};
	\draw [red,dashed] (1.5,0)- -(1.5,1.6);
	\draw [red,dashed] (0,1.6)- -(1.5,1.6);
	\draw  [red,dashed] (0,1.6)node[below,xshift= -1.5em,yshift=0.6em] {\footnotesize $1,6$};
	\draw  [red,dashed] (1.5,0)node[below,xshift= 0em,yshift=-0.3em] {\footnotesize $1,5$};
	\draw  [red,dashed] (0.9,0.9)node[rotate=47,xshift=-1em,yshift=+1em] {\footnotesize $|2,19|$};
	\end{tikzpicture}\\
\end{center}
\subsubsection{Polarkoordinaten komplexe Exponentialfunktion}
\href{https://youtu.be/5P0a9p3mMw8}{Video 10}\\
\textbf{Komplexe  Exponentialfunktion}\\
Es gilt für $ x,y\in \mathbb{R}  $ und der Eulerischen Formel: $ \mathrm{e}^{\mathrm{i}y}=\cos(y)+i\cdot \sin(y) $: \\
\begin{align*}
&\mathrm{e}^{x+\mathrm{i}y}=\mathrm{e}^x(\textcolor{red}{\mathrm{e}^{\mathrm{i}y})}=\mathrm{e}^x\textcolor{red}{(\cos(y)+\mathrm{i} \cdot \sin(y))}\\
&\Rightarrow \mathrm{e}^{\mathrm{i}y}=\cos(y)+\mathrm{i}\cdot \sin(y)\\
&\Rightarrow \mathrm{e}^x \in \mathbb{R} \text{ daher ein Faktor}\\
\end{align*}
Damit können wir $ z\in \mathbb{C} \backslash \{0\} $ darstellen als:\\
\begin{align*}
z=|z|(\cos(\alpha)+\mathrm{i}\;\sin(\alpha))=|z|\mathrm{e}^{\mathrm{i}\alpha}=\mathrm{e}^x\mathrm{e}^{\mathrm{i}\alpha}\;\;\;\; \text{   für } \alpha \in [0,2\pi[
\end{align*}
\begin{center}



\boxed{z=|z|\mathrm{e}^{\mathrm{i}\alpha}}


\end{center}
Bemerkung: $\mathrm{e}^x\mathrm{e}^{\mathrm{i}x}$ wird nicht Null, deshalb ist $ z\neq 0 $.

\subsection{Multiplikation komplexer Zahlen in Polarkoordinaten}
\href{https://youtu.be/4eKWaZxQxvY}{Video 11}\\
Seien $ z_1,z_2 \in  \mathbb{C} \text{ mit } z_1=|z_1|\mathrm{e}^{\mathrm{i}\alpha_1} \text{ und }z_2=|z_2|\mathrm{e}^{\mathrm{i}\alpha_2} $ dann gilt:\\
\begin{center}
	\begin{align*}
	\boxed{z_1\cdot z_2=|z_1|\mathrm{e}^{\mathrm{i}\alpha_1}\cdot |z_2|\mathrm{e}^{\mathrm{i}\alpha_2}=|z_1|\cdot|z_2|\cdot \mathrm{e}^{\mathrm{i}(\alpha_1+\alpha_2)}}
	\end{align*}
\end{center}
\newpage
\noindent
Beispiel:\\
\begin{align*}
-1=\mathrm{i}^2=\bigg(\cos\bigg(\frac{\pi}{2}\bigg)+\mathrm{i}\;\sin\bigg(\frac{\pi}{2}\bigg)\bigg)^2=\big(\mathrm{e}^{\mathrm{i}\frac{\pi}{2}}\big)^2=\mathrm{e}^{\mathrm{i}\pi}
\end{align*}
\begin{center}
	\begin{align*}
	\boxed{-1=\mathrm{e}^{\mathrm{i}\pi}}
	\end{align*}
\end{center}
\underline {Motivation der Polarkoordinaten:}\\
Berechne: $ z^n $ mit $ z \in  \mathbb{C} $\\
\begin{align*}
z^n=(x+y\mathrm{i})^n=\big(|z|\mathrm{e}^{\mathrm{i}\alpha}\big)^n=|z|^n\cdot \mathrm{e}^{\mathrm{i}\alpha n}
\end{align*}
Bemerkung: $ (x+y\mathrm{i})^n $ ist nicht leicht zu berechnen!!\vspace*{2mm}\\
Beispiel:\\
Berechne: $(2+2\mathrm{i})^8 $\\
\begin{align*}
(2+2\mathrm{i})^8=(2(1+\mathrm{i}))^8=2^8(1+\mathrm{i})^8=2^8(|1+\mathrm{i}|\cdot \mathrm{e}^{\mathrm{i}\alpha})^8=2^8(\sqrt{1^2+1^2}\cdot \mathrm{e}^{\mathrm{i}\alpha})^8
\end{align*}
Jetzt benötigt man noch $ \alpha $ :\\
$ \alpha=? $\\
\begin{center}

\begin{tikzpicture}[scale=2.5]

% Gitternetzlinien
\draw[xstep=1cm, ystep=1cm, very thin, color=lightgray] (-2,-1) grid (2,2);
% x-Achse
\draw[->] (-2,0) -- (2,0) node[below,xshift=-1em] {$\mathrm{Re}$}; 
%Zahlen auf x-Achse
\foreach \x in {-1,...,-1}
\draw[shift={(\x,0)},color=black] (0pt,2pt) -- (0pt,-2pt) node[below] {\footnotesize $\x$};
\foreach \x in {1,...,1}
\draw[shift={(\x,0)},color=black] (0pt,2pt) -- (0pt,-2pt) node[below] {\footnotesize $\x$};
% y-Achse 
\draw[->] (0,-1) -- (0,2) node[left,yshift=-1em] {$\mathrm{Im}$};%node[above left]
%Zahlen auf y-Achse
\foreach \y in {-1,...,-1}
\draw[shift={(0,\y)},color=black] (2pt,0pt) -- (-2pt,0pt) node[left] {\footnotesize $\y$};
\foreach \y in {1,...,1}
\draw[shift={(0,\y)},color=black] (2pt,0pt) -- (-2pt,0pt) node[left] {\footnotesize $\y$};
%Ursprung
%\draw[color=black] (0pt,-10pt) node[right] {\footnotesize $0$};
\draw[color=black] (-4pt,-4pt) node[right] {\footnotesize $0$};
%\draw[thick,color=black] plot[samples=200, domain=2:4] 
%(\x,{2}) node[below, xshift=-5em] {$h(x)$};

%\draw[thick,color=black] plot[samples=200, domain=-1:2] 
%(\x,{\x-1}) node[below, xshift=1em]{$g(x)$};

%\draw [blue] (0,0) circle (1);
\filldraw[fill=green!20,draw=green!50!black] (0,0) -- (8mm,0mm) arc (0:45:8mm) -- cycle; 
\draw [blue,->] (0,0)-.(1,1)node[below,xshift= 0.5em,yshift=1em] {\footnotesize $1+i$};
%\draw [blue,dashed] (0,.7)-.(0.7,0.7);
%\draw [blue,dashed] (0.7,0)-.(0.7,0.7);
%\draw [blue] (0.7,0) node[below,xshift= 0em,yshift=-0.2em] {\footnotesize $x$};
\draw [blue] (0.3,0.3) node[below,xshift= -0.5em,yshift=1.8em,rotate=45] {\footnotesize $|\sqrt{2}|$};
%\draw [blue] (0,0.7) node[below,xshift=- 0.51em,yshift=1em] {\footnotesize $y$};


%\draw [red,thick] (0,0)-.(0.7,0) node[midway,xshift= 0em,yshift=-0.5em] {\footnotesize $cos(\alpha)$};
%\draw [red,thick] (0,0)-.(0,0.7) node[midway,xshift= -0.50em,yshift=-0em,rotate=90] {\footnotesize $sin(\alpha)$};
\draw [red] (0.2,0.2) node[xshift= 0.8em,yshift=-0.6em,rotate=0] {\footnotesize $\alpha=\frac{\pi}{4}$};
\end{tikzpicture}\\
\end{center}
dann gilt:\\
\begin{align*}
(2+2\mathrm{i})^8=2^8=2^8\big(\sqrt{2}\big)^8\cdot \mathrm{e}^{\mathrm{i} \frac{\pi}{4}\cdot 8}=2^8 2^4 \mathrm{e}^{\mathrm{i}2\pi}=2^{12} \big(\underbrace {\mathrm{e}^{\mathrm{i}\pi}}_{=-1}\big)^2=4096\cdot (-1)^2=4096
\end{align*}
\subsubsection{ Aufgabe zur Multiplikation komplexer Zahlen in Polarkoordinaten}
\href{https://youtu.be/JumERKyIWbM}{Video 12}\\
Stelle die folgende komplexe Zahl in der Form von $ x+\mathrm{i}y $ dar:\\
\begin{align*}
(2+2\mathrm{i})^n+(2-2\mathrm{i})^n \text{ für } n\in \mathbb{N}
\end{align*}
\newpage
\noindent
\textbullet \;\; Umformung der komplexen Zahl:\\
\begin{quote}
	\begin{align*}
	(2+2\mathrm{i})^n+(2-2\mathrm{i})^n&=(2(1+\mathrm{i}))^n+(2(1-\mathrm{i}))^n\\
	                 &=2^n(1+\mathrm{i})^n+2^n(1-\mathrm{i})^n\\
	                 &=2^n \big((1+\mathrm{i})^n+(1-\mathrm{i})^n  \big)
	\end{align*}
\end{quote}
\textbullet \;\; Darstellung von $ (1+\mathrm{i}) $ und $ (1-\mathrm{i}) $ in Polarkoordinaten.\\
\begin{center}
	
	\begin{tikzpicture}[scale=2.5]
		\filldraw [->,draw=black,fill=red!20] (0,0)--(9mm,0mm) arc (0:315:9mm);
		\filldraw[->,draw=black,fill=green!20] (0,0) -- (7mm,0mm) arc (0:45:7mm); 
		\filldraw[->,fill=orange!20,draw=black,rotate=-45] (0,0) -- (8mm,0mm) arc (0:45:8mm); 
		
	% Gitternetzlinien
	\draw[xstep=1cm, ystep=1cm, very thin, color=lightgray] (-2,-2) grid (2,2);
	% x-Achse
	\draw[->] (-2,0) -- (2,0) node[below,xshift=-1em] {$\mathrm{\mathrm{Re}}$}; 
	%Zahlen auf x-Achse
	\foreach \x in {-1,...,-1}
	\draw[shift={(\x,0)},color=black] (0pt,2pt) -- (0pt,-2pt) node[below] {\footnotesize $\x$};
	\foreach \x in {1,...,1}
	\draw[shift={(\x,0)},color=black] (0pt,2pt) -- (0pt,-2pt) node[below] {\footnotesize $\x$};
	% y-Achse 
	\draw[->] (0,-2) -- (0,2) node[left,yshift=-1em] {$\mathrm{\mathrm{Im}}$};%node[above left]
	%Zahlen auf y-Achse
	\foreach \y in {-1,...,-1}
	\draw[shift={(0,\y)},color=black] (2pt,0pt) -- (-2pt,0pt) node[left] {\footnotesize $\y$};
	\foreach \y in {1,...,1}
	\draw[shift={(0,\y)},color=black] (2pt,0pt) -- (-2pt,0pt) node[left] {\footnotesize $\y$};
	%Ursprung
	\draw[color=black] (-4pt,-4pt) node[right] {\footnotesize $0$};
	

	\draw [blue,->] (0,0)-.(1,1)node[below,xshift= 0.5em,yshift=1em] {\footnotesize $1+\mathrm{i}$};

	\draw [blue,->] (0,0)-.(1,-1)node[below,xshift= 0.5em,yshift=1em] {\footnotesize $1-\mathrm{i}$};

	\draw [blue] (0.2,0.2) node[xshift= 0.8em,yshift=-0.6em,rotate=0] {\footnotesize $\alpha=\frac{\pi}{4}$};
	\draw [blue] (-0.2,-0.2) node[xshift= 3.8em,yshift=0.1em,rotate=0] {\footnotesize $\overset{\AC}{\alpha}=\frac{\pi}{4}$};
    \draw [blue] (-0.6,0.6) node[xshift= 1.3em,yshift=-0.9em,rotate=0] {\footnotesize $\beta=2\pi-\overset{\AC}{\alpha}$};
	 
	\end{tikzpicture}\\
\end{center}
\textbullet \;\; $ 1+\mathrm{i}=|1+\mathrm{i}|e^{\mathrm{i}\alpha}=\sqrt{2}\mathrm{e}^{\mathrm{i}\frac{\pi}{4}}$\\
\textbullet \;\; $ 1-\mathrm{i}=|1-\mathrm{i}|\mathrm{e}^{\mathrm{i}\beta} $\\
\begin{quote}
	$\beta=2\pi-\overset{\AC}{\alpha}=2\pi-\frac{\pi}{4}=\frac{7\pi}{4}  $\\
	$ 1-\mathrm{i}=\sqrt{2}\mathrm{e}^{\mathrm{i}\frac{7\pi}{4}} $
\end{quote}
Dann ist:\\
\begin{align*}
2^n \big((1+\mathrm{i})^n+(1-\mathrm{i})^n  \big)&=2^n \bigg(\big(\sqrt{2}\mathrm{e}^{\mathrm{i}\frac{\pi}{4}}\big)+\big(\sqrt{2} \mathrm{e}^{\mathrm{i}\frac{7\pi}{4}} \big)  \bigg)^n\\
&=2^n\big(2^{\frac{n}{2}}\mathrm{e}^{\mathrm{i}\cdot n\frac{\pi}{4}}+2^{\frac{n}{2}}\mathrm{e}^{\mathrm{i}\cdot n\frac{7\pi}{4}}\big)\\
&=2^{n+\frac{n}{2}}\big(\mathrm{e}^{\mathrm{i}\cdot n\frac{\pi}{4}}+\mathrm{e}^{\mathrm{i}\cdot n\frac{7\pi}{4}}\big)
\end{align*}
Es gilt: $\mathrm{e}^{\mathrm{i}\alpha}=\cos(\alpha)+\mathrm{i}\cdot \sin(\alpha)$\\
Dies eingesetzt ergibt:\\ 
\begin{align*}
2^n \big((1+\mathrm{i})^n+(1-\mathrm{i})^n  \big)&=2^{n+\frac{n}{2}}\big(\cos(\genfrac{}{}{}{2}{ n\pi}{ 4}) +\mathrm{i}\cdot \sin(\genfrac{}{}{}{2}{n\pi}{4})+ \cos(\genfrac{}{}{}{2}{  7 n \pi }{ 4}) +\mathrm{i}\cdot \sin(\genfrac{}{}{}{2}{ 7 n\pi  }{4})  \big)\\
&
\end{align*}
\begin{align*}
(2+2\mathrm{i})^n+(2-2\mathrm{i})^n=\underbrace {2^{n+\frac{n}{2}}\big(\cos(\genfrac{}{}{}{2}{n\pi }{4})+\cos(\genfrac{}{}{}{2}{ 7 n\pi  }{4}) \big)}_{x}+\underbrace {\mathrm{i} }_{\mathrm{i}}
\underbrace {2^{n+\frac{n}{2}}(sin(\genfrac{}{}{}{2}{n\pi}{4})+\sin(\genfrac{}{}{}{2}{ 7 n\pi }{4 }   )}_{y}
\end{align*}
\subsection{Wurzeln von komplexen Zahlen}
\href{https://youtu.be/5IVSZmGtt2Q}{Video 13}\\
Löse $ z^4=1 $:\\


\begin{center}
	Die Gleichung hat 4 Lösungen: $ z_1=1,z_2=\mathrm{i},z_3=-1,z_4=-\mathrm{i} $\vspace*{4mm}\\
	\begin{tikzpicture}[scale=2.5]
	
	% Gitternetzlinien
	\draw[xstep=1cm, ystep=1cm, very thin, color=lightgray] (-2,-2) grid (2,2);
	% x-Achse
	\draw[->] (-2,0) -- (2,0) node[below,xshift=-1em] {$\mathrm{Re}$}; 
	%Zahlen auf x-Achse
	\foreach \x in {-1,...,-1}
	\draw[shift={(\x,0)},color=black] (0pt,2pt) -- (0pt,-2pt) node[below] {\footnotesize $\x$};
	\foreach \x in {1,...,1}
	\draw[shift={(\x,0)},color=black] (0pt,2pt) -- (0pt,-2pt) node[below] {\footnotesize $\x$};
	% y-Achse 
	\draw[->] (0,-2) -- (0,2) node[left,yshift=-1em] {$\mathrm{Im}$};%node[above left]
	%Zahlen auf y-Achse
	\foreach \y in {-1,...,-1}
	\draw[shift={(0,\y)},color=black] (2pt,0pt) -- (-2pt,0pt) node[left] {\footnotesize $\y$};
	\foreach \y in {1,...,1}
	\draw[shift={(0,\y)},color=black] (2pt,0pt) -- (-2pt,0pt) node[left] {\footnotesize $\y$};
	%Ursprung
	%\draw[color=black] (0pt,-10pt) node[right] {\footnotesize $0$};
	\draw[color=black] (-4pt,-4pt) node[right] {\footnotesize $0$};
	%\draw[thick,color=black] plot[samples=200, domain=2:4] 
	%(\x,{2}) node[below, xshift=-5em] {$h(x)$};
	
	%\draw[thick,color=black] plot[samples=200, domain=-1:2] 
	%(\x,{\x-1}) node[below, xshift=1em]{$g(x)$};
	
	\draw [blue] (0,0) circle (1);
	%\filldraw[fill=green!20,draw=green!50!black] (0,0) -- (5mm,0mm) arc (0:45:5mm) -- cycle; 
	%\draw [blue,->] (0,0)-.(0.7,0.7)node[below,xshift= 0.5em,yshift=1em] {\footnotesize $z$};
	%\draw [blue,dashed] (0,.7)-.(0.7,0.7);
	%\draw [blue,dashed] (0.7,0)-.(0.7,0.7);
	%\draw [blue] (0.7,0) node[below,xshift= 0em,yshift=-0.2em] {\footnotesize $x$};
	%\draw [blue] (0.3,0.3) node[below,xshift= -0.5em,yshift=1.8em,rotate=45] {\footnotesize $|z|$};
	\draw [blue] (1.3,0.2) node[below,xshift=0em,yshift=0em] {\footnotesize $z_0=\mathrm{e}^{\mathrm{i}\cdot 0}$};
	\filldraw [blue] (1,0) circle (0.03);
	\draw [blue] (0.3,1.2) node[below,xshift=0em,yshift=0em] {\footnotesize $z_1=\mathrm{e}^{\mathrm{i}\frac{\pi}{2}}$};
	\filldraw [blue] (0,1) circle (0.03);
	\draw [blue] (-1.32,0.2) node[below,xshift=0em,yshift=0em] {\footnotesize $z_2=\mathrm{e}^{\mathrm{i}\pi}$};
	\filldraw [blue] (-1,0) circle (0.03);
	\draw [blue] (0.3,-1.1) node[below,xshift=0em,yshift=0em] {\footnotesize $z_3=\mathrm{e}^{\mathrm{i}\frac{3\pi}{2}}$};
	\filldraw [blue] (0,-1) circle (0.03);
	\end{tikzpicture}\\
\end{center}
es gilt:\\
\begin{align*}
z^4&=1=\mathrm{e}^{\mathrm{i}2\pi}\\
z&=\sqrt[4]{\mathrm{e}^{\mathrm{i}2\pi}}=\mathrm{e}^{\mathrm{i}\frac{2\pi}{4}}\\
\end{align*}
\begin{align*}
&z_0=\mathrm{e}^{\mathrm{i}\frac{2\pi}{4}\cdot 0}=\mathrm{e}^{\mathrm{i}\cdot 0}\\
&z_1=\mathrm{e}^{\mathrm{i}\frac{2\pi}{4}\cdot 1}=\mathrm{e}^{\mathrm{i}\frac{\pi}{2}}\\
&z_2=\mathrm{e}^{\mathrm{i}\frac{2\pi}{4}\cdot 2}=\mathrm{e}^{\mathrm{i}\pi}\\
&z_3=\mathrm{e}^{\mathrm{i}\frac{2\pi}{4}\cdot 3}=\mathrm{e}^{\mathrm{i}\frac{3\pi}{2}}\\
\end{align*}
Allgemein gilt dann:\\
Die Lösungen für $ z^n=1 $ sind:\\
\begin{center}
\boxed{z_k =\mathrm{e}^{\mathrm{i}\frac{2\pi\cdot k}{n}} \text{\;\;\;\;   für } k=0,1,2 \ldots n-1 }
\end{center}
\newpage
\noindent
\href{https://youtu.be/Uqpu2trEjM4}{Video 14}\\
Gegeben:\\$ z^n-a=0 $ für $ a\in  \mathbb{C} $\\
Gesucht: \\Alle Lösungen aus $  \mathbb{C}$\\
Es gilt:\\
\begin{align*}
&z^n=a=|a|e^{\mathrm{i}\alpha}\\
&z=\sqrt[n]{|a|}\mathrm{e}^{\mathrm{i} \frac{\alpha}{n}}
\end{align*}
Dies ist für den Winkel $ \alpha $ eine Lösung.\\
\begin{center}
	Die Gleichung hat k Lösungen an den Stellen: $ \mathrm{i}\frac{\alpha}{n}+\mathrm{i}\frac{2k\pi}{n}$ für $ k=0,1,2\ldots n-1 $\vspace*{4mm}\\
	\begin{tikzpicture}[scale=2.5]
	
	% Gitternetzlinien
	\draw[xstep=1cm, ystep=1cm, very thin, color=lightgray] (-2,-2) grid (2,2);
	% x-Achse
	\draw[->] (-2,0) -- (2,0) node[below,xshift=-1em] {$\mathrm{Re}$}; 
	%Zahlen auf x-Achse
	\foreach \x in {-1,...,-1}
	\draw[shift={(\x,0)},color=black] (0pt,2pt) -- (0pt,-2pt) node[below] {\footnotesize $\x$};
	\foreach \x in {1,...,1}
	\draw[shift={(\x,0)},color=black] (0pt,2pt) -- (0pt,-2pt) node[below] {\footnotesize $\x$};
	% y-Achse 
	\draw[->] (0,-2) -- (0,2) node[left,yshift=-1em] {$\mathrm{Im}$};%node[above left]
	%Zahlen auf y-Achse
	\foreach \y in {-1,...,-1}
	\draw[shift={(0,\y)},color=black] (2pt,0pt) -- (-2pt,0pt) node[left] {\footnotesize $\y$};
	\foreach \y in {1,...,1}
	\draw[shift={(0,\y)},color=black] (2pt,0pt) -- (-2pt,0pt) node[left] {\footnotesize $\y$};
	%Ursprung
	%\draw[color=black] (0pt,-10pt) node[right] {\footnotesize $0$};
	\draw[color=black] (-4pt,-4pt) node[right] {\footnotesize $0$};
	%\draw[thick,color=black] plot[samples=200, domain=2:4] 
	%(\x,{2}) node[below, xshift=-5em] {$h(x)$};
	
	%\draw[thick,color=black] plot[samples=200, domain=-1:2] 
	%(\x,{\x-1}) node[below, xshift=1em]{$g(x)$};
	
	\draw [blue] (0,0) circle (2);
	%\filldraw[fill=green!20,draw=green!50!black] (0,0) -- (5mm,0mm) arc (0:45:5mm) -- cycle; 
	%\draw [blue,->] (0,0)-.(0.7,0.7)node[below,xshift= 0.5em,yshift=1em] {\footnotesize $z$};
	%\draw [blue,dashed] (0,.7)-.(0.7,0.7);
	%\draw [blue,dashed] (0.7,0)-.(0.7,0.7);
	%\draw [blue] (0.7,0) node[below,xshift= 0em,yshift=-0.2em] {\footnotesize $x$};
	%\draw [blue] (0.3,0.3) node[below,xshift= -0.5em,yshift=1.8em,rotate=45] {\footnotesize $|z|$};
	\draw [blue] (1.41,1.41) node[below,xshift=2em,yshift=1em] {\footnotesize $\mathrm{i}\frac{\alpha}{n}$};
	\draw [blue,->,thick] (0,0)--(1.41,1.41);
	\draw [blue] (-1.41,1.41) node[below,xshift=3em,yshift=1em] {\footnotesize $\mathrm{i}\frac{\alpha}{n}+\mathrm{i}\frac{\alpha k 2\pi}{n}$};
	\draw [blue,->,thick] (0,0)--(-1.41,1.41);
	%\filldraw [blue] (1.41,1.41) circle (0.03);
	
	\end{tikzpicture}\\
\end{center}
\begin{align*}
 \boxed{z_k=\sqrt[n]{|a|}\cdot \mathrm{e}^{\mathrm{i}\frac{\alpha}{n}+\mathrm{i}\frac{2k\pi}{n}}\;\;\;\;\text{für } k=0,1,2\ldots n-1}
\end{align*}
Beispiel:\\
Gegeben: $ z^3=2\mathrm{i} $\\
Gesucht: Alle Lösungen aus $  \mathbb{C} $\\
Es gilt:\\
\begin{align*}
z^3=2\mathrm{i}=0+2\mathrm{i}=|2|\mathrm{e}^{\mathrm{i}\alpha}
\end{align*}
\newpage
\noindent
Winkel $ \alpha $ bestimmen:\vspace*{2mm}\\
\begin{center}
	
	\begin{tikzpicture}[scale=2.5]
	
	% Gitternetzlinien
	\draw[xstep=1cm, ystep=1cm, very thin, color=lightgray] (-2,-2) grid (2,2);
	% x-Achse
	\draw[->] (-2,0) -- (2,0) node[below,xshift=-1em] {$\mathrm{Re}$}; 
	%Zahlen auf x-Achse
	\foreach \x in {-1,...,-1}
	\draw[shift={(\x,0)},color=black] (0pt,2pt) -- (0pt,-2pt) node[below] {\footnotesize $\x$};
	\foreach \x in {1,...,1}
	\draw[shift={(\x,0)},color=black] (0pt,2pt) -- (0pt,-2pt) node[below] {\footnotesize $\x$};
	% y-Achse 
	\draw[->] (0,-2) -- (0,2) node[left,yshift=-1em] {$\mathrm{Im}$};%node[above left]
	%Zahlen auf y-Achse
	\foreach \y in {-1,...,-1}
	\draw[shift={(0,\y)},color=black] (2pt,0pt) -- (-2pt,0pt) node[left] {\footnotesize $\y$};
	\foreach \y in {1,...,1}
	\draw[shift={(0,\y)},color=black] (2pt,0pt) -- (-2pt,0pt) node[left] {\footnotesize $\y$};
	%Ursprung
	%\draw[color=black] (0pt,-10pt) node[right] {\footnotesize $0$};
	\draw[color=black] (-4pt,-4pt) node[right] {\footnotesize $0$};
	%\draw[thick,color=black] plot[samples=200, domain=2:4] 
	%(\x,{2}) node[below, xshift=-5em] {$h(x)$};
	
	%\draw[thick,color=black] plot[samples=200, domain=-1:2] 
	%(\x,{\x-1}) node[below, xshift=1em]{$g(x)$};
	
	\draw [blue] (0,0) circle (2);
	\filldraw[->,draw=black,fill=green!20] (0,0) -- (8mm,0mm) arc (0:90:8mm); 
	%\filldraw[->,fill=green!20,draw=green!50!black] (0,0) -- (9mm,0mm) arc (0:90:9mm) -- cycle; 
	\draw [blue] (0.2,0.2) node[xshift= 0.8em,yshift=0.6em,rotate=0] {\footnotesize $\alpha=\frac{\pi}{2}$};
	%\draw [blue,->] (0,0)-.(0.7,0.7)node[below,xshift= 0.5em,yshift=1em] {\footnotesize $z$};
	%\draw [blue,dashed] (0,.7)-.(0.7,0.7);
	%\draw [blue,dashed] (0.7,0)-.(0.7,0.7);
	%\draw [blue] (0.7,0) node[below,xshift= 0em,yshift=-0.2em] {\footnotesize $x$};
	%\draw [blue] (0.3,0.3) node[below,xshift= -0.5em,yshift=1.8em,rotate=45] {\footnotesize $|z|$};
	\draw [blue] (0,2) node[below,xshift=2em,yshift=-1em] {\footnotesize $0+2\mathrm{i}$};
	\draw [blue,->,thick] (0,0)--(0,2);
	%\draw [blue] (-1.41,1.41) node[below,xshift=3em,yshift=1em] {\footnotesize $i\frac{\alpha}{n}+i\frac{\alpha k 2\pi}{n}$};
	%\draw [blue,->] (0,0)--(-1.41,1.41);
	%\filldraw [blue] (1.41,1.41) circle (0.03);
	
	\end{tikzpicture}\\
\end{center}
Winkel $ \alpha=\frac{\pi}{2} $
also gilt:
\begin{align*}
z^3=|2|\mathrm{e}^{\mathrm{i}\frac{\pi}{2}}
\end{align*}
D.h.:\\
\begin{align*}
z_k=\sqrt[3]{2}\cdot \mathrm{e}^{\mathrm{i} \frac{\pi}{2\cdot 3}+\mathrm{i}\frac{2k\pi}{3}}\;\;\;\;\text{für } k=0,1,2 \ldots n-1
\end{align*}
\subsubsection{Aufgabe zum Thema komplexe Lösungen von Gleichungen}
\href{https://youtu.be/UOdKG8lbwiQ}{Video 15}\\
Berechne alle Lösungen von:\\
\begin{tabular}[t]{lll}% l=linksbündig r=rechtsbündig
	(i) & $z^5=3\mathrm{i}+3  $  \\
	(ii) & $z^7=16\mathrm{i}-16  $ \\
	(iii)&$ \mathrm{i}z-\overline{z} =4$	
\end{tabular} \vspace*{6mm}\\
(i)\;\;$ z^5=3\mathrm{i}+3 $\\
\begin{quote}
Es gilt:\\
\begin{align*}
z^5=3\mathrm{i}+3=3(1+\mathrm{i})=3\sqrt{1^2+1^2}\;\;\mathrm{e}^{\mathrm{i}\alpha}
\end{align*}
\begin{center}
	Winkel $ \alpha $ bestimmen:\vspace*{2mm}\\
	\begin{tikzpicture}[scale=2.5]
	\filldraw[->,draw=black,fill=green!20] (0,0) -- (8mm,0mm) arc (0:45:8mm); 
	% Gitternetzlinien
	\draw[xstep=1cm, ystep=1cm, very thin, color=lightgray] (-2,-1) grid (2,2);
	% x-Achse
	\draw[->] (-2,0) -- (2,0) node[below,xshift=-1em] {$\mathrm{Re}$}; 
	%Zahlen auf x-Achse
	\foreach \x in {-1,...,-1}
	\draw[shift={(\x,0)},color=black] (0pt,2pt) -- (0pt,-2pt) node[below] {\footnotesize $\x$};
	\foreach \x in {1,...,1}
	\draw[shift={(\x,0)},color=black] (0pt,2pt) -- (0pt,-2pt) node[below] {\footnotesize $\x$};
	% y-Achse 
	\draw[->] (0,-1) -- (0,2) node[left,yshift=-1em] {$\mathrm{Im}$};%node[above left]
	%Zahlen auf y-Achse
	\foreach \y in {-1,...,-1}
	\draw[shift={(0,\y)},color=black] (2pt,0pt) -- (-2pt,0pt) node[left] {\footnotesize $\y$};
	\foreach \y in {1,...,1}
	\draw[shift={(0,\y)},color=black] (2pt,0pt) -- (-2pt,0pt) node[left] {\footnotesize $\y$};
	%Ursprung
	%\draw[color=black] (0pt,-10pt) node[right] {\footnotesize $0$};
	\draw[color=black] (-4pt,-4pt) node[right] {\footnotesize $0$};
	%\draw[thick,color=black] plot[samples=200, domain=2:4] 
	%(\x,{2}) node[below, xshift=-5em] {$h(x)$};
	
	%\draw[thick,color=black] plot[samples=200, domain=-1:2] 
	%(\x,{\x-1}) node[below, xshift=1em]{$g(x)$};
	
	%\draw [blue] (0,0) circle (1);
	%\filldraw[fill=green!20,draw=green!50!black] (0,0) -- (8mm,0mm) arc (0:45:8mm) -- cycle; 
	\draw [blue,->] (0,0)-.(1,1)node[below,xshift= 0.5em,yshift=1em] {\footnotesize $1+\mathrm{i}$};
	%\draw [blue,dashed] (0,.7)-.(0.7,0.7);
	%\draw [blue,dashed] (0.7,0)-.(0.7,0.7);
	%\draw [blue] (0.7,0) node[below,xshift= 0em,yshift=-0.2em] {\footnotesize $x$};
	\draw [blue] (0.3,0.3) node[below,xshift= -0.5em,yshift=1.8em,rotate=45] {\footnotesize $|\sqrt{2}|$};
	%\draw [blue] (0,0.7) node[below,xshift=- 0.51em,yshift=1em] {\footnotesize $y$};
	
	
	%\draw [red,thick] (0,0)-.(0.7,0) node[midway,xshift= 0em,yshift=-0.5em] {\footnotesize $cos(\alpha)$};
	%\draw [red,thick] (0,0)-.(0,0.7) node[midway,xshift= -0.50em,yshift=-0em,rotate=90] {\footnotesize $sin(\alpha)$};
	\draw [red] (0.2,0.2) node[xshift= 0.8em,yshift=-0.6em,rotate=0] {\footnotesize $\alpha=\frac{\pi}{4}$};
	\end{tikzpicture}\\
\end{center}
\begin{align*}
z^5=3\sqrt{2}\;\;\mathrm{e}^{\mathrm{i}\frac{\pi}{4}}
\end{align*}
Lösungen sind:\\
\begin{align*}
z_k=\sqrt[5]{\scriptscriptstyle 3\sqrt{2}}\;\;\mathrm{e}^{\mathrm{i}\frac{\pi}{5\cdot4}+\mathrm{i}\frac{2\pi k}{5}}=\sqrt[5]{\scriptscriptstyle 3\sqrt{2}}\;\;\mathrm{e}^{\mathrm{i}\frac{\pi}{20}+\mathrm{i}\frac{2\pi k}{5}}\;\;\;\text{für } k=0,1,2\ldots n-1
\end{align*}
\end{quote}
(ii)\;\;$ z^7=16\mathrm{i}-16 $\\
\begin{quote}
	
\begin{align*}
 z^7=16\mathrm{i}-16=16(-1+\mathrm{i})=16\sqrt{2}\;\;\mathrm{e}^{\mathrm{i}\alpha}
\end{align*}
\begin{center}
	Winkel $ \alpha $ bestimmen:\vspace*{2mm}\\
	\begin{tikzpicture}[scale=2.5]
	\filldraw[->,draw=black,fill=green!20] (0,0) -- (8mm,0mm) arc (0:135:8mm); 
	% Gitternetzlinien
	\draw[xstep=1cm, ystep=1cm, very thin, color=lightgray] (-2,-1) grid (2,2);
	% x-Achse
	\draw[->] (-2,0) -- (2,0) node[below,xshift=-1em] {$\mathrm{Re}$}; 
	%Zahlen auf x-Achse
	\foreach \x in {-1,...,-1}
	\draw[shift={(\x,0)},color=black] (0pt,2pt) -- (0pt,-2pt) node[below] {\footnotesize $\x$};
	\foreach \x in {1,...,1}
	\draw[shift={(\x,0)},color=black] (0pt,2pt) -- (0pt,-2pt) node[below] {\footnotesize $\x$};
	% y-Achse 
	\draw[->] (0,-1) -- (0,2) node[left,yshift=-1em] {$\mathrm{Im}$};%node[above left]
	%Zahlen auf y-Achse
	\foreach \y in {-1,...,-1}
	\draw[shift={(0,\y)},color=black] (2pt,0pt) -- (-2pt,0pt) node[left] {\footnotesize $\y$};
	\foreach \y in {1,...,1}
	\draw[shift={(0,\y)},color=black] (2pt,0pt) -- (-2pt,0pt) node[left] {\footnotesize $\y$};
	%Ursprung
	%\draw[color=black] (0pt,-10pt) node[right] {\footnotesize $0$};
	\draw[color=black] (-4pt,-4pt) node[right] {\footnotesize $0$};
	
	
	\draw [blue,->] (0,0)-.(-1,1)node[below,xshift= -0.5em,yshift=1em] {\footnotesize $-1+\mathrm{i}$};

	\draw [blue] (-0.15,0.3) node[below,xshift= -0.5em,yshift=1.8em,rotate=-45] {\footnotesize $|\sqrt{2}|$};

	\draw [red] (0.2,0.2) node[xshift= 0.8em,yshift=-0.6em,rotate=0] {\footnotesize $\alpha=\frac{\pi}{2}+\frac{\pi}{4}$};
	\end{tikzpicture}\\
\end{center}
\begin{align*}
z^7=16\sqrt{2}\;\;\mathrm{e}^{\mathrm{i}\frac{3}{4}\pi}
\end{align*}
Lösungen sind:\\
\begin{align*}
z_k=\sqrt[7]{\scriptscriptstyle 16\sqrt{2}}\;\;\mathrm{e}^{\mathrm{i}\frac{3\pi}{4\cdot7}+\frac{2\pi k}{7}}=\sqrt[7]{\scriptscriptstyle 16\sqrt{2}}\;\;\mathrm{e}^{\mathrm{i}\frac{3\pi}{28}+\frac{2\pi k}{7}}\;\;\;\;\text{für }k=0,1,2\ldots n-1
\end{align*}
\end{quote}
(iii)\;\; $\mathrm{i}z-\mathrm{i} \overline{z}=4$\\
\begin{quote}
Es gilt mit: $ z=x+\mathrm{i}y $\\
	\begin{center}
	
	\begin{tabular}[t]{r@{$=$}ll}% l=linksbündig r=rechtsbündig
		$\mathrm{i}(x+\mathrm{i}y)-\mathrm{i}(x-\mathrm{i}y)$ & $4$\\
		$\mathrm{i}x+\mathrm{i}^2y-\mathrm{i}x+\mathrm{i}^2y$ & $4 $ \\	
		$ -2y $&$ 4 $\\
		$ y $&$-2  $ und $ x $ beliebig\\
	\end{tabular} \\
\end{center}
Lösung:$ \{ x+\mathrm{i}y\;|\;x\in \mathbb{R}\wedge y=-2\} $
\end{quote}

\subsection{Bestimmung des Winkels bei Polarkoordinaten}
\href{https://youtu.be/FokqSBZakd4}{Video 16}\\
Sei $z \in \mathbb{C}\setminus\{0\} $ und $ z=x+\mathrm{i}y=|z|\;\mathrm{e}^{\mathrm{i}\alpha} $\vspace*{2mm}\\
Wie bestimmt man den Winkel $ \alpha $ ?\\
\begin{center}
	\begin{align*}
	z=|z|(\cos(\alpha)+\mathrm{i}\;\sin(\alpha))
	\end{align*}
	\begin{tikzpicture}[scale=2.5]
	
	% Gitternetzlinien
	\draw[xstep=1cm, ystep=1cm, very thin, color=lightgray] (-2,-2) grid (2,2);
	% x-Achse
	\draw[->] (-2,0) -- (2,0) node[below,xshift=-2em] {$\mathrm{Re}$}; 
	%Zahlen auf x-Achse
	\foreach \x in {-1,...,-1}
	\draw[shift={(\x,0)},color=black] (0pt,2pt) -- (0pt,-2pt) node[below] {\footnotesize $\x$};
	\foreach \x in {1,...,1}
	\draw[shift={(\x,0)},color=black] (0pt,2pt) -- (0pt,-2pt) node[below] {\footnotesize $\x$};
	% y-Achse 
	\draw[->] (0,-2) -- (0,2) node[left,yshift=-0.6em] {$\mathrm{Im}$};%node[above left]
	%Zahlen auf y-Achse
	\foreach \y in {-1,...,-1}
	\draw[shift={(0,\y)},color=black] (2pt,0pt) -- (-2pt,0pt) node[left] {\footnotesize $\y$};
	\foreach \y in {1,...,1}
	\draw[shift={(0,\y)},color=black] (2pt,0pt) -- (-2pt,0pt) node[left] {\footnotesize $\y$};
	%Ursprung
	%\draw[color=black] (0pt,-10pt) node[right] {\footnotesize $0$};
	\draw[color=black] (-4pt,-4pt) node[right] {\footnotesize $0$};
	%\draw[thick,color=black] plot[samples=200, domain=2:4] 
	%(\x,{2}) node[below, xshift=-5em] {$h(x)$};
	
	%\draw[thick,color=black] plot[samples=200, domain=-1:2] 
	%(\x,{\x-1}) node[below, xshift=1em]{$g(x)$};
	
	\draw [blue] (0,0) circle (1.8);
	\filldraw[fill=green!20,draw=green!50!black] (0,0) -- (5mm,0mm) arc (0:45:5mm) -- cycle; 
	\draw [blue,->] (0,0)-.(1.27,1.27)node[below,xshift= 2em,yshift=1em] {\footnotesize $z=x+\mathrm{i}y$};
	\draw [blue,dashed] (0,1.27)-.(1.27,1.27);
	\draw [blue,dashed] (1.27,0)-.(1.27,1.27);
	\draw [blue] (1.27,0) node[below,xshift= 0em,yshift=-0.2em] {\footnotesize $x$};
	\draw [blue] (0.5,0.5) node[below,xshift= -0.5em,yshift=1.8em,rotate=45] {\footnotesize $|z|$};
	\draw [blue] (0,1.27) node[below,xshift=- 0.51em,yshift=1em] {\footnotesize $y$};
	
	
	\draw [red,thick] (0,0)-.(1.27,0) node[midway,xshift= 0em,yshift=-0.5em] {\footnotesize $|z|\cos(\alpha)$};
	\draw [red,thick] (0,0)-.(0,1.27) node[midway,xshift= -0.50em,yshift=0em,rotate=90] {\footnotesize $|z|\sin(\alpha)$};
	\draw [red] (0.2,0.2) node[xshift= 0em,yshift=-0.6em,rotate=0] {\footnotesize $\alpha$};
	\end{tikzpicture}\\
\end{center}
Die Umkehrfunktion von $ \cos(\alpha) $ bzw. $ \sin(\alpha) $ sind:\\
\begin{align*}
&\arccos: [-1,1]\longmapsto [0,\pi]\\
&\arcsin: [-1,1]\longmapsto [\scriptscriptstyle -\genfrac{}{}{}{\pi}{\pi}{2 },\genfrac{}{}{}{2}{\pi}{2 }\displaystyle ]
\end{align*}
Dann gilt:\\
\begin{align*}
&x=|z|\cos(\alpha)\Leftrightarrow \cos(\alpha)=\frac{x}{|z|}&y=|z|sin(\alpha)\Leftrightarrow \cos(\alpha)=\frac{y}{|z|}\\
&\alpha=\arccos \bigg(\frac{x}{|z| }\bigg)
&\alpha=\arcsin \bigg(\frac{y }{|z| }\bigg)\\
\end{align*}
Bemerkung:\\
Falsch ist:\\
\begin{align*}
\alpha=\arccos(\genfrac{}{}{}{2}{x}{|z| })=\arcsin(\genfrac{}{}{}{2}{y}{|z| })
\end{align*}
Weil die Abbildung vom arccos $[0,\pi]  $ ist und vom arcsin $ [\scriptscriptstyle -\genfrac{}{}{}{\pi}{\pi}{2 },\genfrac{}{}{}{2}{\pi}{2 }\displaystyle ] $ ist. \vspace*{4mm}\\
\textbullet \;\; für $ \alpha [0,\pi]$ gilt:\\

\begin{center}
$ \alpha=\arccos(\genfrac{}{}{}{2}{x }{|z| }) $ für $[0,\pi]$\\
	\begin{tikzpicture}[scale=2.5]
	\filldraw [->,draw=black,fill=red!20] (0,0)--(18mm,0mm) arc (0:180:18mm);
	% Gitternetzlinien
	\draw[xstep=1cm, ystep=1cm, very thin, color=lightgray] (-2,-2) grid (2,2);
	% x-Achse
	\draw[->] (-2,0) -- (2,0) node[below,xshift=-2em] {$\mathrm{Re}$}; 
	%Zahlen auf x-Achse
	\foreach \x in {-1,...,-1}
	\draw[shift={(\x,0)},color=black] (0pt,2pt) -- (0pt,-2pt) node[below] {\footnotesize $\x$};
	\foreach \x in {1,...,1}
	\draw[shift={(\x,0)},color=black] (0pt,2pt) -- (0pt,-2pt) node[below] {\footnotesize $\x$};
	% y-Achse 
	\draw[->] (0,-2) -- (0,2) node[left,yshift=-0.6em] {$\mathrm{Im}$};%node[above left]
	%Zahlen auf y-Achse
	\foreach \y in {-1,...,-1}
	\draw[shift={(0,\y)},color=black] (2pt,0pt) -- (-2pt,0pt) node[left] {\footnotesize $\y$};
	\foreach \y in {1,...,1}
	\draw[shift={(0,\y)},color=black] (2pt,0pt) -- (-2pt,0pt) node[left] {\footnotesize $\y$};
	%Ursprung
	%\draw[color=black] (0pt,-10pt) node[right] {\footnotesize $0$};
	\draw[color=black] (-4pt,-4pt) node[right] {\footnotesize $0$};
	%\draw[thick,color=black] plot[samples=200, domain=2:4] 
	%(\x,{2}) node[below, xshift=-5em] {$h(x)$};
	
	%\draw[thick,color=black] plot[samples=200, domain=-1:2] 
	%(\x,{\x-1}) node[below, xshift=1em]{$g(x)$};
	
	%\draw [blue] (0,0) circle (1.8);
	\filldraw[fill=green!20,draw=green!50!black] (0,0) -- (5mm,0mm) arc (0:45:5mm) -- cycle; 
	\draw [blue,->] (0,0)-.(1.27,1.27)node[below,xshift= 2em,yshift=1em] {\footnotesize $z=x+\mathrm{i}y$};
	\draw [blue,dashed] (0,1.27)-.(1.27,1.27);
	\draw [blue,dashed] (1.27,0)-.(1.27,1.27);
	\draw [blue] (1.27,0) node[below,xshift= 0em,yshift=-0.2em] {\footnotesize $x$};
	\draw [blue] (0.5,0.5) node[below,xshift= -0.5em,yshift=1.8em,rotate=45] {\footnotesize $|z|$};
	\draw [blue] (0,1.27) node[below,xshift=- 0.51em,yshift=1em] {\footnotesize $y$};
	
	
	
	\draw [red] (0.2,0.2) node[xshift= 0em,yshift=-0.6em,rotate=0] {\footnotesize $\alpha$};
		%\draw[->,draw=green] (0,0) -- (18mm,0mm) arc (0:180:18mm) -- cycle; 
	\draw [blue] (-1.5,.2) node[below,xshift=- 0em,yshift=0em] { $\alpha=\pi$};
	\draw [blue] (1.5,.2) node[below,xshift=- 0em,yshift=0em] { $\alpha=0$};
	\end{tikzpicture}\\
\end{center}
\newpage
\noindent
\textbullet \;\; für $ \alpha ]\pi,2\pi[ $ gilt:\\

\begin{center}
	$ \alpha=2\pi-\arccos(\genfrac{}{}{}{2}{x }{|z| }) $ für $]\pi,2\pi[  $\\
	\begin{tikzpicture}[scale=2.5]
	\filldraw [->,draw=black,fill=red!20] (0,0)--(18mm,0mm) arc (0:315:18mm);
	% Gitternetzlinien
	\draw[xstep=1cm, ystep=1cm, very thin, color=lightgray] (-2,-2) grid (2,2);
	% x-Achse
	\draw[->] (-2,0) -- (2,0) node[below,xshift=-2em] {$\mathrm{Re}$}; 
	%Zahlen auf x-Achse
	\foreach \x in {-1,...,-1}
	\draw[shift={(\x,0)},color=black] (0pt,2pt) -- (0pt,-2pt) node[below] {\footnotesize $\x$};
	\foreach \x in {1,...,1}
	\draw[shift={(\x,0)},color=black] (0pt,2pt) -- (0pt,-2pt) node[below] {\footnotesize $\x$};
	% y-Achse 
	\draw[->] (0,-2) -- (0,2) node[left,yshift=-0.6em] {$\mathrm{Im}$};%node[above left]
	%Zahlen auf y-Achse
	\foreach \y in {-1,...,-1}
	\draw[shift={(0,\y)},color=black] (2pt,0pt) -- (-2pt,0pt) node[left] {\footnotesize $\y$};
	\foreach \y in {1,...,1}
	\draw[shift={(0,\y)},color=black] (2pt,0pt) -- (-2pt,0pt) node[left] {\footnotesize $\y$};
	%Ursprung
	%\draw[color=black] (0pt,-10pt) node[right] {\footnotesize $0$};
	\draw[color=black] (-4pt,-4pt) node[right] {\footnotesize $0$};
	%\draw[thick,color=black] plot[samples=200, domain=2:4] 
	%(\x,{2}) node[below, xshift=-5em] {$h(x)$};
	
	%\draw[thick,color=black] plot[samples=200, domain=-1:2] 
	%(\x,{\x-1}) node[below, xshift=1em]{$g(x)$};
	
	%\draw [blue] (0,0) circle (1.8);
	\filldraw[fill=green!20,draw=green!50!black] (0,0) -- (5mm,0mm) arc (0:-45:5mm) -- cycle; 
	\draw [blue,->] (0,0)-.(1.27,-1.27)node[below,xshift= 2em,yshift=0.8em] {\footnotesize $z=x+\mathrm{i}y$};
	\draw [blue,dashed] (0,-1.27)-.(1.27,-1.27);
	\draw [blue,dashed] (1.27,0)-.(1.27,-1.27);
	%\draw [blue] (1.27,0) node[below,xshift= 0em,yshift=-0.2em] {\footnotesize $x$};
	%\draw [blue] (-0.5,0.9) node[below,xshift= -0.5em,yshift=1.8em,rotate=45] {\footnotesize $|z|$};
	%\draw [blue] (0,1.27) node[below,xshift=- 0.51em,yshift=1em] {\footnotesize $y$};
	
	
	
	\draw [red] (0.2,-0) node[xshift= 0em,yshift=-0.6em,rotate=0] {\footnotesize $\beta$};
	%\draw[->,draw=green] (0,0) -- (18mm,0mm) arc (0:180:18mm) -- cycle; 
	\draw [blue] (-1,0.9) node[below,xshift=- 0em,yshift=0em] { $\alpha=2\pi-\beta$};
	%\draw [blue] (1.5,.2) node[below,xshift=- 0em,yshift=0em] { $\alpha=0$};
	\end{tikzpicture}\\
\end{center}
%	\chapter{Folgen \& Konvergenz}


\setcounter{section}{4}
\setcounter{subsection}{0}
\setcounter{subsubsection}{1}
\setcounter{secnumdepth}{3}
\subsection{Zahlenfolgen }
\href{https://youtu.be/WS8kO3O_byk}{Video 02}\\
Definition:\\
Eine reele bzw. komplexe (Zahlen) Folge ist eine Funktion:\\
\begin{center}
	\boxed{x:\mathbb{N}\longmapsto \mathbb{R} (\text{ bzw.  }  \mathbb{C}) \text{ mit } n\longmapsto x_n} 
\end{center}
\begin{center}
	\begin{tikzpicture}
	%\filldraw [red,opacity=0.5] (0,0)--(-1,2)--(-2,4)--(-4,4)--(-4,0)--(0,0);
	% Gitternetzlinien
	\draw[xstep=1cm, ystep=1cm, very thin, color=lightgray] (0,-3) grid (10,4);
	% x-Achse
	\draw[->] (0,0) -- (10,0) node[below,xshift=-1em] {$\mathbb{N}$}; 
	%Zahlen auf x-Achse
	%\foreach \x in {-3,...,-1}
	%\draw[shift={(\x,0)},color=black] (0pt,2pt) -- (0pt,-2pt) node[below] {\footnotesize $\x$};
	\foreach \x in {1,...,9}
	\draw[shift={(\x,0)},color=black] (0pt,2pt) -- (0pt,-2pt) node[below] {\footnotesize $\x$};
	% y-Achse 
	\draw[->] (0,-3) -- (0,4) node[left,yshift=-1em] {$\mathbb{R}$};%node[above left]
	%Zahlen auf y-Achse
	\foreach \y in {-3,...,3}
	\draw[shift={(0,\y)},color=black] (2pt,0pt) -- (-2pt,0pt) node[left] {\footnotesize $\y$};
	\foreach \y in {1,...,3}
	\draw[shift={(0,\y)},color=black] (2pt,0pt) -- (-2pt,0pt) node[left] {\footnotesize $\y$};
	%Ursprung
	%\draw[color=black] (-1.5pt,-2pt) node[below] {\footnotesize $0$};
	
	%\draw[thick,color=black] plot[samples=200, domain=-2:0] 
	%(\x,{-2*\x}) node[below, xshift=-6.5em, yshift=9em] {$y=-2x$};
	
	
	
	\fill [blue] (0,0.4) circle (.051cm) ;	
	
	\fill [blue] (1,0.5) circle (.051cm) ;
	\fill [blue] (2,2.3) circle (.051cm) ;
	\fill [blue] (3,1.3) circle (.051cm) ;
	\fill [blue] (5,2.1) circle (.051cm) ;
	\fill [blue] (6,3) circle (.051cm) ;
	\fill [blue] (7,-1.1) circle (.051cm) ;
	\fill [blue] (8,-2.2) circle (.051cm) ;
	\fill [blue] (9,-2.4) circle (.051cm) ;	
	\fill [blue] (4,-0.5) circle (.051cm) ;	
	\end{tikzpicture}\\
\end{center}
Stattdessen schreibt man:\\
\begin{center}
	\boxed{x_0,x_1,x_2\ldots = (x_n)_{n\in \mathbb{N}}}
\end{center}
Beispiel:\vspace*{2mm}\\
\textbullet \;\; $ (\frac{1}{2^n})_{n\in \mathbb{N}}\; =1,\frac{1}{2},\frac{1}{4},\frac{1}{8}\ldots$\\
\newpage
\noindent
\textbullet \;\;$ (n^2)_{n\in \mathbb{N}} \;=0,1,4,9\ldots$\\
\begin{center}
	
	
	\begin{tikzpicture} [scale=0.5]
	
	% Gitternetzlinien
	\draw[xstep=1cm, ystep=1cm, very thin, color=lightgray] (0,0) grid (10,10);
	% x-Achse
	\draw[->] (0,0) -- (10,0) node[below,xshift=-1em] {$\mathbb{N}$}; 
	%Zahlen auf x-Achse
	%\foreach \x in {-3,...,-1}
	%\draw[shift={(\x,0)},color=black] (0pt,2pt) -- (0pt,-2pt) node[below] {\footnotesize $\x$};
	\foreach \x in {0,...,8}
	\draw[shift={(\x,0)},color=black] (0pt,2pt) -- (0pt,-2pt) node[below] {\footnotesize $\x$};
	% y-Achse 
	\draw[->] (0,0) -- (0,10) node[left,yshift=-1em] {$\mathbb{R}$};%node[above left]
	%Zahlen auf y-Achse
	%\foreach \y in {-1,...,-1}
	%\draw[shift={(0,\y)},color=black] (2pt,0pt) -- (-2pt,0pt) node[left] {\footnotesize $\y$};
	\foreach \y in {1,...,8}
	\draw[shift={(0,\y)},color=black] (2pt,0pt) -- (-2pt,0pt) node[left] {\footnotesize $\y$};
	%Ursprung
	%\draw[color=black] (0pt,-10pt) node[right] {\footnotesize $0$};
	
	\draw[thick,dashed,color=red] plot[samples=200, domain=0:3] 
	(\x,{\x*\x}) node[below, xshift=4em,yshift=-6] {\tiny{Teilgraph einer Parabel}};
	
	%\draw[thick,color=black] plot[samples=200, domain=-1:2] 
	%(\x,{\x-1}) node[below, xshift=1em]{$g(x)$};
	\fill [blue] (0,0) circle (.1cm) ;	
	\fill [blue] (1,1) circle (.1) ;
	\fill [blue] (2,4) circle (.1) ;
	\fill [blue] (3,9) circle (.1) ;
	%\fill [blue] (3,3) circle (.1cm) ;
	
	
	\end{tikzpicture}\\
\end{center}
\textbullet \;\; $(\frac{1}{n})_{n\in \mathbb{N}\backslash \{0\}} \;=1,\frac{1}{2} ,\frac{1}{3}\ldots$\vspace*{2mm}\\
\textbullet \;\; $ (1)_{n\in \mathbb{N}}=1,1,1,\ldots$\vspace*{2mm}\\
\textbullet \;\; $ (\mathrm{i}^n)_{n\in \mathbb{N}}\;=1,\mathrm{i},\mathrm{i}^2,\mathrm{i}^3,\mathrm{i}^4\ldots =\mathrm{i},-1,-\mathrm{i},1\ldots$
\subsection{Konvergenz von Zahlenfolgen }
\href{https://youtu.be/KBmVpgrY57o}{Video 03}\\
Definition:\\
Eine Folge $ (x_n)_{n\in \mathbb{N}} $ heißt konvergent gegen $ a\in\mathbb{R} $, wenn es zu jeden $ \varepsilon>0 $ ein $ n_0 $ gibt, s. d.\\
\begin{align*}
\boxed{
	|x_n-a|<\varepsilon \;\; \text{ für alle}\;\; n\geq n_0}
\end{align*}
gilt.\vspace*{2mm}\\
Wir schreiben dann:\\
\begin{align*}
\boxed{\lim\limits_ {n\to\infty}x_n=a}
\end{align*}
Eine Folge $ (x_n)_{n\in \mathbb{N}} $ heißt konvergent, wenn es ein $ a\in\mathbb{R} $ gibt, gegen das sie konvergiert.\\
Andersfall heißt $ (x_n)_{n\in \mathbb{N}} $ divergent.\\
\begin{center}
	
	
	\begin{tikzpicture}
	
	% Gitternetzlinien
	\draw[xstep=1cm, ystep=1cm, very thin, color=lightgray] (0,0) grid (8,4);
	% x-Achse
	\draw[->] (0,0) -- (9,0) node[below,xshift=-1em] {$x$}; 
	%Zahlen auf x-Achse
	%\foreach \x in {-3,...,-1}
	%\draw[shift={(\x,0)},color=black] (0pt,2pt) -- (0pt,-2pt) node[below] {\footnotesize $\x$};
	\foreach \x in {0,...,3}
	\draw[shift={(\x,0)},color=black] (0pt,2pt) -- (0pt,-2pt) node[below] {\footnotesize $\x$};
	% y-Achse 
	\draw[->] (0,0) -- (0,4) node[left,yshift=-1em] {$\mathbb{R}$};%node[above left]
	%Zahlen auf y-Achse
	%\foreach \y in {-1,...,-1}
	%\draw[shift={(0,\y)},color=black] (2pt,0pt) -- (-2pt,0pt) node[left] {\footnotesize $\y$};
	%\foreach \y in {1,...,3}
	%\draw[shift={(0,\y)},color=black] (2pt,0pt) -- (-2pt,0pt) node[left] {\footnotesize $\y$};
	%Ursprung
	%\draw[color=black] (0pt,-10pt) node[right] {\footnotesize $0$};
	
	\draw[thick,color=red,dashed] plot[samples=200, domain=0:8] 
	(\x,{2.8}) node[left, xshift=-20em] {$a+\varepsilon$};
	\draw[thick,color=blue] plot[samples=200, domain=0:8] 
	(\x,{2}) node[left, xshift=-20em] {$a$};
	\draw[thick,color=red,dashed] plot[samples=200, domain=0:8] 
	(\x,{1.2}) node[left, xshift=-20em]{$a-\varepsilon$};
	
	\fill [blue] (0,0) circle (.05cm) ;	
	\fill [blue] (1,0.2) circle (.05) ;
	\fill [blue] (2,0.4) circle (.05) ;
	\fill [blue] (3,0.6) circle (.05) ;
	\fill [blue] (3,0.6) circle (.05) ;
	\fill [blue] (4,1.25) circle (.05) ;
	\fill [blue] (5,1.8) circle (.05) ;
	\fill [blue] (6,2.6) circle (0.05) ;
	\fill [blue] (7,1.6) circle (0.05) ;
	\draw[color=black] (4,-3pt) node[below] {\footnotesize $x_{n_0}$};
	\draw[color=black] (5,-3pt) node[below] {\footnotesize $x_{n_0+1}$};
	\draw[color=black] (7,-3pt) node[below] {\footnotesize $x_{n}$};
	\end{tikzpicture}\\
\end{center}
Ab den Wert $ n_0=4 $ sind alle Wert bis $ n $ im Bereich von $ a\pm \varepsilon $, dann konvergiert die Folge $ (x)_{n\in \mathbb{N}} $ gegen den Grenwert $ \lim\limits_ {n\to\infty}x_n=a $.\\
\subsubsection{Betrachtung der Folge $(\frac{1}{n})_{n\in \mathbb{N}\backslash\{0\}}  $}
\href{https://youtu.be/apfeLM-bE6A}{Video 04}\\
Betrachte die Folge $ (\frac{1}{n})_{n\in \mathbb{N}\backslash\{0\}}\;=1,\frac{1}{2},\frac{1}{3}\ldots $.\\
Wir wollen zeigen $ \lim\limits_ {n\to\infty}\frac{1}{n}=0 $.\vspace*{2mm}\\
Zu zeigen: Für jedes $ \varepsilon>0 $ gibt es ein $ n_0 \in \mathbb{N} $, s. d. $ |\frac{1}{n}-0|<\varepsilon $ für $ n\geq n_0 $.\\
Skizze:\\
\begin{center}
	\begin{tikzpicture}[yscale=3]
	
	% Gitternetzlinien
	\draw[xstep=1cm, ystep=0.3cm, very thin, color=lightgray] (0,-1) grid (8,1.2);
	% x-Achse
	\draw[->] (0,0) -- (8,0) node[below,xshift=-1em] {$\mathbb{N}$}; 
	%Zahlen auf x-Achse
	\foreach \x in {1,...,7}
	\draw[shift={(\x,0)},color=black] (0pt,2pt) -- (0pt,-2pt) node[below] {\footnotesize $\x$};
	%\foreach \x in {1,...,3}
	%\draw[shift={(\x,0)},color=black] (0pt,2pt) -- (0pt,-2pt) node[below] {\footnotesize $\x$};
	% y-Achse 
	\draw[->] (0,-1) -- (0,1.2) node[left,yshift=0em] {$\mathbb{R}$};%node[above left]
	%Zahlen auf y-Achse
	\foreach \y in {-1,...,-1}
	\draw[shift={(0,\y)},color=black] (2pt,0pt) -- (-2pt,0pt) node[left] {\footnotesize $\y$};
	\foreach \y in {1,...,1}
	\draw[shift={(0,\y)},color=black] (2pt,0pt) -- (-2pt,0pt) node[left] {\footnotesize $\y$};
	%Ursprung
	%\draw[color=black] (0pt,-10pt) node[right] {\footnotesize $0$};
	
	\draw[thick,color=red,dashed] plot[samples=200, domain=0:8] 
	(\x,{1/4}) node[left, xshift=-20em] {$\varepsilon=\genfrac{}{}{}{2}{1 }{4 }$};
	
	\draw[thick,color=red,dashed] plot[samples=200, domain=0:8] 
	(\x,{-1/4}) node[left, xshift=-20em]{$\varepsilon$};
	
	
	\draw[color=black] (1,1) node {\footnotesize \textbullet};
	\draw[color=black] (2,0.5) node {\footnotesize \textbullet};
	\draw[color=black] (3,0.33) node {\footnotesize \textbullet};
	\draw[color=black] (4,0.25) node {\footnotesize \textbullet};
	\draw[color=black] (5,0.2) node {\footnotesize \textbullet};
	\draw[color=black] (6,0.166) node {\footnotesize \textbullet};
	\draw[color=black] (7,0.142) node {\footnotesize \textbullet};
	\draw[color=black] (-1.5,0) node {\footnotesize Grenzwert=0};
	\draw[color=black] (5,-0.4) node {\footnotesize $ n_0=5 $};
	\end{tikzpicture}\\
\end{center}
Exemplarischen formalen Beweis von $ \lim\limits_ {n\to\infty}\frac{1}{n}=0 $.\\
Betrachte:\\
Sei $ \varepsilon>0 $ beliebig aber fest.
\begin{align}
&\bigg|\frac{1}{n}-0\bigg|=\bigg|\frac{1}{n}\bigg|=\frac{1}{n}<\varepsilon\\
&\Leftrightarrow \frac{1}{\varepsilon}<n
\end{align}
Wähle:\\
\begin{align*}
n_0=\bigg\lceil\frac{1}{\varepsilon}\bigg \rceil+1 \;\;\text{\tiny{aufgerundet $ \lceil0.7\rceil=1,\lceil2.1\rceil=3 $} }
\end{align*}
Dann gilt:\\
\begin{align*}
&n\geq n_0=\bigg\lceil\frac{1}{\varepsilon}\bigg\rceil+1\geq \frac{1}{\varepsilon}+1>\frac{1}{\varepsilon}\\
&\Leftrightarrow \frac{1}{\varepsilon} < n\\
&\Leftrightarrow\bigg|\frac{1}{n}-0\bigg|<\varepsilon_{\;\;_{\Square}} \;\;\;\;\text{wegen (IV.1), (IV.2)}\;\;\\
\end{align*}
\newpage
\noindent
\subsubsection{Beispiel für kovergente Folgen}
\href{https://youtu.be/qHlKFxy1ILw}{Video 06}\\
Zeige mit Hilfe der Definition der Konvergenz von Folgen, dass\\
\begin{align*}
\lim\limits_ {n\to\infty}\frac{1}{\sqrt[3]{n}}=0
\end{align*}
Zu zeigen:\\
Für $ \varepsilon>0 $ fest und beliebig exestiert ein $ n_0\in \mathbb{N}$, s.d. $\big |\frac{1}{\sqrt[3]{n}}\big|<\varepsilon $ für alle $ n>n_0 $\\
Betrachte:\\
\begin{align*}
&\bigg|\frac{1}{\sqrt[3]{n}}-0 \bigg |=\bigg |\frac{1}{\sqrt[3]{n}}\bigg |=\frac{1}{\sqrt[3]{n}}<\varepsilon\\
&\Leftrightarrow\frac{1}{\varepsilon}<\sqrt[3]{n}\\
&\Leftrightarrow \frac{1}{\varepsilon^3}<n\\
\end{align*}
Wähle:\\
\begin{align*}
n_0:=\bigg\lceil\frac{1}{\varepsilon^3}\bigg\rceil+1
\end{align*}
Dann gilt für alle:\\

\begin{align*}
n\geq n_0&\Leftrightarrow n\geq n_0=\bigg\lceil\frac{1}{\varepsilon^3}\bigg\rceil+1\geq\frac{1}{e^3}+1> \frac{1}{\varepsilon^3}_{_{\Square}}\\
& \Leftrightarrow \bigg |\frac{1}{\sqrt[3]{n}}-0 \bigg|<\varepsilon
\end{align*}
\subsection{Divergenz}
\href{https://youtu.be/r8IF-n0EyQg}{Video 07}\\
Eine Folge $(x_n)_{n\in\mathbb{N}}$ heißt divergent, wenn sie nicht konvergent ist.\vspace*{2mm}\\
Beispiel:\\
\textbullet \;\; $(n)_{n \in \mathbb{N}}=0,1,2,\ldots \lim\limits_ {n\to\infty}=\infty\;\;\;\;$\tiny{Bemerkung: $\infty\neq\mathbb{R}$} \\
\normalsize
\begin{center}
	
	
	\begin{tikzpicture} [scale=0.5]
	
	% Gitternetzlinien
	\draw[xstep=1cm, ystep=1cm, very thin, color=lightgray] (0,0) grid (10,10);
	% x-Achse
	\draw[->] (0,0) -- (10,0) node[below,xshift=-1em] {$\mathbb{N}$}; 
	%Zahlen auf x-Achse
	%\foreach \x in {-3,...,-1}
	%\draw[shift={(\x,0)},color=black] (0pt,2pt) -- (0pt,-2pt) node[below] {\footnotesize $\x$};
	\foreach \x in {0,...,8}
	\draw[shift={(\x,0)},color=black] (0pt,2pt) -- (0pt,-2pt) node[below] {\footnotesize $\x$};
	% y-Achse 
	\draw[->] (0,0) -- (0,10) node[left,yshift=-0em] {$\mathbb{R}$};%node[above left]
	%Zahlen auf y-Achse
	%\foreach \y in {-1,...,-1}
	%\draw[shift={(0,\y)},color=black] (2pt,0pt) -- (-2pt,0pt) node[left] {\footnotesize $\y$};
	\foreach \y in {1,...,8}
	\draw[shift={(0,\y)},color=black] (2pt,0pt) -- (-2pt,0pt) node[left] {\footnotesize $\y$};
	%Ursprung
	%\draw[color=black] (0pt,-10pt) node[right] {\footnotesize $0$};
	
	%\draw[thick,dashed,color=red] plot[samples=200, domain=0:3] 
	%(\x,{\x*\x}) node[below, xshift=4em,yshift=-6] {\tiny{Teilgraph einer Parabel}};
	
	%\draw[thick,color=black] plot[samples=200, domain=-1:2] 
	%(\x,{\x-1}) node[below, xshift=1em]{$g(x)$}; 
	\fill [blue] (0,0) circle (.1cm) ;	
	\fill [blue] (1,1) circle (.1) ;
	\fill [blue] (2,2) circle (.1) ;
	\fill [blue] (3,3) circle (.1) ;
	
	\fill [blue] (4,4) circle (.1cm) ;	
	\fill [blue] (5,5) circle (.1) ;
	\fill [blue] (6,6) circle (.1) ;
	\fill [blue] (7,7) circle (.1) ;
	\fill [blue] (8,8) circle (.1) ;
	\fill [blue] (9,9) circle (.1) ;
	
	\end{tikzpicture}\\
\end{center}
\textbullet \;\; $ \big((-1)^n\big)_{n\in \mathbb{N}} =1,-1,1,-1,\ldots \lim\limits_ {n\to\infty}\big((-1)^n\big)$ exestiert nicht!\\

\begin{center}
	
	
	\begin{tikzpicture} [scale=1]
	
	% Gitternetzlinien
	\draw[xstep=1cm, ystep=1cm, very thin, color=lightgray] (0,-3) grid (10,4);
	% x-Achse
	\draw[->] (0,0) -- (10,0) node[below,xshift=-1em] {$\mathbb{N}$}; 
	%Zahlen auf x-Achse
	%\foreach \x in {-3,...,-1}
	%\draw[shift={(\x,0)},color=black] (0pt,2pt) -- (0pt,-2pt) node[below] {\footnotesize $\x$};
	\foreach \x in {0,...,8}
	\draw[shift={(\x,0)},color=black] (0pt,2pt) -- (0pt,-2pt) node[below] {\footnotesize $\x$};
	% y-Achse 
	\draw[->] (0,-3) -- (0,4) node[left,yshift=-0.5em] {$\mathbb{R}$};%node[above left]
	%Zahlen auf y-Achse
	\foreach \y in {-2,...,-1}
	\draw[shift={(0,\y)},color=black] (2pt,0pt) -- (-2pt,0pt) node[left] {\footnotesize $\y$};
	\foreach \y in {1,...,3}
	\draw[shift={(0,\y)},color=black] (2pt,0pt) -- (-2pt,0pt) node[left] {\footnotesize $\y$};
	%Ursprung
	%\draw[color=black] (0pt,-10pt) node[right] {\footnotesize $0$};
	
	%\draw[thick,dashed,color=red] plot[samples=200, domain=0:3] 
	%(\x,{\x*\x}) node[below, xshift=4em,yshift=-6] {\tiny{Teilgraph einer Parabel}};
	
	%\draw[thick,color=black] plot[samples=200, domain=-1:2] 
	%(\x,{\x-1}) node[below, xshift=1em]{$g(x)$}; 
	\fill [blue] (0,1) circle (.1cm) ;	
	\fill [blue] (1,-1) circle (.1) ;
	\fill [blue] (2,1) circle (.1) ;
	\fill [blue] (3,-1) circle (.1) ;
	
	\fill [blue] (4,1) circle (.1cm) ;	
	\fill [blue] (5,-1) circle (.1) ;
	\fill [blue] (6,1) circle (.1) ;
	\fill [blue] (7,-1) circle (.1) ;
	\fill [blue] (8,1) circle (.1) ;
	\fill [blue] (9,-1) circle (.1) ;
	\draw[thick,color=red,dashed] plot[samples=200, domain=0:10] 
	(\x,{3/4}) node[left, xshift=-26em]{$\varepsilon$};
	\draw[thick,color=red,dashed] plot[samples=200, domain=0:10] 
	(\x,{5/4}) node[left, xshift=-26em]{$\varepsilon$};
	\end{tikzpicture}\\
\end{center}
Bei jedem $ \varepsilon $ muß gelten, dass alle Werte größer $ n_0 $ im $ \varepsilon $-Bereich sein mussen. Dies ist hier nicht gegeben. Diese Folge hat keinen Grenzwert.\vspace*{2mm}\\
\textbullet \;\; $ (-n)_{n\in \mathbb{N}} =0,-1,-2,\ldots \lim\limits_ {n\to\infty}=-\infty$\\
\begin{center}
	
	
	\begin{tikzpicture} [scale=0.5]
	
	% Gitternetzlinien
	\draw[xstep=1cm, ystep=1cm, very thin, color=lightgray] (0,-9) grid (10,2);
	% x-Achse
	\draw[->] (0,0) -- (10,0) node[below,xshift=-1em] {$\mathbb{N}$}; 
	%Zahlen auf x-Achse
	%\foreach \x in {-3,...,-1}
	%\draw[shift={(\x,0)},color=black] (0pt,2pt) -- (0pt,-2pt) node[below] {\footnotesize $\x$};
	\foreach \x in {0,...,8}
	\draw[shift={(\x,0)},color=black] (0pt,2pt) -- (0pt,-2pt) node[below] {\footnotesize $\x$};
	% y-Achse 
	\draw[->] (0,0) -- (0,2) node[left,yshift=0em] {$\mathbb{R}$};%node[above left]
	%Zahlen auf y-Achse
	\foreach \y in {-9,...,-1}
	\draw[shift={(0,\y)},color=black] (2pt,0pt) -- (-2pt,0pt) node[left] {\footnotesize $\y$};
	\foreach \y in {1,...,1}
	\draw[shift={(0,\y)},color=black] (2pt,0pt) -- (-2pt,0pt) node[left] {\footnotesize $\y$};
	%Ursprung
	%\draw[color=black] (0pt,-10pt) node[right] {\footnotesize $0$};
	
	%\draw[thick,dashed,color=red] plot[samples=200, domain=0:3] 
	%(\x,{\x*\x}) node[below, xshift=4em,yshift=-6] {\tiny{Teilgraph einer Parabel}};
	
	%\draw[thick,color=black] plot[samples=200, domain=-1:2] 
	%(\x,{\x-1}) node[below, xshift=1em]{$g(x)$}; 
	\fill [blue] (0,0) circle (.1cm) ;	
	\fill [blue] (1,-1) circle (.1) ;
	\fill [blue] (2,-2) circle (.1) ;
	\fill [blue] (3,-3) circle (.1) ;
	
	\fill [blue] (4,-4) circle (.1cm) ;	
	\fill [blue] (5,-5) circle (.1) ;
	\fill [blue] (6,-6) circle (.1) ;
	\fill [blue] (7,-7) circle (.1) ;
	\fill [blue] (8,-8) circle (.1) ;
	\fill [blue] (9,-9) circle (.1) ;
	
	\end{tikzpicture}\\
\end{center}
Definition: (Bestimmte Divergenz)\\
Eine Folge $ (x_n)_{n\in \mathbb{N}} $ ist bestimmt divergent gegen $ +\infty $ bzw. $ -\infty $, wenn zu jedem $ \varepsilon>0 $ ein $ n_0\in\mathbb{N} $ exestiert mit der Eigenschaft s.d. gilt:\\
.
\begin{center}
	\fbox{\parbox  {38mm} {
			$x_n>\varepsilon\; \;\;\;\;\text{für } n>n_0$\\
			$x_n<-\varepsilon \;\;\text{für } n>n_0$}}
\end{center}

\begin{center}
	
	
	\begin{tikzpicture} [scale=0.5]
	
	% Gitternetzlinien
	\draw[xstep=1cm, ystep=1cm, very thin, color=lightgray] (0,0) grid (10,10);
	% x-Achse
	\draw[->] (0,0) -- (10,0) node[below,xshift=-1em] {$\mathbb{N}$}; 
	%Zahlen auf x-Achse
	%\foreach \x in {-3,...,-1}
	%\draw[shift={(\x,0)},color=black] (0pt,2pt) -- (0pt,-2pt) node[below] {\footnotesize $\x$};
	\foreach \x in {0,...,8}
	\draw[shift={(\x,0)},color=black] (0pt,2pt) -- (0pt,-2pt) node[below] {\footnotesize $\x$};
	% y-Achse 
	\draw[->] (0,0) -- (0,10) node[left,yshift=-0em] {$\mathbb{R}$};%node[above left]
	%Zahlen auf y-Achse
	%\foreach \y in {-1,...,-1}
	%\draw[shift={(0,\y)},color=black] (2pt,0pt) -- (-2pt,0pt) node[left] {\footnotesize $\y$};
	\foreach \y in {1,...,8}
	\draw[shift={(0,\y)},color=black] (2pt,0pt) -- (-2pt,0pt) node[left] {\footnotesize $\y$};
	%Ursprung
	%\draw[color=black] (0pt,-10pt) node[right] {\footnotesize $0$};
	
	%\draw[thick,dashed,color=red] plot[samples=200, domain=0:3] 
	%(\x,{\x*\x}) node[below, xshift=4em,yshift=-6] {\tiny{Teilgraph einer Parabel}};
	
	%\draw[thick,color=black] plot[samples=200, domain=-1:2] 
	%(\x,{\x-1}) node[below, xshift=1em]{$g(x)$}; 
	\fill [blue] (0,0) circle (.1cm) ;	
	\fill [blue] (1,1) circle (.1) ;
	\fill [blue] (2,2) circle (.1) ;
	\fill [blue] (3,3) circle (.1) ;
	
	\fill [blue] (4,4) circle (.1cm) ;	
	\fill [blue] (5,5) circle (.1) ;
	\fill [blue] (6,6) circle (.1) ;
	\fill [blue] (7,7) circle (.1) ;
	\fill [blue] (8,8) circle (.1) ;
	\fill [blue] (9,9) circle (.1) ;
	\draw[thick,color=red,dashed] plot[samples=200, domain=0:10] 
	(\x,{5}) node[left, xshift=-13em]{$\varepsilon$};
	\draw[color=black] (7,7) node[right] {\footnotesize $(n)_{n\in\mathbb{N}}$};
	\end{tikzpicture}\\
\end{center}
\subsection{Rechenregeln für konvergente Folgen}
\href{https://youtu.be/m6GeFKpvFQ8}{Video 08}\\
\href{https://youtu.be/wWXpWP5EGCU}{Video 09}\\
Berechnung von Grenzwerten:\\
Beispiel:\\
\textbullet \;\; $ 0=\lim\limits_ {n\to\infty}\bigg(\frac{1}{n}\cdot \frac{1}{n}\bigg)=\lim\limits_ {n\to\infty}\bigg(\frac{1}{n}\bigg)\cdot\lim\limits_ {n\to\infty}\bigg(\frac{1}{n}\bigg)=0\cdot 0=0 $\\
\textbullet \;\; $ 1=\lim\limits_ {n\to\infty}1 =\lim\limits_ {n\to\infty}\bigg(n\cdot \frac{1}{n}\bigg)\overset{falsch}{=}\lim\limits_ {n\to\infty} n\cdot \lim\limits_ {n\to\infty}\bigg(\frac{1}{n}\bigg)=0\cdot\infty=0 $ \text{ falsch!!}\vspace*{4mm}\\
Nicht definiert und falsch ist:\\
\begin{align*}
&0\cdot\infty\\
&\frac{1}{\infty}=0\\
&\text{mit }\;\;\infty\;\; \text{kann man nicht rechnen! }
\end{align*}\vspace*{2mm}\\
\textbf{Rechenregeln für konvergente Folgen:}\\
Seien $ (x_n)_{n\in\mathbb{N}} $,$ (y_n)_{y\in \mathbb{N}} $ konvergente Folgen mit:\\
\begin{align*}
\lim\limits_ {n\to\infty}x_n=a \;\;\text{und }\;\; \lim\limits_ {n\to\infty}y_n=b\;\;\text{sowie }c\in \mathbb{R}
\end{align*}
dann gilt:\\
\begin{center}
	
	\fbox{\parbox  {98mm} {
			\begin{align*}
			1.&\lim\limits_ {n\to\infty}(x_n+y_n)=a+b\\
			2.&\lim\limits_ {n\to\infty}(x_n\cdot y_n)=a\cdot b\\
			3.&\lim\limits_ {n\to\infty}(c\cdot x_n)=c\cdot a\\
			4.&\lim\limits_ {n\to\infty}\bigg(\frac{x_n}{y_n}\bigg)=\frac{a}{b}\;\; \text{falls } b\neq 0\neq y_n\;\;\text{für alle }\;\;n\in \mathbb{N}\\
			5.&\; a\leq b\;\;\text{falls }\;\;x_n\leq y_n\;\;\text{für alle }\;\;n\in\mathbb{N}
			\end{align*}}}
\end{center}
\textbullet \;\; zu 1. :\\
\begin{align*}
\lim\limits_ {n\to\infty}\bigg(\frac{1}{n}+\frac{1}{n^2}\bigg)=\lim\limits_ {n\to\infty}\bigg(\frac{1}{n}\bigg)+\lim\limits_ {n\to\infty}\bigg(\frac{1}{n^2}\bigg)=0+0=0
\end{align*}
\textbullet \;\; zu 2. :\\
\begin{align*}
\lim\limits_ {n\to\infty}\bigg(\frac{1}{n}\cdot \frac{1}{n}\bigg)=\lim\limits_ {n\to\infty}\bigg(\frac{1}{n}\bigg)\cdot\lim\limits_ {n\to\infty}\bigg(\frac{1}{n}\bigg)=0\cdot 0=0\\
\end{align*}
\textbullet \;\; zu 3. :\\
\begin{align*}
\lim\limits_ {n\to\infty}\bigg(7\cdot\frac{1}{n^2}\bigg)=7\cdot \lim\limits_ {n\to\infty}\bigg(\frac{1}{n^2}\bigg)=7\cdot0 =0
\end{align*}
\textbullet \;\; zu 4. :\\
\begin{align*}
\lim\limits_ {n\to\infty}\Biggr(\frac{\lim\limits_ {n\to\infty}\frac{1}{n}}{\lim\limits_ {n\to\infty}\frac{1}{n^2}+1}\Bigg)=\Bigg(\frac{\lim\limits_ {n\to\infty}\frac{1}{n}}{\lim\limits_ {n\to\infty}\frac{1}{n^2}+\lim\limits_ {n\to\infty}1}\Bigg)=\frac{0}{0+1}=0
\end{align*}
\textbullet \;\; zu 5. :\\
\begin{align*}
\lim\limits_ {n\to\infty}\bigg|\frac{\cos(n)}{n}\bigg|\;\;\text{für }\;\;n>0
\end{align*}
Suche Folgen $ a,b,c $ für die gilt:\\
\begin{align*}
&a\leq b \leq c \;\;\text{falls }\;\; x_n\leq y_n \leq z_n \;\; \text{für alle }\;\;n\in \mathbb{N}\backslash {0}\\
&a=0\;\; \text{weil }\;\;0<\bigg|\frac{\cos(n)}{n}\bigg|\;\text{für }\;\;n>0\\
&b=\lim\limits_ {n\to\infty}\bigg|\frac{\cos(x)}{n}\bigg|\\
&c=\lim\limits_ {n\to\infty}\frac{1}{n}>\lim\limits_ {n\to\infty}\bigg|\frac{\cos(n)}{n}\bigg|
\end{align*}
$\bigg|\frac{\cos(n)}{n}\bigg|<\frac{1}{n}  $ weil:\\
\begin{align*}
0<|\cos(n)|<1
\end{align*}
\begin{center}
	\begin{tikzpicture}
	
	% Gitternetzlinien
	\draw[xstep=1cm, ystep=1cm, very thin, color=lightgray] (0,-2) grid (8,2);
	% x-Achse
	\draw[->] (0,0) -- (8,0) node[below,xshift=-1em] {$\mathbb{N}$}; 
	%Zahlen auf x-Achse
	\foreach \x in {1,...,7}
	\draw[shift={(\x,0)},color=black] (0pt,2pt) -- (0pt,-2pt) node[below] {\footnotesize $\x$};
	%\foreach \x in {1,...,3}
	%\draw[shift={(\x,0)},color=black] (0pt,2pt) -- (0pt,-2pt) node[below] {\footnotesize $\x$};
	% y-Achse 
	\draw[->] (0,-2) -- (0,2) node[left,yshift=-1em] {$\mathbb{R}$};%node[above left]
	%Zahlen auf y-Achse
	\foreach \y in {-1,...,-1}
	\draw[shift={(0,\y)},color=black] (2pt,0pt) -- (-2pt,0pt) node[left] {\footnotesize $\y$};
	\foreach \y in {1,...,1}
	\draw[shift={(0,\y)},color=black] (2pt,0pt) -- (-2pt,0pt) node[left] {\footnotesize $\y$};
	%Ursprung
	%\draw[color=black] (0pt,-10pt) node[right] {\footnotesize $0$};
	\draw[thick,color=red,dashed] plot[samples=200, domain=0:3.14/2] 
	(\x,{cos(\x r)});
	\draw[thick,color=red,dashed] plot[samples=200, domain=3.14/2:3.14] 
	(\x,{abs(cos(\x r))});
	\draw[thick,color=red,dashed] plot[samples=200, domain=3.14:(3.14+3.14/2)] 
	(\x,{abs(cos(\x r))});
	\draw[thick,color=red,dashed] plot[samples=200, domain=(3.14+3.14/2):(3.14+3.14/2+3.14)] 
	(\x,{cos(\x r)});
	%\draw[thick,color=black] plot[samples=200, domain=-1:2] 
	%(\x,{\x-1}) node[below, xshift=1em]{$g(x)$};
	\fill [blue] (1,0.54) circle (.1cm) ;	
	\fill [blue] (2,0.42) circle (.1) ;
	\fill [blue] (3,0.98) circle (.1) ;
	\fill [blue] (4,0.65) circle (.1) ;
	\fill [blue] (5,0.28) circle (.1) ;
	\fill [blue] (6,0.96) circle (.1) ;
	\fill [blue] (7,0.75) circle (.1) ;
	\draw[color=blue] (1,1.2) node[right] {\footnotesize $|\cos(n)|$};
	\end{tikzpicture}\\
\end{center}
dann gilt:\\
\begin{align*}
&\lim\limits_ {n\to\infty} a\leq \lim\limits_ {n\to\infty} b\leq \lim\limits_ {n\to\infty} c\\
&\lim\limits_ {n\to\infty}0\leq\lim\limits_ {n\to\infty}\bigg|\frac{\cos(n)}{n}\bigg|\leq\lim\limits_ {n\to\infty}\frac{1}{n}\\
&0\leq\lim\limits_ {n\to\infty}\bigg|\frac{\cos(n)}{n}\bigg|\leq0\\
&\Rightarrow\lim\limits_ {n\to\infty}\bigg|\frac{\cos(n)}{n}\bigg|=0
\end{align*}
\newpage
\noindent
\subsubsection{Tricks zur Berechnung von Grenzwerten I}
\href{https://youtu.be/4DUydSPIVLM}{Video 10}\\
Berechne (falls exestieren) die folgenden Grenzwerte:\\
\begin{align*}
(i)&\lim\limits_ {n\to\infty}\bigg(\frac{n^5+7n}{3n^5+7n^4}\bigg)\\
(ii)&\lim\limits_ {n\to\infty}\bigg(\frac{n^2+4n}{n^3-2n^2}-7n^4\bigg)
\end{align*}
zu(i)$ \;\;\lim\limits_ {n\to\infty}\big(\frac{n^5+7n}{3n^5+7n^4}\big)\\ $
\begin{quote}
	NR:\\
	Man schaut sich die höchsten n- Potenzen vom Zähler und Nenner an.\\
	\begin{align*}
	\lim\limits_ {n\to\infty}\frac{n^5}{3n^5}=\lim\limits_ {n\to\infty} \frac{1}{3}=\frac{1}{3}\\ 
	\end{align*}
	Dies ist nur eine Überlegung und kein Beweis!\\
	\begin{align*}
	\lim\limits_ {n\to\infty}\bigg(\frac{n^5+7n}{3n^5+7n^4}\bigg)=\lim\limits_ {n\to\infty}\bigg(\frac{n^5(1+\frac{7}{n^4})}{n^5(1+\frac{7}{n})}\bigg)=\frac{\lim\limits_ {n\to\infty} 1+\frac{7}{n^4}}{\lim\limits_ {n\to\infty}3+\frac{7}{n}}=\frac{1+0}{3+0}=\frac{1}{3}\\
	\end{align*}
\end{quote}
zu(ii)\;\;$\lim\limits_ {n\to\infty}\big(\frac{n^2+4n}{n^3-2n^2}-7n^4\big)  $\\
\begin{quote}
	Die Rechenregeln für konvergente Folgen kann hier nicht angewendet werden, weil $ -7n^4 $ divergent ist.\\
	Da der Grenzwert von $ \lim\limits_ {n\to\infty}-7n^4=-\infty $ ist, divergiert diese Folge.\\
	\begin{align*}
	\lim\limits_ {n\to\infty}\bigg(\frac{n^2+4n}{n^3-2n^2}-7n^4\bigg)=-\infty
	\end{align*}
	Auch die Überlegung mit den n-Potenzen ergibt:\\
	NR:\\
	\begin{align*}
	\frac{n^2+4n}{n^3-2n^2}-7n=\frac{n^2+4n-(n^3-2n^2)7n}{n^3-2n^2}=\frac{n^2+4-7n^4+14n^3}{n^3-2n^2}
	\end{align*}
	Die höchsten n-Potenzen sind:\\
	\begin{align*}
	\frac{-7n^4}{n^3}=-7n\\
	\lim\limits_ {n\to\infty}-7n=-\infty
	\end{align*}
\end{quote}
\subsubsection{Tricks zur Berechnung von Grenzwerten II}
\href{https://youtu.be/3wf0Q-vtFGo}{Video 11}\\
Berechne (falls exestieren) die folgenden Grenzwerte:\\
\begin{tabular}[t]{lll}% l=linksbündig r=rechtsbündig
	(i) & $\lim\limits_ {n\to\infty}\frac{n^2}{\sqrt{n^5-n}} \;\;\text{für }\;\; n>0   $ \vspace*{2mm} \\
	(ii) & $\lim\limits_ {n\to\infty}\frac{-1^n}{n}\;\;\text{für }\;\; n>0  $\vspace*{2mm}  \\	
	(iii)&$\lim\limits_ {n\to\infty}\frac{cos(n)}{n}\;\;\text{für }\;\; n>0  \vspace*{4mm}$
\end{tabular} \\
Zu (i):\\
\begin{quote}
	\begin{align*}
	\lim\limits_ {n\to\infty}\frac{n^2}{\sqrt{n^5-n}} \;\;\text{für }\;\; n>0 
	\end{align*}
	Der Grenzwert, mit den höchsten n-Potenzen ist $ \approx \lim\limits_ {n\to\infty}\frac{n^2}{\sqrt[2]{n^5}}$ \\
	\begin{align*}
	\lim\limits_ {n\to\infty}\frac{2}{\sqrt[2]{n^5}}=\frac{n^2}{n^{\frac{5}{2}}}=n^2 n^{-\frac{5}{2}}=n^{-\frac{1}{2}}=\frac{1}{\sqrt{n}}=0
	\end{align*}
	Mit Berechnung ergibt sich:\\
	\begin{align*}
	&\lim\limits_ {n\to\infty}\frac{n^2}{\sqrt{n^5-n}} =\frac{n^2}{\sqrt{n^5(1-\frac{1}{n^4})}}=\frac{n^2}{\sqrt{n^5}}\frac{1}{\sqrt{1-\frac{1}{n^4}}}=\frac{n^2}{n^{\frac{5}{2}}}
	\frac{1}{\sqrt{1-\frac{1}{n^4}}}=\lim\limits_ {n\to\infty}\frac{1}{\sqrt{n}}\frac{1}{\lim\limits_ {n\to\infty}\sqrt{1-\frac{1}{n^4}}}=0\cdot 1=0\\
	&\lim\limits_ {n\to\infty}\frac{n^2}{\sqrt{n^5-n}}=0
	\end{align*}
\end{quote}
zu (ii):\\
\begin{quote}
	\begin{align*}
	\lim\limits_ {n\to\infty}\frac{-1^n}{n}\;\;\text{für }\;\;n>0
	\end{align*}
	Der Grenzwert von$\lim\limits_ {n\to\infty}-1^n  $ exestiert nicht, deswegen sind die Rechenregeln nicht anwendbar.\\
	Es gilt:\\
	\begin{align*}
	\frac{-1}{n}\leq\frac{-1^n}{n}\leq \frac{1}{n}
	\end{align*}
	Ebenso gilt auch:\\
	\begin{align*}
	\lim\limits_ {n\to\infty}\frac{-1}{n}&\leq \lim\limits_ {n\to\infty}\frac{-1^n}{n}\leq \lim\limits_ {n\to\infty}\frac{1}{n}\\
	0&\leq\lim\limits_ {n\to\infty}\frac{-1^n}{n}\leq 0
	\end{align*}
	Daraus folgt nach dem Quetschlemma:\\
	\begin{align*}
	\lim\limits_ {n\to\infty}\frac{-1^n}{n}=0
	\end{align*}
\end{quote}
zu (iii):\\
\begin{quote}
	\begin{align*}
	\lim\limits_ {n\to\infty}\frac{cos(n)}{n}\;\;\text{für }\;\; n>0 
	\end{align*}
	Es gilt für $ cos(n) $:\\
	\begin{align*}
	-1<cos(n)<1
	\end{align*}
	Ebenso gilt:\\
	\begin{align*}
	\frac{-1}{n}<\frac{cos(n)}{n}<\frac{1}{n}\\
	\end{align*}
	Daraus folgt:\\
	\begin{align*}
	\lim\limits_ {n\to\infty} \frac{-1}{n}&<\lim\limits_ {n\to\infty}\frac{cos(n)}{n}<\lim\limits_ {n\to\infty}\frac{1}{n}\\
	0&<\lim\limits_ {n\to\infty}\frac{cos(x)}{n}<0
	\end{align*}
	Und nach dem Quetschlemma gilt:\\
	\begin{align*}
	\lim\limits_ {n\to\infty}\frac{cos(n)}{n}=0
	\end{align*}	
\end{quote}
\subsubsection{Tricks zur Berechnung von Grenzwerten III}
\href{https://youtu.be/KqTM73LzkMo}{Video 12}\\
Berechne (falls exestieren) die folgenden Grenzwerte:\\
\begin{align*}
\lim\limits_ {n\to\infty}\sqrt{n^2+n+2}-\sqrt{n^2+2}
\end{align*}
Mit Hilfe der 3. Binomischen Formel erzeugt man einen Bruch:\\
\begin{align*}
&a^2-b^2=(a-b)(a+b)\\
&a-b=\frac{(a-b)(a+b)}{(a+b)}=\frac{a^2-b^2}{(a+b)}
\end{align*}
\begin{align*}
\lim\limits_ {n\to\infty} \underbrace {\sqrt{n^2+n+2}}_{a}-\underbrace {\sqrt{n^2+2}}_{b}=\frac{(n^2+n+2)-(n^2+2)}{\sqrt{n^2+n+2}+\sqrt{n^2+2}}=\frac{n}{\sqrt{n^2+n+2}+\sqrt{n^2+2}}
\end{align*}
Der Grenzwert, mit den höchsten n-Potenzen ist $ \approx \lim\limits_ {n\to\infty}\frac{n}{\sqrt{n^2}+\sqrt{n^2}}$ \\
\begin{align*}
\lim\limits_ {n\to\infty}\frac{n}{\sqrt{n^2}+\sqrt{n^2}}=\frac{n}{n+n}=\frac{n}{2n}=\frac{1}{2}
\end{align*}
Mit Berechnung ergibt sich:\\
\begin{align*}
\lim\limits_ {n\to\infty}\frac{n}{\sqrt{n^2+n+2}+\sqrt{n^2+2}}&=\frac{n}{\sqrt{n^2(1+\frac{1}{n}+\frac{2}{n^2})}+\sqrt{n^2(1+\frac{2}{n^2})}}\\
&=\frac{n}{n\bigg(\sqrt{1+\frac{1}{n}+\frac{2}{n^2}}\bigg)+\bigg(\sqrt{1+\frac{2}{n^2}}\bigg)}\\
&=\frac{1}{\sqrt{1+\frac{1}{n}+\frac{2}{n^2}}+\sqrt{1+\frac{2}{n^2}}}\\
&=\lim\limits_ {n\to\infty}\frac{1}{\sqrt{1+\frac{1}{n}+\frac{2}{n^2}}+\sqrt{1+\frac{2}{n^2}}}=\frac{1}{\sqrt{1+0+0}+\sqrt{1+0}}=\frac{1}{\sqrt{1}+\sqrt{1}}=\frac{1}{2}\\
\end{align*}
Dann ist der Grenzwert von:\\
\begin{align*}
\lim\limits_ {n\to\infty}\sqrt{n^2+n+2}-\sqrt{n^2+2}=\frac{1}{2}
\end{align*}
\subsection{Monotonie von Folgen}
\href{https://youtu.be/W9Xl_hHf18I}{Video 13}\\
Definition:\\
Eine reelle Folge $ (x)_{n\in\mathbb{N}} $ heißt streng monoton wachsend wenn gilt:\\
\begin{align*}
x_n< x_n+1\;\;\text{für alle }\;\; n\in \mathbb{N}\\  
\end{align*}
Eine reelle Folge $ (x)_{n\in\mathbb{N}} $ heißt streng monoton fallend wenn gilt:\\
\begin{align*}
x_n>x_n+1\;\;\text{für alle }\;\; n\in \mathbb{N}\\  
\end{align*}
Eine reelle Folge $ (x)_{n\in\mathbb{N}} $ heißt monoton wachsend wenn gilt:\\
\begin{align*}
x_n\leq x_n+1\;\;\text{für alle }\;\; n\in \mathbb{N}\\  
\end{align*}
Eine reelle Folge $ (x)_{n\in\mathbb{N}} $ heißt  monoton fallend wenn gilt:\\
\begin{align*}
x_n\geq x_n+1\;\;\text{für alle }\;\; n\in \mathbb{N}\\  
\end{align*}
\newpage
\noindent
\begin{center}
	Streng monoton wachsende Folge $ (x)_{n\in\mathbb{N}}$\\
	\begin{tikzpicture}
	%\filldraw [red,opacity=0.5] (0,0)--(-1,2)--(-2,4)--(-4,4)--(-4,0)--(0,0);
	% Gitternetzlinien
	\draw[xstep=1cm, ystep=1cm, very thin, color=lightgray] (0,-3) grid (10,4);
	% x-Achse
	\draw[->] (0,0) -- (10,0) node[below,xshift=-1em] {$\mathbb{N}$}; 
	%Zahlen auf x-Achse
	%\foreach \x in {-3,...,-1}
	%\draw[shift={(\x,0)},color=black] (0pt,2pt) -- (0pt,-2pt) node[below] {\footnotesize $\x$};
	\foreach \x in {1,...,9}
	\draw[shift={(\x,0)},color=black] (0pt,2pt) -- (0pt,-2pt) node[below] {\footnotesize $\x$};
	% y-Achse 
	\draw[->] (0,-3) -- (0,4) node[left,yshift=-1em] {$\mathbb{R}$};%node[above left]
	%Zahlen auf y-Achse
	\foreach \y in {-3,...,3}
	\draw[shift={(0,\y)},color=black] (2pt,0pt) -- (-2pt,0pt) node[left] {\footnotesize $\y$};
	\foreach \y in {1,...,3}
	\draw[shift={(0,\y)},color=black] (2pt,0pt) -- (-2pt,0pt) node[left] {\footnotesize $\y$};
	%Ursprung
	%\draw[color=black] (-1.5pt,-2pt) node[below] {\footnotesize $0$};
	
	%\draw[thick,color=black] plot[samples=200, domain=-2:0] 
	%(\x,{-2*\x}) node[below, xshift=-6.5em, yshift=9em] {$y=-2x$};
	
	
	
	\fill [blue] (0,-2.5) circle (.051cm) ;	
	
	\fill [blue] (1,-2.1) circle (.051cm) ;
	\fill [blue] (2,-2) circle (.051cm) ;
	\fill [blue] (3,-1) circle (.051cm) ;
	\fill [blue] (4,-0.2) circle (.051cm) ;
	\fill [blue] (5,0.5) circle (.051cm) ;
	\fill [blue] (6,1.7) circle (.051cm) ;
	\fill [blue] (7,2) circle (.051cm) ;
	\fill [blue] (8,3) circle (.051cm) ;
	\fill [blue] (9,3.2) circle (.051cm) ;	
	%\fill [blue] (8,3) node{\footnotesize streng monoton wachsend} ;	
	\end{tikzpicture}\\
\end{center}
\begin{center}
	Streng monoton fallende Folge $ (x)_{n\in\mathbb{N}}$\\
	\begin{tikzpicture}
	%\filldraw [red,opacity=0.5] (0,0)--(-1,2)--(-2,4)--(-4,4)--(-4,0)--(0,0);
	% Gitternetzlinien
	\draw[xstep=1cm, ystep=1cm, very thin, color=lightgray] (0,-3) grid (10,4);
	% x-Achse
	\draw[->] (0,0) -- (10,0) node[below,xshift=-1em] {$\mathbb{N}$}; 
	%Zahlen auf x-Achse
	%\foreach \x in {-3,...,-1}
	%\draw[shift={(\x,0)},color=black] (0pt,2pt) -- (0pt,-2pt) node[below] {\footnotesize $\x$};
	\foreach \x in {1,...,9}
	\draw[shift={(\x,0)},color=black] (0pt,2pt) -- (0pt,-2pt) node[below] {\footnotesize $\x$};
	% y-Achse 
	\draw[->] (0,-3) -- (0,4) node[left,yshift=-1em] {$\mathbb{R}$};%node[above left]
	%Zahlen auf y-Achse
	\foreach \y in {-3,...,3}
	\draw[shift={(0,\y)},color=black] (2pt,0pt) -- (-2pt,0pt) node[left] {\footnotesize $\y$};
	\foreach \y in {1,...,3}
	\draw[shift={(0,\y)},color=black] (2pt,0pt) -- (-2pt,0pt) node[left] {\footnotesize $\y$};
	%Ursprung
	%\draw[color=black] (-1.5pt,-2pt) node[below] {\footnotesize $0$};
	
	%\draw[thick,color=black] plot[samples=200, domain=-2:0] 
	%(\x,{-2*\x}) node[below, xshift=-6.5em, yshift=9em] {$y=-2x$};
	
	
	
	\fill [blue] (0,3.5) circle (.051cm) ;	
	
	\fill [blue] (1,3) circle (.051cm) ;
	\fill [blue] (2,2) circle (.051cm) ;
	\fill [blue] (3,1.5) circle (.051cm) ;
	\fill [blue] (4,-0.2) circle (.051cm) ;
	\fill [blue] (5,-0.5) circle (.051cm) ;
	\fill [blue] (6,-1.7) circle (.051cm) ;
	\fill [blue] (7,-2) circle (.051cm) ;
	\fill [blue] (8,-2.6) circle (.051cm) ;
	\fill [blue] (9,-2.8) circle (.051cm) ;	
	%\fill [blue] (8,3) node{\footnotesize streng monoton wachsend} ;	
	\end{tikzpicture}\\
\end{center}
\begin{center}
	Monoton fallende Folge $ (x)_{n\in\mathbb{N}}$\\
	\begin{tikzpicture}
	%\filldraw [red,opacity=0.5] (0,0)--(-1,2)--(-2,4)--(-4,4)--(-4,0)--(0,0);
	% Gitternetzlinien
	\draw[xstep=1cm, ystep=1cm, very thin, color=lightgray] (0,-3) grid (10,4);
	% x-Achse
	\draw[->] (0,0) -- (10,0) node[below,xshift=-1em] {$\mathbb{N}$}; 
	%Zahlen auf x-Achse
	%\foreach \x in {-3,...,-1}
	%\draw[shift={(\x,0)},color=black] (0pt,2pt) -- (0pt,-2pt) node[below] {\footnotesize $\x$};
	\foreach \x in {1,...,9}
	\draw[shift={(\x,0)},color=black] (0pt,2pt) -- (0pt,-2pt) node[below] {\footnotesize $\x$};
	% y-Achse 
	\draw[->] (0,-3) -- (0,4) node[left,yshift=-1em] {$\mathbb{R}$};%node[above left]
	%Zahlen auf y-Achse
	\foreach \y in {-3,...,3}
	\draw[shift={(0,\y)},color=black] (2pt,0pt) -- (-2pt,0pt) node[left] {\footnotesize $\y$};
	\foreach \y in {1,...,3}
	\draw[shift={(0,\y)},color=black] (2pt,0pt) -- (-2pt,0pt) node[left] {\footnotesize $\y$};
	%Ursprung
	%\draw[color=black] (-1.5pt,-2pt) node[below] {\footnotesize $0$};
	
	%\draw[thick,color=black] plot[samples=200, domain=-2:0] 
	%(\x,{-2*\x}) node[below, xshift=-6.5em, yshift=9em] {$y=-2x$};
	
	
	
	\fill [blue] (0,3.5) circle (.051cm) ;	
	
	\fill [blue] (1,3) circle (.051cm) ;
	\fill [blue] (2,2) circle (.051cm) ;
	\fill [blue] (3,1.5) circle (.051cm) ;
	\fill [blue] (4,-1) circle (.051cm) ;
	\fill [blue] (5,-1) circle (.051cm) ;
	\fill [blue] (6,-1) circle (.051cm) ;
	\fill [blue] (7,-1) circle (.051cm) ;
	\fill [blue] (8,-1.6) circle (.051cm) ;
	\fill [blue] (9,-2.8) circle (.051cm) ;	
	%\fill [blue] (8,3) node{\footnotesize streng monoton wachsend} ;	
	\end{tikzpicture}\\
\end{center}
\newpage
\noindent
\begin{center}
	Monoton wachsende Folge $ (x)_{n\in\mathbb{N}}$\\
	\begin{tikzpicture}
	%\filldraw [red,opacity=0.5] (0,0)--(-1,2)--(-2,4)--(-4,4)--(-4,0)--(0,0);
	% Gitternetzlinien
	\draw[xstep=1cm, ystep=1cm, very thin, color=lightgray] (0,-3) grid (10,4);
	% x-Achse
	\draw[->] (0,0) -- (10,0) node[below,xshift=-1em] {$\mathbb{N}$}; 
	%Zahlen auf x-Achse
	%\foreach \x in {-3,...,-1}
	%\draw[shift={(\x,0)},color=black] (0pt,2pt) -- (0pt,-2pt) node[below] {\footnotesize $\x$};
	\foreach \x in {1,...,9}
	\draw[shift={(\x,0)},color=black] (0pt,2pt) -- (0pt,-2pt) node[below] {\footnotesize $\x$};
	% y-Achse 
	\draw[->] (0,-3) -- (0,4) node[left,yshift=-1em] {$\mathbb{R}$};%node[above left]
	%Zahlen auf y-Achse
	\foreach \y in {-3,...,3}
	\draw[shift={(0,\y)},color=black] (2pt,0pt) -- (-2pt,0pt) node[left] {\footnotesize $\y$};
	\foreach \y in {1,...,3}
	\draw[shift={(0,\y)},color=black] (2pt,0pt) -- (-2pt,0pt) node[left] {\footnotesize $\y$};
	%Ursprung
	%\draw[color=black] (-1.5pt,-2pt) node[below] {\footnotesize $0$};
	
	%\draw[thick,color=black] plot[samples=200, domain=-2:0] 
	%(\x,{-2*\x}) node[below, xshift=-6.5em, yshift=9em] {$y=-2x$};
	
	
	
	\fill [blue] (0,-2.5) circle (.051cm) ;	
	
	\fill [blue] (1,-2) circle (.051cm) ;
	\fill [blue] (2,-1.5) circle (.051cm) ;
	\fill [blue] (3,.5) circle (.051cm) ;
	\fill [blue] (4,1) circle (.051cm) ;
	\fill [blue] (5,1) circle (.051cm) ;
	\fill [blue] (6,1) circle (.051cm) ;
	\fill [blue] (7,2) circle (.051cm) ;
	\fill [blue] (8,2.6) circle (.051cm) ;
	\fill [blue] (9,3) circle (.051cm) ;	
	%\fill [blue] (8,3) node{\footnotesize streng monoton wachsend} ;	
	\end{tikzpicture}\\
\end{center}
\begin{center}
	Monoton wachsende oder fallende Folge $ (1)_{n\in\mathbb{N}}$\\
	\begin{tikzpicture}
	%\filldraw [red,opacity=0.5] (0,0)--(-1,2)--(-2,4)--(-4,4)--(-4,0)--(0,0);
	% Gitternetzlinien
	\draw[xstep=1cm, ystep=1cm, very thin, color=lightgray] (0,-3) grid (10,4);
	% x-Achse
	\draw[->] (0,0) -- (10,0) node[below,xshift=-1em] {$\mathbb{N}$}; 
	%Zahlen auf x-Achse
	%\foreach \x in {-3,...,-1}
	%\draw[shift={(\x,0)},color=black] (0pt,2pt) -- (0pt,-2pt) node[below] {\footnotesize $\x$};
	\foreach \x in {1,...,9}
	\draw[shift={(\x,0)},color=black] (0pt,2pt) -- (0pt,-2pt) node[below] {\footnotesize $\x$};
	% y-Achse 
	\draw[->] (0,-3) -- (0,4) node[left,yshift=-1em] {$\mathbb{R}$};%node[above left]
	%Zahlen auf y-Achse
	\foreach \y in {-3,...,3}
	\draw[shift={(0,\y)},color=black] (2pt,0pt) -- (-2pt,0pt) node[left] {\footnotesize $\y$};
	\foreach \y in {1,...,3}
	\draw[shift={(0,\y)},color=black] (2pt,0pt) -- (-2pt,0pt) node[left] {\footnotesize $\y$};
	%Ursprung
	%\draw[color=black] (-1.5pt,-2pt) node[below] {\footnotesize $0$};
	
	%\draw[thick,color=black] plot[samples=200, domain=-2:0] 
	%(\x,{-2*\x}) node[below, xshift=-6.5em, yshift=9em] {$y=-2x$};
	
	
	
	\fill [blue] (0,-2.5) circle (.051cm) ;	
	
	\fill [blue] (1,1) circle (.051cm) ;
	\fill [blue] (2,1) circle (.051cm) ;
	\fill [blue] (3,1) circle (.051cm) ;
	\fill [blue] (4,1) circle (.051cm) ;
	\fill [blue] (5,1) circle (.051cm) ;
	\fill [blue] (6,1) circle (.051cm) ;
	\fill [blue] (7,1) circle (.051cm) ;
	\fill [blue] (8,1) circle (.051cm) ;
	\fill [blue] (9,1) circle (.051cm) ;	
	%\fill [blue] (8,3) node{\footnotesize streng monoton wachsend} ;	
	\end{tikzpicture}\\
\end{center}
Beispiel:\\
\textbullet \;\; Die Folge $ (n)_{n\in\mathbb{N}} $ ist streng monoton wachsend, weil:\\
\begin{align*}
x_n=n< x_{n+1}=n+1 \;\;\text{für alle }\;\;n\in\mathbb{N}
\end{align*}
\textbullet \;\;Die Folge $ (\frac{1}{n})_{n\in\mathbb{N}\backslash\{0\}} $ ist streng monoton fallend, weil:\\
\begin{align*}
x_n=\frac{1}{n}>x_{n+1}=\frac{1}{n+1}\;\;\text{für alle } n\in\mathbb{N}\backslash\{0\}
\end{align*}
\textbullet \;\; Die konstante Folge $ (a)_{n\in \mathbb{N}} $ mit  $ a\in\mathbb{R} $ ist sowohl monoton wachsend als auch monoton fallend.\\
\subsubsection{Aufgabe zur Monotonie von Folgen}
\href{https://youtu.be/WSqcscE6Dhg}{Video 14}\\
Zeige, dass dir rekrusive Folge gegeben durch $a_0=1,\; a_{n+1}=a_n+\frac{1}{a_n} $ streng monoton wachsend ist.\vspace*{2mm}\\
Bemerkung:\\
Bei einer rekrusiven Folge muss man die Monotonie mit Hilfe der vollständigen Induktion beweisen.\\
Die Summanden von der rekrusiven Folge müssen größer Null sein, dass die streng wachsende Monotonie möglich ist.\\
Also zeige:\\
\begin{align*}
a_n>0 \;\;\text{für alle }\;\;n_0\in \mathbb{N}\Rightarrow \frac{1}{a_n}>0
\end{align*}
Damit hätte man:\\
\begin{align*}
&a_{n+1}=a_n+\frac{1}{a_n}>a_n+0=a_n \;\;\text{für alle }\;\;n\in \mathbb{N}\\
&a_{n+1}> a_n
\end{align*}
Beweis von $ a_n>0 $ für alle $ n\in\mathbb{N} $ mittels vollständiger Induktion:\\
\textbf{\textbullet \;\; (I.A.)}:\\
\begin{quote}
	für $ n=0 $ gilt:\\
	\begin{align*}
	a_n=a_0=1>0_{\;\;_{\Square}}
	\end{align*}
\end{quote}
\textbf{\textbullet \;\; (I.V.)}:\\
\begin{quote}
	für ein $ n $ gilt:\\
	\begin{align*}
	a_n>0 \Rightarrow \frac{1}{a_n}>0
	\end{align*}
\end{quote}
\textbf{\textbullet \;\;(I.S.)}:\\
\begin{quote}
	zu zeigen:\\
	\begin{align*}
	a_{n+1}>0
	\end{align*}
	betrachte:\\
	\begin{align*}
	a_{n+1}>a_n+\frac{1}{a_n}\overset{(I.V.)}{>}0+ 0=0_{\;\;_{\Square}}
	\end{align*}
\end{quote}
\newpage
\noindent
\subsection{Beschränktheit von Folgen }
\href{https://youtu.be/ftHUnCyM89E}{Video 15}\\
Die reelle Folge $ (x)_{n\in\mathbb{N}} $ heißt beschränkt, wenn es $ a,b \in \mathbb{R} $ gibt mit $ a\leq x_n\leq b $ für alle $ n\in \mathbb{N} $\\
\begin{center}
	\begin{tikzpicture}[yscale=3]
	
	% Gitternetzlinien
	\draw[xstep=1cm, ystep=0.3cm, very thin, color=lightgray] (0,-1) grid (8,1.2);
	% x-Achse
	\draw[->] (0,0) -- (8,0) node[below,xshift=-1em] {$\mathbb{N}$}; 
	%Zahlen auf x-Achse
	\foreach \x in {1,...,7}
	\draw[shift={(\x,0)},color=black] (0pt,2pt) -- (0pt,-2pt) node[below] {\footnotesize $\x$};
	%\foreach \x in {1,...,3}
	%\draw[shift={(\x,0)},color=black] (0pt,2pt) -- (0pt,-2pt) node[below] {\footnotesize $\x$};
	% y-Achse 
	\draw[->] (0,-1) -- (0,1.2) node[left,yshift=0em] {$\mathbb{R}$};%node[above left]
	%Zahlen auf y-Achse
	\foreach \y in {-1,...,-1}
	\draw[shift={(0,\y)},color=black] (2pt,0pt) -- (-2pt,0pt) node[left] {\footnotesize $\y$};
	\foreach \y in {1,...,1}
	\draw[shift={(0,\y)},color=black] (2pt,0pt) -- (-2pt,0pt) node[left] {\footnotesize $\y$};
	%Ursprung
	%\draw[color=black] (0pt,-10pt) node[right] {\footnotesize $0$};
	
	\draw[thick,color=red,dashed] plot[samples=200, domain=0:8] 
	(\x,{1/2}) node[left, xshift=-20em] {$a$};
	
	\draw[thick,color=red,dashed] plot[samples=200, domain=0:8] 
	(\x,{-1.5/4}) node[left, xshift=-20em]{$b$};
	\draw[color=black] (0,0.1) node {\footnotesize \textbullet};
	\draw[color=blue,thick] (0,-1.5/4) --(0,0.5);
	\draw[color=blue,thick] (0,-1.5/4) node[rotate=90] {\footnotesize $[$};
	\draw[color=blue,thick] (0,0.5) node[rotate=90] {\textbf{$]$}};
	\draw[color=black] (1,0.3) node {\footnotesize \textbullet};
	\draw[color=black] (2,0.5) node {\footnotesize \textbullet};
	\draw[color=black] (3,0.4) node {\footnotesize \textbullet};
	\draw[color=black] (4,0.25) node {\footnotesize \textbullet};
	\draw[color=black] (5,0.2) node {\footnotesize \textbullet};
	\draw[color=black] (6,0.15) node {\footnotesize \textbullet};
	\draw[color=black] (7,-1.5/4) node {\footnotesize \textbullet};
	\draw[color=black] (8,-0.29) node {\footnotesize \textbullet};
	
	%\draw[color=black] (-1.5,0) node {\footnotesize Grenzwert=0};
	%\draw[color=black] (5,-0.4) node {\footnotesize $ n_0=5 $};
	\end{tikzpicture}\\
\end{center}
\begin{align*}
\Leftrightarrow x_n\in [a,b]\;\;\text{für }\;\; n\in \mathbb{N}
\end{align*}
Beispiel:\\
\textbullet \;\; Die Folge $ \big((-1)^n\big)_{n\in\mathbb{N}} $ ist beschränkt weil:\\
\begin{quote}
	\begin{align*}
	-1\leq(-1)^n\leq 1\;\;\text{für alle }\;\; n\in\mathbb{N}\\
	-20\leq(-1)^n\leq 20\;\;\text{für alle }\;\; n\in\mathbb{N}\\
	(-1)^n \in [-200,15]\\
	\end{align*}
	Kleinste obere Schranke oder obere Schanke, bzw. größte untere Schranke oder untere Schranke. Wichtig ist eigentlich ist die Folge beschränkt oder nicht.\\
	Bemerkung:\\
	\begin{align*}
	(-1)^n \;\;\text{ist nicht konvergent, aber beschränkt ! }
	\end{align*}
\end{quote}
\textbullet \;\;Die Folge $ (\frac{1}{n}) $ ist beschränkt weil:\\
\begin{quote}
	\begin{align*}
	1\geq \frac{1}{n}\geq 0 \;\;\text{für alle }\;\;n\in\mathbb{N}_{\backslash\{0\}}
	\end{align*}
\end{quote}
\textbullet \;\; Die Folge $ (n)_{n\in\mathbb{N}} $ ist unbeschränkt, weil:\\
\begin{quote}
	Es gibt eine untere Schranke, aber keine obere Schranke.\\
	\begin{align*}
	x_n=n\in[0,\infty[
	\end{align*}
\end{quote}
\subsubsection{Aufgabe zur Beschränktheit von Folgen }
\href{https://youtu.be/w8bu4mqcZrQ}{Video 16}\\
Zeige, dass die rekrusive Folge, gegeben durch $ a_0=0 $ und $ a_{n+1}=\frac{1}{3}(a_n+1)$ für alle $ n\in \mathbb{N} $ durch $ 0 $ und $ 1 $ beschränkt ist.\vspace*{2mm}\\
Zu zeigen:\\
\begin{align*}
0\leq a_n < 1 \;\;\text{für alle }\;\;n\in \mathbb{N}
\end{align*}
Beweis mittels vollständiger Induktion.\\
\textbf{\textbullet \;\; (I.A.):}\\
\begin{quote}
	für $ n=0 $ gilt:\\
	\begin{align*}
	a_n=0=a_0=0\;\;\in[0,1 [_{\;\;_{\Square}}
	\end{align*}
	
\end{quote}
\textbf{\textbullet \;\; (I.V.):}\\
\begin{quote}
	für ein $ n\in\mathbb{N} $ gilt:\\
	\begin{align*}
	0\leq a_n < 1
	\end{align*}
\end{quote}
\textbf{\textbullet \;\; (I.S.):}\\
\begin{quote}
	zu zeigen:\\
	\begin{align*}
	0\leq a_{n+1} < 1
	\end{align*}
	Betrachte für $ a_n<1 $\\
	\begin{align*}
	a_{n+1}=\frac{1}{3}(a_n+1)=\frac{1}{3}a_n+\frac{1}{3}\overset{(I.V.)}{<}\frac{1}{3}\cdot1 +\frac{1}{3}=\frac{2}{3}<1_{\;\;_{\Square}}\\
	\end{align*}
	Betrachte für $ a_n\geq 0 $\\
	\begin{align*}
	&a_{n+1}=\frac{1}{3}(a_n+1)=\frac{1}{3}a_n+\frac{1}{3}\overset{(I.V.)}{>}\frac{1}{3}\cdot 0 +\frac{1}{3}>0_{\;\;_{\Square}}\\
	&a_{n+1}\geq 0 \;\;\text{weil $ a_{n_{0}}=0=a_0 $, deswegen: $ \geq $ }
	\end{align*}
	\begin{align*}
	\Rightarrow 0\leq a_{n+1} < 1
	\end{align*}
\end{quote}
\subsection{Konvergenzkriterien }
\href{https://youtu.be/KU9z0fJVi-U}{Video 17}\\
Konvergenzkriterium:\\
\begin{center}
	
	
	\fbox{\parbox  {120mm} {
			Jede monotone und beschränkte Folge reeller Zahlen ist konvergent.}}
\end{center}
Beispiel:\\
\textbullet \;\; Die Folge $ (\frac{1}{n})_{n\in\mathbb{N}\backslash\{0\}} $:\\
\begin{quote}
	(i) Monotonie:\\
	\begin{quote}
		Die Folge ist streng monoton fallend, da:\\
		\begin{align*}
		x_n=\frac{1}{n}>\frac{1}{n+1}=x_{n+1} \;\;\text{für alle}\;\;n\in\mathbb{N}{\backslash\{0\}}
		\end{align*}
	\end{quote}
	(ii) Beschränktheit:\\
	\begin{quote}
		Die Folge ist beschränkt, da:\\
		\begin{align*}
		0<\frac{1}{n}\leq 1 \;\;\text{für alle }n\in\mathbb{N}{\backslash\{0\}}
		\end{align*}
	\end{quote}
	Weil die Folge monoton und beschränkt ist, ist die Folge konvergent.\\
	Bemerkung:\\
	Die Folge ist immer größer als Null, aber der Grenzwert $ \lim\limits_ {n\to\infty} \frac{1}{n}=0 $.
	
\end{quote}
\textbullet \;\; Die Folge $ \bigg(\frac{(-1)^n}{n}\bigg)_{n\in\mathbb{N}\backslash\{0\}} $\\
\begin{quote}
	(i) Monotonie\\
	\begin{quote}
		Die Folge ist nicht monoton, da:\\
		\begin{align*}
		&\bigg(\frac{(-1)^n}{n}\bigg)_{n\in\mathbb{N}\backslash\{0\}}=-1,\frac{1}{2},-\frac{1}{3}\ldots\\
		&-1<\frac{1}{2}>-\frac{1}{3}\\&\Rightarrow \;\;\text{Die Folge ist schon nach drei Glieder nicht monoton. }
		\end{align*}
	\end{quote}
	(ii) Beschränktheit:\\
	\begin{quote}
		Die Folge ist beschränkt, da:\\
		\begin{align*}
		-1\leq \frac{(-1)^n}{n}<1\;\;\text{für alle }\;\;n\in\mathbb{N}\backslash\{0\}
		\end{align*}
		Aber: Der Grenzwert der Folge ist:\\
		\begin{align*}
		\lim\limits_ {n\to\infty} \frac{(-1)^n}{n}=0
		\end{align*}
		Deshalb ist diese Folge konvergent.\\
	\end{quote}
\end{quote}
Fazit:\\
Monotonie $ + $ Beschränktheit $ \Rightarrow $ Konvergenz\\
Aber nicht:\\
Konvergenz $ \nRightarrow $ Monotonie $ + $ Beschränktheit\\
\subsubsection{Aufgabe zum Konvergenzkriterien-Teil I}
\href{https://youtu.be/fpaipwCojak}{Video 18}\\
Beweise, dass die rekrusive Folge gegeben durch $ x_0=\sqrt{3} $ und $ x_{n+1}=\sqrt{3+x_n} $ \;für alle $ n\in \mathbb{N} $ konvergiert und berechne, wenn vorhanden, den Grenzwert.\\
Vorgehensweise bei der Lösung:\\
\begin{tabular}[t]{ll}% l=linksbündig r=rechtsbündig
	(i) & Nachweis der Monotonie  \\
	(ii) & Nachweis der Beschränktheit\\
	(iii) & Grenzwert berechnen \\	
\end{tabular} \\
\\
(i)\;\;Nachweis der Monotonie\\
\begin{quote}
	Zuerst einfach Zahlen in die Folge $ x_{n+1}=\sqrt{3+x_n} $ einsetzen.\\
	Dies ergibt:\\
	$x_0=\sqrt{3}$\\
	$x_1=\sqrt{3+\sqrt{3}}$\\
	$x_2=\sqrt{3+\sqrt{3+\sqrt{3}}}$\\
	Da die Wurzel immer positiv ist, wird die Folge streng monoton steigend sein.\\
	Zu zeigen:
	\begin{align*}
	x_n<x_{n+1}\;\;\text{für alle }\;\;n\in\mathbb{N}
	\end{align*}
	Beweis mt vollständiger Induktion:\\
	\textbullet \;\;\textbf{(I.A.):}\\
	\begin{quote}
		für $ n=0$ gilt:\\
		\begin{align*}
		x_0=\sqrt{3}<x_{n+1}=x_1=\sqrt{3+\sqrt{3}}_{\;\;_{\Square}}
		\end{align*}
	\end{quote}
	\textbullet \;\;\textbf{(I.V.):}\\
	\begin{quote}
		für ein $ n\in\mathbb{N} $ gilt:\\
		\begin{align*}
		x_{n+1}>x_{n}
		\end{align*}
	\end{quote}
	\textbullet \;\;\textbf{(I.S.):}\\
	\begin{align*}
	x_{n+2}>x_{n+1}
	\end{align*}
	Betrachte:
	\begin{align*}
	x_{n+2}=\sqrt{3+\sqrt{x_{n+1}}}\overset{(I.V.)}{\geq}\sqrt{3+x_n}=x_{n+1_{\;\;\Square}}
	\end{align*}
	Die Folge ist streng monoton steigend.
\end{quote}





\subsubsection{Aufgabe zum Konvergenzkriterien-Teil II}
\href{https://youtu.be/ldMJJfYpobk}{Video 19}\\
(ii)\;\;Nachweis der Beschränktheit\\
\begin{quote}
	Zu zeigen:\\
	Es gibt $ a,b\in\mathbb{R} $ mit $ a\leq x_n\leq b $ für alle $ n\in\mathbb{N} $\vspace*{2mm}\\
	Es gilt:\\
	$ x_n $ ist streng monoton wachsen, d.h. $ x_n<x_{n+1} $ oder $ x_0<x_1<x_2\ldots $\\
	Es gilt also für $ x_0=\sqrt{3} $:\\
	\begin{align*}
	\sqrt{3}\leq x_n\;\;\text{für alle }\;\;n\in\mathbb{N}\\
	\Rightarrow a=\sqrt{3} \;\;\text{ist untere Schranke }
	\end{align*}
	Die obere Schranke muss man durch Ausprobieren einfach abschätzen.\\
	Es soll gelten: $ x_n<10 $\\
	Beweis mit vollständiger Induktion:\\
	\textbullet \;\;(I.A.):\\
	\begin{quote}
		für $ n=0 $ gilt:\\
		\begin{align*}
		x_0=\sqrt{3}<10_{\;\;_{\Square}}
		\end{align*}
	\end{quote}
	
	\textbullet \;\;(I.V.):\\
	\begin{quote}
		für ein $ n\in\mathbb{N} $ gilt:\\
		\begin{align*}
		x_{n+1}=\sqrt{3+x_n}<10
		\end{align*}
	\end{quote}
	\textbullet \;\;(I.S.):\\
	\begin{quote}
		\begin{align*}
		x_{n+2}=\sqrt{3+x_{n+1}}<10
		\end{align*}
		Betrachte:\\
		\begin{align*}
		x_{n+2}=\sqrt{3+x_{n+1}}\overset{(I.V.)}{<} \sqrt{3+10}=\sqrt{13}<10_{\;\;_{\Square}}
		\end{align*}
		\begin{align*}
		&\Rightarrow b=10\;\;\text{ ist eine obere Schranke}\\
		&\Rightarrow a=\sqrt{3}\;\;\text{ ist eine untere Schranke}\\
		&\text{Es gilt für}\;\;a,b \in \mathbb{R} \;\;a\leq x_n< b\;\;\text{für alle }\;\;n\in\mathbb{N}
		\end{align*}
	\end{quote}
	Fazit:\\
	Die Folge ist beschränkt, außerdem ist sie streng monoton wachsend, daraus folgt mit dem Konverkenzkriterium, dass die Folge konvergent ist.
\end{quote}


\subsubsection{Aufgabe zum Konvergenzkriterien-Teil III}
\href{https://youtu.be/4zIPa5nNbdU}{Video 20}\\
(iii)\;\;Grenzwert berechnen\\
\begin{quote}
	Die Folge mit $ x_0=\sqrt{3} $ und $ x_{n+1}=\sqrt{3+x_n} $ für alle $ n\in\mathbb{N} $ ist konvergent.\vspace*{4mm}\\
	Bemerkung:\\
	\begin{quote}
		Es gilt:\\
		\begin{align*}
		(a_n)_{n\in\mathbb{N}}\;\;\text{konvergiert gegen } \;\;a\\
		\end{align*}
		Dann gilt auch:\\
		\begin{align*}
		\lim\limits_ {n\to\infty}a=\lim\limits_ {n\to\infty}a_{n+1}
		\end{align*}
		Beispiel:\\
		\begin{align*}
		\lim\limits_ {n\to\infty}\frac{1}{n}=\lim\limits_ {n\to\infty}\frac{1}{n+1}=0
		\end{align*}
		Das heißt:\\
		\begin{align*}
		\lim\limits_ {n\to\infty}x_n=a=	\lim\limits_ {n\to\infty}x_{n+1}\;\;\text{für } \;\;a=\in\mathbb{R}\\
		\end{align*}
	\end{quote} 
	Für die Folge gilt dann:\\
	\begin{align*}
	\lim\limits_ {n\to\infty}\sqrt{3+x_n}=a\\
	\Leftrightarrow a=\sqrt{3+a}
	\end{align*}
	Berachte also:\\
	\begin{align*}
	\sqrt{3+a}=a\Leftrightarrow 0&=a^2-a-3\\
	&=a^2-a+\frac{1}{4}-\frac{1}{4}-3\\
	&=\bigg(a-\frac{1}{2}\bigg)^2-\frac{13}{4}\\
	\Leftrightarrow \bigg(a-\frac{1}{2}\bigg)^2&=\frac{13}{4}\\
	\Leftrightarrow a_{1,2}&=\frac{1}{2}\pm \sqrt{\frac{13}{4}}=\frac{1}{2}\pm \frac{\sqrt{13}}{2}=\frac{1\pm\sqrt{13}}{2}
	\end{align*}
	Die Folge ist streng monoton wachsend $ (\sqrt{3}\leq x_n\;\;\text{für alle }\;\;n\in\mathbb{N}) $ , also kann die Lösung nur positiv sein.\\
	Deshalb gilt für den Grenzwert:\\
	\begin{align*}
	a=\frac{1+\sqrt{13}}{2}
	\end{align*}
	
\end{quote}



%	\chapter{Konvergenz von Funktionen}

\setcounter{section}{5}
\setcounter{subsection}{0}
\setcounter{subsubsection}{1}
\setcounter{secnumdepth}{3}
\subsection{Funktionen und Folgen}
\href{https://youtu.be/uTHOemsexgc}{Video 01}\\
Sei $ f:\mathbb{R}\supset \mathbb{D}_f\longmapsto \mathbb{R} $ eine Funktion und $ (x_n)_{n\in\mathbb{N}} $ eine reelle Folge, dann ist:\\
\begin{align*}
\bigg(f(x_n)\bigg)_{n\in\mathbb{N}}=f(x_0), f(x_1)\ldots
\end{align*}
auch eine reelle Zahlenfolge.\vspace*{4mm}\\
Beispiel:\\
Sei:\\
$ f:\mathbb{R}\longmapsto [0,\infty[ $ mit $ f(x)=x^2 $ und $x_n= \big(\frac{1}{n}\big)_{n\in\mathbb{N}\backslash\{0\}} $\\
Dann ist:\\
$ f\big((x_n)\big)_{n\in\mathbb{N}\backslash\{0\}} =\big((\frac{1}{n})^2\big)_{n\in\mathbb{N}\backslash\{0\}}=1,\frac{1}{4},\frac{1}{9},\frac{1}{16}\ldots$\\
\begin{center}

\begin{tikzpicture}[yscale=10,xscale=5]

% Gitternetzlinien
\draw[xstep=0.1cm, ystep=0.1cm, very thin, color=lightgray] (0,1.2) grid (1.2,0);
% x-Achse
\draw[->] (0,0) -- (1.2,0) node[below,xshift=-1em] {$x$}; 
%Zahlen auf x-Achse

% y-Achse 
\draw[->] (0,0) -- (0,1.2) node[left,yshift=-1em] {$f(x)$};%node[above left]
%Zahlen auf y-Achse

%\foreach \y in {1,...,1}
%\draw[shift={(0,\y)},color=black] (2pt,0pt) -- (-2pt,0pt) node[left] {\footnotesize $\y$};
%Ursprung
\draw[color=black] (1,0) node[yshift=-1.5em] {\footnotesize $1$};
\draw[color=black] (1,0) --(1,-0.5pt);
\draw[color=black] (0,0) node[yshift=-1.5em] {\footnotesize $0$};
\draw[color=black] (0,0) --(0,-0.5pt);
\draw[color=black] (0.5,0) node[yshift=-1.5em] {\footnotesize $\frac{1}{2}$};
\draw[color=black] (0.5,0) --(0.5,-0.5pt);

\draw[color=black] (0.25,0) node[yshift=-1.5em] {\footnotesize $\frac{1}{4}$};
\draw[color=black] (0.25,0) --(0.25,-0.5pt);

\draw[color=black] (0.33,0) node[yshift=-1.5em] {\footnotesize $\frac{1}{3}$};
\draw[color=black] (0.33,0) --(0.33,-0.5pt);
\draw[color=black] (0.2,0) node[yshift=-1.5em] {\footnotesize $\frac{1}{5}$};
\draw[color=black] (0.2,0) --(0.2,-0.5pt);
\draw[color=black] (0.16,0) node[yshift=-1.5em] {\footnotesize $\frac{1}{6}$};
\draw[color=black] (0.16,0) --(0.16,-0.5pt);
\draw[color=black] (0,0.25) node[xshift=-3em] {\footnotesize $f(x_2)=\frac{1}{4}$};
\draw[color=black] (0,0.25) --(-1pt,0.25);
\draw[color=black] (0,0.11) node[xshift=-3em] {\footnotesize $f(x_3)=\frac{1}{9}$};
\draw[color=black] (0,0.11) --(-1pt,0.11);
\draw[color=black] (0,0.0625) node[xshift=-3em] {\footnotesize $f(x_4)=\frac{1}{16}$};
\draw[color=black] (0,0.0625) --(-1pt,0.0625);
\draw[color=black] (0,1) node[xshift=-3em] {\footnotesize $f(x_1)=1$};
\draw[color=black] (0,1) --(-1pt,1);

\draw[color=red,thick,dashed] (1,0) --(1,1)--(0,1);
\draw[color=red,thick,dashed] (0.5,0) --(0.5,0.25)--(0,0.25);
\draw[color=red,thick,dashed] (0.25,0) --(0.25,1/16)--(0,1/16);
\draw[color=red,thick,dashed] (1/3,0) --(1/3,1/9)--(0,1/9);
\draw[thick,color=black] plot[samples=200, domain=0:1.1] 
(\x,{\x*\x}) node[below, xshift=1em]{$f(x)$};



\end{tikzpicture}\
\end{center}
Die Zahlenfolge $ f\big((x_n)\big)_{n\in\mathbb{N}\backslash\{0\}} $ konvergiert auf der y-Achse bis zum Grenzwert 0 und die Folge $ (\frac{1}{n})_{n\in\mathbb{N}\backslash\{0\}} $ konvergiert ebenfalls auf der x-Achse gegen den Grenzwert 0.\\
\subsection{Konvergenz von Funktionen}
\href{https://youtu.be/JM5KwRHuPKc}{Video 02}\\
Definition:\\
Sei $ f:\mathbb{R}\supset \mathbb{D}_f\longmapsto \mathbb{R} $ eine Funktion, dann gilt:\\
Die Funktion $ f $ hat für $ x $ gegen $ x_0 $ den Grenzwert $ y_0 $, wenn für jede Folge in $ (x_n)_{n\in\mathbb{N}} $ in $ \mathbb{D}_f\backslash \{x_0\} $, mit $ \lim\limits_ {n\to\infty}x_n=x_0 $, der Grenzwert $ \lim\limits_ {n\to\infty} f(x_n) $ exestiert und gleich $ y_0 $ ist.\vspace*{2mm}\\
Beispiel:\\
Sei $ f: \mathbb{R}\longmapsto \mathbb{R} $ mit $ f(x)=2x $ und $ x_0=2\Rightarrow y_0=4 $ und die Folge $ x_n=(\frac{1}{n})+2 $.\\
Es gilt:
\begin{align*}
\lim\limits_ {n\to\infty} \frac{1}{n}+2=2=x_0
\end{align*}
Ebenso gilt:\\
\begin{align*}
\lim\limits_ {n\to\infty} f(x_n)=2\bigg(\frac{1}{n}+2\bigg)=4=y_0
\end{align*}
Diese Aussage gilt jetzt genau für diese eine Folge, aber nicht für alle!\\
Jede Folge mit dem Grenzwert 2 muss ergeben, dass der Grenzwert existiert und $ \lim\limits_ {n\to\infty} f(x_n)=y_0 $ ist:\\
\begin{center}

\begin{tikzpicture}

% Gitternetzlinien
\draw[xstep=1cm, ystep=1cm, very thin, color=lightgray] (-4,5) grid (4,-2);
% x-Achse
\draw[->] (-4,0) -- (4,0) node[below,xshift=-1em] {$x$}; 
%Zahlen auf x-Achse
\foreach \x in {-3,...,-1}
\draw[shift={(\x,0)},color=black] (0pt,2pt) -- (0pt,-2pt) node[below] {\footnotesize $\x$};
\foreach \x in {1,...,3}
\draw[shift={(\x,0)},color=black] (0pt,2pt) -- (0pt,-2pt) node[below] {\footnotesize $\x$};
% y-Achse 
\draw[->] (0,-2) -- (0,5) node[left,yshift=-1em] {$f(x)$};%node[above left]
%Zahlen auf y-Achse
\foreach \y in {-1,...,-1}
\draw[shift={(0,\y)},color=black] (2pt,0pt) -- (-2pt,0pt) node[left] {\footnotesize $\y$};
\foreach \y in {1,...,4}
\draw[shift={(0,\y)},color=black] (2pt,0pt) -- (-2pt,0pt) node[left] {\footnotesize $\y$};
%Ursprung
\draw[color=black] (0pt,-10pt) node[right] {\footnotesize $0$};



\draw[thick,color=black] plot[samples=200, domain=-1:2.5] 
(\x,{2*\x}) node[below, xshift=1em]{$f(x)$};

\draw[color=blue] (2,0) node[yshift=-20] {\footnotesize $x_0$};
\draw[color=blue] (1,0) node[yshift=-30] {\footnotesize $(x_n)\longrightarrow x_0$};
\draw[color=blue] (0,4) node[xshift=-20] {\footnotesize $y_0$};
\draw[color=blue] (0,2) node[xshift=-20,rotate=90] {\footnotesize $f(x_n)\longrightarrow y_0$};
\draw[color=blue,dashed,thick] (2,0)--(2,4)--(0,4);
\end{tikzpicture}\\
\end{center}
Bemerkung:\\
$ x_0 $ muß nicht notwendigerweise in $ \mathbb{D}_f $ sein.\\
Wir betrachten auch $ \lim\limits_ {x\to\pm\infty}f(x) $. Z.B.($ \lim\limits_ {x\to\infty}f(x)=\frac{1}{\frac{1}{x}+1} $)\\
\subsubsection{Beispiel zum Nachweis der Konvergenz von Funktionen-Teil I}
\href{https://youtu.be/sLrX5Et79Kg}{Video 03}\\
(i)\;\;$f:\mathbb{R}\longmapsto[0,\infty[$ für $ f(x)=x^2 $\\
\begin{quote}
	Zeige:\\
	\begin{align*}
	\lim\limits_ {x\to 0}f(x)=\lim\limits_ {x\to 0}x^2=0\\
	\end{align*}
	\textbf{Fehler: Falsche Betrachtung!}\\
	Betrachte:\\ Nullfolge $ (\frac{1}{n})_{n\in\mathbb{N}} $\\
	Es gilt:
	\begin{align*}
	\lim\limits_ {n\to \infty}f(x_n)=\lim\limits_ {n\to \infty}\frac{1}{n}=0
	\end{align*}
	Dies reicht nicht, weil $ (\frac{1}{n})_{n\in\mathbb{N}} $ eine konkrete Folge ist und nicht alle Nullfolgen repräsentieren!!\vspace*{4mm}\\
	\textbf{Richtig:}
	Man nehme eine beliebige Folge:\\
	Sei $ (x_n) $ eine beliebige Folge für die gilt $ \lim\limits_ {n\to\infty}(x_n)=0 $\\
	Betrachte:\\
	\begin{align*}
	\lim\limits_ {n\to\infty}f(x_n)=\lim\limits_ {n\to\infty}x_{n}^2=\lim\limits_ {n\to\infty}x_n\lim\limits_ {n\to\infty}x_n=0\cdot0=0
	\end{align*}
	D.h.:\\
	Der $ \lim\limits_ {n\to\infty}f(x_n) =0$ für jede beliebige Nullfolge.\\
	\begin{align*}
	\Rightarrow\lim\limits_ {x\to 0}f(x)=\lim\limits_ {x\to 0}x^2=0
	\end{align*}
\end{quote}
(ii)\;\;$ \lim\limits_ {x\to 2}x^2=4 $\\
\begin{quote}
Wähle eine beliebige Folge $(x_n)_{n\in\mathbb{N}}$ für die gilt $ \lim\limits_ {n\to\infty}x_n=2 $ .
\begin{align*}
\lim\limits_ {n\to\infty}(x_n)^2=\lim\limits_ {n\to\infty}x_n\lim\limits_ {n\to\infty}x_n=2\cdot2=4
\end{align*}
\end{quote}
(iii) \;\;$\lim\limits_ {x\to 4}\frac{1}{x}=\frac{1}{4}$\\
\begin{quote}
Wähle eine beliebige Folge $(x_n)_{n\in\mathbb{N}}$ für die gilt $ \lim\limits_ {n\to\infty}x_n=4 $ .\\
\begin{align*}
\lim\limits_ {n\to\infty}\frac{1}{x_n}=\frac{\lim\limits_ {n\to\infty}1}{\lim\limits_ {n\to\infty}x_n}=\frac{1}{4}
\end{align*}
\end{quote}
\subsubsection{Beispiel zum Nachweis der Konvergenz von Funktionen-Teil II}
\href{https://youtu.be/in86heC1gIQ}{Video 04}\\
(i)\;\; $ f_1(x):\mathbb{R}\longmapsto \mathbb{R} $ mit $ f(x)=2x $\\
\begin{quote}
	\begin{align*}
	\lim\limits_ {x\to\infty} f_1(x)=\lim\limits_ {x\to\infty}2x=0\\
	\end{align*}
\end{quote}
(ii)\;\;$ f_2:\mathbb{R}\backslash\{0\}\longmapsto\mathbb{R} $ mit $ f(x) =\frac{2x}{x}$\\
\begin{quote}
	\begin{align*}
	\lim\limits_ {x\to\infty}f_2(x)=\lim\limits_ {x\to\infty}\frac{2x}{x}=0
	\end{align*}
	Null kann hier nicht eingesetzt werden! Der Grenzwert ist keine Funktion!\\
\end{quote}
(iii)\;\;$ \lim\limits_ {x\to\infty}\sin(x) $\\
\begin{quote}
	\begin{center}
		\begin{tikzpicture} [yscale=2]
		
		% Gitternetzlinien
		\draw[xstep=1cm, ystep=0.5cm, very thin, color=lightgray] (-5,1.5) grid (8,-1.5);
		% x-Achse
		\draw[->] (-5,0) -- (8,0) node[below,xshift=-1em] {$x$}; 
		%Zahlen auf x-Achse
		%\foreach \x in {-1,...,-6}
		%\draw[shift={(\x,0)},color=black] (0pt,2pt) -- (0pt,-2pt) node[below] {\footnotesize $\x$};
		%\foreach \x in {1,...,5}
		%\draw[shift={(\x,0)},color=black] (0pt,2pt) -- (0pt,-2pt) node[below] {\footnotesize $\x$};
		% y-Achse 
		\draw[->] (0,-2) -- (0,2) node[left,yshift=-1em] {$f(x)$};%node[above left]
		%Zahlen auf y-Achse
		\foreach \y in {-1,...,-1}
		\draw[shift={(0,\y)},color=black] (2pt,0pt) -- (-2pt,0pt) node[left] {\footnotesize $\y$};
		\foreach \y in {1,...,1}
		\draw[shift={(0,\y)},color=black] (2pt,0pt) -- (-2pt,0pt) node[left] {\footnotesize $\y$};
		%Ursprung
		\draw[color=black] (0pt,-5pt) node[right] {\footnotesize $0$};
		
		\draw[color=black] (3.14/2,-5pt) node {\footnotesize $\frac{\pi}{2}$};
		\draw[color=black] (3.14/2,0) --(3.14/2,-0.05);
		
		\draw[color=black] (3.14,-5pt) node {\footnotesize $\pi$};
		\draw[color=black] (3.14,0) --(3.14,-0.05);
		
		\draw[color=black] (3.14*2,-5pt) node {\footnotesize $2\pi$};
		\draw[color=black] (3.14*2,0) --(3.14*2,-0.05);
		
		\draw[color=black] (3.14*3/2,-5pt) node {\footnotesize $\frac{3}{2}\pi$};
		\draw[color=black] (3.14*3/2,0) --(3.14*3/2,-0.05);
		\draw[thick,color=black] plot[samples=200, domain=-5:8] 
		(\x,{sin(\x r)}) node[below, xshift=1em]{$f(x)$};
		
		
		
		\end{tikzpicture}\\
	\end{center}
Die Frage: Wohin läuft die Funktion?\\
Betrachte:\\
\textbullet \;\;Die Folge $ (n\cdot\pi)_{n\in\mathbb{N}} $\\
\begin{quote}
	Offensichtlich gilt:\\
	\begin{align}
	&\lim\limits_ {n\to\infty} n\cdot\pi=\infty \nonumber\\
	&\lim\limits_ {n\to\infty}(\sin(n\cdot\pi))=0
	\end{align}
\end{quote}
\textbullet \;\;Die Folge $ (\frac{\pi}{2}+n\cdot\pi)_{n\in\mathbb{N}} $\\
\begin{quote}
	Offensichtlich gilt:\\
	\begin{align}
	&\lim\limits_ {n\to\infty}\bigg(\frac{\pi}{2}+n\cdot\pi\bigg)=\infty\nonumber \\
	&\lim\limits_ {n\to\infty}\sin\bigg(\frac{\pi}{2}+n\cdot\pi\bigg)=(-1)^n=+1,-1,+1\ldots
	\end{align}
\end{quote}
Der Grenzwert von $ \lim\limits_ {n\to\infty}(-1)^n $ existiert nicht, also kann auch der Grenzwert $ \lim\limits_ {n\to\infty}\sin(x) $ nicht existieren.
Nach der Definition (Konvergenz von Folgen) existiert kein Grenzwert. Der Grenzwert von (V.1) muß gleich den Grenzwert von (V.2) sein.\vspace*{2mm}\\
\end{quote}
(iv)\;\;$ \lim\limits_ {x\to\infty} x=\infty$\\
\begin{quote}
	Die Folge divergiert in das positive Unendliche.\\
\end{quote}
(v)\;\;$\lim\limits_ {n\to-\infty}x=-\infty$\\
\begin{quote}
	Die Folge divergiert in das negative Unendliche.\\
\end{quote}
\subsection{Einseitige Grenzwerte-Teil I u. II}
\href{https://youtu.be/oxpwECM2DKc}{Video 05}\\
Sei $ f:\mathbb{R}\supset \mathbb{D}_f\backslash\{x_0\}\longmapsto \mathbb{R} $ eine Funktion.\\
$ y_0 $ ist der linksseitige Grenzwert von $ f $ für $ x $ gegen $ x_0 $, wenn gilt: $ \lim\limits_ {n\to\infty}(a_n)=y_0 $ für jede Folge $ (a_n)_{n\in\mathbb{N}} $ mit $ a_n<x_0 $ für alle $ n\in\mathbb{N} $ und rechtseitig, wenn gilt: $ \lim\limits_ {n\to\infty}(a_n)=y_0 $ für jede Folge $ (a_n)_{n\in\mathbb{N}} $ mit $ a_n>x_0 $ für alle $ n\in\mathbb{N} $\\
\pagebreak
\begin{center}
	Die Folge $ (a_n) $ kommt von links zum Wert $ x_0 $.\\
	\begin{tikzpicture}
	
	% Gitternetzlinien
	\draw[xstep=1cm, ystep=1cm, very thin, color=lightgray] (-4,5) grid (4,-2);
	% x-Achse
	\draw[->] (-4,0) -- (4,0) node[below,xshift=-1em] {$x$}; 
	%Zahlen auf x-Achse
	\foreach \x in {-3,...,-1}
	\draw[shift={(\x,0)},color=black] (0pt,2pt) -- (0pt,-2pt) node[below] {\footnotesize $\x$};
	\foreach \x in {1,...,3}
	\draw[shift={(\x,0)},color=black] (0pt,2pt) -- (0pt,-2pt) node[below] {\footnotesize $\x$};
	% y-Achse 
	\draw[->] (0,-2) -- (0,5) node[left,yshift=-1em] {$f(x)$};%node[above left]
	%Zahlen auf y-Achse
	\foreach \y in {-1,...,-1}
	\draw[shift={(0,\y)},color=black] (2pt,0pt) -- (-2pt,0pt) node[left] {\footnotesize $\y$};
	\foreach \y in {1,...,4}
	\draw[shift={(0,\y)},color=black] (2pt,0pt) -- (-2pt,0pt) node[left] {\footnotesize $\y$};
	%Ursprung
	\draw[color=black] (0pt,-10pt) node[right] {\footnotesize $0$};
	
	
	
	\draw[thick,color=black] plot[samples=200, domain=-1:2.5] 
	(\x,{2*\x}) node[below, xshift=1em]{$f(x)$};
	
	\draw[color=blue] (2,0) node[yshift=-20] {\footnotesize $x_0$};
	\draw[color=blue] (1,0) node[yshift=-30] {\footnotesize $(a_n)\longrightarrow x_0$};
	\draw[color=blue] (0,4) node[xshift=-20] {\footnotesize $y_0$};
	\draw[color=blue] (0,2) node[xshift=-20,rotate=90] {\footnotesize $f(a_n)\longrightarrow y_0$};
	\draw[color=blue,dashed,thick] (2,0)--(2,4)--(0,4);
	\end{tikzpicture}\\
\end{center}
Linksseitiger Grenzwert:\\
\begin{align*}
\lim\limits_ {x\to\infty} f(x)=\lim\limits_ {\substack{x\to\infty\\x<x_0}}f(x)=y_0
\end{align*}
Rechtsseitiger Grenzwert:\\
\begin{align*}
\lim\limits_ {x\to\infty} f(x)=\lim\limits_ {\substack{x\to\infty\\x>x_0}}f(x)=y_0
\end{align*}
Beispiel:\\
\textbullet \;\; Linksseitiger Grenzwert von: $ \lim\limits_ {\substack{x\to\infty\\x<0}}\frac{2|x|}{x} $\\
\begin{align*}
\lim\limits_ {\substack{x\to\infty\\x<0}}\frac{2|x|}{x}=\lim\limits_ {x\to\infty}\frac{2(-x)}{x}=-2\\
\end{align*}
\textbullet \;\; Rechtsseitiger Grenzwert von: $ \lim\limits_ {\substack{x\to\infty\\x>0}}\frac{2|x|}{x} $\\
\begin{align*}
\lim\limits_ {\substack{x\to\infty\\x>0}}\frac{2|x|}{x}=\lim\limits_ {x\to\infty}\frac{2(x)}{x}=2\\
\end{align*}
%\subsection{Einseitige Grenzwerte-Teil II}
\href{https://youtu.be/KLn2p_dOASs}{Video 06}\\
\underline{Satz:}\\
\begin{center}

\fbox{\parbox  {128mm} {Existieren $ \lim\limits_ {\substack{x\to x_0\\x<x_0}}=a $ und $ \lim\limits_ {\substack{x\to x_0\\x>x_0}}=b $ und gilt $ a=b $, dann existiert auch:\\
		 \begin{center}
		 	 $ \lim\limits_ {x\to x_0}f(x)=a $
		 \end{center}}}
\end{center}
\vspace*{4mm}
\newpage
\noindent
Beispiel:\\
\textbullet \;\;$\lim\limits_ {x\to 0} \frac{2|x|}{x} $\\
$ \lim\limits_ {\substack{x\to 0\\x<0}}\frac{2(-x)}{x}=-2\neq 2=\lim\limits_ {\substack{x\to 0\\x>0}}\frac{2x}{x}\Rightarrow \lim\limits_ {x\to 0} \frac{2|x|}{x}$ existiert nicht \vspace*{2mm}\\
\textbullet \;\; $  \lim\limits_ {\substack{x\to 0\\x<0}}|x| $\\
\begin{quote}
	$  \lim\limits_ {\substack{x\to 0\\x<0}}|x|=\lim\limits_ {x\to 0}(-x) =0$\\
	$  \lim\limits_ {\substack{x\to 0\\x>0}}|x|=\lim\limits_ {x\ 0}  x=0$\\
	$\Rightarrow \lim\limits_ {x\to 0}|x|=0$\\
\end{quote}
\textbullet \;\; $ \lim\limits_ {x\to 0}\frac{1}{x} $\\
\begin{quote}
	$ \lim\limits_ {\substack{x\to 0\\x<0}}\frac{1}{x}=-\infty $ hat keinen Grenzwert, geht in das negative Unendliche.\\
	$ \lim\limits_ {\substack{x\to 0\\x>0}}\frac{1}{x}=\infty $ geht in das positive Unendliche\\
	$ \Rightarrow \lim\limits_ {x\to\ 0}\frac{1}{x}$ existiert nicht.
	\begin{center}
		\begin{tikzpicture}
		
		% Gitternetzlinien
		\draw[xstep=1cm, ystep=1cm, very thin, color=lightgray] (-4,4) grid (4,-4);
		% x-Achse
		\draw[->] (-4,0) -- (4,0) node[below,xshift=-1em] {$x$}; 
		%Zahlen auf x-Achse
		\foreach \x in {-3,...,-1}
		\draw[shift={(\x,0)},color=black] (0pt,2pt) -- (0pt,-2pt) node[below] {\footnotesize $\x$};
		\foreach \x in {1,...,3}
		\draw[shift={(\x,0)},color=black] (0pt,2pt) -- (0pt,-2pt) node[below] {\footnotesize $\x$};
		% y-Achse 
		\draw[->] (0,-4) -- (0,4) node[left,yshift=-1em] {$f(x)$};%node[above left]
		%Zahlen auf y-Achse
		\foreach \y in {-3,...,-1}
		\draw[shift={(0,\y)},color=black] (2pt,0pt) -- (-2pt,0pt) node[left] {\footnotesize $\y$};
		\foreach \y in {1,...,3}
		\draw[shift={(0,\y)},color=black] (2pt,0pt) -- (-2pt,0pt) node[left] {\footnotesize $\y$};
		%Ursprung
		\draw[color=black] (0pt,-10pt) node[right] {\footnotesize $0$};
		
		\draw[thick,color=blue] plot[samples=200, domain=0.25:4] 
		(\x,{1/\x}) node[xshift=-5em,yshift=90] {$f(x)=\frac{1}{x}$};
		
		\draw[thick,color=blue] plot[samples=200, domain=-4:-0.25] 
		(\x,{1/\x}) ;
		
		
		
		\end{tikzpicture}\\
	\end{center}
\end{quote}
\subsubsection{Aufgabe zum Thema Einseitige Grenzwerte}
\href{https://youtu.be/3QFihiHQzq0}{Video 07}\\
Geben eine Funktion:  $f(x) =     \left\{ \begin{array}{lll}
	x^4-1&\mbox{für} & x>1\\
	x^2 &\mbox{für} & x\leq 1 \\
	
\end{array}\right.$\;\;\;\;$ x\in\mathbb{R} $\\
Überprüfe die Grenzwerte:\\
\begin{tabular}[t]{lll}% l=linksbündig r=rechtsbündig
	(i) & $ \lim\limits_ {x\to -1}f(x) $  \\
	(ii) & $ \lim\limits_ {x\to 1}f(x) $ \\	
\end{tabular} \vspace*{2mm}\\
(i) \;\;$\lim\limits_ {x\to -1}f(x)$\\
\begin{quote}
	\textbullet \;\; $ \lim\limits_ {\substack{x\to -1\\x<-1}}f(x) =\lim\limits_ {x\to -1}x^2=1$\\
	\textbullet \;\;$ \lim\limits_ {\substack{x\to -1\\x>-1}}f(x) =\lim\limits_ {x\to -1}x^2=1$\\
	Da beide einseitigen Grenzwerte exestieren und gleich sind, gilt hier $ \lim\limits_ {x\to -1}f(x)=1 $\\
\end{quote}
(ii) \;\;$\lim\limits_ {x\to 1}f(x)$\\
\begin{quote}
	\textbullet \;\; $ \lim\limits_ {\substack{x\to 1\\x<1}}f(x) =\lim\limits_ {x\to 1}x^2=1$\\
	\textbullet \;\; $ \lim\limits_ {\substack{x\to 1\\x>1}}f(x) =\lim\limits_ {x\to 1}x^4-1=0$\\
	Da beide einseitigen Grenzwerte verschieden sind existiert $ \lim\limits_ {x\to\ 1}f(x) $ nicht.\\
\end{quote}
\begin{center}
 $f(x) =     \left\{ \begin{array}{lll}
x^4-1&\mbox{für} & x>1\\
x^2 &\mbox{für} & x\leq 1 \\

\end{array}\right.$\\
\begin{tikzpicture}

% Gitternetzlinien
\draw[xstep=1cm, ystep=1cm, very thin, color=lightgray] (-4,4) grid (4,-2);
% x-Achse
\draw[->] (-4,0) -- (4,0) node[below,xshift=-1em] {$x$}; 
%Zahlen auf x-Achse
\foreach \x in {-3,...,-1}
\draw[shift={(\x,0)},color=black] (0pt,2pt) -- (0pt,-2pt) node[below] {\footnotesize $\x$};
\foreach \x in {1,...,3}
\draw[shift={(\x,0)},color=black] (0pt,2pt) -- (0pt,-2pt) node[below] {\footnotesize $\x$};
% y-Achse 
\draw[->] (0,-2) -- (0,4) node[left,yshift=-1em] {$f(x)$};%node[above left]
%Zahlen auf y-Achse
\foreach \y in {-1,...,-1}
\draw[shift={(0,\y)},color=black] (2pt,0pt) -- (-2pt,0pt) node[left] {\footnotesize $\y$};
\foreach \y in {1,...,3}
\draw[shift={(0,\y)},color=black] (2pt,0pt) -- (-2pt,0pt) node[left] {\footnotesize $\y$};
%Ursprung
\draw[color=black] (0pt,-10pt) node[right] {\footnotesize $0$};

\draw[thick,color=blue] plot[samples=200, domain=-2:1] 
(\x,{\x*\x});

\draw[thick,color=blue] plot[samples=200, domain=1:1.5] 
(\x,{(\x*\x*\x*\x)-1}) node[below, xshift=1em]{$f(x)$};





\end{tikzpicture}\\
\end{center}

%	\chapter{Stetigkeit}

\setcounter{section}{6}
\setcounter{subsection}{0}
\setcounter{subsubsection}{1}
\setcounter{secnumdepth}{3}
\subsection{Stetigkeit (Definition)}
\href{https://youtu.be/w5zo40vnyRQ}{Video 01}\\
Sei $ f:\mathbb{D}_f\longmapsto \mathbb{W}_f $ Funktion und $ x_0\in \mathbb{D}_f $\\
Definition:\\
Die Funktion $ f $ heißt in $ x_0\in\mathbb{D}_f $ stetig, wenn $ \lim\limits_ {x\to x_0}f(x)=f(x_0) $\\
\begin{align*}
\Leftrightarrow \lim\limits_ {\substack{x\to x_0\\x<0}}f(x)=f(x_0)=\lim\limits_ {\substack{x\to x_0\\x>0}}f(x)=f(x_0)
\end{align*}
\begin{center}
	Die Folge $ (a_n)$ bzw. $f(a_n) $ kommt von links zum Wert $ x_0$ bzw.$f(x_0) $\\ und rechts zum Wert $ x_0$ bzw. $f(x_0) $.\\
	\begin{tikzpicture}
	
	% Gitternetzlinien
	\draw[xstep=1cm, ystep=1cm, very thin, color=lightgray] (-4,8) grid (4,-2);
	% x-Achse
	\draw[->] (-4,0) -- (4,0) node[below,xshift=-1em] {$x$}; 
	%Zahlen auf x-Achse
	\foreach \x in {-3,...,-1}
	\draw[shift={(\x,0)},color=black] (0pt,2pt) -- (0pt,-2pt) node[below] {\footnotesize $\x$};
	\foreach \x in {1,...,3}
	\draw[shift={(\x,0)},color=black] (0pt,2pt) -- (0pt,-2pt) node[below] {\footnotesize $\x$};
	% y-Achse 
	\draw[->] (0,-2) -- (0,8) node[left,yshift=-1em] {$f(x)$};%node[above left]
	%Zahlen auf y-Achse
	\foreach \y in {-1,...,-1}
	\draw[shift={(0,\y)},color=black] (2pt,0pt) -- (-2pt,0pt) node[left] {\footnotesize $\y$};
	\foreach \y in {1,...,7}
	\draw[shift={(0,\y)},color=black] (2pt,0pt) -- (-2pt,0pt) node[left] {\footnotesize $\y$};
	%Ursprung
	\draw[color=black] (0pt,-10pt) node[right] {\footnotesize $0$};
	
	
	
	\draw[thick,color=black] plot[samples=200, domain=-1:4] 
	(\x,{2*\x}) node[below, xshift=-1.5em]{$f(x)$};
	
	\draw[color=blue] (2,0) node[yshift=-20] {\footnotesize $x_0$};
	\draw[color=blue] (1,0) node[yshift=-30] {\footnotesize $(a_n)\longrightarrow x_0$};
	\draw[color=blue] (0,4) node[xshift=-20] {\footnotesize $f(x_0)$};
	\draw[color=blue] (0,2) node[xshift=-20,rotate=90] {\footnotesize $f(a_n)\longrightarrow f(x_0)$};
	\draw[color=blue,dashed,thick] (2,0)--(2,4)--(0,4);
	\draw[color=red] (3,0) node[yshift=-30] {\footnotesize $x_0\longleftarrow (a_n)$};
	\draw[color=red] (1,5.6) node[xshift=-20,rotate=-90] {\footnotesize $f(a_n)\longrightarrow f(x_0)$};
	\end{tikzpicture}\\
\end{center}
Beispiel:\\
Sei $ f:\mathbb{R}\longmapsto[0,\infty[ $ mit $ f(x)=x^2 $\\
\textbullet\;\; Wenn $ f(x) $ ist stetig bei $ x_0=0 $, dann ist $ f(x) $ stetig im Punkt $ x_0 $.\\
\begin{quote}
	\begin{align*}
	\lim\limits_ {\substack{x \to 0\\x<0}}f(x)=	\lim\limits_ {\substack{x \to 0\\x<0}}x^2=\underbrace {0}_{f(0)=0^2}=\lim\limits_ {\substack{x \to 0\\x>0}}f(x)=	\lim\limits_ {\substack{x \to 0\\x>0}}x^2
	\end{align*}
	
\end{quote}
\textbullet\;\; Wenn $ f(x) $ ist stetig bei jeden Wert von $ x_0\in\mathbb{D}_f$, dann heißt $ f(x) $ ist stetig.\\
\begin{quote}
	\begin{align*}
	\lim\limits_ {\substack{x \to x_0\\x<0}}f(x)=	\lim\limits_ {\substack{x \to x_0\\x<0}}x^2=\underbrace {x_0^2}_{f(x_0)=x_0^2}=\lim\limits_ {\substack{x \to x_0\\x>0}}f(x)=	\lim\limits_ {\substack{x \to x_0\\x>0}}x^2
	\end{align*}
	
	
\end{quote}
\subsubsection{Unstetigkeit-Teil I}
\href{https://youtu.be/KsP333nJ9PY}{Video 02}\\
Die Funktion $ f:\mathbb{D}_f\longmapsto\mathbb{W}_f $ heißt in $ x_0 \in \mathbb{D}_f $ unstetig , wenn $ f $ in $ x_0 $ nicht stetig ist.\\
Wenn die Funktion in einem Punkt $ x_0 $ nicht stetig ist, dann ist die ganze Funktion nicht stetig.\\
\begin{align*}
\boxed{\lim\limits_ {\substack{x\to x_0\\x<x_0}}f(x)=\lim\limits_ {\substack{x\to x_0\\x>x_0}}f(x)=f(x_0)}
\end{align*}
\textbf{Typ 1: Lücke}\\
\begin{align*}
&\text{Grenzwerte exestieren }\\
&\boxed{\lim\limits_ {\substack{x\to x_0\\x<x_0}}f(x)=\lim\limits_ {\substack{x\to x_0\\x>x_0}}f(x)\neq f(x_0)}
\end{align*}
Bemerkung: Grenzwert muß ein reeller Wert sein:\vspace*{2mm}\\
Beispiel:\\
Sei: $ f:\mathbb{R}\longmapsto \mathbb{R} $ mit $ f(x) =     \left\{ \begin{array}{rcl}
x^2& \mbox{für} & x\neq0\\
1& \mbox{für} & x= 0 \\

\end{array}\right. $\\
\begin{center}
	\begin{tikzpicture}
	
	% Gitternetzlinien
	\draw[xstep=1cm, ystep=1cm, very thin, color=lightgray] (-4,4) grid (4,-2);
	% x-Achse
	\draw[->] (-4,0) -- (4,0) node[below,xshift=-1em] {$x$}; 
	%Zahlen auf x-Achse
	\foreach \x in {-3,...,-1}
	\draw[shift={(\x,0)},color=black] (0pt,2pt) -- (0pt,-2pt) node[below] {\footnotesize $\x$};
	\foreach \x in {1,...,3}
	\draw[shift={(\x,0)},color=black] (0pt,2pt) -- (0pt,-2pt) node[below] {\footnotesize $\x$};
	% y-Achse 
	\draw[->] (0,-2) -- (0,4) node[left,yshift=-1em] {$f(x)$};%node[above left]
	%Zahlen auf y-Achse
	\foreach \y in {-1,...,-1}
	\draw[shift={(0,\y)},color=black] (2pt,0pt) -- (-2pt,0pt) node[left] {\footnotesize $\y$};
	\foreach \y in {1,...,3}
	\draw[shift={(0,\y)},color=black] (2pt,0pt) -- (-2pt,0pt) node[left] {\footnotesize $\y$};
	%Ursprung
	\draw[color=blue] (-15pt,-20pt) node[right] {\footnotesize Lücke};
	\draw[color=blue] (-15pt,-10pt) node[right] {\footnotesize $x=0$};
	\draw[color=blue] (-20pt,40pt) node[right] {\footnotesize $f(x)=1$};
	\draw[thick,color=blue] plot[samples=200, domain=-2:-0.1] 
	(\x,{\x*\x});
	
	\draw[thick,color=blue] plot[samples=200, domain=0.1:2] 
	(\x,{\x*\x}) node[below, xshift=1em]{$f(x)$};
	
	\draw [blue] (0,0) circle (.081cm) ;
	\fill [blue] (0,1) circle (.081cm) ;	
	\end{tikzpicture}\\
\end{center}
Wenn der linksseitige Grenzwert bei $ x_0 $ und der rechtsseitige Grenzwert bei $ x_0 $ gleich ist und der Funktionswert nicht gleich den Grenzwert ist, dann spricht man von einer Lücke.\\
Es gilt:\\
\begin{align*}
&\lim\limits_ {\substack{x \to 0\\x<0}}f(x)=0=\lim\limits_ {\substack{x \to 0\\x>0}}f(x)\neq f(x_0)=f(0)=1\\
&\Rightarrow\;\;f\;\text{ist in  }\;x_0=0\;\;\text{nicht stetig}
\end{align*}
Da der linksseitige und der rechtsseitige Grenzwert gleich ist und der Funktionswert nicht gleich dem Grenzwert  an der Stelle $ x_0 $  ist, ist die Funktion nicht stetig und besitzt eine behebbare Lücke.\\
Bemerkung:\\
In diesen Fall kann man die Unstetigkeit beheben.\\
Z.B. ist die Funktion $ \overset{\AC}{f}:\mathbb{R}\longmapsto \mathbb{R}$  mit $ f(x) =     \left\{ \begin{array}{rcl}
x^2 & \mbox{für} & x\neq0\\
0 & \mbox{für} & x= 0 \\

\end{array}\right.$ stetig.\\
Damit ist die Unstetigkeitsstelle behoben.\\
Die Funktion $ \overset{\AC}{f}$ ist gleich die stetige Fortsetzung von $ f $.\vspace*{6mm}\\
Es ist eigentlich nicht notwendig, dass $ x_0 $ in $ \mathbb{D}_f $ ist.\\
Sei: $ f:\mathbb{R}\backslash \{0\}\longmapsto\mathbb{R} $ mit $ f(x)=x^2 $\\
\begin{center}
	\begin{tikzpicture}
	
	% Gitternetzlinien
	\draw[xstep=1cm, ystep=1cm, very thin, color=lightgray] (-4,4) grid (4,-2);
	% x-Achse
	\draw[->] (-4,0) -- (4,0) node[below,xshift=-1em] {$x$}; 
	%Zahlen auf x-Achse
	\foreach \x in {-3,...,-1}
	\draw[shift={(\x,0)},color=black] (0pt,2pt) -- (0pt,-2pt) node[below] {\footnotesize $\x$};
	\foreach \x in {1,...,3}
	\draw[shift={(\x,0)},color=black] (0pt,2pt) -- (0pt,-2pt) node[below] {\footnotesize $\x$};
	% y-Achse 
	\draw[->] (0,-2) -- (0,4) node[left,yshift=-1em] {$f(x)$};%node[above left]
	%Zahlen auf y-Achse
	\foreach \y in {-1,...,-1}
	\draw[shift={(0,\y)},color=black] (2pt,0pt) -- (-2pt,0pt) node[left] {\footnotesize $\y$};
	\foreach \y in {1,...,3}
	\draw[shift={(0,\y)},color=black] (2pt,0pt) -- (-2pt,0pt) node[left] {\footnotesize $\y$};
	%Ursprung
	\draw[color=blue] (-15pt,-10pt) node[right] {\footnotesize Lücke};
	%\draw[color=blue] (-20pt,40pt) node[right] {\footnotesize $f(x)=1$};
	\draw[thick,color=blue] plot[samples=200, domain=-2:-0.1] 
	(\x,{\x*\x});
	
	\draw[thick,color=blue] plot[samples=200, domain=0.1:2] 
	(\x,{\x*\x}) node[below, xshift=1em]{$f(x)$};
	
	\draw [blue] (0,0) circle (.081cm) ;
	%\fill [blue] (0,1) circle (.081cm) ;	
	\end{tikzpicture}\\
\end{center}

\begin{align*}
\lim\limits_ {\substack{x \to 0\\x<0}}f(x)=0=\lim\limits_ {\substack{x \to 0\\x>0}}f(x)\neq f(0) \;\;\text{nicht im Definitionbereich, hier ist eine Lücke }
\end{align*}
In diesen Fall kann man die Unstetigkeit beheben mit:\\
\begin{align*}
f(x) =     \left\{ \begin{array}{rcl}
x^2& \mbox{für} & x\neq0\\
0 & \mbox{für} & x= 0 \\
\end{array}\right.
\end{align*}
\textbf{Typ 2: Sprungstelle}\\
\begin{align*}
\text{Grenzwerte exestieren }\\
\boxed{\lim\limits_ {\substack{x\to x_0\\x<x_0}}f(x)\neq\lim\limits_ {\substack{x\to x_0\\x>x_0}}f(x)}
\end{align*}
Beispiel:\\
Sei: $ f:\mathbb{R}\longmapsto\mathbb{R} $ mit$ f(x) =     \left\{ \begin{array}{rcl}
1 & \mbox{für} & x\geq0\\
-1 & \mbox{für} & x< 0 \\

\end{array}\right.$\\
\begin{center}
	\begin{tikzpicture}
	
	% Gitternetzlinien
	\draw[xstep=1cm, ystep=1cm, very thin, color=lightgray] (-4,4) grid (4,-2);
	% x-Achse
	\draw[->] (-4,0) -- (4,0) node[below,xshift=-1em] {$x$}; 
	%Zahlen auf x-Achse
	\foreach \x in {-3,...,-1}
	\draw[shift={(\x,0)},color=black] (0pt,2pt) -- (0pt,-2pt) node[below] {\footnotesize $\x$};
	\foreach \x in {1,...,3}
	\draw[shift={(\x,0)},color=black] (0pt,2pt) -- (0pt,-2pt) node[below] {\footnotesize $\x$};
	% y-Achse 
	\draw[->] (0,-2) -- (0,4) node[left,yshift=-1em] {$f(x)$};%node[above left]
	%Zahlen auf y-Achse
	\foreach \y in {-1,...,-1}
	\draw[shift={(0,\y)},color=black] (2pt,0pt) -- (-2pt,0pt) node[right] {\footnotesize $\y$};
	\foreach \y in {1,...,3}
	\draw[shift={(0,\y)},color=black] (2pt,0pt) -- (-2pt,0pt) node[left] {\footnotesize $\y$};
	%Ursprung
	\draw[color=blue] (0pt,-10pt) node[right] {\footnotesize $0$};
	
	\draw[thick,color=blue] plot[samples=200, domain=0:4] 
	(\x,{1}) node[below, xshift=-5em] {$f(x)$};
	
	\draw[thick,color=blue] plot[samples=200, domain=-4:0] 
	(\x,{-1});
	\draw[color=blue,thick] (0pt,-1) node[xshift=0em] {\footnotesize $]$};
	
	\draw[thick,color=blue] (0pt,1) node[xshift=0em] {\footnotesize $]$};
	\draw[thick,color=blue] (0pt,1) node[xshift=0em,yshift=0.8em] {\footnotesize Sprungstelle};
	\end{tikzpicture}\\
\end{center}
Betrachte: für $ x_0=0 $\\
\begin{align*}
\lim\limits_ {\substack{x \to 0\\x<0}}f(x)=-1\neq\lim\limits_ {\substack{x \to 0\\x\geq0}}f(x)=1
\end{align*}
\subsubsection{Unstetigkeit-Teil II}
\href{https://youtu.be/8bunAaehx4A}{Video 03}\\
\textbf{Typ 3: Polstelle}\\
\begin{align*}
\boxed{\lim\limits_ {\substack{x\to x_0\\x<x_0}}f(x)=\pm\infty\;\;\text{oder }\lim\limits_ {\substack{x\to x_0\\x>x_0}}f(x)=\pm\infty}
\end{align*}
Beispiel:\\
Sei: $ f:\mathbb{R}\longmapsto\mathbb{R} $ mit$ f(x) =     \left\{ \begin{array}{rcl}
\frac{1}{x} & \mbox{für} & x>0\\
0 & \mbox{für} & x\leq0 \\

\end{array}\right.$\\
\begin{center}
	\begin{tikzpicture}
	
	% Gitternetzlinien
	\draw[xstep=1cm, ystep=1cm, very thin, color=lightgray] (-4,4) grid (4,-2);
	% x-Achse
	\draw[->] (-4,0) -- (4,0) node[below,xshift=-1em] {$x$}; 
	%Zahlen auf x-Achse
	\foreach \x in {-3,...,-1}
	\draw[shift={(\x,0)},color=black] (0pt,2pt) -- (0pt,-2pt) node[below] {\footnotesize $\x$};
	\foreach \x in {1,...,3}
	\draw[shift={(\x,0)},color=black] (0pt,2pt) -- (0pt,-2pt) node[below] {\footnotesize $\x$};
	% y-Achse 
	\draw[->] (0,-2) -- (0,4) node[left,yshift=-1em] {$f(x)$};%node[above left]
	%Zahlen auf y-Achse
	\foreach \y in {-1,...,-1}
	\draw[shift={(0,\y)},color=black] (2pt,0pt) -- (-2pt,0pt) node[left] {\footnotesize $\y$};
	\foreach \y in {1,...,3}
	\draw[shift={(0,\y)},color=black] (2pt,0pt) -- (-2pt,0pt) node[left] {\footnotesize $\y$};
	%Ursprung
	\draw[color=black] (0pt,-10pt) node[right] {\footnotesize $0$};
	
	\draw[thick,color=blue] plot[samples=200, domain=0.25:4] 
	(\x,{1/\x}) node[xshift=-1em,yshift=20] {$f(x)$};
	
	\draw[thick,color=blue] plot[samples=200, domain=-4:0] 
	(\x,{0});
	
	
	
	\end{tikzpicture}\\
	\glqq Ist auch eine unendliche Sprungstelle.\grqq\\
\end{center}
\begin{align*}
&\lim\limits_ {\substack{x \to 0\\x<0}}f(x)=0\;\;\;\;\lim\limits_ {\substack{x \to 0\\x>0}}f(x)=\infty\\
&\Rightarrow f\;\;\text{in }\;\; x_0=0\;\;\text{unstetig }\\
&\text{d.h. }\;\;f\;\;\text{ist nicht stetig }\\
\end{align*}
\textbf{Typ 3: Oszillation}\\
\begin{center}
	\fbox{\parbox  {78mm} {
			$\lim\limits_ {\substack{x \to 0\\x<0}}f(x)$ existiert nicht und $\lim\limits_ {\substack{x \to 0\\x<0}}f(x)\neq\infty$\vspace*{4mm}\\
			$\lim\limits_ {\substack{x \to 0\\x>0}}f(x)$ existiert nicht und $\lim\limits_ {\substack{x \to 0\\x>0}}f(x)\neq\infty$\\}}
\end{center}
Beispiel:\\
Sei: $ f:\mathbb{R}n\longmapsto\mathbb{R}n $ mit $ f(x) =     \left\{ \begin{array}{rcl}
\sin(\frac{1}{x}) & \mbox{für} & x\neq0\\
0 & \mbox{für} & x=0 \\

\end{array}\right. $\\
\begin{center}
	\begin{tikzpicture}
	
	% Gitternetzlinien
	\draw[xstep=1cm, ystep=1cm, very thin, color=lightgray] (-4,2) grid (4,-2);
	% x-Achse
	\draw[->] (-4,0) -- (4,0) node[below,xshift=-1em] {$x$}; 
	%Zahlen auf x-Achse
	\foreach \x in {-3,...,-1}
	\draw[shift={(\x,0)},color=black] (0pt,2pt) -- (0pt,-2pt) node[below] {\footnotesize $\x$};
	\foreach \x in {1,...,3}
	\draw[shift={(\x,0)},color=black] (0pt,2pt) -- (0pt,-2pt) node[below] {\footnotesize $\x$};
	% y-Achse 
	\draw[->] (0,-2) -- (0,2) node[left,yshift=-1em] {$f(x)$};%node[above left]
	%Zahlen auf y-Achse
	\foreach \y in {-1,...,-1}
	\draw[shift={(0,\y)},color=black] (2pt,0pt) -- (-2pt,0pt) node[left] {\footnotesize $\y$};
	\foreach \y in {1,...,1}
	\draw[shift={(0,\y)},color=black] (2pt,0pt) -- (-2pt,0pt) node[left] {\footnotesize $\y$};
	%Ursprung
	\draw[color=black] (0pt,-10pt) node[right] {\footnotesize $0$};
	
	\draw[thick,color=blue,rotate=180] plot[samples=200, domain=0.01:4] 
	(\x,{sin(1/\x r)});
	
	\draw[thick,color=blue] plot[samples=200, domain=0.01:4] 
	(\x,{sin(1/\x r)}) node[xshift=-1em,yshift=1.5em]{$f(x)$};
	
	
	\end{tikzpicture}\\
\end{center}
\begin{align*}
\lim\limits_ {\substack{x \to 0\\x<0}}f(x)=\lim\limits_ {\substack{x \to 0\\x<0}} \sin(\frac{1}{x})\text{existiert nicht! };\;\\
\lim\limits_ {\substack{x \to 0\\x>0}}f(x)=\lim\limits_ {\substack{x \to 0\\x<0}} \sin(\frac{1}{x})\text{existiert nicht! };\;\\
\end{align*}
Der Funktionswert ist immer zwischen -1 und +1 (ähnlich wie $ \lim\limits_ {n\to\infty}(-1)^n $)und der Grenzwert von dem Funktionswert existiert nicht.\\
\begin{align*}
&-1\leq\lim\limits_ {\substack{x \to 0\\x<0}}f(x)=\lim\limits_ {\substack{x \to 0\\x<0}}\leq 1 \;\;\text{und }\lim\limits_ {\substack{x \to 0\\x<0}}f(x)\;\;\text{existiert nicht }\\
&\Rightarrow f \;\;\text{in }\;\;x_0=0\;\;\text{nicht stetig ist. }
\end{align*}
Bemerkung:\\

\textbullet\;\;Typ 2-4 sind nicht behebbar\\

\textbullet\;\;Polynome und rationale Funktionen sind auf ihren Definitionsbereich stetig.\\

\textbullet\;\;Stetigkeitsuntersuchungen werden immer auf zusammengesetzte Funktionen angewandt.\\

\subsection{Rechenregeln für stetge Funktionen}
\href{https://youtu.be/hV8uqmwoxiI}{Video 04}\\
Seien $ f,g $ stetige Funktionen, dann sind\\
\begin{center}
	\fbox{\parbox  {70mm} {\textbullet \;\;$ f+g, f\cdot g, f-g $\vspace*{2mm}\\
			\textbullet \;\; $ \frac{f}{g} $ für $ g(x)\neq 0 $\vspace*{2mm}\\\textbullet \;\;$ f\circ g $, wenn $ f\circ g  $ wohldefiniert ist}}\vspace*{4mm}\\
	stetige Funktionen.
\end{center}
Beispiel:\\
Überprüfe ob:\\
\begin{align*}
|.|:\mathbb{R}\longmapsto \mathbb{R} [0,\infty[\;\;\text{mit }\;\;|x| =     \left\{ \begin{array}{rcl}
-x & \mbox{für} & x<0\\
x & \mbox{für} & x\geq 0 \\
\end{array}\right.\;\;\text{stetig ist. }
\end{align*}
Da der Betrag $ |.| $ für $ x\in \;\; ]-\infty,0[ $ und $ x\in \;\; [0,\infty[  $ als Komposition einer stetigen Abbildung ist, bleibt nur noch die Stelle $ x_0 $ zur Prüfung übrig.\\
Bemerkung: Dies in der Prüfung hinschreiben, sonst gibt es Punktabzug!\\
\begin{align*}
\lim\limits_ {\substack{x\to 0\\x<0}}|x|=\lim\limits_ {x\to 0}-x=0=|0|=\lim\limits_ {\substack{x\to 0\\x>0}}|x|=\lim\limits_ {x\to\ 0}x=0\\
\end{align*}
D.h. der $ |.| $ ist stetig.\\
\subsubsection{Aufgabe 1 zu Rechenregeln von stetigen Funktionen}
\href{https://youtu.be/0dB1Z7ZqMlk}{Video 08}\\
Untersuche folgende Funktion auf Stetigkeit:\\
\begin{align*}
f:\mathbb{R}\longmapsto \mathbb{R}\;\;\text{mit }\;\;f(x):=\left\{ \begin{array}{rcl}
|1-x| & \mbox{für} & x\geq 0\\
-x^2 & \mbox{für} & x< 0 \\
\end{array}\right. 
\end{align*}
\begin{center}
	\begin{tikzpicture}
	
	% Gitternetzlinien
	\draw[xstep=1cm, ystep=1cm, very thin, color=lightgray] (-4,4) grid (4,-4);
	% x-Achse
	\draw[->] (-4,0) -- (4,0) node[below,xshift=-1em] {$x$}; 
	%Zahlen auf x-Achse
	\foreach \x in {-3,...,-1}
	\draw[shift={(\x,0)},color=black] (0pt,2pt) -- (0pt,-2pt) node[below] {\footnotesize $\x$};
	\foreach \x in {1,...,3}
	\draw[shift={(\x,0)},color=black] (0pt,2pt) -- (0pt,-2pt) node[below] {\footnotesize $\x$};
	% y-Achse 
	\draw[->] (0,-4) -- (0,4) node[left,yshift=-1em] {$f(x)$};%node[above left]
	%Zahlen auf y-Achse
	\foreach \y in {-1,...,-3}
	\draw[shift={(0,\y)},color=black] (2pt,0pt) -- (-2pt,0pt) node[left] {\footnotesize $\y$};
	\foreach \y in {1,...,3}
	\draw[shift={(0,\y)},color=black] (2pt,0pt) -- (-2pt,0pt) node[left] {\footnotesize $\y$};
	%Ursprung
	\draw[color=black] (0pt,-10pt) node[right] {\footnotesize $0$};
	
	\draw[thick,color=black] plot[samples=200, domain=-2:0] 
	(\x,{-\x*\x});
	
	\draw[thick,color=black] plot[samples=200, domain=0:1] 
	(\x,{-\x+1});
	
	\draw[thick,color=black] plot[samples=200, domain=1:4] 
	(\x,{\x-1}) node[below, xshift=-3em]{$f(x)$};
	
	\end{tikzpicture}\\	
\end{center}
Frage: Welche Stellen muß ich untersuchen?\\
\begin{quote}
	\textbullet \;\; $ |1-x| =     \left\{ \begin{array}{rcl}
	-1+x & \mbox{für} & x\geq 1\\
	1-x & \mbox{für} & x < 1 \\
	
	\end{array}\right. $\\
	\textbullet \;\; $ x_0=0$\\
\end{quote}
Für $ x\neq 1,0 $ ist $ f $ als Komposition stetiger Funktion stetig.\vspace*{2mm}\\
\newpage
\noindent
\textbullet \;\; Bei der Stelle: $ x_0=1 $:
\begin{quote}
	zu zeigen:\\
	\begin{align*}
	&\lim\limits_ {\substack{x \to 1 \\ x>1}} f(x)=	\lim\limits_ {\substack{x \to 1 \\ x>1}}|1-x|=\lim\limits_ {x\to\ 1}(-1+x)=0\\
	&\lim\limits_ {\substack{x \to 1 \\ x<1}} f(x)=	\lim\limits_ {\substack{x \to 1 \\ x<1}}|1-x|=\lim\limits_ {x\to\ 1}1-x=0\\
	&f(1)=|1-1|=0\\
	&\Rightarrow\\
	&\lim\limits_ {\substack{x \to 1 \\ x>1}} f(x)=\lim\limits_ {\substack{x \to 1 \\ x<1}}f(x)=f(x)\\
	&\Rightarrow f\;\; \text{in }\;\;x_0=1\;\;\text{stetig }
	\end{align*}
	
\end{quote}
\textbullet \;\; Bei der Stelle: $ x_0=0 $:
\begin{quote}
	zu zeigen:\\
	\begin{align*}
	&\lim\limits_ {\substack{x \to 0 \\ x>0}}f(x)=\lim\limits_ {\substack{x \to 0 \\ x>0}}1-x=1\\
	&\lim\limits_ {\substack{x \to 0 \\ x<0}}f(x)=\lim\limits_ {\substack{x \to 0 \\ x<0}}-x^2=0\\
	&\Rightarrow\\
	&\lim\limits_ {\substack{x \to 0 \\ x>0}}f(x)\neq\lim\limits_ {\substack{x \to 0 \\ x<0}}f(x)\\
	&\Rightarrow f\;\; \text{nicht in }\;\;x_0=0\;\;\text{stetig }
	\end{align*}
\end{quote}
Fazit:\\
\begin{quote}
	$ f $ ist für $ x\in \mathbb{R}\backslash\{0\} $ stetig.\\
\end{quote}
\subsubsection{Aufgabe 2 zu Rechenregeln von stetigen Funktionen}
\href{https://youtu.be/hae5ZE-L7QU}{Video 09}\\
Wähle $ \alpha ,\beta \in\mathbb{R} $ , s.d. die folgende Funktion stetig ist:\\
\begin{align*}
f(x) =     \left\{ \begin{array}{rcr}
\mathrm{e}^x & \mbox{für} & x\geq 0\\
\alpha \cos(x)+\beta \sin(x) & \mbox{für} & -\frac{\pi}{2}\leq x < 0 \\
2+(x+\frac{\pi}{2})^2 & \mbox{für} & x< -\frac{\pi}{2}\\
\end{array}\right.
\end{align*}
In den Bereichen:\\
\begin{align*}
&x\in \;\; ]0,\infty[\\
&x\in \;\;]-\frac{\pi}{2},0[ \\
&x\in \;\;]-\infty,-\frac{\pi}{2}[\\
\end{align*}
ist $ f $ als Komposition stetiger Abb. stetig.\\
\newpage
\noindent
\textbullet\;\; Bei der Stelle $ x_0=0 $:\\
\begin{quote}
	zu zeigen:\\
	\begin{align*}
	&\lim\limits_ {\substack{x \to 0 \\ x>0}}f(x)=\lim\limits_ {\substack{x \to 0 \\ x<0}}f(x)=f(0)\\
	&\lim\limits_ {\substack{x \to 0 \\ x>0}}f(x)=\lim\limits_ {x\to 0}\mathrm{e}^x=1\\
	&\lim\limits_ {\substack{x \to 0 \\ x<0}}f(x)=\lim\limits_ {x\to 0}(\alpha \cos(x)+\beta \sin(x))=\alpha\cdot 1+\beta \cdot0=f(0)=\mathrm{e}^0=1\\
	&\Rightarrow \text{für }\;\;\alpha=1\;\text{ ist}\;f \;\;\text{in }\;\;x_0=0\;\;\text{stetig }
	\end{align*}
	
\end{quote}
\textbullet\;\; Bei der Stelle $ x_0=-\frac{\pi}{2} $:\\
\begin{quote}
	zu zeigen:
	\begin{align*}
	&\lim\limits_ {\substack{x \to -\frac{\pi}{2}  \\ x<-\frac{\pi}{2}}}f(x)=\lim\limits_ {\substack{x \to -\frac{\pi}{2} \\ x>-\frac{\pi}{2}}}f(x)=f(-\genfrac{}{}{}{2}{\pi}{2 })\\
	&\lim\limits_ {\substack{x \to -\frac{\pi}{2}  \\ x<-\frac{\pi}{2}}}f(x)=\lim\limits_ {x\to -\frac{\pi}{2}}\big(2+(x+\genfrac{}{}{}{2}{\pi}{2})^2\big)=2\\
	&\lim\limits_ {\substack{x \to -\frac{\pi}{2}  \\ x>-\frac{\pi}{2}}}f(x)=\lim\limits_ {x\to -\frac{\pi}{2}}(\alpha \cos(x)+\beta \sin(x)=\alpha\cdot 0+\beta(-1)=-ß\\
	&f(-\genfrac{}{}{}{2}{\pi }{2})=\alpha \cos(-\genfrac{}{}{}{2}{\pi}{2})+\beta \sin(-\genfrac{}{}{}{2}{\pi}{2})=-\beta\\
	&\Rightarrow \text{für }\;\;\beta=-2\;\text{ ist}\;f \;\;\text{in }\;\;x_0=-\frac{\pi}{2}\;\;\text{stetig }
	\end{align*}
\end{quote}
Fazit:\\
Für $ \alpha=1 $ und $ \beta=-2 $ ist $ f $ stetig.\\
\subsection{Epsilon-Delta Definition}
\href{https://youtu.be/4twSaHZ1G3g}{Video 05}\\
Definition (Epsilon-Delta Kriterium):\vspace*{2mm}\\
Die Funktion $ f:\mathbb{D}\subset \mathbb{R}\longmapsto\mathbb{R} $ ist stetig in $ x_0\in \mathbb{D} $, wenn zu jedem $ \varepsilon > 0 $ ein $ \delta>0 $ existiert, so dass für alle $ x\in \mathbb{D} $ mit $ |x-x_0|<\delta $ gilt $ |f(x)-f(x_0)<\varepsilon $.\\
Oder in math. Kurzschreibweise:\\
\begin{align*}
\forall\varepsilon>0\;\exists\delta>0\;\forall x\in \mathbb{D}:(|x-x_0|<\delta\Rightarrow|f(x)-f(x_0|<\varepsilon)
\end{align*}
\begin{center}
	
	
	\begin{tikzpicture}
	
	% Gitternetzlinien
	\draw[xstep=1cm, ystep=1cm, very thin, color=lightgray] (-4,5) grid (5,-2);
	% x-Achse
	\draw[->] (-4,0) -- (4,0) node[below,xshift=-1em] {$x$}; 
	%Zahlen auf x-Achse
	\foreach \x in {-3,...,-1}
	\draw[shift={(\x,0)},color=black] (0pt,2pt) -- (0pt,-2pt) node[below] {\footnotesize $\x$};
	%\foreach \x in {1,...,3}
	%\draw[shift={(\x,0)},color=black] (0pt,2pt) -- (0pt,-2pt) node[below] {\footnotesize $\x$};
	% y-Achse 
	\draw[->] (0,-2) -- (0,5) node[left,yshift=-1em] {$f(x)$};%node[above left]
	%Zahlen auf y-Achse
	\foreach \y in {-1,...,-1}
	\draw[shift={(0,\y)},color=black] (2pt,0pt) -- (-2pt,0pt) node[left] {\footnotesize $\y$};
	%foreach \y in {1,...,3}
	%draw[shift={(0,\y)},color=black] (2pt,0pt) -- (-2pt,0pt) node[left] {\footnotesize $\y$};
	%Ursprung
	\draw[color=black] (0pt,-10pt) node[right] {\footnotesize $0$};
	
	
	
	\draw[thick,color=blue] plot[samples=200, domain=-4:3.3] 
	(\x,{0.25*(\x+1.5)^2-1}) node[below, xshift=1em]{$f(x)$};
	\draw[color=blue,dashed,thick] (2,0)--(2,2);
	\draw[color=blue] (2,-0.3) node{\footnotesize $x_0$};
	\draw[color=blue,dashed,thick] (0,2)--(2,2);
	\draw[color=blue] (-0.5,2) node{\footnotesize $f(x_0)$};
	\draw[color=red,thick] (0,0.5)--(0,3.5);
	\draw[color=red,thick] (-1,3.5)node{\footnotesize $f(x_0)+\varepsilon$};
	\draw[color=red,thick] (-1,0.5)node{\footnotesize $f(x_0)-\varepsilon$};
	%\draw[color=red,thick,rotate=90] (0,3)node{\footnotesize $[$};
	\node at (0,3.5) [rotate=90,color=red] {$[$};
	\node at (0,0.5) [rotate=-90,color=red] {$[$};
	%\draw[color=red,thick,dashed] (2.5,0)--(2.5,3);
	%\draw[color=red,thick,dashed] (2.5,3)--(0,3);
	%\draw[color=red,thick,dashed] (0,1)--(1.35,1);
	%\draw[color=red,thick,dashed] (1.35,1)--(1.35,0);
	\draw[color=red,thick] (1.3,0)--(2.7,0);
	\node at (1.3,0) [rotate=0,color=red] {$]$};
	\node at (2.7,0) [rotate=0,color=red] {$[$};
	\node at (2.8,-1) [rotate=-90,color=red] {\footnotesize $x_0+\delta$};
	\node at (1.3,-1) [rotate=-90,color=red] {\footnotesize $x_0-\delta$};
	\draw[color=black] (2.5,-0.3) node{\footnotesize $x^\prime$};
	\draw[color=black,thick,dashed] (2.5,0)--(2.5,3);
	\draw[color=black,thick,dashed] (2.5,3)--(0,3);
	\draw[color=black] (-0.5,3) node{\footnotesize $f(x^\prime )$};
	\end{tikzpicture}\\
	Hier gilt:\\
	\begin{align*}
	\forall\varepsilon>0\;\exists\;\delta>0\;\forall x^\prime \in \mathbb{D} :(|x^\prime -x_0|<\delta \Rightarrow
	|f(x^ \prime)-f(x_0)|<\varepsilon) \\
	\end{align*}
\end{center}
\begin{center}
	
	
	\begin{tikzpicture}
	
	% Gitternetzlinien
	\draw[xstep=1cm, ystep=1cm, very thin, color=lightgray] (-4,6) grid (6,-2);
	% x-Achse
	\draw[->] (-4,0) -- (4,0) node[below,xshift=-1em] {$x$}; 
	%Zahlen auf x-Achse
	\foreach \x in {-3,...,-1}
	\draw[shift={(\x,0)},color=black] (0pt,2pt) -- (0pt,-2pt) node[below] {\footnotesize $\x$};
	%\foreach \x in {1,...,3}
	%\draw[shift={(\x,0)},color=black] (0pt,2pt) -- (0pt,-2pt) node[below] {\footnotesize $\x$};
	% y-Achse 
	\draw[->] (0,-2) -- (0,6) node[left,yshift=-1em] {$f(x)$};%node[above left]
	%Zahlen auf y-Achse
	\foreach \y in {-1,...,-1}
	\draw[shift={(0,\y)},color=black] (2pt,0pt) -- (-2pt,0pt) node[left] {\footnotesize $\y$};
	%foreach \y in {1,...,3}
	%draw[shift={(0,\y)},color=black] (2pt,0pt) -- (-2pt,0pt) node[left] {\footnotesize $\y$};
	%Ursprung
	\draw[color=black] (0pt,-10pt) node[right] {\footnotesize $0$};
	\draw[thick,color=blue] plot[samples=200, domain=-4:2] 
	(\x,{0.25*(\x+1.5)^2-1});
	
	
	\draw[thick,color=blue] plot[samples=200, domain=2:3.4] 
	(\x,{0.15*(\x+1.5)^2+2.4}) node[below, xshift=1em]{$f(x)$};
	\draw[color=blue,dashed,thick] (2,0)--(2,2);
	\draw[color=blue] (2,-0.3) node{\footnotesize $x_0$};
	\draw[color=blue,dashed,thick] (0,2)--(2,2);
	\draw[color=blue] (-0.5,2) node{\footnotesize $f(x_0)$};
	\draw[color=red,thick] (0,0.5)--(0,3.5);
	\draw[color=red,thick] (-1,3.5)node{\footnotesize $f(x_0)+\varepsilon$};
	\draw[color=red,thick] (-1,0.5)node{\footnotesize $f(x_0)-\varepsilon$};
	%\draw[color=red,thick,rotate=90] (0,3)node{\footnotesize $[$};
	\node at (0,3.5) [rotate=90,color=red] {$[$};
	\node at (0,0.5) [rotate=-90,color=red] {$[$};
	%\draw[color=red,thick,dashed] (1.3,0)--(1.3,0.8);
	%\draw[color=red,thick,dashed] (2.5,3)--(0,3);
	%\draw[color=red,thick,dashed] (0,1)--(1.35,1);
	%\draw[color=red,thick,dashed] (1.35,1)--(1.35,0);
	\draw[color=red,thick] (1.3,0)--(2.7,0);
	\node at (1.3,0) [rotate=0,color=red] {$]$};
	\node at (2.7,0) [rotate=0,color=red] {$[$};
	\node at (2.8,-1) [rotate=-90,color=red] {\footnotesize $x_0+\delta$};
	\node at (1.3,-1) [rotate=-90,color=red] {\footnotesize $x_0-\delta$};
	\draw[color=black] (2.5,-0.3) node{\footnotesize $x^\prime$};
	\draw[color=black,thick,dashed] (2.5,0)--(2.5,4.8);
	\draw[color=black,thick,dashed] (0,4.8)--(2.5,4.8);
	\draw[color=black] (-0.5,4.8) node{\footnotesize $f(x^\prime )$};
	\end{tikzpicture}\\
	Es gilt nicht mehr:\\
	\begin{align*}
	\forall\varepsilon>0\;\exists\;\delta>0\;\forall x^\prime \in \mathbb{D} :(|x^\prime -x_0|<\delta \Rightarrow
	|f(x^ \prime)-f(x_0)|<\varepsilon) \\
	\end{align*}
	
\end{center}
Bemerkung:\\
Eine Funktion $ f $ heißt stetig in $ \mathbb{D} $, wenn $ f $ stetig in allen $ x_0\in \mathbb{D} $ ist.\\
\subsubsection{Beispiel Epsilon-Delta Definition}
\href{https://youtu.be/1SoRceatdgo}{Video 06}\\
Zeige, dass $ f:\mathbb{R}\longmapsto\mathbb{R} $ mit $ f(x)=x^2 $ in $ x_0=0 $ stetig ist.\vspace*{2mm}\\
Zu zeigen:\\
\begin{align*}
\forall\;\varepsilon>0\;\exists\;\delta >0\;\forall x\in \mathbb{R}:(|x-x_0|<\varepsilon\Rightarrow|f(x)-f(x_0)|<\delta)
\end{align*}
Sei $ \varepsilon>0 $ beliebig:\\
Es gilt für $ x_0=0 $ :\\
\begin{align}
|x-x_0|<\delta \nonumber\\
|x|<\delta 
\end{align}
Und es gilt für $ f(x)=x^2 $ und $ f(x_0)=x_0^2 $\\
\begin{align}
|f(x)-f(x_0)|<\varepsilon \nonumber\\
|x^2-x_0^2|=|x^2|<\varepsilon
\end{align}
In Gleichung (VI.2) Gleichung (VI.1) einsetzen:\\
\begin{align*}
|\delta^2|<\varepsilon\\
\delta<\sqrt{\varepsilon}\\
\end{align*}
Wähle: $ 0<\delta<\sqrt{\varepsilon} $\\
Beweis mit: $ \varepsilon>0 $ mit $ \delta=\sqrt{\varepsilon}\;\;\forall x\in \mathbb{R} $ mit $ |x|<\delta $\\
\begin{align*}
|f(x)-f(x_0)|&=|x^2|\\
&<|\delta^2| \;\;\text{mit } \delta=\sqrt{\varepsilon}\\
&<|\sqrt{\varepsilon}^2|\\|f(x)-f(x_0)|&<\varepsilon_{_{\Square}}
\end{align*}
\subsubsection{Aufgabe Epsilon-Delta Definition}
\href{https://youtu.be/Qnf4afRSK38}{Video 07}\\
Zeige, dass $ f:\mathbb{R}\longmapsto\mathbb{R} $ mit $ f(x)=\frac{x-1}{x^2+1} $ in $ x_0=-1 $ stetig ist.\vspace*{2mm}\\
Zu zeigen:\\
\begin{align*}
\forall\;\varepsilon>0\;\exists\;\delta >0\;\forall x\in \mathbb{R}:(|x-x_0|<\varepsilon\Rightarrow|f(x)-f(x_0)|<\delta)
\end{align*}
Es gilt für $ x_0=-1 $ :\\
\begin{align*}
|x-x_0&|<\delta\\
|x-(-1)|&<\delta\\
|x+1|&<\delta
\end{align*}
\newpage
\noindent
Und es gilt:\\
\begin{align*}
|f(x)-f(x_0)&<\varepsilon\\
\bigg|\frac{x-1}{x^2+1}-\frac{x_0-1}{x_0^2+1}\bigg|&<\varepsilon\\
\bigg|\frac{x-1}{x^2+1}-\frac{-2}{2}\bigg|&<\varepsilon\\
\bigg|\frac{x-1}{x^2+1}+1\bigg|&<\varepsilon\\
\bigg|\frac{x-1+x^2+1}{x^2+1}\bigg|&<\varepsilon\\
\bigg|\frac{x^2+x}{x^2+1}\bigg|&<\varepsilon\\
\bigg|\frac{x(x+1)}{x^2+1}\bigg|&<\varepsilon\\
\bigg|\frac{x}{x^2+1}\bigg|\bigg|x+1\bigg|&<\varepsilon
\end{align*}
Behauptung:\\
\begin{align*}
\bigg|\frac{x}{x^2+1}\bigg|<1\;\;\forall x\in\mathbb{R}
\end{align*}
Beweis:\\
\begin{align*}
&\bigg|\frac{x}{x^2+1}\bigg|<1\\
&\Leftrightarrow |x|<|x^2+1|\;\;\forall x\in\mathbb{R}_{_{\;\;\Square}}
\end{align*}
Zur Stetigkeit:\\
Sei $ \varepsilon>0 $ beliebig und wähle $ \delta=\varepsilon $ dann gilt für:\\
\begin{align*}
|x-x_0|=|x+1|<\delta=\varepsilon
\end{align*}
dass
\begin{align*}
|f(x)-f(x_0)|=\underbrace {\bigg|\frac{x}{x^2+1}\bigg|}_{<1}\underbrace {\bigg|x+1\bigg|}_{<\delta=\varepsilon}<\varepsilon_{\;\;_{\Square}}
\end{align*}
\subsection{Zwischenwertsatz}
\href{https://youtu.be/qJcBXQWDClk}{Video 10}\\
\begin{center}
	\fbox{\parbox  {84mm} {
			Sei $ I\subseteq \mathbb{R} $ ein Intervall und $ f:I\longmapsto\mathbb{R} $ stetig.\\
			Seien $ a,b \in I $ und $ y_0 $ zwischen $ f(a) $ und $ f(b) $.\\
			Dann existiert ein $ x_0 \in [ a,b]   $ mit $ f(x_0)=y_0 $}}
\end{center}
\begin{center}
	
	
	\begin{tikzpicture}
	
	% Gitternetzlinien
	\draw[xstep=1cm, ystep=1cm, very thin, color=lightgray] (-4,5) grid (5,-2);
	% x-Achse
	\draw[->] (-4,0) -- (4,0) node[below,xshift=-1em] {$x$}; 
	%Zahlen auf x-Achse
	\foreach \x in {-3,...,-1}
	\draw[shift={(\x,0)},color=black] (0pt,2pt) -- (0pt,-2pt) node[below] {\footnotesize $\x$};
	%\foreach \x in {1,...,3}
	%\draw[shift={(\x,0)},color=black] (0pt,2pt) -- (0pt,-2pt) node[below] {\footnotesize $\x$};
	% y-Achse 
	\draw[->] (0,-2) -- (0,5) node[left,yshift=-1em] {$f(x)$};%node[above left]
	%Zahlen auf y-Achse
	\foreach \y in {-1,...,-1}
	\draw[shift={(0,\y)},color=black] (2pt,0pt) -- (-2pt,0pt) node[left] {\footnotesize $\y$};
	%foreach \y in {1,...,3}
	%draw[shift={(0,\y)},color=black] (2pt,0pt) -- (-2pt,0pt) node[left] {\footnotesize $\y$};
	%Ursprung
	\draw[color=black] (0pt,-10pt) node[right] {\footnotesize $0$};
	
	
	
	\draw[thick,color=blue] plot[samples=200, domain=-4:3.3] 
	(\x,{0.25*(\x+1.5)^2-1}) node[below, xshift=1em]{$f(x)$};
	\draw[color=blue,dashed,thick] (2,0)--(2,2);
	\draw[color=blue] (2,-0.3) node{\footnotesize $x_0$};
	\draw[color=blue,dashed,thick] (0,2)--(2,2);
	\draw[color=blue] (-1,2) node{\footnotesize $f(x_0)=y_0$};
	\draw[color=red,thick] (0,0.5)--(0,3.5);
	\draw[color=red,thick] (-0.5,3.5)node{\footnotesize $f(b)$};
	\draw[color=red,thick] (-0.5,0.5)node{\footnotesize $f(a)$};
	%\draw[color=red,thick,rotate=90] (0,3)node{\footnotesize $[$};
	
	\draw[color=red,thick,dashed] (0,0.5)--(1.3,0.5);
	\draw[color=red,thick,dashed] (1.3,0.5)--(1.3,0);
	\draw[color=red,thick,dashed] (0,3.5)--(2.7,3.5)--(2.7,0);
	%\draw[color=red,thick,dashed] (1.35,1)--(1.35,0);
	\draw[color=red,thick] (1.3,0)--(2.7,0);
	\node at (1.3,0) [rotate=180,color=red] {$]$};
	\node at (2.7,0) [rotate=180,color=red] {$[$};
	\node at (2.8,-0.5) [rotate=0,color=red] {\footnotesize $b$};
	\node at (1.3,-0.5) [rotate=0,color=red] {\footnotesize $a$};
	
	\end{tikzpicture}\\
	
\end{center}
Liegt ein Wert zwischen $ f(a) $ und $ f(b) $, dann findet man einen Wert der zwischen $ a $ und $ b $ liegt. Natürlich nur bei einer stetigen Abbildung.\\
Den Zwischenwertsatz (ZWS) benötigt man zur Berechnung von Nullstellen einer Funktion, wenn diese Nullstelle nicht algebraisch berechenbar ist.\vspace*{2mm}\\
Beispiel:\\
Betrachte: $ f:\mathbb{R}\longmapsto \mathbb{R} $ mit $ f(x)=x-1 $\\
Zeige, dass $ f(x) $ eine Nullstelle hat.\\
\begin{center}
	
	
	\begin{tikzpicture}
	
	% Gitternetzlinien
	\draw[xstep=1cm, ystep=1cm, very thin, color=lightgray] (-4,4) grid (4,-2);
	% x-Achse
	\draw[->] (-4,0) -- (4,0) node[below,xshift=-1em] {$x$}; 
	%Zahlen auf x-Achse
	\foreach \x in {-3,...,-1}
	\draw[shift={(\x,0)},color=black] (0pt,2pt) -- (0pt,-2pt) node[below] {\footnotesize $\x$};
	\foreach \x in {1,...,3}
	\draw[shift={(\x,0)},color=black] (0pt,2pt) -- (0pt,-2pt) node[below] {\footnotesize $\x$};
	% y-Achse 
	\draw[->] (0,-2) -- (0,4) node[left,yshift=-1em] {$f(x)$};%node[above left]
	%Zahlen auf y-Achse
	\foreach \y in {-1,...,-1}
	\draw[shift={(0,\y)},color=black] (2pt,0pt) -- (-2pt,0pt) node[left] {\footnotesize $\y$};
	\foreach \y in {1,...,3}
	\draw[shift={(0,\y)},color=black] (2pt,0pt) -- (-2pt,0pt) node[left] {\footnotesize $\y$};
	%Ursprung
	\draw[color=blue] (-4pt,-10pt) node[below] {\footnotesize $a$};
	\draw[color=blue] (2,-10pt) node[below] {\footnotesize $b$};
	\draw[color=black] (-4pt,-2pt) node[below] {\footnotesize $0$};
	\draw[color=blue] (-34pt,1) node {\footnotesize $f(b)=$};
	\draw[thick,color=black] plot[samples=200, domain=-1:4] 
	(\x,{\x-1}) node[below, xshift=1em]{$f(x)$};
	\draw[color=blue,thick, dashed] (0,1)--(2,1)--(2,0);
	\draw[color=white,thick] (0,0)--(0,-1);
	\draw[color=blue,thick, dashed] (0,0)--(0,-1);
	\draw[color=blue] (-34pt,-1) node {\footnotesize $f(a)=$};
	\end{tikzpicture}\\
\end{center}
(i)
\begin{quote}
	Offensichtlich ist $ f $ stetig.
\end{quote}
(ii)
\begin{quote}
	Betrachte: $ f(2)=2-1+1 >0$ 
\end{quote}
(ii)
\begin{quote}
	Betrachte: $ f(0)=0-1 <0$ 
\end{quote}
Nach dem ZWS gilt:\\
\begin{align*}
-1=f(0)<y_0=0<f(2)=1
\end{align*}
Es existiert ein $ x_0\in [0,2] $ mit $ f(x_0)=y_0=0 $, d.h. $ x_0 $ ist eine Nullstelle von $ f $.\\
Die Nullstelle algebraisch gerechnet ist $ x_0=1 $\\
\subsubsection{Aufgabe zum Zwischenwertsatz}
\href{https://youtu.be/hsNWvy0z_VE}{Video 11}\\
Zeige, dass die Gleichung $ \sin(x)=1-\frac{\sin(x)}{\cos(x)} $ eine Lösung im Intervall $ [-\frac{\pi}{4},\frac{\pi}{4}] $ hat.\vspace*{2mm}\\
Betrachte mit der Hilfsfunktion $ h(x) $:\\
\begin{align*}
h(x)=1-\frac{\sin(x)}{\cos(x)}-\sin(x)\;\;\text{wenn }\;\;h(x)=0
\end{align*}
Dann gilt:
\begin{align*}
\sin(x)=1-\frac{\sin(x)}{\cos(x)}
\end{align*}
D.h., wir suchen $ x_0 \in [-\frac{\pi}{4},\frac{\pi}{4}]$  mit $ h(x) =0.$\\
(i)
\begin{quote}
	
	\begin{align*}
	h(- \genfrac{}{}{}{2}{\pi }{4})=1-\frac{\sin(-\genfrac{}{}{}{2}{\pi}{4 })}{\cos(-\genfrac{}{}{}{2}{\pi }{4 })}-\sin(-\genfrac{}{}{}{2}{\pi }{4 })=1-\frac{-\genfrac{}{}{}{2}{1}{\sqrt{2}}}{\genfrac{}{}{}{2}{1 }{\sqrt{2} }}-(-\genfrac{}{}{}{2}{1 }{\sqrt{2} })=2+\frac{1}{\sqrt{2}}>0
	\end{align*}
\end{quote}
(ii)
\begin{quote}
	
	\begin{align*}
	h( \genfrac{}{}{}{2}{\pi }{4})=1-\frac{\sin(\genfrac{}{}{}{2}{\pi}{4 })}{\cos(\genfrac{}{}{}{2}{\pi }{4 })}-\sin(\genfrac{}{}{}{2}{\pi }{4 })=1-\frac{\genfrac{}{}{}{2}{1}{\sqrt{2}}}{\genfrac{}{}{}{2}{1 }{\sqrt{2} }}-\frac{1}{\sqrt{2}}=-\frac{1}{\sqrt{2}}<0
	\end{align*}
\end{quote}


Skizze von Cosinus und Sinus:\\
\begin{center}
	
	
	\begin{tikzpicture}[scale=2]
	
	% Gitternetzlinien
	\draw[xstep=1cm, ystep=1cm, very thin, color=lightgray] (-3.5,2) grid (3.5,-2);
	% x-Achse
	\draw[->] (-3.5,0) -- (3.5,0) node[below,xshift=-1em] {$x$}; 
	%Zahlen auf x-Achse
	\foreach \x in {-2,...,-1}
	\draw[shift={(\x,0)},color=black] (0pt,2pt) -- (0pt,-2pt) node[below] {\footnotesize $\x$};
	\foreach \x in {1,...,2}
	\draw[shift={(\x,0)},color=black] (0pt,2pt) -- (0pt,-2pt) node[below] {\footnotesize $\x$};
	% y-Achse 
	\draw[->] (0,-2) -- (0,2) node[left,yshift=-1em] {$f(x)$};%node[above left]
	%Zahlen auf y-Achse
	\foreach \y in {-1,...,-1}
	\draw[shift={(0,\y)},color=black] (2pt,0pt) -- (-2pt,0pt) node[left] {\footnotesize $\y$};
	\foreach \y in {1,...,1}
	\draw[shift={(0,\y)},color=black] (2pt,0pt) -- (-2pt,0pt) node[left] {\footnotesize $\y$};
	%Ursprung
	\draw[color=black] (0pt,-10pt) node[right] {\footnotesize $0$};
	
	\draw[thick,color=blue] plot[samples=200, domain=-3.5:3.5] 
	(\x,{sin(\x r)}) node[xshift=-1em,yshift=4.5em]{$\sin(x)$};
	\draw[thick,color=red] plot[samples=200, domain=-3.5:3.5] 
	(\x,{cos(\x r)}) node[xshift=1em,yshift=1.5em]{$\cos(x)$};
	\draw[color=green] (-3.14/4,-1.2) node {\footnotesize $-\genfrac{}{}{}{2}{\pi }{4 }$};
	\draw[color=green] (3.14/4,-1.2) node {\footnotesize $\genfrac{}{}{}{2}{\pi }{4 }$};
	\draw[color=green] (3.14/4,1)--(3.14/4,-1);
	\draw[color=green] (-3.14/4,1)--(-3.14/4,-1);
	\end{tikzpicture}\\
	\begin{align*}
	&\cos(\genfrac{}{}{}{2}{\pi }{4})=\sin(\genfrac{}{}{}{2}{\pi }{4})=\frac{1}{\sqrt{2}}\\
	&\cos(-\genfrac{}{}{}{2}{\pi }{4})=\frac{1}{\sqrt{2}}\\
	&\sin(-\genfrac{}{}{}{2}{\pi }{4})=-\frac{1}{\sqrt{2}}\\
	\end{align*}
\end{center}
(iii)
\begin{quote}
	$ h(x) $ ist als Komposition stetiger Funktionen stetig auf $ [-\frac{\pi}{4},\frac{\pi}{4}] $\\
\end{quote}
Nach dem ZWS gilt im Intervall $x_0\in [-\frac{\pi}{4},\frac{\pi}{4}]  $ mit $ h(x_0)=y_0=0 $:\\
\begin{align*}
h(-\genfrac{}{}{}{2}{\pi }{4})<y_0=0<h(\genfrac{}{}{}{2}{\pi }{4})\\
\end{align*}
Dann ist $ x_0 $ auch eine Lösung der Ausgangsgleichung.\\
\subsection{Extrema von Funktionen}
\href{https://youtu.be/la8s-zLqf5w}{Video 12}\\
\underline{Definition(Maximum, Minimum):}\vspace*{2mm}\\
(i) $ f $ nimmt in $ x_0 $ sein globales Maximum $ y_0 $ an, wenn gilt:\\
\begin{align*}
f(x)=y_0\geq f(x)\;\;\text{für alle  } \;\;x\in\mathbb{D}_f
\end{align*}
(ii) $ f $ nimmt in $ x_0 $ sein lokales Minimum $ y_0 $ an, wenn es:\\
\begin{align*}
\varepsilon>0 \;\;\text{gibt mit }\;\;f(x_0)\geq f(x)\;\;\text{für alle }\;\;x\in\mathbb{D}_f\;\;\text{mit }\;\;|x-x_0|<\varepsilon
\end{align*}
Beispiel:\\
\begin{center}
	Globales Maximum\vspace*{2mm}\\	
	
	\begin{tikzpicture}
	
	% Gitternetzlinien
	\draw[xstep=1cm, ystep=1cm, very thin, color=lightgray] (-4,4) grid (4,-2);
	% x-Achse
	\draw[->] (-4,0) -- (4,0) node[below,xshift=-1em] {$x$}; 
	%Zahlen auf x-Achse
	\foreach \x in {-3,...,-1}
	\draw[shift={(\x,0)},color=black] (0pt,2pt) -- (0pt,-2pt) node[below] {\footnotesize $\x$};
	\foreach \x in {1,...,3}
	\draw[shift={(\x,0)},color=black] (0pt,2pt) -- (0pt,-2pt) node[below] {\footnotesize $\x$};
	% y-Achse 
	\draw[->] (0,-2) -- (0,4) node[left,yshift=-1em] {$f(x)$};%node[above left]
	%Zahlen auf y-Achse
	\foreach \y in {-1,...,-1}
	\draw[shift={(0,\y)},color=black] (2pt,0pt) -- (-2pt,0pt) node[left] {\footnotesize $\y$};
	\foreach \y in {1,...,3}
	\draw[shift={(0,\y)},color=black] (2pt,0pt) -- (-2pt,0pt) node[left] {\footnotesize $\y$};
	%Ursprung
	%\draw[color=black] (0pt,-10pt) node[right] {\footnotesize $0$};
	
	\draw[thick,color=blue] plot[samples=200, domain=-2:2] 
	(\x,{-\x*\x +2});
	\draw[color=blue] (-1.5,2.5) node[right] {\footnotesize $globales Maximum f(x_0)=y_0$};
	\draw[color=blue] (-0.5,-10pt) node[right] {\footnotesize $x_0=0$};
	
	
	
	
	\end{tikzpicture}\\
\end{center}

\begin{center}
	Lokales Maximum\vspace*{2mm}\\	
	
	\begin{tikzpicture}
	
	% Gitternetzlinien
	\draw[xstep=1cm, ystep=1cm, very thin, color=lightgray] (-4,4) grid (4,-2);
	% x-Achse
	\draw[->] (-4,0) -- (4,0) node[below,xshift=-1em] {$x$}; 
	%Zahlen auf x-Achse
	\foreach \x in {-3,...,-1}
	\draw[shift={(\x,0)},color=black] (0pt,2pt) -- (0pt,-2pt) node[below] {\footnotesize $\x$};
	\foreach \x in {1,...,3}
	\draw[shift={(\x,0)},color=black] (0pt,2pt) -- (0pt,-2pt) node[below] {\footnotesize $\x$};
	% y-Achse 
	\draw[->] (0,-2) -- (0,4) node[left,yshift=-1em] {$f(x)$};%node[above left]
	%Zahlen auf y-Achse
	\foreach \y in {-1,...,-1}
	\draw[shift={(0,\y)},color=black] (2pt,0pt) -- (-2pt,0pt) node[left] {\footnotesize $\y$};
	\foreach \y in {1,...,3}
	\draw[shift={(0,\y)},color=black] (2pt,0pt) -- (-2pt,0pt) node[left] {\footnotesize $\y$};
	%Ursprung
	%\draw[color=black] (0pt,-10pt) node[right] {\footnotesize $0$};
	
	\draw[thick,color=blue] plot[samples=200, domain=-4:0] 
	(\x,{sin(2*\x r)});
	\draw[thick,color=blue] plot[samples=200, domain=0:3] 
	(\x,{2*sin(3*\x r)});
	\draw[color=blue] (-1.5,2.5) node[right] {\footnotesize $lokales Maximum f(x_0)=y_0$};
	\draw[color=red] (-3,1) node[left] {\footnotesize $x_0-\varepsilon$};
	\draw[color=red] (-1.8,1) node[right] {\footnotesize $x_0+\varepsilon$};
	\draw[color=red,thick,dashed] (-3,-1)--(-3,1) ;
	\draw[color=red,thick,dashed] (-1.8,-1)--(-1.8,1) ;
	\draw[color=red,thick,dashed] (-2.4,-1)--(-2.4,1) ;
	\draw[color=red] (-2.4,1.2) node {\footnotesize $x_0$};
	\end{tikzpicture}\\
\end{center}
Analog definiert man globales und lokales Minimum.\\
\begin{center}
	
	
	\begin{tikzpicture}
	
	% Gitternetzlinien
	\draw[xstep=1cm, ystep=1cm, very thin, color=lightgray] (-4,4) grid (4,-3);
	% x-Achse
	\draw[->] (-4,0) -- (4,0) node[below,xshift=-1em] {$x$}; 
	%Zahlen auf x-Achse
	\foreach \x in {-3,...,-1}
	\draw[shift={(\x,0)},color=black] (0pt,2pt) -- (0pt,-2pt) node[below] {\footnotesize $\x$};
	\foreach \x in {1,...,3}
	\draw[shift={(\x,0)},color=black] (0pt,2pt) -- (0pt,-2pt) node[below] {\footnotesize $\x$};
	% y-Achse 
	\draw[->] (0,-3) -- (0,4) node[left,yshift=-1em] {$f(x)$};%node[above left]
	%Zahlen auf y-Achse
	\foreach \y in {-1,...,-1}
	\draw[shift={(0,\y)},color=black] (2pt,0pt) -- (-2pt,0pt) node[left] {\footnotesize $\y$};
	\foreach \y in {1,...,3}
	\draw[shift={(0,\y)},color=black] (2pt,0pt) -- (-2pt,0pt) node[left] {\footnotesize $\y$};
	%Ursprung
	%\draw[color=black] (0pt,-10pt) node[right] {\footnotesize $0$};
	
	\draw[thick,color=blue] plot[samples=200, domain=-4:0] 
	(\x,{sin(2*\x r)});
	\draw[thick,color=blue] plot[samples=200, domain=0:3] 
	(\x,{2*sin(3*\x r)});
	\draw[color=red] (-2.4,-1.2) node {\footnotesize $lokale \; Minima$};
	\draw[color=red] (-2.4,1.2) node {\footnotesize $lokales\; Maximum$};
	\draw[color=red] (1.8,2.5) node {\footnotesize $globale\; Maxima$};
	\draw[color=red] (2,-2.7) node {\footnotesize $globales\; Minimum$};
	\end{tikzpicture}\\
\end{center}
\begin{center}
	
	kein globales Maximum\\
	\begin{tikzpicture}
	
	% Gitternetzlinien
	\draw[xstep=1cm, ystep=1cm, very thin, color=lightgray] (-4,4) grid (4,-3);
	% x-Achse
	\draw[->] (-4,0) -- (4,0) node[below,xshift=-1em] {$x$}; 
	%Zahlen auf x-Achse
	\foreach \x in {-3,...,-1}
	\draw[shift={(\x,0)},color=black] (0pt,2pt) -- (0pt,-2pt) node[below] {\footnotesize $\x$};
	\foreach \x in {1,...,3}
	\draw[shift={(\x,0)},color=black] (0pt,2pt) -- (0pt,-2pt) node[below] {\footnotesize $\x$};
	% y-Achse 
	\draw[->] (0,-3) -- (0,4) node[left,yshift=-1em] {$f(x)$};%node[above left]
	%Zahlen auf y-Achse
	\foreach \y in {-1,...,-1}
	\draw[shift={(0,\y)},color=black] (2pt,0pt) -- (-2pt,0pt) node[left] {\footnotesize $\y$};
	\foreach \y in {1,...,3}
	\draw[shift={(0,\y)},color=black] (2pt,0pt) -- (-2pt,0pt) node[left] {\footnotesize $\y$};
	%Ursprung
	%\draw[color=black] (0pt,-10pt) node[right] {\footnotesize $0$};
	\draw[thick,color=blue] plot[samples=200, domain=-2.85:-1.45] 
	(\x,{-3*\x -4.6)});
	\draw[thick,color=blue] plot[samples=200, domain=-1.5:0] 
	(\x,{sin(2*\x r)});
	\draw[thick,color=blue] plot[samples=200, domain=0:3] 
	(\x,{2*sin(3*\x r)});
	\draw[color=red] (-2.4,-1.2) node {\footnotesize $lokale \; Minima$};
	%\draw[color=red] (-2.4,1.2) node {\footnotesize $lokales\; Maximum$};
	\draw[color=red] (1.8,2.5) node {\footnotesize $lokale\; Maxima$};
	\draw[color=red] (2,-2.7) node {\footnotesize $globales\; Minimum$};
	\end{tikzpicture}\\
\end{center}
\begin{center}
	
	kein globales Maximum und Minimum\\
	\begin{tikzpicture}
	
	% Gitternetzlinien
	\draw[xstep=1cm, ystep=1cm, very thin, color=lightgray] (-4,4) grid (4,-3);
	% x-Achse
	\draw[->] (-4,0) -- (4,0) node[below,xshift=-1em] {$x$}; 
	%Zahlen auf x-Achse
	\foreach \x in {-3,...,-1}
	\draw[shift={(\x,0)},color=black] (0pt,2pt) -- (0pt,-2pt) node[below] {\footnotesize $\x$};
	\foreach \x in {1,...,3}
	\draw[shift={(\x,0)},color=black] (0pt,2pt) -- (0pt,-2pt) node[below] {\footnotesize $\x$};
	% y-Achse 
	\draw[->] (0,-3) -- (0,4) node[left,yshift=-1em] {$f(x)$};%node[above left]
	%Zahlen auf y-Achse
	\foreach \y in {-1,...,-1}
	\draw[shift={(0,\y)},color=black] (2pt,0pt) -- (-2pt,0pt) node[left] {\footnotesize $\y$};
	\foreach \y in {1,...,3}
	\draw[shift={(0,\y)},color=black] (2pt,0pt) -- (-2pt,0pt) node[left] {\footnotesize $\y$};
	%Ursprung
	%\draw[color=black] (0pt,-10pt) node[right] {\footnotesize $0$};
	\draw[thick,color=blue] plot[samples=200, domain=-2.85:-1.45] 
	(\x,{-3*\x -4.6)});
	\draw[thick,color=blue] plot[samples=200, domain=-1.5:0] 
	(\x,{sin(2*\x r)});
	\draw[thick,color=blue] plot[samples=200, domain=0:3] 
	(\x,{2*sin(3*\x r)});
	\draw[thick,color=blue] plot[samples=200, domain=3:4] 
	(\x,{-3*\x +9.9)});
	\draw[color=red] (-2.4,-1.2) node {\footnotesize $lokales \; Minimum$};
	%\draw[color=red] (-2.4,1.2) node {\footnotesize $lokales\; Maximum$};
	\draw[color=red] (1.8,2.5) node {\footnotesize $lokale\; Maxima$};
	\draw[color=red] (2,-2.7) node {\footnotesize $lokales\; Minimum$};
	\end{tikzpicture}\\
\end{center}
\pagebreak

\begin{center}
	kein globales bzw. lokales Maximum/Minimum\\	
	\begin{tikzpicture}
	
	% Gitternetzlinien
	\draw[xstep=1cm, ystep=1cm, very thin, color=lightgray] (-4,4) grid (4,-2);
	% x-Achse
	\draw[->] (-4,0) -- (4,0) node[below,xshift=-1em] {$x$}; 
	%Zahlen auf x-Achse
	\foreach \x in {-3,...,-1}
	\draw[shift={(\x,0)},color=black] (0pt,2pt) -- (0pt,-2pt) node[below] {\footnotesize $\x$};
	\foreach \x in {1,...,3}
	\draw[shift={(\x,0)},color=black] (0pt,2pt) -- (0pt,-2pt) node[below] {\footnotesize $\x$};
	% y-Achse 
	\draw[->] (0,-2) -- (0,4) node[left,yshift=-1em] {$f(x)$};%node[above left]
	%Zahlen auf y-Achse
	\foreach \y in {-1,...,-1}
	\draw[shift={(0,\y)},color=black] (2pt,0pt) -- (-2pt,0pt) node[left] {\footnotesize $\y$};
	\foreach \y in {1,...,3}
	\draw[shift={(0,\y)},color=black] (2pt,0pt) -- (-2pt,0pt) node[left] {\footnotesize $\y$};
	%Ursprung
	\draw[color=black] (0pt,-10pt) node[right] {\footnotesize $0$};
	
	%\draw[thick,color=black] plot[samples=200, domain=2:4] 
	%(\x,{2}) node[below, xshift=-5em] {$h(x)$};
	
	\draw[thick,color=blue] plot[samples=200, domain=-4:4] 
	(\x,{0.5*\x+0.5}) node[below, xshift=1em]{$f(x)$};
	
	
	
	\end{tikzpicture}\\
\end{center}
\begin{center}
	jeder Punkt nimmt globales bzw. lokales Max./Min. ein\\	
	
	\begin{tikzpicture}
	
	% Gitternetzlinien
	\draw[xstep=1cm, ystep=1cm, very thin, color=lightgray] (-4,4) grid (4,-2);
	% x-Achse
	\draw[->] (-4,0) -- (4,0) node[below,xshift=-1em] {$x$}; 
	%Zahlen auf x-Achse
	\foreach \x in {-3,...,-1}
	\draw[shift={(\x,0)},color=black] (0pt,2pt) -- (0pt,-2pt) node[below] {\footnotesize $\x$};
	\foreach \x in {1,...,3}
	\draw[shift={(\x,0)},color=black] (0pt,2pt) -- (0pt,-2pt) node[below] {\footnotesize $\x$};
	% y-Achse 
	\draw[->] (0,-2) -- (0,4) node[left,yshift=-1em] {$f(x)$};%node[above left]
	%Zahlen auf y-Achse
	\foreach \y in {-1,...,-1}
	\draw[shift={(0,\y)},color=black] (2pt,0pt) -- (-2pt,0pt) node[left] {\footnotesize $\y$};
	\foreach \y in {1,...,3}
	\draw[shift={(0,\y)},color=black] (2pt,0pt) -- (-2pt,0pt) node[left] {\footnotesize $\y$};
	%Ursprung
	\draw[color=black] (0pt,-10pt) node[right] {\footnotesize $0$};
	
	\draw[thick,color=blue] plot[samples=200, domain=-4:4] 
	(\x,{2}) node[below, xshift=-5em] {$f(x)$};
	
	
	
	
	\end{tikzpicture}\\
\end{center}
\begin{center}
	\begin{align*}
	f:[-2,2]\longmapsto[-2,2]\;\;\text{mit }\;\;f(x)=x
	\end{align*}
	\begin{tikzpicture}
	
	% Gitternetzlinien
	\draw[xstep=1cm, ystep=1cm, very thin, color=lightgray] (-4,4) grid (4,-4);
	% x-Achse
	\draw[->] (-4,0) -- (4,0) node[below,xshift=-1em] {$x$}; 
	%Zahlen auf x-Achse
	\foreach \x in {-3,...,-1}
	\draw[shift={(\x,0)},color=black] (0pt,2pt) -- (0pt,-2pt) node[below] {\footnotesize $\x$};
	\foreach \x in {1,...,3}
	\draw[shift={(\x,0)},color=black] (0pt,2pt) -- (0pt,-2pt) node[below] {\footnotesize $\x$};
	% y-Achse 
	\draw[->] (0,-4) -- (0,4) node[left,yshift=-1em] {$f(x)$};%node[above left]
	%Zahlen auf y-Achse
	\foreach \y in {-1,...,-3}
	\draw[shift={(0,\y)},color=black] (2pt,0pt) -- (-2pt,0pt) node[left] {\footnotesize $\y$};
	\foreach \y in {1,...,3}
	\draw[shift={(0,\y)},color=black] (2pt,0pt) -- (-2pt,0pt) node[left] {\footnotesize $\y$};
	%Ursprung
	\draw[color=black] (0pt,-10pt) node[right] {\footnotesize $0$};
	
	
	
	\draw[thick,color=blue] plot[samples=200, domain=-2:2] 
	(\x,{\x}) node[below, xshift=1em]{$f(x)$};
	
	\draw[color=red] (-1,-2.7) node {\footnotesize $globales,\;\;lokales\; Minimum$};
	\draw[color=red] (1.5,2.7) node {\footnotesize $globales,\;\;lokales\; Maximum$};
	\end{tikzpicture}\\
\end{center}
\pagebreak
\begin{center}
	\begin{align*}
	f:]-2,2]\longmapsto\;]-2,2]\;\;\text{mit }\;\;f(x)=x
	\end{align*}
	\begin{tikzpicture}
	
	% Gitternetzlinien
	\draw[xstep=1cm, ystep=1cm, very thin, color=lightgray] (-4,4) grid (4,-4);
	% x-Achse
	\draw[->] (-4,0) -- (4,0) node[below,xshift=-1em] {$x$}; 
	%Zahlen auf x-Achse
	\foreach \x in {-3,...,-1}
	\draw[shift={(\x,0)},color=black] (0pt,2pt) -- (0pt,-2pt) node[below] {\footnotesize $\x$};
	\foreach \x in {1,...,3}
	\draw[shift={(\x,0)},color=black] (0pt,2pt) -- (0pt,-2pt) node[below] {\footnotesize $\x$};
	% y-Achse 
	\draw[->] (0,-4) -- (0,4) node[left,yshift=-1em] {$f(x)$};%node[above left]
	%Zahlen auf y-Achse
	\foreach \y in {-1,...,-3}
	\draw[shift={(0,\y)},color=black] (2pt,0pt) -- (-2pt,0pt) node[left] {\footnotesize $\y$};
	\foreach \y in {1,...,3}
	\draw[shift={(0,\y)},color=black] (2pt,0pt) -- (-2pt,0pt) node[left] {\footnotesize $\y$};
	%Ursprung
	\draw[color=black] (0pt,-10pt) node[right] {\footnotesize $0$};
	
	
	
	\draw[thick,color=blue] plot[samples=200, domain=-2:2] 
	(\x,{\x}) node[below, xshift=1em]{$f(x)$};
	
	\draw[color=red] (0,-2.7) node {\footnotesize $kein\;\; globales,\;\;lokales\; Minimum$};
	\draw[color=red] (1.5,2.7) node {\footnotesize $globales,\;\;lokales\; Maximum$};
	\end{tikzpicture}\\
\end{center}
\subsection{ Extrema stetiger Funktionen}
\href{https://youtu.be/mIDbM-qesZs}{Video 13}\\
\underline{Satz}\\
Seien $ a,b\in\mathbb{R} $ und $ f:[a,b] \longmapsto \mathbb{W}_f\subseteq \mathbb{R}$, so besitzt $ f $ ein Maximum und ein Minimum.\vspace*{2mm}\\
Wenn eine stetige Funktion mit einem kompakten Intervall definiert ist, dann hat es mindestens ein Minumum und mindestens ein Maximum.\vspace*{2mm}\\
Beispiel:\\
$\blacktriangleright \;\; f:[-2,2] $ mit $ x \longmapsto x $\\
\begin{quote}
	Da $ f $ stetig ist und $[-2,2]  $ ein kompaktes Intervall besitzt hat die Funktion ein Maximum und ein Minimum.\\
	
	\pagebreak
	\begin{center}
		\begin{align*}
		f:[-2,2]\longmapsto\;[-2,2]\;\;\text{mit }\;\;f(x)=x
		\end{align*}
		\begin{tikzpicture}
		
		% Gitternetzlinien
		\draw[xstep=1cm, ystep=1cm, very thin, color=lightgray] (-4,4) grid (4,-4);
		% x-Achse
		\draw[->] (-4,0) -- (4,0) node[below,xshift=-1em] {$x$}; 
		%Zahlen auf x-Achse
		\foreach \x in {-3,...,-1}
		\draw[shift={(\x,0)},color=black] (0pt,2pt) -- (0pt,-2pt) node[below] {\footnotesize $\x$};
		\foreach \x in {1,...,3}
		\draw[shift={(\x,0)},color=black] (0pt,2pt) -- (0pt,-2pt) node[below] {\footnotesize $\x$};
		% y-Achse 
		\draw[->] (0,-4) -- (0,4) node[left,yshift=-1em] {$f(x)$};%node[above left]
		%Zahlen auf y-Achse
		\foreach \y in {-1,...,-3}
		\draw[shift={(0,\y)},color=black] (2pt,0pt) -- (-2pt,0pt) node[left] {\footnotesize $\y$};
		\foreach \y in {1,...,3}
		\draw[shift={(0,\y)},color=black] (2pt,0pt) -- (-2pt,0pt) node[left] {\footnotesize $\y$};
		%Ursprung
		\draw[color=black] (0pt,-10pt) node[right] {\footnotesize $0$};
		
		\draw[color=red,thick] (-2,0)--(2,0);
		\draw[color=red,thick] (2,0) node {]};
		\draw[color=red,thick] (-2,0) node {[};
		\draw[thick,color=blue] plot[samples=200, domain=-2:2] 
		(\x,{\x}) node[below, xshift=1em]{$f(x)$};
		
		\draw[color=red] (-1.5,-2.7) node {\footnotesize $globales,\;\; Minimum$};
		\draw[color=red] (1.5,2.7) node {\footnotesize $globales\;\; Maximum$};
		\end{tikzpicture}\\
	\end{center}
\end{quote}
$\blacktriangleright \;\; f:[-2,2] $ mit $ x \longmapsto 1 $\\
\begin{quote}
	Jeder Punkt nimmt globales bzw. lokales Max./Min. ein.\\	
	\begin{center}
		
		\begin{align*}
		f:[-2,2]\longmapsto\;[1]\;\;\text{mit }\;\;f(x)=1
		\end{align*}
		\begin{tikzpicture}
		
		% Gitternetzlinien
		\draw[xstep=1cm, ystep=1cm, very thin, color=lightgray] (-4,4) grid (4,-2);
		% x-Achse
		\draw[->] (-4,0) -- (4,0) node[below,xshift=-1em] {$x$}; 
		%Zahlen auf x-Achse
		\foreach \x in {-3,...,-1}
		\draw[shift={(\x,0)},color=black] (0pt,2pt) -- (0pt,-2pt) node[below] {\footnotesize $\x$};
		\foreach \x in {1,...,3}
		\draw[shift={(\x,0)},color=black] (0pt,2pt) -- (0pt,-2pt) node[below] {\footnotesize $\x$};
		% y-Achse 
		\draw[->] (0,-2) -- (0,4) node[left,yshift=-1em] {$f(x)$};%node[above left]
		%Zahlen auf y-Achse
		\foreach \y in {-1,...,-1}
		\draw[shift={(0,\y)},color=black] (2pt,0pt) -- (-2pt,0pt) node[left] {\footnotesize $\y$};
		\foreach \y in {1,...,3}
		\draw[shift={(0,\y)},color=black] (2pt,0pt) -- (-2pt,0pt) node[left] {\footnotesize $\y$};
		%Ursprung
		\draw[color=black] (0pt,-10pt) node[right] {\footnotesize $0$};
		
		\draw[thick,color=blue] plot[samples=200, domain=-2:2] 
		(\x,{1}) node[below, xshift=-5em] {$f(x)$};
		
		\draw[color=red,thick] (-2,0)--(2,0);
		\draw[color=red,thick] (2,0) node {]};
		\draw[color=red,thick] (-2,0) node {[};
		
		
		\end{tikzpicture}\\
	\end{center}
\end{quote}
$\blacktriangleright \;\; f:[-1,0[ $ mit $ x \longmapsto \frac{1}{x} $\\
\begin{quote}
	Hat ein globales Maximum und kein Minimum.\\
	\pagebreak
	
	\begin{center}
		
		\begin{align*}
		f:]-1,0[\;\longmapsto\;]-1,-\infty[\;\;\text{mit }\;\;f(x)=\frac{1}{x}
		\end{align*}
		\begin{tikzpicture}
		
		% Gitternetzlinien
		\draw[xstep=1cm, ystep=1cm, very thin, color=lightgray] (-4,2) grid (4,-5);
		% x-Achse
		\draw[->] (-4,0) -- (4,0) node[below,xshift=-1em] {$x$}; 
		%Zahlen auf x-Achse
		\foreach \x in {-3,...,-1}
		\draw[shift={(\x,0)},color=black] (0pt,2pt) -- (0pt,-2pt) node[below] {\footnotesize $\x$};
		\foreach \x in {1,...,3}
		\draw[shift={(\x,0)},color=black] (0pt,2pt) -- (0pt,-2pt) node[below] {\footnotesize $\x$};
		% y-Achse 
		\draw[->] (0,-5) -- (0,2) node[left,yshift=-1em] {$f(x)$};%node[above left]
		%Zahlen auf y-Achse
		\foreach \y in {-1,...,-4}
		\draw[shift={(0,\y)},color=black] (2pt,0pt) -- (-2pt,0pt) node[left] {\footnotesize $\y$};
		\foreach \y in {1,...,1}
		\draw[shift={(0,\y)},color=black] (2pt,0pt) -- (-2pt,0pt) node[left] {\footnotesize $\y$};
		%Ursprung
		\draw[color=black] (0pt,-10pt) node[right] {\footnotesize $0$};
		
		\draw[thick,color=blue] plot[samples=200, domain=-1:-0.2] 
		(\x,{1/ \x}) node[left, xshift=1, yshift=10] {$f(x)$};
		
		\draw[color=red,thick] (-1,0)--(0,0);
		\draw[color=red,thick] (0,0) node {[};
		\draw[color=red,thick] (-1,0) node {[};
		\draw[color=red,thick] (-2,-1) node {\footnotesize glob. Max.};
		
		\end{tikzpicture}\\
	\end{center}
\end{quote}
Kein Minimum weil:\\
\begin{align*}
\lim\limits_ {\substack{x \to 0 \\ x<0}}\frac{1}{x}=-\infty\;\;\text{es existiert kein Minimum }
\end{align*}
Maximum weil:\\
\begin{align*}
\text{ Maximum ist bei} \;\;x_0=-1
\end{align*}
$\blacktriangleright \;\; f:]-1,0[ $ mit $ x \longmapsto \frac{1}{x} $\\
\begin{quote}
	Hat kein globales Maximum und kein Minimum.\\
	%\pagebreak
	
	\begin{center}
		
		\begin{align*}
		f:]-1,0[\;\longmapsto\;]-1,-\infty[\;\;\text{mit }\;\;f(x)=\frac{1}{x}
		\end{align*}
		\begin{tikzpicture}
		
		% Gitternetzlinien
		\draw[xstep=1cm, ystep=1cm, very thin, color=lightgray] (-4,2) grid (4,-5);
		% x-Achse
		\draw[->] (-4,0) -- (4,0) node[below,xshift=-1em] {$x$}; 
		%Zahlen auf x-Achse
		\foreach \x in {-3,...,-1}
		\draw[shift={(\x,0)},color=black] (0pt,2pt) -- (0pt,-2pt) node[below] {\footnotesize $\x$};
		\foreach \x in {1,...,3}
		\draw[shift={(\x,0)},color=black] (0pt,2pt) -- (0pt,-2pt) node[below] {\footnotesize $\x$};
		% y-Achse 
		\draw[->] (0,-5) -- (0,2) node[left,yshift=-1em] {$f(x)$};%node[above left]
		%Zahlen auf y-Achse
		\foreach \y in {-1,...,-4}
		\draw[shift={(0,\y)},color=black] (2pt,0pt) -- (-2pt,0pt) node[left] {\footnotesize $\y$};
		\foreach \y in {1,...,1}
		\draw[shift={(0,\y)},color=black] (2pt,0pt) -- (-2pt,0pt) node[left] {\footnotesize $\y$};
		%Ursprung
		\draw[color=black] (0pt,-10pt) node[right] {\footnotesize $0$};
		
		\draw[thick,color=blue] plot[samples=200, domain=-1:-0.2] 
		(\x,{1/ \x}) node[left, xshift=1, yshift=10] {$f(x)$};
		
		\draw[color=red,thick] (-1,0)--(0,0);
		\draw[color=red,thick] (0,0) node {[};
		\draw[color=red,thick] (-1,0) node {]};
		%\draw[color=red,thick] (-2,-1) node {\footnotesize glob. Max.};
		
		\end{tikzpicture}\\
	\end{center}
	
	Kein Minimum weil:\\
	\begin{align*}
	\lim\limits_ {\substack{x \to 0 \\ x<0}}\frac{1}{x}=-\infty\;\;\text{es existiert kein Minimum }
	\end{align*}
	Kein Maximum weil:\\
	\begin{quote}
		
		\begin{align*}
		\lim\limits_ {\substack{x \to -1 \\ x>-1}}\frac{1}{x}=-1
		\end{align*}
		Aber $ -1\neq \mathbb{W}_f $ , deshalb existiert kein Maximun.\\
	\end{quote}
\end{quote}
\subsubsection{Aufgabe zu Extrema stetiger Funktionen}
\href{https://youtu.be/CkySz6nT0R8}{Video 14}\\
Haben folgende Funktionen globale Extrema?\\
(i)\;\;$ f:[2,17]\longmapsto \mathbb{R} $ mit $ f(x)=\frac{2x}{\mathrm{e}^x}-\cos(x) $\vspace*{2mm}\\
(ii)\;$ g:]0,1[\;\cup\;]1,2]\longmapsto \mathbb{R} $ mit $ g(x)= \frac{12}{1-x} $\vspace*{2mm}\\
zu (i):\\
\begin{quote}
	Die Funktion $ f $ ist als Komposition stetiger Funktionen in dem kompakten Intervall $[2,17]  $ stetig, also hat $ f $ ein globales Maximum und globales Minimum.
\end{quote}
zu (ii):\\
\begin{quote}
	Betrachte einseitige Grenzwerte:\\
	\begin{align*}
	\lim\limits_ {\substack{x \to 1 \\ x<1}}g(x)=\lim\limits_ {\substack{x \to 1 \\ x<1}} \frac{12}{1-x}=\infty
	\end{align*}
	Wenn $ x $ gegen 1 geht wird der Nenner ganz klein und der Bruch geht in das Unendliche.\\$ \Rightarrow $kein globales Maximum.
	\begin{align*}
	\lim\limits_ {\substack{x \to 1 \\ x>1}}g(x)=\lim\limits_ {\substack{x \to 1 \\ x>1}} \frac{12}{1-x}=-\infty
	\end{align*}
	Wenn $ x $ gegen gößer 1 zu 1 geht wird der Nenner negativ und ganz klein und der Bruch geht dannn in das negative Unendliche.\\
	$\Rightarrow $ kein globales Minimum.
\end{quote}
\subsection{Infimum und Supremum}
\href{https://youtu.be/z6nV00vtIVY}{Video 15}\\
\underline{Definition:}\\
Sei $ f:\mathbb{D}_f\longmapsto \mathbb{W}_f $ eine Funktion, $ x_0\in \mathbb{D}_f $ sowie $ y_0\in \mathbb{W}_f\cup \{+\infty\} $ (bzw.$ \mathbb{W}_f\cup \{-\infty\} $).\vspace*{2mm}\\
$ y_0 $ heißt Supremum (bzw. Infimum) von $ f $, wenn gilt:\\
(i) \;$ f(x)\leq y_0 $ (bzw. $ f(x) \geq y_0$) für alle $ x\in\mathbb{D}_f $\\
(ii) Es gibt eine Folge $ (x_n)_{n\in\mathbb{N}} $ mit $ \lim\limits_ {n\to\infty}=y_0$\vspace*{2mm}\\
Man schreibt dann:\\
\begin{align*}
\sup_{x\in\mathbb{D}_f}=y_0\\
\inf_{x\in\mathbb{D}_f}=y_0
\end{align*}
Beispiel:\\
(i) $ f:\mathbb{R}\longmapsto [0,\infty[ $ mit $f(x)=x^2$
\begin{quote}
	\textbullet \;\; $ \underset{x\in\mathbb{R}}{\sup}f(x)=\infty $ wegen $ \lim\limits_ {x\to\infty}f(x) =\infty$\\
	\textbullet	\;\; $ \underset{x\in\mathbb{R}}{\inf}f(x)=0 $ \;\;\;\; ist auch ein globales Minimum\\
\end{quote}
\begin{center}
	\begin{tikzpicture}
	
	% Gitternetzlinien
	\draw[xstep=1cm, ystep=1cm, very thin, color=lightgray] (-4,4) grid (4,-2);
	% x-Achse
	\draw[->] (-4,0) -- (4,0) node[below,xshift=-1em] {$x$}; 
	%Zahlen auf x-Achse
	\foreach \x in {-3,...,-1}
	\draw[shift={(\x,0)},color=black] (0pt,2pt) -- (0pt,-2pt) node[below] {\footnotesize $\x$};
	\foreach \x in {1,...,3}
	\draw[shift={(\x,0)},color=black] (0pt,2pt) -- (0pt,-2pt) node[below] {\footnotesize $\x$};
	% y-Achse 
	\draw[->] (0,-2) -- (0,4) node[left,yshift=-1em] {$f(x)$};%node[above left]
	%Zahlen auf y-Achse
	\foreach \y in {-1,...,-1}
	\draw[shift={(0,\y)},color=black] (2pt,0pt) -- (-2pt,0pt) node[left] {\footnotesize $\y$};
	\foreach \y in {1,...,3}
	\draw[shift={(0,\y)},color=black] (2pt,0pt) -- (-2pt,0pt) node[left] {\footnotesize $\y$};
	%Ursprung
	%\draw[color=black] (0pt,-10pt) node[right] {\footnotesize $0$};
	
	
	\draw[thick,color=blue] plot[samples=200, domain=-2:2] 
	(\x,{\x*\x}) node[below, xshift=1em]{$f(x)$};
	\draw[color=blue] (0.1,-0.5) node[right] {\footnotesize $\inf$};
	
	
	\end{tikzpicture}\\
\end{center}
(ii)\;\;$ f:\mathbb{R}\backslash\{0\} \longmapsto \mathbb{R}$ mit $ f(x)=\frac{1}{x} $\\
\begin{quote}
	\textbullet \;\;$ \underset{x\in \mathbb{R}\backslash\{0\}}{\sup} f(x)=\infty$ wegen $ \lim\limits_ {\substack{x \to 0\\x>0}}\frac{1}{x}=\infty $\\
	\textbullet \;\; $\underset{x\in \mathbb{R}\backslash\{0\}}{\inf} f(x)=-\infty$ wegen $ \lim\limits_ {\substack{x \to 0\\x<0}}\frac{1}{x}=-\infty $\\
\end{quote}
\begin{center}
	\begin{tikzpicture}
	
	% Gitternetzlinien
	\draw[xstep=1cm, ystep=1cm, very thin, color=lightgray] (-4,4) grid (4,-4);
	% x-Achse
	\draw[->] (-4,0) -- (4,0) node[below,xshift=-1em] {$x$}; 
	%Zahlen auf x-Achse
	\foreach \x in {-3,...,-1}
	\draw[shift={(\x,0)},color=black] (0pt,2pt) -- (0pt,-2pt) node[below] {\footnotesize $\x$};
	\foreach \x in {1,...,3}
	\draw[shift={(\x,0)},color=black] (0pt,2pt) -- (0pt,-2pt) node[below] {\footnotesize $\x$};
	% y-Achse 
	\draw[->] (0,-4) -- (0,4) node[left,yshift=-1em] {$f(x)$};%node[above left]
	%Zahlen auf y-Achse
	\foreach \y in {-3,...,-1}
	\draw[shift={(0,\y)},color=black] (2pt,0pt) -- (-2pt,0pt) node[left] {\footnotesize $\y$};
	\foreach \y in {1,...,3}
	\draw[shift={(0,\y)},color=black] (2pt,0pt) -- (-2pt,0pt) node[left] {\footnotesize $\y$};
	%Ursprung
	\draw[color=black] (0pt,-10pt) node[right] {\footnotesize $0$};
	
	\draw[thick,color=blue] plot[samples=200, domain=0.25:4] 
	(\x,{1/\x}) node[below, xshift=-2em,yshift=2em]{$f(x)$};
	
	\draw[thick,color=blue] plot[samples=200, domain=-4:-0.25] 
	(\x,{1/\x});
	
	
	
	\end{tikzpicture}\\	
\end{center}
\subsubsection{Aufgabe zu Infimum und Supremum}
\href{https://youtu.be/HnmGV8AtkOc}{Video 16}\\
Bestimme bei folgenden Funktionen das Infimum und Supremum.\vspace*{2mm}\\
(i)\;\;\;$ \tan:\; ]-\frac{\pi}{2},\frac{\pi}{2}[\;\;\longmapsto \mathbb{R} $ mit $ \tan(x)=\frac{\sin(x)}{\cos(x)} $\vspace*{2mm}\\
(ii)\; $ g:\mathbb{R}\backslash\{\pm 1\} \longmapsto \mathbb{R}$ mit $ g(x)=\frac{7x}{x^2-1} $\vspace*{2mm}\\
Zu (i):\\
\begin{quote}
	Skizze von Cosinus und Sinus:\\
	\begin{center}
		
		
		\begin{tikzpicture}[scale=2]
		
		% Gitternetzlinien
		\draw[xstep=1cm, ystep=1cm, very thin, color=lightgray] (-3.5,2) grid (3.5,-2);
		% x-Achse
		\draw[->] (-3.5,0) -- (3.5,0) node[below,xshift=-1em] {$x$}; 
		%Zahlen auf x-Achse
		\foreach \x in {-2,...,-1}
		\draw[shift={(\x,0)},color=black] (0pt,2pt) -- (0pt,-2pt) node[below] {\footnotesize $\x$};
		\foreach \x in {1,...,2}
		\draw[shift={(\x,0)},color=black] (0pt,2pt) -- (0pt,-2pt) node[below] {\footnotesize $\x$};
		% y-Achse 
		\draw[->] (0,-2) -- (0,2) node[left,yshift=-1em] {$f(x)$};%node[above left]
		%Zahlen auf y-Achse
		\foreach \y in {-1,...,-1}
		\draw[shift={(0,\y)},color=black] (2pt,0pt) -- (-2pt,0pt) node[left] {\footnotesize $\y$};
		\foreach \y in {1,...,1}
		\draw[shift={(0,\y)},color=black] (2pt,0pt) -- (-2pt,0pt) node[left] {\footnotesize $\y$};
		%Ursprung
		\draw[color=black] (0pt,-10pt) node[right] {\footnotesize $0$};
		
		\draw[thick,color=blue] plot[samples=200, domain=-3.5:3.5] 
		(\x,{sin(\x r)}) node[xshift=-1em,yshift=4.5em]{$\sin(x)$};
		\draw[thick,color=red] plot[samples=200, domain=-3.5:3.5] 
		(\x,{cos(\x r)}) node[xshift=1em,yshift=1.5em]{$\cos(x)$};
		\draw[color=green] (-3.14/2,-1.2) node {\footnotesize $-\genfrac{}{}{}{2}{\pi }{2 }$};
		\draw[color=green] (3.14/2,-1.2) node {\footnotesize $\genfrac{}{}{}{2}{\pi }{2 }$};
		\draw[color=green] (3.14/2,1)--(3.14/2,-1);
		\draw[color=green] (-3.14/2,1)--(-3.14/2,-1);
		\end{tikzpicture}\\
		
	\end{center}
	\textbullet \;\;$ \lim\limits_ {\substack{x \to \frac{\pi}{2}  \\ x< \frac{\pi}{2}}} \tan(x)=\lim\limits_ {\substack{x \to \frac{\pi}{2}  \\ x< \frac{\pi}{2}}}\frac{\sin(x)}{\cos(x)}=+\infty $\\
	\begin{quote}
		Bemerkung: $ +\infty $, weil der Zähler bei $ \frac{\pi}{2}$ positiv ist, wenn der $ \cos(\frac{\pi}{2}) $ Null wird.\\
		
		$\Rightarrow  \underset{x\in \mathbb{D}_f}{\sup}\tan(x) =+\infty $
		
	\end{quote}
	\textbullet \;\;$ \lim\limits_ {\substack{x \to -\frac{\pi}{2}  \\ x> -\frac{\pi}{2}}} \tan(x)=\lim\limits_ {\substack{x \to -\frac{\pi}{2}  \\ x> -\frac{\pi}{2}}}\frac{\sin(x)}{\cos(x)}=-\infty $\\
	\begin{quote}
		Bemerkung: $ -\infty $, weil der Zähler bei $ -\frac{\pi}{2}$ negativ ist, wenn der $ \cos(-\frac{\pi}{2}) $ Null wird.\\
		
		$\Rightarrow  \underset{x\in \mathbb{D}_f}{\inf}\tan(x) =-\infty $
		
	\end{quote}
	
\end{quote}
\pagebreak
Zu (ii):\\
\begin{quote}
	\textbullet\;\; $\lim\limits_ {\substack{x \to 1 \\ x>1}}g(x)=\lim\limits_ {\substack{x \to 1 \\ x>1}}\frac{7x}{x^2-1}=+\infty$\\
	\begin{quote}
		Bemerkung:$ +\infty $, weil Nenner negativ bei $ x>1 $wird und der Zähler durch 7 beschränkt ist  und positiv ist.\\
		$\Rightarrow  \underset{x\in \mathbb{D}_f}{\sup}g(x) =+\infty $
		
	\end{quote}
	\textbullet\;\; $\lim\limits_ {\substack{x \to 1 \\ x<1}}g(x)=\lim\limits_ {\substack{x \to 1 \\ x<1}}\frac{7x}{x^2-1}=-\infty$\\
	\begin{quote}
		Bemerkung:$ -\infty $, weil Nenner negativ bei $ x<1 $wird und der Zähler durch 7 beschränkt ist  und positiv ist.\\
		$\Rightarrow  \underset{x\in \mathbb{D}_f}{\inf}g(x) =-\infty $
		
	\end{quote}
	Den Fall für -1 braucht man ja nicht mehr zu betrachten, weil schon ein Infimum um Supremum gefunden wurde.\\
\end{quote}
%	\chapter{Elementare Funktionen I}

\setcounter{section}{7}
\setcounter{subsection}{0}
\setcounter{subsubsection}{1}
\setcounter{secnumdepth}{4}
\subsection{Polynome}
\href{https://youtu.be/b0zNTRgnXPM}{Video 01}\\
Ein Polynom $ p: \mathbb{C}\longmapsto  \mathbb{C} $ (oder $\mathbb{R}\longmapsto \mathbb{R}  $) mit:\\
\begin{align*}
p(x)=a_0+a_1\cdot x+a_2\cdot x^2+a_3\cdot x^3+\ldots+a_n\cdot x^n
\end{align*}
wobei:
\begin{quote}
	\textbullet\;\;$ a_0,\ldots,a_n \in \mathbb{C}$ (oder $ \mathbb{R} $) heißen Koeffizienten \\
	\textbullet\;\;$ a_n\neq 0 $ heißt $ a_n $ der höchste Koeffizient\\
	\textbullet\;\;Ist $ a_n\neq 0 $ dann hat $ p $ den Grad $ n $, man schreibt: $ \deg(p)=n\in\mathbb{N} $\\
	\textbullet\;\;Für $ a_0=a_1=\ldots=a_n=0 $ ($ \Leftrightarrow p(x)=0\;\forall x\in \mathbb{C}$) definieren wir: $ \deg(P)=-\infty $\\
	\textbullet \;\;$ x_i\in  \mathbb{C} $ heißt Nullstelle von $ p $, wenn $ p(x_i)=0 $ dann gilt: $ (x-x_i) $ teilt $ p(x) $
\end{quote} 
zum Beispiel:\\
\begin{tabular}[t]{ll}% l=linksbündig r=rechtsbündig
	$\blacktriangleright$ & $ \deg(7+\mathrm{i})=\deg((7+\mathrm{i})\cdot x^0)=0 $  \\
	& Dies ist eine konstante Funktion und auch ein Polynom mit $ x^0 $\vspace*{2mm} \\
	$ \blacktriangleright $	& $ p(x)=7x^2+2 $\\
	&Höchste Koeffizient ist 7\\
	&Der Grad von $ p=\deg(2) $\vspace*{2mm}\\
	$\blacktriangleright$ & $\overset{\AC}{p}(x)=x^2-3x+2$\\
	&Es gilt  $\overset{\AC}{p}(1)=1^2-3\cdot 1+2=0$\\
	&$\Rightarrow x_1=1$ ist Nullstelle von $\overset{\AC}{p}(x)  $\\
	&$ (x-x_1)=(x-1)$ teilt das Polynom $\overset{\AC}{p}(x)  $\\
	&Und es gilt  $\overset{\AC}{p}(2)=2^2-3\cdot 2+2=0$\\
    &$\Rightarrow x_2=2$ ist Nullstelle von $\overset{\AC}{p}(x)  $\\
  &$ (x-x_2)=(x-2)$ teilt das Polynom $\overset{\AC}{p}(x)  $\\
& Warum teilt $ (x-x_i) $ das Polynom $\overset{\AC}{p}(x)  $?\\
&Es gilt: $ \overset{\AC}{p}(x)=(x-1)(x-2)=x^2+3x+2 $\\
&d.h. $ \frac{\overset{\AC}{p}(x)}{(x-1)}=\frac{(x-1)(x-2)}{(x-1)}=(x-2) $
\end{tabular} \\
\subsubsection{Beispiel zum Thema Polynome}
\href{https://youtu.be/cbCiEF4aO6E}{Video 02}\\
Mit Polynome kann man wie mit ganzen Zahlen  rechnen.\\
\newpage
\noindent
Beispiele:\\
\textbullet\;\;Sei $ p_1(x) =2x^2+7\mathrm{i}x+5$ und $\;\; p_2(x)=4x^{12}+7x$\\
\begin{quote}


\begin{align*}
p_1(x)+p_2(x)=(2x^2+7\mathrm{i}x+5)+(4x^{12}+7x)&=4x^{12}+2x^2+7\mathrm{i}x+7x+5\\&=4x^{12}+2x^2+(7\mathrm{i}+7)x+5\\
\end{align*}
Bemerkung: \\
\begin{align*}
&\deg(p_1(x))=2\\
&\deg(p_2(x))=12\\
&\deg(p_1(x)+p_2(x))=12
\end{align*}
\end{quote}
\textbullet \;\; Sei $ p(x)=(7x+1) $ und $ q(x)=(2x^2+2x) $\\
\begin{quote}
	\begin{align*}
	(7x+1)(2x^2+2x)=14x^3+14x^2+2x^2+2x=14x^3+16x^2+2x\\
	\end{align*}
	Bemerkung:\\
	\begin{align*}
&	\deg(p(x))=1\\
&\deg(q(x))=2\\
&\deg(q(x)\cdot p(x))=3\\
&\text{Der Grad des Produktes ist die Summe der Grade der Polynome $ p(x), q(x) $}
	\end{align*}
\end{quote}
\textbullet\;\; Sei $ p(x)=\frac{(x^2+1)}{(x-\mathrm{i})} $\\
\begin{align*}
\frac{(x^2+1)}{(x-\mathrm{i})}=\frac{(x-\mathrm{i})(x+\mathrm{i})}{(x-\mathrm{i})}=x+1
\end{align*}
 
 Bemerkung: $ (x-\mathrm{i})(x+\mathrm{i})=x^2+\mathrm{i}x-\mathrm{i}x-\mathrm{i}^2=x^2+1 $\\
\subsection{Faktorisierung von Polynomen}
\subsubsection{Komplexe Polynome}
\href{https://youtu.be/hK5Ggr4f8As}{Video 03}\\
\begin{center}
	

\fbox{\parbox  {95mm} {
	\begin{center}
		Fundamentalsatz der Algebra\\
		\begin{align*}
		& \text{Sei }\\
		&p: \mathbb{C}\longmapsto  \mathbb{C}\;\;\text{,mit }\;\;p(x)=a_0+a_1 x+\ldots +a_nx^n \\
		&\text{ein Polynom, dann gibt es: }\\
		&x_1,\ldots,x_n\in  \mathbb{C}\;\;\text{s.d. }\\
		&p(x)=a_n(x-x_1)(x-x_2)\cdotp\ldots\cdotp (x-x_n)\;\;\text{wenn }\;\;\deg(p)\geq 1
		\end{align*}
	\end{center}
	}}

\end{center}
Bemerkung:\\
Dies wird Faktorisierung von $ p $ genannt.\\
$ (x-x_n) $ nennt man Linearfaktoren.\vspace*{2mm}\\
Offensichtlich gilt:\\
\begin{align*}
p(x_i)=0\;\;\text{für }\;\;i=1,\ldots,n
\end{align*}
Also sind $ x_i $ mit $ i\in\{1,\ldots,n\} $\;\;Nullstellen von $ p(x) $\vspace*{2mm}\\
Bemerkung:\\
Nullstellen müssen nicht verschieden sein. (z.B: doppelte Nullstelle).\vspace*{2mm}\\
Beispiele:\\
\textbullet\;\;Sei $ p(x)=x^2+1$
\begin{align*}
&x^2+1=(x-\mathrm{i})(x+\mathrm{i})\\
&\text{Nullstellen sind:}\;\;x_1=\mathrm{i}, x_2=-\mathrm{i}\\
\end{align*}
\textbullet\;\;Sei $p(x)=(x^2+1)^2$\\
\begin{align*}
(x^2+1)^2=(x^2+1)(x^2+1)&=(x-\mathrm{i})(x+\mathrm{i})(x-\mathrm{i})(x+\mathrm{i})\\
&=(x-\mathrm{i})^2(x+\mathrm{i})^2
\end{align*}

\subsubsection{Reelle Polynome}
\href{https://youtu.be/yTjK64ypuI0}{Video 04}\\
Sei $ p:\mathbb{R}\longmapsto \mathbb{R} $ mit $ p(x)=a_0+a_1 x+\ldots +a_n x $ ein Polynom mit $ \deg(p)\geq 1 $.\\
Dann existieren:\\
$ x_1,\ldots.x_n\in\mathbb{R} $ mit $ k\leq n $ und ein Polynom $ \overset{\AC}{p}(x) $ mit:\\
\begin{align*}
p(x)=\overset{\AC}{p}(x)\cdot(x-x_1)(x-x_k)
\end{align*}
$\blacktriangleright\;\; x_1,\ldots,x_n $ sind reelle Nullstellen von p.\\
$ \blacktriangleright $\;\; Es kann sein, dass ein Polynom übrig bleibt.\\
$ \blacktriangleright $\;\; Dabei gilt: $ \deg( \overset{\AC}{p}(x)) =n-k$\vspace*{4mm}\\ 
Beispiele:\\
\textbullet \;\;Sei $p:\mathbb{R}\longmapsto\mathbb{R}$ mit $ p(x)=x^2+1 $\\
\begin{align*}
 p(x)=x^2+1=x^2+1\\
\end{align*}
Bemerkung:\\
$ p(x)=x^2+1=(x-\mathrm{i})(x+\mathrm{i}) $ hat zwei komplexe Nullstellen!\\
Deswegen ist dies die faktorisierte Form. Das Polynom $ p(x) $ kann nicht weiter faktorisiert werden.\vspace*{4mm}\\
\textbullet\;\; Sei $p:\mathbb{R}\longmapsto\mathbb{R}$ mit $ p(x) =x^4-1$\\
\begin{align*}
p(x) =x^4-1=(x^2-1)(x^2+1)=(x-1)(x+1)(x^2+1)\\
\end{align*}
Faktorisieren heißt:\\
Mit reellen Polynomen die reelen Nullstellen und mit komplexen Polynomen alle Nullstellen bestimmen.\\
\subsubsection{Nullstellenberechnung von Polynomen}
\paragraph{Nullstellenberechnung von Polynomen vom Grad 2}$  $\vspace*{3mm}\\
\href{https://youtu.be/l27JVKX78Kc}{Video 05}\\
Betrachte:\;\;$ p(x)=2x^2-6\mathrm{i}x-4 $\\
Gesucht: $ x_1 $ und $x_2$ mit $ p(x_1)=p(x_2)=0 $\vspace*{2mm}\\
Mit $ p(x)=0 $ erhalten wir:\\
\begin{align*}
2x^2-6\mathrm{i}x-4&=0\\
2(x^2-3\mathrm{i}x-2)&=0
\end{align*}
Mit der $ p-q- $Formel gilt:\\
\begin{align*}
x_{1,2}=-\frac{p}{q}\pm\sqrt{\bigg(\frac{p}{2}\bigg)^2-q}
\end{align*}
\begin{align*}
x_{1,2}&=-\frac{-3\mathrm{i}}{2}\pm\sqrt{\bigg(\frac{3\mathrm{i}}{2}\bigg)^2-(-2)}=-\frac{-3\mathrm{i}}{2}\pm\sqrt{\bigg(\frac{3\mathrm{i}}{2}\bigg)^2+2}\\
&=+\frac{3\mathrm{i}}{2}\pm\sqrt{\frac{-9+8}{4}}=\frac{3\mathrm{i}}{2}\pm\sqrt{-\frac{1}{4}}\\
\end{align*}
Bemerkung:\\
$ x_{1,2} \in  \mathbb{C}\backslash \mathbb{R}\Rightarrow$ Es existiert keine reelle Nullstelle!\\
$ i^2=-1\Leftrightarrow i=\sqrt{-1} $\\
\begin{align*}
x_{1,2}&=\frac{3\mathrm{i}}{2}\pm\sqrt{-\frac{1}{4}}=\frac{3\mathrm{i}}{2}\pm\sqrt{\frac{1}{4}\cdot(-1)}=\frac{3\mathrm{i}}{2}\pm\sqrt{\genfrac{}{}{}{2}{1}{4 }}\sqrt{-1}\\
&=\frac{3\mathrm{i}}{2}\pm\frac{1}{2}\mathrm{i}\\
\end{align*}
\begin{align*}
&x_1=\frac{3}{2}\mathrm{i}+\frac{1}{2}\mathrm{i}=2\mathrm{i}\\
&x_2=\frac{3}{2}\mathrm{i}-\frac{1}{2}\mathrm{i}=\mathrm{i}\\
\end{align*}
Oder mit quatratischer Ergänzung:\\
\begin{align*}
x^2-3\mathrm{i}x-2&=0\\
x^2-3\mathrm{i}x+\bigg(\frac{-3\mathrm{i}}{2}\bigg)^2-\bigg(\frac{-3\mathrm{i}}{2}\bigg)^2-2&=0\\
x^2-3\mathrm{i}x+\frac{9\mathrm{i}^2}{4}-\frac{9\mathrm{i}^2}{4}-2&=0\\
\bigg(x-\frac{3}{2}\mathrm{i}\bigg)^2&=-\frac{9}{4}+2\\
x-\frac{3}{2}\mathrm{i}&=\pm\sqrt{-\frac{1}{4}}\\
x_{1,2}&=\frac{3}{2}i\pm\sqrt{-\frac{1}{4}}\\
x_1&=2\mathrm{i}\\
x_2&=\mathrm{i}\\
\end{align*}
Fazit:\\
\begin{align*}
p(x)=2(x^2-3\mathrm{i}x-2)=2(x-x_0)(x-x_1)=2(x-2\mathrm{i})(x-\mathrm{i}) 
\end{align*}
\paragraph{Nullstellenberechnung von Polynomen vom Grad größer 2}$  $\vspace*{3mm}\\
\href{https://youtu.be/Phr9F9GFmws}{Video 06}\\
Sei ein Polynom $ p $  vom Grad $ n>2 $ gegeben.\\Gesucht sind die Nullstellen.\vspace*{2mm}\\
Vorgehensweise:\\
Bestimme mittels raten oder ausprobieren alle Nullstellen $ n-2 $ ($x_1, x_2,\ldots,x_{n-2}$).\\
Dann gilt:\\
\begin{align*}
p(x)=\overset{\AC}{p}(x)\cdot (x-x_1)\cdot\ldots\cdot(x-x_{n-2})
\end{align*}
Wegen $ \deg(p)=n $ muss $ \deg(\overset{\AC}{p})=n-(n-2)=2 $ sein.\\
Dann können wir auf $\overset{\AC}{p}  $ die p-q-Formel anwenden.\\
Dabei ist:\\
\begin{align*}
\overset{\AC}{p}(x)=\frac{p(x)}{(x-x_1)\cdot\ldots\cdot(x-x_{n-2})}
\end{align*}
Beispiel:\\
Gegeben das Polynom:   \\
\begin{align*}
p(x)=2x^3-4\mathrm{i}x^2+2x-4\mathrm{i}
\end{align*}
Bestimme die Nullstellen.\\
1. Rate eine Nst.:
\begin{quote}
	$x_1=-\mathrm{i}$\\
	$ x_1$ einsetzen \\
	\begin{align*}
	&2(-\mathrm{i})^3-4\mathrm{i}(-\mathrm{i})^2+2(-\mathrm{i})-4\mathrm{i}=2\mathrm{i}+4\mathrm{i}-2\mathrm{i}-4\mathrm{i}=0\\
	&\Rightarrow \frac{p(x)}{(x+\mathrm{i})}
	\end{align*}
\end{quote}
2. Polynomdivision:\\
\begin{quote}
	\begin{align*}
	& (2x^3-4\mathrm{i}x^2+2x-4\mathrm{i})\div(x+\mathrm{i})=2x^2-6\mathrm{i}x-4\\
	-& \underline{(2x^3+2\mathrm{i}x^2)}\\
	& \hspace{0.7cm}-6ix^2+2x\\
	&\hspace{0.2cm}-\underline{(-6\mathrm{i}x^2+6x)}\\
	&\hspace{1.9cm}-4x-4\mathrm{i}\\
	&\hspace{1.4cm}-\underline{(-4x-4\mathrm{i})}\\
	&\hspace{3.5cm} 0
	\end{align*}
	D.h.:\\
	\begin{align*}
	p(x)&=(2x^2-6\mathrm{i}x-4)\cdot(x+\mathrm{i}) \;\; \text{\tiny mit p-q-Formel \normalsize }\\
	&=2(x-\mathrm{i})(x-2\mathrm{i})(x+\mathrm{i})
	\end{align*}
	Die Nullstellen sind:\\
	$ x_1=-\mathrm{i} $\\
	$ x_2=\mathrm{i} $\\
	$ x_3=2\mathrm{i} $\\
	
	
	
\end{quote}
\subsubsection{Komplexe Nullstellen reeller Polynome}
\href{https://youtu.be/uVsYwAYrcUs}{Video 07}\\
Wenn ein Polynom mit reellen Koeffizienten eine komplexe Nullstelle $ x_1 $ hat, dann ist das komplex konjungierte von $ x_1 $ auch eine Nullstelle.\\
Beispiel:\\
\begin{align*}
p(x)=x^2+1=(x-\mathrm{i})(x+1)\\
\end{align*}
Also Nullstellen sind:\\
\begin{align*}
 x_1&=\mathrm{i} \\
 x_2&=-\mathrm{i}=\overline{x}_1=\overline{\mathrm{i}} \\
\end{align*}
Bemerkung:\\
Komplex konjungiert heißt:\\
\begin{align*}
\overline{a+\mathrm{i}b}&=a-ib\\
\overline{i}&=\overline{0+1\cdot \mathrm{i}}=0-\mathrm{i}=-\mathrm{i}
\end{align*}
Allgemein gilt:\\
Sei $ p: \mathbb{C}\longmapsto  \mathbb{C} $ mit $ p(x)=a_0+a_1x+\ldots +a_nx^n $ mit $ a_0, a_1,\ldots, a_n\in \mathbb{R} $ und ist $ x_1 $ eine komplexe Nullstelle von $ p $, dann ist $x_2= \overline{x}_1 $ auch eine Nullstelle von $ p $.\vspace*{2mm}\\
Beispiel:\\
\begin{align*}
p(x)=x^2+4&=(x-2\mathrm{i})(x+2\mathrm{i})\\
x_1&=2\mathrm{i}\\
x_2&=\overline{x}_1=-2\mathrm{i}\\
\end{align*}
\paragraph{Aufgabe zur Nullstellenberechnung Teil I}$  $\vspace*{3mm}\\
\href{https://youtu.be/ujBUQtLfSI4}{Video 08}\\
Faktorisiere das Polynom $ P(x) $:\\
\begin{align*}
&p(x)=x^7-4x^6+8x^5-11x^4+11x^3-8x^2+4x-1\\
&\text{mit Nst.: }\;\;x_1=1\;\;\text{und }\;\;x_2=\mathrm{i}\\
\end{align*}
Da die Koeffizienten von $ P(x) $ reell sind, ist $ x_3=\overline{x}_2=-\mathrm{i} $ auch eine Nullstelle.\\
Also teilt $ P(x) $:\\
\begin{align*}
(x-1)(x-\mathrm{i})(x+\mathrm{i})=(x-1)(x^2+1)=x^3-x^2+x-1\\
\end{align*}
\begin{align*}
&(x^7-4x^6+8x^5-11x^4+11x^3-8x^2+4x-1)\div(x^3-x^2+x-1)=x^4-3x^3+4x^2-3x+1\\
-&\underline{(x^7-\;\;x^6+\;\;x^5-\;\;x^4)}\\
&\hspace{0.35cm}0-\;3x^6+7x^5-10x^4+11x^3\\
&\hspace{0.22cm}-(\underline{-3x^6+3x^5-\;\;3x^4+\;3x^3})\\
&\hspace{1.62cm} 0+4x^5-\;\;7x^4+\;8x^3-8x^2\\
&\hspace{1.35cm}-\underline{(-4x^5-\;\;4x^4+\;4x^3-4x^2)}\\
&\hspace{2.8cm} 0-\;\;3x^4+\;4x^3-4x^2+4x\\
&\hspace{2.75cm}-(\underline{-3x^4+\;3x^3-3x^2+3x})\\
&\hspace{4.2cm} 0+\;\;x^3-\;\;x^2+\;\;x-1\\
&\hspace{3.85cm}-\underline{(+\;\;\;x^3-\;\;x^2+\;\;x-1)}\\
&\hspace{8.2cm}0
\end{align*}
Dann gilt:\\
\begin{align*}
\frac{P(x)}{(x^3-x^2+x-1)}=x^4-3x^3+4x^2-3x+1
\end{align*}
\begin{align*}
\Leftrightarrow P(x)=(x^4-3x^3+4x^2-3x+1)(x-1)(x-i)(x+i)\\
\end{align*}
\paragraph{Aufgabe zur Nullstellenberechnung Teil II}$  $\vspace*{3mm}\\
\href{https://youtu.be/iHXUOPQJfgM}{Video 09}\\
Faktorisiere $ P_2(x)=x^4-3x^3+4x^2-3x+1$\\
Suche Nullstelle:\\
\begin{align*}
P_2(1)=(1-3+4-3+1)=0
\end{align*}
Also ist $ (x-1) $ teiler von $ P_2(x) $\\
\begin{align*}
&(x^4-3x^3+4x^2-3x+1)\div(x-1)=x^3-2x^2+2x-1\\
-&\underline{(x^4-x^3)}\\
&\hspace{0.35cm}0-2x^3+4x^2\\
&\hspace{0.10cm}-\underline{(-2x^3+2x^2)}\\
&\hspace{1.5cm}0+2x^2-3x\\
&\hspace{1.6cm}-\underline{(2x^2-2x)}\\
&\hspace{2.7cm}0-\;\;x+1\\
&\hspace{2.56cm}-\underline{(-\;x+1)}\\
&\hspace{4.4cm}0\\
\end{align*}
Dann gilt:\\
\begin{align*}
\frac{P_2(x)}{(x-1)}=x^3-2x^2+2x-1
\end{align*}
\begin{align*}
\Leftrightarrow P_2(x)=(x^3-2x^2+2x-1)(x-1)\\
\end{align*}
\paragraph{Aufgabe zur Nullstellenberechnung Teil III}$  $\vspace*{3mm}\\
\href{https://youtu.be/I40ZrXwe9rc}{Video 10}\\
Faktorisiere $ P_3(x)=x^3-2x^2+2x-1 $\\
Suche Nullstelle:\\
\begin{align*}
P_3(1)=1-2+2-1=0\\
\end{align*}
Also ist $ (x-1) $ teiler von $ P_2(x) $\\
\begin{align*}
&(x^3-2x^2+2x-1)\div (x-1)=x^2-x+1\\
-&\underline{(x^3-\;\;x^2)}\\
&\hspace{0.35cm}0-\;\;x^2+2x\\
&\hspace{0.3cm}-\underline{(-x^2+\;\;x)}\\
&\hspace{1.55cm}0+\;\;x-1\\
&\hspace{1.9cm}-\underline{(x-1)}\\
&\hspace{3.4cm}0\\
\end{align*}
Dann gilt:\\
\begin{align*}
\frac{P_3(x)}{(x-1)}=x^2-x+1
\end{align*}
\begin{align*}
\Leftrightarrow P_3(x)=(x^2-x+1)(x-1)\\
\end{align*}
Faktorisiere zuletzt: $ P_4(x)=x^2-x+1 $\\
Mit quatratischer Ergänzung oder p-q-Formel:\\
\begin{align*}
x^2-x+1&=0\\
x^2-x+\frac{1}{4}-\frac{1}{4}+1&=0\\
\Leftrightarrow\bigg(x-\frac{1}{2}\bigg)^2&=-1+\frac{1}{4}\\
\Leftrightarrow x-\frac{1}{2}&=\pm\sqrt{-\frac{3}{4}}=\sqrt{-1}\pm\sqrt{\frac{3}{4}}\;\;\notin \mathbb{R}\\
\Leftrightarrow x_{6,7}&=\frac{1}{2}\pm \mathrm{i} \sqrt{\frac{3}{4}}\\
x_6&=\frac{1+\mathrm{i}\sqrt{3}}{2}\\
x_7&=\frac{1-\mathrm{i}\sqrt{3}}{2}\\
\end{align*}
Zusammenfassung der Faktorisierung von $ P(x) $\\
Faktorisierung von $ P(x) $ mit komplexen Argument:\\
\begin{align*}
p(x)&=x^7-4x^6+8x^5-11x^4+11x^3-8x^2+4x-1\\
&=(x^4-3x^3+4x^2-3x+1)(x-1)(x-\mathrm{i})(x+\mathrm{i})\\
&=(x^3-2x^2+2x-1)(x-1)(x-1)(x-\mathrm{i})(x+\mathrm{i})\\
&=(x^2-x+1)(x-1)(x-1)(x-1)(x-i)(x+i)\\
&=\bigg(x-\bigg(\frac{1+\mathrm{i}\sqrt{3}}{2}\bigg)\bigg)\bigg(x-\bigg(\frac{1-\mathrm{i}\sqrt{3}}{2}\bigg)\bigg)(x-1)^3(x-\mathrm{i})(x+\mathrm{i})\\
\end{align*}
Faktorisierung von $ P(x) $ mit reellen Argument:\\
\begin{align*}
p(x)&=x^7-4x^6+8x^5-11x^4+11x^3-8x^2+4x-1\\
&=(x^2-x+1)(x^2+1)(x-1)^3\\
\end{align*}
\subsection{Koeffizientenvergleich}
\href{https://youtu.be/Zipnj1HU4aU}{Video 11}\\
Seien $ p(x)=a_0+a_1x+\ldots+a_nx^n$ und $q(x) =b_0+b_1x+\ldots +b_kx^k $ sowie $ k\leq n $.\\
Stimmen $ p $ und $ q $ auf $ n $ Stellen überein, so gilt $ n=k $ und $ a_i=b_i $ für $ i=0, \cdots ,n $.\vspace*{2mm}\\
Wenn gilt $ p(x_i)=q(x_i) $ für $ i\in\{0,\cdots ,n\} $, dann gilt $ p(x)=q(x) $\\
Die $ x_i $ und $ x_j $ sollen unterschiedlich sein, weil man sonst immer den gleichen Wert einnehmen kann und dies bringt keine Aussage bis $ n $. Deswegen $ x_i\neq x_j $\vspace*{2mm}\\
Beispiel:\\
Schreibe das Polynom $ 3x^2+12 $ als Summe der Polynome $ x^2-3x $ und $ x^2+4 $.\\
Bestimme:\\
\begin{align*}
&A,B \in \mathbb{R}\;\;\text{ mit}\\
&3x^2+12=A(x^2-3x)+B(x^2+4)\\
\end{align*}
Dies ist:\\
\begin{align*}
\Leftrightarrow 3x^2+0\cdot x+12&=Ax^2-3Ax+Bx^2+4B\\
&=(A+B)x^2-3Ax+4B\\
\end{align*}
Es muß übereinstimmen:\\
\begin{align*}
\textcolor{red}{3}x^2\textcolor{blue}{+0}\cdot x \textcolor{green}{+12}=\textcolor{red}{(A+B)}x^2\textcolor{blue}{-3A}x+\textcolor{green}{4B}\\
\end{align*}
Man bekommt dann 3 Relationen und bestimmt dann $ A $ und $ B $.\\
Dies führt auf 3 LGS (lineares Gleichungssystem):\\
\begin{align*}
3&=A+B&\Rightarrow B=3\\
0&=-3A&\Rightarrow A=0\\
12&=4B&\Rightarrow 12=12\\
\end{align*}
Also gilt:\\
\begin{align*}
3x^2+12=0(x^2-3x)+3(x^2+4)\\
\end{align*}
\subsubsection{Aufgabe zum Koeffizientenvergleich}
\href{https://youtu.be/lLfxDN2AW4k}{Video 12}\\
Bestimme zu folgenden Eigenschaften das entsprechende Polynom $ p(x) $.\\
\begin{quote}
	(i)\;\; $ p(x)=p(-x)$ für alle $ x\in \mathbb{C} $\\
	(ii)\;\;$ p(2\mathrm{i})=p(7)=0 $ und $ p(0)=2 $\\
	(iii) $ \deg(p)\leq 4 $\\
\end{quote}
Bestimmung von $ p(x) $:\\
Wegen:\\
\begin{align*}
p(x)=p(-x)
\end{align*}
gilt:\\
\begin{align*}
p(2i)=p(-2i)=p(7)=p(-7)=0
\end{align*}
d.h.: Die Nullstellen von $ p(x) $ sind:\\
\begin{align*}
&x_1=2i\\
&x_2=-2i\\
&x_3=7\\
&x_4=-7\\
\end{align*}
Da $ \deg(p)\leq 4 $ hat $ p $ max. 4 Nullstellen.\\
Dann git wegen den Fundamentalsatz der Algebra:\\
\begin{align*}
p(x)&=\mathrm{k}(x-2\mathrm{i})(x+2\mathrm{i})(x+7)(x-7)\\
&=\mathrm{k}(x^2+4)(x^2-49)=\mathrm{k}(x^4-49x^2+4x^2-196)\\
&=\mathrm{k}(x^4-45x^2-196)\\
&=\mathrm{k}x^4-45\mathrm{k}x^2-\mathrm{k}196\\
\end{align*}
Bestimmung von k:\\
Es gilt:\\
\begin{align*}
&p(0)=2\Rightarrow \mathrm{k}\cdot 0^4 -45\mathrm{k}\cdot 0^2-\mathrm{k}196=2\\
&\Leftrightarrow -\mathrm{k}196=2\\
&\Leftrightarrow \mathrm{k}= -\frac{1}{98}\\
\end{align*}
Fazit:\\
\begin{align*}
p(x)=-\frac{1}{98}x^4+\frac{45}{98}x^2+2\\
\end{align*}
\subsection{Partialbruchzerlegung rationaler Funktionen}
\href{https://youtu.be/imQZqfJI8yI}{Video 13}\\
Eine rationale Funktion $ f(x)=\frac{p(x)}{q(x)} $ mit $ x\in \mathbb{C} $, wobei $ p $ und $ q $ Polynome sind und  $ \deg(q)>\deg(p) $  ist.\\
Wir wollen $ f(x) $ als Summe einfacher Summanden der Form:\\
\begin{align*}
f(x)=\frac{ax+b}{(x-x_i)^n}
\end{align*}
mit $ x_i $ sind Nullstellen von $ q(x) $ schreiben.\\
Das Ziel ist eine rationale Funktion in eine einfache rationale Funktion zu schreiben.\vspace*{2mm}\\
Beispiel:\\
\begin{align*}
f(x)=\frac{5x^2-4x+7}{x^3+x^2-5x+3}=\frac{5x^2-4x+7}{(x-1)^2(x+3)}=\frac{1}{(x-1)}+\frac{2}{(x-1)^2}+\frac{4}{(x+3)}\\
\end{align*}
Es gibt drei Fälle:\\
\underline{1.Fall:}\\
\begin{quote}
	$ q(x) $ zerfällt vollständig $ q(x)=(x-x_1)\cdot\ldots\cdot(x-x_n) $ und $ x_i $ verschieden.\\
\end{quote}
\underline{2.Fall:}\\
\begin{quote}
	$ q(x) $ zerfällt vollständig $ q(x)=(x-x_1)\cdot\ldots\cdot(x-x_n) $ aber $ x_i $ nicht verschieden.\\
\end{quote}
\underline{3.Fall:}\\
\begin{quote}
	$ q(x) $ zerfällt nicht vollständig. \\
\end{quote}
\subsubsection{PBZ: Nenner zerfällt vollständig mit verschiedenen Nullstellen.}
\href{https://youtu.be/MYsd2kAmsk0}{Video 14}\\
$ f(x)=\frac{p(x)}{q(x)} $, $ q(x) $ zerfällt vollständig in verschiedenen Linearfaktoren.\\
Beispiel:\\
Bestimme die PBZ (Partialbruchzerlegung) für:\\
\begin{align*}
f(x)=\frac{7x}{x^3-x^2+x-1}
\end{align*}
Nullstellen von $ x^3-x^2+x-1=0 $ sind:\\
\begin{align*}
&x_1=1\\
&x_2=\mathrm{i}\\
&x_3=-\mathrm{i}\\
\end{align*}
Es muß gelten:\\
\begin{align*}
f(x)=\frac{7x}{x^3-x^2+x-1}=\frac{7x}{(x-1)(x-\mathrm{i})(x+\mathrm{i})}=\frac{A}{(x-1)}+\frac{B}{(x-\mathrm{i})}+\frac{C}{(x+\mathrm{i})}\\
\end{align*}
Bestimme A,B und C:\\
Vorgehensweise:\\
Erklärung der Zuhaltemethode:\\
Mulitipliziere die Gleichung mit $ (x-1) $:\\
\begin{align*}
\frac{7x}{(x-1)(x-\mathrm{i})(x+\mathrm{i})}&=\frac{A}{(x-1)}+\frac{B}{(x-\mathrm{i})}+\frac{C}{(x+\mathrm{i})}\;\;|\cdot (x-1)\\
\frac{7x(x-1)}{(x-1)(x-\mathrm{i})(x+\mathrm{i})}&=\frac{A(x-1)}{(x-1)}+\frac{B(x-1)}{(x-\mathrm{i})}+\frac{C(x-1)}{(x+\mathrm{i})}\\
\frac{7x}{(x-\mathrm{i})(x+\mathrm{i})}&=A+\frac{B(x-1)}{(x-\mathrm{i})}+\frac{C(x-1)}{(x+\mathrm{i})}\\
\end{align*}
Setze die Nullstelle $ x_1=1 $ ein.\\
\begin{align*}
\frac{7\cdot 1}{(1-\mathrm{i})(1+\mathrm{i})}&=A+\frac{B(1-1)}{(1-\mathrm{i})}+\frac{C(1-1)}{(1+\mathrm{i})}\\
\frac{7}{(1-\mathrm{i})(1+\mathrm{i})}&=A\\
\end{align*}
Berechnung von A,B und C mit der Zuhaltemethode.\\
\begin{align*}
\frac{7x}{\textcolor{red}{(x-1)}\textcolor{blue}{(x-\mathrm{i})}\textcolor{green}{(x+\mathrm{i})}}&=\frac{A}{(x-1)}+\frac{B}{(x-\mathrm{i})}+\frac{C}{(x+\mathrm{i})}\;\;|\cdot (x-1)\\
\end{align*}
Durch Zuhalten von $ \textcolor{red}{(x-1)} $ und für $ x=1 $ einsetzen ergibt dann A:\\
\begin{align*}
\frac{7\cdot1}{(1-\mathrm{i})(1+\mathrm{i})}=\frac{7}{(1-\mathrm{i})^2}=\frac{7}{1+1}=\frac{7}{2}=A\\
\end{align*}
Durch Zuhalten von $ \textcolor{blue}{(x-\mathrm{i})} $ und für $ x=\mathrm{i} $ einsetzen ergibt dann B:\\
\begin{align*}
\frac{7\cdot \mathrm{i}}{(\mathrm{i}-1)(\mathrm{i}+1)}=\frac{7\mathrm{i}}{(\mathrm{i}-1)2\mathrm{i}}=\frac{7\mathrm{i}}{-2-2\mathrm{i}}=B\\
\end{align*}
Durch Zuhalten von $ \textcolor{green}{(x+\mathrm{i})} $ und für $ x=-\mathrm{i} $ einsetzen ergibt dann C:\\
\begin{align*}
\frac{7\cdot(-\mathrm{i})}{(-\mathrm{i}-1)(-\mathrm{i}-\mathrm{i})}=\frac{-7\mathrm{i}}{(-\mathrm{i}-1)(-2\mathrm{i})}=\frac{-7\mathrm{i}}{-2+2\mathrm{i}}=C\\
\end{align*}
Fazit:\\
\begin{align*}
\frac{7x}{x^3-x^2+x-1}=\frac{\frac{7}{2}}{(x-1)}+\frac{\frac{7\mathrm{i}}{-2-2\mathrm{i}}}{(x-\mathrm{i})}+\frac{\frac{-7\mathrm{i}}{-2+2\mathrm{i}}}{(x+\mathrm{i})}\\
\end{align*}
\subsubsection{PBZ: Nenner hat mehrfache Nullstellen}
\href{https://youtu.be/4sYdi6gqWZo}{Video 15}\\
$ f(x)=\frac{p(x)}{q(x)} $, $ q(x) $ zerfällt vollständig in Linearfaktoren mit mehrfachen Nullstellen.\\
Beispiel:\\
Bestimme die PBZ für:\\
\begin{align*}
f(x)=\frac{1}{x^3(x+1)^2}\\
\end{align*}
Die Nullstellen von $ x^3(x+1)^2 =0$ sind:\\
\begin{align*}
x_1, x_2, x_3& = 0\\
x_4, x_5& = -1\\
\end{align*}
Es muß gelten:\\
\begin{align*}
\frac{1}{x^3(x+1)^2}=\frac{A_1}{x}+\frac{A_2}{x^2}+\frac{A_3}{x^3}+\frac{B_1}{(x+1)}+\frac{B_2}{(x+1)^2}\\
\end{align*}
Bestimmung von $ A_1, A_2, A_3, B_1$ und $B_2$:\\
Für die höchsten Potenzen ($ A_3, B_2 $) funktioniert die Zuhaltemethode:\\
\begin{align*}
A_3=\frac{1}{(0+1)^2}=1\\
\end{align*}
\begin{align*}
B_2=\frac{1}{(-1)^3}=-1\\
\end{align*}
Bleibt noch $ A_1, A_2 $ und $ B_1 $ zu bestimmen:\\
Es gilt nun:\\
\begin{align*}
\frac{1}{x^3(x+1)^2}=\frac{A_1}{x}+\frac{A_2}{x^2}+\frac{1}{x^3}+\frac{B_1}{(x+1)}-\frac{1}{(x+1)^2}\\
\end{align*}
Die Gleichung mit dem Nenner: $ x^3(x+1)^2 $ mulitiplizieren.\\
\begin{align*}
1&=\frac{A_1x^3(x+1)^2}{x}+\frac{A_2x^3(x+1)^2}{x^2}+\frac{x^3(x+1)^2}{x^3}+\frac{B_1x^3(x+1)^2}{(x+1)}-\frac{x^3(x+1)^2}{(x+1)^2}\\
&=A_1x^2(x+1)^2 +A_2x(x+1)^2+(x+1)^2+B_1x^3(x+1)-x^3\\
\end{align*}
Es gilt:\\
\begin{align*}
0x^4+0x^3+0x^2+0x+1&=A_1x^2(x+1)^2 +A_2x(x+1)^2+(x+1)^2+B_1x^3(x+1)-x^3\\
&=A_1x^2(x^2+2x+1)+A_2x(x^2+2x+1)+x^2+2x+1+B_1x^3(x+1)-x^3\\
&=A_1(x^4+2x^3+x^2)+A_2(x^3+2x^2+x)+x^2+2x+1+B_1(x^4+x^3)-x^3\\
\end{align*}
Jetzt die Gleichung nach $ x $ Potenzen sortieren:\\
\begin{align*}
0x^4+0x^3+0x^2+0x+1&=(A_1+B_1)x^4+(2A_1+A_2+B_1-1)x^3+(A_1+2A_2+1)x^2+1\\
\end{align*}
%\subsubsection{ PBZ: Nenner hat mehrfache Nullstellen Teil II}
\href{https://youtu.be/_I2IcH8-JMw}{Video 16}\\
Mit Koeffizientvergleich bekommt man folgendes LGS:\\
\begin{align}
0&=A_1+B_1 \label{a1}\\
0&=2A_1+A_2+B_1-1 \label{a2}\\
0&=A_1+2A_2+1 \label{a3}\\
0&=A_2+2\label{a4}\\
1&=1
\end{align}
Mit \ref{a4} :\\
\begin{align*}
&0=A_2+2\\
&\Leftrightarrow A_2=-2\\
\end{align*}
$ A_2=-2 $ in \ref{a3} einsetzen:\\
\begin{align*}
&0=A_1+2A_2+1=A_1+2(-2)+1\\
&\Leftrightarrow A_1=3\\
\end{align*}
$ A_1=3 $ in\ref{a1} einsetzen:\\
\begin{align*}
&0=A_1+B_1=3+B_1\\
&\Leftrightarrow B_1=-3\\
\end{align*}
Fazit:\\
\begin{align*}
\frac{1}{x^3(x+1)^2}=\frac{3}{x}-\frac{2}{x^2}+\frac{1}{x^3}-\frac{3}{(x+1)}-\frac{1}{(x+1)^2}\\
\end{align*}

\subsubsection{PBZ: Nenner zerfällt nicht vollständig }
\href{https://youtu.be/EepB1ozV-fQ}{Video 17}\\
$ f(x)=\frac{p(x)}{q(x)} $, $ q(x) $ zerfällt nicht vollständig in Linearfaktoren.\\
Kann wegen des Fundamentalsatzes der Algebra nur auftreten, wenn nur reelle Nullstellen für $ q(x) $ zugelassen sind.\\
Beispiel:\\
Bestimme die PBZ für:\\
\begin{align*}
f(x)=\frac{2x-1}{(x-2)(x^2+2)}\;\;\text{mit }\;\;x\in\mathbb{R}\\
\end{align*}
Die Nullstellen von $ (x-2)(x^2+2)=0 $ sind:\\
\begin{align*}
x_1=2
\end{align*}
Bemerkung:\\
Die Nullstellen von $ x^2+2=0 $ sind:\\
\begin{align*}
x_2=\mathrm{i}\sqrt{2}\\
x_3=-\mathrm{i}\sqrt{2}\\
\end{align*}
D.h.: $ x^2+2 $ zerfällt nicht weiter, da $ x_2, x_3 $ komplexe Zahlen sind ($x_1,x_2 \notin\mathbb{R} $). \\
Wenn $ x\in \mathbb{C} $, dann zerfällt das Polynom komplett.\vspace*{2mm}\\
Es muß gelten:\\
\begin{align*}
f(x)=\frac{2x-1}{(x-2)(x^2+2)}=\frac{A}{(x-2)}+\frac{Bx+C}{(x^2+2)}\\
\end{align*}
Bemerkung:\\
$ Bx+c $ weil $ x^2+2 $ nicht weiter zerfällt.\\
%\subsubsection{PBZ: Nenner zerfällt nicht vollständig Teil II }
\href{https://youtu.be/FvWtzBpsLvk}{Video 18}\\
Berechnung von $ A,B $ und $ C $.\\
Bestimmung von $ A $ mit der Zuhaltemethode:\\
\begin{align*}
A=\frac{2\cdot 2-1}{2^2+2}=\frac{3}{6}=\frac{1}{2}\\
\end{align*}
Bestimmung von $ B $ und $ C $:\\
Es gilt nun:\\
\begin{align*}
\frac{2x-1}{(x-2)(x^2+2)}=\frac{\frac{1}{2}}{(x-2)}+\frac{Bx+C}{(x^2+2)}\\
\end{align*}
Die Gleichung mit dem Nenner $ (x-2)(x^2+2) $ multiplizieren.\\
\begin{align*}
\frac{(2x-1)(x-2)(x^2+2)}{(x-2)(x^2+2)}&=\frac{\frac{1}{2}(x-2)(x^2+2)}{(x-2)}+\frac{(Bx+C)(x-2)(x^2+2)}{(x^2+2)}\\
2x-1&=\frac{1}{2}(x^2+2)+(Bx+C)(x-2)\\
&=\frac{1}{2}x^2 +1+Bx^2-2Bx+Cx-2C\\
0x^2+2x-1&=\bigg(\frac{1}{2}+B\bigg)x^2+(C-2B)x+1-2C\\
\end{align*}
Mit Koeffizientvergleich bekommt man folgendes LGS:\\
\begin{align*}
&0=\frac{1}{2}+B\\
&\Leftrightarrow B=-\frac{1}{2}\\
&-1=1-2C\\
&\Leftrightarrow C=1\\
\end{align*}
Fazit:\\
\begin{align*}
&A=\frac{1}{2}\\
&B=-\frac{1}{2}\\
&C=1\\
\end{align*}
Daraus folgt:\\
\begin{align*}
f(x)=\frac{2x-1}{(x-2)(x^2+2)}=\frac{\frac{1}{2}}{(x-2)}+\frac{(-\frac{1}{2})x+1}{(x^2+2)}
\end{align*}
\subsubsection{PBZ: Nennergrad ist kleiner gleich als der Zählergrad }
\href{https://youtu.be/ovgTWyr47GY}{Video 19}\\
\begin{align*}
f(x)\frac{p(x)}{q(x)}\;\;\text{ mit}\;\;\deg(p)\geq \deg (q)\\
\end{align*}
In diesen Fall verwendet wir die Polynomdivision.\\
Wähle ein Polynom $ h(x) $ und $ r(x) $ mit.\\
\begin{align*}
p(x)=h(x)\cdot q(x)+r(x)\;\;\text{mit } \deg(r)<\deg(q)\\
\end{align*}
\newpage
\noindent
Bemerkung:\\
\begin{quote}
	

Das Polynom: $ \frac{7}{2} $\\
Addition und Multiplikation bei Polynomen gilt genau so wie bei Zahlen.\\
Wenn man den Bruch $ \frac{7}{2} $ teilt, dann ist der Rest immer kleiner als der Nenner.\\
\begin{align*}
7=3\cdot 2+1
\end{align*}

 Wähle $h(x)=3  $ und $ p(x)=7 $ dann ergibt sich $ r(x)=1 $ mit der Bedingung $\deg(r)<\deg(q)  $.\\
\begin{align*}
p(x)&=h(x)q(x)+r(x)\\
7&=\;\;3\;\;\cdot\; 2 \;\;+\;1\\
\end{align*}
Rest=1 ist imme kleiner als der Nenner=2.
\end{quote}
Nun teilt man die das Polynom $ p(x)=h(x)q(x)+r(x) $ mit $ q(x) $.\\
\begin{align*}
f(x)=\frac{p(x)}{q(x)}=h(x)+\frac{r(x)}{q(x)}\\
\end{align*}
Nun, da $ \deg(r)<\deg(q) $ gilt können wir die PBZ auf $ \frac{r(x)}{q(x)}$ anwenden.\vspace*{2mm}\\
Beispiel:\\
\begin{align*}
f(x)=\frac{p(x)}{q(x)}=\frac{x^2+x+1}{x-2}\\
\end{align*}
Wegen $ \deg(p)=2>1=\deg(q) $ muß  zunächst die Polynomdivision durchgeführt werden.\\
\begin{align*}
&(x^2+x+1)\div(x-2)=x+3\\
-&\underline{(x^2-2x)}\\
&\hspace{1.1cm} 3x+1\\
&\hspace{0.5cm}-\underline{(3x-6)}\\
&\hspace{2.2cm}7\\
\end{align*}
\begin{align*}
&h(x)=x+3\\
&r(x)=7\\
\end{align*}
Fazit:\\
\begin{align*}
f(x)=\frac{p(x)}{q(x)}=\frac{x^2+x+1}{x-2}=x+3+\frac{7}{x-2}\\
\end{align*}
\subsubsection{Partialbruchzerlegung: Anleitung }
\href{https://youtu.be/ADKDURyGOmg}{Video 20}\\
Anleitung für eine Partialbruchzerlegung von $ f(x)=\frac{p(x)}{q(x)} $\\
1. Bestimmung von $ \deg(p) $ und $ \deg(q) $\vspace*{2mm}\\
2.1 $ \deg(p)< \deg(q) $\\
\begin{quote}
	$\blacktriangleright$ Faktorisierung von $ q(x) $\\
	\begin{quote}
		a. $ q(x) $ zerfällt vollständig in Linearfaktoren\\
		\begin{quote}
			\textbullet \;\;mit einfacher Nullstelle $ \Rightarrow $ Zuhaltemethode\\
			\textbullet \;\;mit mehrfacher Nullstelle $ \Rightarrow $ mit Koeffizientenvergleich $ \Rightarrow $ LGS\\
		\end{quote}
	    b. $ q(x) $ zerfällt nicht vollständig\\
        \begin{quote}
	    Betrachte Ansatz mit $ \frac{Bx+C}{x-x_i} \Rightarrow$ Koefizientenvergleich $ \Rightarrow $ LGS\\
	    \end{quote}
	\end{quote}
\end{quote}
2.2 $ \deg(p)> \deg(q) $\\
\begin{quote}
	$\blacktriangleright$ Bestimme mit Polynomdivision
	 \begin{align*}
	  &h(x) \;\;\text{und }\;\;r(x)\\
	  &\frac{p(x)}{q(x)}=h(x)+\frac{r(x)}{q(x)}\;\;\text{mit }  \deg(r)<\deg(q) \\
	  &\text{dann weiter mit 2.1 }
	 \end{align*}
	
\end{quote}
	%\chapter{Differentation}

\setcounter{section}{8}
\setcounter{subsection}{0}
\setcounter{subsubsection}{1}
\setcounter{secnumdepth}{4}
\subsection{Geometrische Interpretation der Ableitung Teil I u. II}
\href{https://youtu.be/iWT2SWmVfuw}{Video 01}\\
Sei $ f:\mathbb{R}\supset \mathbb{D}_f\longmapsto \mathbb{W}_f\subset\mathbb{R} $ und $ x_0\in \mathbb{D}_f $\\
\begin{center}
	
		\begin{tikzpicture}
	
	% Gitternetzlinien
	\draw[xstep=1cm, ystep=1cm, very thin, color=lightgray] (-5,5) grid (5,-2);
	% x-Achs
	
	\draw[thick,color=black] plot[samples=200, domain=0.3:4] 
	(\x,{4.1-\x}) node[yshift=2.2em,xshift=-0.5em] {$g(x)$};
	\draw[color=blue] (2,2.1) node {\textbullet};
	
	\draw[thick,color=blue] plot[samples=200, domain=-4:4] 
	(\x,{-0.23*\x*\x+3});
	
	\draw[thick,color=blue,dashed] (2,0) -- (2,2);
	
	\draw[color=blue] (1.8,2.4) node[right] {\footnotesize $P_0=(x_0,f(x_0))$};
	\draw[color=blue] (1.7,-0.8) node[right] {\footnotesize $x_0$};
	\draw[color=blue] (-5,-.5) node[right] {\footnotesize $f(x)$};
	\end{tikzpicture}\\

\end{center}
Die Tangente von $ g $ an $ f $ im Punkt $ P_0 $ soll mit der Geradengleichung dargestellt werden.\\
\begin{align*}
g(x)=mx+n\\
\end{align*}
Die Ableitung von $ f $ in $ x_0 $ ist $ f' (x_0)=m $. Daraus folgt, dass die Steigung der Geraden  bestimmt werden soll.\vspace*{2mm}\\
\newpage
\noindent
Nun wählt man einen zweiten, beliebigen Punkt $P_1\neq P_0 $ auf $ f $ und bildet nun damit eine Sekante.\\
\begin{center}
	\begin{tikzpicture}
	
	% Gitternetzlinien
	\draw[xstep=1cm, ystep=1cm, very thin, color=lightgray] (-5,5) grid (5,-2);
	% x-Achse
	\draw[->] (-5,0) -- (5,0) node[below,xshift=-1em] {$x$}; 
	%Zahlen auf x-Achse
	\foreach \x in {-3,...,-1}
	\draw[shift={(\x,0)},color=black] (0pt,2pt) -- (0pt,-2pt) node[below] {\footnotesize $\x$};
	\foreach \x in {1,...,3}
	\draw[shift={(\x,0)},color=black] (0pt,2pt) -- (0pt,-2pt) node[below] {\footnotesize $\x$};
	% y-Achse 
	\draw[->] (0,-2) -- (0,5) node[left,yshift=-1em] {$f(x)$};%node[above left]
	%Zahlen auf y-Achse
	\foreach \y in {-1,...,-1}
	\draw[shift={(0,\y)},color=black] (2pt,0pt) -- (-2pt,0pt) node[left] {\footnotesize $\y$};
	\foreach \y in {1,...,3}
	\draw[shift={(0,\y)},color=black] (2pt,0pt) -- (-2pt,0pt) node[left] {\footnotesize $\y$};
	%Ursprung
	\draw[color=black] (0pt,-10pt) node[right] {\footnotesize $0$};
	
	\draw[thick,color=black] plot[samples=200, domain=0.3:4] 
	(\x,{4-\x}) node[yshift=2.2em,xshift=-0.8em] {$g(x)$};
	
	\draw[thick,color=blue] plot[samples=200, domain=-4:4] 
	(\x,{-0.25*\x*\x+3});
	
	\draw[thick,color=blue,dashed] (2,0) -- (2,2);
	\draw[color=blue] (1.8,2.4) node[right] {\footnotesize $P_0=(x_0,f(x_0))$};
	\draw[color=blue] (1.7,-0.8) node[right] {\footnotesize $x_0$};
	
	\draw[color=blue] (2,2) node {\textbullet};
	\draw[thick,color=blue,dashed] (-3,0) -- (-3,0.75);
	\draw[color=blue] (-3.1,-0.8) node[right] {\footnotesize $x_1$};
	\draw[color=blue] (-3,0.75) node {\textbullet};
	
	\draw[color=blue] (-3,0.55) node[right] {\footnotesize $P_1=(x_1,f(x_1))$};
	\draw[thick,color=red,dashed] (-3,0.75) -- (2,2);
	\draw[color=red] (0.5,1.4) node[right] {\footnotesize $s(x)$};
	\draw[color=blue] (-5,-.5) node[right] {\footnotesize $f(x)$};
	\end{tikzpicture}\\	
\end{center}
Bestimmung der Sekante: $ s(x)=mx+n $\\
Es gilt mit $ P_1=(x_1,f(x_1)) $ und $ P_2=(x_0,f(x_0)) $:\\
\begin{align}
f(x_0)=mx_0+n \tag{I}\\
f(x_1)=mx_1+n \tag{II}
\end{align}
Lösung des LGS's.\\
Mit II-I erhält man $ m $.\\
\begin{align*}
&\;\;\;\;\;\;\;\;\;\;\;\;f(x_1)=mx_1+n \\
&\;\;\;\;\;\;\;-(\underline{f(x_0)=mx_0+n })\\
&f(x_1)-f(x_0)=mx_1-mx_0\\
&f(x_1)-f(x_0)=m(x_1-x_0)\\
&\Leftrightarrow\\
&m=\frac{f(x_1)-f(x_0)}{x_1-x_0}\\
\end{align*}
Durch Einsetzen von $ m $ in I erhält man $ n $.\\
\begin{align*}
&f(x_0)=\frac{f(x_1)-f(x_0)}{x_1-x_0}\cdot x_0+n\\
&\Leftrightarrow\\
&n=f(x_0)-\frac{f(x_1)-f(x_0)}{x_1-x_0}\cdot x_0\\
\end{align*}
Durch einsetzen von $ n,m $ in $ s(x)=mx+n $ erhält man $ s(x) $:
\begin{align*}
s(x)&=\frac{f(x_1)-f(x_0)}{x_1-x_0}\cdot x+f(x_0)-\frac{f(x_1)-f(x_0)}{x_1-x_0}\cdot x_0\\
s(x)&=f(x_0)+\frac{f(x_1)-f(x_0)}{x_1-x_0}\cdot(x-x_0)
\end{align*}
Das Ziel ist jetzt, eine geschlossene Formel für die Steigung $ m $ zu finden.\\
Die Idee ist:\\
Den Punkt $ P_1 $ entlang der Kurve von $ f(x) $ zum Punkt $ P_0 $ laufen zu lassen.\\
\begin{center}
	\begin{tikzpicture}
	
	% Gitternetzlinien
	\draw[xstep=1cm, ystep=1cm, very thin, color=lightgray] (-5,5) grid (5,-2);
	% x-Achse
	\draw[->] (-5,0) -- (5,0) node[below,xshift=-1em] {$x$}; 
	%Zahlen auf x-Achse
	\foreach \x in {-3,...,-1}
	\draw[shift={(\x,0)},color=black] (0pt,2pt) -- (0pt,-2pt) node[below] {\footnotesize $\x$};
	\foreach \x in {1,...,3}
	\draw[shift={(\x,0)},color=black] (0pt,2pt) -- (0pt,-2pt) node[below] {\footnotesize $\x$};
	% y-Achse 
	\draw[->] (0,-2) -- (0,5) node[left,yshift=-1em] {$f(x)$};%node[above left]
	%Zahlen auf y-Achse
	\foreach \y in {-1,...,-1}
	\draw[shift={(0,\y)},color=black] (2pt,0pt) -- (-2pt,0pt) node[left] {\footnotesize $\y$};
	\foreach \y in {1,...,3}
	\draw[shift={(0,\y)},color=black] (2pt,0pt) -- (-2pt,0pt) node[left] {\footnotesize $\y$};
	%Ursprung
	\draw[color=black] (0pt,-10pt) node[right] {\footnotesize $0$};
	
	\draw[thick,color=black] plot[samples=200, domain=0.3:4] 
	(\x,{4-\x}) node[yshift=2.2em,xshift=-0.8em] {$g(x)$};
	
	\draw[thick,color=blue] plot[samples=200, domain=-4:4] 
	(\x,{-0.25*\x*\x+3});
	
	\draw[thick,color=blue,dashed] (2,0) -- (2,2);
	\draw[color=blue] (1.8,2.4) node[right] {\footnotesize $P_0=(x_0,f(x_0))$};
	\draw[color=blue] (1.7,-0.8) node[right] {\footnotesize $x_0$};
	
	\draw[color=blue] (2,2) node {\textbullet};
	\draw[thick,color=blue,dashed] (-3,0) -- (-3,0.75);
	\draw[color=blue] (-3.1,-0.8) node[right] {\footnotesize $x_1$};
	\draw[color=blue] (-3,0.75) node {\textbullet};
	
	\draw[color=blue] (-3,0.55) node[right] {\footnotesize $P_1=(x_1,f(x_1))$};
	\draw[thick,color=red,dashed] (-3,0.75) -- (2,2);
	\draw[color=red] (0.5,1.4) node[right] {\footnotesize $s(x)$};
	\draw[thick,color=red,dashed] (-2,2) -- (2,2);
	\draw[color=blue] (-2,2) node {\textbullet};
	\draw[color=blue] (-2.5,2.45) node[right] {\footnotesize $P_1$};
	\draw[thick,color=red,->] (-3,1) -- (-2.2,2);	
	\draw[thick,color=red,dashed] (-1,2.7) -- (2,2);
	\draw[color=blue] (-1,2.7) node {\textbullet};
	\draw[color=blue] (-1.5,3) node[right] {\footnotesize $P_1$};
	\draw[thick,color=red,->] (-0.8,3) -- (0,3.1);	
	\draw[thick,color=red,->] (-2,2.2) -- (-1.2,2.8);
	\draw[color=blue] (-5,-.5) node[right] {\footnotesize $f(x)$};
	\end{tikzpicture}\\	
\end{center}
Dann gilt, wenn $ f(x) $ bis zu $ f(x_0) $ geht:\\
\begin{align*}
\lim\limits_ {x_1\to x_0}\frac{f(x_1)-f(x_0)}{x_1-x_0}=m=f'(x_0)\\
\end{align*}
\href{https://youtu.be/TlYjfC4L394}{Video 02}\\
Natürlich kann man ein beliebiges $ x_1 $ wählen, also gilt dann:\\
\begin{align*}
\lim\limits_ {x\to x_0}\frac{f(x)-f(x_0)}{x-x_0}=m=f'(x_0)\\
\end{align*}
Überlegungen:\\
\begin{center}
	
	
	\begin{tikzpicture}
	
	% Gitternetzlinien
	\draw[xstep=1cm, ystep=1cm, very thin, color=lightgray] (-4,4) grid (4,-2);
	% x-Achse
	\draw[->] (-4,0) -- (4,0) node[below,xshift=-1em] {$x$}; 
	%Zahlen auf x-Achse
	\foreach \x in {-3,...,-1}
	\draw[shift={(\x,0)},color=black] (0pt,2pt) -- (0pt,-2pt) node[below] {\footnotesize $\x$};
	\foreach \x in {1,...,3}
	\draw[shift={(\x,0)},color=black] (0pt,2pt) -- (0pt,-2pt) node[below] {\footnotesize $\x$};
	% y-Achse 
	\draw[->] (0,-2) -- (0,4) node[left,yshift=-1em] {$f(x)$};%node[above left]
	%Zahlen auf y-Achse
	\foreach \y in {-1,...,-1}
	\draw[shift={(0,\y)},color=black] (2pt,0pt) -- (-2pt,0pt) node[left] {\footnotesize $\y$};
	\foreach \y in {1,...,3}
	\draw[shift={(0,\y)},color=black] (2pt,0pt) -- (-2pt,0pt) node[left] {\footnotesize $\y$};
	%Ursprung
	\draw[color=black] (0pt,-10pt) node[right] {\footnotesize $0$};
	
	\draw[thick,color=blue] plot[samples=200, domain=1:4] 
	(\x,{sqrt(\x)}) node[below, xshift=-2.5em,yshift=2em] {$g(x)$};
	
	\draw[thick,color=blue] plot[samples=200, domain=-2:1] 
	(\x,{-\x*\x+2}) node[below, xshift=-75]{$f(x)$};
	\draw[thick,color=red] plot[samples=200, domain=-0.5:2.5] 
	(\x,{\x});
	\draw[thick,color=red] plot[samples=200, domain=-0.5:2.5] 
	(\x,{-\x+2});
	\draw[thick,color=red] plot[samples=200, domain=-0.5:2.5] 
	(\x,{1});
	\end{tikzpicture}\\
\end{center}
An dieser Spitze können beliebig viele Tangenten angelegt werden. Der Grenzwert ist hier nicht eindeutig. Er existiert hier nicht.\vspace*{4mm}\\
Polstelle:\\
\begin{center}
	\begin{tikzpicture}
	
	% Gitternetzlinien
	\draw[xstep=1cm, ystep=1cm, very thin, color=lightgray] (-4,4) grid (4,-2);
	% x-Achse
	\draw[->] (-4,0) -- (4,0) node[below,xshift=-1em] {$x$}; 
	%Zahlen auf x-Achse
	\foreach \x in {-3,...,-1}
	\draw[shift={(\x,0)},color=black] (0pt,2pt) -- (0pt,-2pt) node[below] {\footnotesize $\x$};
	\foreach \x in {1,...,3}
	\draw[shift={(\x,0)},color=black] (0pt,2pt) -- (0pt,-2pt) node[below] {\footnotesize $\x$};
	% y-Achse 
	\draw[->] (0,-2) -- (0,4) node[left,yshift=-1em] {$f(x)$};%node[above left]
	%Zahlen auf y-Achse
	\foreach \y in {-1,...,-1}
	\draw[shift={(0,\y)},color=black] (2pt,0pt) -- (-2pt,0pt) node[left] {\footnotesize $\y$};
	\foreach \y in {1,...,3}
	\draw[shift={(0,\y)},color=black] (2pt,0pt) -- (-2pt,0pt) node[left] {\footnotesize $\y$};
	%Ursprung
	\draw[color=blue] (0pt,-10pt) node {\footnotesize $x_0$};
	
	\draw[thick,color=blue] plot[samples=200, domain=-4:-0.25] 
	(\x,{1/abs(\x)}) node[below, yshift=-6.6em,xshift=-8em] {$h(x)$};
	
	\draw[thick,color=blue] plot[samples=200, domain=0.25:4] 
	(\x,{1/\x}) node[below, yshift=2em,xshift=-1em]{$g(x)$};
	
	
	
	\end{tikzpicture}\\
\end{center}
Bei $ x_0 $ kann keine Tangente angelegt werden. Die Funkton ist hier nicht differenzierbar.\\
Der Grenzwert geht gegen $ +\infty $. Geometrisch könnte man sich eine Tangente parallel zur y-Achse vorstellen.\\
\subsection{Differenzierbarkeit}
\href{https://youtu.be/RvEYe71qW_w}{Video 03}\\
Definition:\\
$f:\mathbb{R}\supset I\longmapsto\mathbb{R}$ heißt in $ x_0\in I $ differenzierbar, wenn dieser Differentialquotient existiert.\\
\begin{align*}
\lim\limits_ {x\to x_0}\frac{f(x)-f(x_0)}{x-x_0}=f'(x_0)\in \mathbb{R}\\
\end{align*}
Man nennt $ f'(x_0) $ die Ableitung von $ f $ in $ x_0\in I $.\\
Die Funktion $ f $ ist auf der Menge $I$ diffbar, wenn $ f $ in jedem $ x\in I $ diffbar ist.\vspace*{2mm}\\
Beispiel:\\
Die Funktion $ f:\mathbb{R}\longmapsto [0,\infty[ $ mit $ x\longmapsto x^2 $ ist in $ x_0 $ diffbar.\\
Denn der\\
\begin{align*}
\lim\limits_ {x\to 0}\frac{f(x)-f(x_0)}{x-x_0}=\lim\limits_ {x\to 0}\frac{x^2-0^2}{x-0}=\lim\limits_ {x\to 0}\frac{x^2}{x}=\lim\limits_ {x\to 0}x=0=f'(0)\\
\end{align*}
Allgemein für $ x_0\in \mathbb{R} $ beliebig gilt:\\
\begin{align*}
\lim\limits_ {x\to 0}\frac{f(x)-f(x_0)}{x-x_0}=\lim\limits_ {x\to 0}\frac{x^2-x_0^2}{x-x_0}=\lim\limits_ {x\to 0}\frac{(x-x_0)(x+x_0)}{x-x_0}=\lim\limits_ {x\to x_0}x+x_0=x_0+x_0=2x_0=f'(x_0)\\
\end{align*}
Da die $ 2\in \mathbb{R} $ (von $2x_0=f'(x_0)$ )ist, gilt:\\
$ f $ ist für alle $ x\in \mathbb{R} $ diffbar mit $ f' (x)=(x^2)'=2\cdot x$\\
\newpage
\noindent
\subsubsection{Aufgabe zur Differenzierbarkeit}
\href{https://youtu.be/xyC4rGxDFrM}{Video 04}\\
Zeige, dass $ f:[0,\infty[ \;\;\longmapsto[0,\infty[$ mit $ f(x)=\sqrt{x} $ auf $ ]0,\infty[ $ diffbar ist und berechne $ f'(x) $.\vspace*{2mm}\\
Betrachte den Grenzwert: \\
\begin{align*}
\lim\limits_ {x\to x_0}\frac{f(x)-f(x_0)}{x-x_0}\;\;\text{für }\;\; x_0\in \;]0,\infty[ \\
\end{align*}
\begin{align*}
\lim\limits_ {x\to x_0}\frac{f(x)-f(x_0)}{x-x_0} &=\lim\limits_ {x\to x_0}\frac{\sqrt{x}-\sqrt{x_0}}{x-x_0}\\
&=\lim\limits_ {x\to x_0}\frac{(\sqrt{x}+\sqrt{x_0})(\sqrt{x}-\sqrt{x_0})}{(\sqrt{x}+\sqrt{x_0})(x-x_0)}\\
&=\lim\limits_ {x\to x_0}\frac{x-x_0}{(\sqrt{x}+\sqrt{x_0})(x-x_0)}\\
&=\lim\limits_ {x\to x_0}\frac{1}{\sqrt{x}+\sqrt{x_0}}\\
&=\frac{1}{\sqrt{x_0}+\sqrt{x_0}}\\
&=\underbrace{\frac{1}{2\sqrt{x_0}}}_{\in \mathbb{R}}=f'(x_0)
\end{align*}
D.h. der Grenzwert für $\lim\limits_ {x\to x_0}\frac{f(x)-f(x_0)}{x-x_0}  $ existiert und so gilt:\\
\begin{align*}
f'=(\sqrt{x})'=\frac{1}{2\sqrt{x}}\;\;\text{für }\;\; x\in \;]0,\infty[\\
\end{align*}
Bei Null existiert kein Grenzwert, weil von links kein Grenzwert berechnet werden kann.\\
\begin{align*}
\lim\limits_ {\substack{x \to 0 \\ x<0}}=\frac{1}{2\sqrt{x}}\;\;\text{geht nicht, weil nicht definiert!! }
\end{align*}
Geometrisch gedeutet:\\
\begin{center}
	
	
	\begin{tikzpicture}
	
	% Gitternetzlinien
	\draw[xstep=1cm, ystep=1cm, very thin, color=lightgray] (-4,4) grid (4,-2);
	% x-Achse
	\draw[->] (-4,0) -- (4,0) node[below,xshift=-1em] {$x$}; 
	%Zahlen auf x-Achse
	\foreach \x in {-3,...,-1}
	\draw[shift={(\x,0)},color=black] (0pt,2pt) -- (0pt,-2pt) node[below] {\footnotesize $\x$};
	\foreach \x in {1,...,3}
	\draw[shift={(\x,0)},color=black] (0pt,2pt) -- (0pt,-2pt) node[below] {\footnotesize $\x$};
	% y-Achse 
	\draw[->] (0,-2) -- (0,4) node[left,yshift=-1em] {$f(x)$};%node[above left]
	%Zahlen auf y-Achse
	\foreach \y in {-1,...,-1}
	\draw[shift={(0,\y)},color=black] (2pt,0pt) -- (-2pt,0pt) node[left] {\footnotesize $\y$};
	\foreach \y in {1,...,3}
	\draw[shift={(0,\y)},color=black] (2pt,0pt) -- (-2pt,0pt) node[left] {\footnotesize $\y$};
	%Ursprung
	\draw[color=black] (0pt,-10pt) node[right] {\footnotesize $0$};
	
	
	
	\draw[thick,color=blue] plot[samples=200, domain=0:4] 
	(\x,{sqrt(\x)}) node[below, xshift=1em]{$f(x)$};
	
	\draw[thick,color=red] plot[samples=200, domain=-1.5:1.5] 
	(\x,{\x});
	\draw[thick,color=red] plot[samples=200, domain=-1.5:1] 
	(\x,{-2*\x});
	\draw[thick,color=red] plot[samples=200, domain=-1.5:1.5] 
	(\x,{-\x});
	\draw[thick,color=red] plot[samples=200, domain=-1:1.5] 
	(\x,{2*\x});
	\end{tikzpicture}\\
\end{center}
Bei Null können beliebig viele Tangenten angelegt werden.\\
\subsection{Stetigkeit und Differenzierbarkeit}
\href{https://youtu.be/Dx6XH3DJ3GA}{Video 05}\\
Satz:\\
\begin{center}
	\boxed{\text{Jede differenzierbare Funktion ist stetig. }}\\
\end{center}
Beispiel:\\
\textbullet\;\;$ f:\mathbb{R}\longmapsto [o,\infty[ $ mit $ f(x)=x^2 $\\
\begin{quote}
	Diese Funktion ist für alle $ x\in \mathbb{R} $ diffbar also auch in allen $ x\in \mathbb{R} $ stetig.\vspace*{2mm}\\
	Bemerkung:\\
	Ist eine Funktion diffbar dann ist sie auch stetig. So kann auch eine Stetigkeit überprüft werden.\vspace*{2mm}\\
	
\end{quote}
\textbullet\;\;$ f:\mathbb{R}\longmapsto[o,\infty[ $ mit $ f(x) = |x|    \left\{ \begin{array}{rcl}
x & \mbox{für} & x\geq0\\
-x & \mbox{für} & x<0 \\

\end{array}\right. $
\begin{quote}
	$\blacktriangleright$ Grenzwert in $ x_0=0 $ bestimmen:\\
	\begin{align*}
	&\lim\limits_ {\substack{x \to 0 \\ x<0}}|x|=\lim\limits_ {x\to 0}=0=\lim\limits_ {x\to 0}=\lim\limits_ {\substack{x \to 0 \\ x>0}}|x|\\
	&\Rightarrow |x| \;\;\text{in }\;\;x_0=0\;\;\text{stetig }\\
	\end{align*}
	$\blacktriangleright$ linksseitigen Grenzwert in $ x_0=0 $ vom Differentialquotienten bestimmen:\\
	\begin{align*}
	&\lim\limits_ {x\to x_0}=\frac{f(x)-f(x_0)}{x-x_0}\\
	&\lim\limits_ {\substack{x \to 0 \\ x<0}}\frac{|x|-|0|}{x-0}=\lim\limits_ {\substack{x \to 0 \\ x<0}}\frac{|x|}{x}=\lim\limits_ {x\to 0}\frac{-x}{x}=\lim\limits_ {x\to 0}-1=-1\\
	\end{align*}
	$\blacktriangleright$ rechtsseitigen Grenzwert in $ x_0=0 $ vom Differentialquotienten bestimmen:\\
	\begin{align*}
	&\lim\limits_ {\substack{x \to 0 \\ x>0}}\frac{|x|}{x}=\lim\limits_ {x\to 0}\frac{x}{x}=\lim\limits_ {x\to 0}1=1\\
	\end{align*}
	\begin{quote}
		Der Grenzwert existiert nicht weil der linkseitige und der rechtsseitige Grenzwert nicht gleich ist. \\
	\end{quote}
	\begin{align*}
	&\lim\limits_ {\substack{x \to 0 \\ x<0}}\frac{-x}{x}=-1\neq1=\lim\limits_ {\substack{x \to 0 \\ x>0}}\frac{x}{x}\\
	&\Rightarrow\text{ $ f $ ist nicht diffbar. }\\
	\end{align*}
	
\end{quote}
Fazit:\\
Nicht jede stetige Funktion ist differenzierbar.
\subsubsection{Stetigkeit und Differenzierbarkeit Aufgabe}
\href{https://youtu.be/5wELQFVR_XA}{Video 06}\\
Für welche $ a,b\in \mathbb{R} $ ist $ f $ in $x_0=0$ diffbar?\\
\begin{align*}
f:\mathbb{R}\longmapsto \mathbb{R}\;\;\text{mit }\;\;f(x) =     \left\{ \begin{array}{rcl}
x^3\cdot \mathrm{e}^{\sin^2(x)}+a & \mbox{für} & x>0\\
xb+b & \mbox{für} & x\leq 0 \\
\end{array}\right.
\end{align*}
Zuerst Stetigkeitsuntersuchung durchführen, weil jede diffbare Funktion stetig ist. Dann die Differenzierbarbarkeit mittels dem Differentialquotienten überprüfen.\vspace*{2mm}\\
$\blacktriangleright$ Zur Stetigkeit an der Stelle $ x_0=0 $:\\
\begin{quote}
	Es muß gelten:\\
	\begin{align*}
	\lim\limits_ {x\to\ 0}f(x)=f(x_0)=f(0)\\
	\end{align*}
	\begin{align*}
	&\lim\limits_ {\substack{x \to 0 \\ x<0}}f(x)=\lim\limits_ {x\to 0}xb+b=\lim\limits_ {x\to0}b=b\overset{!}{=}\lim\limits_ {\substack{x \to 0 \\ x>0}}f(x)=\lim\limits_ {x\to 0}(x^3\cdot \mathrm{e}^{\sin^2(x)}+a)=\lim\limits_ {x\to 0}a=a\\
	&\Rightarrow a=b\\
	\end{align*}
	Bemerkung zu der Notation \grqq$\overset{!}{=}  $\glqq :  Dies bedeutet die Forderung: \grqq  soll gleich sein\glqq\\
\end{quote}
$\blacktriangleright$ Zur Diffbarkeit an der Stelle $ x_0=0 $:\\
\begin{quote}
	Es muß gelten:\\
	\begin{align*}
	\lim\limits_ {x\to 0}\frac{f(x)-f(x_0)}{x-x_0}=f'(x_0)\in\mathbb{R}\\
	\end{align*}
	\begin{align*}
	&\lim\limits_ {\substack{x \to 0 \\ x<0}}\frac{f(x)-f(x_0)}{x-x_0}=\lim\limits_ {x\to 0}\frac{xb+b-b}{x-0}=b\overset{!}{=}	\lim\limits_ {\substack{x \to 0 \\ x>0}}\frac{f(x)-f(x_0)}{x-x_0}=\lim\limits_ {x\to 0}\frac{x^3\cdot \mathrm{e}^{\sin^2(x)}+a-a}{x-0}=0\\
	&\Rightarrow b=0\\
	\end{align*}
\end{quote}
Fazit:\\
Die Funktion ist in $ x_0=0 $ differenzierbar für $ a,b=0 $.\\
\newpage
\noindent
\subsection{Elementare Ableitungen}
\href{https://youtu.be/V_SmItBir8U}{Video 07}\\
(i) Für $n\in \mathbb{Z}  $ und $ x $ beliebig gilt:\\
\begin{quote}
	\begin{align*}
	(x^n)'=n\cdot x^{n-1}
	\end{align*}
	(a) $ (x^4)'=4\cdot x^{4-1}=4\cdot x^3 $\\
	(b) $ (\frac{1}{x^3})'=(x^{-3})'=-3\cdot x^{-3-1}=-3\cdot x^{-4}=\frac{-3}{x^4} $\\
\end{quote}
(ii) Für $ a\in\mathbb{R} $ und $ x>0 $ gilt:\\
\begin{quote}
	\begin{align*}
	(x^{a})'=a\cdot x^{a-1}\\
	\end{align*}
	Bemerkung: $ x^{\frac{1}{2}}=\sqrt{x}\Rightarrow x>0 $\vspace*{4mm}\\
	 $ (\sqrt{x})'=(x^{\frac{1}{2}})' = \frac{1}{2}\cdot x^{\frac{1}{2}-1}=\frac{1}{2}\cdot x^{-\frac{1}{2}}=\frac{1}{2}\cdot \frac{1}{x^{\frac{1}{2}}}=\frac{1}{2\sqrt{x}}$\\
\end{quote}
(iii) Es gilt $ (\mathrm{e}^x)' =\mathrm{e}^x$ ,für  $ a\in \mathbb{R}\;\;(\mathrm{e}^{ax})'=a\cdot \mathrm{e}^{ax} $ \\
\begin{quote}
  $ (a^x)'=(\mathrm{e}^{x\cdot ln(a)})'=\ln(a)\cdot e^{x\cdot \ln(a)}=\ln(a)\cdot a^x$ mit $ a>0 $\\
	Bemerkung:\\
	$ \mathrm{e}^{x\cdot \ln(a)}=(\underbrace {\mathrm{e}^{\ln(a)}}_{=a})^x=a^x $\\
	oder:\\
	$a=10^x$\\
	$ \Leftrightarrow x=\log_{10}(a) $\\
	$\Rightarrow 10^x=a=10^{\log_{10}(a)}$
\end{quote}
(IV) $ (\sin(x))'=\cos(x) $\\
(V)\;\;$ (\cos(x))'=-\sin(x) $\\
\subsection{Rechenregeln für diffbare Funktionen}
\subsubsection{Die Summenregel}
\href{https://youtu.be/hjYhkf3M6wA}{Video 08}\\
Sind $ f $ und $ g $ Funktionen und in $ x\in\mathbb{R} $ diffbar , dann ist $ \lambda f+\mu g $ in $ x\in\mathbb{R} $ mit $ \lambda,\mu \in \mathbb{R} $ diffbar mit:\\
\begin{center}
	\fbox{\parbox  {60mm} {$(\lambda f+\mu g)'(x)=\lambda f'(x)+ \mu g'(x)$}}
\end{center}
\newpage
\noindent
Beispiel:\\
$\blacktriangleright$ Es sei $ f(x)=x^4$ und $ g(x)=x^2+1 $. 
Differenziere: $ \big(3f(x)+2g(x)\big)' $\\
\begin{align*}
\big(3f(x)+2g(x)\big)'&=(3\cdot x^4+2(x^2+1))'=3(x^4)'+2(x^2+1)'\\
&=3\cdot 4 x^3+2((x^2)'+(1)')=12x^3+2(2x+0)\\
&=12x^3+4x\\
\end{align*}
$ \blacktriangleright $	 Differenziere: $ (2\cdot 4^x+7\sin(x))' $\\
\begin{align*}
(2\cdot 4^x+7\cdot \sin(x))'&=2\cdot(4^x)'+7\cdot(\sin(x))'=2\cdot 4^x\cdot \ln(4)+7\cdot \cos(x)\\
\end{align*}
\subsubsection{Die Produktregel}
\href{https://youtu.be/F_hkOScmyQQ}{Video 09}\\
Sind $ f $ und $ g $ Funktionen und in $ x\in\mathbb{R} $ diffbar, dann ist $ f\cdot g $ in $ x\in\mathbb{R} $ diffbar mit:\\
\begin{center}
	\fbox{\parbox  {70mm} {$(f\cdot g)'(x)=f'(x)\cdot g(x) + f(x)\cdot g'(x)$}}\\	
\end{center}
Beispiel:\\
$\blacktriangleright$ Es sei $ f(x)=x^2$ und $ g(x)=\frac{1}{x} $. Bilde: $ (f\cdot g)'(x) $\\
\begin{align*}
(f\cdot g)'(x)&=\bigg(x^2\cdot \frac{1}{x}\bigg)'=(x^2)'\cdot \frac{1}{x}+x^2\cdot\bigg(\frac{1}{x}\bigg)'\\
&=2x\cdot \frac{1}{x}+x^2\cdot \bigg(-\frac{1}{x^2}\bigg)\\
&=2-1\\
&=1\\
\end{align*}
\begin{quote}
	oder:\\
\end{quote}
\begin{align*}
(f\cdot g)'(x)&=\bigg(x^2\cdot \frac{1}{x}\bigg)'=(x)'=1\\
\end{align*}
\subsubsection{Die Quotientenregel}
\href{https://youtu.be/tA1GHaDpCYE}{Video 10}\\
Sind die Funktionen $ f $ und $ g $ und in $ x\in\mathbb{R} $ diffbar sowie $ g(x)\neq 0 $, dann ist $ \frac{f}{g} $ in $ x\in\mathbb{R} $ diffbar mit:\\
\begin{center}
	\fbox{\parbox  {50mm} {$\big(\genfrac{}{}{}{2}{f }{g }\big)'(x)=\frac{f'(x)g(x)-f(x)g'(x)}{g(x)^2} $}}	
\end{center}
\newpage
\noindent
Beispiel:\\
$\blacktriangleright$ Es sei $ f(x)=\sin(x) $ und $ g(x) =\mathrm{e}^x$ und bilde damit $ \bigg( \frac{f(x)}{g(x)}\bigg)' $.\\
\begin{quote}
	Da $ f $ und $ g $ für alle $ x\in\mathbb{R} $ diffbar ist und $ g(x)>0 $ ist. Gilt $\frac{f}{g}  $ ist für alle $ x\in\mathbb{R} $ diffbar.\\
	Es gilt:\\
	\begin{align*}
	\bigg(\frac{f}{g}\bigg)'&=\frac{f'(x)g(x)-f(x)g'(x)}{g(x)^2}\\
	&=\frac{\cos(x)\mathrm{e}^x-\sin(x)\mathrm{e}^x}{(\mathrm{e}^x)^2}\\
	&=\frac{\mathrm{e}^x(\cos(x)-\sin(x))}{\mathrm{e}^{2x}}\\
	&=\frac{\cos(x)-\sin(x)}{\mathrm{e}^x}\\
	\end{align*}
\end{quote}
\subsubsection{Die Kettenregel}
\href{https://youtu.be/tA1GHaDpCYE}{Video 11}\\
Sei $ g $ in $ x $ diffbar und $ f $ in $ g(x) $, dann ist auch $ f\circ g $ in $ x $ diffbar und es gilt:
\begin{center}
	\fbox{\parbox  {50mm} {$(f\circ g)'(x)=f'(g(x))\cdot g'(x)$}}	
\end{center}
Beispiel:\\
$ f(x)=x^2 $ und $ g(x)=x^2+1 $\\
Dann ist:\\
\begin{align*}
&(f\circ g)'(x)=(f(g(x)))'=(f(x^2+1))'=((x^2+1)^2)'=f'(g(x))\cdot g'(x)\\
& f'=2x \Rightarrow f'(g(x))=2(g(x))=2(x^2+1)\\
& g'=2x\\
&\Rightarrow\\
&f'(g(x))\cdot g'(x)=2(x^2+1)\cdot 2x=4x^3+4x\\
\end{align*}
Was macht man eigentlich?\\
Man bildet $ f'(x) $ und da $ x = g(x) $ ist, setzt man $ g(x) $ in $ f'(x) $ ein und multipliziert $f'(g(x))  $ mit $ g'(x) $.\vspace*{2mm}\\
Beispiel:\\
Bilde $ (\mathrm{e}^{7x})' $.\\
\begin{quote}
	Es gilt:\\
	\begin{align*}
	f(x) &= \mathrm{e}^x\\
	f'(x)&=\mathrm{e}^x\\
	g(x)&=7x\\
	g'(x)&=7\\
	\end{align*}
	\begin{align*}
	(\mathrm{e}^{7x})'=f'(g(x))\cdot g'(x)=(\mathrm{e}^{7x})\cdot 7\\
	\end{align*}
	
	
\end{quote}
\subsubsection{Ableitung der Umkehrfunktion}
\href{https://youtu.be/1dFQ3JSPOpk}{Video 12}\\
Sei $ f:\mathbb{R}\supset \mathbb{I} \longmapsto \mathbb{J}\;\subseteq \mathbb{R}$ diffbar und umkehrbar sowie $ f':\mathbb{J}\longmapsto \mathbb{I} $ die Umkehrfunktion auch diffbar, dann ist:\\

\begin{center}
	\fbox{\parbox  {80mm} {$f'(f^{-1}(x))=\frac{1}{f'(f^{-1}(x))}$ mit $ f'(f^{-1}(x))\neq 0 $}}\\	
\end{center}
Es gilt:\\
\begin{align*}
x=f(f^{-1}(x))
\end{align*}
Beispiel: $ f(x)=x^2 $ dann ist $ f^{-1}=\sqrt{x} $\\
\begin{align*}
x=f(f^{-1}(x))=(\sqrt{x})^2=x\\
\end{align*}
Durch die  Ableitung von $ x=f(f^{-1}(x)) $ erhält man:\\
\begin{align*}
1=f'(f^{-1}(x))\cdot (f^{-1}(x))'\\
\Leftrightarrow \frac{1}{f'(f^{-1}(x))}=(f^{-1}(x))'\\
\end{align*}
Beispiel:\\
Betrachte $ f:\mathbb{R}\longmapsto [o,\infty[ $ mit $ x\longmapsto x^2 $\\
Dann ist:\\
\begin{align*}
&f^{-1}(x)=\sqrt{x}\;\;\text{mit }\;\;f^{-1}:[0,\infty[\;\longmapsto [0,\infty[\\
&f'(x)=(x^2)'=2x\\
&\Rightarrow f'(f^{-1}(x))=2\cdot \sqrt{x}\\
\end{align*}
Also, dann ist die Ableitung der Umkehrfunktion $ f^{-1}(x) $:\\
\begin{align*}
(f^{-1}(x))'&=\frac{1}{f'(f^{-1}(x))}\\
(\sqrt{x})'&=\frac{1}{2\cdot \sqrt{x}}\;\;\text{gilt nur wenn }\;\;2\sqrt{x}\neq 0\\
\end{align*}
\newpage
\noindent
\paragraph{Ableitung der Umkehrfunktion Aufgabe}$  $\vspace*{3mm}\\
\href{https://youtu.be/krjiEulDq-g}{Video 13}\\
Leite die Ableitung von $ \ln(x) $ her.\\
Es gilt mit:\\
\begin{align*}
f(x)&=\mathrm{e}^x\\
f^{-1}(x)&=\ln(x)\\
f'(x)&=\mathrm{e}^x\\
\end{align*}
Dann ist:\\
\begin{align*}
(f^{-1}(x))'=\frac{1}{f'(f^{-1}(x))}=\frac{1}{\mathrm{e}^{\ln(x)}}=\frac{1}{x}\;\;\text{mit } x\neq 0
\end{align*}
Fazit:\\
\begin{align*}
(\ln(x))'=\frac{1}{x}\;\;\text{mit }\;\;x\in \;]0,\infty[\\
\end{align*}
Bemerkung:\\
Für $ \ln(x) $ gilt $ \mathbb{D}_f=]0,\infty[ $
\subsubsection{Aufgaben zum Thema: Ableitungsregeln}
\paragraph{Aufgabe 1 zum Thema: Ableitungsregeln}$  $\vspace*{3mm}\\
\href{https://youtu.be/XE6Zkj9KDfg}{Video 14}\\
Berechne zu folgender Funktion die Ableitung: $ f(x)=x^3\cdot \mathrm{e}^{-2x} $\\
\begin{align*}
f'(x)=(x^3\cdot \mathrm{e}^{-2x})'&=(x^3)'\cdot \mathrm{e}^{-2x}+ x^3\cdot (\mathrm{e}^{-2x})'\;\;\text{\tiny Produktregel }\\
&=3x^2\cdot \mathrm{e}^{-2x}+x^3\cdot \mathrm{e}^{-2x}\cdot (-2x)'\;\;\text{\tiny Kettenrregel }\\
&=3x^2\cdot \mathrm{e}^{-2x}+ x^3\cdot \mathrm{e}^{-2x}\cdot (-2)\\
&=3x^2 \mathrm{e}^{-2x}-2 x^3 \mathrm{e}^{-2x} \\
&=x^2 \mathrm{e}^{-2x}(3-2x)\\
\end{align*}
\newpage
\noindent
\paragraph{Aufgabe 2 zum Thema: Ableitungsregeln}$  $\vspace*{3mm}\\
\href{https://youtu.be/RJ5ZqJzbv_I}{Video 15}\\
Berechne zu folgender Funktion die Ableitung: $ f(x)=\frac{\sin(x)}{\sqrt{1-x^2}} $\\
\begin{align*}
f'(x)=\bigg(\frac{\sin(x)}{\sqrt{1-x^2}}\bigg)'&=\frac{(\sin(x))'\sqrt{1-x^2}-\sin(x)(\sqrt{1-x^2})'}{(\sqrt{1-x^2})^2}\;\;\text{\tiny Quotientenregel }\\
&=\frac{\cos(x)\sqrt{1-x^2}-\sin(x)\bigg((1-x^2)^{-\frac{1}{2}}\bigg)'}{1-x^2}\\
&=\frac{\cos(x)\sqrt{1-x^2}-\sin(x)\frac{1}{2}(1-x^2)^{\frac{1}{2}-1}(-2x)}{1-x^2}\\
&=\frac{\cos(x)\sqrt{1-x^2}-\sin(x)\frac{-2x}{2\sqrt{1-x^2}}}{1-x^2}\\
&=\frac{\cos(x)\sqrt{1-x^2}+\sin(x)\frac{x}{\sqrt{1-x^2}}}{1-x^2}\\
\end{align*}
\paragraph{Aufgabe 3 zum Thema: Ableitungsregeln}$  $\vspace*{3mm}\\
\href{https://youtu.be/uPo2hsEsAdU}{Video 16}\\
Berechne zu folgender Funktion die Ableitung: $ f(x)=x^{\cos(x)} $\\
Es gilt:\\
\begin{align*}
&x=\mathrm{e}^{\ln(x)}\\
&\text{d.h. }\\
&x^{\cos(x)}=\big(\mathrm{e}^{\ln(x)}\big)^{\cos(x)}=\mathrm{e}^{\ln(x)\cos(x)}\\
\end{align*}
Die Ableitung von $\mathrm{e}^{\ln(x)\cos(x)}  $ ist dann:\\
\begin{align*}
f'(x)=\bigg(\mathrm{e}^{\ln(x)\cos(x)}\bigg)'&=\mathrm{e}^{\ln(x)\cos(x)}\cdot (\ln(x)\cos(x)'\;\;\text{\tiny Kettenrregel dann Produktregel }\\
&=\mathrm{e}^{\ln(x)\cos(x)}\cdot \frac{1}{x}\cos(x)+\ln(x)(-\sin(x))\;\;|\;\;\text{\tiny $ x=\mathrm{e}^{\ln(x)} $ }\\
&=x^{\cos(x)}\bigg(\frac{\cos(x)}{x}-\ln(x)\sin(x)\bigg)\\
\end{align*}
\paragraph{Aufgabe 1 Ableitungsregeln/maximaler Definitionsbereich Teil I u.II}$  $\vspace*{3mm}\\
\href{https://youtu.be/fz6oN4-i_84}{Video 17}\\
Bestimme für folgende Funktion den maximalen Definitionsbereich und berechne deren Ableitung.\\
\begin{align*}
f(x)=\frac{e^{(\sin(x^2)+\frac{1}{x})}}{x^2+3x+2}\\
\end{align*}
\newpage
\noindent
$\blacktriangleright$ Zum max. Definitionsbereich:\\
\begin{quote}
	Beim Zähler muß $ \frac{1}{x}\neq 0 $ sein.\\
	Der Nenner darf nicht Null werden, deshalb muß man $ x^2+3x+2 $ auf Nullstellen untersuchen, denn $ f $ ist bei den Nullstellen nicht definiert.\\
	Bestimmung der Nullstellen mit der $ p-q- $Formel:\\
	\begin{align*}
	&x_{1,2}=-\frac{p}{2}\pm\sqrt{\bigg(\frac{p}{q}\bigg)^2-q}=-\frac{3}{2}\pm \sqrt{\bigg(\frac{3}{2}\bigg)^2-2}=-\frac{3}{2}\pm \sqrt{\frac{9}{4}-\frac{8}{4}}=-\frac{3}{2}\pm \sqrt{\frac{1}{4}}\\
	&x_1=-2\\
	&x_2=-1\\
	\end{align*}
	Fazit:\\
	Der Max. Definitionsbereich von $ f $ ist:\\
	\begin{align*}
	\mathbb{D}_f=\mathbb{R}\setminus \{-2,-1,0\}\\
	\end{align*}
\end{quote}
\href{https://youtu.be/2hB1iRS_O8w}{Video 18}\\
$ \blacktriangleright $ Ableitung von $ f(x)=\frac{e^{(\sin(x^2)+\frac{1}{x})}}{x^2+3x+2} $\\
\begin{align*}
f'(x)=\frac{\bigg(e^{(\sin(x^2)+\frac{1}{x})}\bigg)'(x^2+3x+2)-\mathrm{e}^{(\sin(x^2)+\frac{1}{x})}(x^2+3x+2)'}{(x^2+3x+2)^2}\\
\end{align*}
\begin{quote}
	NR:\\
	\begin{align*}
	\bigg(\mathrm{e}^{(\sin(x^2)+\frac{1}{x})}\bigg)'&=(\mathrm{e}^{(\sin(x^2)+\frac{1}{x})})\cdot \bigg(\sin(x^2)+\frac{1}{x}\bigg)'\;\;\text{\tiny $ \bigg(\frac{1}{x}\bigg)'=\bigg(x^{-1}\bigg)'=-1x^{-2}=-\frac{1}{x^2} $ }\\
	&=(\mathrm{e}^{(\sin(x^2)-\frac{1}{x})})\cdot \bigg(\cos(x^2)2x+\frac{-1}{x^2}\bigg)\\
	\bigg(\mathrm{e}^{(\sin(x^2)+\frac{1}{x})}\bigg)'&=(\mathrm{e}^{(\sin(x^2)-\frac{1}{x})})\cdot 	\bigg(\cos(x^2)2x-\frac{1}{x^2}\bigg)\\
		\end{align*}
	NR:\\
	\begin{align*}
	(x^2+3x+2)'&=2x+3\\
	\end{align*}
	Nun ergibt sich:\\
	\begin{align*}
	f'(x)&=\frac{(\mathrm{e}^{(\sin(x^2)+\frac{1}{x})})\cdot (\cos(x^2)2x-\frac{1}{x^2})(x^2+3x+2)-\mathrm{e}^{(\sin(x^2)+\frac{1}{x})}(2x+3)}{(x^2+3x+2)^2}\\
	&=\mathrm{e}^{(\sin(x^2)+\frac{1}{x})} \cdot \bigg[\frac{ (\cos(x^2)2x-\frac{1}{x^2})(x^2+3x+2)-(2x+3)}{(x^2+3x+2)^2}\bigg]\\
	\end{align*}
\end{quote}
\paragraph{Aufgabe 2 Ableitungsregeln/maximaler Definitionsbereich Teil I u.II}$  $\vspace*{3mm}\\
\href{https://youtu.be/_vTSb95KkUg}{Video 19}\\
Bestimme für folgende Funktion den maximalen Definitionsbereich und berechne deren Ableitung.\\
\begin{align*}
f(x)=\ln\bigg(\frac{x^2+1}{x^2-1}\bigg)-\ln\bigg(\frac{1}{x^2}\bigg)\\
\end{align*}
$ \blacktriangleright $ Zum max. Definitionsbereich:\\
\begin{quote}
	Da $ \ln(x):\;]0,\infty[\;\longmapsto \mathbb{R}$ muß gelten:\vspace*{2mm}\\
	(i) $ \frac{x^2+1}{x^2-1}> 0 $ und $ x^2-1\neq 0 $ \vspace*{2mm}\\
	(ii) $\frac{1}{x^2}>0$ und $ x\neq 0 $\vspace*{2mm}\\
	zu (i):\\
	\begin{quote}
		Der Zähler $ x^2+1 $ ist immer positiv, also kann nur noch der Nenner kleiner Null werden.\\
		Bestimme $ x\in \mathbb{R} $ mit $ \frac{1}{x^2-1}>0 $\\
		Die Nullstellen sind +1 und -1.\\ Dann gilt:\\
		\begin{quote}
			Für $ x>1 $ ist $ \frac{1}{x^2-1} $  immer positiv.\\
			Für $ x<1 $ ist $ \frac{1}{x^2-1} $  immer positiv.\\
			Für $ 1>x>-1 $ ist $ \frac{1}{x^2-1} $  immer negativ.\\
		\end{quote}
		
		\begin{center}
			Mit den Nullstellen -1,1 auf dem Zahlenstrahl:\\
			\begin{tikzpicture}
			
			% Gitternetzlinien
			%\draw[xstep=1cm, ystep=1cm, very thin, color=lightgray] (-4,4) grid (4,-2);
			% x-Achse
			\draw[->] (-4,0) -- (4,0) node[below,xshift=-1em] {$x$}; 
			%Zahlen auf x-Achse
			\foreach \x in {-3,...,3}
			\draw[shift={(\x,0)},color=black] (0pt,2pt) -- (0pt,-2pt) node[below] {\footnotesize $\x$};
			
			%\draw[color=black] (0pt,-10pt) node[right] {\footnotesize $0$};
			
			\draw [color=blue] (-2,0.5) node[right] {\footnotesize $(+)$};;
			\draw [color=blue] (-0.60,0.5) node[right] {\footnotesize $(-)$};;
			\draw [color=blue] (1.2,.5) node[right] {\footnotesize $(+)$};;
			\draw [color=blue] (-1,0) -- (-1,1);
			\draw [color=blue] (1,0) -- (1,1);
			
			\end{tikzpicture}\\
		\end{center}
		
		Der Definitionsbereich für $ g(x)=\ln(\frac{x^2+1}{x^2-1}) $ ist:\\
		\begin{align*}
		\mathbb{D}_g= \; ]-\infty,-1[\;\;\cup \;\;]1,+\infty[
		\end{align*}
	\end{quote}
	zu (ii):\\
	\begin{quote}
		Bestimme $ x\in \mathbb{R} $ mit $ \frac{1}{x^2}>0 $\\
		Der Definitionsbereich für $h(x)=\frac{1}{x^2}$ ist:\\
		\begin{align*}
		\mathbb{D}_h=\mathbb{R}\setminus \{0\}\\
		\end{align*}	
	\end{quote}
\newpage
\noindent
	Der Definitionsbereich von $ f(x)=\ln\bigg(\frac{x^2+1}{x^2-1}\bigg)-\ln\bigg(\frac{1}{x^2}\bigg) $ ist nun die Schnittmenge von $ \mathbb{D}_g $ und $ \mathbb{D}_h $.\\
	\begin{align*}
	\mathbb{D}_f&=\mathbb{D}_g\cap\mathbb{D}_h\\
	&=\big(]-\infty,-1[\;\;\cup \;\;]1,+\infty[\big)\;\;\cap\;\;\mathbb{R}\setminus \{0\}\\
	\end{align*}
	Fazit:\\
	\begin{align*}
	\mathbb{D}_f&=]-\infty,-1[\;\;\cup \;\;]1,+\infty[
	\end{align*}
\end{quote}
\href{https://youtu.be/3Xab3Mt4204}{Video 20}\\
$ \blacktriangleright $ Ableitung von $ f(x)=\ln\bigg(\frac{x^2+1}{x^2-1}\bigg)-\ln\bigg(\frac{1}{x^2}\bigg) $\\
\begin{quote}
	\begin{align*}
	f'(x)&=\bigg (\ln\bigg(\frac{x^2+1}{x^2-1}\bigg)-\ln\bigg(\frac{1}{x^2}\bigg)\bigg)'\\
	&=\bigg(\ln\bigg(\frac{x^2+1}{x^2-1}\bigg)\bigg)'-\bigg(\ln\bigg(\frac{1}{x^2}\bigg)\bigg)'\;\;\text{\tiny Summenregel }\\
	\end{align*}
	NR.:\\
	\begin{align*}
	\bigg(\ln\bigg(\frac{x^2+1}{x^2-1}\bigg)\bigg)'&=\frac{1}{\frac{x^2+1}{x^2-1}}\cdot \bigg(\frac{x^2+1}{x^2-1}\bigg)'=\frac{x^2-1}{x^2+1}\cdot\frac{(x^2+1)'(x^2-1)-(x^2+1)(x^2-1)'}{(x^2-1)^2}\\
	&=\frac{x^2-1}{x^2+1}\cdot\frac{2x(x^2-1)-2x(x^2+1)}{(x^2-1)^2}\\
	&=\frac{x^2-1}{x^2+1}\cdot\frac{2x^3-2x-2x^3-2x}{(x^2-1)^2}=\frac{(x^2-1)(-4x)}{(x^2+1)(x^2-1)^2}\\
	&=\frac{-4x}{(x^2+1)(x^2-1)}=\frac{-4x}{x^4-1}
	\end{align*}
	Nun ergibt sich:\\
	\begin{align*}
	f'(x)&=\frac{-4x}{x^4-1}-\bigg(\ln\bigg(\frac{1}{x^2}\bigg)\bigg)'\\
	&=\frac{-4x}{x^4-1}-\frac{1}{\frac{1}{x^2}}\cdot\bigg(\frac{1}{x^2}\bigg)'\;\;\text{\tiny NR.: $\bigg(\frac{1}{x^2}\bigg)'=\frac{-2}{x^3} $}\\
	&=\frac{-4x}{x^4-1}-\frac{1}{\frac{1}{x^2}}\cdot\bigg(\frac{-2}{x^3}\bigg)=\frac{-4x}{x^4-1}-x^2\bigg(\frac{-2}{x^3}\bigg)\\
	&=\frac{-4x}{x^4-1}+\frac{2}{x}
	\end{align*}
\end{quote}
\subsection{Das notwendige Kriterium für Extrema}
\href{https://youtu.be/IkfZZjapfWk}{Video 21}\\
Sei $ f: I\longmapsto \mathbb{R} $ diffbar auf $ I\subseteq \mathbb{R} $ und $ x_0 $ ein innerer Punkt (kein Randpunkt vom Intervall $ [a,b] $).\\
Hat $ f  $ ein Maximum oder Minimum in $ x_0 $, dann gilt:\\
\begin{align*}
f'(x_0)=0
\end{align*}
Beispiel:\\
\textbullet\;\; $ f:\mathbb{R}\longmapsto[0,\infty[ $ mit $ x\longmapsto x^2 $\\
\begin{quote}
	\begin{center}
		\begin{tikzpicture}
		
		% Gitternetzlinien
		\draw[xstep=1cm, ystep=1cm, very thin, color=lightgray] (-4,4) grid (4,-2);
		% x-Achse
		\draw[->] (-4,0) -- (4,0) node[below,xshift=-1em] {$x$}; 
		%Zahlen auf x-Achse
		\foreach \x in {-3,...,-1}
		\draw[shift={(\x,0)},color=black] (0pt,2pt) -- (0pt,-2pt) node[below] {\footnotesize $\x$};
		\foreach \x in {1,...,3}
		\draw[shift={(\x,0)},color=black] (0pt,2pt) -- (0pt,-2pt) node[below] {\footnotesize $\x$};
		% y-Achse 
		\draw[->] (0,-2) -- (0,4) node[left,yshift=-1em] {$y$};%node[above left]
		%Zahlen auf y-Achse
		\foreach \y in {-1,...,-1}
		\draw[shift={(0,\y)},color=black] (2pt,0pt) -- (-2pt,0pt) node[left] {\footnotesize $\y$};
		\foreach \y in {1,...,3}
		\draw[shift={(0,\y)},color=black] (2pt,0pt) -- (-2pt,0pt) node[left] {\footnotesize $\y$};
		%Ursprung
		\draw[color=black] (0pt,-22pt) node[right] {\footnotesize $x_0=0$};
		
		%\draw[thick,color=black] plot[samples=200, domain=2:4] 
		%(\x,{2}) node[below, xshift=-5em] {$h(x)$};
		
		\draw[thick,color=black] plot[samples=200, domain=-2:2] 
		(\x,{\x*\x}) node[below, xshift=1em]{$f(x)$};
		
		
		
		\end{tikzpicture}\\
	\end{center}
	Es gilt:\\
	\begin{align*}
	f'(x_0)=f'(0)=2x=2\cdot 0=0\\
	\end{align*}
	Warum muß hier bei $ x_0=0 $ die Ableitung Null sein?\\
	\begin{center}
		\begin{tikzpicture}
		
		% Gitternetzlinien
		\draw[xstep=1cm, ystep=1cm, very thin, color=lightgray] (-4,4) grid (4,-2);
		% x-Achse
		\draw[->] (-4,0) -- (4,0) node[below,xshift=-1em] {$x$}; 
		%Zahlen auf x-Achse
		\foreach \x in {-3,...,-1}
		\draw[shift={(\x,0)},color=black] (0pt,2pt) -- (0pt,-2pt) node[below] {\footnotesize $\x$};
		\foreach \x in {1,...,3}
		\draw[shift={(\x,0)},color=black] (0pt,2pt) -- (0pt,-2pt) node[below] {\footnotesize $\x$};
		% y-Achse 
		\draw[->] (0,-2) -- (0,4) node[left,yshift=-1em] {$y$};%node[above left]
		%Zahlen auf y-Achse
		\foreach \y in {-1,...,-1}
		\draw[shift={(0,\y)},color=black] (2pt,0pt) -- (-2pt,0pt) node[left] {\footnotesize $\y$};
		\foreach \y in {1,...,3}
		\draw[shift={(0,\y)},color=black] (2pt,0pt) -- (-2pt,0pt) node[left] {\footnotesize $\y$};
		%Ursprung
		\draw[color=black] (0pt,-10pt) node[right] {\footnotesize $0$};
		\draw[color=blue] (-1.6,1) node[right] {\footnotesize $x_1$};
		\draw[color=blue] (1,1) node[right] {\footnotesize $x_2$};
		\draw[thick,color=blue] plot[samples=200, domain=0.2:2] 
		(\x,{\x*2-1}) node[below, xshift=2em] {$g_2(x)$};
		\draw[thick,color=blue] plot[samples=200, domain=-2:-0.2] 
		(\x,{-\x*2-1}) node[yshift=80, xshift=-5em] {$g_1(x)$};
		\draw[thick,color=black] plot[samples=200, domain=-2:2] 
		(\x,{\x*\x}) node[below, xshift=1em]{$f(x)$};
		
		
		
		\end{tikzpicture}\\
	\end{center}
	Die Steigung von $ g_1(x_1) $ und $ g_2(x_2) $ ist gleich der Ableitung von $ f'(x_1)=2x_1 $ oder $f'(x_2)=2x_2  $ und für $ x_1<0 $ ist die Steigung negativ und für $ x_2>0 $ ist die Steigung positiv, also muß die Steigung zwischen $ x_1 $ und $ x_2 $ einmal Null werden.\\
\end{quote}
\newpage
\noindent
\textbullet \;\;$ g:\mathbb{R}\longmapsto\mathbb{R} $ mit $ x\longmapsto x^3 $\\
\begin{quote}
	\begin{center}
		\begin{tikzpicture}
		
		% Gitternetzlinien
		\draw[xstep=1cm, ystep=1cm, very thin, color=lightgray] (-4,4) grid (4,-2);
		% x-Achse
		\draw[->] (-4,0) -- (4,0) node[below,xshift=-1em] {$x$}; 
		%Zahlen auf x-Achse
		\foreach \x in {-3,...,-1}
		\draw[shift={(\x,0)},color=black] (0pt,2pt) -- (0pt,-2pt) node[below] {\footnotesize $\x$};
		\foreach \x in {1,...,3}
		\draw[shift={(\x,0)},color=black] (0pt,2pt) -- (0pt,-2pt) node[below] {\footnotesize $\x$};
		% y-Achse 
		\draw[->] (0,-2) -- (0,4) node[left,yshift=-1em] {$y$};%node[above left]
		%Zahlen auf y-Achse
		\foreach \y in {-1,...,-1}
		\draw[shift={(0,\y)},color=black] (2pt,0pt) -- (-2pt,0pt) node[left] {\footnotesize $\y$};
		\foreach \y in {1,...,3}
		\draw[shift={(0,\y)},color=black] (2pt,0pt) -- (-2pt,0pt) node[left] {\footnotesize $\y$};
		%Ursprung
		\draw[color=black] (0pt,-10pt) node[right] {\footnotesize $0$};
		
		%\draw[thick,color=black] plot[samples=200, domain=-1.2:4] 
		%(\x,{2}) node[below, xshift=-5em] {$h(x)$};
		
		\draw[thick,color=black] plot[samples=200, domain=-1.2:1.5] 
		(\x,{\x*\x*\x }) node[below, xshift=1em]{$g(x)$};
		
		
		
		\end{tikzpicture}\\	
	\end{center}
	Es gilt:\\
	\begin{align*}
	g'(x_0)=3x^2_0=3\cdot 0^2=0\;\;\text{für }x_0=0\\
	\end{align*}
	$ g $ hat weder Maximum noch Minimum !\\
\end{quote}
Fazit: Die Umkehrung ist falsch. Die Aussage: Die Ableitung ist Null folgert nicht, dass hier ein Maximum oder Minimum ist.\\
\subsubsection{Das notwendige Kriterium für Extrema Aufgabe Teil I u. II}
\href{https://youtu.be/s-gd2f4tDws}{Video 22}\\
Bestimme alle globalen Extrema von der stetigen Funktion $ f:[-2,0]\longmapsto \mathbb{R} $ mit $ f(x)=\sqrt{x^3+x^2-2x} $.\\
Untersuche die inneren Punkten vom kompakten Intervall $x\in\; ]-2,0[  $ .\vspace*{4mm}\\
Berechne die Nullstellen von $ f'(x) $.\\
\begin{align*}
f'(x)=\bigg(\sqrt{x^3+x^2-2x}\bigg)'&=\frac{1}{2\sqrt{x^3+x^2-2x}}\cdot \bigg(x^3+x^2-2x\bigg)'\\
f'(x)&=\frac{3x^2+2x-2}{2\sqrt{x^3+x^2-2x}}\overset{!}{=} 0\\
\end{align*}
Bestimmung der Nullstellen vom Nenner. (Nenner darf nicht Null werden!)\\
\begin{align*}
x^3+x^2-2x=x(x^2+x-2)=x(x-1)(x+2)=0\\
\end{align*}
Die Nullstellen wären -2, 0 und 1, diese sind aber nicht im Intervall: $ \in \;]-2,0[  $ .\vspace*{4mm}\\
Bestimmung der Nullstellen vom Zähler $3x^2+2x-2=0 $ mit der p-q-Formel.\\
\begin{align*}
&3x^2+2x-2=0\\
&\Leftrightarrow x^2+\frac{2}{3}x-\frac{2}{3}x=0\\
\end{align*}
\begin{align*}
x_{1,2}&=-\frac{p}{2}\pm \sqrt{\bigg(\frac{p}{2}\bigg)^2-q}\\
&p=\frac{2}{3}\\
&q=\frac{-2}{3}\\
x_{1,2}&=-\frac{2}{3\cdot 2}\pm\sqrt{\bigg(\frac{2}{3\cdot2}\bigg)^2-(\frac{-2}{3})}\\
&=-\frac{1}{3}\pm\sqrt{\frac{1}{9}+\frac{2}{3}}\\
&=-\frac{1}{3}\pm\sqrt{\frac{7}{9}}\\
&=-\frac{1}{3}\pm \frac{1}{3}\sqrt{7}\\
&=\frac{-1\pm \sqrt{7}}{3}\\
x_1&=\frac{-1- \sqrt{7}}{3}\\
x_2&=\frac{-1+ \sqrt{7}}{3}\\
\end{align*}
Da $ x_2 $ offensichtlich größer Null ist, ist $ x\notin [-2,0] $.\\Nr: $ \sqrt{7} $ ist zwischen $ \sqrt{4} $ und $\sqrt{9}$, also zwischen $ 2 $ und $ 3 \Rightarrow x_2>0 $\\
Fazit:\\
Die einzige potentielle Extremstelle aus dem kompakten Intervall von  $ ]-2,0[  $ ist  $x_1=\frac{-1- \sqrt{7}}{3}  $.\\
\href{https://youtu.be/RrEGA0egEKQ}{Video 23}\\
Randpunktuntersuchung bei $ x_3=-2 $ und $ x_4=0$\\
\begin{align*}
&f(x)=\sqrt{x^3+x^2-2x}=\sqrt{x(x^2+x-2)}=\sqrt{x(x-1)(x+2)}\\
&\Rightarrow\\
&f(x_3)=f(-2)=0\\
&f(x_4)=f(0)=0\\
\end{align*}
Für $ x_1=\frac{-1-\sqrt{7}}{3} $ gilt:\\
\begin{align*}
&f(x_1)=\sqrt{\underbrace {x_1}_{<0}\underbrace {(x_1-1)}_{<0}\underbrace {(x_1-2)}_{>0}}\\
&\Rightarrow\\
&f(x_1)=f\bigg(\frac{-1-\sqrt{7}}{3}\bigg)>0\\
\end{align*}
In $ x_3 $ und $ x_4 $ hat $ f $ ein globales Minimum, weil $ f(x_1)>0 $ ist.\\
In $ x_1 $ hat f ein Maximum.\\
\newpage
\noindent
Bemerkung:\\
Da $ f $ stetig auf einen kompakten Intevall ist nimmt  $ f $ ein Maximum und ein Minimun an.\\
\begin{center}
	$  f:[-2,0]\longmapsto \mathbb{R} $ mit $ f(x)=\sqrt{x^3+x^2-2x} $\\
	\begin{tikzpicture}
	
	% Gitternetzlinien
	\draw[xstep=1cm, ystep=1cm, very thin, color=lightgray] (-4,4) grid (4,-2);
	% x-Achse
	\draw[->] (-4,0) -- (4,0) node[below,xshift=-1em] {$x$}; 
	%Zahlen auf x-Achse
	\foreach \x in {-3,...,-1}
	\draw[shift={(\x,0)},color=black] (0pt,2pt) -- (0pt,-2pt) node[below] {\footnotesize $\x$};
	\foreach \x in {1,...,3}
	\draw[shift={(\x,0)},color=black] (0pt,2pt) -- (0pt,-2pt) node[below] {\footnotesize $\x$};
	% y-Achse 
	\draw[->] (0,-2) -- (0,4) node[left,yshift=-1em] {$f(x)$};%node[above left]
	%Zahlen auf y-Achse
	\foreach \y in {-1,...,-1}
	\draw[shift={(0,\y)},color=black] (2pt,0pt) -- (-2pt,0pt) node[left] {\footnotesize $\y$};
	\foreach \y in {1,...,3}
	\draw[shift={(0,\y)},color=black] (2pt,0pt) -- (-2pt,0pt) node[left] {\footnotesize $\y$};
	%Ursprung
	\draw[color=black] (0pt,-10pt) node[right] {\footnotesize $0$};
	
	
	
	\draw[thick,color=blue] plot[samples=200, domain=-2:0] 
	(\x,{sqrt(\x*\x*\x+\x*\x-2*\x)}) node[below, xshift=-3em,yshift=5em]{$f(x)$};
	
	
	
	\end{tikzpicture}\\
\end{center}
\subsection{Der Mittelwertsatz (der Differentialrechnung)}
\href{https://youtu.be/VGZF-dpPEDY}{Video 24}\\
Sei $ f: \mathbb{R}\supset \mathbb{I}\longmapsto \mathbb{R} $ diffbar auf dem Intevall $ I $ und $ a,b\in \mathbb{I} $ mit $ a\neq b $, dann existiert $ x_0 $ zwischen $ a $ und $ b $ mit:\\
\begin{align*}
f'(x_0)=\frac{f(b)-f(a)}{b-a}\\
\end{align*}
\begin{center}
	\begin{tikzpicture}
	
	% Gitternetzlinien
	\draw[xstep=1cm, ystep=1cm, very thin, color=lightgray] (-4,4) grid (4,-2);
	% x-Achse
	\draw[->] (-4,0) -- (4,0) node[below,xshift=-1em] {$x$}; 
	%Zahlen auf x-Achse
	%\foreach \x in {-3,...,-1}
	%\draw[shift={(\x,0)},color=black] (0pt,2pt) -- (0pt,-2pt) node[below] {\footnotesize $\x$};
	%\foreach \x in {1,...,3}
	%\draw[shift={(\x,0)},color=black] (0pt,2pt) -- (0pt,-2pt) node[below] {\footnotesize $\x$};
	% y-Achse 
	\draw[->] (0,-2) -- (0,4) node[left,yshift=-1em] {$f(x)$};%node[above left]
	%Zahlen auf y-Achse
	%\foreach \y in {-1,...,-1}
	%\draw[shift={(0,\y)},color=black] (2pt,0pt) -- (-2pt,0pt) node[left] {\footnotesize $\y$};
	%\foreach \y in {1,...,3}
	%\draw[shift={(0,\y)},color=black] (2pt,0pt) -- (-2pt,0pt) node[left] {\footnotesize $\y$};
	%Ursprung
	\draw[color=blue] (1.5,-0.2) node[right] {\footnotesize $b$};
	\draw[color=blue] (-1,1.75) node[right] {\footnotesize $f(b)$};
	\draw[color=blue,dashed] (1.65,0)--(1.65,1.75);
	\draw[color=blue,dashed] (1.65,1.75)--(0,1.75);
	
	
	
	\draw[color=blue] (-1.7,-0.2) node[right] {\footnotesize $a$};
	\draw[color=blue] (0.1,0.55) node[right] {\footnotesize $f(a)$};
	\draw[color=blue,dashed] (-1.4,0)--(-1.4,0.6);
	\draw[color=blue,dashed] (-1.4,0.55)--(0,0.55);
	\draw[thick,color=blue] plot[samples=200, domain=-3:3] 
	(\x,{0.4*\x+1.1}) node[below, xshift=1em]{$s$};
	\draw[thick,color=black] plot[samples=200, domain=-2.5:2.5] 
	(\x,{\x*\x*\x/6+1}) node[below, xshift=1em]{$f(x)$};
	
	
	
	\end{tikzpicture}\\
\end{center}
Die Formel der Sekante $ s $ ist :\\
\begin{align*}
&s(x)=f(a)+\frac{f(b)-f(a)}{b-a}(x-a)\;\;\text{\tiny Herleitung: siehe Kapitel 8.1} \\
&\Rightarrow \text{die Steigung der Geraden ist: }\;\;\frac{f(b)-f(a)}{b-a}
\end{align*}
Also existiert mindestens eine Steigung von $ f(x_0) $ mit $ m=\frac{f(b)-f(a)}{b-a} $ zwischen $ a $ und $ b $.\\
Eine Parallelverschiebung von $ s(x) $ ergibt $ g(x_0)$ mit der Steigung $ f'(x_0) $.\\
\begin{center}
	\begin{tikzpicture}
	
	% Gitternetzlinien
	\draw[xstep=1cm, ystep=1cm, very thin, color=lightgray] (-4,4) grid (4,-2);
	% x-Achse
	\draw[->] (-4,0) -- (4,0) node[below,xshift=-1em] {$x$}; 
	%Zahlen auf x-Achse
	%\foreach \x in {-3,...,-1}
	%\draw[shift={(\x,0)},color=black] (0pt,2pt) -- (0pt,-2pt) node[below] {\footnotesize $\x$};
	%\foreach \x in {1,...,3}
	%\draw[shift={(\x,0)},color=black] (0pt,2pt) -- (0pt,-2pt) node[below] {\footnotesize $\x$};
	% y-Achse 
	\draw[->] (0,-2) -- (0,4) node[left,yshift=-1em] {$f(x)$};%node[above left]
	%Zahlen auf y-Achse
	%\foreach \y in {-1,...,-1}
	%\draw[shift={(0,\y)},color=black] (2pt,0pt) -- (-2pt,0pt) node[left] {\footnotesize $\y$};
	%\foreach \y in {1,...,3}
	%\draw[shift={(0,\y)},color=black] (2pt,0pt) -- (-2pt,0pt) node[left] {\footnotesize $\y$};
	%Ursprung
	
	\draw[thick,color=red] plot[samples=200, domain=-3:3] 
	(\x,{0.4*\x+0.75}) node[below, xshift=1em,yshift=0em]{$g(x)$};
	
	
	\draw[thick,color=blue] plot[samples=200, domain=-3:3] 
	(\x,{0.4*\x+1.1}) node[below, xshift=1em,yshift=1em]{$s(x)$};
	\draw[thick,color=black] plot[samples=200, domain=-2.5:2.5] 
	(\x,{\x*\x*\x/6+1}) node[below, xshift=1em]{$f(x)$};
	
	
	
	\end{tikzpicture}\\
\end{center}
\underline{Der Satz von Rolle}\\
Wenn gilt:\\
\begin{align*}
f(a)=f(b)\\
\end{align*}
Dann existiert ein $ x_0 $ zwischen $ a $ und $ b $ mit:\\
\begin{align*}
f'(x_0)=0\\
\end{align*}
\begin{center}
	\begin{tikzpicture}
	
	% Gitternetzlinien
	\draw[xstep=1cm, ystep=1cm, very thin, color=lightgray] (-4,4) grid (4,-2);
	% x-Achse
	\draw[->] (-4,0) -- (4,0) node[below,xshift=-1em] {$x$}; 
	%Zahlen auf x-Achse
	%\foreach \x in {-3,...,-1}
	%\draw[shift={(\x,0)},color=black] (0pt,2pt) -- (0pt,-2pt) node[below] {\footnotesize $\x$};
	%\foreach \x in {1,...,3}
	%\draw[shift={(\x,0)},color=black] (0pt,2pt) -- (0pt,-2pt) node[below] {\footnotesize $\x$};
	% y-Achse 
	\draw[->] (0,-2) -- (0,4) node[left,yshift=-1em] {$y$};%node[above left]
	%Zahlen auf y-Achse
	%\foreach \y in {-1,...,-1}
	%\draw[shift={(0,\y)},color=black] (2pt,0pt) -- (-2pt,0pt) node[left] {\footnotesize $\y$};
	%\foreach \y in {1,...,3}
	%\draw[shift={(0,\y)},color=black] (2pt,0pt) -- (-2pt,0pt) node[left] {\footnotesize $\y$};
	%Ursprung
	
	\draw[thick,color=red] plot[samples=200, domain=-3:3] 
	(\x,{2.5}) node[below,color=blue,yshift=0em]{$f'(x_0)=0$};
	
	\draw[color=blue] (0,-0.4) node {\footnotesize $x_0$};
	
	\draw[thick,color=blue] plot[samples=200, domain=-1.65:1.65] 
	(\x,{-\x*\x+2.5});
	\draw[color=red] (0,2.75) node[right] {\footnotesize $f(a)=f(b)$};
	\draw[color=red] (-1.8,-0.4) node[right] {\footnotesize $a$};
	\draw[color=red,dashed] (-1.6,0)--(-1.6,2.5);
	\draw[color=red] (1.36,-0.4) node[right,] {\footnotesize $b$};
	\draw[color=red,dashed] (1.6,0)--(1.6,2.5);
	\draw[color=white] (0,0)--(0,2.5);
	\draw[color=blue,dashed] (0,0)--(-0,2.5);
	\draw[color=blue] (1,1.9) node {\footnotesize $f$};
	\draw[color=red] (-2,2.75) node {\footnotesize $g$};
	\end{tikzpicture}\\
\end{center}
\subsubsection{Aufgabe 1 zum Mittelwertsatz}
\href{https://youtu.be/EaH9h63wHm0}{Video 25}\\
Zeige mit dem MWS, dass für $ x\geq 0: \ln(2x+1)\leq 2x $ gilt:\\
Betrachte die Funktion: $ f(x)=\ln(2x+1) $ und wende hier den MWS an und versuche durch Umformung und Abschätzung der Gleichung die Ungleichung zu beweisen.\\\\
\begin{align*}
f(x)&=\ln(2x+1)\\
f'(x)&=\frac{2}{2x+1}\\
\end{align*}
Wählt man $ a=0 $ und $ x>0 $, dann existiert nach MWS ein $ x_0 $ zwischen $ a $ und $ x $ mit:\\
\begin{align*}
&f'(x_0)=\frac{f(x)-f(a)}{x-a}\\
&\Leftrightarrow\\
&\big(\ln(2x+1)\big)'_{(x_0)}=\frac{\ln(2x+1)-\ln(1)}{x}\\
&\Leftrightarrow\\
&\frac{2}{2x_0+1}=\frac{\ln(2x+1)}{x}\\
\end{align*}
Nun gilt, für $ x_0>0 $ und durch Abschätzung von $ 2x_0+1 $:\\
\begin{align*}
&2x_0+1>1\;\;\text{für }\;x_0>0\\
&\Rightarrow\\
&\frac{2}{2x_0+1}<2\\
\end{align*}
Dann gilt für $ x>0 $:\\
\begin{align*}
&\frac{\ln(2x+1)}{x}=\frac{2}{2x_0+1}<2\\
&\Rightarrow\\
&\frac{\ln(2x+1)}{x}<2\;\;|\text{\tiny $\cdot x$ }\\
&\ln(2x+1)<2x
\end{align*}
Für $ x=0 $ gilt:\\
\begin{align*}
\ln(2x+1)&\leq 2x\\
\ln(2\cdot0+1)&\leq 2\cdot 0\\
0&=0
\end{align*}
Nun gilt für $ x\geq0 $:\\
\begin{align*}
\ln(2x+1)\leq 2x_{_{\Square}}\\
\end{align*}
\subsubsection{Aufgabe 2 zum Mittelwertsatz}
\href{https://youtu.be/uzNiQL4wJ-M}{Video 26}\\
Zeige mit dem MWS, dass für alle $ x\geq 0 $ gilt $ \sqrt{1+2x}\leq 1+x $\\
Für $ x=0 $ gilt:\\
\begin{align*}
\sqrt{1+2x}&\leq 1+x\\
\sqrt{1+2\cdot 0}&\leq 1+0\\
1&=1\\
\end{align*}
Betrachte nun nur die linke Seite der Ungleichung und wende dann den MWS an.\\
\begin{align*}
f(x)&=\sqrt{1+2x}\\
f'(x)&=\frac{1}{2\sqrt{1+2x}}\cdot \big(1+2x\big)'=\frac{1}{\sqrt{1+2x}}\\
\end{align*}
Wähle $ a=0 $ und $ x>0 $.\\
Dann existiert ein $ x_0 $ zwischen $ a $ und $ x $ mit:\\
\begin{align*}
f'(x_0)=\frac{f(x)-f(a)}{x-a}\\
\end{align*}
Dann gilt mit $ x_0 $ und $ a=0 $ :\\
\begin{align*}
\frac{1}{\sqrt{1+2x_0}}=\frac{\sqrt{1+2x}-\sqrt{1-0}}{x-0}=\frac{\sqrt{1+2x}-1}{x}\\
\end{align*}
Nun wird $ \frac{1}{\sqrt{1+2x_0}} $ abgeschätzt.\\
Es gilt für $ x_0>0 $\\
\begin{align*}
&\sqrt{1+2x_0}>1\\
&\Leftrightarrow\\
&\frac{1}{\sqrt{1+2x_0}}<1\\
\end{align*}
Dann gilt für $ x>0 $:\\
\begin{align*}
1&>\frac{1}{\sqrt{1+2x_0}}=\frac{\sqrt{1+2x}-1}{x}\\
&\Rightarrow\\
1&>\frac{\sqrt{1+2x}-1}{x}\\
1&>\frac{\sqrt{1+2x}}{x}-\frac{1}{x}\;\;|\text{\tiny $ +\frac{1}{x} $ }\\
1+\frac{1}{x}&>\frac{\sqrt{1+2x}}{x}-\frac{1}{x}+\frac{1}{x}\;\;|\text{\tiny $ \cdot x $ }\\
1+x&>\sqrt{1+2x}\\
\end{align*}
Nun gilt für $ x\geq  0 $ :\\
\begin{align*}
\sqrt{1+2x}\leq 1+x_{_{\Square}}
\end{align*}
\subsection{Folgerung des Mittelwertsatzes}
\subsubsection{Monotonoie von Folgen}
\href{https://youtu.be/M2WZh7_cc0s}{Video 27}\\
Eine Funktion $ f:\mathbb{R}\supseteq \mathbb{D}_f\longmapsto \mathbb{W}_f\subseteq\mathbb{R} $ heißt monoton steigend, wenn für alle $ a,b\in \mathbb{D}_f $ mit $ a<b $ folgt $ f(a)\leq f(b) $.\\
Analog definieren wir monoton fallend, wenn für alle $ a,b\in \mathbb{D}_f $ mit $ a<b $ folgt $ f(a)\geq f(b) $.\\
\begin{center}
	
	
	\begin{tikzpicture}
	
	% Gitternetzlinien
	\draw[xstep=1cm, ystep=1cm, very thin, color=lightgray] (-4,4) grid (4,-2);
	% x-Achse
	\draw[->] (-4,0) -- (4,0) node[below,xshift=-1em] {$x$}; 
	%Zahlen auf x-Achse
	\foreach \x in {-3,...,-1}
	\draw[shift={(\x,0)},color=black] (0pt,2pt) -- (0pt,-2pt) node[below] {\footnotesize $\x$};
	\foreach \x in {1,...,3}
	\draw[shift={(\x,0)},color=black] (0pt,2pt) -- (0pt,-2pt) node[below] {\footnotesize $\x$};
	% y-Achse 
	\draw[->] (0,-2) -- (0,4) node[left,yshift=-1em] {$f(x)$};%node[above left]
	%Zahlen auf y-Achse
	\foreach \y in {-1,...,-1}
	\draw[shift={(0,\y)},color=black] (2pt,0pt) -- (-2pt,0pt) node[left] {\footnotesize $\y$};
	\foreach \y in {1,...,3}
	\draw[shift={(0,\y)},color=black] (2pt,0pt) -- (-2pt,0pt) node[left] {\footnotesize $\y$};
	%Ursprung
	\draw[color=black] (0pt,-10pt) node[right] {\footnotesize $0$};
	
	\draw[thick,color=black] plot[samples=200, domain=-2.2:3.3] 
	(\x,{-(\x-0.2)^3/9+2.0}) node[below, xshift=-2.6em,yshift=0.6em] {\tiny monoton fallend};
	
	\draw[thick,color=blue] plot[samples=200, domain=-1.6:3.3] 
	(\x,{(\x-0.2)^3/8-1}) node[below, xshift=-0.7em,yshift=1.5em]{\footnotesize \tiny monoton steigend};
	
	
	
	\end{tikzpicture}\\
\end{center}
Bei streng monoton gilt:\\
\begin{align*}
f(a)<f(b)\\
f(a)>f(b)\\
\end{align*}
\begin{center}
	
	\begin{tikzpicture}
	
	% Gitternetzlinien
	\draw[xstep=1cm, ystep=1cm, very thin, color=lightgray] (-4,4) grid (4,-2);
	% x-Achse
	\draw[->] (-4,0) -- (4,0) node[below,xshift=-1em] {$x$}; 
	%Zahlen auf x-Achse
	\foreach \x in {-3,...,-1}
	\draw[shift={(\x,0)},color=black] (0pt,2pt) -- (0pt,-2pt) node[below] {\footnotesize $\x$};
	\foreach \x in {1,...,3}
	\draw[shift={(\x,0)},color=black] (0pt,2pt) -- (0pt,-2pt) node[below] {\footnotesize $\x$};
	% y-Achse 
	\draw[->] (0,-2) -- (0,4) node[left,yshift=-1em] {$f(x)$};%node[above left]
	%Zahlen auf y-Achse
	\foreach \y in {-1,...,-1}
	\draw[shift={(0,\y)},color=black] (2pt,0pt) -- (-2pt,0pt) node[left] {\footnotesize $\y$};
	\foreach \y in {1,...,3}
	\draw[shift={(0,\y)},color=black] (2pt,0pt) -- (-2pt,0pt) node[left] {\footnotesize $\y$};
	%Ursprung
	\draw[color=black] (0pt,-10pt) node[right] {\footnotesize $0$};
	
	\draw[thick,color=black] plot[samples=200, domain=-3:0.9] 
	(\x,{-(\x-5.2)^2/9+6.0}) node[below, xshift=2.6em,yshift=-0.6em] {\tiny streng monoton steigend};
	
	\draw[thick,color=blue] plot[samples=200, domain=-1.6:3.3] 
	(\x,{(\x-5.2)^2/8-2}) node[below, xshift=-1.8em,yshift=0.2em]{\footnotesize \tiny streng monoton fallend};
	
	
	
	\end{tikzpicture}\\
\end{center}
Beispiel:\\
Zeige, dass $ f:\mathbb{R}\longmapsto \mathbb{R} $ mit $ f(x)=2x $ streng monoton steigend ist.\\
Beweis:\\
Seien $ a,b\in \mathbb{R} $ und $ a<b $\\
Zu zeigen:\\
\begin{align*}
&f(a)<f(b)\\
&a<b\Rightarrow f(a)=2a<2b=f(b)_{_{\Square}}\\
\end{align*}
\paragraph{Monotonie von Funktionen Aufgabe 1}$  $\vspace*{3mm}\\
\href{https://youtu.be/BN9OWtPFY7g}{Video 28}\\
Zeige, dass die Funktion $ f:\;]0,\infty[ \;\longmapsto \mathbb{R}$ mit $ f(x) =\frac{1}{x}$ streng monoton fallend ist.\vspace*{4mm}\\
Für $ a,b\in \left] 0,\infty\right]  $ beliebig und mit $ a<b $ ist z.z. $ f(a)>f(b) $.\\
Beweis:\\
\begin{align*}
&a<b\\
&\Leftrightarrow\\
&\frac{1}{a}>\frac{1}{b}\\
&\Leftrightarrow\\
&f(b)=\frac{1}{b}<\frac{1}{a}=f(a)_{_{\Square}}\\
\end{align*}
Berachten wir $ \overset{\AC}{f}:\mathbb{R}\backslash \{0\}\longmapsto \mathbb{R}  $ mit $ f(x)=\frac{1}{x} $.
Diese Funktion ist nicht monoton!\\
Beweis:\\
Für $ a=-1 $ und $ b=1 $ gilt:\\
\begin{align*}
&a=-1<b=1\\
&\Leftrightarrow\\
&f(a)=-1<f(b)=1\Rightarrow  \text{monoton steigend }\\
\end{align*}
\begin{center}
	\begin{tikzpicture}
	
	% Gitternetzlinien
	\draw[xstep=1cm, ystep=1cm, very thin, color=lightgray] (-4,4) grid (4,-4);
	% x-Achse
	\draw[->] (-4,0) -- (4,0) node[below,xshift=-1em] {$x$}; 
	%Zahlen auf x-Achse
	\foreach \x in {-3,...,-1}
	\draw[shift={(\x,0)},color=black] (0pt,2pt) -- (0pt,-2pt) node[below] {\footnotesize $\x$};
	\foreach \x in {1,...,3}
	\draw[shift={(\x,0)},color=black] (0pt,2pt) -- (0pt,-2pt) node[below] {\footnotesize $\x$};
	% y-Achse 
	\draw[->] (0,-4) -- (0,4) node[left,yshift=-1em] {$f(x)$};%node[above left]
	%Zahlen auf y-Achse
	\foreach \y in {-1,...,-3}
	\draw[shift={(0,\y)},color=black] (2pt,0pt) -- (-2pt,0pt) node[left] {\footnotesize $\y$};
	\foreach \y in {1,...,3}
	\draw[shift={(0,\y)},color=black] (2pt,0pt) -- (-2pt,0pt) node[left] {\footnotesize $\y$};
	%Ursprung
	\draw[color=black] (0pt,-10pt) node[right] {\footnotesize $0$};
	
	\draw[thick,color=blue] plot[samples=200, domain=-4:-0.25] 
	(\x,{1/\x});% node[below, xshift=-5em] {$h(x)$};
	
	\draw[thick,color=blue] plot[samples=200, domain=0.25:4] 
	(\x,{1/\x});% node[below, xshift=1em]{$g(x)$};
	
	\draw[color=blue] (-2.5,-1.2) node {\tiny streng monoton fallend};
	\draw[color=blue] (2.3,1.2) node {\tiny streng monoton fallend};
	\end{tikzpicture}\\	
\end{center}
\paragraph{Monotonie von Funktionen Aufgabe 2}$  $\vspace*{3mm}\\
\href{https://youtu.be/EiJkuy1dVWc}{Video 29}\\
Zeige, dass die Funktion $ f:[0,\infty [\;  \longmapsto \;[0,\infty [  $ mit $ f(x)=\sqrt{x} $ streng monoton steigend ist.\vspace*{4mm}\\
Für $ a,b\in [0,\infty[$ beliebig, mit $ a<b $ ist z.z. $ f(a)<f(b) $\\
Beweis:
\begin{align*}
a<b\Leftrightarrow \frac{a}{b}<1\Leftrightarrow \sqrt{\frac{a}{b}}<1\Leftrightarrow \frac{\sqrt{a}}{\sqrt{b}}<1\Leftrightarrow \sqrt{a}<\sqrt{b}_{_{\Square}}\\
\end{align*}
Bemerkung (b kann nicht Null werden): \\
Da $a,b\in [0,\infty[ $ und $ a<b   \Rightarrow a=0\Rightarrow b>0$\\
\\
Also ist $ f(x)=\sqrt{x} $ streng monoton steigend.\\
\begin{center}
	\begin{tikzpicture}
	
	% Gitternetzlinien
	\draw[xstep=1cm, ystep=1cm, very thin, color=lightgray] (-4,4) grid (4,-2);
	% x-Achse
	\draw[->] (-4,0) -- (4,0) node[below,xshift=-1em] {$x$}; 
	%Zahlen auf x-Achse
	\foreach \x in {-3,...,-1}
	\draw[shift={(\x,0)},color=black] (0pt,2pt) -- (0pt,-2pt) node[below] {\footnotesize $\x$};
	\foreach \x in {1,...,3}
	\draw[shift={(\x,0)},color=black] (0pt,2pt) -- (0pt,-2pt) node[below] {\footnotesize $\x$};
	% y-Achse 
	\draw[->] (0,-2) -- (0,4) node[left,yshift=-1em] {$f(x)$};%node[above left]
	%Zahlen auf y-Achse
	\foreach \y in {-1,...,-1}
	\draw[shift={(0,\y)},color=black] (2pt,0pt) -- (-2pt,0pt) node[left] {\footnotesize $\y$};
	\foreach \y in {1,...,3}
	\draw[shift={(0,\y)},color=black] (2pt,0pt) -- (-2pt,0pt) node[left] {\footnotesize $\y$};
	%Ursprung
	\draw[color=black] (0pt,-10pt) node[right] {\footnotesize $0$};
	
	
	
	\draw[thick,color=blue] plot[samples=200, domain=0.001:4] 
	(\x,{(\x)^0.5}) node[below, xshift=-2em,yshift=-1em]{$f(x)$};
	
	
	
	\end{tikzpicture}\\
\end{center}
\subsubsection{Monotoniekriterium}
\href{https://youtu.be/K_BfNKGb46I}{Video 30}\\
Sei $ f: \mathbb{I}\longmapsto \mathbb{R} $ differenzierbar auf dem Intervall $ \mathbb{I} $.\\
Für alle $ x\in \mathbb{I} $ gilt:\\
\begin{center}
	Monotoniekriterium\\
	\fbox{\parbox  {100mm} {
			1. $ f'(x)\geq 0 $, dann ist $ f $ auf $ \mathbb{I} $ monoton steigend.\\
			2. $ f'(x)> 0$, dann ist $ f $ auf $ \mathbb{I} $ streng monoton steigend.\\
			3. $ f'(x)\leq 0$, dann ist $ f $ auf $ \mathbb{I} $ monoton fallend.\\
			4. $ f'(x)< 0$, dann ist $ f $ auf $ \mathbb{I} $ streng monoton fallend.\\}}
\end{center}
Dies sind Folgerungen aus dem Mittelwertsatz.\\
Beweis zu 1.)\\
Sei $ f'(x)\geq 0 $ für alle $ x\in \mathbb{I} $.\\
Zz.: $ a,b\in \mathbb{I} $ und $ a<b\Rightarrow f(a)\leq f(b) $\\
Nach dem MWS existiert ein $ x_0 $ zwischen $ a $ und $ b $ mit $ f'(x)=\frac{f(b)-f(a)}{b-a} $\\
Mit der Vorraussetzung $ f'(x)\geq 0 $ gilt:\\
\begin{align*}
0&\leq f'(x)=\frac{f(b)-f(a)}{b-a} \\
0&\leq\frac{f(b)-f(a)}{b-a} \;\;|\text{\tiny $ \cdot (b-a) $ }\\
&\leq\frac{f(b)-f(a)(b-a)}{(b-a)} \\
&\leq f(b)-f(a)\\
&\Leftrightarrow f(a)\leq f(b)_{_{\Square}}
\end{align*}
Beispiel:\\
\textbullet\;\; $ f:[0,\infty[\;\longmapsto \;[0,\infty[ $ mit $ f(x)=x^2 $ ist monoton wachsend,weil gilt:\\
\begin{align*}
f'(x)=2x\geq 0 \;\;\text{für }\;\;x\in [0,\infty[\\
\end{align*}
\textbullet\;\; $ g:\;]0,\infty[\;\longmapsto \;]0,\infty[ $ mit $ g(x)=x^2 $ ist streng monoton wachsend,weil gilt:\\
\begin{align*}
g'(x)=2x > 0 \;\;\text{für }\;\;x\in ]0,\infty[\\
\end{align*}
\paragraph{Aufgabe zum Monotoniekriterium}$   $\vspace*{3mm}\\
\href{https://youtu.be/E2T-2jXem2Y}{Video 31}\\
Sei $ f:\mathbb{R}\longmapsto \mathbb{R} $ mit $ f(x)=2x^3+3x^2-12x+1 $ gegeben\\
Bestimme die Bereiche auf denen $ f $ streng monoton ist.\vspace*{2mm}\\
Betrachte:\\
\begin{align*}
f'(x)=(2x^3+3x^2-12x+1)'=6x^2+6x-12=6(x^2+x-2)\\
\end{align*}
Die Nullstellen von $ f'(x) =6(x^2+x-2)$ sind $ -2 $ und $ 1 $ \\
\begin{align*}
f'(x)=6(x-1)(x+2)\\
\end{align*}
Bei den Nullstellen hat man ein Vorzeichenwechsel und dies zeigt den Wechsel zwischen  monoton steigend und fallend. Ist $ f $ streng monoton, dann darf $ f'(x) $ nicht Null werden.\vspace*{2mm}\\
Bei $ x=-3 $ ist $ f'(x)=6(-3-1)(-3+2)\Rightarrow $ positiv $ f'(x)>0 $\\
Bei $ x=0 $ ist $ f'(x)=6(0-1)(0+2)\Rightarrow $ negativ $ f'(x)<0 $\\
Bei $ x=3 $ ist $ f'(x)=6(3-1)(3+2)\Rightarrow $ positiv $ f'(x)>0 $\\
\begin{center}
	
	\begin{tikzpicture}
	
	
	% x-Achse
	\draw[->] (-4,0) -- (4,0) node[below,xshift=-1em] {$x$}; 
	%Zahlen auf x-Achse
	\foreach \x in {-3,...,3}
	\draw[shift={(\x,0)},color=black] (0pt,2pt) -- (0pt,-2pt) node[below] {\footnotesize $\x$};
	\draw [dashed,color=blue] (-2,1)--(-2,-1);
	\draw [dashed,color=blue] (1,1)--(1,-1);
	\draw [color=blue] (-3.1,0.3) node {$ f'(x)>0 $};
	\draw [color=blue] (2.2,0.3) node {$ f'(x)>0 $};
	\draw [color=blue] (-0.6,0.3) node {$ f'(x)<0 $};
	\end{tikzpicture}\\
\end{center}
Fazit:\\
(i)\;\,\,Für $ x\in \;]-\infty,-2[ $ ist $ f'(x) >0$\\
(ii) Für $ x\in \;]-2,1[ $ ist $ f'(x) <0$\\
(iii)\,Für $ x\in \;]1,\infty[ $ ist $ f'(x) >0$\vspace*{4mm}\\
$f(x) $ ist streng monoton steigend im Intervall:\\
\begin{align*}
x\in\;]-\infty,-2[\;\;\cup\;\;]1,\infty[\\
\end{align*}
$f(x) $ ist streng monoton fallend im Intervall:\\
\begin{align*}
x\in\;]-2,1[\\
\end{align*}
\subsubsection{Konstanzkriterium}
\href{https://youtu.be/SUuhqlbPbGI}{Video 32}\\
Sei $f:\mathbb{R}\supset \mathbb{I}\longmapsto \mathbb{R}$ auf $ \mathbb{I} $ diffbar.\\
Gelte $f'(x)=0$ für $x\in\mathbb{I}  $ dann ist $ f $ konstant, also von der Form $ f'(x)=0 $ für $ x\in \mathbb{I}, k\in\mathbb{R} $.\vspace*{2mm}\\
Nach dem MWS für alle $ a,b\in \mathbb{I} $ mit $ a\neq b $ existiert ein $ x_0 $ zwischen $ a $ und $ b $ mit $ f'(x)=\frac{f(b)-f(a)}{b-a} $.\vspace*{4mm}\\
Für $f$ konstant gilt dann:\\
\begin{align*}
&0=f'(x)=\frac{f(b)-f(a)}{b-a}\;\;|\text{\tiny $\cdot (b-a)$ } \\
&\Leftrightarrow f(b)-f(a)=0\\
&\Leftrightarrow f(b)=f(a)\\
\end{align*}
D.h. $ f(x)=k $ für alle $ x\in\mathbb{I} $ mit $ k\in \mathbb{R} $\vspace*{4mm}\\
Beispiel:\\
\textbullet\;\;Ist $ f $ konstant, dann gilt $ f'(x)=0 $ denn:\\
\begin{align*}
(f(x))'=(k)'=(k\cdot x^0)'=0\cdot k\cdot x^{-1}=0\\
\end{align*}
\textbullet\;\;$ f(x)=10\Rightarrow f'(x)=0 $ für alle $ x\in\mathbb{R} $\\
\begin{quote}
	Rückschluß nicht möglich, wenn $ f'(x)=10 $ erhält man nicht $ f(x) =10$!\\
\end{quote}
\paragraph{Aufgabe zum Konstanzkriterium}$  $\vspace*{3mm}\\
\href{https://youtu.be/mMXOmLVoRmo}{Video 33}\\
Bestimme alle diffbaren Funktionen $ f $ mit folgenden Eigenschaften:\\
(i)\; $ (f-p)'_{(x)}=0 $ für alle $ x\in \mathbb{R} $\\
(ii) $f(1)=p(1)$,
wobei $ p(x)=x^2+x-2 $ ist.\\
zu(i):\\
\begin{quote}
	Mit dem Konstanzkriterium folgt, dass $ f-p $ konstant ist:\\
	Es gilt:\\
	\begin{align*}
	&f(x)-p(x)=k\\
	&\Leftrightarrow f(x)=p(x)+k=x^2+x-2+k\\
	\end{align*}
\end{quote}
zu (ii):\\
\begin{quote}
	\begin{align*}
	f(1)=1^2+1-2+k=k=p(1)=0\Rightarrow k=0\\
	\end{align*}
\end{quote}
Dann sind alle diffbaren Funktionen mit den Eigenschften (i) und (ii) von der Form:\\
\begin{align*}
f(x)=x^2+x-2=p(x)\\
\end{align*}
\subsection{Regel von L’Hospital}
\href{https://youtu.be/wBqlghVLs04}{Video 34}\\
Seien $ f,g: \mathbb{I} \longmapsto \mathbb{R}$ diffbare Funktionen auf dem Intervall $\mathbb{I}\subset \mathbb{R}$ sowie $ x_0\in \mathbb{R}\cup \{\pm\infty\} $.\\
Es gelte weiterhin $ g'\neq 0 $ für alle $ x\in\mathbb{I}$ und es existiert der Limes\\
\begin{align*}
\lim\limits_ {x\to x_0} \frac{f'(x)}{g'(x)}=G\;\;\text{mit }\;\;  G\in \mathbb{R}\\
\end{align*}
dann folgt falls:\\
\textbullet \;\; $ \lim\limits_ {x\to x_0}f(x)= \lim\limits_ {x\to x_0}g(x)= 0 $\\
\begin{align*}
\lim\limits_ {x\to x_0} \frac{f'(x)}{g'(x)}=G
\end{align*}
\textbullet \;\; $ \lim\limits_ {x\to x_0}f(x)= \lim\limits_ {x\to x_0}g(x)= \pm\infty $\\
\begin{align*}
\lim\limits_ {x\to x_0} \frac{f'(x)}{g'(x)}=G
\end{align*}
Beispiel:\\
\begin{align*}
\lim\limits_ {x\to 0}\frac{\sin(x)}{x}\overset{L'H}{=}\lim\limits_ {x\to\ 0}\frac{(\sin(x))'}{(x)'}=\lim\limits_ {x\to 0} \frac{\cos(x)}{1}=\lim\limits_ {x\to 0} \cos(x)=1\\
\end{align*}
\subsubsection{Trick 1 zur Regel von L’Hospital}
\href{https://youtu.be/Chm62OX6xbE}{Video 35 }\\
Berechne falls existieren die folgenden Genzwerte.\\
(i)\;\;$ \lim\limits_ {x\to 0}\frac{\cos(x)-1}{x\cdot \sin(x)} $\\
(ii) $ \lim\limits_ {x\to 0}\frac{\ln(x+1)}{\ln(x-1)} $\vspace*{2mm}\\
Zu (i):\\
\begin{align*}
\lim\limits_ {x\to 0}\frac{\cos(x)-1}{x\cdot \sin(x)} \\
\end{align*}
Ist die Regel anwendbar?\\
\begin{align*}
\cos(0)-1&=0\\
0\cdot \sin(0)&=0\\
&\Rightarrow \frac{0}{0}\\
\end{align*}
\begin{align*}
\lim\limits_ {x\to 0}\frac{\cos(x)-1}{x\cdot \sin(x)}\overset{L'H}{=}\lim\limits_ {x\to 0}\frac{(\cos(x)-1)'}{(x\cdot \sin(x))'}=\lim\limits_ {x\to 0}\bigg[\frac{-\sin(x)}{\sin(x)+x\cdot \cos(x)}\bigg]
\end{align*}
\begin{align*}
-\sin(0)&=0\\
\sin(x)+x\cdot \cos(x)&=0\\
&\Rightarrow \frac{0}{0}\\
\end{align*}
Die Regel L'Hospital nochmals anwenden.\\
\begin{align*}
\lim\limits_ {x\to 0}\bigg[\frac{-\sin(x)}{\sin(x)+x\cdot \cos(x)}\bigg]\overset{L'H}{=}\lim\limits_ {x\to 0}\bigg[\frac{-\cos(x)}{\cos(x)+\cos(x)+x\cdot (-\sin(x))}\bigg]=\lim\limits_ {x\to 0}\bigg[\frac{-\cos(x)}{2\cdot \cos(x)-x\cdot \sin(x)}\bigg]=-\frac{1}{2}\\
\end{align*}
Der Grenzwert von:\\
\begin{align*}
\lim\limits_ {x\to 0}\frac{\cos(x)-1}{x\cdot \sin(x)}=-\frac{1}{2} \\
\end{align*}
Zu (ii):\\
\begin{align*}
\lim\limits_ {x\to \infty}\frac{\ln(x+1)}{\ln(x-1)}\overset{L'H}{=}\lim\limits_ {x\to \infty}\frac{(\ln(x+1))'}{(\ln(x-1))'}&=\lim\limits_ {x\to \infty}\frac{\frac{1}{x+1}}{\frac{1}{x-1}}=\lim\limits_ {nx\to 0}\bigg[\frac{(x-1)}{(x+1)}\bigg]=\lim\limits_ {x\to \infty}\bigg[\frac{x(1-\frac{1}{x})}{x(1+\frac{1}{x})}\bigg]=1\\
\end{align*}
\subsubsection{Trick 2 zur Regel von L’Hospital}
\href{https://youtu.be/MozzXCc8xNM}{Video 36}\\
Berechne falls existieren die folgenden Grenzwerte\\
(i)\;\;$ \lim\limits_ {x\to 0}=\frac{\cos(\sin(x))-1}{\sin(x)} $\\
(ii) $ \lim\limits_ {x\to 0}\frac{x}{\mathrm{e}^x-1} $\vspace*{2mm}\\
Zu (i):\\
\begin{align*}
\lim\limits_ {x\to 0}\bigg[\frac{\cos(\sin(x))-1}{\sin(x)}\bigg]&\overset{L'H}{=}\lim\limits_ {x\to 0}\bigg[\frac{\cos'(\sin(x))\cdot \sin'(x)}{\cos(x)}\bigg]=\lim\limits_ {x\to 0}\bigg[\frac{-\sin(\sin(x))\cdot \cos(x)}{\cos(x)}\bigg]=\lim\limits_ {x\to 0}[-\sin(\sin(x))]=0\\
\end{align*}
Zu (ii):\
\begin{align*}
\lim\limits_ {x\to 0}\frac{x}{\mathrm{e}^x-1}\overset{L'H}{=}\lim\limits_ {x\to 0}\frac{1}{\mathrm{e}^x}=1\\
\end{align*}
\subsubsection{Trick 3 zur Regel von L’Hospital}
\href{https://youtu.be/ODertm5pPaU}{Video 37 }\\
Berechne für $ m\in \mathbb{N} $ und $ n\in \mathbb{N}\backslash\{0\} $ den Grenzwert von: $\lim\limits_ {x\to\infty}\frac{\ln(x^m+1)}{\ln(x^n)}$\\
\begin{align*}
\lim\limits_ {x\to\infty}\frac{\ln(x^m+1)}{\ln(x^n)}&\overset{L'H}{=}\lim\limits_ {x\to\infty}\frac{(\ln(x^m+1))'}{(\ln(x^n))'}=\lim\limits_ {n\to\infty}\bigg[\frac{\genfrac{}{}{}{2}{mx^{m-1} }{x^m+1 }}{\genfrac{}{}{}{2}{nx^{n-1}}{x^{n} }}\bigg]\\
&=\lim\limits_ {n\to\infty}\bigg[\frac{x^n\cdot mx^{m-1}}{(x^{m}+1)\cdot nx^{n-1}}\bigg]\\
&=\lim\limits_ {n\to\infty}\bigg[\frac{x^n\cdot mx^m\cdot x^{-1}}{(x^m +1)\cdot nx^nx^{-1}}\bigg]=\frac{m}{n}\lim\limits_ {n\to\infty}\bigg[\frac{x^m}{x^{m}+1}\bigg]=\frac{m}{n}\lim\limits_ {n\to\infty}\bigg[\frac{x^m}{x^m(1+\frac{1}{x^m})}\bigg]\\
&=\frac{m}{n}\lim\limits_ {n\to\infty}\bigg[\frac{1}{1+\frac{1}{x^m}}\bigg]=\frac{m}{n}\cdot 1=\frac{m}{n}\\
\end{align*}
\subsubsection{Trick 4 zur Regel von L’Hospital}
\href{https://youtu.be/moiaqTqBnkI}{Video 38 }\\
Berechne $ \lim\limits_ {n\to\infty}\frac{x^2}{\mathrm{e}^x} $ und zeige, dass $ \lim\limits_ {n\to\infty}\frac{x^n}{\mathrm{e}^x} $ für jedes $ n\in \mathbb{N} $ gleich ist:\vspace*{4mm}\\
\textbullet \;\;$\lim\limits_ {n\to\infty}\frac{x^2}{\mathrm{e}^x}\overset{L'H}{=}\lim\limits_ {n\to\infty}\frac{(x^2)'}{(\mathrm{e}^x)'}=\lim\limits_ {n\to\infty}\frac{2x}{\mathrm{e}^x}\overset{L'H}{=}\lim\limits_ {n\to\infty}\frac{2}{\mathrm{e}^x}=0$\vspace*{4mm}\\
\textbullet\;\;$\lim\limits_ {n\to\infty}\frac{x^n}{\mathrm{e}^x}\overset{L'H}{=}\lim\limits_ {n\to\infty}\frac{nx^{n-1}}{\mathrm{e}^x}\overset{L'H}{=}\lim\limits_ {n\to\infty}\frac{n(n-1)x^{n-2}}{\mathrm{e}^x}\overset{L'H}{=}\dots\overset{L'H}{=}\lim\limits_ {n\to\infty}\frac{n(n-1(n-2)\dots 1\cdot x^{n-n}}{\mathrm{e}^x}=\lim\limits_ {n\to\infty}\frac{n!}{\mathrm{e}^x}=0$\\
\begin{quote}
	Bemerkung: n! ist eine Konstante.\\
	Fazit: Die Exponentialfunktion wächst schneller wie $ x^n $.\\
	\begin{align*}
	\lim\limits_ {n\to\infty}\frac{x^n}{\mathrm{e}^x}=0\;\;\text{für alle }\;\; n\in \mathbb{N} 
	\end{align*}
\end{quote}


%	\chapter{Höhere Ableitungen}

\setcounter{section}{9}
\setcounter{subsection}{0}
\setcounter{subsubsection}{1}
\setcounter{secnumdepth}{4}
\subsection{Mehrfaches Ableiten}
\href{https://youtu.be/CA2VJcrvSVo}{Video 01}\\
Sei $ f:\mathbb{R}\supset\mathbb{I}\longmapsto\mathbb{R} $ auf $ \mathbb{I} $ diffbar, dann ist $ f':\mathbb{I}\longmapsto\mathbb{R} $ mit $ x\longmapsto f'(x) $ eine Funktion.\\
Ist $ f' $ auf $ \mathbb{I} $ diffbar, dann heißt\\
\begin{align*}
f'':=(f')'\\
\end{align*}
die zweite Ableitung von $ f $.\vspace*{4mm}\\
Beispiel:\\
\begin{align*}
f(x)&=x^2\\
f'(x)&=2x\\
f''(x)&=(f'(x))'=(2x)'=2\\
\end{align*}
Für $ k\in \mathbb{N} $ definieren wir die $ k- $te Ableitung:\\
\begin{align*}
f^{(k)}:=(f^{k-1})'\;\;\text{, wobei } f^{(0)}:=f\;\;\text{ist. }\\
\end{align*}
Beispiel:\\
\begin{align*}
f(x)&=x^3\\
f'(x)&=3x\\
f''(x)&=3\\
f^{(3)}(x)&=0\\
f^{(123)}(x)&=0\\
\end{align*}
Hier gibt es unendlich viele Ableitungen.\\
\newpage
\noindent
Bemerkung:\\
Beispiel:\\
\textbullet \;\;$ f(x)=|x| $ \\
\begin{quote}
	$ f(x)=|x| $ ist in $ x_0=0 $ nicht ableitbar, weil in $ x_0 $ nicht diffbar!\\
\end{quote}
\textbullet \;\;$ g(x) =     \left\{ \begin{array}{rcl}
\frac{1}{2 }x^2 & \mbox{für} & x\geq0\\
-\frac{1}{2}x^2 & \mbox{für} & x< 0 \\

\end{array}\right.$ \\
\begin{quote}
Überprüfung der Differenzierbarkeit in $ x_0=0 $.\\
Grenzwertbetrachtung:\\
\begin{align*}
\lim\limits_ {\substack{x \to x_0 \\ x<0}}\frac{g(x)-g(x_0)}{x-x_0}=\lim\limits_ {x\to x_0}\frac{-\frac{1}{2}x^2}{x}=\lim\limits_ {n\to x_0}-\frac{1}{2}x=0=\lim\limits_ {x\to x_0}\frac{1}{2}x=\lim\limits_ {x\to x_0}\frac{\frac{1}{2}x^2}{x}=\lim\limits_ {\substack{x \to x_0 \\ x> 0}}\frac{g(x)-g(x_0)}{x-x_0}
\end{align*}
Der links- und rechtsseitige Grenzwert ist gleich und deshalb ist die Funktion
$ g(x) $ in $ x_0=0 $ differenzierbar. Bei $ x_0\neq 0 $ ist $ g(x) $ als Komposition einer differenzierbaren Funktion differenzierbar.\\

\begin{align*}
&g'(x) =     \left\{ \begin{array}{rcl}
(\frac{1}{2}x^2)' & \mbox{für} & x\geq 0\\
(-\frac{1}{2}x^2)' & \mbox{für} & x<0 \\
\end{array}\right.\\
&g'(x) =     \left\{ \begin{array}{rcl}
x & \mbox{für} & x\geq 0\\
-x & \mbox{für} & x<0 \\
\end{array}\right.=|x|\\
\end{align*}
$ g'(x)=|x| $ ist nicht diffbar.\\
Fazit:\\
$ g(x) $ ist nur einmal differenzierbar.\\
\end{quote}
\subsubsection{Mehrfaches Ableiten - Aufgabe}
\href{https://youtu.be/1SbIqyzxLeg}{Video 02}\\
Sei $ f(x)=x^2+2x+7 $.\\
Berechne:\\
\begin{align*}
g(x)=f(1)+f'(1)(x-1)+\frac{f''(x)}{2}(x-2)^2
\end{align*}
und vergleiche $ g(x) $ mit $ f(x) $.\\
\begin{align*}
f(x)&=x^2+2x+7\\
f(1)&=1^2+2+7=10\\
f'(x)&=2x+2\\
f'(1)&=2+2=4\\
f''(x)&=2\\
f''(2)&=2\\
g(x)&=10+4(x-1)+(x-1)^2\\
g(x)&=10+4x-4+x^2-2x+1=x^2+2x+7=f(x)\\
\end{align*}
Die Frage ist nun: Gilt dies für alle Polynome?\\
\subsection{ Herleitung des Taylorpolynoms}
\href{https://youtu.be/EFkwp2qckNU}{Video 03}\\
Sei $ x_0\in\mathbb{R} $ und betrachte $ f(x)=a_0+a_1(x-x_0)+\ldots+a_k(x-x_0)^k\ldots+a_n(x-x_0)^n $.\vspace*{4mm}\\
Untersuchung von $ h(x)=a_k(x-x_0)^k $\\
\begin{align*}
h(x)&=a_k(x-x_0)^k\\
h'(x)&=a_k\cdot k(x-x_0)^{k-1}\\
h''(x)&=a_k\cdot k\cdot (k-1)(x-x_0)^{k-2}\\
&\;\;\cdot\\
&\;\;\cdot\\
&\;\;\cdot\\
h^{(k)}(x)&=a_k\cdot \underbrace {k(k-1)\ldots 3\cdot 2 \cdot 1}_{=k!}(x-x_0)^0=a_k k!\\
h^{k+1}(x)&=0\\
\end{align*}
Durch einsetzen von $ x=x_0 $ in $ f(x) =a_0+a_1(x-x_0)+\ldots+a_k(x-x_0)^k\ldots+a_n(x-x_0)^n$ kann man die Koeffizienten $ a_k $ berechnen.\\
Dann gilt für $ x=x_0 $ eingesetzt in die Funktion $ f(x) $:\\
\begin{align*}
f(x_0)&=a_0+a_1(x-x_0)+\cdots +a_k(x-x_o)^k\\
f(x_0)&=a_0
\end{align*}
Alle anderen Summanten werden Null $ (x_0-x_0)=0 $.\vspace*{2mm}\\
Dann gilt für $ x=x_0 $ eingesetzt in die Funktion $ f'(x) $:\\
\begin{align*}
f'(x_0)&=a_1+\cdots+a_k\cdot k(x-x_0)^{k-1}\\
f'(x_0)&=a_1
\end{align*}
Dann gilt für: $ f^k(x_0) $
\begin{align*}
f^k (x_0)=a_k\cdot k!\Leftrightarrow a_k=\frac{f^k(x_0)}{k!}\\
\end{align*}
Dann ist:\\
\begin{align*}
a_0=\frac{f(x_0)}{0!}\\
a_1=\frac{f'(x_0)}{1!}\\
a_2=\frac{f''(x_0)}{2!}\\
a_n=\frac{f^n(x_0)}{n!}\\
\end{align*}
Daraus ergibt sich dann für f(x):\\
\begin{align*}
f(x)&=\frac{f(x_0)}{0!}+\frac{f'(x_0)(x-x_0)^1}{1!}+\frac{f''(x_0)(x-x_0)^2}{2!}+\ldots +\frac{f^n(x_0)-(x-x_0)^n}{n!}\\
&=\sum\limits_{k=0}^{n}\frac{f^k(x_0)}{k!}(x-x_0)^k\\
\end{align*}
Nun stellen sich die Fragen:\\
Was ist, wenn $ f $ kein Polynom ist ?\\
Welcher Zusammenhang besteht zwischen $ f $ und dem Polynom ?\\
\subsubsection{Das Taylorpolynom (Definition) }
\href{https://youtu.be/bSzpFWy4gTE}{Video 04}\\
Sei $ f: \mathbb{I}\longmapsto\mathbb{R} $ n-mal diffbar und $ x_0\in  \mathbb{I} $, dann heißt
\begin{align*}
T_{n,x_0}(x)=\sum\limits_{k=0}^{n}\frac{f^k(x_0)}{k!}(x-x_0)^k
\end{align*}
das n-te Taylorpolynom von $ f $ im Entwicklungspunkt $ x_0 $.\vspace*{2mm}\\
Beispiel: $ f(x)=\ln(x) $ mit $ x_0=1 $ und $ n=2 $\\
\begin{align*}
&f^{(0)}(1)=\ln(1)=0\\
&f^{(1)}(x)=(\ln(x))'=\frac{1}{x}\Rightarrow f^{(1)}(1)=1\\
&f^{(2)}(x)=\bigg(\frac{1}{x}\bigg)'=\frac{-1}{x^2}\Rightarrow f^{(2)}(1)=-1\\
\end{align*}
Daraus ergibt sich das Taylorpolynom:\\
\begin{align*}
T_{2,x_0}(x)&=\frac{f^{(0)}(x_0)}{0!}+\frac{f^{(1)}(x_0)}{1!}(x-x_0)+\frac{f^{(2)}(x_0)}{2!}(x-x_0)^2\\
T_{2,1}(x)&=0+1(x-1)+\frac{-1}{2}(x-1)^2\\
&=(x-1)-\frac{1}{2}(x-1)^2\\
\end{align*}
\subsubsection{Aufgabe 1 zum Taylorpolynom}
\href{https://youtu.be/YGIDHu8E8S4}{Video 05}\\
Berechne für $ f(x)=\sqrt{x} $ das Taylorpolynom 3. Grades im Entwicklungspunkt $ x_0=1 $.\\
\begin{align*}
T_{3,1}(x)=\sum\limits_{k=0}^{3}\frac{f^k(1)}{k!}(x-1)^k
\end{align*}
\newpage
\noindent
\textbullet \;\;Bestimmung der Ableitungen von  $ f'(x),f''(x)\;\text{ und}\; f^{(3)}(x)$.\\
\begin{quote}
\begin{align*}
f'(x)&=\frac{1}{2\sqrt{x}}\\
f''(x)&=\bigg(\frac{1}{2\sqrt{x}}\bigg)'=\frac{1}{2}\bigg(x^{-\frac{1}{2}}\bigg)'=\frac{1}{2}\bigg(-\frac{1}{2}\bigg)x^{-\frac{3}{2}}=-\frac{1}{4}\frac{1}{\sqrt{x^3}}\\
f^{(3)}(x)&=\bigg(-\frac{1}{4}\frac{1}{\sqrt{x^3}}\bigg)'=-\frac{1}{4}\bigg(x^{-\frac{3}{2}}\bigg)'=-\frac{1}{4}\bigg(-\frac{3}{2}\bigg)x^{-\frac{5}{2}}=\frac{3}{8}\frac{1}{\sqrt{x^5}}\\
\end{align*}
\end{quote}
\textbullet\;\; Bestimmung von $ f(x),\; f'(1),\; f''(1),\; f^{(3)}(1) $.\\
\begin{quote}
\begin{align*}
f(1)&=1\\
f'(1)&=\frac{1}{2\sqrt{1}}=\frac{1}{2}\\
f''(1)&=-\frac{1}{4}\frac{1}{\sqrt{1^3}}=-\frac{1}{4}\\
f^{(3)}(1)&=\frac{3}{8}\frac{1}{\sqrt{x^5}}=\frac{3}{8}\\
\end{align*}
\end{quote}
\textbullet\;\; Einsetzen in das Taylorpolynom.\\
\begin{quote}
\begin{align*}
T_{3,x_0}(x)&=\sum\limits_{k=0}^{3}\frac{f^k(x_0)}{k!}(x-x_0)^k=f(x_0)+\frac{f'(x_0)}{1!}(x-x_0)+\frac{f''(x_0)}{2!}(x-x_0)^2+\frac{f^{(3)}(x_0)}{3!}(x-x_0)^3\\
T_{3,1}(x)&=\sum\limits_{k=0}^{3}\frac{f^k(1)}{k!}(x-1)^k=f(x_0)+\frac{f'(1)}{1!}(x-1)+\frac{f''(1)}{2!}(x-1)^2+\frac{f^{(3)}(1)}{3!}(x-1)^3\\
\end{align*}
Einsetzen von $ f(x)=1,\; f'(1)=\frac{1}{2},\; f''(1)=-\frac{1}{4},\; f^{(3)}(1)=\frac{3}{8} $.
\begin{align*}
T_{3,1}(x)=1+\frac{1}{2}(x-1)-\frac{1}{8}(x-1)^2+\frac{3}{48}(x-1)^3\\
\end{align*}
\end{quote}
\subsubsection{Aufgabe 2 zum Taylorpolynom}
\href{https://youtu.be/uHJAzrQ-4-I}{Video 06}\\
Sei $ f $ die Lösung der Differentialgleichung $ f'(x)=2f(x)+\frac{3}{f(x)} $ mit $ f(0)=1 $.\\
Bestimme Taylorpolynom von 2. Grades in $ x_0=0 $.\vspace*{2mm}\\
\newpage
\noindent
\textbullet \;\;Bestimmung der Ableitungen von  $ f'(x),\;f''(x)$.\\
\begin{quote}
\begin{align*}
f'(x)&=2f(x)+\frac{3}{f(x)}\\
f''(x)&=\bigg(2f(x)+\frac{3}{f(x)}\bigg)'=2f'(x)+\bigg(\frac{3}{f(x)}\bigg)'\\
\end{align*}
NR.: $\bigg(\frac{3}{f(x)}\bigg)'  $\\
Es gilt:\\
\begin{align*}
\bigg(\frac{1}{x}\bigg)'&=\bigg(x^{-1}\bigg)'=-x^{-2}=-\frac{1}{x^2}\\
\bigg(\frac{1}{f(x)}\bigg)'&=\bigg((f(x))^{-1}\bigg)'=-f'(x)(f(x))^{-2}=-\frac{f'(x)}{(f(x))^2}\\
\end{align*}
\begin{align*}
&\bigg(\frac{3}{f(x)}\bigg)'=3\bigg(\frac{1}{f(x)}\bigg)'=3\big(f(x)^{-1}\big )'=-3f'(x)\big ( f(x)\big)^{-2}=-3\frac{f'(x)}{(f(x))^2}\\
&\text{mit Quotientenregel }\\
&\bigg(\frac{3}{f(x)}\bigg)'=\frac{0\cdot f(x)-3\cdot f'(x)}{(f(x))^2}=-3\frac{f'(x)}{(f(x))^2}\\
\end{align*}
\begin{align*}
\Rightarrow\\
f''(x)&=\bigg(2f(x)+\frac{3}{f(x)}\bigg)'=\bigg(2f(x)\bigg)'+\bigg(\frac{3}{f(x)}\bigg)'\\
&=2f'(x)-3\frac{f'(x)}{(f(x))^2}\\
&=f'(x)\bigg (2-\frac{3}{(f(x))^2}\bigg)
\end{align*}
\end{quote}
\textbullet\;\; Bestimmung von $ f(0),\; f'(0),\; f''(0) $.\\
\begin{quote}
\begin{align*}
f(0)&=1\\
f'(0)&=2f(0)+\frac{3}{f(0)}=2+3=5\\
f''(0)&=f'(0)\bigg (2-\frac{3}{(f(0))^2}\bigg)=5\bigg (2-\frac{3}{1}\bigg)=-5\\
\end{align*}
\end{quote}
\textbullet\;\; Einsetzen in das Taylorpolynom.\\
\begin{quote}
\begin{align*}
T_{2,0}&=\sum\limits_{k=0}^{2}\frac{f^k(0)}{k!}(x-0)^k=\frac{f(0)}{0!}(x-0)^0+\frac{f'(0)}{1!}(x-0)^1+\frac{f''(0)}{2!}(x-0)^2\\
&=1+5x-\frac{5}{2}x^2\\
\end{align*}
\end{quote}
\subsection{Geometrische Interpretation der Taylorapproximation Teil I+II}
\href{https://youtu.be/fznq-R-W_NQ}{Video 07}\\
Die Funktion $ f(x)=\mathrm{e}^x $ kann durch ein Taylorpolynom im Entwicklungspunkt $ x_0=0 $ dargestellt werden. Dabei gilt dann:\\
\begin{align*}
f(x)=\mathrm{e}^x\approx T_{n,x_0}\;\;\text{ für}\;\;x\;\;\text{in der Nähe von }\;\;x_0\\
\end{align*}
In der Nähe von $ x_0 $ sieht das Taylorpolynom $T_{n,x_0}  $ wie $ f(x)=\mathrm{e}^x $ aus.\\
Es gilt auch:\\
\begin{align*}
f(x)=\mathrm{e}^x=T_{n,x_0}+\text{Fehler }\\
\end{align*}
Es gilt für $ f(x)=\mathrm{e}^x $:\\
\begin{align*}
f(x)&=\mathrm{e}^x\\
f'(x)&=\mathrm{e}^x\\
f^{k}(x)&=\mathrm{e}^x\\
f^{k}(0)&=1\\
\end{align*}
Das Taylorpolynom für $ f(x)=\mathrm{e}^x $ im Entwicklungspunkt $ x_0=0 $ heißt dann:\\
\begin{align*}
T_{n,0}(x)=\sum\limits_{k=0}^{n}\frac{f^k(0)}{k!}(x-0)^k=\frac{f(0)}{0!}(x-0)^0+\frac{f'(0)}{1!}(x-0)^1+\frac{f''(0)}{2!}(x-0)^2+\frac{f^{3}(0)}{3!}(x-0)^3+\ldots
\end{align*}
\underline{Für $ n=1 $ gilt:}\\
\begin{align*}
T_{1,0}(x)0&=f(0)+\frac{f'(0)}{1!}(x-0)\\
&=1+x\\
\end{align*}
\begin{center}
\begin{tikzpicture}


% Gitternetzlinien
\draw[xstep=1cm, ystep=1cm, very thin, color=lightgray] (-4,4) grid (4,-2);
% x-Achse
\draw[->] (-4,0) -- (4,0) node[below,xshift=-1em] {$x$}; 
%Zahlen auf x-Achse
\foreach \x in {-3,...,-1}
\draw[shift={(\x,0)},color=black] (0pt,2pt) -- (0pt,-2pt) node[below] {\footnotesize $\x$};
\foreach \x in {1,...,3}
\draw[shift={(\x,0)},color=black] (0pt,2pt) -- (0pt,-2pt) node[below] {\footnotesize $\x$};
% y-Achse 
\draw[->] (0,-2) -- (0,4) node[left,yshift=-1em] {$y$};%node[above left]
%Zahlen auf y-Achse
\foreach \y in {-1,...,-1}
\draw[shift={(0,\y)},color=black] (2pt,0pt) -- (-2pt,0pt) node[left] {\footnotesize $\y$};
\foreach \y in {1,...,3}
\draw[shift={(0,\y)},color=black] (2pt,0pt) -- (-2pt,0pt) node[left] {\footnotesize $\y$};
%Ursprung
\draw[color=black] (0pt,-10pt) node[right] {\footnotesize $0$};

\draw[thick,color=blue] plot[samples=200, domain=-3:3] 
(\x,{\x +1}) node[below, xshift=1em,yshift=-0.8em] {$T_{1,0}(x)$};

\draw[thick,color=black] plot[samples=200, domain=-3:1.35] 
(\x,{exp(\x)}) node[below, xshift=1em]{$f(x)$};



\end{tikzpicture}\\
Dies ist die Tangente im Punkt (0,1)\\
\end{center}
\underline{Für $ n=2 $ gilt:}\\
\begin{align*}
T_{2,0}(x)0&=f(0)+\frac{f'(0)}{1!}(x-0)+\frac{f''(0)}{2!}(x-0)^2\\
&=T_{1,0}+\frac{f''(0)}{2!}(x-0)^2\\
&=1+x+\frac{1}{2}x^2\\
\end{align*}
\begin{center}
\begin{tikzpicture}

% Gitternetzlinien
\draw[xstep=1cm, ystep=1cm, very thin, color=lightgray] (-4,4) grid (4,-2);
% x-Achse
\draw[->] (-4,0) -- (4,0) node[below,xshift=-1em] {$x$}; 
%Zahlen auf x-Achse
\foreach \x in {-3,...,-1}
\draw[shift={(\x,0)},color=black] (0pt,2pt) -- (0pt,-2pt) node[below] {\footnotesize $\x$};
\foreach \x in {1,...,3}
\draw[shift={(\x,0)},color=black] (0pt,2pt) -- (0pt,-2pt) node[below] {\footnotesize $\x$};
% y-Achse 
\draw[->] (0,-2) -- (0,4) node[left,yshift=-1em] {$y$};%node[above left]
%Zahlen auf y-Achse
\foreach \y in {-1,...,-1}
\draw[shift={(0,\y)},color=black] (2pt,0pt) -- (-2pt,0pt) node[left] {\footnotesize $\y$};
\foreach \y in {1,...,3}
\draw[shift={(0,\y)},color=black] (2pt,0pt) -- (-2pt,0pt) node[left] {\footnotesize $\y$};
%Ursprung
\draw[color=black] (0pt,-10pt) node[right] {\footnotesize $0$};

\draw[thick,color=blue] plot[samples=200, domain=-3:1.5] 
(\x,{\x +1+0.5*\x*\x}) node[below, xshift=1em,yshift=-0.8em] {$T_{2,0}(x)$};

\draw[thick,color=black] plot[samples=200, domain=-3:1.35] 
(\x,{exp(\x)}) node[below, xshift=-9em,yshift=-7.5em]{$f(x)$};



\end{tikzpicture}\\
Die Annäherung um den Punkt $ x_0 $ ist bei $ n=2 $ besser als bei $ n=1 $.\\
\end{center}
\href{https://youtu.be/rOYhaQnvv98}{Video 08}\\
\underline{Für $ n=3 $ gilt:}\\
\begin{align*}
T_{3,0}&=\underbrace {1+x+\frac{1}{2}x^2}_{T_{2,0}}+\frac{f^{3}(0)}{3!}(x-0)^3\\
&=1+x+\frac{1}{2}x^2+\frac{1}{6}x^3\\
\end{align*}
\begin{center}
\begin{tikzpicture}

% Gitternetzlinien
\draw[xstep=1cm, ystep=1cm, very thin, color=lightgray] (-4,4) grid (4,-2);
% x-Achse
\draw[->] (-4,0) -- (4,0) node[below,xshift=-1em] {$x$}; 
%Zahlen auf x-Achse
\foreach \x in {-3,...,-1}
\draw[shift={(\x,0)},color=black] (0pt,2pt) -- (0pt,-2pt) node[below] {\footnotesize $\x$};
\foreach \x in {1,...,3}
\draw[shift={(\x,0)},color=black] (0pt,2pt) -- (0pt,-2pt) node[below] {\footnotesize $\x$};
% y-Achse 
\draw[->] (0,-2) -- (0,4) node[left,yshift=-1em] {$y$};%node[above left]
%Zahlen auf y-Achse
\foreach \y in {-1,...,-1}
\draw[shift={(0,\y)},color=black] (2pt,0pt) -- (-2pt,0pt) node[left] {\footnotesize $\y$};
\foreach \y in {1,...,3}
\draw[shift={(0,\y)},color=black] (2pt,0pt) -- (-2pt,0pt) node[left] {\footnotesize $\y$};
%Ursprung
\draw[color=black] (0pt,-10pt) node[right] {\footnotesize $0$};

\draw[thick,color=blue] plot[samples=200, domain=-2.9:1.4] 
(\x,{\x +1+0.5*\x*\x+0.17*\x*\x*\x}) node[below, xshift=1em,yshift=-0.8em] {$T_{3,0}(x)$};

\draw[thick,color=black] plot[samples=200, domain=-3:1.35] 
(\x,{exp(\x)}) node[below, xshift=-9em,yshift=-7.5em]{$f(x)$};



\end{tikzpicture}\\
Hier ist die Approximation in der Umgebung von $ x_0=0 $ noch besser.
\end{center}
Graphischer Vergleich der Taylorpolynome.\\
\begin{center}
\begin{tikzpicture}

% Gitternetzlinien
\draw[xstep=1cm, ystep=1cm, very thin, color=lightgray] (-4,4) grid (4,-2);
% x-Achse
\draw[->] (-4,0) -- (4,0) node[below,xshift=-1em] {$x$}; 
%Zahlen auf x-Achse
\foreach \x in {-3,...,-1}
\draw[shift={(\x,0)},color=black] (0pt,2pt) -- (0pt,-2pt) node[below] {\footnotesize $\x$};
\foreach \x in {1,...,3}
\draw[shift={(\x,0)},color=black] (0pt,2pt) -- (0pt,-2pt) node[below] {\footnotesize $\x$};
% y-Achse 
\draw[->] (0,-2) -- (0,4) node[left,yshift=-1em] {$y$};%node[above left]
%Zahlen auf y-Achse
\foreach \y in {-1,...,-1}
\draw[shift={(0,\y)},color=black] (2pt,0pt) -- (-2pt,0pt) node[left] {\footnotesize $\y$};
\foreach \y in {1,...,3}
\draw[shift={(0,\y)},color=black] (2pt,0pt) -- (-2pt,0pt) node[left] {\footnotesize $\y$};
%Ursprung
\draw[color=black] (0pt,-10pt) node[right] {\footnotesize $0$};

\draw[thick,color=blue] plot[samples=200, domain=-2.9:1.4] 
(\x,{\x +1+0.5*\x*\x+0.17*\x*\x*\x}) node[below, xshift=-10.9em,yshift=-10.6em] {$T_{3,0}(x)$};

\draw[thick,color=black] plot[samples=200, domain=-3:1.35] 
(\x,{exp(\x)}) node[below, xshift=-9em,yshift=-7.5em]{$f(x)$};

\draw[thick,color=red] plot[samples=200, domain=-3:1.5] 
(\x,{\x +1+0.5*\x*\x}) node[below, xshift=-9em,yshift=-.8em] {$T_{2,0}(x)$};
\draw[thick,color=green] plot[samples=200, domain=-3:3] 
(\x,{\x +1}) node[below, xshift=1em,yshift=-0.8em] {$T_{1,0}(x)$};
\draw[thick,color=blue,dashed] (1,0)--(1,3) ;
\end{tikzpicture}\\
\end{center}
Bei $ x=1 $ kann man sehr gut die unterschiedlichen Fehler der Approximation ablesen.\\
Fehlerberechnung:\\
\begin{align*}
|T_{1,0}-f(1)|=|1+1-2,718|=0.718\\
|T_{2,0}-f(1)|=|1+1+\frac{1}{2}\cdot 1^2-2,718|=0.281\\
|T_{3,0}-f(1)|=|1+1+\frac{1}{2}\cdot 1^2+\frac{1}{6}\cdot 1^3-2,718|=0.051\\
\end{align*}
Den Fehler bei dem Taylorpolynom  $ T_{3,0} $ kann man graphisch bei $ x=1 $ schon nicht mehr ablesen.\\
\newpage
\noindent
\subsubsection{Die Taylorapproximation}
\href{https://youtu.be/sZeivAKBMa0}{Video 09}\\
Sei $ f: \mathbb{I}\longmapsto \mathbb{R} $ n-mal diffbar auf dem Intevall $  \mathbb{I} $ und sei $ x_0\in  \mathbb{I} $ dann gilt:\\
\begin{align*}
f(x)\approx T_{n.x_0}(x)\\
\end{align*}
Dieses Polynom entspricht genau der  Funktion $ f(x) $, wenn man den Fehler $ R_{n}(x) $ dem Polynom dazu addiert. Dann gilt:\\
\begin{align*}
f(x)=T_{n,x_0}(x)+R_{n}(x)=\sum\limits_{k=0}^{n}\frac{f^{k}(x_0)}{k!}(x-x_0)^k+R_n(x)\\
\end{align*}
Die Tayloapproximation entspricht: Polynom + Fehler.\\
Für diesen Fehler gilt:\\
\begin{align*}
\lim\limits_ {x\to x_0}\frac{R_n(x)}{(x-x_0)^n}=0\\
\end{align*}
Dies bedeutet, obwohl $ \frac{1}{(x-x_0)^n} $ schon sehr schnell gegen Null geht (wegen den hoch n), dass der Zähler (Fehler)  noch schneller gegen Null läuft.\vspace*{4mm}\\
Ist $ f^{n+1} $ n-mal diffbar, dann gilt, wenn $ a $ zwischen $ x $ und $ x_0 $ liegt :\\
\begin{center}
Restglied von Lagrange\\
$\boxed{R_n(x)=\frac{f^{n+1}(a)}{(n+1)!}(x-x_0)^{n+1}\;\;a\in[x,x_0]}$
\end{center}
\subsubsection{Beispiel zur Taylorapproximation}
\href{https://youtu.be/T6jBJ9BVPg8}{Video 10}\\
Bestimme die Taylorapproximation $ T_{1,x_0} $ von  $ f(x)=\mathrm{e}^x $ mit $ x_0=0 $.\\
Es gilt:\\
\begin{align*}
f(x)=\mathrm{e}^x&=T_{1,x_0}+R_1(x)\\
&=f(0)+\frac{f'(0)}{1!}(x-0)^1+R_1(x)\\
&=1+x+R_1(x)\\
\end{align*}
Restglied $ R_{1}(x) $ bestimmen:\\
Es gilt:\\
\begin{align*}
R_n(x)=\frac{f^{n+1}(a)}{(n+1)!}(x-x_0)^{n+1}\\
\end{align*}
\begin{align*}
R_1(x)=\frac{f''(a)}{2!}(x-x_0)^2=\frac{\mathrm{e}^a}{2}x^2\\
\end{align*}
\textbullet\;\;Es muß gelten:\\
\begin{align*}
\lim\limits_ {x\to x_0}\frac{R_1(x)}{(x-x_0)^1}=0
\end{align*}
\begin{align*}
\lim\limits_ {x\to 0}=\frac{\frac{\mathrm{e}^a}{2}x^2}{x}=\lim\limits_ {x\to 0}\frac{\mathrm{e}^a}{2}x=0
\end{align*}
\textbullet\;\;Restglied abschätzen:\\
\begin{quote}
1. Für $ x\in \;[0,2] $:\\
\begin{align*}
R_1(x)=\frac{\mathrm{e}^a}{2}x^2\leq \frac{\mathrm{e}^a}{2}\cdot 2=\mathrm{e}^a\cdot 2\\
\end{align*}
Da $ a $ zwischen $ x_0 $ und $  x $ liegt, gilt:\\
Abschätzung nach oben:\\
\begin{align*}
R_1(x)=\frac{\mathrm{e}^a}{2}x^2\leq \mathrm{e}^a\cdot 2 < \mathrm{e}^2\cdot 2
\end{align*}
Abschätzung nach unten:\\
\begin{align*}
0=\mathrm{e}^0\cdot 0<R_1(x)&=\frac{\mathrm{e}^a}{2}x^2\leq \mathrm{e}^a\cdot 2 < \mathrm{e}^2\cdot 2\\
0<R_1(x)& < \mathrm{e}^2\cdot 2\\
\end{align*}
2. Für $ x\in \;[0,200] $:\\
\begin{align*}
R_1(x)=\frac{\mathrm{e}^a}{2}x^2\leq \frac{\mathrm{e}^a}{2}\cdot 200^2<\frac{\mathrm{e}^{200}}{2}\cdot 200^2
\end{align*}
Hier ist der Fehler riesengroß!\\

\end{quote}

\subsubsection{Aufgabe 1 zur Taylorapproximation Teil I u. Teil II}
\href{https://youtu.be/NbAlqm36CTo}{Video 11}\\
Berechne das Taylorpolynom 3. Grades von $ f(x)=\cos(x) $ in $ x_0=\pi $ und approximiere $ f(3)=\cos(3) $ damit. Schätze danach den Fehler ab.\\
\textbullet\;\; Bestimmung der Ableitung $ f'(x),\;f''(x),\;f^{(3)}(x) $.\\
\begin{align*}
f(x)&=\cos(x)\\
f'(x)&=-\sin(x)\\
f''(x)&=-\cos(x)\\
f^{(3)}(x)&=\sin(x)\\
\end{align*}
\textbullet\;\; Bestimmung von $ f(\pi),\; f'(\pi),\; f''(\pi),\; f^{(3)}(\pi) $.\\
\begin{align*}
f(\pi)&=-1\\
f'(\pi)&=0\\
f''(\pi)&=1\\
f^{(3)}(\pi)&=0
\end{align*}
\textbullet\;\; Einsetzen in das Taylorpolynom.\\
\begin{quote}
\begin{align*}
T_{3,\pi}&=\sum\limits_{k=0}^{3}\frac{f^k(\pi)}{k!}(x-\pi)^k=\frac{f(\pi)}{0!}(x-\pi)^0+\frac{f'(\pi)}{1!}(x-\pi)^1+\frac{f''(\pi)}{2!}(x-\pi)^2+\frac{f^{(3)}(\pi)}{3!}(x-\pi)^3\\
&=-1+\frac{1}{2}(x-\pi)^2\\
\end{align*}
\end{quote}
\textbullet\;\;Für $ \cos(3) $ approximiert gilt dann:\\
\begin{align*}
&T_{3,\pi}\approx-1+\frac{1}{2}(3-\pi)^2=-0.9899757\\ 
&\text{mit Taschenrechner }\;\;\cos(3)=0.9899924\\
\end{align*}
\href{https://youtu.be/sWuEka8oIQY}{Video 12}\\
\textbullet\;\;Fehlerabschätzung:\\
\begin{quote}
Es gilt:\\
\begin{align*}
&f(x)=T_{3,\pi}(x)+R_3(x)\\
&\Leftrightarrow\\
&R_3(x)=f(x)-T_{3,\pi}(x)\\
\end{align*}
Mit dem Restglied von Lagrange wird der Fehler abgeschätzt.\\
\begin{align*}
R_3(x)=\frac{f^{(4)}(a)}{4!}(x-x_0)^{4}\;\;a\in[x,x_0]\\
\end{align*}
\begin{align*}
f^{(4)}&=\cos(x)\\
x&=3\\
x_0&=\pi\\
\end{align*}
Bestimmung von $ a $:\\
\begin{align*}
&a\in[3,\pi]\\
&\cos(3)=-0.989\\
&\cos(\pi)=-1
\end{align*}
Da der Fehler absolut ist, wählt man den $ \cos(a) $ so, dass er den größten Fehler hat.\\
Dies gilt bei $ \cos(a)=\cos(\pi)=-1 $ also $ a=\pi $.\\
Dann gilt:\\
\begin{align*}
\big|R_{3}(x)\big|<\bigg|\frac{\cos(a)}{4!}(3-\pi)^4\bigg|<\bigg|\frac{\cos(\pi)}{4!}(3-\pi)^4\bigg|=\bigg|\frac{-(3-\pi)^4}{4!}\bigg|
\end{align*}
Dieser Fehler ist hier schon sehr klein!\\
\end{quote}
\subsubsection{Aufgabe 1 zur Taylorapproximation Teil I u. Teil II}
\href{https://youtu.be/iHpIwbXhDLs}{Video 13}\\
Bestimme für $ f(x)=\sqrt{2+x} $ das Taylorpolynom 3.Grades in $ x_0=-1 $.\\
Schätze den Fehler für $ x\in\;]-\frac{1}{4},\frac{1}{4}[ $ ab.\\
\textbullet\;\; Bestimmung der Ableitung $ f'(x),\;f''(x),\;f^{(3)}(x) $.\\
\begin{align*}
f(x)&=\sqrt{2+x}\\
f'(x)&=\frac{1}{2\sqrt{2+x}}\\
f''(x)&=\bigg(\frac{1}{2}(2+x)^{-\frac{1}{2}}\bigg)'=-\frac{1}{4}(2+x)^{-\frac{3}{2}}=-\frac{1}{4}\frac{1}{\sqrt{(2+x)^3}}\\
f^{(3)}(x)&=\bigg(-\frac{1}{4}(2+x)^{-\frac{3}{2}}\bigg)'=\frac{3}{8}(2+x)^{-\frac{5}{2}}=\frac{3}{8}\frac{1}{\sqrt{(x+2)^{5}}}\\
\end{align*}
\textbullet\;\; Bestimmung von $ f(-1),\; f'(-1),\; f''(-1),\; f^{(3)}(-1) $.\\
\begin{align*}
f(-1)&=1\\
f'(-1)&=\frac{1}{2}\\
f''(-1)&=-\frac{1}{4}\\
f^{3}(-1)&=\frac{3}{8}\\
\end{align*}
\textbullet\;\; Einsetzen in das Taylorpolynom.\\
\begin{align*}
T_{3,-1}&\approx \frac{f(-1)}{0!}+\frac{f'(-1)}{1!}(x+1)+\frac{f''(-1)}{2!}(x+1)^2+\frac{f^{3}(-1)}{3!}(x+1)^3\\
&\approx 1+\frac{1}{2}(x+1)-\frac{1}{8}(x+1)^2+\frac{1}{16}(x+1)^3\\
\end{align*}
\href{https://youtu.be/KGQM0psy20I}{Video 14}\\
\textbullet\;\;Mit dem Restglied von Lagrange wird der Fehler abgeschätzt.\\
\begin{align*}
\big|R_3(x)\big|=\bigg|\frac{f^{(4)}(a)}{4!}(x-x_0)^{4}\bigg|\;\;a\in[x,x_0], \; x\in\;]-\genfrac{}{}{}{2}{1 }{4 },\genfrac{}{}{}{2}{1 }{4 }[\\
\end{align*}
\begin{align*}
f^{(4)}(a)=\bigg(\frac{3}{8}\frac{1}{\sqrt{(x+2)^{5}}}\bigg)'=-\frac{15}{16}\frac{1}{\sqrt{(2+a)^7}}
\end{align*}
\begin{align*}
\big|R_3(x)\big|=\bigg|-\frac{15}{16}\frac{1}{\sqrt{(2+a)^7}4!}(x+1)^4\bigg|
\end{align*}
\begin{quote}
Abschätzen des Terms $a\in [-1,x]$ und $ x\in\;]-\frac{1}{4},\frac{1}{4}[ $\\
Für $ x=\frac{1}{4} $ wird der Term $ (x+1)^4 $ am größten.\\
Für $ a=-1 $ wird der Nenner $ \sqrt{(2+a)^7} $ am kleinsten und der Bruch am größten.\\
Der größtmögliche Fehler ist dann:\\
\begin{align*}
\big|R_3(x)\big|<\bigg|-\frac{15}{16}\frac{1}{\sqrt{(2-1)^7}4!}\bigg(\frac{1}{4}+1\bigg)^4\bigg|=\frac{5\cdot3}{4\cdot4}\cdot\frac{1}{4\cdot3\cdot 2}\cdot\frac{5^4}{4^4}=\frac{5^5}{2\cdot4^7}\\
\end{align*}


\end{quote}
\subsection{Das hinreichende Kriterium für Extrema}
\href{https://youtu.be/qZN36TCHIVw}{Video 15}\\
Sei $ f:  \mathbb{I}\longmapsto \mathbb{R} $ auf dem Intervall $  \mathbb{I} $ n- mal differenzierbar.\\
Sei $ x_0\in  \mathbb{I} $ ein innerer Punkt.\\
Ist $ f'(x_0)=f''(x_0) =\ldots =f^{n-1}(x_0)=0$ und $ f^{n}(x_0)\neq 0 $ dann gilt:\\
\begin{quote}
(i)\;\; Ist $ n $ ungerade $ \Rightarrow  f$ kein Extrema in $ x_0 $\\
(ii)\;\;Ist $ n $ gerade dann gilt für:
\begin{quote}
1. $f^n > 0 \Rightarrow f$ lokales Minimum in $ x_0 $\\
2. $f^n < 0 \Rightarrow f$ lokales Maximum in $ x_0 $\\
\end{quote}
\end{quote}
Beispiel:\\
\textbullet \;\;$ f:\mathbb{R}\longmapsto[0,\infty[\;\;\text{mit }\;\; f(x)=x^2 $\\
\begin{quote}
$ f $ hat in $ x_0=0 $ ein lokales Minimum, denn:\\
\begin{align*}
f'(x_0)&=2x_0=0\\
f''(x_0)&=2>0\\
\end{align*}
\end{quote}
\newpage
\noindent
\textbullet\;\;$\overset{\AC}{f}: \mathbb{R}\longmapsto\mathbb{R}\;\;\text{mit }\;\;\overset{\AC}{f}(x)=x^3$\\
\begin{quote}
$\overset{\AC}{f}$ hat in $ x_0=0 $ kein Extrema, denn:\\
\begin{align*}
\overset{\AC}{f'}(x_0)&=3x_0^2=0=6x_0=\overset{\AC}{f''}(x_0)\\
\overset{\AC}{f^{3}}(x_0)&=6\\
\end{align*}
Da $ \overset{\AC}{f^{3}}(x_0)=6\neq 0 $ ungerade ist hat $\overset{\AC}{f}  $ in $ x_o $ kein Extrema.\\
\end{quote}
\subsubsection{Aufgabe 1 zum hinreichenden Kriterium für Extrema Teil I u. Teil II}
\href{https://youtu.be/Anjy9d6Qu2Y}{Video 16}\\
Bestimme alle lokalen und globalen Extrema von $ f:\;]-\frac{1}{2},3]\;\longmapsto \mathbb{R} $ mit $ f(x)=\mathrm{e}^{2x^3+3x^2} $.\\
\\
(I) Untersuchung der inneren Punkte: $ x\in\;]-\frac{1}{2},3] $\\
\begin{quote}
1. Bestimmung von $ f'(x) $ und die Nullstellen von $ f'(x) $.\\
\begin{quote}
\begin{align*}
f'(x)=\mathrm{e}^{2x^3+3x^2}(6x^2+6x)=0\\
\end{align*}
$ \mathrm{e}^{2x^3+3x^2} $ ist immer positiv.\\
$ f'(x)=0 $, wenn$ (6x^2+6x)=0 $\\
\begin{align*}
&6x^2+6x=0\\
&\Leftrightarrow\\
&\Leftrightarrow 6(x^2+x)=0\\
&\Leftrightarrow 6x(x+1)=0\\
&x_1=0\\
&x_2=-1\\
\end{align*}
$ x_2 \notin \;]-\genfrac{}{}{}{2}{1 }{2 },3] $\\
Die einzige potentielle Extremstelle im Intervall $ ]-\frac{1}{2},3] $ ist $x_1=0  $\\
\end{quote}
2. $ x_1=0 $ in 2. Ableitung einsetzen.\\
\begin{quote}
\begin{align*}
f''(x)&=\mathrm{e}^{2x^3+3x^2}(6x^2+6x)^2+\mathrm{e}^{2x^3+3x^2}(12x+6)\\
f''(x_1)&=f''(0)=0+\mathrm{e}^0\cdot 6=6>0
\end{align*}
Also hat $ f $ in $ x_1 $ ein lokales Minimum.\\
\end{quote}
\end{quote}
\href{https://youtu.be/xKTPSfSrtGo}{Video 17}\\
(II) Randwertbetrachtung und Grenzwertbetrachtung an den Rändern vom Intervall  $ x\in\;]-\frac{1}{2},3] $ .\\
\begin{quote}
1. An der Stelle $ x=3 $.
\begin{align*}
&f(x)=\mathrm{e}^{2x^3+3x^2}\\
&f(3)=\mathrm{e}^{2\cdot3^3+3\cdot3^2}=\mathrm{e}^{3\cdot27}\\
\end{align*}
2. Grenzwert an der Stelle $ x\mapsto -\frac{1}{2} $
\begin{align*}
\lim\limits_ {\substack{x \to -\frac{1}{2} \\ x>-\frac{1}{2}}}\mathrm{e}^{-2\cdot \frac{1}{8}+3\cdot\frac{1}{4}}=\mathrm{e}^{\frac{1}{2}}=\sqrt{e}\\
\end{align*}
\end{quote}
Fazit:\\
\begin{align*}
f(0)=\mathrm{e}^{2\cdot 0+3\cdot 0}=\mathrm{e}^0=1<\sqrt{e}<\mathrm{e}^{3\cdot 27}\\
\end{align*}
D.h.:\\
In $ x_1=0 $ hat $ f $ globales Minimum.\\
In $ x_2=3 $ hat $ f $ globales Maximum.
\subsubsection{Aufgabe 2 zum hinreichenden Kriterium für Extrema}
\href{https://youtu.be/JRMKUC_oVdM}{Video 18}\\
Bestimme alle lokalen und globalen Extrema von: $ \mathbb{R}\longmapsto \mathbb{R} $ mit 
$f(x) =     \left\{ \begin{array}{rcl}
\mathrm{e}^{6x^2} & \mbox{für} &  x\geq 0\\
6x^2+1 & \mbox{für} & x<0 \\

\end{array}\right.$\\
$ f(x) $ ist stetig und bis zur 2.Ableitung diffbar.\\
\\
Nullstellen der 1. Ableitung.\\
\begin{center}
$f'(x) =     \left\{ \begin{array}{rcl}
\mathrm{e}^{6x^2}\cdot12x & \mbox{für} & x\geq0\\
12x & \mbox{für} & x< 0 \\

\end{array}\right.$\\
\end{center}
\begin{align*}
f'(x)=0\Leftrightarrow \mathrm{e}^{6x^2}\cdot 12x\Rightarrow x_1=0\\
\end{align*}
Potenzielle Extremstelle ist $ x_1=0 $\\
$ x_1=0 $ in $ f''(x) $ einsetzen.\\
\begin{center}
$f''(x) =     \left\{ \begin{array}{rcl}
\mathrm{e}^{6x^2}(\cdot12x)^2+\mathrm{e}^{6x^2}\cdot12 & \mbox{für} & x\geq0\\
12 & \mbox{für} & x< 0 \\

\end{array}\right.$\\
\end{center}
\begin{align*}
&f''(0)=\mathrm{e}^{6\cdot 0^2}(\cdot12\cdot 0)^2+\mathrm{e}^{6\cdot 0^2}\cdot12=12>0\\
&\Rightarrow f\;\;\text{hat in  }\;\;x_1\;\;\text{ein lokales Minimum }\\
\end{align*}
Grenzwertbetrachtung:\\
\begin{align*}
&\text{\textbullet} \;\;\lim\limits_ {x\to\infty}f(x)=\lim\limits_ {x\to\infty}\mathrm{e}^{6x^2}=\infty\\
&\text{\textbullet }\;\;\lim\limits_ {x\to -\infty}f(x)=\lim\limits_ {x\to -\infty}6x^2+1=\infty\\
\end{align*}
Fazit: 
$ f $ hat in $ x_1 $ ein globales Minimum.
\\
\\
Bemerkung:
Bei der Grenzwert- oder Randwertbetrachten immer die $ \mathbb{D}_f $ und $ \mathbb{W}_f $ beachten, damit man richtig zwischen globalen und lokalen Maximun oder Minimum beurteilt.\\
Z.B:
\begin{align*}
\mathbb{R}&\longmapsto \mathbb{R}\\
[a,b]&\longmapsto ]a,b[\\
\mathbb{R}&\longmapsto ]a,b]\\
\end{align*}



%	\chapter{Elementare Funktionen}

\setcounter{section}{10}
\setcounter{subsection}{0}
\setcounter{subsubsection}{1}
\setcounter{secnumdepth}{4}
\subsection{Sinus und Kosinus}
\href{https://youtu.be/8ApuQ6_PooA}{Video 01}\\
Es gibt eine eindeutige Funktion $ f:\mathbb{R}\longmapsto \mathbb{R} $ mit $ f+f''=0 $ und $ f(0)=0 $ sowie $ f'(0) =1$.\\
Wir schreiben $ f=\sin $ und $ f'=\cos $ und sagen Sinus bzw. Kosinus.\\
\begin{center}
	\begin{tikzpicture}[xscale=1,yscale=2][xscale=1,yscale=2]
	
	% Gitternetzlinien
	\draw[xstep=1cm, ystep=1cm, very thin, color=lightgray] (-4,2) grid (7,-2);
	% x-Achse
	\draw[->] (-4,0) -- (7,0) node[below,xshift=-1em] {$x$}; 
	%Zahlen auf x-Achse
	\foreach \x in {-3.14,-3.14/2}
	\draw[shift={(\x,0)},color=black] (0pt,2pt) -- (0pt,-2pt);
	\foreach \x in {3.14/2,3.14,4.71,6.28}
	\draw[shift={(\x,0)},color=black] (0pt,2pt) -- (0pt,-2pt);
	% y-Achse 
	\draw[->] (0,-2) -- (0,2) node[left,yshift=-1em] {$f(x)$};%node[above left]
	%Zahlen auf y-Achse
	\foreach \y in {-2,...,-1}
	\draw[shift={(0,\y)},color=black] (2pt,0pt) -- (-2pt,0pt) node[left] {\footnotesize $\y$};
	\foreach \y in {1,...,1}
	\draw[shift={(0,\y)},color=black] (2pt,0pt) -- (-2pt,0pt) node[left] {\footnotesize $\y$};
	%Ursprung
	\draw[color=black] (0pt,-10pt) node[right] {\footnotesize $0$};
	\draw[color=black] (3.14,-2pt) node[below] {\footnotesize $\pi$};
	\draw[color=black] (6.28,-2pt) node[below] {\footnotesize $2\pi$};
	\draw[color=black] (1.57,-2pt) node[below] {\footnotesize $\frac{\pi}{2}$};
	\draw[color=black] (-1.57,-2pt) node[below] {\footnotesize $-\frac{\pi}{2}$};
	\draw[color=black] (-3.14,-2pt) node[below] {\footnotesize $-\pi$};
	\draw[thick,color=blue] plot[samples=200, domain=-4:7] 
	(\x,{sin(\x r)}) node[below, xshift=-2em]{$\sin$};
	
	\draw[thick,color=red] plot[samples=200, domain=-4:7] 
	(\x,{cos(\x r)}) node[ xshift=-5em]{$\cos$};
	
	\end{tikzpicture}\\	
\end{center}
\underline{Eigenschaften}:\\
Für alle $ x\in \mathbb{R} $ gilt:\\
\begin{quote}
	\textbullet\;\;$ \sin(-x)=-\sin(x) $
	\begin{quote}
		ungerade Funktion \\
		punktsymmetrisch zum Koordinatenursprung\\
	\end{quote}
	\textbullet\;\;$ \cos(-x)=\cos(x) $
	\begin{quote}
		gerade Funktion\\
		symmetrisch zur y-Achse\\
	\end{quote}
	\newpage
	\noindent
	\textbullet\;\;$ \sin, \cos: \mathbb{R}\longmapsto[-1,1] $
	\begin{quote}
		$|\cos(x)|\leq 1$\\
		$|\sin(x)|\leq 1$\\
	\end{quote}
	\textbullet\;\;$\sin(k\cdot\pi)=0$ für $ k\in \mathbb{Z} $\\ \\
	\textbullet\;\;$\sin(\frac{\pi}{2}+k\frac{\pi}{2})=(-1)^k$ für $ k\in \mathbb{Z} $\\ \\
	\textbullet\;\;$\cos(\frac{\pi}{2}+k\frac{\pi}{2})=0 $ für $ k\in \mathbb{Z} $\\ \\
	\textbullet\;\;$ \cos(k\cdot\pi)=(-1)^k $ für $ k\in \mathbb{Z} $\\
\end{quote}
\subsubsection{Additionstheoreme des Sinus und Kosinus}\label{sec:Bezeichnung}
\href{https://youtu.be/Q-yJWE6P4uU}{Video 02}\\
Für $ x_1, x_2 \in\mathbb{R} $ gilt:\\ \\
(i)\;\;$ \sin(x_1+x_2)=\sin(x_1)\cdot \cos (x_2)+\sin(x_2)\cdot \cos(x_1)$\\ \\
(ii)\;$ \cos(x_1+x_2)=\cos(x_1)\cdot \cos(x_2) - \sin(x_1)\cdot \
sin(x_2)$\\ \\
Für $ x\in\mathbb{R} $ gilt:\\
\textbullet \;\;$\sin(x+\frac{\pi}{2})=\cos(x)$\\
\begin{align*}
\sin(x+\genfrac{}{}{}{2}{\pi }{2 })&\overset{(\mathrm{i})}{=}\sin(x)\cdot \underbrace {\cos(\genfrac{}{}{}{2}{\pi }{2 })}_{=0}+\underbrace {\sin(\genfrac{}{}{}{2}{\pi}{2 })}_{1}\cdot \cos(x)\\
&=\cos(x)
\end{align*}
\textbullet\;\;$\cos(2x)=\cos^2(x)-\sin^2(x)$\\
\begin{align*}
\cos(2x)=\cos(x+x)&\overset {(\mathrm{ii})}{=}\cos(x)\cdot\cos(x)-\sin(x)\cdot\sin(x)\\
&=\cos^2(x)-\sin^2(x)\\
\end{align*}
\textbullet \;\;$ \sin(2x)=2\sin(x)\cdot\cos(x) $\\
\begin{align*}
\sin(2x)=\sin(x+x)&\overset{(\mathrm{i})}{=}\sin(x)\cdot \cos(x)+\sin(x)\cdot \cos(x)\\
&=2\sin(x)\cdot\cos(x) \\
\end{align*}
\textbullet\;\;$\cos^2(x)+\sin^2(x)=1$\\
\newpage
\noindent
\subsubsection{ Aufgabe 1 zu den Additionstheoremen des Sinus und Kosinus}
\href{https://youtu.be/fVSPz-IaCqA}{Video 03}\\
Zeige für alle $ x\in\mathbb{R} $ gilt $ \cos^2(x)+\sin^2(x) =1$.\\ \\
Es gilt für alle $ x\in\mathbb{R} $:\\
\begin{align*}
\mathrm{e}^{\mathrm{i}x}=\cos(x)+\mathrm{i}\cdot\sin(x)
\end{align*}
Betrachte:\\
\begin{align*}
\cos^2(x)+\sin^2(x)&= \cos^2(x)\underbrace {-\mathrm{i^2}}_{=1}\cdot \sin^2(x)\;\;|\text{\tiny $\mathrm{i}^2=-1\Rightarrow -\mathrm{i}^2=1$ }\\
&=\cos^2(x)-(\mathrm{i}\cdot \sin(x))^2\;\;|\text{\tiny mit 3. binomischer Formel }\\
&=\bigg(\cos(x)-\mathrm{i}\cdot \sin(x\big)\bigg)\bigg(\cos(x)+\mathrm{i}\cdot \sin(x)\bigg)\;\;|\text{\tiny $ \cos(x)=\cos(-x),-\mathrm{i}\cdot \sin(x)=+\mathrm{i}\cdot\sin(-x) $ }\\
&=\bigg(\cos(-x)+\mathrm{i}\cdot\sin(-x)\bigg)\bigg(\cos(x)+\mathrm{i}\cdot \sin(x)\bigg)\;\;|\text{\tiny $\mathrm{e}^{\mathrm{i}x}=\cos(x)+\mathrm{i}\cdot\sin(x)  $}\\
&=\mathrm{e}^{-\mathrm{i}x}\cdot\mathrm{e}^{\mathrm{i}x}=\mathrm{e}^0=1_{_{\Square}}
\end{align*}
\subsubsection{ Aufgabe 2 zu den Additionstheoremen des Sinus und Kosinus}
\href{https://youtu.be/Ibl3wG8lgyE}{Video 04}\\
Bestimme alle $ x\in\mathbb{R} $ mit $ \sqrt{2}\sin(x)-\sin(2x)=0 $.\\
Bemerkung:\\
\begin{center}
	\renewcommand{\arraystretch}{1.8}
	\( \begin{array}{l|c|c|c|c|c} 
	x&0\overset{\wedge}{=}0^\circ&\frac{\pi}{6}\overset{\wedge}{=}30^\circ&\frac{\pi}{4}\overset{\wedge}{=}45^\circ&\frac{\pi}{3}\overset{\wedge}{=}60^\circ&\frac{\pi}{2}\overset{\wedge}{=}90^\circ\\
	\hline
	\sin(x)&\frac{\sqrt{0}}{2}&\frac{\sqrt{1}}{2}&\frac{\sqrt{2}}{2}&\frac{\sqrt{3}}{2}&\frac{\sqrt{4}}{2}\\
	\hline
	\cos(x)&\frac{\sqrt{4}}{2}&\frac{\sqrt{3}}{2}&\frac{\sqrt{2}}{2}&\frac{\sqrt{1}}{2}&\frac{\sqrt{0}}{2}
	\end{array} \)	
\end{center}
\begin{align*}
&\sqrt{2}\sin(x)-\sin(2x)=0\\
&\Leftrightarrow\\
&\sqrt{2}\sin(x)=\sin(2x)\;\;|\text{\tiny $ \sin(2x)=2\sin(x)\cos(x) $ }\\
&\sqrt{2}\sin(x)=2\sin(x)\cos(x)\\
\end{align*}
\underline{1.Fall:} $ \sin(x)\neq 0 $\\
\begin{align*}
\sqrt{2}\sin(x)&=2\sin(x)\cos(x)\;\;|\text{\tiny $\div \sin(x)  $  }\\
\sqrt{2}&=2\cos(x)\\
\frac{\sqrt{2}}{2}&=\cos(x)\\
\end{align*}
\begin{quote}
	D.h. eine Lösung bei $ x=\frac{\pi}{4}=\cos(\frac{\sqrt{2}}{2}) $ und alle Lösungen bei:\\
	\begin{align*}
	\text{Es gilt: }&\cos(x+k2\pi)=\cos(x)\;\;\text{mit }\;\; k\in\mathbb{Z}\\
	&\Rightarrow \cos(\genfrac{}{}{}{2}{\pi }{4 }+k2\pi) \;\;\text{mit }\;\;k\in\mathbb{Z}\\
	\end{align*}
\end{quote}
\underline{1.Fall:} $ \sin(x)= 0 $\\
\begin{align*}
\sqrt{2}\sin(x)=2\sin(x)\cos(x)=0\\
\end{align*}
\begin{quote}
	Immer erfüllt für alle Nullstellen von $ \sin(x) $ mit $ x=j\cdot\pi $ mit $ j\in  \mathbb{Z} $
\end{quote}
Die Lösungsmenge ist:\\
\begin{align*}
\mathbb{L}=\{\genfrac{}{}{}{2}{\pi }{4 }+k \cdot 2 \pi, j\cdot \pi | k,j\in  \mathbb{Z}\}
\end{align*}
oder:\
\begin{align*}
\mathbb{L}=\underset{k\in \mathbb{Z}}{\bigcup}\{\genfrac{}{}{}{2}{\pi}{4}+k2\pi,k\pi\}\\
\end{align*}
\subsection{Der Tangens und Kotangens}
\subsubsection{Der Tangens}
\href{https://youtu.be/zwQ_9S4sgro}{Video 05}\\
Defintion des Tangens:\\
\begin{align*}
\tan(x)=\frac{\sin(x)}{\cos(x)}\;\;\text{für }\;\;x\neq \frac{\pi}{2}+k\cdot\pi,\; k\in \mathbb{Z}
\end{align*}
Diese Funktion ist $ \pi- $periodisch.\\
Es gilt:\\
\begin{align*}
\tan(x+\pi)=\tan(x)
\end{align*}
\begin{center}
	\begin{tikzpicture}[xscale=1,yscale=2]
	
	% Gitternetzlinien
	\draw[xstep=1cm, ystep=1cm, very thin, color=lightgray] (-4,2) grid (7,-2);
	% x-Achse
	\draw[->] (-4,0) -- (7,0) node[below,xshift=-1em] {$x$}; 
	%Zahlen auf x-Achse
	\foreach \x in {-3.14,-3.14/2}
	\draw[shift={(\x,0)},color=black] (0pt,2pt) -- (0pt,-2pt);
	\foreach \x in {3.14/2,3.14,4.71,6.28}
	\draw[shift={(\x,0)},color=black] (0pt,2pt) -- (0pt,-2pt);
	% y-Achse 
	\draw[->] (0,-2) -- (0,2) node[left,yshift=-1em] {$f(x)$};%node[above left]
	%Zahlen auf y-Achse
	\foreach \y in {-2,...,-1}
	\draw[shift={(0,\y)},color=black] (2pt,0pt) -- (-2pt,0pt) node[left] {\footnotesize $\y$};
	\foreach \y in {1,...,1}
	\draw[shift={(0,\y)},color=black] (2pt,0pt) -- (-2pt,0pt) node[left] {\footnotesize $\y$};
	%Ursprung
	\draw[color=black] (0pt,-10pt) node[right] {\footnotesize $0$};
	\draw[color=black] (3.14,-2pt) node[below] {\footnotesize $\pi$};
	\draw[color=black] (6.28,-2pt) node[below] {\footnotesize $2\pi$};
	\draw[color=black] (1.57,-2pt) node[below] {\footnotesize $\frac{\pi}{2}$};
	\draw[color=black] (-1.57,-2pt) node[below] {\footnotesize $-\frac{\pi}{2}$};
	\draw[color=black] (-3.14,-2pt) node[below] {\footnotesize $-\pi$};
	\draw[color=black] (3*3.14/2,-2pt) node[below] {\footnotesize $\frac{3}{2}\pi$};
	\draw[thick,color=black] plot[samples=200, domain=-2.2/2:2.2/2] 
	(\x,{tan(\x r)}) node[below, xshift=-1em]{$\tan$};
	\draw[thick,color=black] plot[samples=200, domain=4.07/2:4.24] 
	(\x,{tan(\x r)}) node[below, xshift=-1em]{$\tan$};
	\draw[dashed,color=red] plot[samples=200, domain=-3.14/2:3.14/2] 
	(\x,{cos(\x r)}) node[ xshift=-9em,yshift=1em]{$\cos$};
	\draw[dashed,color=blue] plot[samples=200, domain=-3.14/2:3.14/2] 
	(\x,{sin(\x r)}) node[ xshift=1em]{$\sin$};
	\draw[dashed,color=black] (-3.14/2,-2)--(-3.14/2,2);
	\draw[dashed,color=black] (3.14/2,-2)--(3.14/2,2);
	\draw[dashed,color=black] (3.14/2+3.14,-2)--(3.14/2+3.14,2);
	\end{tikzpicture}\\	
\end{center}
Für die Abbildung der Hauptperiode gilt:\\
\begin{align*}
\tan:\;]\genfrac{}{}{}{2}{-\pi }{2},\genfrac{}{}{}{2}{\pi}{2}[\;\;\text{mit }\;\; \longmapsto \mathbb{R}\;\;x\longmapsto \frac{\sin(x)}{\cos(x)}
\end{align*}
Der Tangens ist als Komposition diffbarer Funktionen diffbar.\\
\begin{align*}
\tan'(x)=\bigg(\frac{\sin(x)}{\cos(x)}\bigg)'&=\frac{\cos(x)\cos(x)-\sin(x)(-\sin(x))}{\cos^2(x)}\\
&=\frac{\cos^2(x)+\sin^2(x)}{\cos^2(x)}\;\;|\text{\tiny $ \cos^2(x)+\sin^2(x)=1 $ }\\
&=\frac{1}{\cos^2(x)}\\
&=\frac{\cos^2(x)}{\cos^2(x)}+\frac{\sin^2(x)}{\cos^2(x)}=1+\tan^2(x)>0\\
\end{align*}
Nach dem Monotoniekriterium ist der Tangens streng monoton steigend.\\
\subsubsection{Der Kotangens}
\href{https://youtu.be/3N8L_ze4y2o}{Video 06}\\
Definition des Kotangens:\\
\begin{align*}
\cot(x+k\pi)=\cot(x)\;\;\text{für }\;\; x\neq k\pi, k\in \mathbb{Z}
\end{align*}
\begin{center}
	\begin{tikzpicture}[xscale=1,yscale=2]
	
	% Gitternetzlinien
	\draw[xstep=1cm, ystep=1cm, very thin, color=lightgray] (-4,2) grid (7,-2);
	% x-Achse
	\draw[->] (-4,0) -- (7,0) node[below,xshift=-1em] {$x$}; 
	%Zahlen auf x-Achse
	\foreach \x in {-3.14,-3.14/2}
	\draw[shift={(\x,0)},color=black] (0pt,2pt) -- (0pt,-2pt);
	\foreach \x in {3.14/2,3.14,4.71,6.28}
	\draw[shift={(\x,0)},color=black] (0pt,2pt) -- (0pt,-2pt);
	% y-Achse 
	\draw[->] (0,-2) -- (0,2) node[left,yshift=-1em] {$f(x)$};%node[above left]
	%Zahlen auf y-Achse
	\foreach \y in {-2,...,-1}
	\draw[shift={(0,\y)},color=black] (2pt,0pt) -- (-2pt,0pt) node[left] {\footnotesize $\y$};
	\foreach \y in {1,...,1}
	\draw[shift={(0,\y)},color=black] (2pt,0pt) -- (-2pt,0pt) node[left] {\footnotesize $\y$};
	%Ursprung
	\draw[color=black] (0pt,-10pt) node[right] {\footnotesize $0$};
	\draw[color=black] (3.14,-2pt) node[below] {\footnotesize $\pi$};
	\draw[color=black] (6.28,-2pt) node[below] {\footnotesize $2\pi$};
	\draw[color=black] (1.57,-2pt) node[below] {\footnotesize $\frac{\pi}{2}$};
	\draw[color=black] (-1.57,-2pt) node[below] {\footnotesize $-\frac{\pi}{2}$};
	\draw[color=black] (-3.14,-2pt) node[below] {\footnotesize $-\pi$};
	\draw[color=black] (3*3.14/2,-2pt) node[below] {\footnotesize $\frac{3}{2}\pi$};
	\draw[thick,color=black] plot[samples=200, domain=0.47:2.68] 
	(\x,{cot(\x r)}) node[below, xshift=-1em,yshift=15]{$\cot$};
	\draw[thick,color=black] plot[samples=200, domain=-2.68:-0.47] 
	(\x,{cot(\x r)}) node[below, xshift=-1em,yshift=15]{$\cot$};
	\draw[dashed,color=red] plot[samples=200, domain=0:3.14] 
	(\x,{cos(\x r)}) node[ xshift=-9em,yshift=10.5em]{$\cos$};
	\draw[dashed,color=blue] plot[samples=200, domain=0:3.14] 
	(\x,{sin(\x r)}) node[ xshift=-0.5em,yshift=3em]{$\sin$};
	\draw[dashed,color=black] (-3.14,-2)--(-3.14,2);
	\draw[dashed,color=black] (3.14,-2)--(3.14,2);
	%\draw[dashed,color=black] (3.14/2+3.14,-2)--(3.14/2+3.14,2);
	\end{tikzpicture}\\	
\end{center}
Für die Abbildung der Hauptperiode gilt:\\
\begin{align*}
\cot:\;]0,\pi[ \;\longmapsto \mathbb{R}\;\;\text{für }\;\;x\longmapsto\frac{\cos(x)}{\sin(x)}
\end{align*}
Der Kotangens ist als Komposition diffbarer Funktionen diffbar.\\
\begin{align*}
\cot'(x)=\bigg(\frac{\cos(x)}{\sin(x)}\bigg)'&=\frac{-\sin^2(x)-\cos^2(x)}{\sin^2(x)}\\
&=\frac{-(\sin^2(x)+\cos^2(x))}{\sin^2(x)}=-\frac{1}{\sin^2(x)}\\
&=-1-\cot^2(x)<0\\
\end{align*}
Nach dem Monotoniekriterium ist der Kotangens streng monoton fallend.\\
\subsubsection{Arcussinus und Arcuskosinus}
\href{https://youtu.be/WWMBFDeShK4}{Video 07}\\
$\blacktriangleright$ Umkehrfunktion vom Kosinus.\\
\begin{quote}
	
	\begin{center}
		
		
		\begin{tikzpicture}[xscale=1,yscale=2]
		
		% Gitternetzlinien
		\draw[xstep=1cm, ystep=1cm, very thin, color=lightgray] (-4,2) grid (7,-2);
		% x-Achse
		\draw[->] (-4,0) -- (7,0) node[below,xshift=-1em] {$x$}; 
		%Zahlen auf x-Achse
		\foreach \x in {-3.14,-3.14/2}
		\draw[shift={(\x,0)},color=black] (0pt,2pt) -- (0pt,-2pt);
		\foreach \x in {3.14/2,3.14,4.71,6.28}
		\draw[shift={(\x,0)},color=black] (0pt,2pt) -- (0pt,-2pt);
		% y-Achse 
		\draw[->] (0,-2) -- (0,2) node[left,yshift=-1em] {$f(x)$};%node[above left]
		%Zahlen auf y-Achse
		\foreach \y in {-2,...,-1}
		\draw[shift={(0,\y)},color=black] (2pt,0pt) -- (-2pt,0pt) node[left] {\footnotesize $\y$};
		\foreach \y in {1,...,1}
		\draw[shift={(0,\y)},color=black] (2pt,0pt) -- (-2pt,0pt) node[left] {\footnotesize $\y$};
		%Ursprung
		\draw[color=black] (0pt,-10pt) node[right] {\footnotesize $0$};
		\draw[color=black] (3.14,-2pt) node[below] {\footnotesize $\pi$};
		\draw[color=black] (6.28,-2pt) node[below] {\footnotesize $2\pi$};
		\draw[color=black] (1.57,-2pt) node[below] {\footnotesize $\frac{\pi}{2}$};
		\draw[color=black] (-1.57,-2pt) node[below] {\footnotesize $-\frac{\pi}{2}$};
		\draw[color=black] (-3.14,-2pt) node[below] {\footnotesize $-\pi$};
		\draw[color=black] (3*3.14/2,-2pt) node[below] {\footnotesize $\frac{3}{2}\pi$};
		
		\draw[color=blue] plot[samples=200, domain=0:3.14] 
		(\x,{cos(\x r)}) node[ xshift=-9em,yshift=10.5em]{$\cos$};
		\draw[color=blue,thick] (0pt,0pt) node {$[$};
		\draw[color=blue,thick] (3.14,0pt) node { $]$};
		\draw[color=blue,thick] (0,0pt)--(3.14,0);
		%\draw[dashed,color=black] (-3.14,-2)--(-3.14,2);
		%\draw[dashed,color=black] (3.14,-2)--(3.14,2);
		%\draw[dashed,color=black] (3.14/2+3.14,-2)--(3.14/2+3.14,2);
		\end{tikzpicture}\\	
		Umkehrbahr nur in diesen Intervall, wegen der Bijektivität \\
		$\cos:[0,\pi]\longmapsto[-1,1] $\\ 
	\end{center}
	\newpage
	\noindent
	\begin{align*}
	\arccos:=\cos^{-1}:[-1,1]\longmapsto[0,\pi]
	\end{align*}
	\begin{center}
		\begin{tikzpicture}[xscale=3, yscale=2]
		
		% Gitternetzlinien
		\draw[xstep=1cm, ystep=1cm, very thin, color=lightgray] (-2,4) grid (2,-1);
		% x-Achse
		\draw[->] (-2,0) -- (2,0) node[below,xshift=-1em] {$x$}; 
		%Zahlen auf x-Achse
		\foreach \x in {-1,...,-1}
		\draw[shift={(\x,0)},color=black] (0pt,2pt) -- (0pt,-2pt) node[below] {\footnotesize $\x$};
		\foreach \x in {1,...,1}
		\draw[shift={(\x,0)},color=black] (0pt,2pt) -- (0pt,-2pt) node[below] {\footnotesize $\x$};
		% y-Achse 
		\draw[->] (0,-1) -- (0,4) node[left,yshift=-1em] {$f(x)$};%node[above left]
		%Zahlen auf y-Achse
		\foreach \y in {-1,...,-1}
		\draw[shift={(0,\y)},color=black] (2pt,0pt) -- (-2pt,0pt) node[left] {\footnotesize $\y$};
		\foreach \y in {1,...,3}
		\draw[shift={(0,\y)},color=black] (2pt,0pt) -- (-2pt,0pt) node[left] {\footnotesize $\y$};
		%Ursprung
		\draw[color=black] (0pt,-10pt) node[right] {\footnotesize $0$};
		
		\draw[thick,color=blue] plot[samples=200, domain=-0.999:0.999] 
		(\x,{rad(acos(\x ))}) node[yshift=1, xshift=2em] {$\arccos$};
		
		\draw[color=black] (2pt,3.14) -- (-2pt,3.14) node[left] {\footnotesize $\pi$};
		
		\draw[color=black] (2pt,3.14/2) -- (-2pt,3.14/2) node[left] {\footnotesize $\frac{\pi}{2}$};
		
		\end{tikzpicture}\\
		\begin{align*}
		\arccos'(x)=\frac{-1}{\sqrt{1-x^2}}
		\end{align*}
		
	\end{center}
\end{quote}
$\blacktriangleright$ Umkehrfunktion vom Sinus.\\
\begin{quote}
	\begin{center}
		\begin{tikzpicture}[xscale=1,yscale=2]
		
		% Gitternetzlinien
		\draw[xstep=1cm, ystep=1cm, very thin, color=lightgray] (-4,2) grid (4,-2);
		% x-Achse
		\draw[->] (-4,0) -- (4,0) node[below,xshift=-1em] {$x$}; 
		%Zahlen auf x-Achse
		\foreach \x in {-3.14,3.14}
		\draw[shift={(\x,0)},color=black] (0pt,2pt) -- (0pt,-2pt);
		\foreach \x in {3.14/2,3.14}
		\draw[shift={(\x,0)},color=black] (0pt,2pt) -- (0pt,-2pt);
		% y-Achse 
		\draw[->] (0,-2) -- (0,2) node[left,yshift=-1em] {$f(x)$};%node[above left]
		%Zahlen auf y-Achse
		\foreach \y in {-2,...,-1}
		\draw[shift={(0,\y)},color=black] (2pt,0pt) -- (-2pt,0pt) node[left] {\footnotesize $\y$};
		\foreach \y in {1,...,1}
		\draw[shift={(0,\y)},color=black] (2pt,0pt) -- (-2pt,0pt) node[left] {\footnotesize $\y$};
		%Ursprung
		\draw[color=black] (0pt,-10pt) node[right] {\footnotesize $0$};
		\draw[color=black] (3.14,-2pt) node[below] {\footnotesize $\pi$};
		%\draw[color=black] (6.28,-2pt) node[below] {\footnotesize $2\pi$};
		\draw[color=black] (1.57,-2pt) node[below] {\footnotesize $\frac{\pi}{2}$};
		\draw[color=black] (-1.57,-2pt) node[below] {\footnotesize $-\frac{\pi}{2}$};
		\draw[color=black] (-3.14,-2pt) node[below] {\footnotesize $-\pi$};
		%\draw[color=black] (3*3.14/2,-2pt) node[below] {\footnotesize $\frac{3}{2}\pi$};
		
		\draw[color=blue] plot[samples=200, domain=-3.14/2:3.14/2] 
		(\x,{sin(\x r)}) node[ xshift=-7em,yshift=-8.5em]{$\sin$};
		\draw[color=blue,thick] (-3.14/2,0pt) node {$[$};
		\draw[color=blue,thick] (3.14/2,0pt) node { $]$};
		\draw[color=blue,thick] (-3.14/2,0)--(3.14/2,0);
		%\draw[dashed,color=black] (-3.14,-2)--(-3.14,2);
		%\draw[dashed,color=black] (3.14,-2)--(3.14,2);
		%\draw[dashed,color=black] (3.14/2+3.14,-2)--(3.14/2+3.14,2);
		\end{tikzpicture}\\	
		Umkehrbar nur in diesen Intervall, wegen der Bijektivität.\\
		$ \sin:[-\frac{\pi}{2},+\frac{\pi}{2}]\longmapsto[-1,1] $\\
	\end{center}
	\newpage
	\noindent
	\begin{align*}
	\arcsin:=\sin^{-1}:[-1,1]\longmapsto[-\genfrac{}{}{}{1.7}{\pi}{2 },\genfrac{}{}{}{1.7}{\pi }{2}]
	\end{align*}
	\begin{center}
		\begin{tikzpicture}[xscale=3, yscale=2]
		
		% Gitternetzlinien
		\draw[xstep=1cm, ystep=1cm, very thin, color=lightgray] (-2,2) grid (2,-2);
		% x-Achse
		\draw[->] (-2,0) -- (2,0) node[below,xshift=-1em] {$x$}; 
		%Zahlen auf x-Achse
		\foreach \x in {-1,...,-1}
		\draw[shift={(\x,0)},color=black] (0pt,2pt) -- (0pt,-2pt) node[below] {\footnotesize $\x$};
		\foreach \x in {1,...,1}
		\draw[shift={(\x,0)},color=black] (0pt,2pt) -- (0pt,-2pt) node[below] {\footnotesize $\x$};
		% y-Achse 
		\draw[->] (0,-1) -- (0,2) node[left,yshift=-1em] {$f(x)$};%node[above left]
		%Zahlen auf y-Achse
		\foreach \y in {-1,...,-1}
		\draw[shift={(0,\y)},color=black] (2pt,0pt) -- (-2pt,0pt) node[left] {\footnotesize $\y$};
		\foreach \y in {1,...,1}
		\draw[shift={(0,\y)},color=black] (2pt,0pt) -- (-2pt,0pt) node[left] {\footnotesize $\y$};
		%Ursprung
		\draw[color=black] (0pt,-10pt) node[right] {\footnotesize $0$};
		
		\draw[thick,color=blue] plot[samples=200, domain=-0.999:0.999] 
		(\x,{rad(asin(\x ))}) node[yshift=1, xshift=2em] {$\arcsin$};
		
		\draw[color=black] (2pt,-3.14/2) -- (-2pt,-3.14/2) node[left] {\footnotesize $-\frac{\pi}{2}$};
		
		\draw[color=black] (2pt,3.14/2) -- (-2pt,3.14/2) node[left] {\footnotesize $\frac{\pi}{2}$};
		
		\end{tikzpicture}\\
		\begin{align*}
		\arcsin'(x)=\frac{1}{\sqrt{1-x^2}}
		\end{align*}
		
	\end{center}	
\end{quote}
\subsubsection{Arcustangens und Arcuskotangens}
\href{https://youtu.be/D5xXaAdpNfQ}{Video 08}\\
$ \blacktriangleright $ Umkehrfunktion vom Tangens.\\

\begin{center}
	\begin{tikzpicture}[xscale=1,yscale=2]
	
	% Gitternetzlinien
	\draw[xstep=1cm, ystep=1cm, very thin, color=lightgray] (-4,2) grid (4,-2);
	% x-Achse
	\draw[->] (-4,0) -- (4,0) node[below,xshift=-1em] {$x$}; 
	%Zahlen auf x-Achse
	\foreach \x in {-3.14,-3.14/2}
	\draw[shift={(\x,0)},color=black] (0pt,2pt) -- (0pt,-2pt);
	\foreach \x in {3.14/2,3.14}
	\draw[shift={(\x,0)},color=black] (0pt,2pt) -- (0pt,-2pt);
	% y-Achse 
	\draw[->] (0,-2) -- (0,2) node[left,yshift=-1em] {$f(x)$};%node[above left]
	%Zahlen auf y-Achse
	\foreach \y in {-2,...,-1}
	\draw[shift={(0,\y)},color=black] (2pt,0pt) -- (-2pt,0pt) node[left] {\footnotesize $\y$};
	\foreach \y in {1,...,1}
	\draw[shift={(0,\y)},color=black] (2pt,0pt) -- (-2pt,0pt) node[left] {\footnotesize $\y$};
	%Ursprung
	\draw[color=black] (0pt,-10pt) node[right] {\footnotesize $0$};
	\draw[color=black] (3.14,-2pt) node[below] {\footnotesize $\pi$};
	%\draw[color=black] (6.28,-2pt) node[below] {\footnotesize $2\pi$};
	\draw[color=black] (1.57,-2pt) node[below] {\footnotesize $\frac{\pi}{2}$};
	\draw[color=black] (-1.57,-2pt) node[below] {\footnotesize $-\frac{\pi}{2}$};
	\draw[color=black] (-3.14,-2pt) node[below] {\footnotesize $-\pi$};
	%\draw[color=black] (3*3.14/2,-2pt) node[below] {\footnotesize $\frac{3}{2}\pi$};
	\draw[thick,color=blue] plot[samples=200, domain=-2.2/2:2.2/2] 
	(\x,{tan(\x r)}) node[below, xshift=-1em]{$\tan$};
	
	%\draw[dashed,color=black] (-3.14/2,-2)--(-3.14/2,2);
	%\draw[dashed,color=black] (3.14/2,-2)--(3.14/2,2);
	%\draw[dashed,color=black] (3.14/2+3.14,-2)--(3.14/2+3.14,2);
	\draw[color=blue,thick] (-3.14/2,0pt) node {$[$};
	\draw[color=blue,thick] (3.14/2,0pt) node { $]$};
	\draw[color=blue,thick] (-3.14/2,0)--(3.14/2,0);
	\end{tikzpicture}\\	
	Umkehrbar nur in diesen Intervall, wegen der Bijektivität.\\
	$ \tan:\;]-\frac{\pi}{2},+\frac{\pi}{2}[\;\longmapsto \mathbb{R} $\\
\end{center}
\newpage
\noindent
\begin{align*}
\arctan:=\tan^{-1}:\mathbb{R}\longmapsto\;]-\genfrac{}{}{}{1.7}{\pi}{2 },\genfrac{}{}{}{1.7}{\pi }{2}[
\end{align*}
\begin{center}
	\begin{tikzpicture}[xscale=1, yscale=2]
	
	% Gitternetzlinien
	\draw[xstep=1cm, ystep=1cm, very thin, color=lightgray] (-4,2) grid (4,-2);
	% x-Achse
	\draw[->] (-4,0) -- (4,0) node[below,xshift=-1em] {$x$}; 
	%Zahlen auf x-Achse
	\foreach \x in {-3,...,-1}
	\draw[shift={(\x,0)},color=black] (0pt,2pt) -- (0pt,-2pt) node[below] {\footnotesize $\x$};
	\foreach \x in {1,...,3}
	\draw[shift={(\x,0)},color=black] (0pt,2pt) -- (0pt,-2pt) node[below] {\footnotesize $\x$};
	% y-Achse 
	\draw[->] (0,-1) -- (0,2) node[left,yshift=-1em] {$f(x)$};%node[above left]
	%Zahlen auf y-Achse
	\foreach \y in {-1,...,-1}
	\draw[shift={(0,\y)},color=black] (2pt,0pt) -- (-2pt,0pt) node[left] {\footnotesize $\y$};
	\foreach \y in {1,...,1}
	\draw[shift={(0,\y)},color=black] (2pt,0pt) -- (-2pt,0pt) node[left] {\footnotesize $\y$};
	%Ursprung
	\draw[color=black] (0pt,-10pt) node[right] {\footnotesize $0$};
	
	\draw[thick,color=blue] plot[samples=200, domain=-4:4] 
	(\x,{rad(atan(\x ))}) node[yshift=1.5em, xshift=-2em] {$\arctan$};
	
	\draw[color=black] (2pt,-3.14/2) -- (-2pt,-3.14/2) node[left] {\footnotesize $-\frac{\pi}{2}$};
	
	\draw[color=black] (2pt,3.14/2) -- (-2pt,3.14/2) node[left] {\footnotesize $\frac{\pi}{2}$};
	\end{tikzpicture}\\
\end{center}
\begin{align*}
\arctan'(x)=\frac{1}{1+x^2}
\end{align*}
$ \blacktriangleright $ Umkehrfunktion vom Kotangens.\\
\begin{center}
	\begin{tikzpicture}[xscale=1,yscale=2]
	
	% Gitternetzlinien
	\draw[xstep=1cm, ystep=1cm, very thin, color=lightgray] (-4,2) grid (4,-2);
	% x-Achse
	\draw[->] (-4,0) -- (4,0) node[below,xshift=-1em] {$x$}; 
	%Zahlen auf x-Achse
	\foreach \x in {-3.14,3.14}
	\draw[shift={(\x,0)},color=black] (0pt,2pt) -- (0pt,-2pt);
	\foreach \x in {3.14/2,3.14}
	\draw[shift={(\x,0)},color=black] (0pt,2pt) -- (0pt,-2pt);
	% y-Achse 
	\draw[->] (0,-2) -- (0,2) node[left,yshift=-1em] {$f(x)$};%node[above left]
	%Zahlen auf y-Achse
	\foreach \y in {-2,...,-1}
	\draw[shift={(0,\y)},color=black] (2pt,0pt) -- (-2pt,0pt) node[left] {\footnotesize $\y$};
	\foreach \y in {1,...,1}
	\draw[shift={(0,\y)},color=black] (2pt,0pt) -- (-2pt,0pt) node[left] {\footnotesize $\y$};
	%Ursprung
	\draw[color=black] (0pt,-10pt) node[right] {\footnotesize $0$};
	\draw[color=black] (3.14,-2pt) node[below] {\footnotesize $\pi$};
	%\draw[color=black] (6.28,-2pt) node[below] {\footnotesize $2\pi$};
	\draw[color=black] (1.57,-2pt) node[below] {\footnotesize $\frac{\pi}{2}$};
	\draw[color=black] (-1.57,-2pt) node[below] {\footnotesize $-\frac{\pi}{2}$};
	\draw[color=black] (-3.14,-2pt) node[below] {\footnotesize $-\pi$};
	%\draw[color=black] (3*3.14/2,-2pt) node[below] {\footnotesize $\frac{3}{2}\pi$};
	\draw[thick,color=blue] plot[samples=200, domain=0.47:2.68] 
	(\x,{cot(\x r)}) node[below, xshift=-1em,yshift=15]{$\cot$};
	
	\draw[color=blue,thick] (0/2,0pt) node {$[$};
	\draw[color=blue,thick] (3.14,0pt) node { $]$};
	\draw[color=blue,thick] (0,0)--(3.14,0);
	%\draw[dashed,color=black] (-3.14,-2)--(-3.14,2);
	%\draw[dashed,color=black] (3.14,-2)--(3.14,2);
	%\draw[dashed,color=black] (3.14/2+3.14,-2)--(3.14/2+3.14,2);
	\end{tikzpicture}\\	
	Umkehrbar nur in diesen Intervall, wegen der Bijektivität.\\
	$ \cot:\;]0,\pi[\;\longmapsto\mathbb{R} $\\
\end{center}
\begin{align*}
\mathrm{arccot}:=\cot^{-1}:\mathbb{R}\longmapsto\;]0,\pi[
\end{align*}
\begin{align*}
\mathrm{arccot}'=-\frac{1}{1+x^2}
\end{align*}
\newpage
\noindent
\subsubsection{Aufgabe 1 zu trigonometrischen Funktionen}
\href{https://youtu.be/voAHLA8F08c}{Video 09}\\
Bestimme den maximalen Definitionsbereich und die Ableitung von:\\
\begin{align*}
f(x)=(\arccos(4-x))^5
\end{align*}
Es gilt:\\
\begin{align*}
&\arccos: [-1,1]\longmapsto [0,\pi]\\
&\arccos'(x)=-\frac{1}{\sqrt{1-x^2}}
\end{align*}
(i)\;\;Zum max. Definitionsbereich.\\
\begin{quote}
	\begin{align*}
	&f(x)=(\arccos(4-x))^5\\
	&\Rightarrow\\
	&\text{Für }\;\;-1\leq 4-x\leq 1\;\;\text{ist }\;\;f(x)\;\;\text{definiert. }
	\end{align*}
	D.h.:\\
	\begin{align*}
	&-1\leq 4-x\leq 1\\
	&\overset{(-4)}{\Leftrightarrow}\\
	&-5\leq -x\leq -3\\
	&\overset{(\cdot -1)}{\Leftrightarrow}\\
	&5\geq x\geq 3\\
	&\Rightarrow\mathbb{D}_f=[3,5]\\
	\end{align*}
\end{quote}
(ii)\;\;Bestimmung von $ f'(x) $.\\
\begin{quote}
	$ f(x) $ ist diffbar, weil es eine Komposition von diff. Funktionen ist.\\
	\begin{align*}
	f'(x)=(\arccos^5(4-x))'&=5\cdot \arccos^4(4-x)\cdot(\arccos(4-x))'\\
	&=5\cdot \arccos^4(4-x)\cdot\bigg(-\frac{1}{\sqrt{1-(4-x)^2}}\bigg)\cdot(4-x)'\\
	&=5\cdot \arccos^4(4-x)\cdot\frac{1}{\sqrt{1-(4-x)^2}}\\
	&=\frac{5\cdot \arccos^4(4-x)}{\sqrt{1-(4-x)^2}}\\
	\end{align*}
\end{quote}
\subsubsection{Aufgabe 2 zu trigonometrischen Funktionen}
\href{https://youtu.be/dyxaIIiA6r4}{Video 10}\\
Bestimme den maximalen Definitionsbereich von:\\
\begin{align*}
g(x)=\arcsin\bigg(\frac{1-x^2}{1+x^2}\bigg)
\end{align*}
Wenn möglich, bestimme $ g'(x) $.\\
Es gilt:\\
\begin{align*}
&\arcsin: [-1,1]\longmapsto [-\genfrac{}{}{}{1.7}{\pi }{2 },\genfrac{}{}{}{1.7}{\pi }{2 }]\\
&\arcsin'(x)=\frac{1}{\sqrt{1-x^2}}
\end{align*}
(i)\;\;Zum max. Definitionsbereich.\\
\begin{quote}
	\begin{align*}
	&g(x)=\arcsin\bigg(\frac{1-x^2}{1+x^2}\bigg)\\
	&\Rightarrow\\
	&-1\leq\frac{1-x^2}{1+x^2}\leq 1
	\end{align*}
	D.h:\\
	Für $ -1\leq\frac{1-x^2}{1+x^2} $ gilt:\\
	\begin{align*}
	&-1\leq\frac{1-x^2}{1+x^2}\\
	&\overset{(1+x^2>0)}{\Leftrightarrow}\\
	&-1(1+x^2)\leq 1-x^2\\
	&-1-x^2\leq 1-x^2\\
	&-1\leq 1\\
	&\Rightarrow \frac{1-x^2}{1+x^2}\geq -1
	\end{align*}
	Für $ 1\geq\frac{1-x^2}{1+x^2} $ gilt:\\
	\begin{align*}
	&\frac{1-x^2}{1+x^2}\leq 1\\
	&\overset{(1+x^2>0)}{\Leftrightarrow}\\
	&1-x^2\leq 1(1+x^2)\\
	&\overset{(x^2)}{\Leftrightarrow}\\
	&1\leq 1+2x^2\\
	&\Rightarrow \frac{1-x^2}{1+x^2}\leq 1
	\end{align*}	
	D.h: Maximaler Definitionsbereich ist $ \mathbb{D}_g =\mathbb{R}$\\
\end{quote}
(ii)\;\;Bestimmung von $ g'(x) $:\\
\begin{quote}
	$ g(x) $ ist diffbar, weil es eine Komposition von diff. Funktionen ist.\\
	\begin{align*}
	g'(x)=\bigg(\arcsin\bigg(\frac{1-x^2}{1+x^2}\bigg)\bigg)'&=\frac{1}{\sqrt{1-(\frac{1-x^2}{1+x^2})^2}}\cdot\bigg(\frac{1-x^2}{1+x^2}\bigg)'\\
	&=\frac{1}{\sqrt{1-(\frac{1-x^2}{1+x^2})^2}}\cdot \frac{-2x(1+x^2)-(1-x^2)2x}{(1+x^2)^2}\\
	&=\frac{1}{\sqrt{1-(\frac{1-x^2}{1+x^2})^2}}\cdot\frac{-2x-2x^2-2x+2x^2}{(1+x^2)^2}\\
	&=\frac{-4x}{\sqrt{1-(\frac{1-x^2}{1+x^2})^2}}\cdot \frac{1}{(1+x^2)^2}
	\end{align*}
\end{quote}
\subsubsection{Aufgabe 3 zu trigonometrischen Funktionen}
\href{https://youtu.be/jgmWhGejYJI}{Video 11}\\
Beweise: $ \arccos(x) +\arcsin(x)=\frac{\pi}{2} $ für $ ]-1,1[ $ .\\
\\
Definiere:\\
\begin{align*}
f(x)=\arccos(x) +\arcsin(x)
\end{align*}
Es gilt:\\
\begin{align*}
f'(x)=\frac{1}{\sqrt{1-x^2}}-\frac{1}{\sqrt{1-x^2}}=0 \;\;\text{für }\;\;x\in\;\;]-1,1[
\end{align*}
Wegen dem Konstanzkriterium gilt, dass $ f(x) $ konstant ist.\\
D.h.:\\
\begin{align*}
&f(x)=k\;\;\text{ für}\;\;k=\in\mathbb{R}\\
&f(0)=\arccos(0) +\arcsin(0)\\
&f(0)=0+\frac{\pi}{2}\\
&\Rightarrow \arccos(x) +\arcsin(x)=\frac{\pi}{2}
\end{align*}
\begin{center}
	\begin{tikzpicture}[xscale=3, yscale=2]
	
	% Gitternetzlinien
	\draw[xstep=1cm, ystep=1cm, very thin, color=lightgray] (-2,4) grid (2,-2);
	% x-Achse
	\draw[->] (-2,0) -- (2,0) node[below,xshift=-1em] {$x$}; 
	%Zahlen auf x-Achse
	\foreach \x in {-1,...,-1}
	\draw[shift={(\x,0)},color=black] (0pt,2pt) -- (0pt,-2pt) node[below] {\footnotesize $\x$};
	\foreach \x in {1,...,1}
	\draw[shift={(\x,0)},color=black] (0pt,2pt) -- (0pt,-2pt) node[below] {\footnotesize $\x$};
	% y-Achse 
	\draw[->] (0,-2) -- (0,4) node[left,yshift=-1em] {$f(x)$};%node[above left]
	%Zahlen auf y-Achse
	\foreach \y in {-1,...,-1}
	\draw[shift={(0,\y)},color=black] (2pt,0pt) -- (-2pt,0pt) node[left] {\footnotesize $\y$};
	\foreach \y in {1,...,3}
	\draw[shift={(0,\y)},color=black] (2pt,0pt) -- (-2pt,0pt) node[left] {\footnotesize $\y$};
	%Ursprung
	\draw[color=black] (0pt,-10pt) node[right] {\footnotesize $0$};
	
	\draw[thick,color=blue] plot[samples=200, domain=-0.999:0.999] 
	(\x,{rad(acos(\x ))}) node[yshift=1, xshift=2.5em] {$\arccos(x)$};
	\draw[thick,color=blue] plot[samples=200, domain=-0.999:0.999] 
	(\x,{rad(asin(\x ))}) node[yshift=1, xshift=2.5em] {$\arcsin(x)$};
	\draw[color=black] (2pt,3.14) -- (-2pt,3.14) node[left] {\footnotesize $\pi$};
	
	\draw[color=black] (2pt,3.14/2) -- (-2pt,3.14/2) node[left] {\footnotesize $\frac{\pi}{2}$};
	%		\draw[thick,color=blue] plot[samples=200, domain=-2:2] 
	%	(\x,{sinh(\x )}) node[yshift=1, xshift=2.5em] {$\arcsin(x)$};
	%	\draw[color=black] (2pt,3.14) -- (-2pt,3.14) node[left] {\footnotesize $\pi$};
	\end{tikzpicture}\\
	
	
\end{center}
\subsection{Sinus Hyperbolicus und Kosinus Hyperbolicus}
\href{https://youtu.be/FTkh0mVQCD0}{Video 12}\\
\underline{Analogie zum Cosinus und Sinus:}\\
\\
Es gilt mit eulerischer Formel: $ \mathrm{e}^{\mathrm{i}x}=\cos(x)+\mathrm{i}\cdot \sin(x) $ :\\
\begin{align*}
\frac{\mathrm{e}^{\mathrm{i}x}+\mathrm{e}^{-\mathrm{i}x}}{2}&=\frac{1}{2}\bigg[(\cos\big(x)+\mathrm{i}\cdot \sin(x)\big)+\big(\cos(-x)+\mathrm{i}\cdot \sin(-x)\big)\bigg]\\
&=\frac{1}{2}\big(\cos(x)+\mathrm{i}\cdot \sin(x)+\cos(x)-\mathrm{i}\cdot \sin(x)\big)\\
&=\frac{1}{2}(2\cdot \cos(x))\\
\frac{\mathrm{e}^{\mathrm{i}x}+\mathrm{e}^{-\mathrm{i}x}}{2}&=\cos(x)\\
\end{align*}
Analog gilt:\\
\begin{align*}
\frac{\mathrm{e}^{\mathrm{i}x}-\mathrm{e}^{-\mathrm{i}x}}{2\mathrm{i}}&=\sin(x)\\
\end{align*}
\underline {Kosinius Hyperbolikus}\\
\begin{align*}
\cosh(x)=\frac{\mathrm{e}^x+\mathrm{e}^{-x}}{2}\;\;\text{mit: } \cosh:\mathbb{R}\longmapsto[1,\infty]\\
\end{align*}
\begin{align*}
\cosh'(x)&=\bigg(\frac{\mathrm{e}^x+\mathrm{e}^{-x}}{2}\bigg)'=\frac{1}{2}\bigg[(\mathrm{e}^x)'+(\mathrm{e}^{-x})'\bigg]\\
\cosh'(x)&=\frac{\mathrm{e}^x-\mathrm{e}^{-x}}{2}=\sinh(x)
\end{align*}
Die Funktion $f(x) =\cosh(x) $ ist gerade (symmetrisch zur y-Achse),\\ weil gilt: $ f(x)=f(-x) $.\\
\begin{align*}
\cosh(-x)=\frac{\mathrm{e}^{-x}+\mathrm{e}^{-(-x)}}{2}=\frac{\mathrm{e}^x+\mathrm{e}^{-x}}{2}=\cosh(x)\\
\end{align*}
\begin{center}
	\begin{tikzpicture}[xscale=1, yscale=1]
	
	% Gitternetzlinien
	\draw[xstep=1cm, ystep=1cm, very thin, color=lightgray] (-4,4) grid (4,-2);
	% x-Achse
	\draw[->] (-4,0) -- (4,0) node[below,xshift=-1em] {$x$}; 
	%Zahlen auf x-Achse
	\foreach \x in {-3,...,-1}
	\draw[shift={(\x,0)},color=black] (0pt,2pt) -- (0pt,-2pt) node[below] {\footnotesize $\x$};
	\foreach \x in {1,...,3}
	\draw[shift={(\x,0)},color=black] (0pt,2pt) -- (0pt,-2pt) node[below] {\footnotesize $\x$};
	% y-Achse 
	\draw[->] (0,-2) -- (0,4) node[left,yshift=-1em] {$f(x)$};%node[above left]
	%Zahlen auf y-Achse
	\foreach \y in {-1,...,-1}
	\draw[shift={(0,\y)},color=black] (2pt,0pt) -- (-2pt,0pt) node[left] {\footnotesize $\y$};
	\foreach \y in {1,...,3}
	\draw[shift={(0,\y)},color=black] (2pt,0pt) -- (-2pt,0pt) node[left] {\footnotesize $\y$};
	%Ursprung
	
	
	
	\draw[thick,color=blue] plot[samples=200, domain=-2.05:2.05] 
	(\x,{cosh(\x )}) node[yshift=-2em, xshift=2em] {$\cosh(x)$};
	
	
	\end{tikzpicture}\\
	
\end{center}
\underline{Sinus Hyperbolikus}
\begin{align*}
\sinh(x)=\frac{\mathrm{e}^x-\mathrm{e}^{-x}}{2}\;\;\text{mit: } \sinh:\mathbb{R}\longmapsto\mathbb{R}\\
\end{align*}
\begin{align*}
\sinh'(x)&=\bigg(\frac{\mathrm{e}^x-\mathrm{e}^{-x}}{2}\bigg)'=\frac{1}{2}\bigg[(\mathrm{e}^x)'-(\mathrm{e}^{-x})'\bigg]\\
\sinh'(x)&=\frac{\mathrm{e}^x+\mathrm{e}^{-x}}{2}=\cosh(x)
\end{align*}
Die Funktion $f(x) =\sinh(x) $ ist ungerade (Punktsymmetrisch zum Nullpunkt),\\ weil gilt: $ f(-x)=-f(x) $.\\
\begin{align*}
\sinh(-x)=\frac{\mathrm{e}^{-x}-\mathrm{e}^x}{2}=\frac{-\mathrm{e}^x+\mathrm{e}^{-x}}{2}=-\bigg(\frac{\mathrm{e}^x-\mathrm{e}^{-x}}{2}\bigg)=-\sinh(x)
\end{align*}
\begin{center}
	\begin{tikzpicture}[xscale=1, yscale=1]
	
	% Gitternetzlinien
	\draw[xstep=1cm, ystep=1cm, very thin, color=lightgray] (-4,4) grid (4,-4);
	% x-Achse
	\draw[->] (-4,0) -- (4,0) node[below,xshift=-1em] {$x$}; 
	%Zahlen auf x-Achse
	\foreach \x in {-3,...,-1}
	\draw[shift={(\x,0)},color=black] (0pt,2pt) -- (0pt,-2pt) node[below] {\footnotesize $\x$};
	\foreach \x in {1,...,3}
	\draw[shift={(\x,0)},color=black] (0pt,2pt) -- (0pt,-2pt) node[below] {\footnotesize $\x$};
	% y-Achse 
	\draw[->] (0,-4) -- (0,4) node[left,yshift=-1em] {$f(x)$};%node[above left]
	%Zahlen auf y-Achse
	\foreach \y in {-3,...,-1}
	\draw[shift={(0,\y)},color=black] (2pt,0pt) -- (-2pt,0pt) node[left] {\footnotesize $\y$};
	\foreach \y in {1,...,3}
	\draw[shift={(0,\y)},color=black] (2pt,0pt) -- (-2pt,0pt) node[left] {\footnotesize $\y$};
	%Ursprung
	
	
	
	\draw[thick,color=blue] plot[samples=200, domain=-2.05:2.05] 
	(\x,{sinh(\x )}) node[yshift=-2em, xshift=2em] {$\sinh(x)$};
	
	
	
	\end{tikzpicture}\\
	
\end{center}
\underline{Hyperbelgleichung}
\begin{align*}
\cosh^2(x)-\sinh^2(x)=1
\end{align*}
Beweis:\\
\begin{align*}
\cosh^2(x)-\sinh^2(x)&=\bigg(\frac{\mathrm{e}^x+\mathrm{e}^{-x}}{2}\bigg)^2-\bigg(\frac{\mathrm{e}^x-\mathrm{e}^{-x}}{2}\bigg)^2\\
&=\bigg(\frac{1}{2}\bigg)^2\bigg[(\mathrm{e}^x+\mathrm{e}^{-x})^2-(\mathrm{e}^x-\mathrm{e}^{-x})^2\bigg]\\
&=\frac{1}{4}\cdot \bigg[\mathrm{e}^{2x}+\underbrace {2\mathrm{e}^x\mathrm{e}^{-x}}_{=2}+\mathrm{e}^{-2x}-(\mathrm{e}^{2x}-\underbrace {2\mathrm{e}^x\mathrm{e}^{-x}}_{=-2}+\mathrm{e}^{-2x})\bigg]\\
&=\frac{1}{4}(2+2)=1_{_{\Square}}
\end{align*}
\underline{Einheitshyperbel}\\
Für die gilt:\\
\begin{align*}
\{(x,y)|x^2-y^2=1\}
\end{align*}
Der Kosinus Hyperbolikus und der Sinus Hyperolikus sind auf der Einheitshyperbel.\\
\begin{center}
	
	
	\begin{tikzpicture}
	
	\begin{axis}[grid = major,xticklabels={0,0,1,2,3,{x}},yticklabels={0,-2,-1,0,1,{f(x)}},
	%ylabel near ticks,
	%xlabel = {x},
	%xlabel = {x},
	%ylabel = {y},
	xmin=0,xmax=5,
	ymin=-5,ymax=5]
	\node[blue,right] at (axis cs:3,2) {\footnotesize Hyperbel };
	\addplot [blue,thick,domain=-2:2] ({cosh(x)}, {sinh(x)});
	\addplot[red,dashed] expression {x};
	\addplot[red,dashed] expression {-x};
	\addplot[black] expression {0};
	\draw[blue,thick] ({axis cs:6,1}) node[above,right] {Formel 2};
	
	\end{axis}
	\end{tikzpicture}
\end{center}
\underline{Einheitskreis}\\
Für den gilt wegen:
\begin{align*}
\cos^2(x)+\sin^2(x)=1 \\
\end{align*}
\begin{align*}
\{(x,y)|x^2+y^2=1\}
\end{align*}
\begin{center}
	
	
	\begin{tikzpicture}
	
	\begin{axis}[grid = major,
	%ylabel near ticks,
	%xlabel = {x},
	%xlabel = {x},
	%ylabel = {y},
	xmin=-2,xmax=2,
	ymin=-2,ymax=2,xticklabels={0,0,-1,0,{x}},yticklabels={0,-2,-1,0,1,{f(x)}},]
	\node[blue,left of=eft] at (axis cs:1,1.2) {\footnotesize Einheitskreis };
	
	
	
	\draw[blue, thick] (axis cs:0,0) circle (100);
	
	
	\end{axis}
	\end{tikzpicture}
\end{center}
Der Sinus und der Cosinus sind auf dem Einheitskreis.\\
Bemerkung:\\
Der Sinus Hyperbolikus und der Kosinus Hyperbolikus sind der ungerade und gerade Anteil der Exponentialfunktion.\\
\begin{align*}
\mathrm{e}^x=\cosh(x)+\sinh(x)=\frac{\mathrm{e}^{x}}{2}+\frac{\mathrm{e}^{-x}}{2}+\frac{\mathrm{e}^{x}}{2}-\frac{\mathrm{e}^{-x}}{2}=\mathrm{e}^x\\
\end{align*}
\subsubsection{Tangens Hyperbolicus}
\href{https://youtu.be/cUg7vCA_ggw}{Video 13}\\
Es gilt für den Tangens:\\
\begin{align*}
\tan(x)=\frac{\sin(x)}{\cos(x)}
\end{align*}
Analog gilt für den Tangens Hyperbolikus:\\
\begin{align*}
\tanh(x)=\frac{\sinh(x)}{\cosh(x)}=\frac{\mathrm{e}^x-\mathrm{e}^{-x}}{\mathrm{e}^x+\mathrm{e}^{-x}}\;\;\text{mit }\;\;\tanh:\mathbb{R}\longmapsto\; ]-1,1[
\end{align*}
\newpage
\noindent
Warum gilt: $\tanh:\mathbb{R}\longmapsto\; ]-1,1[  $ ?
\begin{quote}
	\textbullet\;\;Es gilt: \\
	\begin{align*}
	&\frac{\mathrm{e}^x-\mathrm{e}^{-x}}{\mathrm{e}^x+\mathrm{e}^{-x}}\overset{!}{<}1\\
	&\mathrm{e}^x-\mathrm{e}^{-x}<\mathrm{e}^x+\mathrm{e}^{-x}\boxed{}\;\;\text{ $ \mathrm{e}^x $ ist immer positiv}\\
	&\Rightarrow\\
	&\mathrm{e}^x-\mathrm{e}^{-x}\neq\mathrm{e}^x+\mathrm{e}^{-x}\\
	&\Rightarrow\\
	&\frac{\mathrm{e}^x-\mathrm{e}^{-x}}{\mathrm{e}^x+\mathrm{e}^{-x}}<1
	\end{align*}
	\textbullet \;\;Ebenso gilt:\\
	\begin{align*}
	&\frac{\mathrm{e}^x-\mathrm{e}^{-x}}{\mathrm{e}^x+\mathrm{e}^{-x}}\overset{!}{>}-1\\
	&\mathrm{e}^x-\mathrm{e}^{-x}>-1(\mathrm{e}^x+\mathrm{e}^{-x})\\
	&\mathrm{e}^x-\mathrm{e}^{-x}>-\mathrm{e}^x-\mathrm{e}^{-x}\boxed{}\\
	&\Rightarrow\\
	&\mathrm{e}^x-\mathrm{e}^{-x}\neq-1(\mathrm{e}^x+\mathrm{e}^{-x})\\
	&\Rightarrow\\
	&\frac{\mathrm{e}^x-\mathrm{e}^{-x}}{\mathrm{e}^x+\mathrm{e}^{-x}}>-1
	\end{align*}
	\begin{align*}
	\Rightarrow \tanh:\mathbb{R}\longmapsto\; ]-1,1[
	\end{align*}
	
	
	
\end{quote}
Ableitung vom Tangens Hyperbolikus:\\
\begin{align*}
\tanh'(x)=\bigg(\tanh(x)=\frac{\sinh(x)}{\cosh(x)}\bigg)'=\frac{\cosh^2(x)-\sinh^2(x)}{\cosh^2(x)}=1-\frac{\sinh^2(x)}{\cosh^2(x)}\\
\end{align*}
\begin{quote}
	Ebenso gilt:\\
	\begin{align*}
	\cosh^2(x)-\sinh^2(x)=1
	\end{align*}
	Also gilt:\\
	\begin{align*}
	\tanh'(x)=\frac{1}{\cosh^2(x)}=1-\frac{\sinh^2(x)}{\cosh^2(x)}=1-\tanh^2(x)\\
	\end{align*}
\end{quote}

\begin{center}
	\begin{tikzpicture}[xscale=3, yscale=2]
	
	% Gitternetzlinien
	\draw[xstep=1cm, ystep=1cm, very thin, color=lightgray] (-2,2) grid (2,-2);
	% x-Achse
	\draw[->] (-2,0) -- (2,0) node[below,xshift=-1em] {$x$}; 
	%Zahlen auf x-Achse
	\foreach \x in {-1,...,-1}
	\draw[shift={(\x,0)},color=black] (0pt,2pt) -- (0pt,-2pt) node[below] {\footnotesize $\x$};
	\foreach \x in {1,...,1}
	\draw[shift={(\x,0)},color=black] (0pt,2pt) -- (0pt,-2pt) node[below] {\footnotesize $\x$};
	% y-Achse 
	\draw[->] (0,-2) -- (0,2) node[left,yshift=-1em] {$f(x)$};%node[above left]
	%Zahlen auf y-Achse
	\foreach \y in {-1,...,-1}
	\draw[shift={(0,\y)},color=black] (2pt,0pt) -- (-2pt,0pt) node[left] {\footnotesize $\y$};
	\foreach \y in {1,...,2}
	\draw[shift={(0,\y)},color=black] (2pt,0pt) -- (-2pt,0pt) node[left] {\footnotesize $\y$};
	%Ursprung
	\draw[color=black] (0pt,-10pt) node[right] {\footnotesize $0$};
	\draw[thick,color=blue] plot[samples=200, domain=-2.0:2.0] 
	(\x,{tanh(\x )}) node[yshift=-2em, xshift=-2em] {$\tanh(x)$};
	\draw[color=blue,dashed] (-2,1)--(2,1);
	\draw[color=blue,dashed] (-2,-1)--(2,-1);
	%		\draw[thick,color=blue] plot[samples=200, domain=-2:2] 
	%	(\x,{sinh(\x )}) node[yshift=1, xshift=2.5em] {$\arcsin(x)$};
	%	\draw[color=black] (2pt,3.14) -- (-2pt,3.14) node[left] {\footnotesize $\pi$};
	\end{tikzpicture}\\
	
	
\end{center}
Der Tangens Hyperbolikus ist streng monoton wachsend wegen $ \tanh'(x)=\frac{1}{\cosh^2(x)}>0 $.\\
\subsubsection{Kotangens Hyperbolicus}
\href{https://youtu.be/4mRuIkIRMCY}{Video 14}\\
Es gilt für den Kotangens:\\
\begin{align*}
\cot(x)=\frac{\cos(x)}{\sin(x)}
\end{align*}
Analog gilt für den Kotangens Hyperbolikus:\\
\begin{align*}
\coth(x)=\frac{\cosh(x)}{\sinh(x)}=\frac{\mathrm{e}^x+\mathrm{e}^{-x}}{\mathrm{e}^x-\mathrm{e}^{-x}}\;\;\text{mit }\;\;\coth:\mathbb{R}\backslash\{0\}\longmapsto \mathbb{R}\backslash \{-1,1\}
\end{align*} 
Warum gilt: $\coth:\mathbb{R}\backslash\{0\}\longmapsto \mathbb{R}\backslash \{-1,1\}  $ ?\\
\begin{quote}
	\textbullet \;\;$ \mathbb{R} \backslash\{0\}$ weil $ \mathrm{e}^x-\mathrm{e}^{-x}\neq 0 $ sein muß.\\
	\newpage
	\noindent
	\textbullet \;\; Es gilt:\\
	\begin{align*}
	&|\mathrm{e}^x-\mathrm{e}^{-x}|<|\mathrm{e}^x +\mathrm{e}^{-x}|\\
	&\Leftrightarrow\\
	&1<\bigg|\frac{\mathrm{e}^x +\mathrm{e}^{-x}}{\mathrm{e}^x-\mathrm{e}^{-x}}\bigg|\;\;\;\;|\text{\tiny $\mathrm{e}^x +\mathrm{e}^{-x}  $  immer positv }\\
	&1<\frac{\mathrm{e}^x +\mathrm{e}^{-x}}{|\mathrm{e}^x-\mathrm{e}^{-x}|}\\
	&\text{Für }\;\;\mathrm{e}^x-\mathrm{e}^{-x}<0\;\;\text{gilt: }\\
	&1<\frac{\mathrm{e}^x+\mathrm{e}^{-x}}{-(\mathrm{e}^x-\mathrm{e}^{-x})}\\
	&-1>\frac{\mathrm{e}^x+\mathrm{e}^{-x}}{\mathrm{e}^x-\mathrm{e}^{-x}}=\coth(x)\\
	&\text{Für }\;\;\mathrm{e}^x-\mathrm{e}^{-x}>0\;\;\text{gilt: }\\
	&1<\frac{\mathrm{e}^x+\mathrm{e}^{-x}}{\mathrm{e}^x-\mathrm{e}^{-x}}=\coth(x)\\
	&\Rightarrow\\
	&\coth(x)>1 \wedge \coth<-1\\
	&\Rightarrow\\
	&\coth:\mathbb{R}\backslash\{0\}\longmapsto \mathbb{R}\backslash [-1,1 ]
	\end{align*}
\end{quote}
Ableitung vom Kotangens Hyperbolikus:
\begin{align*}
\coth'(x)&=\bigg(\frac{\coth(x)}{\sinh(x)}\bigg)'=\frac{\sinh^2(x)-\cosh^2(x)}{\sinh^2(x)}=1-\coth^2(x)\\
&=\frac{-(-\sinh^2(x)+\cosh^2(x))}{\sinh^2(x)}\\
&=\frac{-(\cosh^2(x)-\sinh^2(x))}{\sinh^2(x)}=-\frac{1}{\sinh^2(x)}<0
\end{align*}
D.h. der Kotangens Hyperbolikus ist streng monoton fallend in den Intervallen $ ]-\infty,0[ $ und  $ ]0,\infty[ $. \\
Für die Monotoniebetrachtung muß jedes einzelne Intervall betrachtet werden.\\
\begin{center}
	\begin{tikzpicture}[xscale=1.5, yscale=1.5]
	
	% Gitternetzlinien
	\draw[xstep=1cm, ystep=1cm, very thin, color=lightgray] (-4,4) grid (4,-4);
	% x-Achse
	\draw[->] (-4,0) -- (4,0) node[below,xshift=-1em] {$x$}; 
	%Zahlen auf x-Achse
	\foreach \x in {-3,...,-1}
	\draw[shift={(\x,0)},color=black] (0pt,2pt) -- (0pt,-2pt) node[below] {\footnotesize $\x$};
	\foreach \x in {1,...,3}
	\draw[shift={(\x,0)},color=black] (0pt,2pt) -- (0pt,-2pt) node[below] {\footnotesize $\x$};
	% y-Achse 
	\draw[->] (0,-4) -- (0,4) node[left,yshift=-1em] {$f(x)$};%node[above left]
	%Zahlen auf y-Achse
	\foreach \y in {-3,...,-1}
	\draw[shift={(0,\y)},color=black] (2pt,0pt) -- (-2pt,0pt) node[left] {\footnotesize $\y$};
	\foreach \y in {1,...,3}
	\draw[shift={(0,\y)},color=black] (2pt,0pt) -- (-2pt,0pt) node[left] {\footnotesize $\y$};
	%Ursprung
	\draw[color=black] (0pt,-10pt) node[right] {\footnotesize $0$};
	\draw[thick,color=blue] plot[samples=200, domain=-4.0:-.26] 
	(\x,{cosh(\x )/sinh(\x)}) node[yshift=2em, xshift=-2em] {$\coth(x)$};
	\draw[thick,color=blue] plot[samples=200, domain=.26:4] 
	(\x,{cosh(\x )/sinh(\x)}) node[yshift=-2em, xshift=-2em] {$\coth(x)$};
	
	\draw[color=blue,dashed] (-4,1)--(4,1);
	\draw[color=blue,dashed] (-4,-1)--(4,-1);
	
	
	%		\draw[thick,color=blue] plot[samples=200, domain=-2:2] 
	%	(\x,{sinh(\x )}) node[yshift=1, xshift=2.5em] {$\arcsin(x)$};
	%	\draw[color=black] (2pt,3.14) -- (-2pt,3.14) node[left] {\footnotesize $\pi$};
	\end{tikzpicture}\\
	
	
\end{center}
Der Kotangens Hyperbolikus ist punktsymmetrisch.\\
Es gilt für Punktsymmetrie:\\
\begin{align*}
f(-x)=-f(x)\\
\end{align*}
\begin{align*}
\coth(-x)=\frac{\cosh(-x)}{\sinh(-x)}=\frac{\cosh(x)}{-\sinh(x)}=-\coth(x)
\end{align*}
\subsubsection{Areasinus Hyperbolicus und Areakosinus Hyperbolicus}
\href{https://youtu.be/LC1HLiQZvc8}{Video 15}\\
Die Area-Funktionen sind die Umkehrfunktionen der Hyperbelfunktionen.\\
$\blacktriangleright$ Umkehrfunktion vom Kosinus Hyperbolikus.
\begin{center}
	\begin{tikzpicture}[xscale=1, yscale=1]
	
	% Gitternetzlinien
	\draw[xstep=1cm, ystep=1cm, very thin, color=lightgray] (-4,4) grid (4,-2);
	% x-Achse
	\draw[->] (-4,0) -- (4,0) node[below,xshift=-1em] {$x$}; 
	%Zahlen auf x-Achse
	\foreach \x in {-3,...,-1}
	\draw[shift={(\x,0)},color=black] (0pt,2pt) -- (0pt,-2pt) node[below] {\footnotesize $\x$};
	\foreach \x in {1,...,3}
	\draw[shift={(\x,0)},color=black] (0pt,2pt) -- (0pt,-2pt) node[below] {\footnotesize $\x$};
	% y-Achse 
	\draw[->] (0,-2) -- (0,4) node[left,yshift=-1em] {$f(x)$};%node[above left]
	%Zahlen auf y-Achse
	\foreach \y in {-1,...,-1}
	\draw[shift={(0,\y)},color=black] (2pt,0pt) -- (-2pt,0pt) node[left] {\footnotesize $\y$};
	\foreach \y in {1,...,3}
	\draw[shift={(0,\y)},color=black] (2pt,0pt) -- (-2pt,0pt) node[left] {\footnotesize $\y$};
	%Ursprung
	
	
	
	\draw[thick,color=blue] plot[samples=200, domain=0:2.05] 
	(\x,{cosh(\x )}) node[yshift=-2em, xshift=2em] {$\cosh(x)$};
	\draw[color=blue,thick] (0pt,0pt) node {$[$};
	%\draw[color=blue,thick] (3.14,0pt) node { $]$};
	\draw[color=blue,thick] (0,0pt)--(4,0);
	
	\end{tikzpicture}\\
	Umkehrbahr nur in diesen Intervall, wegen der Bijektivität \\
	$\cosh:[0,\infty]\longmapsto[1,\infty] $\\ 
\end{center}
\begin{align*}
&\mathrm{arcosh}: [1,\infty[ \;\longmapsto [0,\infty[\\
&x \longmapsto \cosh^{-1}(x)=\mathrm{arcosh(x)}
\end{align*}
\begin{center}
	\begin{tikzpicture}[xscale=1, yscale=1]
	
	% Gitternetzlinien
	\draw[xstep=1cm, ystep=1cm, very thin, color=lightgray] (-4,4) grid (4,-2);
	% x-Achse
	\draw[->] (-4,0) -- (4,0) node[below,xshift=-1em] {$x$}; 
	%Zahlen auf x-Achse
	\foreach \x in {-3,...,-1}
	\draw[shift={(\x,0)},color=black] (0pt,2pt) -- (0pt,-2pt) node[below] {\footnotesize $\x$};
	\foreach \x in {1,...,3}
	\draw[shift={(\x,0)},color=black] (0pt,2pt) -- (0pt,-2pt) node[below] {\footnotesize $\x$};
	% y-Achse 
	\draw[->] (0,-2) -- (0,4) node[left,yshift=-1em] {$f(x)$};%node[above left]
	%Zahlen auf y-Achse
	\foreach \y in {-1,...,-1}
	\draw[shift={(0,\y)},color=black] (2pt,0pt) -- (-2pt,0pt) node[left] {\footnotesize $\y$};
	\foreach \y in {1,...,3}
	\draw[shift={(0,\y)},color=black] (2pt,0pt) -- (-2pt,0pt) node[left] {\footnotesize $\y$};
	%Ursprung
	\draw[thick,color=blue] plot[samples=200, domain=1.0:4] 
	(\x,{(ln(\x+sqrt(\x*\x-1) ))}) node[yshift=2, xshift=-3em] {$\mathrm{arcosh}(x)$};
	\fill[blue] (1,0) circle (0.045);
	
	%\draw[thick,color=blue] plot[samples=200, domain=0:2.05] 
	%	(\x,{cosh(\x )}) node[yshift=-2em, xshift=2em] {$\cosh(x)$};
	%	\draw[color=blue,thick] (0pt,0pt) node {$[$};
	%\draw[color=blue,thick] (3.14,0pt) node { $]$};
	%	\draw[color=blue,thick] (0,0pt)--(4,0);
	
	\end{tikzpicture}\\
	Umkehrbahr nur in diesen Intervall, wegen der Bijektivität \\
	$\cosh:[0,\infty]\longmapsto[1,\infty] $\\ 
\end{center}
Ableitung von Areakosinus Hyperbolikus.\\
\begin{align*}
\mathrm{arcosh}'(x)=\frac{1}{\sqrt{x^2-1}}
\end{align*}
Analog zum Arcuskosinus.\\
\begin{align*}
\arccos'(x)=\frac{-1}{\sqrt{1-x^2}}
\end{align*}
$\blacktriangleright$ Umkehrfunktion vom Sinus Hyperbolikus.
\begin{center}
	\begin{tikzpicture}[xscale=1, yscale=1]
	
	% Gitternetzlinien
	\draw[xstep=1cm, ystep=1cm, very thin, color=lightgray] (-4,4) grid (4,-4);
	% x-Achse
	\draw[->] (-4,0) -- (4,0) node[below,xshift=-1em] {$x$}; 
	%Zahlen auf x-Achse
	\foreach \x in {-3,...,-1}
	\draw[shift={(\x,0)},color=black] (0pt,2pt) -- (0pt,-2pt) node[below] {\footnotesize $\x$};
	\foreach \x in {1,...,3}
	\draw[shift={(\x,0)},color=black] (0pt,2pt) -- (0pt,-2pt) node[below] {\footnotesize $\x$};
	% y-Achse 
	\draw[->] (0,-4) -- (0,4) node[left,yshift=-1em] {$f(x)$};%node[above left]
	%Zahlen auf y-Achse
	\foreach \y in {-3,...,-1}
	\draw[shift={(0,\y)},color=black] (2pt,0pt) -- (-2pt,0pt) node[left] {\footnotesize $\y$};
	\foreach \y in {1,...,3}
	\draw[shift={(0,\y)},color=black] (2pt,0pt) -- (-2pt,0pt) node[left] {\footnotesize $\y$};
	%Ursprung
	
	
	
	\draw[thick,color=blue] plot[samples=200, domain=-2.1:2.1] 
	(\x,{sinh(\x )}) node[yshift=-2em, xshift=2em] {$\sinh(x)$};
	%\draw[color=blue,thick] (0pt,0pt) node {$[$};
	%\draw[color=blue,thick] (3.14,0pt) node { $]$};
	%\draw[color=blue,thick] (0,0pt)--(4,0);
	
	\end{tikzpicture}\\
	Umkehrbahr im ganzen Intervall. \\
	$\sinh:\mathbb{R}\longmapsto \mathbb{R}$\\ 
\end{center}
\begin{align*}
&\mathrm{arsinh: \mathbb{R}\longmapsto \mathbb{R}}\\
&x\longmapsto \sinh^{-1}(x)=\mathrm{arsinh(x)}
\end{align*}
\begin{center}
	\begin{tikzpicture}[xscale=1, yscale=1]
	
	% Gitternetzlinien
	\draw[xstep=1cm, ystep=1cm, very thin, color=lightgray] (-4,4) grid (4,-4);
	% x-Achse
	\draw[->] (-4,0) -- (4,0) node[below,xshift=-1em] {$x$}; 
	%Zahlen auf x-Achse
	\foreach \x in {-3,...,-1}
	\draw[shift={(\x,0)},color=black] (0pt,2pt) -- (0pt,-2pt) node[below] {\footnotesize $\x$};
	\foreach \x in {1,...,3}
	\draw[shift={(\x,0)},color=black] (0pt,2pt) -- (0pt,-2pt) node[below] {\footnotesize $\x$};
	% y-Achse 
	\draw[->] (0,-4) -- (0,4) node[left,yshift=-1em] {$f(x)$};%node[above left]
	%Zahlen auf y-Achse
	\foreach \y in {-3,...,-1}
	\draw[shift={(0,\y)},color=black] (2pt,0pt) -- (-2pt,0pt) node[left] {\footnotesize $\y$};
	\foreach \y in {1,...,3}
	\draw[shift={(0,\y)},color=black] (2pt,0pt) -- (-2pt,0pt) node[left] {\footnotesize $\y$};
	%Ursprung
	
	\draw[thick,color=blue] plot[samples=200, domain=-4:4] 
	(\x,{(ln(\x+sqrt(\x*\x+1) ))}) node[yshift=2, xshift=-3em] {$\mathrm{arsinh}(x)$};
	%\draw[color=blue,thick] (0pt,0pt) node {$[$};
	%\draw[color=blue,thick] (3.14,0pt) node { $]$};
	%\draw[color=blue,thick] (0,0pt)--(4,0);
	
	\end{tikzpicture}\\
	
\end{center}
Ableitung von Areasinus Hyperbolikus.\\
\begin{align*}
\mathrm{arsinh}'(x)=\frac{1}{\sqrt{x^2+1}}
\end{align*}
Analog zum Arcussinus.\\
\begin{align*}
\arcsin'(x)=\frac{1}{\sqrt{1-x^2}}
\end{align*}
\subsubsection{ Areatangens Hyperbolicus und Areakotangens Hyperbolicus}
\href{https://youtu.be/B5OvQbfafVQ}{Video 16}\\
$\blacktriangleright$ Die Umkehrfunktion von Tangens Hyperbolikus ist der Areatangens Hyperbolikus.\\
\begin{center}
	\begin{tikzpicture}[xscale=3, yscale=2]
	
	% Gitternetzlinien
	\draw[xstep=1cm, ystep=1cm, very thin, color=lightgray] (-2,2) grid (2,-2);
	% x-Achse
	\draw[->] (-2,0) -- (2,0) node[below,xshift=-1em] {$x$}; 
	%Zahlen auf x-Achse
	\foreach \x in {-1,...,-1}
	\draw[shift={(\x,0)},color=black] (0pt,2pt) -- (0pt,-2pt) node[below] {\footnotesize $\x$};
	\foreach \x in {1,...,1}
	\draw[shift={(\x,0)},color=black] (0pt,2pt) -- (0pt,-2pt) node[below] {\footnotesize $\x$};
	% y-Achse 
	\draw[->] (0,-2) -- (0,2) node[left,yshift=-1em] {$f(x)$};%node[above left]
	%Zahlen auf y-Achse
	\foreach \y in {-1,...,-1}
	\draw[shift={(0,\y)},color=black] (2pt,0pt) -- (-2pt,0pt) node[left] {\footnotesize $\y$};
	\foreach \y in {1,...,2}
	\draw[shift={(0,\y)},color=black] (2pt,0pt) -- (-2pt,0pt) node[left] {\footnotesize $\y$};
	%Ursprung
	\draw[color=black] (0pt,-10pt) node[right] {\footnotesize $0$};
	\draw[thick,color=blue] plot[samples=200, domain=-2.0:2.0] 
	(\x,{tanh(\x )}) node[yshift=-2em, xshift=-2em] {$\tanh(x)$};
	\draw[color=blue,dashed] (-2,1)--(2,1);
	\draw[color=blue,dashed] (-2,-1)--(2,-1);
	%		\draw[thick,color=blue] plot[samples=200, domain=-2:2] 
	%	(\x,{sinh(\x )}) node[yshift=1, xshift=2.5em] {$\arcsin(x)$};
	%	\draw[color=black] (2pt,3.14) -- (-2pt,3.14) node[left] {\footnotesize $\pi$};
	\end{tikzpicture}\\
	\begin{align*}
	\tanh: \mathbb{R}\longmapsto \;]-1,1[\\
	\end{align*}
	
	
\end{center}
\begin{align*}
&\mathrm{artanh}:\;]-1,1[\;\longmapsto \mathbb{R}\\
&x\longmapsto \tanh^{-1}=\mathrm{artanh}(x)\\
\end{align*}
\begin{center}
	\begin{tikzpicture}[xscale=3, yscale=3]
	
	% Gitternetzlinien
	\draw[xstep=1cm, ystep=1cm, very thin, color=lightgray] (-2,2) grid (2,-2);
	% x-Achse
	\draw[->] (-2,0) -- (2,0) node[below,xshift=-1em] {$x$}; 
	%Zahlen auf x-Achse
	\foreach \x in {-2,...,-1}
	\draw[shift={(\x,0)},color=black] (0pt,2pt) -- (0pt,-2pt) node[below] {\footnotesize $\x$};
	\foreach \x in {1,...,2}
	\draw[shift={(\x,0)},color=black] (0pt,2pt) -- (0pt,-2pt) node[below] {\footnotesize $\x$};
	% y-Achse 
	\draw[->] (0,-2) -- (0,2) node[left,yshift=-1em] {$f(x)$};%node[above left]
	%Zahlen auf y-Achse
	\foreach \y in {-1,...,-1}
	\draw[shift={(0,\y)},color=black] (2pt,0pt) -- (-2pt,0pt) node[left] {\footnotesize $\y$};
	\foreach \y in {1,...,1}
	\draw[shift={(0,\y)},color=black] (2pt,0pt) -- (-2pt,0pt) node[left] {\footnotesize $\y$};
	%Ursprung
	\draw[thick,color=blue] plot[samples=200, domain=-0.997:0.9997] 
	(\x,{(0.5*ln(sqrt((1+\x)/(1-\x)))}) node[yshift=2, xshift=-3em] {$\mathrm{artanh}(x)$};
	
	
	%\draw[thick,color=blue] plot[samples=200, domain=0:2.05] 
	%	(\x,{cosh(\x )}) node[yshift=-2em, xshift=2em] {$\cosh(x)$};
	%	\draw[color=blue,thick] (0pt,0pt) node {$[$};
	%\draw[color=blue,thick] (3.14,0pt) node { $]$};
	%	\draw[color=blue,thick] (0,0pt)--(4,0);
	
	\end{tikzpicture}\\	
\end{center}
Die Ableitung von Areatangens Hyperbolicus:\\
\begin{align*}
\mathrm{artanh}'(x)=\frac{1}{1-x^2}
\end{align*}
\newpage
\noindent
$\blacktriangleright$ Die Umkehrfunktion von Kotangens Hyperbolikus ist der Areakotangens Hyperbolikus.\\
\begin{center}
	\begin{tikzpicture}[xscale=1.5, yscale=1.5]
	
	% Gitternetzlinien
	\draw[xstep=1cm, ystep=1cm, very thin, color=lightgray] (-4,4) grid (4,-4);
	% x-Achse
	\draw[->] (-4,0) -- (4,0) node[below,xshift=-1em] {$x$}; 
	%Zahlen auf x-Achse
	\foreach \x in {-3,...,-1}
	\draw[shift={(\x,0)},color=black] (0pt,2pt) -- (0pt,-2pt) node[below] {\footnotesize $\x$};
	\foreach \x in {1,...,3}
	\draw[shift={(\x,0)},color=black] (0pt,2pt) -- (0pt,-2pt) node[below] {\footnotesize $\x$};
	% y-Achse 
	\draw[->] (0,-4) -- (0,4) node[left,yshift=-1em] {$f(x)$};%node[above left]
	%Zahlen auf y-Achse
	\foreach \y in {-3,...,-1}
	\draw[shift={(0,\y)},color=black] (2pt,0pt) -- (-2pt,0pt) node[left] {\footnotesize $\y$};
	\foreach \y in {1,...,3}
	\draw[shift={(0,\y)},color=black] (2pt,0pt) -- (-2pt,0pt) node[left] {\footnotesize $\y$};
	%Ursprung
	\draw[color=black] (0pt,-10pt) node[right] {\footnotesize $0$};
	\draw[thick,color=blue] plot[samples=200, domain=-4.0:-.26] 
	(\x,{cosh(\x )/sinh(\x)}) node[yshift=2em, xshift=-2em] {$\coth(x)$};
	\draw[thick,color=blue] plot[samples=200, domain=.26:4] 
	(\x,{cosh(\x )/sinh(\x)}) node[yshift=-2em, xshift=-2em] {$\coth(x)$};
	
	\draw[color=blue,dashed] (-4,1)--(4,1);
	\draw[color=blue,dashed] (-4,-1)--(4,-1);
	
	
	%		\draw[thick,color=blue] plot[samples=200, domain=-2:2] 
	%	(\x,{sinh(\x )}) node[yshift=1, xshift=2.5em] {$\arcsin(x)$};
	%	\draw[color=black] (2pt,3.14) -- (-2pt,3.14) node[left] {\footnotesize $\pi$};
	\end{tikzpicture}\\
	\begin{align*}
	\coth:\mathbb{R}\backslash\{0\}\longmapsto \mathbb{R}\backslash \{-1,1\}
	\end{align*}
	
\end{center}
\begin{align*}
\mathrm{arcoth}:\;]-\infty,-1[\; \cup\;]1,\infty[\;\longmapsto \mathbb{R}\setminus \{0\}\\
\end{align*}
\begin{center}
	\begin{tikzpicture}[xscale=2, yscale=2]
	
	% Gitternetzlinien
	\draw[xstep=1cm, ystep=1cm, very thin, color=lightgray] (-4,2) grid (4,-2);
	% x-Achse
	\draw[->] (-4,0) -- (4,0) node[below,xshift=-1em] {$x$}; 
	%Zahlen auf x-Achse
	\foreach \x in {-3,...,-1}
	\draw[shift={(\x,0)},color=black] (0pt,2pt) -- (0pt,-2pt) node[below] {\footnotesize $\x$};
	\foreach \x in {1,...,3}
	\draw[shift={(\x,0)},color=black] (0pt,2pt) -- (0pt,-2pt) node[below] {\footnotesize $\x$};
	% y-Achse 
	\draw[->] (0,-2) -- (0,2) node[left,yshift=-1em] {$f(x)$};%node[above left]
	%Zahlen auf y-Achse
	\foreach \y in {-1,...,-1}
	\draw[shift={(0,\y)},color=black] (2pt,0pt) -- (-2pt,0pt) node[left] {\footnotesize $\y$};
	\foreach \y in {1,...,1}
	\draw[shift={(0,\y)},color=black] (2pt,0pt) -- (-2pt,0pt) node[left] {\footnotesize $\y$};
	%Ursprung
	\draw[thick,color=blue] plot[samples=200, domain=1.0009:4] 
	(\x,{(0.5*ln(sqrt((1+\x)/(\x-1)))});
	\draw[thick,color=blue] plot[samples=200, domain=-4:-1.00009] 
	(\x,{(0.5*ln(sqrt((1+\x)/(\x-1)))}) node[yshift=2, xshift=-3em] {$\mathrm{arcoth}(x)$};
	
	%\draw[thick,color=blue] plot[samples=200, domain=0:2.05] 
	%	(\x,{cosh(\x )}) node[yshift=-2em, xshift=2em] {$\cosh(x)$};
	%	\draw[color=blue,thick] (0pt,0pt) node {$[$};
	%\draw[color=blue,thick] (3.14,0pt) node { $]$};
	%	\draw[color=blue,thick] (0,0pt)--(4,0);
	
	\end{tikzpicture}\\	
\end{center}
Die Ableitung von Areakotangens Hyperbolicus:\\
\begin{align*}
\mathrm{arcoth}'(x)=-\frac{1}{1-x^2}
\end{align*}
\subsubsection{Aufgabe 1 zu Area-Funktionen}
\href{https://youtu.be/fkV0enqPc9E}{Video 17}\\
Zeige für $ x\in \mathbb{R} $ gilt $\mathrm{arsinh}=\ln(x+\sqrt{x^2+1})$.\\
Beweis:\\
Es gilt:\\
\begin{align*}
f(x)=\mathrm{arsinh}(x)-\ln(x+\sqrt{x^2+1})=0
\end{align*}
Wenn dies gelten soll, muß die Ableitung von $ f(x) $ konstant sein , denn $ f(x) $ muss ja Null sein wenn $\mathrm{arsinh}=\ln(x+\sqrt{x^2+1})  $ gelten soll.\\
\begin{align*}
f'(x)&=\mathrm{arsinh}'(x)-\big (\ln(x+\sqrt{x^2+1})\big )'\\
&=\frac{1}{\sqrt{x^2+1}}-\frac{1}{x+\sqrt{x^2+1}}\big(x+\sqrt{x^2+1}\big)'\\
&=\frac{1}{\sqrt{x^2+1}}-\frac{1}{x+\sqrt{x^2+1}}\bigg(1+\frac{1}{2\sqrt{x^2+1}}\cdot 2x\bigg)\\
&=\frac{1}{\sqrt{x^2+1}}-\frac{1}{x+\sqrt{x^2+1}}\bigg(1+\frac{x}{\sqrt{x^2+1}}\bigg)\\
&=\frac{1}{\sqrt{x^2+1}}-\frac{1}{x+\sqrt{x^2+1}}\bigg(\frac{\sqrt{x^2+1}+x}{\sqrt{x^2+1}}\bigg)\\
&=\frac{1}{\sqrt{x^2+1}}-\frac{x+\sqrt{x^2+1}}{(x+\sqrt{x^2+1})(\sqrt{x^2+1})}\\
&=\frac{x+\sqrt{x^2+1}-(x+\sqrt{x^2+1})}{(\sqrt{x^2+1})(x+\sqrt{x^2+1})}=0\\
\end{align*}
D.h.: $ f'(x)=0 $ für $ x\in \mathbb{R} \overset{\tiny Konstanzkriterium}{\longmapsto} f(x)=k\in\mathbb{R} $ für alle $ x\in \mathbb{R} $\\
Es gilt für $ f(0) $:\\
\begin{align*}
f(0)=\mathrm{arsinh}(0)-\ln(0+\sqrt{0^2+1})=0
\end{align*}
\begin{align*}
&0\overset{\mathrm{sinh}}{\longmapsto} 0\\
&0\overset{\mathrm{arsinh}}{\longmapsto} 0\\
\end{align*}
\subsubsection{Aufgabe 2 zu Area-Funktionen}
\href{https://youtu.be/y-3Ng9HCIEo}{Video 18}\\
Zeige für $ x\geq 1 $ gilt $\mathrm{arcosh}=\ln(x+\sqrt{x^2-1})$.\\
Beweis:\\
Es gilt:\\
\begin{align*}
g(x)=\mathrm{arcosh}(x)-\ln(x+\sqrt{x^2-1})=0
\end{align*}
Für $ x>1 $ gilt:
\begin{align*}
g'(x)&=\mathrm{arcosh}'(x)-\big(\ln(x+\sqrt{x^2-1})\big)'\\
&=\frac{1}{\sqrt{x^2-1}}-\frac{1}{x+\sqrt{x^2-1}}\big(x+\sqrt{x^2-1}\big)'\\
&=\frac{1}{\sqrt{x^2-1}}-\frac{1}{x+\sqrt{x^2-1}}\bigg(1+\frac{x}{\sqrt{x^2-1}}\bigg)\\
&=\frac{1}{\sqrt{x^2-1}}-\frac{\sqrt{x^2-1}+x}{(x+\sqrt{x^2-1})(\sqrt{x^2-1})}\\
&=\frac{x+\sqrt{x^2-1}-(\sqrt{x^2-1}+x)}{(x+\sqrt{x^2-1})(\sqrt{x^2-1})}=0
\end{align*}
D.h.: $ g'(x)=0 $ für $ x>1 \overset{\tiny Konstanzkriterium}{\longmapsto} g(x)=k\in\mathbb{R} $ für alle $ x \geq 1 $\\
Es gilt für $ g(1) $:\\
\begin{align*}
k=g(1)=\mathrm{arcosh}(1)-\ln(x+\sqrt{1^2-1})=0
\end{align*}
\begin{align*}
&0\overset{\mathrm{cosh}}{\longmapsto} 1\\
&1\overset{\mathrm{arcosh}}{\longmapsto} 0\\
\end{align*}
\subsection{Die Exponentialfunktion}
\href{https://youtu.be/au-TAWRpb3Y}{Video 19}\\
Es gibt genau eine diffbar. Funktion $ \exp:\mathbb{R}\longmapsto \mathbb{R}_{>0} $ mit \\
\begin{align*}
\exp'(x)=\exp(x)\;\;\text{und }\;\;\exp(0)=1 .
\end{align*}
Diese Funktion heißt Eponentialfunktion, wobei gilt:\\
\begin{align*}
&\mathrm{e}^x:=\exp(x)\\
&\mathrm{e}=\exp(1)=2.7182\ldots\in\mathbb{R}\backslash \mathbb{Q}\\
\end{align*}
Bemerkung:\\
$\mathbb{R}\backslash \mathbb{Q}  $ sind die irrationalen Zahlen . Ist die Menge aller Elemente von $ \mathbb{R} $, die nicht in  $\mathbb{Q}$ liegen.\\
\\
Wegen $ \big(\mathrm{e}^x\big)'=\mathrm{e}^x>0 \overset{Monotoniekr.}{\Longrightarrow} \mathrm{e}^x$ ist streng monoton wachsend\\
\begin{center}
	\begin{tikzpicture}[xscale=1, yscale=1]
	
	% Gitternetzlinien
	\draw[xstep=1cm, ystep=1cm, very thin, color=lightgray] (-4,4) grid (4,-2);
	% x-Achse
	\draw[->] (-4,0) -- (4,0) node[below,xshift=-1em] {$x$}; 
	%Zahlen auf x-Achse
	\foreach \x in {-3,...,-1}
	\draw[shift={(\x,0)},color=black] (0pt,2pt) -- (0pt,-2pt) node[below] {\footnotesize $\x$};
	\foreach \x in {1,...,3}
	\draw[shift={(\x,0)},color=black] (0pt,2pt) -- (0pt,-2pt) node[below] {\footnotesize $\x$};
	% y-Achse 
	\draw[->] (0,-2) -- (0,4) node[left,yshift=-1em] {$f(x)$};%node[above left]
	%Zahlen auf y-Achse
	\foreach \y in {-1,...,-1}
	\draw[shift={(0,\y)},color=black] (2pt,0pt) -- (-2pt,0pt) node[left] {\footnotesize $\y$};
	\foreach \y in {1,...,3}
	\draw[shift={(0,\y)},color=black] (2pt,0pt) -- (-2pt,0pt) node[left] {\footnotesize $\y$};
	%Ursprung
	\draw[color=black] (0pt,-10pt) node[right] {\footnotesize $0$};
	
	
	\draw[thick,color=blue] plot[samples=200, domain=-4:1.39] 
	(\x,{e^\x}) node[below, xshift=2em]{$\exp(x)$};
	
	
	
	\end{tikzpicture}\\
\end{center}
Es gilt das Taylorpolynom:\\
\begin{align*}
\exp(x)=\mathrm{e}^x=\sum\limits_{k=0}^{n}\frac{x^k}{k!}+\mathrm{R}_n(x)
\end{align*}
\underline {Grenzwertbetrachtungen:}
\begin{quote}
	\begin{align*}
	&\lim\limits_ {n\to\infty}\frac{\mathrm{e}^x}{x^n}=\infty\;\;\text{für }\;\;n\in\mathbb{N}\\
	&\lim\limits_ {n\to-\infty}|x^n|\cdot\mathrm{e}^x=0\;\;\text{für }\;\;n\in\mathbb{N}\\
	\end{align*}
	Die Exponentialfunktion wächst schneller gegen Unendlich als jede Potenz von $ x $.\\
	Ebenso geht sie bei Minusunendlich schneller gegen Null als der Betrag von der Potenz von $ x $ wächst.\\
\end{quote}
\underline{Eigenschaften:}\\
Für $ x,y\in\mathbb{R} $ gilt:\\
\begin{align*}
&\mathrm{e}^x\cdot \mathrm{e}^y=\mathrm{e}^{x+y}\\
&\mathrm{e}^{-x}=\frac{1}{\mathrm{e}^x}
\end{align*}
\underline{Für die allgemeine Exponentialfunktion gilt:}\\
Für $ a>0  $ definiert man:\\
\begin{align*}
a^x:=\mathrm{e}^{\ln(a)\cdot x}
\end{align*}
$ a $ nennt man Basis\\
$ x $ nennt man Exponent\\
Wegen $ (-1)^{\genfrac{}{}{}{2}{1 }{2 }} =\sqrt{-1}$ muß $ a>0 $ sein. Negative Wurzel! \\
Für $ 0<a<1 $ ist $ a^x $ streng monoton fallend.\\
Für $ a>1 $ ist $ a^x $ streng monoton wachsend.\\
\begin{center}
	\begin{tikzpicture}[xscale=1, yscale=1]
	
	% Gitternetzlinien
	\draw[xstep=1cm, ystep=1cm, very thin, color=lightgray] (-4,4) grid (4,-2);
	% x-Achse
	\draw[->] (-4,0) -- (4,0) node[below,xshift=-1em] {$x$}; 
	%Zahlen auf x-Achse
	\foreach \x in {-3,...,-1}
	\draw[shift={(\x,0)},color=black] (0pt,2pt) -- (0pt,-2pt) node[below] {\footnotesize $\x$};
	\foreach \x in {1,...,3}
	\draw[shift={(\x,0)},color=black] (0pt,2pt) -- (0pt,-2pt) node[below] {\footnotesize $\x$};
	% y-Achse 
	\draw[->] (0,-2) -- (0,4) node[left,yshift=-1em] {$f(x)$};%node[above left]
	%Zahlen auf y-Achse
	\foreach \y in {-1,...,-1}
	\draw[shift={(0,\y)},color=black] (2pt,0pt) -- (-2pt,0pt) node[left] {\footnotesize $\y$};
	\foreach \y in {1,...,3}
	\draw[shift={(0,\y)},color=black] (2pt,0pt) -- (-2pt,0pt) node[left] {\footnotesize $\y$};
	%Ursprung
	\draw[color=black] (0pt,-10pt) node[right] {\footnotesize $0$};
	
	
	\draw[thick,color=blue] plot[samples=200, domain=-4:1.99] 
	(\x,{2^\x}) node[below, xshift=1em]{$2^x$};
	\draw[thick,color=red] plot[samples=200, domain=-2:4] 
	(\x,{0.5^\x}) node[below, yshift=2em,xshift=-2em]{$\big(\frac{1}{2}\big)^x$};
	
	
	\end{tikzpicture}\\
\end{center}
Die Achsen-Symmetrie: $ f(x)=f(-x) $ gilt hier wegen:\\
\begin{align*}
\bigg(\frac{1}{2}\bigg)^x=2^{-x}
\end{align*}
\underline{Ableitung}\\
\begin{align*}
(a^x)'=\bigg(\mathrm{e}^{\ln(a)x}\bigg)'=\mathrm{e}^{\ln(a)x}\cdot(\ln(a)x)'=a^x\ln(a)\\
\end{align*}
Bemerkung:\
$ \ln(a) $ ist für $ a<1 $ negativ daraus folgt $ a^x $ ist monoton fallend.\\
\underline{Rechenregeln:}
\begin{align*}
&a^x=\mathrm{e}^{\ln(a)x}\\
&\big(a^b\big)^c=a^{b\cdot c}\\
&\sqrt{x^7}=\big(x^7\big)^{\frac{1}{2}}=x^{\frac{7}{2}}=(\sqrt{x})^7
\end{align*}
\subsubsection{Der Logarithmus}
\href{https://youtu.be/gCVOSyUEYIk}{Video 20}\\
Der Logarithmus ist das Gegenstück zur Exponentialfunktion.\\
Der natürliche Logarithmus ist die Umkehrfunktion von der Eponentialfunktion.
\begin{align*}
\exp:\mathbb{R}\longmapsto \;]0,\infty[ 
\end{align*}\\
D.h.:
\begin{align*}
\ln:\;]0,\infty[ \;\longmapsto \mathbb{R} \;\;\text{mit }\;\;x\longmapsto \exp^{-1}(x)=\ln(x)
\end{align*}
\begin{center}
	
	
	\begin{tikzpicture}[xscale=1, yscale=1]
	
	% Gitternetzlinien
	\draw[xstep=1cm, ystep=1cm, very thin, color=lightgray] (-4,4) grid (4,-4);
	% x-Achse
	\draw[->] (-4,0) -- (4,0) node[below,xshift=-1em] {$x$}; 
	%Zahlen auf x-Achse
	\foreach \x in {-3,...,-1}
	\draw[shift={(\x,0)},color=black] (0pt,2pt) -- (0pt,-2pt) node[below] {\footnotesize $\x$};
	\foreach \x in {1,...,3}
	\draw[shift={(\x,0)},color=black] (0pt,2pt) -- (0pt,-2pt) node[below] {\footnotesize $\x$};
	% y-Achse 
	\draw[->] (0,-4) -- (0,4) node[left,yshift=-1em] {$f(x)$};%node[above left]
	%Zahlen auf y-Achse
	\foreach \y in {-1,...,-3}
	\draw[shift={(0,\y)},color=black] (2pt,0pt) -- (-2pt,0pt) node[left] {\footnotesize $\y$};
	\foreach \y in {1,...,3}
	\draw[shift={(0,\y)},color=black] (2pt,0pt) -- (-2pt,0pt) node[left] {\footnotesize $\y$};
	%Ursprung
	\draw[color=black] (0pt,-10pt) node[right] {\footnotesize $0$};
	
	\draw[thick,color=blue] plot[samples=200, domain=-4:1.37] 
	(\x,{e^\x}) node[below, xshift=2em] {$\exp(x)$};
	
	\draw[thick,color=blue] plot[samples=200, domain=0.04:4] 
	(\x,{ln(\x)}) node[yshift=1em,xshift=-2em]{$\ln(x)$};
	
	
	
	\end{tikzpicture}\\
\end{center}
\pagebreak
Eigenschaften:\\
\begin{quote}
	Sei $ x,y\in\;]0,\infty[ $\\
	\begin{align*}
	&\text{(i) }\;\;\ln(x\cdot y)=\ln(x)+\ln(y)\\
	\\
	&\text{(ii) }\;\ln(x^\alpha) =\alpha\cdot \ln(x)\;\;\text{für }\;\;\alpha\in\mathbb{R})\\
	\\
	&\text{(iii) }\lim\limits_ {x\to\infty}\ln(x)=\infty  \\
	&\;\;\;\;\;\;\;\;\lim\limits_ {\substack{x\mapsto 0\\x>0}}\ln(x)=-\infty\\
	\end{align*}
\end{quote}
Für den Logarithmus zur Basis $ a $ mit $ a>0 $ gilt:\\
\begin{align*}
\log_a(x)=\frac{\ln(x)}{\ln(a)}=y
\end{align*}
Die Umkehrfunktion von der allgemeinen Exponentialfunktion: $ a^x $ ist der Logarithmus: $\log_a(x) $.\\
Beweis:\\
Es gilt die Eigenschaft der Umkehrabbildung:\\
\begin{align*}
f(x)\circ f^{-1}(x)=x=f^{-1}(x)\circ f(x)\\
\end{align*}
Dann gilt für: $ f(x)=a^x=\mathrm{e}^{\ln(a)x}$  und $f^{-1}(x)=\log_a(x) $:\\
\begin{align*}
f(x)\circ f^{-1}(x)&=\mathrm{e}^{\ln(a)(\log_a(x)}=\mathrm{e}^{\ln(a)\frac{\ln(x)}{\ln(a)}}=\mathrm{e}^{\ln(x)}=x_{_{\Square}}\\
\end{align*}
Eigenschaften:\\
\begin{quote}
	Sei $ a,b,c>0\in \mathbb{R}  $ und $ \alpha\in\;]0,\infty[ $\\
	\begin{align*}
	&\text{(i) }\;\;\log_a(b)\cdot \log_b(c)=\frac{\ln(b)}{\ln(a)}\cdot\frac{\ln(c)}{\ln(b)}=\frac{\ln(c)}{\ln(a)}=\log_a(c)\\
	\\
	&\text{(ii) }\lim\limits_ {x\to\infty}\frac{\ln(x)}{x^\alpha}=0 \;\;\text{für }\;\;\alpha=\in\;]0,\infty[\\
	&\;\;\;\;\;\;\lim\limits_ {x\to\infty}\frac{\ln(x)}{\sqrt{x}}=0\;\;\text{für }\;\;\alpha=\frac{1}{2}\\
	\end{align*}
\end{quote}
\underline{Kategorien des Wachstums von Funktionen:}\\
\begin{quote}
	1. Exponentielle $ \mathrm{e}^x,\; a^x $\\
	2. Polynominelle $ x^n,\;\; n\in\mathbb{N} $ (für $ n=1  $ : Funktion ist linear.)\\
	3. Wurzel $ \sqrt[k]{x},\;\; k\in\mathbb{N}\backslash \{0\} $
\end{quote}
\subsubsection{Aufgabe zum Thema Exponentialfunktion und Logarithmus Teil 1+2}
\href{https://youtu.be/g6oiy5HrFNo}{Video 21}\\
Bestimme den maximalen Definitionsbereich, die Ableitung und das Bild von:\\
\begin{align*}
f(x)=x^{\big(\sqrt{x}^x\big)}\\
\end{align*}
\underline{Maximaler Definitionsbereich:}\\
\begin{align*}
\mathbb{D}_f=[0,\infty[
\end{align*}
\underline{Das Bild von f:}\\
\begin{quote}
	(i)\;\;$ f(0)=0^{\big(\sqrt{0}^0\big)}=0^{(0^0)}=0^1=0 $\\
	\\
	(ii)\;$ f(x)\geq 0 $ für alle $ x\in \mathbb{D}_f $\\
	\\
	(iii)$ \lim\limits_ {x\to\infty}x^{\big(\sqrt{x}^x\big)}=\;\;? $\\
	\begin{quote}
		$ f(x)=x^{\big(\sqrt{x}^x\big)}=\mathrm{e}^{\ln(x)\cdot \sqrt{x}^x}=\mathrm{e}^{\ln(x)\cdot x^{\genfrac{}{}{}{2}{x }{2 }}} $\\
		$ \lim\limits_ {x\to\infty}f(x)=\lim\limits_ {x\to\infty}\mathrm{e}^{\ln(x)\cdot x^{\frac{x}{2}}} =\infty $	
	\end{quote}
	Fazit: Wegen (i), (ii) und (iii) folgt:\\
	\begin{align*}
	\mathrm{Bild}(f)=[0,\infty[
	\end{align*}
\end{quote}
\newpage
\noindent
\href{https://youtu.be/qDxy5Th1pfc}{Video 22}\\
\underline{Die Ableitung von $ f(x) $:}\\
\begin{align*}
f'(x)&=\bigg(x^{(\sqrt{x}^x)}\bigg)'=\bigg(\mathrm{e}^{\ln(x)\cdot x^{\frac{x}{2}}}\bigg)'\\
&=\mathrm{e}^{\ln(x)\cdot x^{\frac{x}{2}}}\cdot\big(\ln(x)\cdot x^{\frac{x}{2}}\big)'\\
&=\mathrm{e}^{\ln(x)\cdot x^{\frac{x}{2}}}\cdot \big[\ln'(x)\cdot x^{\frac{x}{2}}+\big(x^{\frac{x}{2}}\big)'\cdot \ln(x)\big]\\
&=\mathrm{e}^{\ln(x)\cdot x^{\frac{x}{2}}}\cdot\bigg[\frac{1}{x}\cdot x^{\frac{x}{2}}+\big(\mathrm{e}^{\ln(x)\cdot \frac{x}{2}}\big)'\ln(x)\bigg]\\
&=\mathrm{e}^{\ln(x)\cdot x^{\frac{x}{2}}}\cdot\bigg[\frac{1}{x}\cdot x^{\frac{x}{2}}+\big[\mathrm{e}^{\ln(x)\cdot \frac{x}{2}}(\genfrac{}{}{}{2}{1 }{x }\cdot \genfrac{}{}{}{2}{x }{2 }+\ln(x)\cdot \genfrac{}{}{}{2}{1}{2 })\big]\ln(x)\bigg]\\
&=\mathrm{e}^{\ln(x)\cdot x^{\frac{x}{2}}}\cdot \bigg[\frac{1}{x}\cdot x^{\frac{x}{2}}+\big[ x^{\frac{x}{2}}(\genfrac{}{}{}{2}{1 }{2 }+ \genfrac{}{}{}{2}{\ln(x)}{2 })\big]\ln(x)\bigg]\\
&=\mathrm{e}^{\ln(x)\cdot x^{\frac{x}{2}}}\cdot\bigg[\frac{1}{x}\cdot x^{\frac{x}{2}}+\big[\genfrac{}{}{}{2}{1 }{2} x^{\frac{x}{2}}(1+ \ln(x))\big]\ln(x)\bigg]\\
&=\mathrm{e}^{\ln(x)\cdot x^{\frac{x}{2}}}\cdot x^{\frac{x}{2}}\bigg[\genfrac{}{}{}{2}{1 }{x }+\big[\genfrac{}{}{}{2}{\ln(x) }{2 }(1+\ln(x)\big ]\bigg ]\\
&=x^{(\sqrt{x}^x)}\cdot x^{\frac{x}{2}}\bigg[\genfrac{}{}{}{2}{1 }{x }+\big[\genfrac{}{}{}{2}{\ln(x) }{2 }(1+\ln(x)\big ]\bigg ]\\
&=x^{\big(\sqrt{x}^x+x^{\frac{x}{2}}\big)}\bigg[\genfrac{}{}{}{2}{1 }{x }+\big[\genfrac{}{}{}{2}{\ln(x) }{2 }(1+\ln(x)\big ]\bigg ]\\
\end{align*}
\subsubsection{ Lösen von Gleichungen mittels der Exponentialfunktion und dem Logarithmus}
\href{https://youtu.be/w04eRCcWpHU}{Video 23}\\
Bestimme alle Lösungen von:\\
\begin{align*}
4^{4x}=2^{6x}\cdot 4^x\cdot 256^{x-1}
\end{align*}
Gleiche Basen erzeugen:\\
\begin{align*}
\Leftrightarrow \big(2^2\big)^{4x}=2^{6x}\cdot \big(2^2\big)^{x}\cdot\big(2^8\big)^{x-1}
\end{align*}
Anwendung von: $ \big(a^b\big)^c=a^{a\cdot c} $\\
\begin{align*}
\Leftrightarrow  2^{2\cdot 4x}=2^{6x}\cdot 2^{2x}\cdot 2^{8(x-1)}
\end{align*}
Anwendung von: $ a^b\cdot a^c=a^{a+b} $\\
\begin{align*}
\Leftrightarrow 2^{8x}=2^{6x+2x+8(x-1)}=2^{16x-8}\\
\end{align*}
Mit $ 2^{8x} $ teilen.\\
\begin{align*}
\Leftrightarrow 1&=2^{16x-8}\cdot 2^{-8x}\\
&=2^{8x-8}
\end{align*}
Mit $ 1=2^0 $ ergibt:\\
\begin{align*}
\Leftrightarrow2^0=2^{8x-8}\\
\end{align*}
Mit $ \log_2(...) $ ergibt:\\
\begin{align*}
\Leftrightarrow \log_2(2^0)&=\log_2\big(2^{(8x-8)}\big)\\
\end{align*}
Anwendung von: $ \log(a^b)=b\log(a) $
\begin{align*}
\log_2(2^0)&=(8x-8)\log_2(2)\\
0&=8x-8\\
\Leftrightarrow 8x&=8\\
\Leftrightarrow x&=1
\end{align*}
Fazit: $ x=1 $ ist die einzigste Lösung.\\
%	\chapter{Integration}
\setcounter{section}{11}
\setcounter{subsection}{0}
\setcounter{subsubsection}{1}
\setcounter{secnumdepth}{4}
\subsection{Integration}
\href{https://youtu.be/MOoeV2ewtj8}{Video 01}\\
\begin{center}
	
	Sei $ f:[a,b] \longmapsto \mathbb{R}$
	
	\begin{tikzpicture}[xscale=1, yscale=1]
	
	% Gitternetzlinien
	\draw[xstep=1cm, ystep=1cm, very thin, color=lightgray] (-1,6) grid (8,-1);
	% x-Achse
	\draw[->] (-1,0) -- (8,0) node[below,xshift=-1em] {$x$}; 
	%Zahlen auf x-Achse
	%\foreach \x in {-3,...,-1}
	%\draw[shift={(\x,0)},color=black] (0pt,2pt) -- (0pt,-2pt) node[below] {\footnotesize $\x$};
	%\foreach \x in {a,...,7}
	%\draw[shift={(\x,0)},color=black] (0pt,2pt) -- (0pt,-2pt) node[below] {\footnotesize $\x$};
	
	% y-Achse 
	\draw[->] (0,-1) -- (0,6) node[left,yshift=-1em] {$f(x)$};%node[above left]
	%Zahlen auf y-Achse
	%\foreach \y in {-1,...,-1}
	%\draw[shift={(0,\y)},color=black] (2pt,0pt) -- (-2pt,0pt) node[left] {\footnotesize $\y$};
	\foreach \y in {1,...,5}
	\draw[shift={(0,\y)},color=black] (2pt,0pt) -- (-2pt,0pt) node[left] {\footnotesize $\y$};
	%Ursprung
	\draw[color=black] (0pt,-10pt) node[right] {\footnotesize $0$};
	
	\draw[thick,color=black] plot[samples=200, domain=1:7] 
	(\x,{0.1*\x^2}) node[below, xshift=-5em] {$f(x)$};
	
	\fill[color =blue,fill opacity=0.3] plot[domain=1:7] (\x,{(0.1*\x^2 })% 
	-- plot[domain=7:1] (\x,{0});% 
	
	\draw[shift={(1,0)},color=black] (0pt,2pt) -- (0pt,-2pt) node[below] {\footnotesize $a$};
	\draw[shift={(7,0)},color=black] (0pt,2pt) -- (0pt,-2pt) node[below] {\footnotesize $b$};
	\draw[color=black] (5.3,1) node {\footnotesize $\int\limits_{a}^{b} f(x)\mathrm{d}x$};
	\end{tikzpicture}\\
\end{center}
Das bestimmte Integral $\int\limits_{a}^{b} f(x)\mathrm{d}x  $ ist die orientierte Fläche zwischen $ a,b $, die zwischen $ f(x) $ und der x-Achse ist. Oberhalb der x-Achse ist die Fläche positiv und unterhalb negativ.\\
\pagebreak

Durch Zerlegung von $ [a,b] $ in gleichgroße Teile kann man die Flächen berechnen.\\
\begin{center}
	$ a=x_0<x_1<x_3\ldots<x_n=3 $\\
	
	\begin{tikzpicture}[xscale=1, yscale=1]
	
	% Gitternetzlinien
	\draw[xstep=1cm, ystep=1cm, very thin, color=lightgray] (-1,6) grid (8,-1);
	% x-Achse
	\draw[->] (0,0) -- (8,0) node[below,xshift=-1em] {$x$}; 
	%Zahlen auf x-Achse
	%\foreach \x in {-3,...,-1}
	%\draw[shift={(\x,0)},color=black] (0pt,2pt) -- (0pt,-2pt) node[below] {\footnotesize $\x$};
	%\foreach \x in {a,...,7}
	%\draw[shift={(\x,0)},color=black] (0pt,2pt) -- (0pt,-2pt) node[below] {\footnotesize $\x$};
	
	% y-Achse 
	\draw[->] (0,-1) -- (0,5) node[left,yshift=-1em] {$f(x)$};%node[above left]
	%Zahlen auf y-Achse
	%\foreach \y in {-1,...,-1}
	%\draw[shift={(0,\y)},color=black] (2pt,0pt) -- (-2pt,0pt) node[left] {\footnotesize $\y$};
	%\foreach \y in {1,...,5}
	%\draw[shift={(0,\y)},color=black] (2pt,0pt) -- (-2pt,0pt) node[left] {\footnotesize $\y$};
	\draw[shift={(0,0.1)},color=black] (2pt,0pt) -- (-2pt,0pt) node[left] {\footnotesize $f(x_0)$};
	\draw[color=black,dashed] (0,0.1)--(1,0.1);
	\draw[shift={(0,0.9)},color=black] (2pt,0pt) -- (-2pt,0pt) node[left] {\footnotesize $f(x_1)$};
	\draw[color=black,dashed] (0,0.9)--(3,0.9);
	\draw[shift={(0,2.5)},color=black] (2pt,0pt) -- (-2pt,0pt) node[left] {\footnotesize $f(x_2)$};
	\draw[color=black,dashed] (0,2.5)--(5,2.5);
	
	%Ursprung
	\draw[color=black] (0pt,-10pt) node[right] {\footnotesize $0$};
	
	\draw[thick,color=black] plot[samples=200, domain=1:7] 
	(\x,{0.1*\x^2}) node[below, xshift=-5em] {$f(x)$};
	
	\fill[color =blue,fill opacity=0.3] plot[domain=1:7] (\x,{(0.1*\x^2 })% 
	-- plot[domain=7:1] (\x,{0});% 
	\draw[shift={(5,0)},color=black] (0pt,2pt) -- (0pt,-2pt) node[below] {\footnotesize $x_2$};
	\draw[shift={(3,0)},color=black] (0pt,2pt) -- (0pt,-2pt) node[below] {\footnotesize $x_1$};
	\draw[shift={(1,0)},color=black] (0pt,2pt) -- (0pt,-2pt) node[below] {\footnotesize $x_0$};
	\draw[shift={(7,0)},color=black] (0pt,2pt) -- (0pt,-2pt) node[below] {\footnotesize $x_3$};
	%\draw[color=black] (5.3,1) node {A};
	%\draw[color=red] (1,0) -- (3,0)--(3,0.1)--(1,0.1)--(1,0);
	%\draw[color=red] (3,0) -- (5,0)--(5,0.9)--(3,0.9)--(3,0);
	%\draw[color=red] (5,0) -- (7,0)--(7,2.5)--(5,2.5)--(5,0);
	\fill[color =red,fill opacity=0.3] plot[domain=1:3] (\x,{(0.1 })% 
	-- plot[domain=3:1] (\x,{0});% 
	\fill[color =red,fill opacity=0.3] plot[domain=3:5] (\x,{(0.9 })% 
	-- plot[domain=5:3] (\x,{0});% 
	\fill[color =red,fill opacity=0.3] plot[domain=5:7] (\x,{(2.5 })% 
	-- plot[domain=7:5] (\x,{0});% 
	\end{tikzpicture}\\	
	
\end{center}
Die Fläche unter der Funktion $ f(x) $ ist dannn die Summe von A:\\
\begin{center}
	
	
	\begin{tikzpicture}[xscale=1, yscale=1]
	\draw[color=black] (0.5,0.2) node {A =};
	\draw[color=black] (3.5,0.2) node {+};
	\draw[color=black] (6.5,0.2) node {+};
	\draw[color=black] (9.5,0.2) node {+};
	%\draw[color=red] (1,0) -- (3,0)--(3,0.1)--(1,0.1)--(1,0);
	%\draw[color=red] (4,0) -- (6,0)--(6,0.9)--(4,0.9)--(4,0);
	%\draw[color=red] (7,0) -- (9,0)--(9,2.5)--(7,2.5)--(7,0);
	\fill[color =blue,fill opacity=0.3] plot[domain=1:3] (\x,{(0.1 })% 
	-- plot[domain=3:1] (\x,{0});% 
	\fill[color =blue,fill opacity=0.3] plot[domain=4:6] (\x,{(0.9 })% 
	-- plot[domain=6:4] (\x,{0});% 
	\fill[color =blue,fill opacity=0.3] plot[domain=7:9] (\x,{(2.5 })% 
	-- plot[domain=9:7] (\x,{0});% 
	
	\fill[color =red,fill opacity=0.3] plot[domain=1:3] (\x,{(0.1 })% 
	-- plot[domain=3:1] (\x,{0});% 
	\fill[color =red,fill opacity=0.3] plot[domain=4:6] (\x,{(0.9 })% 
	-- plot[domain=6:4] (\x,{0});% 
	\fill[color =red,fill opacity=0.3] plot[domain=7:9] (\x,{(2.5 })% 
	-- plot[domain=9:7] (\x,{0});% 
	\fill[color =blue,fill opacity=0.3] plot[domain=10:11.5] (\x,{(2.5 })% 
	-- plot[domain=11.5:10] (\x,{0});% 
	\draw[color=black] (10.7,0.2) node {Fehler};
	\end{tikzpicture}\\
\end{center}
\begin{align*}
\mathrm{A}=f(x_0)\cdot (x_1-x_0)+f(x_1)\cdot (x_2-x_1)+f(x_2)\cdot (x_3-x_2)+\mathrm{Fehler}
\end{align*}
Bei Verwendung von viele Flächen wird der Fehler immer kleiner.\\
\begin{center}
	\begin{tikzpicture}[xscale=1, yscale=1]
	
	% Gitternetzlinien
	\draw[xstep=1cm, ystep=1cm, very thin, color=lightgray] (-1,6) grid (8,-1);
	% x-Achse
	\draw[->] (0,0) -- (8,0) node[below,xshift=-1em] {$x$}; 
	%Zahlen auf x-Achse
	%\foreach \x in {-3,...,-1}
	%\draw[shift={(\x,0)},color=black] (0pt,2pt) -- (0pt,-2pt) node[below] {\footnotesize $\x$};
	%\foreach \x in {a,...,7}
	%\draw[shift={(\x,0)},color=black] (0pt,2pt) -- (0pt,-2pt) node[below] {\footnotesize $\x$};
	
	% y-Achse 
	\draw[->] (0,-1) -- (0,5) node[left,yshift=-1em] {$f(x)$};%node[above left]
	%Zahlen auf y-Achse
	%\foreach \y in {-1,...,-1}
	%\draw[shift={(0,\y)},color=black] (2pt,0pt) -- (-2pt,0pt) node[left] {\footnotesize $\y$};
	%\foreach \y in {1,...,5}
	%\draw[shift={(0,\y)},color=black] (2pt,0pt) -- (-2pt,0pt) node[left] {\footnotesize $\y$};
	
	
	%Ursprung
	\draw[color=black] (0pt,-10pt) node[right] {\footnotesize $0$};
	
	\draw[thick,color=black] plot[samples=200, domain=1:7] 
	(\x,{0.1*\x^2}) node[below, xshift=-5em] {$f(x)$};
	
	\fill[color =blue,fill opacity=0.3] plot[domain=1:7] (\x,{(0.1*\x^2 })% 
	-- plot[domain=7:1] (\x,{0});% 
	
	
	\fill[color =red,fill opacity=0.3] plot[domain=1:2] (\x,{(0.1 })% 
	-- plot[domain=2:1] (\x,{0});% 
	\fill[color =red,fill opacity=0.3] plot[domain=2:3] (\x,{(0.4 })% 
	-- plot[domain=3:2] (\x,{0});% 
	\fill[color =red,fill opacity=0.3] plot[domain=3:4] (\x,{(0.9 })% 
	-- plot[domain=4:3] (\x,{0});% 
	\fill[color =red,fill opacity=0.3] plot[domain=4:5] (\x,{(1.6 })% 
	-- plot[domain=5:4] (\x,{0});% 
	\fill[color =red,fill opacity=0.3] plot[domain=5:6] (\x,{(2.5 })% 
	-- plot[domain=6:5] (\x,{0});% 
	\fill[color =red,fill opacity=0.3] plot[domain=6:7] (\x,{(3.6 })% 
	-- plot[domain=7:6] (\x,{0});% 
	\end{tikzpicture}\\	
	
\end{center}
Bei $ n- $Flächen gilt dann allgemein:\\
\begin{align*}
\mathrm{A}&=f(x_0)(x_1-x_0)+f(x_1)(x_2-x_1)+\ldots +f(x_{x_{n-1}})(x_n-x_{x_{n-1}})+\mathrm{Fehler}(n)\\
&=\sum\limits_{k=1}^{n}f(x_{x_{k-1}})(x_k-x_{x_{k-1}})+\mathrm{Fehler}(n)\\
\end{align*}
Man definiert:\\
\begin{align*}
\Delta_{x_k}=(x_k-x_{k-1})
\end{align*}
Dann gilt für die Summe: $ \mathrm{S}_n $\\
\begin{align*}
\mathrm{S}_n=\mathrm{A}=\sum\limits_{k=1}^{n}f(x_{k-1})\Delta_{x_k}+\mathrm{Fehler}(n)
\end{align*}
Dadurch, dass man die Summentteile unendlich klein wählt wird der Fehler Null, dann gilt für $ n\rightarrow\infty $ das bestimmte Integral:\\
\begin{align*}
\int\limits_{a}^{b} f(x)\mathrm{d}x=\lim\limits_ {n\to\infty} \mathrm{S}_n
\end{align*}
Ist die Funktion $ f:[a,b] \longmapsto \mathbb{R}$ monoton oder stetig dann existiert:\\
\begin{align*}
\lim\limits_ {n\to\infty} \mathrm{S}_n=\int\limits_{a}^{b} f(x)\mathrm{d}x
\end{align*}
Beispiel:\\
(i) $ f:[0,1] \longmapsto \mathbb{R} $ mit $ x\longmapsto 1 $\\
\begin{center}
	
	
	\begin{tikzpicture}[xscale=2.5, yscale=2.5]
	
	% Gitternetzlinien
	\draw[xstep=1cm, ystep=1cm, very thin, color=lightgray] (0,3) grid (3,0);
	% x-Achse
	\draw[->] (0,0) -- (3,0) node[below,xshift=-1em] {$x$}; 
	%Zahlen auf x-Achse
	%\foreach \x in {-2,...,-1}
	%\draw[shift={(\x,0)},color=black] (0pt,2pt) -- (0pt,-2pt) node[below] {\footnotesize $\x$};
	\foreach \x in {1,...,2}
	\draw[shift={(\x,0)},color=black] (0pt,2pt) -- (0pt,-2pt) node[below] {\footnotesize $\x$};
	% y-Achse 
	\draw[->] (0,0) -- (0,3) node[left,yshift=-1em] {$f(x)$};%node[above left]
	%Zahlen auf y-Achse
	%\foreach \y in {-1,...,-1}
	%\draw[shift={(0,\y)},color=black] (2pt,0pt) -- (-2pt,0pt) node[left] {\footnotesize $\y$};
	\foreach \y in {1,...,2}
	\draw[shift={(0,\y)},color=black] (2pt,0pt) -- (-2pt,0pt) node[left] {\footnotesize $\y$};
	%Ursprung
	\draw[color=black] (0pt,-5.3pt) node[right] {\footnotesize $0$};
	
	\draw[thick,color=blue] plot[samples=200, domain=0:1] 
	(\x,{1 }) node[below, yshift=2em] {$f(x)$};
	
	\fill[color =red,fill opacity=0.3] plot[domain=0:1] (\x,{(1 })% 
	-- plot[domain=1:0] (\x,{0});% 
	
	
	
	\end{tikzpicture}\\
\end{center}
\begin{align*}
\int\limits_{0}^{1} f(x)=1\cdot 1=1
\end{align*}
\newpage
\noindent
(ii) $ f:[0,2] \longmapsto \mathbb{R} $ mit $ x\longmapsto x $\\
\begin{center}
	\begin{tikzpicture}[xscale=2.5, yscale=2.5]
	
	% Gitternetzlinien
	\draw[xstep=1cm, ystep=1cm, very thin, color=lightgray] (0,3) grid (3,0);
	% x-Achse
	\draw[->] (0,0) -- (3,0) node[below,xshift=-1em] {$x$}; 
	%Zahlen auf x-Achse
	%\foreach \x in {-2,...,-1}
	%\draw[shift={(\x,0)},color=black] (0pt,2pt) -- (0pt,-2pt) node[below] {\footnotesize $\x$};
	\foreach \x in {1,...,2}
	\draw[shift={(\x,0)},color=black] (0pt,2pt) -- (0pt,-2pt) node[below] {\footnotesize $\x$};
	% y-Achse 
	\draw[->] (0,0) -- (0,3) node[left,yshift=-1em] {$f(x)$};%node[above left]
	%Zahlen auf y-Achse
	%\foreach \y in {-1,...,-1}
	%\draw[shift={(0,\y)},color=black] (2pt,0pt) -- (-2pt,0pt) node[left] {\footnotesize $\y$};
	\foreach \y in {1,...,2}
	\draw[shift={(0,\y)},color=black] (2pt,0pt) -- (-2pt,0pt) node[left] {\footnotesize $\y$};
	%Ursprung
	\draw[color=black] (0pt,-5.3pt) node[right] {\footnotesize $0$};
	
	\draw[thick,color=blue] plot[samples=200, domain=0:2] 
	(\x,{\x }) node[below, yshift=2em] {$f(x)$};
	
	\fill[color =red,fill opacity=0.3] plot[domain=0:2] (\x,{(\x })% 
	-- plot[domain=2:0] (\x,{0});% 
	
	
	
	\end{tikzpicture}\\	
\end{center}
\begin{align*}
\int\limits_{0}^{2}f(x) \mathrm{d}x=\int\limits_{0}^{2} x\cdot \mathrm{d}x=\frac{1}{2}2^2=2
\end{align*}
\subsection{Das unbestimmte Integral}
\href{https://youtu.be/D5fcNNAZP2A}{Video 02}\\
Sei $ f:[a,b]\longmapsto \mathbb{R} $ integrierbar. Weiter sei $ x\in[a,b] $.\\
Wir definieren als unbestimmtes Integral von $ f $ auf $ [a,b] $:\\
\begin{align*}
\mathrm{F}(x):=\int\limits_{a}^{x} f(t)\mathrm{d}t
\end{align*}
Wir erhalten eine Abbildung:\\
\begin{align*}
\mathrm{F}:[a,b]\longmapsto \mathbb{R}\;\;\text{mit }\;\;x\longmapsto \int\limits_{a}^{x} f(t)\mathrm{d}t
\end{align*}
Beispiel:\\
Wir betrachten:\\
\begin{align*}
f:[0,10]\longmapsto \mathbb{R}\;\;\text{mit }\;\; f(t)=t\\
\end{align*}
Dann ist:\\
\begin{align*}
\mathrm{F}:[0,10]\longmapsto \mathbb{R}\;\;\text{mit }\;\; \mathrm{F}(x)=\int\limits_{0}^{x}f(t)\mathrm{d}t=\int\limits_{0}^{x} t\cdot \mathrm{d}t=\frac{x^2}{2}
\end{align*}
\begin{center}
	
	
	\begin{tikzpicture}[xscale=0.6, yscale=0.6]
	
	% Gitternetzlinien
	\draw[xstep=1cm, ystep=1cm, very thin, color=lightgray] (0,11) grid (11,0);
	% x-Achse
	\draw[->] (0,0) -- (11,0) node[below,xshift=-1em] {\footnotesize$10]\;\; x$}; 
		\draw (0,0)  node[below,xshift=0 em] {\footnotesize$[ 0$}; 
	%Zahlen auf x-Achse
	%\foreach \x in {-3,...,-1}
	%\draw[shift={(\x,0)},color=black] (0pt,2pt) -- (0pt,-2pt) node[below] {\footnotesize $\x$};
	\foreach \x in {1,...,5}
	\draw[shift={(\x,0)},color=black] (0pt,2pt) -- (0pt,-2pt) node[below] {\footnotesize $\x$};
	\foreach \x in {7,...,9}
	\draw[shift={(\x,0)},color=black] (0pt,2pt) -- (0pt,-2pt) node[below] {\footnotesize $\x$};
	% y-Achse 
	\draw[->] (0,0) -- (0,11) node[left,yshift=-1em] {$f(x)$};%node[above left]
	%Zahlen auf y-Achse
	%\foreach \y in {-1,...,-1}
	%\draw[shift={(0,\y)},color=black] (2pt,0pt) -- (-2pt,0pt) node[left] {\footnotesize $\y$};
	\foreach \y in {1,...,9}
	\draw[shift={(0,\y)},color=black] (2pt,0pt) -- (-2pt,0pt) node[left] {\footnotesize $\y$};
	%Ursprung
	%\draw[color=black] (0pt,-10pt) node[right] {\footnotesize $0$};
	
	\draw[shift={(6,-0.1)},color=black] (0pt,2pt) -- (0pt,-2pt) node[below] {\footnotesize $x$};
	
	\draw[color=black] (3.2,1.2) node {\footnotesize $\int\limits_{0}^{x} f(t)\mathrm{d}t$};
	
	\draw[thick,color=blue] plot[samples=200, domain=0:10] 
	(\x,{\x }) node[below, xshift=-4em] {$f(t)=t$};
	\fill[color =red,fill opacity=0.3] plot[domain=0:6] (\x,{(\x })% 
	-- plot[domain=6:0] (\x,{0});% 
\draw[color=black,thick,dashed] (10,0)--(10,10);
	\end{tikzpicture}\\
	\begin{align*}
	\int\limits_{0}^{x} f(t)\mathrm{d}t=\frac{1}{2}x^2
	\end{align*}
\end{center}
Bemerkung:\\
\begin{align*}
\mathrm{F}'(x)=\bigg(\frac{x^2}{2}\bigg)'=\frac{2x}{2}=x=f(x)\\
\end{align*}
\subsection{Der Hauptsatz der Differential- und Integralrechnung}
\href{https://youtu.be/2v6KdCjGNc0}{Video 03}\\
Sei $ f: \mathbb{I}\longmapsto \mathbb{R} $ stetig auf dem Intervall $  \mathbb{I} $ und $ a\in \mathbb{I} $, dann gilt:\\
(i) Ist F eine Stammfunktion von $ f $, dann gilt:\\
\begin{align*}
\mathrm{F}'(x)=\bigg(\int\limits_{a}^{x}f(t)\mathrm{d}t \bigg)'=f(x)
\end{align*}
(ii)\;Sei $ \mathrm{G}:[a,b]\longmapsto \mathbb{R} $ diffbar mit $ \mathrm{G}'=f(x) $, dann gilt:\\
\begin{align*}
\mathrm{G}=\int\limits_{a}^{x} f(t)\mathrm{d}t+\alpha\;\;\text{mit }\;\;\alpha\in\mathbb{R}
\end{align*}
(iii)\; Sei G eine Stammfunktion von $ f $ und $  \mathbb{I}=[a,b] $, dann ist das Integral:\\
\begin{align*}
\int\limits_{a}^{b} f(x)\mathrm{d}x=\mathrm{G}(b)-G(a)=\bigg[\mathrm{G}(x)\bigg]^b_a
\end{align*}
\newpage
Betracht: $ f:[0,2]\longmapsto\mathbb{R} $ mit $ f(x)=x^2 $ \\
\begin{quote}
	$\blacktriangleright$ Stammfunktion von $ f $\\
	\begin{quote}
		\begin{align*}
		\mathrm{F}'_1=\bigg(\frac{x^3}{3}\bigg)'=\frac{3}{3}\cdot x^2=f(x)
		\end{align*}
		D.h. $ \mathrm{F}_1 $ ist Stammfunktion von $ f $.
		\begin{align*}
		\mathrm{F}'_2=\bigg(\frac{x^3}{3}+12\bigg)'=\frac{3}{3}\cdot x^2=f(x)
		\end{align*}
		$ \mathrm{F}_2 $ ist ebenfals Stammfunktion von $ f $.
	\end{quote}
	$\blacktriangleright\;\;\mathrm{G}_\alpha(x)=\frac{x^3}{3}+\alpha$ 
	\begin{align*}
	\int\limits_{0}^{2} f(x)\mathrm{d}x=\mathrm{F}_2(2)-\mathrm{F}_2(0)=\frac{2^3}{3}+12-\bigg(\frac{0^3}{3}+12\bigg)=\frac{2^3}{3}=\frac{8}{3}
	\end{align*}
	Die Konstante $ \alpha $ hebt sich durch die Differenz auf.\\
	D.h.: Bei der Stammfunktion ist $ \alpha $ bei der Flächenberechnung nicht notwendig.\\
	\begin{align*}
	\int\limits_{0}^{2} f(x)\mathrm{d}x=\int\limits_{0}^{2} x^2\mathrm{d}x=\bigg[\frac{x^3}{3}\bigg]^2_0=\frac{2^3}{3}-\frac{0^3}{3}=\frac{8}{3}
	\end{align*}
\end{quote}
Bemerkung:\\
Aus dem Hauptsatz folgt: Integrieren und Differenzieren sind Gegenoperationen.\\
\subsubsection{Aufgabe 1 zum Hauptsatz der Differential- und Integralrechnung}
\href{https://youtu.be/yEC_tlNWT7E}{Video 04}\\
Berechne das folgende Integral:\\
\begin{align*}
\int\limits_{-1}^{1} (x^2+2x-7x)\mathrm{d}x
\end{align*}
Gesucht: $ \mathrm{F}(x) $ mit $ \mathrm{F}'(x)=f(x)=x^2+2x-7 $\\
\\
NR.:\\
Es gilt für die Ableitung von $ x^n \;\;\text{mit }\;\;n\in\mathbb{N} $:\\
\begin{align*}
\big(x^n\big)'=n\cdot x^{n-1}
\end{align*}
\newpage
Es gilt dann in die andere Richtung: Die Ableitung von F muss $ f(x) $ ergeben.\\
\begin{center}
	
	
	\renewcommand{\arraystretch}{1.4}
	\( \begin{array}{l|c} f(x)&\mathrm{F}\\
	\hline x^2&\frac{x^3}{3}\\
	x^3&\frac{x^4}{4}\\
	x^n&\frac{x^{n+1}}{n+1}
	\end{array} \)
	\begin{align*}
	\Rightarrow \mathrm{F}(x)=\frac{x^{n+1}}{n+1}
	\end{align*}
\end{center}
\begin{align*}
\mathrm{F}(x)= \mathrm{F}'(x)=f(x)&=x^2+2x-7\\
\mathrm{F}(x)&=\frac{x^{2+1}}{2+1}+2\cdot\frac{x^{1+1}}{1+1}-7\cdot\frac{x^{0+1}}{0+1}+\alpha\;\;\text{mit }\;\; \alpha\in\mathbb{R}\\
&=\frac{x^3}{3}+x^2-7x+\alpha\\
&=\bigg[\frac{x^3}{3}+x^2-7x\bigg]^1_{-1}=\frac{1}{3}+1-7-\bigg(-\frac{1}{3}+1+7\bigg)=-\frac{17}{3}-\frac{23}{3}=-\frac{40}{3}\\
\end{align*}
\subsubsection{ Aufgabe 2 zum Hauptsatz der Differential- und Integralrechnung}
\href{https://youtu.be/O-x6XCl7v-c}{Video 05}\\
Berechne folgende Integrale:\\
(i)\;\;$ \int\limits_{0}^{\pi} (\cos(x)+1) $\\
(ii)\;$ \int\limits_{1}^{2} \big(\sqrt{x}+\frac{4}{x}\big)$ \\
\\
zu (i):\\
\begin{align*}
&\int\limits_{0}^{\pi} (\cos(x)+1)\\
&\mathrm{F}(x)=\bigg[\sin(x)+x\bigg]^\pi_0=(\sin(\pi)+\pi)-(\sin(0)+0)=\pi\\
\end{align*}
zu (ii):\\
\begin{quote}
	
	\begin{align*}
	\int\limits_{1}^{2} \bigg(\sqrt{x}+\frac{4}{x}\bigg)
	\end{align*}
	NR.:\\
	\begin{align*}
	\sqrt{x}=x^{\frac{1}{2}}\\
	\end{align*}
	Exponent um 1 erhöhen:\\
	\begin{align*}
	x^{\frac{3}{2}}
	\end{align*}
	Davon die Ableitung:\\
	\begin{align*}
	\bigg(x^{\frac{3}{2}}\bigg)'=\frac{3}{2}\cdot x^{\frac{1}{2}}\\
	\end{align*}
	Die Ableitung soll aber $ x^{\frac{1}{2}} $ sein, deswegen $ x^{\frac{2}{3}} $mit $ \frac{2}{3} $ multiplitieren und dies ergibt dann:\\
	\begin{align*}
	\bigg(\frac{2}{3}x^{\frac{3}{2}}\bigg)'=x^{\frac{1}{2}}
	\end{align*}
\end{quote}
\begin{align*}
\mathrm{F}(x)&=\bigg[\frac{2}{3}x^{\frac{3}{2}}+4\ln(x)\bigg]^2_1\\
&=\frac{2}{3}\sqrt[2]{2^3}+4\ln(2)-\frac{2}{3}+0\\
&=\frac{2}{3}\sqrt{8}+4\ln(2)-\frac{2}{3}
\end{align*}
\subsubsection{ Aufgabe 1 zur Berechnung der Stammfunktion}
\href{https://youtu.be/bMi4SyWZKcg}{Video 06}\\
Bestimme die Stammfunktionen von:\\
\begin{align*}
f(x)=\frac{x^2-2x-3}{x-3}
\end{align*}
D.h.: Berechne alle $\mathrm{F}_\alpha$ mit $ \mathrm{F}'=f(x) $\\
Vom Zähler die Nullstellenberechnen (einfach probieren oder pq-Formel):\\
\begin{align*}
x^2-2x-3&=0\\
x_1&=3\\
x_2&=-1
\end{align*}
Die 3 ist eine Nullstelle und $ (x-3) $ teilt den Zähler.\\
Die -1 ist die zweite Nullstelle. Daraus folgt dann:\\
\begin{align*}
f(x)=\frac{x^2-2x-3}{x-3}=\frac{(x-3)(x+1)}{x-3}=x+1
\end{align*}
Die Stammfunktionen sind dann von $ f(x)=x+1 $:\\
\begin{align*}
\mathrm{F}_\alpha(x)=\frac{x^2}{2}+x+\alpha\;\;\text{mit }\;\; \alpha\in\mathbb{R}
\end{align*}
Prüfe:\
\begin{align*}
\mathrm{F}_\alpha'(x)=\bigg(\frac{x^2}{2}+x+\alpha\bigg)'=\frac{2}{2}x+1=x+1=f(x)
\end{align*}
\subsubsection{ Aufgabe 2 zur Berechnung der Stammfunktion}
\href{https://youtu.be/hunCkb-iq2Q}{Video 07}\\
Bestimme alle Stammfunktionen von:\\
\begin{align*}
g(x) =\sin(x)+\cos(x)
\end{align*}
Wir wissen:\\
\begin{align*}
\cos'(x)&=-\sin(x)\\
\sin'(x)&=\cos(x)
\end{align*}
Die Stammfunktionen sind dann von $ g(x)=\cos(x)+\sin(x) $:\\
\begin{align*}
\mathrm{G}_\alpha=-\cos(x)+\sin(x)+\alpha\;\;\text{mit }\;\; \alpha\in\mathbb{R}
\end{align*}
Prüfe:\\
\begin{align*}
\mathrm{G}'_\alpha=\bigg(-\cos(x)+\sin(x)+\alpha\bigg)'=\sin(x)+\cos(x)=g(x)
\end{align*}
\subsubsection{ Aufgabe 3 zur Berechnung der Stammfunktion}
\href{https://youtu.be/JBw8qcuJTRA}{Video 08}\\
Bestimme alle Stammfunktionen von:\\ 
\begin{align*}
h(x)=\frac{4}{x^3}+x^7\sqrt{x}
\end{align*}
Funktion $ h(x) $ in Potenzen umformen:\\
\begin{align*}
h(x)=4\cdot x^{-3}+x^7\cdot x^{\frac{1}{2}}=4x^{-3}+x^{\frac{15}{2}}
\end{align*}
Gesucht alle $ \mathrm{H}(x) $ mit $ \mathrm{H}'=h(x) $\\
Es gilt mit $ \mathrm{H}(x)=\frac{x^{n+1}}{n+1} $ oder durch Überlegung:\\
\begin{align*}
\mathrm{H}(x)&=4\cdot \frac{x^{(-3+1)}}{-3+1}+\frac{x^{(\frac{15}{2}+1)}}{\frac{15}{2}+1}\\
&=-2x^{-2}+\genfrac{}{}{}{2}{2 }{17}x^{\frac{17}{2}}
\end{align*}
\subsection{Rechenregeln/Eigenschaften   }
\subsubsection{Rechenregeln/Eigenschaften für das Integral I }
\href{https://youtu.be/aP327zIKmvI}{Video 09}\\
Seien $ f,g:[a,b] \longmapsto \mathbb{R}$ integrierbar, dann gilt für alle $ \lambda,\alpha\in\mathbb{R} $:\\
\begin{align*}
\int\limits_{a}^{b}\bigg( \lambda\cdot f(x)+\beta\cdot g(x)\bigg)=\lambda\int\limits_{a}^{b} f(x)\mathrm{d}x+\beta\int\limits_{a}^{b} g(x)\mathrm{d}x
\end{align*}
Beispiel:\\
\begin{align*}
\int\limits_{-1}^{1}(3x^2+7x) \mathrm{d}x&=3\cdot \int\limits_{-1}^{1} x^2\mathrm{d}x+7\cdot\int\limits_{-1}^{1} x\mathrm{d}x\\
&=3\bigg[\frac{x^3}{3}\bigg]^1_{-1}+7\bigg[\frac{x^2}{2}\bigg]^1_{-1}\\
&=3\bigg(\frac{1}{3}-\frac{-1}{3}\bigg)+7\bigg(\frac{1}{2}-\frac{1}{2}\bigg)\\
&=2
\end{align*}
oder:\\
\begin{align*}
\int\limits_{-1}^{1}(3x^2+7x) \mathrm{d}x&=\bigg[x^3+\frac{7}{2}x^2\bigg]^1_{-1}=2
\end{align*}
\newpage
\noindent
\subsubsection{Rechenregeln/Eigenschaften für das Integral II }
\href{https://youtu.be/bCqPDYJzXw0}{Video 10}\\
Seien $ f,g:[a,b]\longmapsto \mathbb{R} $.\\
Sei $ f(x)\leq g(x) $ für alle $ x\in [a,b] $, dann gilt:\\
\begin{align*}
\int\limits_{a}^{b}f(x)\mathrm{d}x\leq\int\limits_{a}^{b} g(x)\mathrm{d}x 
\end{align*}
Beispiel:\\
Es gilt:\\
\begin{align*}
f(x)=x^2\leq x = g(x)\;\;\text{mit }\;\;x\in[0,1]
\end{align*}
Für das Integral gilt dann:\\
\begin{align*}
\int\limits_{0}^{1} f(x)\mathrm{d}x&=\int\limits_{0}^{1} x^2\mathrm{d}x=\bigg[\frac{x^3}{3}\bigg]^1_0=\frac{1}{3}\\
\int\limits_{0}^{1}g(x)\mathrm{d}x&=\int\limits_{0}^{1} x\mathrm{d}x=\bigg[\frac{x^2}{2}\bigg]^1_0=\frac{1}{2}\\
\Rightarrow&\frac{1}{2}>\frac{1}{3}
\end{align*}
\newpage
\subsubsection{Rechenregeln/Eigenschaften für das Integral III }
\href{https://youtu.be/3EPUbiOtG2g}{Video 11}\\
Seien $ f,g:[a,b]\longmapsto \mathbb{R} $.\\
Es gilt:\\
\begin{align*}
\bigg|\int\limits_{a}^{b} f(x)\mathrm{d}x\bigg|\leq \int\limits_{a}^{b}\big|f(x)\big| 
\end{align*}
Beispiel: $ f(x)=\sin(x) $\\

$\bullet$ Für das Integral: $ \bigg|\int\limits_{0}^{2\pi} \sin(x)\mathrm{d}x\bigg| $ gilt:\\

\begin{align*}
\bigg|\int\limits_{0}^{2\pi} \sin(x)\mathrm{d}x\bigg|=\bigg|\bigg[-\cos(x)\bigg]^{2\pi}_0\bigg|=(-\cos(2\pi)-\big(-cos(0)\big)=\big|-1+1\big|=0\\
\end{align*}
\begin{center}
	\begin{tikzpicture}[xscale=1, yscale=2]
	
	% Gitternetzlinien
	\draw[xstep=1cm, ystep=1cm, very thin, color=lightgray] (0,2) grid (7,-2);
	% x-Achse
	\draw[->] (0,0) -- (7,0) node[below,xshift=-1em] {$x$}; 
	%Zahlen auf x-Achse
	%\foreach \x in {-2,...,-6}
	%\draw[shift={(\x,0)},color=black] (0pt,2pt) -- (0pt,-2pt) node[below] {\footnotesize $\x$};
	\foreach \x in {1,...,6}
	\draw[shift={(\x,0)},color=black] (0pt,2pt) -- (0pt,-2pt) node[below] {\footnotesize $\x$};
	% y-Achse 
	\draw[->] (0,-2) -- (0,2) node[left,yshift=-1em] {$f(x)$};%node[above left]
	%Zahlen auf y-Achse
	\foreach \y in {-1,...,-1}
	\draw[shift={(0,\y)},color=black] (2pt,0pt) -- (-2pt,0pt) node[left] {\footnotesize $\y$};
	\foreach \y in {1,...,1}
	\draw[shift={(0,\y)},color=black] (2pt,0pt) -- (-2pt,0pt) node[left] {\footnotesize $\y$};
	%Ursprung
	\draw[color=black] (0pt,-5.3pt) node[right] {\footnotesize $0$};
	
	\draw[thick,color=blue] plot[samples=200, domain=0:6.28] 
	(\x,{(sin(\x r)  }) node[below, yshift=2em] {$\sin(x)$};
	
	\fill[color =red,fill opacity=0.3] plot[domain=0:6.28] (\x,{(sin(\x r) })% 
	-- plot[domain=6.28:0] (\x,{0});% 
	
	
	
	\end{tikzpicture}\\	
\end{center}

$\bullet$ Für das Integral: $ \int\limits_{0}^{2\pi}\big| \sin(x)\mathrm{d}x\big| $ gilt:\\
\begin{align*}
\int\limits_{0}^{2\pi}\big| \sin(x)\mathrm{d}x\big| =2\int\limits_{0}^{\pi} \sin(x)\mathrm{d}x=\bigg[-\cos(x)\bigg]^{\pi}_0=2\big(-\cos(\pi)-(-\cos(0)\big)=2\cdot 2=4
\end{align*}
\begin{center}
	\begin{tikzpicture}[xscale=1, yscale=2]
	
	% Gitternetzlinien
	\draw[xstep=1cm, ystep=1cm, very thin, color=lightgray] (0,2) grid (7,0);
	% x-Achse
	\draw[->] (0,0) -- (7,0) node[below,xshift=-1em] {$x$}; 
	%Zahlen auf x-Achse
	%\foreach \x in {-2,...,-6}
	%\draw[shift={(\x,0)},color=black] (0pt,2pt) -- (0pt,-2pt) node[below] {\footnotesize $\x$};
	\foreach \x in {0,...,6}
	\draw[shift={(\x,0)},color=black] (0pt,2pt) -- (0pt,-2pt) node[below] {\footnotesize $\x$};
	% y-Achse 
	\draw[->] (0,0) -- (0,2) node[left,yshift=-1em] {$f(x)$};%node[above left]
	%Zahlen auf y-Achse
	%\foreach \y in {-1,...,-1}
	%\draw[shift={(0,\y)},color=black] (2pt,0pt) -- (-2pt,0pt) node[left] {\footnotesize $\y$};
	\foreach \y in {1,...,1}
	\draw[shift={(0,\y)},color=black] (2pt,0pt) -- (-2pt,0pt) node[left] {\footnotesize $\y$};
	%Ursprung
	%\draw[color=black] (0pt,-5.3pt) node {\footnotesize $0$};
	
	\draw[thick,color=blue] plot[samples=200, domain=0:3.14] 
	(\x,{(sin(\x r)  }) node[below, yshift=6em] {$|\sin(x)|$};
	\draw[thick,color=blue] plot[samples=200, domain=3.14:6.28] 
	(\x,{(cos((\x+1.57) r)  });
	\fill[color =red,fill opacity=0.3] plot[domain=0:3.14] (\x,{(sin(\x r) })% 
	-- plot[domain=6.28:0] (\x,{0});% 
	
	\fill[color =red,fill opacity=0.3] plot[domain=3.14:6.28] (\x,{(cos((\x+1.57) r) })% 
	-- plot[domain=6.28:3.14] (\x,{0});% 
	
	\end{tikzpicture}\\	
\end{center}
\subsubsection{Rechenregeln/Eigenschaften für das Integral IV u. V }
\href{https://youtu.be/O5sAZWsINYo}{Video 12}\\
Seien $ f,g:[a,b]\longmapsto \mathbb{R} $.\\
Es gilt:\\
$ \bullet \int\limits_{a}^{b} f(x)=\int\limits_{a}^{\alpha}f(x)+\int\limits_{\alpha}^{b}f(x)\mathrm{d}x \;\;\;\; $ mit\;\; $ \alpha\in[a,b] $\\
\begin{center}
	
	
	\begin{tikzpicture}[xscale=1, yscale=1]
	
	% Gitternetzlinien
	\draw[xstep=1cm, ystep=1cm, very thin, color=lightgray] (-1,6) grid (8,-1);
	% x-Achse
	\draw[->] (-1,0) -- (8,0) node[below,xshift=-1em] {$x$}; 
	%Zahlen auf x-Achse
	%\foreach \x in {-3,...,-1}
	%\draw[shift={(\x,0)},color=black] (0pt,2pt) -- (0pt,-2pt) node[below] {\footnotesize $\x$};
	%\foreach \x in {a,...,7}
	%\draw[shift={(\x,0)},color=black] (0pt,2pt) -- (0pt,-2pt) node[below] {\footnotesize $\x$};
	
	% y-Achse 
	\draw[->] (0,-1) -- (0,6) node[left,yshift=-1em] {$f(x)$};%node[above left]
	%Zahlen auf y-Achse
	%\foreach \y in {-1,...,-1}
	%\draw[shift={(0,\y)},color=black] (2pt,0pt) -- (-2pt,0pt) node[left] {\footnotesize $\y$};
	\foreach \y in {1,...,5}
	\draw[shift={(0,\y)},color=black] (2pt,0pt) -- (-2pt,0pt) node[left] {\footnotesize $\y$};
	%Ursprung
	\draw[color=black] (0pt,-10pt) node[right] {\footnotesize $0$};
	
	\draw[thick,color=black] plot[samples=200, domain=1:7] 
	(\x,{0.1*\x^2+1}) node[below, xshift=-5em,yshift=-2em] {$f(x)$};
	
	\fill[color =blue,fill opacity=0.3] plot[domain=1:4] (\x,{(0.1*\x^2+1 })% 
	-- plot[domain=4:1] (\x,{0});% 
	\fill[color =red,fill opacity=0.3] plot[domain=4:7] (\x,{(0.1*\x^2+1 })% 
	-- plot[domain=7:4] (\x,{0});% 
	\draw[shift={(1,0)},color=black] (0pt,2pt) -- (0pt,-2pt) node[below] {\footnotesize $a$};
	\draw[shift={(4,0)},color=black] (0pt,2pt) -- (0pt,-2pt) node[below] {\footnotesize $\alpha$};
	\draw[shift={(7,0)},color=black] (0pt,2pt) -- (0pt,-2pt) node[below] {\footnotesize $b$};
	\draw[color=black] (5.3,1) node {\footnotesize $\int\limits_{\alpha}^{b} f(x)\mathrm{d}x$};
	\draw[color=black] (2.8,1) node {\footnotesize $\int\limits_{a}^{\alpha} f(x)\mathrm{d}x$};
	\end{tikzpicture}\\
\end{center}
Beispiel:\\
\begin{align*}
\int\limits_{0}^{2}x\mathrm{d}x=\int\limits_{0}^{1} x\mathrm{d}x +\int\limits_{1}^{2}x\mathrm{d}x
\end{align*}
\newpage
\noindent
Es gilt:\\
$\bullet  \int\limits_{a}^{b}f(x) \mathrm{d}x=-\int\limits_{b}^{a}f(x)\mathrm{d}x  $ 
Beispiel:\\
\begin{align*}
\int\limits_{0}^{1} 1x=\big[x\big]^1_0=1\\
\int\limits_{1}^{0} 1x=\big[x\big]^0_1=-1\\
\end{align*}
\subsection{Integrationsregeln}
\subsubsection{Partielle Integration }
\href{https://youtu.be/bsTdvvGm7Ic}{Video 13}\\
Seien $ u,v:[a,b]\longmapsto \mathbb{R} $ diffbar und $ u', v' $ stetig, dann gilt:\\
\begin{align*}
\big(u(x)\cdot v(x)\big)'&=u'(x)\cdot v(x) +u(x)\cdot v'(x)\\
\Leftrightarrow \int\limits_{}^{}\big(u(x)\cdot v(x)\big)'\mathrm{d}x  &=\int\limits_{}^{}u'(x)\cdot v(x)+\int\limits_{}^{}u(x)\cdot v'(x)  \\
\end{align*}
$ \Leftrightarrow $ Achtung: Die Äquifallenz ist hier wegen den Konstanten nicht korrekt.\\
\begin{align*}
\int\limits_{}^{}\big(u(x)\cdot v(x)\big)'\mathrm{d}x  &=\int\limits_{}^{}u'(x)\cdot v(x)+\int\limits_{}^{}u(x)\cdot v'(x)  \\
u(x)\cdot v(x)  &=\int\limits_{}^{}u'(x)\cdot v(x)+\int\limits_{}^{}u(x)\cdot v'(x)  \\
\Leftrightarrow u(x)\cdot v(x)-\int\limits_{}^{}u'(x)\cdot v(x)&=\int\limits_{}^{}u(x)\cdot v'(x)  \\
\end{align*}
Für die partielle Integration gilt dann mit Integrationsgrenzen:\\
\begin{center}
	\fbox{\parbox  {80mm} {$\big[u(x)\cdot v(x)\big]^b_a-\int\limits_{a}^{b}u'(x)\cdot v(x)=\int\limits_{a}^{b}u(x)\cdot v'(x)   $}}
\end{center}
\subsubsection{Beispiel zur partiellen Integration }
\href{https://youtu.be/92SY1mE5hrMvv}{Video 14}\\
Berechne das Integral:\\
\begin{align*}
\int\limits_{0}^{1} x \mathrm{e}^x\mathrm{d}x
\end{align*}
Man wählt vorteilhaft $ u(x), v'(x) $:\\
\begin{align*}
u(x)=x&\Rightarrow u'(x)=1\\
v'(x)=\mathrm{e}^x&\Rightarrow v(x)=e^x
\end{align*}
Nun gilt:\\
\begin{align*}
\int\limits_{a}^{b}u(x)\cdot v'(x) =\big[u(x)\cdot v(x)\big]^b_a-\int\limits_{a}^{b}u'(x)\cdot v(x) 
\end{align*}
Dann einsetzen von $ u(x), u'(x), v(x), v'(x) $:\\
\begin{align*}
\int\limits_{0}^{1} x\mathrm{e}^x&=\bigg[x\cdot \mathrm{e}^x\bigg]^1_0-\int\limits_{a}^{b} 1\cdot\mathrm{e}^x \mathrm{d}x\\
&=\bigg[x\cdot \mathrm{e}^x\bigg]^1_0-\bigg[\mathrm{e}^x\bigg]_0^1\\
&=\mathrm{e}-0\mathrm{e}^0-(\mathrm{e}-1)\\
&=\mathrm{e}-\mathrm{e}+1=1
\end{align*}
Damit man sich die Formel nicht merken muß, kann man sich merken:\\
Bilde $ u(x) $ und $ v'(x) $ und notiere dies tabelarisch:\\
\begin{align*}
u(x)&\Rightarrow u'(x)\\
v'(x)&\Rightarrow v(x)
\end{align*}
Die Diagonale ergibt den Term: $\big[u(x)\cdot v(x)\big]^b_a  $\\
Die zweit Spalte ergibt den Term:$ \int\limits_{a}^{b}u'(x)\cdot v(x) $
\subsubsection{Trick I zur partiellen Integration }
\href{https://youtu.be/EM79gyqz1yo}{Video 15}\\
Bestimme alle Stammfunktionen von:\\
\begin{align*}
f(x)=x^2\cdot \mathrm{e} ^x
\end{align*}
Berechne:\\
\begin{align*}
\mathrm{F}(x)=\int\limits_{}^{} f(x)\mathrm{d}x=\int\limits_{}^{} x^2\cdot \mathrm{e}^x\mathrm{d}x
\end{align*}
NR.:\\
\begin{align*}
u(x)=x^2&\Rightarrow u'(x)=2x\\
v'(x)=\mathrm{e}^x&\Rightarrow v(x)=e^x\\
\end{align*}
\begin{align*}
\int\limits_{}^{}x^2\cdot \mathrm{e}^x\mathrm{d}x&= x^2\cdot \mathrm{e}^x-\int\limits_{}^{} 2x\mathrm{e}^x\mathrm{d}x\\
&=x^2\cdot \mathrm{e}^x-2\int\limits_{}^{} x\cdot \mathrm{e}^x\mathrm{d}x\\
\end{align*}
Berechnung vom Integral $\int\limits_{}^{}x\cdot \mathrm{e}^x\mathrm{d}x  $:\\
\begin{align*}
u(x)=x&\Rightarrow u'(x)=1\\
v'(x)=\mathrm{e}^x&\Rightarrow v(x)=\mathrm{e}^x
\end{align*}
\begin{align*}
\int\limits_{}^{} x\cdot \mathrm{e}^x=x \cdot \mathrm{e}^x-\int\limits_{}^{}\mathrm{e}^x\mathrm{d}x=x\cdot \mathrm{e}^x-\mathrm{e}^x 
\end{align*}
Dies ergibt dann eingesetzt:\\
\begin{align*}
\int\limits_{}^{} x^2\cdot \mathrm{e}^x\mathrm{d}x&=x^2\cdot \mathrm{e}^x-2(x\cdot \mathrm{e}^x-\mathrm{e}^x)\\
&=x^2\cdot \mathrm{e}^x-2\cdot(x-1)\\
&=x^2\cdot \mathrm{e}^2-2\cdot\mathrm{e}^x\cdot x+2\cdot \mathrm{e}^x\\
&=\mathrm{e}^x(x^2-2x+2)
\end{align*}
Es reicht wenn man jetzt erst mit den Konstanten rechnet um alle Stammfunktionen zu bekommen:\\
Fazit:\\
Alle Stammfunktionen von $ f(x) $ sind:\\
\begin{align*}
\mathrm{F}_\alpha=\mathrm{e}^x(x^2-2x+2) +\alpha\;\;\text{mit }\;\;\alpha\in\mathbb{R}
\end{align*}





\subsubsection{Trick II zur partiellen Integration}
\href{https://youtu.be/TbMWPRG4Ffw}{Video 16}\\
Bestimme alle Stammfunktionen von:\\
\begin{align*}
\int\limits_{}^{}\ln(x)\mathrm{d}x
\end{align*}
\underline{Trick:}\\
\begin{align*}
\int\limits_{}^{} \ln(x)\mathrm{d}x=\int\limits_{}^{} 1\cdot \ln(x)\mathrm{d}x
\end{align*}
NR.:\
\begin{align*}
u(x)=\ln(x)&\Rightarrow u'(x)=\frac{1}{x}\\
v(x)=1&\Rightarrow v'(x)=x
\end{align*}
Dies ergibt dann:\\
\begin{align*}
\int\limits_{}^{} \ln(x)&=u(x) \cdot v'(x)-\int\limits_{}^{} u'(x)\cdot v'(x)\mathrm{d}x\\
&=x\cdot \ln(x)-\int\limits_{}^{} \frac{1}{x}\cdot x\mathrm{d}x\\
&=x\cdot \ln(x)-\int\limits_{}^{}1\mathrm{d}x\\
&=x\cdot \ln(x)-x\\
&=x(\ln(x)-1)
\end{align*}
Probe:\\
\begin{align*}
\bigg(x(\ln(x)-1)\bigg)'=1\cdot (\ln(x)-1)+x\cdot \frac{1}{x}=\ln(x)-1+1=\ln(x)\\
\end{align*}
\subsubsection{Trick III zur partiellen Integration}
\href{https://youtu.be/sKeTwFd3Gfw}{Video 17}\\
Bestimme alle Stammfunktionen von:\\
\begin{align*}
h(x)=\cos(x)\cdot \sin(x)
\end{align*}
Berechne das Integral von:\\
\begin{align*}
\int\limits_{}^{} h(x)\mathrm{d}x=\int\limits_{}^{} \cos(x)\cdot \sin(x)\mathrm{d}x
\end{align*}
NR.:\\
\begin{align*}
u(x)=\cos(x)&\Rightarrow u'(x)=-sin(x)\\
v'(x)=\sin(x)&\Rightarrow v(x)=-\cos(x)\\
\end{align*}
Dies ergibt dann:\\
\begin{align*}
\int\limits_{}^{} h(x)\mathrm{d}x=\int\limits_{}^{}\cos(x)\cdot \sin(x)\mathrm{d}x=\cos(x)\cdot -\cos(x)-\int\limits_{}^{} \cos(x)\cdot\sin(x)\mathrm{d}x\\
\end{align*}
Jetzt ist man nicht viel weiter gekommen. Mit Addition von $\int\limits_{}^{} \cos(x)\cdot\sin(x)\mathrm{d}x  $ links und rechts der Gleichung ergibt dann:\\
\begin{align*}
\int\limits_{}^{}\cos(x)\cdot \sin(x)\mathrm{d}x+\int\limits_{}^{} \cos(x)\cdot\sin(x)\mathrm{d}x& =-\cos^2(x)-\int\limits_{}^{} \cos(x)\cdot\sin(x)\mathrm{d}x+\int\limits_{}^{} \cos(x)\cdot\sin(x)\mathrm{d}x\\
2\int\limits_{}^{} \cos(x)\cdot \sin(x)\mathrm{d}x&=-\cos^2(x)\\
\Leftrightarrow \int\limits_{}^{} \cos(x)\cdot \sin(x)\mathrm{d}x &=-\frac{1}{2} \cos^2(x)
\end{align*}
Fazit:\\
Alle Stammfunktionen sind in der Form:\\
\begin{align*}
\int\limits_{}^{} \cos(x)\cdot \sin(x)\mathrm{d}x &=-\frac{1}{2} \cos^2(x)+\alpha \;\;\text{mit }\;\;\alpha\in\mathbb{R}
\end{align*}
\subsection{Substitution}
\subsubsection{Substitutionsregel}
\href{https://youtu.be/Gfg8adsiy84}{Video 18}\\
Sei F eine Stammfunktion der stetigen Funktion $ f $ und $ x(t) $ eine diffbare Funktion sowie $ x'(t) $ stetig.\\
\begin{align*}
\mathrm{F}(x(t))
\end{align*}
Bemerkung:\\
Die Ableitung von$ f'(x) =\frac{d}{\mathrm{d}x}f(x)$\\
Dann ergibt sich, wenn man $ \mathrm{F}(x(t)) $ ableitet:\\
\begin{align*}
\frac{d}{\mathrm{d}t}\big(\mathrm{F}(x(t))\big)=\mathrm{F}'(x(t))\cdot x'(t)=f(x(t))\cdot x'(t)
\end{align*}
Integriert man nun beide Seiten:\\
\begin{align*}
\int\limits_{}^{} \frac{d}{\mathrm{d}t}(\mathrm{F}(x(t)))\mathrm{d}t=\int\limits_{}^{} f(x(t))\cdot x'(t)\mathrm{d}t
\end{align*}
Das Integrieren und Differenzieren neutralisieren sich.\\
\begin{align*}
\mathrm{F}(x(t))=\int\limits_{}^{} f(x(t))\cdot x'(t)\mathrm{d}t
\end{align*}
Mit Grenzen gilt dann:\\
\begin{align*}
\mathrm{F}(x(b))-\mathrm{F}(x(a))=\int\limits_{a}^{b} f(x(t))\cdot x'(t)\mathrm{d}t
\end{align*}
\subsubsection{Beispiel zur Substitution}
\href{https://youtu.be/JLzCF3VnqkU}{Video 19}\\
Berechne:\\
\begin{align*}
\int\limits_{0}^{\frac{\pi}{2}}  \sin(2x) \mathrm{d}x
\end{align*}
NR.: Substition mit:
\begin{align*}
t=2x
\end{align*}
Dann erhält man:\\
\begin{align*}
\int\limits_{}^{} \sin(t)\mathrm{d}x
\end{align*}
Nun muß noch $ \mathrm{d}x $ von $ \mathrm{d}t $ abhängig werden.\\
\begin{align*}
t=2x\Rightarrow 
\end{align*}
Durch Ableiten nach $ \frac{\mathrm{d}t}{\mathrm{d}x} $ erhält man $ \mathrm{d}x $:\\
\begin{align*}
\frac{\mathrm{d}t}{\mathrm{d}x}&=\frac{\mathrm{d}2x}{\mathrm{d}x}=2\\
\Leftrightarrow\frac{\mathrm{d}t}{2}&=\mathrm{d}x
\end{align*}
So bekommt man das lösbare Integral:\\
\begin{align*}
\int\limits_{}^{} \sin(t)\frac{\mathrm{d}t}{2}
\end{align*}
Die neuen Grenzen sind dann:\\
Obergrenze mit $ x=\frac{\pi}{2} $:\\
\begin{align*}
t&=2x\\
t=&2\cdot \frac{\pi}{2}=\pi
\end{align*}
Untergrenze mit $x= 0 $:\\
\begin{align*}
t&=2x\\
t=&2\cdot 0=0
\end{align*}
Die neue Stammfunktion ist dann:\\
\begin{align*}
\int\limits_{0}^{\pi}\sin(t) \frac{\mathrm{d}t}{2}=\frac{1}{2}\bigg[-\cos(t)\bigg]^\pi_0=\frac{1}{2}\big(-\cos(\pi)-(-\cos(0)\big)=\frac{1}{2}\big(+1-(-1)\big)=\frac{2}{2}=1\\
\end{align*}
Jetzt die Berechnung, so dass man die Substitutionsregel erkennt:\\
NR.: Substitution mit:\\
\begin{align*}
x(t)=\frac{t}{2}\Rightarrow& \frac{\mathrm{d}x}{\mathrm{d}t}=\frac{\mathrm{d}t}{2\mathrm{d}t}=\frac{1}{2}\\
\Leftrightarrow&\mathrm{d}x(t)=\frac{\mathrm{d}t}{2}
\end{align*}
Dann erhält man durch einsetzen in $\int\limits_{}^{} \sin(2\cdot x(t))\mathrm{d}x(t) \mathrm{d}t=\int\limits_{}^{} \sin(2\cdot x(t))x' (t) \mathrm{d}t $:\\
\begin{align*}
\int\limits_{}^{} \sin\bigg(2\cdot \frac{t}{2}\bigg)\cdot \frac{\mathrm{d}t}{2}=\int\limits_{}^{} \sin(t)\frac{\mathrm{d}t}{2}\\
\end{align*}
Untere Grenze:\\
\begin{align*}
x=0=\frac{t}{2}=0
\end{align*}
Obere Grenze:\\
\begin{align*}
x=\frac{\pi}{2}=\frac{t}{2}\Rightarrow t=\pi
\end{align*}
Bemerkung:\\
Zuerst wurde $ t=2x $ substituiert und dies ist die Umkehrfunktion von $ x(t)=\frac{t}{2} $.
\subsubsection{Trick 1 zur Substitutionsregel Aufgabe 1}
\href{https://youtu.be/IjIGyA63qtU}{Video 20}\\
Berechne:\\
\begin{align*}
\int\limits_{1}^{2} \frac{x}{(4x-1)^3}\mathrm{d}x
\end{align*}
(Lineare) Substitution:\\
\begin{align*}
t=4x-1
\end{align*}
Nach $ x $ ableiten:\\
\begin{align*}
\frac{\mathrm{d}t}{\mathrm{d}x}&=\frac{d(4x-1)}{\mathrm{d}x}=4\\
\Leftrightarrow \frac{\mathrm{d}t}{4}&=\mathrm{d}x
\end{align*}
Dann erhält man das Integral:\\
\begin{align*}
\int\limits_{}^{} \frac{x}{t^3}\cdot \frac{\mathrm{d}t}{4}
\end{align*}
Das $ x $ durch $ t $ ersetzen:\\
\begin{align*}
t&=4x-1\\
\Leftrightarrow x&=\frac{t+1}{4}
\end{align*}
Dann erhält man:\\
\begin{align*}
\int\limits_{}^{} \frac{\frac{t+1}{4}}{t^3}\cdot \frac{\mathrm{d}t}{4}=\frac{1}{16}\int\limits_{}^{} \frac{t+1}{t^3}\mathrm{d}t
\end{align*}
Die neuen Grenzen berechnen:\\
\begin{quote}
	Untere Grenze:\\
	\begin{align*}
	x=1\Rightarrow t=4x-1=4\cdot 1-1=3
	\end{align*}
	Obere Grenze:\\
	\begin{align*}
	x=2\Rightarrow t=4x-1=4\cdot 2-1=7
	\end{align*}
\end{quote}
Das Integral ist dann:\\
\begin{align*}
\frac{1}{16}\int\limits_{3}^{7} \frac{t+1}{t^3}\mathrm{d}t&=\frac{1}{16}\int\limits_{3}^{7} \bigg(\frac{t}{t^3}+\frac{1}{t^3}\bigg)=\frac{1}{16}\int\limits_{3}^{7} \bigg(\frac{1}{t^2}+\frac{1}{t^3}\bigg)\\
&=\frac{1}{16}\bigg[-\frac{1}{t}-\frac{1}{2t^2}\bigg]^7_3\\
&=\frac{1}{16}\bigg[-\frac{1}{7}-\frac{1}{2\cdot 49}-\bigg(-\frac{1}{3}-\frac{1}{2\cdot 9}\bigg)\bigg]\\
&\cdots
\end{align*}
\subsubsection{Trick 1 zur Substitutionsregel Aufgabe 2}
\href{https://youtu.be/tlmzM6q92Lc}{Video 21}\\
Berechne:\\
\begin{align*}
\int\limits_{0}^{1} \frac{1}{(x-2)^4}\mathrm{d}x 
\end{align*}
Substitution:\\
\begin{align*}
t=x-2
\end{align*}
Nach $ x $ ableiten:\\
\begin{align*}
\frac{\mathrm{d}t}{\mathrm{d}x}&=\frac{d(x-2)}{\mathrm{d}x}=1\\
\Leftrightarrow \mathrm{d}t&=\mathrm{d}x
\end{align*}
Dann erhält man:\\
\begin{align*}
\int\limits_{}^{} \frac{1}{t^4}\mathrm{d}t
\end{align*}
Die neuen Grenzen berechnen:\\
\begin{quote}
	Untere Grenze:\\
	\begin{align*}
	x=0\Rightarrow t=x-2= 0-2=-2
	\end{align*}
	Obere Grenze:\\
	\begin{align*}
	x=1\Rightarrow t=x-2= 1-2=-1
	\end{align*}
\end{quote}
Das Integral ist dann:\\
\begin{align*}
\int\limits_{-2}^{-1} \frac{1}{t^4}\mathrm{d}t&=\bigg[-\frac{1}{3}\cdot\frac{1}{t^3}\bigg]^{-1}_2\\
&=-\frac{1}{3}\cdot\frac{1}{(-1)^3}-\bigg(-\frac{1}{3}\cdot\frac{1}{(-2)^3}\bigg)=\frac{1}{3}-\frac{1}{24}=\frac{8-1}{24}=\frac{7}{24}
\end{align*}
\subsubsection{Trick 2 zur Substitutionsregel}
\href{https://youtu.be/ECFaj8YJlws}{Video 22}\\
Berechne:\\
\begin{align*}
\int\limits_{0}^{\frac{\pi}{2}}\frac{\cos^3(x)}{1-\sin(x)} \mathrm{d}x
\end{align*}
Substitution:\\
\begin{align*}
t=&\sin(x)\\
\frac{\mathrm{d}t}{\mathrm{d}x}&=\frac{\mathrm{d}(\sin(x))}{\mathrm{d}x}=\cos(x)\\
\Leftrightarrow&\frac{\mathrm{d}t}{\cos(x)}=\mathrm{d}x\\
\end{align*}
Dann erhält man:\\
\begin{align*}
\int\limits_{}^{} \frac{\cos^3(x)}{1-t}\cdot\frac{\mathrm{d}t}{\cos(x)}=\int\limits_{}^{} \frac{\cos^2(x)}{1-t}\mathrm{d}t
\end{align*}
Es gilt:\\
\begin{align*}
&\cos^2(x)+\sin^2(x)=1\\
&\Leftrightarrow \cos^2(x)=1-\sin^2(x)
\end{align*}
Mit $ t=\sin(x) $ gilt dann:\\
\begin{align*}
\cos^2(x)=1-t^2
\end{align*}
Dann erhält man:\\
\begin{align*}
\int\limits_{}^{} \frac{1-t^2}{1-t}\mathrm{d}t
\end{align*}
Die neuen Grenzen berechnen:\\
\begin{quote}
	Untere Grenze:\\
	\begin{align*}
	x=0\Rightarrow t=\sin(0)=0
	\end{align*}
	Obere Grenze:\\
	\begin{align*}
	x=\frac{\pi}{2}\Rightarrow t=\sin\bigg(\frac{\pi}{2}\bigg)=1
	\end{align*}
\end{quote}
Das Integral ist dann:\\
\begin{align*}
\int\limits_{0}^{1} \frac{1-t^2}{1-t}\mathrm{d}t&=\int\limits_{0}^{1} \frac{(1-t)(1+t)}{1-t}\mathrm{d}t=\int\limits_{0}^{1}(1+t)\mathrm{d}t\\ 
&=\bigg[t+\frac{t^2}{2}\bigg]^1_0=\bigg(1+\frac{1}{2}\bigg)=\frac{3}{2}
\end{align*}
\subsubsection{Trick 3 zur Substitutionsregel}
\href{https://youtu.be/wque3ru-qwM}{Video 23}\\
Berechne:\\
\begin{align*}
\int\limits_{0}^{2} \frac{\sqrt{x^2+16}}{2}\cdot \mathrm{d}x
\end{align*}
Konstante herausziehen:\\
\begin{align*}
=\frac{1}{2}\int\limits_{0}^{2} \sqrt{x^2+16}\cdot \mathrm{d}x
\end{align*}
Substitution:\\
\begin{align*}
x&=4\sinh(t)\\
\frac{\mathrm{d}x}{\mathrm{d}t}&=\frac{\mathrm{d}(4\sinh(t))}{\mathrm{d}t}=4\cosh(t)\\
\Leftrightarrow&\mathrm{d}x=4\cosh(t)\mathrm{d}t
\end{align*}
Dann erhält man:\\
\begin{align*}
&=\frac{1}{2}\int\limits_{}^{}\sqrt{4^2\sinh^2(t)+4^2}\cdot 4\cosh(t)\mathrm{d}t\\
&=\frac{1}{2}\int\limits_{}^{} \sqrt{4^2(\sinh^2(t)+1)}\cdot 4\cosh(t)\mathrm{d}t\\
&=\frac{1}{2}\int\limits_{}^{} 4\sqrt{\sinh^2(t)+1}\cdot 4\cosh(t)\mathrm{d}t\\
&=8\int\limits_{}^{} \sqrt{\sinh^2(t)+1}\cdot \cosh(t)\mathrm{d}t\\
\end{align*}
Es gilt:\\
\begin{align*}
&\cosh^2(t)-\sinh^2(t)=1\\
\Leftrightarrow & \cosh^2(t)=\sinh^2(t)+1
\end{align*}
Dann erhält man:\\
\begin{align*}
&=8\int\limits_{}^{} \sqrt{\cosh^2(t)}\cdot \cosh(t)\mathrm{d}t\\
&=8\int\limits_{}^{} \cosh^2(t)\mathrm{d}t
\end{align*}
Die neuen Grenzen berechnen:\\
\begin{quote}
	Untere Grenze:\\
	\begin{align*}
	x=0\Rightarrow 0=4\sinh(t) \Rightarrow t=0\\
	\end{align*}
	Obere Grenze:\\
	\begin{align*}
	&x=2\Rightarrow 2=4\sinh(t)\\
	&\Leftrightarrow \frac{1}{2}=\sinh(t)\\
	&\Leftrightarrow t=\arcsin\bigg(\frac{1}{2}\bigg)
	\end{align*}
\end{quote}
Das Integral ist dann:\\
\begin{align*}
&=8\int\limits_{0}^{\arcsin\big(\frac{1}{2}\big)} \cosh^2(t)\mathrm{d}t
\end{align*}
Es gilt:\\
\begin{align*}
\cosh(t)=\frac{\mathrm{e}^t+\mathrm{e}^{-t}}{2}
\end{align*}
Dann ergibt sich:\\
\begin{align*}
=8\int\limits_{0}^{\arcsin\big(\frac{1}{2}\big)}\bigg(\frac{e^t+\mathrm{e}^{-t}}{2}\bigg )^2 \mathrm{d}t&=8\int\limits_{0}^{\arcsin\big(\frac{1}{2}\big)}\frac{1}{4}\bigg(\mathrm{e}^t+\mathrm{e}^{-t}\bigg )^2\mathrm{d}t\\
&=8\int\limits_{0}^{\arcsin\big(\frac{1}{2}\big)}\frac{1}{4}\big(\mathrm{e}^{2t}+2\underbrace {\mathrm{e}^t\mathrm{e}^{-t}}_{=1}+\mathrm{e}^{-2t}\big)\mathrm{d}t\\
&=8\int\limits_{0}^{\arcsin\big(\frac{1}{2}\big)}\bigg(\frac{\mathrm{e}^{2t}}{4}+\frac{1}{2}+\frac{\mathrm{e}^{-2t}}{4}\bigg)\mathrm{d}t\\
&=\int\limits_{0}^{\arcsin\big(\frac{1}{2}\big)}\big(2\mathrm{e}^{2t}+4+2\mathrm{e}^{-2t}\big)\mathrm{d}t\\
&=\bigg[\mathrm{e}^{2t}+4t-\mathrm{e}^{-2t}\bigg]^{\arcsin\big(\frac{1}{2}\big)}_0\\
&=\mathrm{e}^{2\arcsin\big(\frac{1}{2}\big)}+4\arcsin\bigg(\frac{1}{2}\bigg)-\mathrm{e}^{-2\arcsin\big(\frac{1}{2}\big)}-0\\
&\text{mit Taschenrechner ausrechnen. }
\end{align*}
\subsubsection{Trick 4 zur Substitutionsregel}
\href{https://youtu.be/xmWyiH_de9k}{Video 24}\\
Berechne:\\
\begin{align*}
\int\limits_{2}^{3}\sqrt{x^2-1}\cdot \mathrm{d}x 
\end{align*}
Substitution:\\
\begin{align*}
x&=\cosh(t)\\
\frac{\mathrm{d}x}{\mathrm{d}t}&=\frac{\mathrm{d}(\cosh(t))}{\mathrm{d}t}=\sinh(t)\\
\Leftrightarrow&\mathrm{d}x=\sinh(t)\mathrm{d}t
\end{align*}
Dann erhält man:\\
\begin{align*}
=\int\limits_{}^{} \sqrt{\cosh^2(t)}\cdot\sinh(t)\mathrm{d}t
\end{align*}
Es gilt:\\
\begin{align*}
&\cosh^2(t)-\sinh^2(t)=1\\
\Leftrightarrow & \cosh^2(t)-1=\sinh^2(t)
\end{align*}
Dann erhält man:\\
\begin{align*}
&=\int\limits_{}^{} \sqrt{\sinh^2(t)}\cdot \sinh(t)\mathrm{d}t\\
&=\int\limits_{}^{} \sinh^2(t)\mathrm{d}t
\end{align*}
Die neuen Grenzen berechnen:\\
\begin{quote}
	Untere Grenze:\\
	\begin{align*}
	&x=2\Rightarrow 2=\cosh(t)\\
	&\Leftrightarrow t=\arccos(2)
	\end{align*}
	Obere Grenze:\\
	\begin{align*}
	&x=3\Rightarrow 3=\cosh(t)\\
	&\Leftrightarrow t=\arccos(3)
	\end{align*}
\end{quote}
Das Integral ist dann:\\
\begin{align*}
=\int\limits_{\arccos(2)}^{\arccos(3)}\sinh(t) \mathrm{d}t
\end{align*}
Es gilt:\\
\begin{align*}
\sinh(t)=\frac{\mathrm{e}^t-\mathrm{e}^{-t}}{2}
\end{align*}
Dann ergibt sich:\\
\begin{align*}
=\int\limits_{\arccos(2)}^{\arccos(3)}\bigg(\frac{e^t-e^{-t}}{2}\bigg)^2&=\frac{1}{4}\int\limits_{\arccos(2)}^{\arccos(3)}\big(e^t-e^{-t}\big)^2\\
&=\frac{1}{4}\int\limits_{\arccos(2)}^{\arccos(3)}\bigg(e^{2t}-2e^te^{-t}+e^{-2t}\bigg)\mathrm{d}t\\
&=\frac{1}{4}\int\limits_{\arccos(2)}^{\arccos(3)}\bigg(e^{2t}-2+e^{-2t}\bigg)\mathrm{d}t\\
&=\frac{1}{4}\bigg[\frac{e^{2t}}{2}-2t-\frac{e^{-2t}}{2}\bigg]^{\cosh(3)}_{\cosh(2)}\\
&=\frac{1}{4}\bigg[\frac{e^{2t}-e^{-2t}}{2}-2t\bigg]^{\cosh(3)}_{\cosh(2)}\\
&=\frac{1}{4}\bigg[\sinh(2t)-2t\bigg]^{\cosh(3)}_{\cosh(2)}\\
&=\bigg[\frac{\sinh(2t)}{4}-\frac{t}{2}\bigg]^{\cosh(3)}_{\cosh(2)}\\
\end{align*}
\subsubsection{Trick 5 zur Substitutionsregel}
\href{https://youtu.be/AvuezkKpOus}{Video 25}\\

Berechne:\\
\begin{align*}
\int\limits_{-1}^{1} \sqrt{1-x^2}\mathrm{d}x
\end{align*}
Substitution:\\
\begin{align*}
x&=\sin(t)\\
\frac{\mathrm{d}x}{\mathrm{d}t}&=\frac{\mathrm{d}(\sin(t))}{\mathrm{d}t}=\cos(t)\\
\Leftrightarrow&\mathrm{d}x=\cos(t)\mathrm{d}t
\end{align*}
Dann erhält man:\\
\begin{align*}
&=\int\limits_{}^{} \sqrt{1-\sin^2(t)}\cdot \cos(t)\mathrm{d}t\\
\end{align*}
Die neuen Grenzen berechnen:\\
\begin{quote}
	Untere Grenze:\\
	\begin{align*}
	&x=-1\Rightarrow -1=\sin(t)\Rightarrow t=-\frac{\pi}{2}\\
	\end{align*}
	Obere Grenze:\\
	\begin{align*}
	&x=1\Rightarrow 1=\sin(t)\Rightarrow t=\frac{\pi}{2}\\
	\end{align*}
\end{quote}
Das Integral ist dann:\\
\begin{align*}
=\int\limits_{-\frac{\pi}{2}}^{\frac{\pi}{2}} \sqrt{1-\sin^2(t)}\cdot \cos(t)\mathrm{d}t
\end{align*}
Es gilt:\\
\begin{align*}
&\sin^2(t)+\cos^2 (t)=1\\
&\Leftrightarrow1-\sin(t)=\cos^2(t)
\end{align*}
Dann erhält man:\\
\begin{align*}
&=\int\limits_{-\frac{\pi}{2}}^{\frac{\pi}{2}}\sqrt{\cos^2(t)}\cdot \cos(t)\mathrm{d}t\\
&=\int\limits_{-\frac{\pi}{2}}^{\frac{\pi}{2}} \cos^2(t)\mathrm{d}t
\end{align*}
Es gilt mit dem Additionstheorem:\\
\begin{align*}
\cos (2t)&=\cos^2(t)-\sin^2(t)\;\;\;\;\;\;\;\;\;\;\lvert \text{ \tiny +1}\\
\cos (2t)+1&=\cos^2(t)-\sin^2(t)+1\;\;\;\;\lvert \text{\tiny$ \sin^2(t)+\cos^2 (t)=1 $ }\\
\cos (2t)+1&=\cos^2(t)-\sin^2(t)+\sin^2(t)+\cos^2 (t)\\
\cos (2t)+1&=2\cos^2(t)\\
\frac{\cos(2t)+1}{2}&=\cos^2(t)
\end{align*}
Dann ergibt sich :\\
\begin{align*}
&=\int\limits_{-\frac{\pi}{2}}^{\frac{\pi}{2}}\bigg( \frac{1}{2}\cos(2t)+1\bigg)\mathrm{d}t\\
&=\frac{1}{2}\int\limits_{-\frac{\pi}{2}}^{\frac{\pi}{2}}\big(\cos(2t)+1\big)\mathrm{d}t\\
&=\frac{1}{2}\bigg[\frac{1}{2}\sin(2t)+t\bigg]^{\frac{\pi}{2}}_{-\frac{\pi}{2}}\\
&=\frac{1}{2}\bigg(\frac{1}{2}\sin\big(\genfrac{}{}{}{2}{2\pi}{2 }\big)+\frac{\pi}{2}-\bigg(\frac{1}{2}\sin\big(\genfrac{}{}{}{2}{2\pi}{2 })-\frac{\pi}{2}\bigg)\bigg)\\
&=\frac{1}{2}\bigg(0+\frac{\pi}{2}-\bigg(0-\frac{\pi}{2}\bigg)\bigg)=\frac{1}{2}\big(\pi\big)=\frac{\pi}{2}
\end{align*}
Bemerkung:\\
Die Funktion $ f(x)=\sqrt{1-x^2} $ entspricht einen Halbkreis mit Radius 1 und mit dem Flächeninhalt von $ \frac{\pi}{2} $.\\
\begin{center}
	
	
	\begin{tikzpicture}[xscale=2, yscale=2]
	
	% Gitternetzlinien
	\draw[xstep=1cm, ystep=1cm, very thin, color=lightgray] (-2,2) grid (2,-1);
	% x-Achse
	\draw[->] (-2,0) -- (2,0) node[below,xshift=-1em] {$x$}; 
	%Zahlen auf x-Achse
	\foreach \x in {-1,...,-1}
	\draw[shift={(\x,0)},color=black] (0pt,2pt) -- (0pt,-2pt) node[below] {\footnotesize $\x$};
	\foreach \x in {1,...,1}
	\draw[shift={(\x,0)},color=black] (0pt,2pt) -- (0pt,-2pt) node[below] {\footnotesize $\x$};
	% y-Achse 
	\draw[->] (0,-1) -- (0,2) node[left,yshift=-1em] {$f(x)$};%node[above left]
	%Zahlen auf y-Achse
	%\foreach \y in {-1,...,-1}
	%\draw[shift={(0,\y)},color=black] (2pt,0pt) -- (-2pt,0pt) node[left] {\footnotesize $\y$};
	\foreach \y in {1,...,1}
	\draw[shift={(0,\y)},color=black] (2pt,0pt) -- (-2pt,0pt) node[left] {\footnotesize $\y$};
	%Ursprung
	\draw[color=black] (0pt,-10pt) node[right] {\footnotesize $0$};
	
	
	
	\draw[thick,color=blue] plot[samples=200, domain=-1:1] 
	(\x,{sqrt(1-\x*\x)}) node[below,yshift=3em,xshift=1em]{$f(x)$};
	
	
	
	\end{tikzpicture}\\
\end{center}
\subsubsection{Trick 6 zur Substitutionsregel}
\href{https://youtu.be/b8ivYOZW2XE}{Video 26}\\
Berechne:\\
\begin{align*}
\int\limits_{}^{} \frac{1}{(x^2+1)^2}\cdot \mathrm{d}x
\end{align*}
Substitution:\\
\begin{align*}
x&=\tan(t)\\
\frac{\mathrm{d}x}{\mathrm{d}t}&=\frac{\mathrm{d}(\tan(t))}{\mathrm{d}t}=1+\tan^2(t)\\
\mathrm{d}x&=(1+\tan^2(t))\mathrm{d}t
\end{align*}
Dann erhält man:\\
\begin{align*}
&=\int\limits_{}^{} \frac{1}{(1+\tan^2(t)+1)^2}\cdot (1+\tan^2(t))\mathrm{d}t\\
&=\int\limits_{}^{} \frac{1}{\tan^2(t)+1}\mathrm{d}t\\
&=\int\limits_{1}^{}\frac{1}{\frac{\sin^2(t)}{\cos^2(t)}+1}\mathrm{d}t \\
&=\int\limits_{}^{} \frac{1}{\frac{\sin^2(t)+\cos^2(t)}{\cos^2(t)}}\mathrm{d}t\\
&=\int\limits_{}^{} \frac{\cos^2(t)}{\sin^2(t)+\cos^2(t)}\mathrm{d}t\\
&=\int\limits_{}^{} \cos^2(t)\mathrm{d}t
\end{align*}
Es gilt mit dem Additionstheorem:\\
\begin{align*}
\cos (2t)&=\cos^2(t)-\sin^2(t)\;\;\;\;\;\;\;\;\;\;\lvert \text{ \tiny +1}\\
\cos (2t)+1&=\cos^2(t)-\sin^2(t)+1\;\;\;\;\lvert \text{\tiny$ \sin^2(t)+\cos^2 (t)=1 $ }\\
\cos (2t)+1&=\cos^2(t)-\sin^2(t)+\sin^2(t)+\cos^2 (t)\\
\cos (2t)+1&=2\cos^2(t)\\
\frac{\cos(2t)+1}{2}&=\cos^2(t)
\end{align*}
Dann ergibt sich:\\
\begin{align*}
&=\int\limits_{}^{} \frac{\cos(2t)+1}{2}\cdot \mathrm{d}t\\
&=\frac{1}{2}\bigg(\frac{1}{2}\sin(2t)+t\bigg)+\alpha,\;\; \alpha\in\mathbb{R}\\
&=\frac{\sin(2t)}{4}+\frac{t}{2}+\alpha,\;\; \alpha\in\mathbb{R}\\
\end{align*}
Rücksubstitution:\\
Es gilt:\\
\begin{align*}
x&=\tan(t)\\
\Leftrightarrow&t=\arctan(x)
\end{align*}
Eingesetzt ergibt sich dann:\\
\begin{align*}
=\frac{\sin(2\arctan(x))}{4}+\frac{\arctan(x)}{2}+\alpha,\;\;\;\;\alpha\in\mathbb{R}
\end{align*}
\subsection{Integration komplexer Funktionen}
\subsubsection{Integration komplexer Funktionen}
\href{https://youtu.be/ms9HZdijOKI}{Video 27}\\
Sei $ f:\mathbb{R}\supset[a,b ] \longmapsto \mathbb{C}$ eine Funktion mit reellem Argument.\\
Dann gilt:\\
\begin{align*}
f(x)=u(x)+\mathrm{i}v(x)\;\;\text{mit }\;\;u,v:[a,b]\longmapsto \mathbb{R}
\end{align*}
Damit definieren wir:\\
\begin{align*}
\int\limits_{a}^{b} f(x)\mathrm{d}x=\int\limits_{a}^{b}u(x) \mathrm{d}x+\mathrm{i}\int\limits_{a}^{b} v(x)\mathrm{d}x
\end{align*}
Somit bleiben alle Integrationsregeln gültig.\\
Beispiel:\\
\begin{align*}
f(x)=x^2+\mathrm{i}x+2=x^2+2+\mathrm{i}x\\
\end{align*}
Das Integral ist dann:\\
\begin{align*}
\int\limits_{}^{} f(x)\mathrm{d}x=\int\limits_{}^{} (x^2+2)\mathrm{d}x+\mathrm{i}\int\limits_{}^{} x\mathrm{d}x=\frac{x^3}{3}+2x+\mathrm{i}\frac{x^2}{2}+\alpha,\;\;\alpha\in \mathbb{C}
\end{align*}
\subsubsection{Aufgabe zur Integration komplexer Funktionen}
\href{https://youtu.be/jG6L4TwLSK0}{Video 28}\\
Berechne:\\
(i)\;\;\;$ \int\limits_{}^{} \frac{1}{(x-\mathrm{i})^2}\cdot \mathrm{d}x $\\
(ii)\;\;$ \int\limits_{}^{} \frac{1}{x-\mathrm{i}}\cdot \mathrm{d}x $\\
Zu (i):\\
\begin{align*}
&\int\limits_{}^{} \frac{1}{(x-\mathrm{i})^2}\cdot \mathrm{d}x=\int\limits_{}^{} (x-1)^{-2}\cdot\mathrm{d}x\\
&\text{Wird wie im reellen Bereich integriert. }
\end{align*}
Zu (ii):\\
\begin{align*}
\int\limits_{}^{} \frac{1}{x-\mathrm{i}}\cdot \mathrm{d}x&=\int\limits_{}^{} \bigg(\frac{x+\mathrm{i}}{x+1}\cdot \frac{1}{x-1}\bigg)\mathrm{d}x\\
&=\int\limits_{}^{}\bigg( \frac{x+\mathrm{i}}{x^2+1}\bigg)\mathrm{d}x\\
&=\int\limits_{}^{} \frac{x}{x^2+1}\cdot \mathrm{d}x+\mathrm{i}\int\limits_{}^{} \frac{1}{x^2+1}\cdot \mathrm{d}x
\end{align*}
Es gilt:\\
\begin{align*}
(\arctan(x))'=\frac{1}{x^2+1}
\end{align*}
Die Ableitung von :\\
\begin{align*}
\big(\ln(x^2+1)\big)'= \frac{2x}{x^2+1}
\end{align*}
Dann ist die Ableitung von:\\
\begin{align*}
\bigg(\frac{1}{2}\ln(x^2+1)\bigg)'=\frac{x}{x^2+1}
\end{align*}
So ergibt sich die Stammfunktion:\\
\begin{align*}
\int\limits_{}^{} \frac{1}{x-\mathrm{i}}\cdot \mathrm{d}x=\frac{1}{2}\ln(x^2+1)+\arctan(x)+\alpha,\;\;\alpha\in \mathbb{C}
\end{align*}
Bemerkung:\\
Die Stammfunktion von $ \int\limits_{}^{} \frac{1}{x-1} \cdot \mathrm{d}x=\ln(x-1)+\alpha$ ist nicht identisch mit der Stammfunktion von $ \int\limits_{}^{} \frac{1}{x-\mathrm{i}} $.\\
\subsection{Integration rationaler Funktionen}
\subsubsection{ Integration rationaler Funktionen/Integration mittels Partialbruchzerlegung}
\href{https://youtu.be/QJQBrXdNNSA}{Video 29}\\
Integration mittels Partialbruchzerlegung (PBZ).
Betrachte:\\
\begin{align*}
\int\limits_{}^{} \frac{x}{x^2-1}\cdot \mathrm{d}x
\end{align*}
\underline{PBZ:}\\
\begin{align*}
(x^2-1)=(x-1)(x+1)
\end{align*}
Es gilt:\\
\begin{align*}
\frac{x}{(x-1)(x+1)}=\frac{\mathrm{A}}{x+1}+\frac{\mathrm{B}}{x-1}
\end{align*}
Bestimmung von A und B mittels Zuhaltemethode:\\
Bemerkung zur Zuhaltemethode: (siehe Skript 7.4.5)\\
Zur Bestimmung von A die Gleichung mit $ (x+1) $ multiplizieren.\\
\begin{align*}
\frac{x(x+1)}{(x-1)(x+1)}=\frac{\mathrm{A(x+1)}}{x+1}+\frac{\mathrm{B(x+1)}}{x-1}\\
\end{align*}
Für $ x=-1 $ gilt dann:\\
\begin{align*}
\frac{x}{(x-1)}=\frac{-1}{(-1-1)}=A=\frac{1}{2}\\
\end{align*}
Für $ x=1 $ gilt dann:\\
\begin{align*}
\frac{x}{x+1}=\frac{1}{(1+1)}=B=\frac{1}{2}
\end{align*}
Dies ergibt dann:\\
\begin{align*}
\frac{x}{x^2-1}=\frac{1}{2(x+1)}+\frac{1}{2(x-1)}
\end{align*}
Die Stammfunktion kann nun leicht berechnet werden.\\
\begin{align*}
\int\limits_{}^{} \frac{x}{x^2-1}\cdot \mathrm{d}x&=\frac{1}{2}\int\limits_{}^{} \frac{1}{x+1}\cdot \mathrm{d}x+\frac{1}{2}\int\limits_{}^{} \frac{1}{x-1}\cdot \mathrm{d}x\\
&=\frac{1}{2}\ln(x+1)+\frac{1}{2}\ln(x-1)+\alpha,\;\;\alpha\in \mathbb{R}
\end{align*}
Mittels der Substitution kann man auch die Stammfunktion: $ \int\limits_{}^{} \frac{x}{x^2-1}\cdot \mathrm{d}x $ bestimmen:\\
Substitution:\\
\begin{align*}
t&=x^2+1\\
\frac{\mathrm{d}t}{\mathrm{d}x}&=\frac{d(x^2-1)}{\mathrm{d}x}=2x\\
\Leftrightarrow \mathrm{d}x&=\frac{\mathrm{\mathrm{d}}t}{2x}
\end{align*}
Dann erhält man:\\
\begin{align*}
&=\int\limits_{}^{} \frac{x}{t}\cdot \frac{\mathrm{d}t}{2x}\\
&=\frac{1}{2}\int\limits_{}^{} \frac{1}{t}\cdot \mathrm{d}t=\frac{1}{2}\ln(t)+\alpha,\;\;\alpha\in\mathbb{R}
\end{align*}
Rücksubstitution mit $ t=x^2-1 $:\\
\begin{align*}
\int\limits_{}^{} \frac{x}{x^2-1}\cdot \mathrm{d}x&=\frac{1}{2}\ln(x^2-1)+\alpha,\;\;\alpha\in\mathbb{R}\\
&=\frac{1}{2}\ln((x+1)(x-1))+\alpha,\;\;\alpha\in\mathbb{R}\\
&=\frac{1}{2}\ln(x+1)+\frac{1}{2}\ln(x-1)+\alpha,\;\;\alpha\in\mathbb{R}\\
\end{align*}
\subsubsection{Aufgabe Integration rationaler Funktionen/Integration mittels Partialbruchzerlegung }
\href{https://youtu.be/CS1aHixYoYU}{Video 30}\\
Berechne:\
\begin{align*}
\int\limits_{2}^{4} \frac{2}{x(x-1)^2}\cdot \mathrm{d}x
\end{align*}
Konstante ausklammern:\\
\begin{align*}
2\int\limits_{}^{}  \frac{1}{x(x-1)^2}\cdot \mathrm{d}x
\end{align*}
PBZ:\\
\begin{align*}
\frac{1}{x(x-1)^2}=\frac{\mathrm{A}}{x}+\frac{\mathrm{B}}{(x-1)}+\frac{C}{(x-1)^2}
\end{align*}
Bemerkung: B und C wegen doppelter Nullstelle.\\
\\
Bestimmung von A,B und C:\\
\begin{quote}
	A und C mit Zuhaltemethode bestimmen:\\
	\begin{align*}
	&\mathrm{A}=\frac{1}{(0-1)^2}=1\\
	&\mathrm{C}=\frac{1}{1}=1
	\end{align*}
	Bestimmung von B:\\
	Es gilt nun:\\
	\begin{align*}
	\frac{1}{x(x-1)^2}=\frac{\mathrm{1}}{x}+\frac{\mathrm{B}}{(x-1)}+\frac{1}{(x-1)^2}
	\end{align*}
	Gleichung mit $ x(x-1)^2 $ multiplizieren ergibt:.\\
	\begin{align*}
	1&=(x-1)^2+\mathrm{B}x(x-1)+x\\
	&=x^2-2x+1+\mathrm{B}x^2-\mathrm{B}x+x\\
	&=(1+\mathrm{B})x^2-x(2+\mathrm{B}-1)+1\\
	0x^2+0x^2+1&=(1+\mathrm{B})x^2-x(2+\mathrm{B}-1)+1\\
	\end{align*}
	Mit Koeffizientenvergleich ergibt sich:\\
	\begin{align*}
	0=1+B\Leftrightarrow B=-1\\
	\end{align*}
	Es gilt nun:\\
	\begin{align*}
	\frac{1}{x(x-1)^2}=\frac{\mathrm{1}}{x}+\frac{\mathrm{-1}}{(x-1)}+\frac{1}{(x-1)^2}
	\end{align*}
\end{quote}
So ergibt sich ein lösbares Integral:\\
\begin{align*}
2\int\limits_{2}^{4} \frac{1}{x(x-1)^2}\cdot \mathrm{d}x&=2\bigg(\int\limits_{2}^{4}\frac{1}{x}\cdot \mathrm{d}x+\int\limits_{2}^{4} \frac{-1}{x-1}\cdot \mathrm{d}x+\int\limits_{0}^{2}\frac{1}{(x-1)^2}\cdot \mathrm{d}x  \bigg)\\
&=2\bigg(\bigg[\ln(x)\bigg]^4_2-\bigg[\ln(x-1)\bigg]^4_2+\bigg[\frac{-1}{x-1}\bigg]^4_2\bigg)\\
&=\text{....weiter mit dem Taschenrechner }
\end{align*}
\subsection{Uneigentliche Integrale }
\subsubsection{Uneigentliche Integrale Teil 1}
\href{https://youtu.be/RO8cL5MA2Bo}{Video 31}\\
Sei $ f:[a,\infty[ \longmapsto \mathbb{R}$ eine Funktion und für jedes Intervall $ [a,b] $ integrierbar.\\
Dann gilt:\\
\begin{align*}
\int\limits_{a}^{\infty} f(x) \mathrm{d}x:=\lim\limits_ {b\to\infty}\int\limits_{a}^{b} f(x)\mathrm{d}x
\end{align*}
$\bullet\;\; \int\limits_{0}^{\infty}=1\mathrm{d}x=\infty $\\
\begin{center}
	
	
	\begin{tikzpicture}[xscale=1, yscale=1]
	
	% Gitternetzlinien
	\draw[xstep=1cm, ystep=1cm, very thin, color=lightgray] (-4,4) grid (4,-2);
	% x-Achse
	\draw[->] (-4,0) -- (4,0) node[below,xshift=-1em] {$x$}; 
	%Zahlen auf x-Achse
	\foreach \x in {-3,...,-1}
	\draw[shift={(\x,0)},color=black] (0pt,2pt) -- (0pt,-2pt) node[below] {\footnotesize $\x$};
	\foreach \x in {1,...,3}
	\draw[shift={(\x,0)},color=black] (0pt,2pt) -- (0pt,-2pt) node[below] {\footnotesize $\x$};
	% y-Achse 
	\draw[->] (0,-2) -- (0,4) node[left,yshift=-1em] {$f(x)$};%node[above left]
	%Zahlen auf y-Achse
	\foreach \y in {-1,...,-1}
	\draw[shift={(0,\y)},color=black] (2pt,0pt) -- (-2pt,0pt) node[left] {\footnotesize $\y$};
	\foreach \y in {1,...,3}
	\draw[shift={(0,\y)},color=black] (2pt,0pt) -- (-2pt,0pt) node[left] {\footnotesize $\y$};
	%Ursprung
	\draw[color=black] (0pt,-10pt) node[right] {\footnotesize $0$};
	
	\draw[thick,color=black] plot[samples=200, domain=0:4] 
	(\x,{1}) node[below, xshift=-5em,yshift=18] {$f(x)$};
	\draw[color=black] (1.5,0.5) node[right] {\footnotesize $\mathrm{A}=\infty$};
	\fill[color =red,fill opacity=0.3] plot[domain=0:4] (\x,{1 })% 
	-- plot[domain=4:0] (\x,{0});% 
	
	
	
	\end{tikzpicture}\\
\end{center}
$ \bullet \;\;\int\limits_{a}^{\infty}\mathrm{e}^x\cdot \mathrm{d}x  $\\

\begin{align*}
\int\limits_{0}^{\infty} \mathrm{e}^{-x}&=\lim\limits_ {b\to\infty}\int\limits_{0}^{b} \mathrm{e}^{-x}\\
&=\lim\limits_ {b\to\infty}\bigg[-\mathrm{e}^{-x}\bigg]^b_0\\
&=\lim\limits_ {b\to\infty}\bigg(-\mathrm{e}^{-b}-(-\mathrm{e}^{-0})\bigg)\\
&=\lim\limits_ {b\to\infty}-\mathrm{e}^{-b}+1\\
&=\lim\limits_ {b\to\infty}\frac{-1}{\mathrm{e}^b}+1=0+1=1
\end{align*}
\begin{center}
	
	
	\begin{tikzpicture}[xscale=2, yscale=2]
	
	% Gitternetzlinien
	\draw[xstep=1cm, ystep=1cm, very thin, color=lightgray] (0,2) grid (4,0);
	% x-Achse
	\draw[->] (0,0) -- (4,0) node[below,xshift=-1em] {$x$}; 
	%Zahlen auf x-Achse
	%\foreach \x in {-1,...,-1}
	%\draw[shift={(\x,0)},color=black] (0pt,2pt) -- (0pt,-2pt) node[below] {\footnotesize $\x$};
	\foreach \x in {1,...,3}
	\draw[shift={(\x,0)},color=black] (0pt,2pt) -- (0pt,-2pt) node[below] {\footnotesize $\x$};
	% y-Achse 
	\draw[->] (0,0) -- (0,2) node[left,yshift=-1em] {$f(x)$};%node[above left]
	%Zahlen auf y-Achse
	%\foreach \y in {0,...,-1}
	%\draw[shift={(0,\y)},color=black] (2pt,0pt) -- (-2pt,0pt) node[left] {\footnotesize $\y$};
	\foreach \y in {1,...,1}
	\draw[shift={(0,\y)},color=black] (2pt,0pt) -- (-2pt,0pt) node[left] {\footnotesize $\y$};
	%Ursprung
	%\draw[color=black] (0pt,-10pt) node[right] {\footnotesize $0$};
	
	\draw[thick,color=black] plot[samples=200, domain=0:4] 
	(\x,{e^(-\x)}) node[below, xshift=-8em,yshift=30] {$f(x)=\mathrm{e}^{-x}$};
	\draw[color=black] (.5,0.2) node[right] {\footnotesize $\mathrm{A}=1$};
	\fill[color =red,fill opacity=0.3] plot[domain=0:4] (\x,{e^(-\x) })% 
	-- plot[domain=4:0] (\x,{0});% 
	
	
	
	\end{tikzpicture}\\
\end{center}
Ebenso gilt auch analog für $ f:\;\;]\infty,a]\longmapsto \mathbb{R} $:\\
\begin{align*}
\int\limits_{-\infty}^{a}f(x)\mathrm{d}x =\lim\limits_ {b\to\infty}\int\limits_{b}^{a} f(x)\mathrm{d}x
\end{align*}
\subsubsection{Beispiel zu uneigentlichen Integralen Teil 1}
\href{https://youtu.be/GQxkZxjBx88}{Video 32}\\
Bestimme die Integrale:\\
$ \bullet\;\; \int\limits_{1}^{\infty} \frac{1}{\sqrt{x}}\cdot \mathrm{d}x $\\
\begin{align*}
\int\limits_{1}^{\infty} \frac{1}{\sqrt{x}}\cdot \mathrm{d}x&=\int\limits_{0}^{\infty}x^{-\frac{1}{2}} \cdot \mathrm{d}x\\
&=\lim\limits_ {b\to \infty}\int\limits_{1}^{b} x^{-\frac{1}{2}}\cdot \mathrm{d}x\\
&=\lim\limits_ {b\to\infty}\bigg[2\cdot x^{\frac{1}{2}}\bigg]^b_1\\
&=\lim\limits_ {b\to\infty}(2\sqrt{b}-2)=\infty
\end{align*}
$ \bullet\;\; \int\limits_{1}^{\infty} \frac{1}{x}\cdot \mathrm{d}x $\\
\begin{align*}
\int\limits_{1}^{\infty} \frac{1}{x}\cdot \mathrm{d}x&=\lim\limits_ {b\to\infty}\int\limits_{1}^{b} \frac{1}{x}\cdot \mathrm{d}x\\
&=\lim\limits_ {b\to\infty}\bigg[\ln(x)\bigg]^b_1\\
&=\lim\limits_ {b\to\infty}\big(\ln(b)-\ln(1)\big)=\lim\limits_ {b\to\infty}\ln(b)=\infty
\end{align*}
\begin{center}
	
	
	\begin{tikzpicture}[xscale=2, yscale=2]
	
	% Gitternetzlinien
	\draw[xstep=1cm, ystep=1cm, very thin, color=lightgray] (0,2) grid (4,0);
	% x-Achse
	\draw[->] (0,0) -- (4,0) node[below,xshift=-1em] {$x$}; 
	%Zahlen auf x-Achse
	%\foreach \x in {-1,...,-1}
	%\draw[shift={(\x,0)},color=black] (0pt,2pt) -- (0pt,-2pt) node[below] {\footnotesize $\x$};
	\foreach \x in {1,...,3}
	\draw[shift={(\x,0)},color=black] (0pt,2pt) -- (0pt,-2pt) node[below] {\footnotesize $\x$};
	% y-Achse 
	\draw[->] (0,0) -- (0,2) node[left,yshift=-1em] {$f(x)$};%node[above left]
	%Zahlen auf y-Achse
	%\foreach \y in {0,...,-1}
	%\draw[shift={(0,\y)},color=black] (2pt,0pt) -- (-2pt,0pt) node[left] {\footnotesize $\y$};
	\foreach \y in {1,...,1}
	\draw[shift={(0,\y)},color=black] (2pt,0pt) -- (-2pt,0pt) node[left] {\footnotesize $\y$};
	%Ursprung
	%\draw[color=black] (0pt,-10pt) node[right] {\footnotesize $0$};
	
	\draw[thick,color=black] plot[samples=200, domain=1:4] 
	(\x,{1/sqrt(\x)}) node[below, xshift=-8em,yshift=35] {$f(x)=\frac{1}{\sqrt{x}}$};
	\draw[color=black] (1.5,0.2) node[right] {\footnotesize $\mathrm{A}=\infty$};
	\fill[color =red,fill opacity=0.3] plot[domain=1:4] (\x,{1/sqrt(\x) })% 
	-- plot[domain=4:1] (\x,{0});% 
	
	\draw[thick,color=blue] plot[samples=200, domain=1:4] 
	(\x,{1/\x}) node[below, xshift=-2em,yshift=22] {$f(x)=\frac{1}{x}$};
	%\draw[color=blue] (.5,0.2) node[right] {\footnotesize $\mathrm{A}=1$};
	\fill[color =green,fill opacity=0.3] plot[domain=1:4] (\x,{1/\x })% 
	-- plot[domain=4:1] (\x,{0});% 
	
	\end{tikzpicture}\\
\end{center}
$ \bullet\;\; \int\limits_{1}^{\infty} \frac{1}{\sqrt{x^3}}\cdot \mathrm{d}x $\\
\begin{align*}
\int\limits_{1}^{\infty} \frac{1}{\sqrt{x^3}}\cdot \mathrm{d}x&=\lim\limits_ {b\to\infty}x^{-\frac{3}{2}}\cdot \mathrm{d}x\\
&=\lim\limits_ {b\to\infty}\bigg[-2x^{-\frac{1}{2}}\bigg]^b_1\\
&=\lim\limits_ {b\to\infty}\bigg(-\frac{2}{\sqrt{b}}-(-2)\bigg)\\
&=\lim\limits_ {b\to\infty}\bigg(\underbrace {\frac{-2}{\sqrt{b}}}_{\mapsto 0}+2\bigg)=2
\end{align*}
Wenn die Kurve noch steiler wird, wird der Flächeninhalt=2.\\
\begin{center}
	
	
	\begin{tikzpicture}[xscale=2, yscale=2]
	
	% Gitternetzlinien
	\draw[xstep=1cm, ystep=1cm, very thin, color=lightgray] (0,2) grid (4,0);
	% x-Achse
	\draw[->] (0,0) -- (4,0) node[below,xshift=-1em] {$x$}; 
	%Zahlen auf x-Achse
	%\foreach \x in {-1,...,-1}
	%\draw[shift={(\x,0)},color=black] (0pt,2pt) -- (0pt,-2pt) node[below] {\footnotesize $\x$};
	\foreach \x in {1,...,3}
	\draw[shift={(\x,0)},color=black] (0pt,2pt) -- (0pt,-2pt) node[below] {\footnotesize $\x$};
	% y-Achse 
	\draw[->] (0,0) -- (0,2) node[left,yshift=-1em] {$f(x)$};%node[above left]
	%Zahlen auf y-Achse
	%\foreach \y in {0,...,-1}
	%\draw[shift={(0,\y)},color=black] (2pt,0pt) -- (-2pt,0pt) node[left] {\footnotesize $\y$};
	\foreach \y in {1,...,1}
	\draw[shift={(0,\y)},color=black] (2pt,0pt) -- (-2pt,0pt) node[left] {\footnotesize $\y$};
	%Ursprung
	%\draw[color=black] (0pt,-10pt) node[right] {\footnotesize $0$};
	
	\draw[thick,color=black] plot[samples=200, domain=1:4] 
	(\x,{1/sqrt(\x)});
	\draw[color=black] (1.5,0.2) node[right] {\footnotesize $\mathrm{A}=2$};
	\fill[color =red,fill opacity=0.3] plot[domain=1:4] (\x,{1/sqrt(\x) })% 
	-- plot[domain=4:1] (\x,{0});% 
	
	\draw[thick,color=blue] plot[samples=200, domain=1:4] 
	(\x,{1/\x});
	%\draw[color=blue] (.5,0.2) node[right] {\footnotesize $\mathrm{A}=1$};
	\fill[color =green,fill opacity=0.3] plot[domain=1:4] (\x,{1/\x })% 
	-- plot[domain=4:1] (\x,{0});% 
	
	\draw[thick,color=red] plot[samples=200, domain=1:4] 
	(\x,{1/sqrt(\x^3)}) node[below, xshift=-8em,yshift=72] {$f(x)=\frac{1}{\sqrt{x^3}}$};
	%\draw[color=blue] (.5,0.2) node[right] {\footnotesize $\mathrm{A}=1$};
	\fill[color =red,fill opacity=0.3] plot[domain=1:4] (\x,{1/sqrt(\x^3) })% 
	-- plot[domain=4:1] (\x,{0});% 
	\end{tikzpicture}\\
\end{center}
\subsubsection{Uneigentliche Integrale Teil 2}
\href{https://youtu.be/ObwW-BSHLz4}{Video 33}\\
Sei $ f:\;]a,b ]$ eine Funktion und existieren für alle $ c\in ]a,b ] $ das $ \int\limits_{c}^{b}f(x) \mathrm{d}x $. Wobei bei $ a $ eine kritische Stelle, wie z.B eine Polstelle ist. \\
Dann definieren wir:\\
\begin{align*}
\int\limits_{a}^{b} f(x)\mathrm{d}x:=\lim\limits_ {\substack{t\mapsto a\\t>a}}\int\limits_{t}^{b} f(x)\mathrm{d}x
\end{align*}
\begin{center}
	\begin{tikzpicture}[xscale=1, yscale=1]
	
	% Gitternetzlinien
	\draw[xstep=1cm, ystep=1cm, very thin, color=lightgray] (-4,4) grid (4,-2);
	% x-Achse
	\draw[->] (-4,0) -- (4,0) node[below,xshift=-1em] {$x$}; 
	%Zahlen auf x-Achse
	\foreach \x in {-3,...,-1}
	\draw[shift={(\x,0)},color=black] (0pt,2pt) -- (0pt,-2pt) node[below] {\footnotesize $\x$};
	\foreach \x in {1,...,3}
	\draw[shift={(\x,0)},color=black] (0pt,2pt) -- (0pt,-2pt) node[below] {\footnotesize $\x$};
	% y-Achse 
	\draw[->] (0,-2) -- (0,4) node[left,yshift=-1em] {$f(x)$};%node[above left]
	%Zahlen auf y-Achse
	\foreach \y in {-1,...,-1}
	\draw[shift={(0,\y)},color=black] (2pt,0pt) -- (-2pt,0pt) node[left] {\footnotesize $\y$};
	\foreach \y in {1,...,3}
	\draw[shift={(0,\y)},color=black] (2pt,0pt) -- (-2pt,0pt) node[left] {\footnotesize $\y$};
	%Ursprung
	\draw[color=black] (0pt,-10pt) node[right] {\footnotesize $ t$};
	
	\draw[thick,color=blue] plot[samples=200, domain=0.25:4] 
	(\x,{1/\x}) node[below, xshift=-2em,yshift=2em]{$f(x)$};
	
	\fill[color =red,fill opacity=0.3] plot[domain=0.25:4] (\x,{1/\x) })% 
	-- plot[domain=4:0.25] (\x,{0});% 
	\fill[color =red,fill opacity=0.3] plot[domain=0:0.25] (\x,{4 })% 
	-- plot[domain=0.25:0] (\x,{0});% 
	
	
	\end{tikzpicture}\\
\end{center}
Analog definieren wie für $ f:[a, b[ $ das Folgende:\\
\begin{align*}
\int\limits_{a}^{b} f(x)\mathrm{d}x:=\lim\limits_ {\substack{t\mapsto b\\t<b}}\int\limits_{a}^{t} f(x)\mathrm{d}x
\end{align*}
\begin{center}
	\begin{tikzpicture}[xscale=1, yscale=1]
	
	% Gitternetzlinien
	\draw[xstep=1cm, ystep=1cm, very thin, color=lightgray] (-4,4) grid (4,-2);
	% x-Achse
	\draw[->] (-4,0) -- (4,0) node[below,xshift=-1em] {$x$}; 
	%Zahlen auf x-Achse
	\foreach \x in {-3,...,-1}
	\draw[shift={(\x,0)},color=black] (0pt,2pt) -- (0pt,-2pt) node[below] {\footnotesize $\x$};
	\foreach \x in {1,...,3}
	\draw[shift={(\x,0)},color=black] (0pt,2pt) -- (0pt,-2pt) node[below] {\footnotesize $\x$};
	% y-Achse 
	\draw[->] (0,-2) -- (0,4) node[left,yshift=-1em] {$f(x)$};%node[above left]
	%Zahlen auf y-Achse
	\foreach \y in {-1,...,-1}
	\draw[shift={(0,\y)},color=black] (2pt,0pt) -- (-2pt,0pt) node[left] {\footnotesize $\y$};
	\foreach \y in {1,...,3}
	\draw[shift={(0,\y)},color=black] (2pt,0pt) -- (-2pt,0pt) node[left] {\footnotesize $\y$};
	%Ursprung
	\draw[color=black] (-0.25,-10pt) node[right] {\footnotesize $ t$};
	
	\draw[thick,color=blue] plot[samples=200, domain=-4:-0.25] 
	(\x,{1/-\x}) node[below, xshift=-2em,yshift=-4em]{$f(x)$};
	
	\fill[color =red,fill opacity=0.3] plot[domain=-4:-0.25] (\x,{1/-\x) })% 
	-- plot[domain=-0.25:-4] (\x,{0});% 
	\fill[color =red,fill opacity=0.3] plot[domain=-0.25:0] (\x,{4 })% 
	-- plot[domain=0:-0.25] (\x,{0});% 
	\end{tikzpicture}\\
\end{center}
Beispiel:\\
\begin{align*}
\int\limits_{0}^{2} \frac{1}{x^2}\mathrm{d}x&=\lim\limits_ {\substack{t\mapsto 0\\t>0}}\int\limits_{t}^{2} \frac{1}{x^2}\mathrm{d}x\\
&=\lim\limits_ {\substack{t\mapsto 0\\t>0}}\bigg[-\frac{1}{x}\bigg]^2_t\\
&=\lim\limits_ {\substack{t\mapsto 0\\t>0}}\bigg(-\frac{1}{2}+\underbrace {\frac{1}{t}}_{+\infty}\bigg)=\infty
\end{align*}
\subsubsection{Aufgabe zu uneigentlichen Integralen Teil 2}
\href{https://youtu.be/mUNTuWgMLKA}{Video 34}\\
Für welches $ \alpha \in \;]0,\infty]$ existiert das Integral: $ \int\limits_{0}^{1}\frac{1}{x^\alpha}\mathrm{d}x  $ ?\\
\\
Wie steil muß die Kurve an der Polstelle bei $ x=0 $ werden, dass das Integral existiert und nicht in das Unendliche geht?\\

\begin{align*}
\int\limits_{0}^{1}\frac{1}{x^\alpha} \mathrm{d} x=\lim\limits_{\substack{t\mapsto 0\\t>0}}\int\limits_{t}^{1}\frac{1}{x^\alpha}\mathrm{d}x=\lim\limits_ {\substack{t\mapsto 0\\t>0}} \bigg[\frac{1}{1-\alpha}x^{-\alpha+1}\bigg]^1_t=\lim\limits_ {\substack{t\mapsto 0\\t>0}}\bigg(\frac{1}{1-\alpha}-\frac{1}{1-\alpha}t^{1-\alpha}\bigg)
\end{align*}
Die Frage: Wann wird der $ \lim\limits_ {t\mapsto 0} t^{1-\alpha}$ unendlich oder Null? Für $ 1-\alpha<0 $ ist die Hochzahl negativ $ \Rightarrow \lim\limits_ {t\mapsto 0} \frac{1}{t^\beta}=\infty$ mit $ \beta>0 $ und für $1-\alpha>0  $ ist die Hochzahl positiv $ \Rightarrow \lim\limits_ {t\mapsto 0}t^\beta=0 $ mit $ \beta>0 $.\\
Fallunterscheidung:\\
1. Fall: $ 1-\alpha<0$\\
\begin{quote}
	
	
	\begin{align*}
	\alpha>1\Leftrightarrow 1-\alpha<0
	\end{align*}
	Dann gilt:\\
	\begin{align*}
	&t^{1-\alpha}=t^{-\beta}=\frac{1}{t^\beta} \;\;\text{mit }\beta>0\\
	&\lim\limits_ {t\to 0}\frac{1}{t^\beta}=\infty\\
	&\Rightarrow \int\limits_{0}^{1}\frac{1}{x^\alpha} \mathrm{d} x=\lim\limits_ {\substack{t\mapsto 0\\t>0}}\bigg(\frac{1}{1-\alpha}-\frac{1}{1-\alpha}t^{1-\alpha}\bigg)=\infty
	\end{align*}
	Das Integral geht gegen unendlich für $ \alpha>1 $. Es existiert kein Grenzwert.\\
\end{quote}
2. Fall: $ 1-\alpha>0$\\
\begin{quote}
	\begin{align*}
	\alpha<1\Leftrightarrow 1-\alpha>0
	\end{align*}
	Dann git:\\
	\begin{align*}
	&t^{1-\alpha}=t^\beta \;\;\text{mit }\beta>0\\
	&\lim\limits_ {t\to 0}t^\beta=0\\
	&\Rightarrow \int\limits_{0}^{1}\frac{1}{x^\alpha} \mathrm{d} x=\lim\limits_ {\substack{t\mapsto 0\\t>0}}\bigg(\frac{1}{1-\alpha}-\frac{1}{1-\alpha}t^{1-\alpha}\bigg)=\frac{1}{1-\alpha}
	\end{align*}
	Das Integral geht gegen einen Grenzwert für $ \alpha<1 $. Es existiert ein Grenzwert.\\
\end{quote}
Beispiel:\\
Ein Integral existiert für: $ \alpha=\frac{1}{2} $\\
\begin{align*}
\int\limits_{0}^{1} \frac{1}{\sqrt{x}}\mathrm{d}x
\end{align*}
Es existiert kein Integral für: $ \alpha=2 $\\
\begin{align*}
\int\limits_{0}^{1} \frac{1}{x^2}\mathrm{d}x
\end{align*}
\begin{center}
	
	
	\begin{tikzpicture}[xscale=4, yscale=1]
	
	% Gitternetzlinien
	\draw[xstep=1cm, ystep=1cm, very thin, color=lightgray] (0,5) grid (2,0);
	% x-Achse
	\draw[->] (0,0) -- (2,0) node[below,xshift=-1em] {$x$}; 
	%Zahlen auf x-Achse
	%\foreach \x in {-1,...,-1}
	%\draw[shift={(\x,0)},color=black] (0pt,2pt) -- (0pt,-2pt) node[below] {\footnotesize $\x$};
	\foreach \x in {0,...,1}
	\draw[shift={(\x,0)},color=black] (0pt,2pt) -- (0pt,-2pt) node[below] {\footnotesize $\x$};
	% y-Achse 
	\draw[->] (0,0) -- (0,5) node[left,yshift=-1em] {$f(x)$};%node[above left]
	%Zahlen auf y-Achse
	%\foreach \y in {0,...,-1}
	%\draw[shift={(0,\y)},color=black] (2pt,0pt) -- (-2pt,0pt) node[left] {\footnotesize $\y$};
	\foreach \y in {1,...,4}
	\draw[shift={(0,\y)},color=black] (0.5pt,0pt) -- (-0.5pt,0pt) node[left] {\footnotesize $\y$};
	%Ursprung
	%\draw[color=black] (-1pt,-10pt) node[right] {\footnotesize $0$};
	
	%\draw[thick,color=black] plot[samples=200, domain=0.20:1] 
	%(\x,{1/\x}) node[below, xshift=-8em,yshift=102] {$f(x)=\frac{1}{x }$};
	%\draw[color=black] (1.5,0.2) node[right] {\footnotesize $\mathrm{A}=2$};
	%\fill[color =red,fill opacity=0.3] plot[domain=0.28:4] (\x,{1/\x })% 
	%-- plot[domain=4:0.28] (\x,{0});% 
	\draw[thick,color=white] plot[samples=200, domain=0.:0.446] 
	(\x,{5});
	
	\fill[color =green,fill opacity=0.3] plot[domain=0:0.446] (\x,{5 })% 
	-- plot[domain=0.446:0] (\x,{0});% 
	
	
	\draw[thick,color=blue] plot[samples=200, domain=0.446:2] 
	(\x,{1/\x^2}) node[below, xshift=-10em,yshift=92] {$f(x)=\frac{1}{x^2 }$};
	\draw[color=blue] (.3,2.2) node[right] {\footnotesize $\mathrm{A}=\infty$};
	\fill[color =green,fill opacity=0.3] plot[domain=0.446:1] (\x,{1/\x^2 })% 
	-- plot[domain=1:0.446] (\x,{0});% 
	
	\draw[thick,color=red] plot[samples=200, domain=0.041:2] 
	(\x,{1/(sqrt(\x))}) node[below, xshift=-5em,yshift=32] {$f(x)=\frac{1}{\sqrt{x}}$};
	\draw[color=red] (.3,0.3) node[right] {\footnotesize $\mathrm{A}=\frac{1}{1-\alpha}=2$};
	\fill[color =red,fill opacity=0.3] plot[domain=0.041:1] (\x,{1/(sqrt(\x)) })% 
	-- plot[domain=1:0.041] (\x,{0});% 
	
	
	\fill[color =red,fill opacity=0.3] plot[domain=0:0.042] (\x,{5 })% 
	-- plot[domain=0.041:0] (\x,{0});% 
	\end{tikzpicture}\\
\end{center}
\subsubsection{Majorantenkriterium für uneigentliche }
\href{https://youtu.be/qe4opJMGitY}{Video 35}\\
Seien $ a\in\mathbb{R} $ und $ b\in\mathbb{R}\cup\{\infty\} $ und $ f,g:[a,b[\;\longmapsto \mathbb{R} $ Funktionen, die auf allen Intervallen $ [a,c] $ mit $ a<c<b $ integrierbar sind, dann gilt das Majorantenkriterium.\\
\\
\underline{Majorantenkriterium:}\\
Ist $ \lvert f(x)\lvert\leq g(x) $ für alle $ x\in[a,b[ $ und existiert das Integral: $ \int\limits_{a}^{b}g(x) \mathrm{d}x $, dann existiert auch das Integral: $ \int\limits_{a}^{b}f(x)\mathrm{d}x  $.\\
\begin{center}
	\begin{align*}
	\lvert f(x)\lvert\leq g(x)\;\;\text{für alle }\;\;x\in[-2,2]
	\end{align*}
	\begin{tikzpicture}[xscale=1, yscale=1]
	
	% Gitternetzlinien
	\draw[xstep=1cm, ystep=1cm, very thin, color=lightgray] (-4,4) grid (4,-2);
	% x-Achse
	\draw[->] (-4,0) -- (4,0) node[below,xshift=-1em] {$x$}; 
	%Zahlen auf x-Achse
	\foreach \x in {-3,...,-1}
	\draw[shift={(\x,0)},color=black] (0pt,2pt) -- (0pt,-2pt) node[below] {\footnotesize $\x$};
	\foreach \x in {1,...,3}
	\draw[shift={(\x,0)},color=black] (0pt,2pt) -- (0pt,-2pt) node[below] {\footnotesize $\x$};
	% y-Achse 
	\draw[->] (0,-2) -- (0,4) node[left,yshift=-1em] {$f(x)$};%node[above left]
	%Zahlen auf y-Achse
	\foreach \y in {-1,...,-1}
	\draw[shift={(0,\y)},color=black] (2pt,0pt) -- (-2pt,0pt) node[left] {\footnotesize $\y$};
	\foreach \y in {1,...,3}
	\draw[shift={(0,\y)},color=black] (2pt,0pt) -- (-2pt,0pt) node[left] {\footnotesize $\y$};
	%Ursprung
	\draw[color=black] (0pt,-10pt) node[right] {\footnotesize $0$};
	
	\draw[thick,color=blue] plot[samples=200, domain=-2:2] 
	(\x,{-\x*\x*0.2+1.8}) node[yshift=3em, xshift=1em] {$g(x)=0.2x^2+1.8$};
	
	\draw[thick,color=red] plot[samples=200, domain=0:2] 
	(\x,{0.5*\x}) node[yshift=-4em, xshift=1em]{$\lvert f(x)\lvert=0.5x $};
	\draw[thick,color=red] plot[samples=200, domain=-2:0] 
	(\x,{-0.5*\x});
	
	\fill[color =red,fill opacity=0.3] plot[domain=-2:0] (\x,{-0.5*\x })% 
	-- plot[domain=0:-2] (\x,{0});% 
	\fill[color =red,fill opacity=0.3] plot[domain=0:2] (\x,{0.5*\x })% 
	-- plot[domain=2:0] (\x,{0});% 
	\fill[color =red,fill opacity=0.3] plot[domain=-2:2] (\x,{-\x*\x*0.2+1.8 })% 
	-- plot[domain=2:-2] (\x,{0});% 
	\end{tikzpicture}\\	
\end{center}
Beispiel:\\
Existiert das Integral:\\
\begin{align*}
\int\limits_{1}^{\infty}\frac{\mathrm{e}^{-x}}{x} \mathrm{d}x
\end{align*}
Sei:\\
\begin{align*}
f(x)=\frac{\mathrm{e}^{-x}}{x}
\end{align*}
dann ist:\\
\begin{align*}
\lvert f(x)\lvert=\frac{\mathrm{e}^{-x}}{x}\leq \mathrm{e}^{-x}
\end{align*}
Betrachte:\\
\begin{align*}
\int\limits_{0}^{\infty} \mathrm{e}^{-x}\mathrm{d}x=\lim\limits_ {b\to\infty}\int\limits_{1}^{b} \mathrm{e}^{-x}\mathrm{d}x=\lim\limits_ {b\to\infty}\bigg[-\mathrm{e}^{-x}\bigg]^b_1&=\lim\limits_ {b\to\infty}\bigg(-\mathrm{e}^{-b}-(-e^{-1})\bigg)\\
&=\lim\limits_ {b\to\infty}\bigg(\underbrace {-\frac{1}{\mathrm{e}^b}}_{\mapsto 0}+\mathrm{e}^{-1}\bigg)=\mathrm{e}^{-1}=\frac{1}{\mathrm{e}}\\
\int\limits_{0}^{\infty} \mathrm{e}^{-x}\mathrm{d}x&=\frac{1}{\mathrm{e}}
\end{align*}
Nach dem Majorantenkriterium gilt:\\
\begin{align*}
\lvert f(x)\lvert\leq g(x) 
\end{align*}
Da die Funktion $ \int\limits_{1}^{\infty} g(x)\mathrm{d}x $ existiert folgt nach dem Majorantenkriterium, dass das Integral:\\
\begin{align*}
\int\limits_{1}^{\infty}\frac{\mathrm{e}^{-x}}{x} \mathrm{d}x
\end{align*}
auch existiert.\\
Bemerkung:\\
Existiert kein Grenzwert, so kann keine Aussage getroffen werden, ob das Integral existiert.\\
Das Integral $ \int\limits_{1}^{\infty}\frac{1}{\sqrt{x^3}} \mathrm{d}x=2 $ und das Integral $ \int\limits_{1}^{\infty} \frac{1}{x} \mathrm{dx}=\infty$.\\
Aber $ |f(x)|=\frac{1}{\sqrt{x^3}}<g(x)=\frac{1}{x} $ im Bereich von $ x\in   [1,\infty[ $ .
\subsubsection{Minorantenkriterium für uneigentliche Integrale}
\href{https://youtu.be/wBYcVF1449w}{Video 36}\\
Seien $ a\in\mathbb{R} $ und $ b\in\mathbb{R}\cup\{\infty\} $ und $ f,g:[a,b[\;\longmapsto \mathbb{R} $ Funktionen, die auf allen Intervallen $ [a,c] $ mit $ a<c<b $ integrierbar sind, dann gilt das Minorantenkriterium.\\
Gelte:\\
\begin{align*}
0\leq g(x) \leq f(x)\;\;\text{für alle } x\in[a,c[
\end{align*}
und existieren nicht:\\
\begin{align*}
\int\limits_{a}^{b} g(x)\mathrm{d}x=\lim\limits_ {\substack{t\mapsto b \\ t<b}}\int\limits_{a}^{b} g(x)\mathrm{d}x
\end{align*}
dann existiert auch nicht:
\begin{align*}
\int\limits_{a}^{b} f(x)\mathrm{d}x
\end{align*}
Mit dem Minorantenkriterium zeigt man, dass kein Grenzwert existiert. Denn wenn die kleinere Funktion $ g(x) $ divergiert, dann divergiert auch $ f(x) $.\\
Mit dem Majorantenkriterium zeigt man, dass eine Konvergenz existiert.\\
\\
Beispiel:\\
Existiert das Integral:\\
\begin{align*}
\int\limits_{1}^{\infty} \cosh(x)\mathrm{d}x
\end{align*}
Es gilt:\\
\begin{align*}
0\leq 1 \leq \cosh(x)
\end{align*}
Bemerkung:\\
\begin{align*}
&\cosh(x)=\frac{e^x+e^{-1}}{2}\\
&x=0\Rightarrow \cosh(x)=1
\end{align*}
Daraus folgt dann:\\
\begin{align*}
&g(x)=1\\
&\int\limits_{1}^{\infty} 1\mathrm{d}x=\lim\limits_ {t\to\infty}\int\limits_{1}^{t} 1\mathrm{d}x=\lim\limits_ {t\to\infty}\bigg[x\bigg]^t_1=\lim\limits_ {t\to\infty} t-1=\infty
\end{align*}
Mit dem Minorantenkriterium gilt dann:\\
\begin{align*}
\int\limits_{1}^{\infty} 1\mathrm{d}x=\int\limits_{1}^{\infty} =\infty\overset{Minorantenkr.}{\longmapsto}\int\limits_{1}^{\infty} \cosh(x)\mathrm{d}x\;\;\text{existiert nicht }
\end{align*}
\subsubsection{Aufgabe 1 zum Minoranten- und Majorantenkriterium für uneigentliche Integrale}
\href{https://youtu.be/xJGJ5_nopx0}{Video 37}\\
Existiert das Integral:\\
\begin{align*}
\int\limits_{0}^{1} \frac{1}{(2+x^3)^{\frac{3}{2}}}\mathrm{d}x
\end{align*}
Betrachte mit dem Majoratenkriterium:\\
\begin{align*}
&\lvert f(x)\lvert\leq g(x) \\
&f(x)=\bigg|\frac{1}{(2+x^3)^{\frac{3}{2}}}\bigg|\\
&\text{für }\\x>0\Rightarrow \text{Betrag fällt weg }
\end{align*}
Abschätzen von $ g(x) $:\\
\begin{align*}
f(x)=\frac{1}{(2+x^3)^{\frac{3}{2}}}\leq \frac{1}{(x^3)^{\frac{3}{2}}}=\frac{1}{x^{\frac{9}{2}}}=x^{-\frac{9}{2}}=g(x)\;\;\text{ für}\;\;x\in[1,\infty[
\end{align*}
Integral von $\int\limits_{1}^{\infty}  g(x)\mathrm{d}x $ berechnen:\\
\begin{align*}
\int\limits_{0}^{\infty} x^{-\frac{9}{2}}\mathrm{d}x=\lim\limits_ {b\to\infty}\int\limits_{1}^{b} x^{-\frac{9}{2}}\mathrm{d}x&=\lim\limits_ {b\to\infty}\bigg[-\frac{2}{7}x^{-\frac{7}{2}}\bigg]^b_1\\
&=\lim\limits_ {b\to\infty}\bigg(-\frac{2}{7}\frac{1}{b^{\frac{7}{2}}}-\big(-\frac{2}{7}\big)\bigg)=\frac{2}{7}\\
&\overset{Majork.}{\Rightarrow}\int\limits_{0}^{\infty} f(x)\mathrm{d}x=\int\limits_{1}^{\infty} \frac{1}{(2+x^3)^{\frac{3}{2}}}\mathrm{d}x\;\;\text{existiert. }
\end{align*}
\subsubsection{Aufgabe 2 zum Minoranten- und Majorantenkriterium für uneigentliche Integrale}
\href{https://youtu.be/XJG5oZmOEVE}{Video 38}\\
Existiert das Integral:\\
\begin{align*}
\int\limits_{2}^{\infty} \frac{2+\cos(x)}{\sqrt{x^2-1}}\mathrm{d}x
\end{align*}
Die Frage:\\
Welches Kriterium verwendet man?\\
Wenn das Interal nicht existiert, dann verwenden wir das Minorantekriterium, sonst das Majorantenkriterium.\\
Durch Abschätzen und Überlegung bekommt man für $ x\in[2,\infty[ $:\\
Der $ \cos(x) $ liegt zwischen -1 und 1. Der Zähler wird nicht größer als 3 und nicht kleiner als 1.\\
$ \sqrt{x^2-1} $ ist kleiner als $ \sqrt{x^2}=x $ daraus folgt dann:\\
\begin{align*}
\frac{2+\cos(x)}{\sqrt{x^2-1}}>\frac{1}{\sqrt{x^2-1}}>\frac{1}{x}
\end{align*}
Nun nehmen wir an, dass kein Integral existiert.\\
Es gilt mit dem Minorantenkriterium:\\
\begin{align*}
0\leq g(x)\leq f(x)
\end{align*}
Es gilt:\\
\begin{align*}
0\leq \frac{1}{x}\leq \frac{2+\cos(x)}{\sqrt{x^2-1}}
\end{align*}
Nun bestimmen wir das Integal: $ \int\limits_{2}^{\infty} g(x)\mathrm{d}x=\int\limits_{2}^{\infty} \frac{1}{x}\mathrm{d}x $\\
\begin{align*}
\int\limits_{2}^{\infty} g(x)=\lim\limits_ {t\to\infty}\int\limits_{2}^{t}\frac{1}{x}\mathrm{d} x=\lim\limits_ {t\to\infty}\bigg[\ln(x)\bigg]^t_2=\lim\limits_ {t\to\infty}\bigg(\ln(t)-\ln(2)\bigg)=\infty
\end{align*}
\begin{align*}
\overset{Minork.}{\Rightarrow}\int\limits_{2}^{\infty} f(x)\mathrm{d}x=\int\limits_{2}^{\infty} \frac{2+\cos(x)}{\sqrt{x^2-1}}\mathrm{d}x\;\;\text{existiert nicht. }
\end{align*}

%	\chapter{Fourieranalysis}
\setcounter{section}{12}
\setcounter{subsection}{-1}
\setcounter{subsubsection}{-1}
\setcounter{secnumdepth}{4}
\subsection{Einführung in die Fourieranalysis}
\href{https://youtu.be/7Dge4dbGn78}{Video 01}\\
$\bullet\;\;$Taylorapproximation\\
\begin{quote}
	Mittels der Taylorapproximationen werden diffbare Funktionen mittels Polynome lokal in der Nähe von $ x_0 $ approximieren.\\


\begin{center}
		$ f(x)=\ln(x) \approx \mathrm{T}_{n,x_0}$ \\
	\begin{tikzpicture}[xscale=1, yscale=1]
	
	% Gitternetzlinien
	\draw[xstep=1cm, ystep=1cm, very thin, color=lightgray] (-1,4) grid (4,-2);
	% x-Achse
	\draw[->] (-1,0) -- (4,0) node[below,xshift=-1em] {$x$}; 
	%Zahlen auf x-Achse
	%\foreach \x in {-3,...,-1}
	%\draw[shift={(\x,0)},color=black] (0pt,2pt) -- (0pt,-2pt) node[below] {\footnotesize $\x$};
	\foreach \x in {2,...,2}
	\draw[shift={(\x,0)},color=black] (0pt,2pt) -- (0pt,-2pt) node[below] {\footnotesize $\x$};
	% y-Achse 
	\draw[->] (0,-2) -- (0,4) node[left,yshift=-1em] {$f(x)$};%node[above left]
	%Zahlen auf y-Achse
	\foreach \y in {-1,...,-1}
	\draw[shift={(0,\y)},color=black] (2pt,0pt) -- (-2pt,0pt) node[left] {\footnotesize $\y$};
	\foreach \y in {1,...,3}
	\draw[shift={(0,\y)},color=black] (2pt,0pt) -- (-2pt,0pt) node[left] {\footnotesize $\y$};
	%Ursprung
	\draw[color=black] (0pt,-10pt) node[right] {\footnotesize $0$};
	
\draw[color=black] (0.8,-10pt) node[right] {\footnotesize $x_0$};
	
	\draw[thick,color=blue] plot[samples=200, domain=0.2:3] 
	(\x,{ln(\x)}) node[below, xshift=1em]{$\ln(x)$};
	
	\draw[thick,color=red] plot[samples=200, domain=0.2:3] 
	(\x,{\x-1}) node[below, xshift=1em]{$\mathrm{T}_{1,x_0}$};
		\draw[thick,color=red] plot[samples=200, domain=0.2:3] 
	(\x,{\x-1-0.5*(\x-1)^2}) node[below, xshift=1em]{$\mathrm{T}_{2,x_0}$};
	
	\end{tikzpicture}\\
	
\end{center}
\end{quote}
$\bullet\;\;$Fourieranalysis\\
\begin{quote}
	Mittels der Fourieranalysis werden periodische Funktionen mittels Sinus und Cosinus global approximiert.\\
	\begin{center}
	\begin{tikzpicture}[xscale=0.5, yscale=0.5]
	
	% Gitternetzlinien
	\draw[xstep=1cm, ystep=1cm, very thin, color=lightgray] (-14,7) grid (14,-2);
	% x-Achse
	\draw[->] (-14,0) -- (14,0) node[below,xshift=-1em] {$x$}; 
	%Zahlen auf x-Achse
	%\foreach \x in {-3,...,-1}
	%\draw[shift={(\x,0)},color=black] (0pt,2pt) -- (0pt,-2pt) node[below] {\footnotesize $\x$};
	%\foreach \x in {1,...,3}
	%\draw[shift={(\x,0)},color=black] (0pt,2pt) -- (0pt,-2pt) node[below] {\footnotesize $\x$};
	% y-Achse 
	\draw[->] (0,-2) -- (0,7) node[left,yshift=-1em] {$f(x)$};%node[above left]
	%Zahlen auf y-Achse
	%\foreach \y in {-1,...,-1}
	%\draw[shift={(0,\y)},color=black] (2pt,0pt) -- (-2pt,0pt) node[left] {\footnotesize $\y$};
	%\foreach \y in {1,...,3}
	%\draw[shift={(0,\y)},color=black] (2pt,0pt) -- (-2pt,0pt) node[left] {\footnotesize $\y$};
	%Ursprung
	\draw[color=black] (0pt,-10pt) node[right] {\footnotesize $0$};
	

	\draw[shift={(6.28,0)},color=black] (0pt,2pt) -- (0pt,-2pt) node[below] {\footnotesize $\mathrm{T}$};
	\draw[shift={(-2*pi,0)},color=black] (0pt,2pt) -- (0pt,-2pt) node[below] {\footnotesize $\mathrm{T}$};
	
	\draw[shift={(-4*pi,0)},color=black] (0pt,2pt) -- (0pt,-2pt) node[below] {\footnotesize $-2\mathrm{T}$};
	\draw[shift={(4*pi,0)},color=black] (0pt,2pt) -- (0pt,-2pt) node[below] {\footnotesize $2\mathrm{T}$};
	\draw[thick,color=red] plot[samples=200, domain=-14:14] 
	(\x,{pi-2*(sin(\x r)+0.5*sin((2*\x )r)+0.33*sin((3*\x )r)+0.25*sin((4*\x )r)) }) node[below, xshift=-3em]{$\phi_n(x)$};
	
		\draw[thick,color=blue] plot[samples=200, domain=0:2*pi] 
	(\x,{\x } );
		\draw[thick,color=blue] plot[samples=200, domain=2*pi:4*pi] 
	(\x,{\x-2*pi } );
		\draw[thick,color=blue] plot[samples=200, domain=-2*pi:0] 
	(\x,{\x +2*pi} );
	 	\draw[thick,color=blue] plot[samples=200, domain=-4*pi:-2*pi] 
	 (\x,{\x +4*pi} ) node[below, xshift=1em]{$g(x)$};
	\end{tikzpicture}\\	
	\end{center}
Die Funktion $ g(x)=x $, die  um T immer wieder periodisch wiederholt wird kann mittels einer periodischen Funktion $\phi_n(x)$ approximiert werden.\\
\end{quote}
Definition (Periodizität):\\
Eine Funktion hat die Periode T mit $ \mathrm{T}>0 $, wenn gilt:\\
\begin{align*}
f(x+T)=f(x)\;\;\text{ für alle}\; x
\end{align*}
Dann nennet man auch $ f $: T-periodisch.
\\
\\
Beispiel:\\
Der Sinus und der Cosinus sind $2\pi$-periodisch, denn es gilt:\\
\begin{align*}
\sin(x+2\pi)=\sin(x)&&\cos(x+2\pi)=\cos(x)
\end{align*}
Der Sinus und der Cosinus sind $4\pi$-periodisch, denn es gilt:\\
\begin{align*}
\sin(x+2\pi+2\pi)=\sin(x+2\pi)=\sin(x)&&\cos(x+2\pi+2\pi)=\cos(x+2\pi)=\cos(x)
\end{align*}
Allgemein gilt dann: Der Sinus und Cosinus sind $ k2\pi $-periodisch für $ k\in\mathbb{N}\setminus \{0\} $\\
\\
Betrachte: $ \mathrm{T}=\frac{2\pi}{\omega} $\\
Dann sind $ \sin(k\omega x) $ und $\cos( k\omega x) $ T-periodisch.\\
Es gilt:\\
\begin{align*}
\sin(k\omega(x+\mathrm{T}))&=\sin (k\omega x+k\omega \mathrm{T})\;\;\arrowvert\text{\tiny $\mathrm{T}=\frac{2\pi}{\omega} $ }\\
&=\sin (k\omega x+k2\pi)=\sin (k\omega x) \;\;\text{für }\;k\in \mathbb{Z}\\
&\Rightarrow \text{ T-periodisch}
\end{align*}
Bemerkung:\\
Die Funktion $ \sin(k\omega x) $ ist für $ k=1 $ die Grundschwingung.\\
Für $ k=2,3,4... $ enstehen somit alle Oberschwingungen. Diese Funktionen: $ \sin(k\omega x)$ und $ \cos(k\omega x)  $ sind dann die Grundbausteine für die Fourieranalysis.\\
\\
Ebenfalls gilt:\\
\begin{align*}
\mathrm{e}^{\mathrm{i}x}=\cos (x)+\mathrm{i}\sin (x)
\end{align*}
Dann ist die Euler-Form auch T-periodisch:\\
\begin{align*}
\mathrm{e}^{\mathrm{i}k\omega x}= \cos(k\omega x)+\mathrm{i}\sin(k\omega x)
\end{align*}
\newpage
\noindent
\subsection{ Das reelle Fourierpolynom}
\href{https://youtu.be/kBD9vmv_Gio}{Video 02}\\
Sei $ f:\vee\longmapsto\mathbb{R} $ oder $  \mathbb{C} $ eine T-periodische Funktion und stückweise monoton.\\
\begin{center}
Jeder Teilbereich ist monoton\\
\begin{tikzpicture}[xscale=1, yscale=1]

% Gitternetzlinien
\draw[xstep=1cm, ystep=1cm, very thin, color=lightgray] (0,4) grid (8,0);
% x-Achse
\draw[->] (0,0) -- (8,0) node[below,xshift=-1em] {$x$}; 
%Zahlen auf x-Achse
%\foreach \x in {-3,...,-1}
%\draw[shift={(\x,0)},color=black] (0pt,2pt) -- (0pt,-2pt) node[below] {\footnotesize $\x$};
%\foreach \x in {1,...,3}
%\draw[shift={(\x,0)},color=black] (0pt,2pt) -- (0pt,-2pt) node[below] {\footnotesize $\x$};
% y-Achse 
\draw[->] (0,0) -- (0,4) node[left,yshift=-1em] {$f(x)$};%node[above left]
%Zahlen auf y-Achse
%\foreach \y in {-1,...,-1}
%\draw[shift={(0,\y)},color=black] (2pt,0pt) -- (-2pt,0pt) node[left] {\footnotesize $\y$};
%\foreach \y in {1,...,3}
%\draw[shift={(0,\y)},color=black] (2pt,0pt) -- (-2pt,0pt) node[left] {\footnotesize $\y$};
%Ursprung

\draw[shift={(7,0)},color=black] (0pt,2pt) -- (0pt,-2pt) node[below] {\footnotesize $\mathrm{T}$};
\draw[shift={(0,0)},color=black] (0pt,2pt) -- (0pt,-2pt) node[below] {\footnotesize $\mathrm{0}$};
\draw[thick,color=blue] plot[samples=200, domain=0:1] 
(\x,{2});
\draw[thick,color=blue] plot[samples=200, domain=1:3] 
(\x,{\x});

\draw[thick,color=blue] plot[samples=200, domain=3:5] 
(\x,{-\x+6});
\draw[thick,color=blue] plot[samples=200, domain=5:7] 
(\x,{1});
\end{tikzpicture}\\	

\end{center}
Sei $ \omega =\frac{2\pi}{T} $, dann definieren wir für $ k\in\mathbb{N} $ die Fourierkoeffizienten durch:\\
\begin{center}
	\fbox{\parbox  {50mm} {$a_0=\frac{2}{T}\int\limits_{0}^{\mathrm{T}}f(x)\mathrm{d}x \\a_k=\frac{2}{\mathrm{T}}\int\limits_{0}^{\mathrm{T}}f(x)\cos(k\omega x)\mathrm{d}x   $\\ $b_k=\frac{2}{\mathrm{T}}\int\limits_{0}^{\mathrm{T}}f(x)\sin(k\omega x)\mathrm{d}x   $}}
\end{center}
Dann ist das n-Fourierpolynom (n-ter Ordnung) von der Form:\\
\begin{center}
\boxed{ \phi_n(x)=\frac{a_0}{2} +\sum\limits_{k=1}^{n}a_k\cos(k\omega x)+b_k\sin(k\omega x)}	
\end{center}
Nach der Approximation gilt dann:\\
\begin{align*}
&f(x)\approx \phi_n(x)\\
&f(x)=\lim\limits_ {n\to\infty}\phi_n(x)
\end{align*}
\addtocounter{subsubsection}{1}
\subsubsection{ Beispiel zum reellen Fourierpolynom}
\href{https://youtu.be/EsfssnLs8To}{Video 03}\\
Berachte die T-periodische Funktion $ f(x)=\cos(x)+2\sin(x) $ mit $ 2\pi=\mathrm{T}=\frac{2\pi}{\omega} $ und berechne das 1.-Fourierpolynom.\\
Bemerkung:\\
\begin{align*}
\text{für } \omega=1 \;\;\text{gilt: } \mathrm{T}=\frac{2\pi}{\omega}\Rightarrow \mathrm{T}=2\pi
\end{align*}
Es gilt:\\
\begin{align*}
& \phi_n(x)=\frac{a_0}{2} +\sum\limits_{k=1}^{n}a_k\cos(k\omega x)+b_k\sin(k\omega x)\\
&\text{mit }\;\;\omega=1 \;\;\text{und }\;\;k=1\;\;\text{ergibt sich das 1.-Fourierpolynom: }\\
& \phi_n(x)=\frac{a_0}{2} +a_1\cos(x)+b_1\sin( x)\\
\end{align*}
\underline{Berechnung von $ a_0 $:}\\
\begin{align*}
a_0&=\frac{2}{T}\int\limits_{0}^{\mathrm{T}}f(x)\mathrm{d}x\\
&=\frac{2}{2\pi}\int\limits_{0}^{2\pi} \cos(x)+2\sin(x)\mathrm{d}x\\
&=\frac{1}{\pi}\bigg[\sin(x)-2\cos(x)\bigg]^{2\pi}_0\\
&=\frac{1}{\pi}\bigg(\sin(2\pi)-2\cos(2\pi)-\big(\sin(0)-2\cos(0)\big)\bigg)\\
a_0&=\frac{1}{\pi}\cdot0=0
\end{align*}
Bemerkung:\\
Das Integral von $ \cos(x), \sin(x) $ mit den Grenzen über genau einer Periode ergibt die Fläche: A=0.\\
\begin{center}
	
Fläche A=0\\
\begin{tikzpicture}[xscale=1, yscale=1]

% Gitternetzlinien
\draw[xstep=1cm, ystep=1cm, very thin, color=lightgray] (0,2) grid (7,-2);
% x-Achse
\draw[->] (0,0) -- (7,0) node[below,xshift=-1em] {$x$}; 
%Zahlen auf x-Achse
%\foreach \x in {-3,...,-1}
%\draw[shift={(\x,0)},color=black] (0pt,2pt) -- (0pt,-2pt) node[below] {\footnotesize $\x$};
\foreach \x in {1,...,6}
\draw[shift={(\x,0)},color=black] (0pt,2pt) -- (0pt,-2pt) node[below] {\footnotesize $\x$};
% y-Achse 
\draw[->] (0,-2) -- (0,2) node[left,yshift=-1em] {$f(x)$};%node[above left]
%Zahlen auf y-Achse
\foreach \y in {-1,...,-1}
\draw[shift={(0,\y)},color=black] (2pt,0pt) -- (-2pt,0pt) node[left] {\footnotesize $\y$};
\foreach \y in {1,...,1}
\draw[shift={(0,\y)},color=black] (2pt,0pt) -- (-2pt,0pt) node[left] {\footnotesize $\y$};
%Ursprung
\draw[color=black] (0pt,-10pt) node[right] {\footnotesize $0$};

\draw[thick,color=blue] plot[samples=200, domain=0:6.28] 
(\x,{sin(\x r) }) node[ xshift=-6em,yshift=5] {$\sin(x)$};

\draw[thick,color=red] plot[samples=200, domain=0:6.28] 
(\x,{cos(\x r}) node[ yshift=1em,xshift=-4]{$\cos(x)$};


\end{tikzpicture}\\

	
\end{center}
\underline{Berechnung von $ a_1 $:}\\
\begin{align*}
a_1&=\frac{2}{2\pi}\int\limits_{0}^{\mathrm{2\pi}}f(x)\cos( x)\mathrm{d}x  \\
&=\frac{1}{\pi}\int\limits_{0}^{2\pi} \big(\cos(x)+2\sin(x)\big)\cos(x)\mathrm{d}x\\
&=\frac{1}{\pi}\int\limits_{0}^{2\pi} \big(\cos^2(x)+2\sin(x)\cos(x)\big)\mathrm{d}x\\
\end{align*}
\\
\\
\\
\\
\\
\\
\begin{quote}
NR.: Stammfunktionen von $ \int\limits_{}^{}\big(\cos(x)\sin(x)\big) \mathrm{d}x $ und $ \int\limits_{}^{} \cos^2(x)\mathrm{d}x $ bestimmen.\\
	$ \bullet\;\;  \int\limits_{}^{}\big(\cos(x)\sin(x)\big)\mathrm{d}x$ (Substitution)
	\begin{align*}
	&\int\limits_{}^{}\big(\cos(x)\sin(x)\big)\mathrm{d}x\\
	&u=\sin(x)\\
	&\frac{\mathrm{d}u}{\mathrm{d}x}=\frac{\sin(x)}{\mathrm{dx}}=\cos(x)\\
	&\Leftrightarrow \mathrm{d}x=\frac{du}{\cos(x)}\\
	&\Rightarrow\\
	&\int\limits_{}^{} \big(\sin(x)\cos(x)\big)\mathrm{d}x=\int\limits_{}^{} u \cdot \cos(x)\frac{\mathrm{d}u}{\cos(x)}=\int\limits_{}^{} u\cdot \mathrm{d}u=\frac{1}{2}u^2\\
	&\text{Rücksubstitution: }\;u=\sin(x)\\
	&\int\limits_{}^{} \big(\sin(x)\cos(x)\big)\mathrm{d}x=\frac{1}{2}\sin^2(x)+\mathrm{c}
	\end{align*}
		$ \bullet\;\;  \int\limits_{}^{}\cos^2(x)\mathrm{d}x$ (Partialintegration)\\
		\begin{align*}
		&\int\limits_{}^{} \cos^2(x)\mathrm{d}x=\int\limits_{}^{} \big(\cos(x)\cos(x)\big)\mathrm{d}x\\
		&u=\cos(x)\Rightarrow u'=-\sin(x)\\
		&v'=\cos(x)\Rightarrow v=\sin(x)\\
		&\int\limits_{}^{} \cos^2(x)\mathrm{d}x=\cos(x)\cdot \sin(x)-\int\limits_{}^{} -\sin(x)^2(x)\mathrm{d}x\;\;\arrowvert\text{\tiny $ \cos^2(x)+\sin^2(x)=1\Leftrightarrow \sin^2(x)=1-\cos^2(x) $ }\\
		&\int\limits_{}^{} \cos^2(x)\mathrm{d}x=\cos(x)\sin(x)+\int\limits_{}^{} \big(1-\cos^2(x)\big)\mathrm{d}x\\
		&\int\limits_{}^{} \cos^2(x)\mathrm{d}x=\cos(x)\sin(x)+\int\limits_{}^{} 1\mathrm{d}x-\int\limits_{}^{} \cos^2(x)\mathrm{d}x\\
		\Leftrightarrow&2\int\limits_{}^{} \cos^2(x)\mathrm{d}x=\cos(x)\sin(x)+x\\
		\Leftrightarrow&\int\limits_{}^{} \cos^2(x)\mathrm{d}x=\frac{1}{2}\big(x+\cos(x)\sin(x)\big)+\mathrm{c}\\
		&\text{Es gilt: }\;\;\sin(x+x)=\sin(x)\cos(x)+\sin(x)\cos(x)\;\;|\text{ \tiny Additiontheorem}\\
		&\sin(2x)=2\sin(x)\cos(x)\\
		&\Leftrightarrow \cos(x)\sin(x)=\frac{1}{2}\sin(2x)\\
		&\text{Dies eingesetzt ergibt dann die Stammfunktion: }\\
		&\int\limits_{}^{} \cos^2(x)\mathrm{d}x=\frac{1}{2}\bigg(x+\frac{1}{2}\sin(2x)\bigg)+\mathrm{c}\\
		\end{align*}
\end{quote}
Nun weiter mit der Berechnung von $ a_1 $:\\
\begin{align*}
a_1&=\frac{1}{\pi}\int\limits_{0}^{2\pi} \big(\cos^2(x)+2\sin(x)\cos(x)\big)\mathrm{d}x\\
&=\frac{1}{\pi}\bigg[\frac{x}{2}+\frac{\sin(2x)}{4}+\sin^2(x)\bigg]^{2\pi}_0\\
&=\frac{1}{\pi}\big(\pi+0-(0)\big)=\frac{\pi}{\pi}=1\\
a_1&=1
\end{align*}
\underline{ Berechnung von $ b_1 $:}\\
\begin{align*}
b_1&=\frac{2}{2\pi}\int\limits_{0}^{\mathrm{2\pi}}f(x)\sin( x)\mathrm{d}x  \\
&=\frac{1}{\pi}\int\limits_{0}^{2\pi} \big(\cos(x)+2\sin(x)\big)\sin(x)\mathrm{d}x\\
&=\frac{1}{\pi}\int\limits_{0}^{2\pi} \big(\cos(x)\sin(x)+2\sin^2(x)\big)\mathrm{d}x\\
\end{align*}
\begin{quote}
NR.: Stammfunktionen von $ \int\limits_{}^{}\sin^2(x) \mathrm{d}x $ bestimmen.\\
\begin{align*}
&\int\limits_{}^{} \sin(x)\sin(x)\mathrm{d}x\\
&u=\sin(x)\Rightarrow u'=\cos(x)\\
&v'=\sin(x)\Rightarrow v=-\cos(x)\\
&\int\limits_{}^{} \sin^2(x)\mathrm{d}x=\sin(x)\cdot -\cos(x)-\int\limits_{}^{} -\cos^2(x)\mathrm{d}x\;\;\lvert\text{\tiny $\cos^2(x)+\sin^2(x)=1\Leftrightarrow \cos^2(x)=1-\sin^2(x)$ }\\
&\int\limits_{}^{} \sin^2(x)\mathrm{d}x=-\sin(x)\cos(x)+\int\limits_{}^{} 1-\sin^2(x)\mathrm{d}x\\
&\int\limits_{}^{} \sin^2(x)\mathrm{d}x=-\sin(x)\cos(x)+\int\limits_{}^{} 1\mathrm{d}x-\int\limits_{}^{} \sin^2(x)\mathrm{d}x\\
\Leftrightarrow&2\int\limits_{}^{} \sin^2(x)\mathrm{d}x=-\sin(x)\cos(x)+x\\
\Leftrightarrow&\int\limits_{}^{} \sin^2(x)\mathrm{d}x=\frac{1}{2}\big(x-\sin(x)\cos(x)\big)\;\;\text{\tiny$\cos(x)\sin(x)=\frac{1}{2}\sin(2x)  $ }\\
&\int\limits_{}^{} \sin^2(x)\mathrm{d}x=\frac{1}{2}\bigg(x-\frac{1}{2}\sin(2x)\bigg)+\mathrm{c}
\end{align*}
\end{quote}
Nun weiter mit der Berechnung von $ b_1 $:\\
\begin{align*}
&=\frac{1}{\pi}\int\limits_{0}^{2\pi} \big(\cos(x)\sin(x)+2\sin^2(x)\big)\mathrm{d}x\\
&=\frac{1}{\pi}\bigg[\frac{1}{2}\sin^2(x)+2\bigg(\frac{x}{2}-\frac{\sin(2x)}{4}\bigg)\bigg]^{2\pi}_0\\
&=\frac{1}{\pi}\bigg[0+2\bigg(\frac{2\pi}{4}\bigg)\bigg]=\frac{1}{\pi}\cdot 2\pi=2\\
b_1&=2
\end{align*}
Das 1.-Fourierpolynom lautet nun:\\
\begin{align*}
\phi_1(x)&=\frac{a_0}{2} +a_1\cos(kx)+b_1\sin(k x)\\
&=0+\cos(x)+2\sin(x)\\
\phi_1(x)&=\cos(x)+2\sin(x)=f(x)
\end{align*}
Das Fourierpolynom ist gleich der Funktion: $ f(x) $, weil $ f(x) $ eine Linearkombination von trigometrischen Funktion ist und das Fourierpolynom aus trigometrischen Funktionen erzeugt wird.\\
\\
Für das k.-Fourierpolynom $ k>1 $ mit $ f(x)=\cos(x)+2\sin(x) $ ergibt sich:\\
Es gilt:\\
\begin{align*}
& \phi_n(x)=\frac{a_0}{2} +\sum\limits_{k=2}^{n}a_k\cos(k\omega x)+b_k\sin(k\omega x)\\
\end{align*}
Die Fourierkoeffizenten $a_0, a_k, b_k $ für $ k>1 $ werden Null, weil ja schon $ f(x)=\phi(x) $ ist und daraus folgt, dass gilt:\\
\begin{align*}
& \phi_n(x)=\frac{a_0}{2} +\sum\limits_{k=2}^{n}a_k\cos(k\omega x)+b_k\sin(k\omega x)=0\\
&\Rightarrow \phi_n(x)=\cos(x)+2\sin(x)=f(x)\\
\end{align*}
\subsubsection{Konvergenz des Fourierpolynoms}
\href{https://youtu.be/UCREQw2VLo8}{Video 04}\\
Sei $ f $ eine stückweise, monotone Funktion mit der Periode $ \mathrm{T}=\frac{2\pi}{\omega} $ und $ \phi_n(x) $ das n-te Fourierpolynom von $ f $.\\
\\
\\
\newpage
Bemerkung: Stückweise monoton bedeutet:\\
\begin{center}
	Jeder Teilbereich ist monoton\\
	\begin{tikzpicture}[xscale=1, yscale=1]
	
	% Gitternetzlinien
	\draw[xstep=1cm, ystep=1cm, very thin, color=lightgray] (0,4) grid (8,0);
	% x-Achse
	\draw[->] (0,0) -- (8,0) node[below,xshift=-1em] {$x$}; 
	%Zahlen auf x-Achse
	%\foreach \x in {-3,...,-1}
	%\draw[shift={(\x,0)},color=black] (0pt,2pt) -- (0pt,-2pt) node[below] {\footnotesize $\x$};
	%\foreach \x in {1,...,3}
	%\draw[shift={(\x,0)},color=black] (0pt,2pt) -- (0pt,-2pt) node[below] {\footnotesize $\x$};
	% y-Achse 
	\draw[->] (0,0) -- (0,4) node[left,yshift=-1em] {$f(x)$};%node[above left]
	%Zahlen auf y-Achse
	%\foreach \y in {-1,...,-1}
	%\draw[shift={(0,\y)},color=black] (2pt,0pt) -- (-2pt,0pt) node[left] {\footnotesize $\y$};
	%\foreach \y in {1,...,3}
	%\draw[shift={(0,\y)},color=black] (2pt,0pt) -- (-2pt,0pt) node[left] {\footnotesize $\y$};
	%Ursprung
	
	\draw[shift={(7,0)},color=black] (0pt,2pt) -- (0pt,-2pt) node[below] {\footnotesize $\mathrm{T}$};
	\draw[shift={(0,0)},color=black] (0pt,2pt) -- (0pt,-2pt) node[below] {\footnotesize $\mathrm{0}$};
	\draw[thick,color=blue] plot[samples=200, domain=0:1] 
	(\x,{2});
	\draw[thick,color=blue] plot[samples=200, domain=1:3] 
	(\x,{\x});
	
	\draw[thick,color=blue] plot[samples=200, domain=3:5] 
	(\x,{-\x+6});
	\draw[thick,color=blue] plot[samples=200, domain=5:7] 
	(\x,{1});
	\end{tikzpicture}\\	
	
\end{center}
Dann gilt (ohne Beweis):\\
\begin{align*}
&f(x)=\phi_n(x)+\mathrm{R}_n(x)=\frac{a_0}{2}+\sum\limits_{k=1}^{n}a_k\cos(k\omega x)+b_k\sin(k\omega x)+\mathrm{R}_n(x)\\
&\mathrm{R}_n(x) = \;\;\text{Fehler }\\
&a_k, b_k \;=\;\;\text{Fourierkoeffizienten }\\
\end{align*}
Die Funktion: $ \mathrm{R}_n(x)=f(x) -\phi_n(x)$ erfüllt das Folgende:\\
\begin{align*}
\frac{2}{\mathrm{T}}\int\limits_{0}^{\mathrm{T}} \big(\mathrm{R}_n(x)\big)^2 \mathrm{d}x&=\frac{2}{\mathrm{T}}\int\limits_{0}^{\mathrm{T}} \bigg(f(x) -\phi_n(x)\bigg)^2 \mathrm{d}x\\
&\;\;\cdots\text{nach Umformung ergibt sich dann: }\\
&=\frac{2}{T}\int\limits_{0}^{\mathrm{T}} f^2(x)\mathrm{d}x-\bigg(\frac{a_0^2}{2}+\sum\limits_{k=1}^{n}(a^2_k+b^2_k)\bigg)\\
&\Leftrightarrow \;\;\text{für $ n\mapsto \infty  $ konvergiert dies gegen 0 und wird zur Pasevalschen Gleichung}\\
\frac{2}{T}\int\limits_{0}^{\mathrm{T}} f^2(x)\mathrm{d}x&=\bigg(\frac{a_0^2}{2}+\sum\limits_{k=1}^{\infty}(a^2_k+b^2_k)\bigg)\\
\end{align*}
Daraus folgt, dass der Fehler im quatratischen Mittel gegen 0 konvergiert.\\
\\
Die Parcevalsche Gleichung lautet:\\
\begin{center}
	\boxed{\frac{2}{T}\int\limits_{0}^{\mathrm{T}} f^2(x)\mathrm{d}x=\bigg(\frac{a_0}{2}^2+\sum\limits_{k=1}^{\infty}(a^2_k+b^2_k)\bigg)}
\end{center}
Ist $ f $ stetig in $ x $, so gilt:\\
\begin{align*}
\lim\limits_ {n\to\infty}\phi(x)=f(x)\Leftrightarrow\lim\limits_ {n\to\infty}\mathrm{R}(x)=0
\end{align*}
\newpage
\noindent
\subsubsection{ Beispiel zur Fourierapproximation}
\href{https://youtu.be/JnXjrvjlXg0}{Video 05}\\
Gegeben ist die $ 2\pi- $periodische Funktion $ f(x)=x $ für $ 0\leq x\leq 2\pi $.\\
Berechne das n-te Fourierpolynom.\\
\\
Wie sieht die Funktion aus ?\\
	\begin{center}
	\begin{tikzpicture}[xscale=0.5, yscale=0.5]
	
	% Gitternetzlinien
	\draw[xstep=1cm, ystep=1cm, very thin, color=lightgray] (-14,7) grid (14,-2);
	% x-Achse
	\draw[->] (-14,0) -- (14,0) node[below,xshift=-1em] {$x$}; 
	%Zahlen auf x-Achse
	%\foreach \x in {-3,...,-1}
	%\draw[shift={(\x,0)},color=black] (0pt,2pt) -- (0pt,-2pt) node[below] {\footnotesize $\x$};
	%\foreach \x in {1,...,3}
	%\draw[shift={(\x,0)},color=black] (0pt,2pt) -- (0pt,-2pt) node[below] {\footnotesize $\x$};
	% y-Achse 
	\draw[->] (0,-2) -- (0,7) node[left,yshift=-1em] {$f(x)$};%node[above left]
	%Zahlen auf y-Achse
	%\foreach \y in {-1,...,-1}
	%\draw[shift={(0,\y)},color=black] (2pt,0pt) -- (-2pt,0pt) node[left] {\footnotesize $\y$};
	%\foreach \y in {1,...,3}
	%\draw[shift={(0,\y)},color=black] (2pt,0pt) -- (-2pt,0pt) node[left] {\footnotesize $\y$};
	%Ursprung
	\draw[color=black] (0pt,-10pt) node[right] {\footnotesize $0$};
	
	
	\draw[shift={(6.28,0)},color=black] (0pt,2pt) -- (0pt,-2pt) node[below] {\footnotesize $\mathrm{2\pi}$};
	\draw[shift={(-2*pi,0)},color=black] (0pt,2pt) -- (0pt,-2pt) node[below] {\footnotesize $\mathrm{2\pi}$};
	
	\draw[shift={(-4*pi,0)},color=black] (0pt,2pt) -- (0pt,-2pt) node[below] {\footnotesize $\mathrm{-4\pi}$};
	\draw[shift={(4*pi,0)},color=black] (0pt,2pt) -- (0pt,-2pt) node[below] {\footnotesize $\mathrm{4\pi}$};
	\draw[thick,color=red] plot[samples=200, domain=-14:14] 
	(\x,{pi-2*(sin(\x r)+0.5*sin((2*\x )r)+0.33*sin((3*\x )r)+0.25*sin((4*\x )r)) }) node[below, xshift=-3em]{$\phi_n(x)$};
	
	\draw[thick,color=blue] plot[samples=200, domain=0:2*pi] 
	(\x,{\x } );
	\draw[thick,color=blue] plot[samples=200, domain=2*pi:4*pi] 
	(\x,{\x-2*pi } );
	\draw[thick,color=blue] plot[samples=200, domain=-2*pi:0] 
	(\x,{\x +2*pi} );
	\draw[thick,color=blue] plot[samples=200, domain=-4*pi:-2*pi] 
	(\x,{\x +4*pi} ) node[below, xshift=1em]{$f(x)$};
	\end{tikzpicture}\\	
\end{center}
Berechnung der Koeffizienten $ a_0, a_k, b_k $ vom n-ten Fourierpolynoms:\\
\begin{align*}
\phi_n(x)=\frac{a_0}{2}+\sum\limits_{k=1}^{n}\big(a_k\cos(k\omega x)+b_k\sin(k\omega x)\big)
\end{align*}
Eine Periode ist zwischen 0 und $ 2\pi $ daraus folgt:\\
\begin{align*}
2\pi=\mathrm{T}=\frac{2\pi}{\omega}\Rightarrow \omega=1
\end{align*}
$ \bullet $\; Bestimmung von $ a_0 $:\\
\begin{align*}
a_0=\frac{2}{\mathrm{T}}\int\limits_{0}^{\mathrm{T}} f(x)\mathrm{d}x&=\frac{2}{2\pi}\int\limits_{0}^{2\pi} x\mathrm{d}x=\frac{1}{\pi}\bigg[\frac{x^2}{2}\bigg]^{2\pi}_0\\
&=\frac{1}{\pi}\bigg(\frac{(2\pi)^2}{2}-\frac{0^2}{2}\bigg)=\frac{1}{\pi}\frac{(2\pi)^2}{2}=\frac{4\pi^2}{\pi2}=2\pi\\
a_0&=2\pi
\end{align*}
\href{https://youtu.be/i4RYWjoYFAA}{Video 06}\\
$ \bullet $\;Bestimmung von $ a_k $:
\begin{align*}
a_k=\frac{2}{\mathrm{T}}\int\limits_{0}^{\mathrm{T}} f(x)\cdot \cos(kx)\mathrm{d}x=\frac{2}{2\pi}\int\limits_{0}^{2\pi} x\cdot \cos(kx)\mathrm{d}x
\end{align*}
\begin{quote}
	NR.: partielle Integration von: $ \int\limits_{}^{} x\cdot +\cos(x)\mathrm{d}x $\\
	\begin{align*}
	&u(x)=x\;\;\;\;\;\;\;\;\;\;\Rightarrow u'(x)=1\\
	&v'(x)=\cos(kx)\Rightarrow v(x)=\frac{1}{k}\sin(kx)\\
	&\int\limits_{}^{}x\cdot \cos(kx)\mathrm{d}x=\frac{1}{k}x\cdot \sin(kx)-\int\limits_{}^{} \frac{1}{k}\sin(kx)\mathrm{d}x\\
	&\int\limits_{}^{}x\cdot \cos(kx)\mathrm{d}x=\frac{1}{k}x\cdot \sin(kx)+\frac{1}{k^2}\cos(kx)
	\end{align*}
\end{quote}
\begin{align*}
a_k&=\frac{1}{\pi}\int\limits_{0}^{\pi} x\cdot \cos(kx)\mathrm{d}x\\
&=\frac{1}{\pi}\bigg( \frac{1}{k}\bigg[x\cdot \sin(kx)\bigg]^{2\pi}_0+\frac{1}{k^2}\bigg[\cos(kx)\bigg]^{2\pi}_0\bigg)\\
&=\frac{1}{\pi}\bigg(\frac{1}{k}(2\pi\cdot 0-0\cdot 0)+\frac{1}{k^2}(1-1)\bigg)=0\\
a_k&=0\;\;\text{für }\;\;k>0
\end{align*}
\href{https://youtu.be/Awz46xIlbyc}{Video 07}\\
$ \bullet $\;Bestimmung von $ b_k $:\\
\begin{align*}
b_k=\frac{2}{\mathrm{T}}\int\limits_{0}^{\mathrm{T}} f(x)\cdot \sin(kx)\mathrm{d}x=\frac{2}{2\pi}\int\limits_{0}^{2\pi} x\cdot \sin(kx)\mathrm{d}x
\end{align*}
\begin{quote}
	N.R.: partielle Integration von: $ \int\limits_{}^{} x\cdot \sin(x)\mathrm{d}x $\\
	\begin{align*}
	&u(x)=x\;\;\;\;\;\;\;\;\;\;\Rightarrow u'(x)=1\\
	&v'(x)=\sin(kx)\Rightarrow v(x)=-\frac{1}{k}\cos(kx)\\
	&\int\limits_{}^{}x\cdot \sin(kx)\mathrm{d}x=-\frac{1}{k}x\cdot \cos(kx)-\int\limits_{}^{} -\frac{1}{k}\cos(kx)\mathrm{d}x\\
	&\int\limits_{}^{}x\cdot \sin(kx)\mathrm{d}x=-\frac{1}{k}x\cdot \cos(kx)+\frac{1}{k^2}\sin(kx)
	\end{align*}
\end{quote}
\begin{align*}
b_k&=\frac{1}{\pi}\int\limits_{0}^{2\pi} x\cdot \sin(kx)\mathrm{d}x\\
&=\frac{1}{\pi}\bigg(-\frac{1}{k}\bigg[x\cdot \cos(kx)\bigg]^{2\pi}_0+\frac{1}{k^2}\bigg[\sin(kx)\bigg]^{2\pi}_0\bigg)\\
&=\frac{1}{\pi}\bigg(-\frac{1}{k}(2\pi-0)+\frac{1}{k^2}(0-0)\bigg)=\frac{-2\pi}{\pi k}\\
b_k&=-\frac{2}{k}
\end{align*}
Bemerkung:\\ Das Integral: $ \int\limits_{0}^{2\pi} \cos(x)\mathrm{d}x $ muß ja Null sein, weil die Summe der Flächen unter den Graphen Null über eine Periode ergeben muß. Analog auch der Sinus. \\
\\
\\
\\
\\
\\
\\
	\begin{center}
		Summe der Flächen ist Null.\\
	\begin{tikzpicture}[xscale=2, yscale=0.5]
	
	% Gitternetzlinien
	\draw[xstep=1cm, ystep=1cm, very thin, color=lightgray] (0,7) grid (7,-7);
	% x-Achse
	\draw[->] (0,0) -- (7,0) node[below,xshift=-1em] {$x$}; 
	%Zahlen auf x-Achse
	%\foreach \x in {-3,...,-1}
	%\draw[shift={(\x,0)},color=black] (0pt,2pt) -- (0pt,-2pt) node[below] {\footnotesize $\x$};
	%\foreach \x in {1,...,3}
	%\draw[shift={(\x,0)},color=black] (0pt,2pt) -- (0pt,-2pt) node[below] {\footnotesize $\x$};
	% y-Achse 
	\draw[->] (0,-7) -- (0,7) node[left,yshift=-1em] {$f(x)$};%node[above left]
	%Zahlen auf y-Achse
	%\foreach \y in {-1,...,-1}
	%\draw[shift={(0,\y)},color=black] (2pt,0pt) -- (-2pt,0pt) node[left] {\footnotesize $\y$};
	%\foreach \y in {1,...,3}
	%\draw[shift={(0,\y)},color=black] (2pt,0pt) -- (-2pt,0pt) node[left] {\footnotesize $\y$};
	%Ursprung
	\draw[color=black] (0pt,-10pt) node[right] {\footnotesize $0$};
	
	
	\draw[shift={(6.28,0)},color=black] (0pt,2pt) -- (0pt,-2pt) node[below] {\footnotesize $\mathrm{2\pi}$};
%	\draw[shift={(-2*pi,0)},color=black] (0pt,2pt) -- (0pt,-2pt) node[below] {\footnotesize $\mathrm{2\pi}$};
	
%	\draw[shift={(-4*pi,0)},color=black] (0pt,2pt) -- (0pt,-2pt) node[below] {\footnotesize $\mathrm{-4\pi}$};
%	\draw[shift={(4*pi,0)},color=black] (0pt,2pt) -- (0pt,-2pt) node[below] {\footnotesize $\mathrm{4\pi}$};
	\draw[thick,color=blue] plot[samples=200, domain=0:6.28] 
	(\x,{6*(cos(\x r)) }) node[below, xshift=-6em]{$\cos(x)$};
	\fill[color =blue,fill opacity=0.3] plot[domain=0:1.57] (\x,{6*(cos(\x r))})% 
	-- plot[domain=1.57:0] (\x,{0});% 
		\fill[color =blue,fill opacity=0.3] plot[domain=1.57:3.14] (\x,{6*(cos(\x r))})% 
	-- plot[domain=3.14:1.57] (\x,{0});% 
		\fill[color =green,fill opacity=0.3] plot[domain=3.14:4.71] (\x,{6*(cos(\x r))})% 
	-- plot[domain=4.71:3.14] (\x,{0});% 
		\fill[color =green,fill opacity=0.3] plot[domain=4.71:6.28] (\x,{6*(cos(\x r))})% 
	-- plot[domain=6.28:4.71] (\x,{0});% 

	\end{tikzpicture}\\	
\end{center}
Fazit: \\
Für das Fourierplynom mit $ \omega =1,\; a_0=2\pi,\; a_k=0 $ und $ b_k=-\frac{2}{k} $ gilt dann:\\
\begin{align*}
\phi_n(x)&=\frac{a_0}{2}+\sum\limits_{k=1}^{n}\big(a_k\cos(k\omega x)+b_k\sin(k\omega x)\big)\\
&=\pi+\sum\limits_{k=1}^{n}-\frac{2}{k}\sin(kx)\\
\phi_n(x)&=\pi-\sum\limits_{k=1}^{n}\frac{2}{k}\sin(kx)\\
\end{align*}
\href{https://youtu.be/7VbdjMYAqLM}{Video 08}\\
\underline {I. Konvergenzbetrachtung an der Stelle: $ x=\frac{\pi}{2} $.}\\
Ist $ f $ in $ x $ stetig, dann gilt:\\
\begin{align*}
\lim\limits_ {n\to\infty}\phi_n(x)=f(x)
\end{align*}
Für $ x=\frac{\pi}{2} $ und ist die Funktion: $ f(x)=x $ stetig, dann gilt:\\
\begin{align*}
&f(\genfrac{}{}{}{2}{\pi}{2 })=\lim\limits_ {n\to\infty}\phi_n(\genfrac{}{}{}{2}{\pi }{2 })\\
&\Rightarrow\\
&\frac{\pi}{2}=\lim\limits_ {n\to\infty}\bigg(\pi-\sum\limits_{k=1}^{n}\frac{2}{k}\sin(k\genfrac{}{}{}{2}{\pi }{2 })\bigg)\\
&\frac{\pi}{2}=\pi-\sum\limits_{k}^{\infty}\frac{2}{k}\cdot \sin(k\genfrac{}{}{}{2}{\pi }{2 })
\end{align*}
NR.:\\
Für $ \sum\limits_{k=1}^{\infty} \sin(k\genfrac{}{}{}{2}{\pi }{2})$ gilt:\\
	\begin{center}

	\begin{tikzpicture}[xscale=0.5, yscale=0.5]
	
	% Gitternetzlinien
	\draw[xstep=1cm, ystep=1cm, very thin, color=lightgray] (0,7) grid (20,-7);
	% x-Achse
	\draw[->] (0,0) -- (20,0) node[below,xshift=-1em] {$x$}; 
	%Zahlen auf x-Achse
	%\foreach \x in {-3,...,-1}
	%\draw[shift={(\x,0)},color=black] (0pt,2pt) -- (0pt,-2pt) node[below] {\footnotesize $\x$};
	%\foreach \x in {1,...,3}
	%\draw[shift={(\x,0)},color=black] (0pt,2pt) -- (0pt,-2pt) node[below] {\footnotesize $\x$};
	% y-Achse 
	\draw[->] (0,-7) -- (0,7) node[left,yshift=-1em] {$f(x)$};%node[above left]
	%Zahlen auf y-Achse
	%\foreach \y in {-1,...,-1}
	%\draw[shift={(0,\y)},color=black] (2pt,0pt) -- (-2pt,0pt) node[left] {\footnotesize $\y$};
	%\foreach \y in {1,...,3}
	%\draw[shift={(0,\y)},color=black] (2pt,0pt) -- (-2pt,0pt) node[left] {\footnotesize $\y$};
	%Ursprung
	\draw[shift={(0,5)},color=black] (2pt,0pt) -- (-2pt,0pt) node[left] {\footnotesize $1$};
	\draw[shift={(0,-5)},color=black] (2pt,0pt) -- (-2pt,0pt) node[left] {\footnotesize $-1$};
	\draw[color=black] (0pt,-10pt) node[right] {\footnotesize $0$};
	\draw[shift={(1.57,0)},color=black] (0pt,2pt) -- (0pt,-2pt) node[below] {\footnotesize $\frac{\pi}{2}$};
	\draw[shift={(3*1.57,0)},color=black] (0pt,2pt) -- (0pt,-2pt) node[below] {\footnotesize $3\frac{\pi}{2}$};

	\draw[shift={(5*1.57,0)},color=black] (0pt,2pt) -- (0pt,-2pt) node[below] {\footnotesize $5\frac{\pi}{2}$};
		\draw[shift={(7*1.57,0)},color=black] (0pt,2pt) -- (0pt,-2pt) node[below] {\footnotesize $7\frac{\pi}{2}$};
			\draw[shift={(9*1.57,0)},color=black] (0pt,2pt) -- (0pt,-2pt) node[below] {\footnotesize $9\frac{\pi}{2}$};
				\draw[shift={(11*1.57,0)},color=black] (0pt,2pt) -- (0pt,-2pt) node[below] {\footnotesize $11\frac{\pi}{2}$};
	%	\draw[shift={(-2*pi,0)},color=black] (0pt,2pt) -- (0pt,-2pt) node[below] {\footnotesize $\mathrm{2\pi}$};
	
	%	\draw[shift={(-4*pi,0)},color=black] (0pt,2pt) -- (0pt,-2pt) node[below] {\footnotesize $\mathrm{-4\pi}$};
	%	\draw[shift={(4*pi,0)},color=black] (0pt,2pt) -- (0pt,-2pt) node[below] {\footnotesize $\mathrm{4\pi}$};
	\draw[thick,color=blue] plot[samples=200, domain=0:19] 
	(\x,{5*(sin(\x r)) });

	
	\end{tikzpicture}\\	
\end{center}
Es gilt für k gerade:\\
\begin{align*}
&k=2=\sin(2\genfrac{}{}{}{2}{\pi }{2 })=0\\
&k=4=\sin(4\genfrac{}{}{}{2}{\pi }{2 })=0\\
&k=6=\sin(6\genfrac{}{}{}{2}{\pi }{2 })=0\\
&\colon
\end{align*}
Es gilt für k ungerade:\\
\begin{align*}
&k=1\Rightarrow\sin(1\genfrac{}{}{}{2}{\pi }{2 })=1\\
&k=3\Rightarrow\sin(3\genfrac{}{}{}{2}{\pi }{2 })=-1\\
&k=5\Rightarrow\sin(5\genfrac{}{}{}{2}{\pi }{2 })=1\\
&k=7\Rightarrow\sin(7\genfrac{}{}{}{2}{\pi }{2 })=-1\\
&\colon
\end{align*}
Ebenso gilt :\\
\begin{align*}
&\sum\limits_{n=1}^{\infty}\big(-1)^{n-1}=\sum\limits_{k=1}^{\infty}\sin(k\genfrac{}{}{}{2}{\pi }{2 })\;\;\text{für $ k $ ungerade } \\
&n=1\Rightarrow 1\\
&n=2\Rightarrow -1\\
&n=3\Rightarrow 1\\
&\colon
\end{align*}

\newpage
Also gilt dann für: $ \frac{\pi}{2}=\pi-\sum\limits_{k}^{\infty}\frac{2}{k}\cdot \sin(k\genfrac{}{}{}{2}{\pi }{2 }) $:\\
\begin{align*}
&\frac{\pi}{2}=\pi-\sum\limits_{n=1}^{\infty}\bigg(\frac{2}{(2n-1)}\big(-1\big)^{n-1}\bigg)\\
&\Leftrightarrow\\
&\sum\limits_{n=1}^{\infty}\bigg(\frac{2}{(2n-1)}\big(-1\big)^{n-1}\bigg)=\frac{\pi}{2}\\
&\Leftrightarrow\\
&\sum\limits_{n=1}^{\infty}\bigg(\frac{1}{(2n-1)}\big(-1\big)^{n-1}\bigg)=\frac{\pi}{4}\\
\end{align*}
Diese Reihe ist die Leibnitz-Reihe.\\
\href{https://youtu.be/81emk4QNvwc}{Video 09}\\
\underline{II. Konvergenzverhalten mit Pasevalscher Gleichung:}\\
\\
Pasevalsche Gleichung:\\
\begin{align*}
\frac{2}{T}\int\limits_{0}^{\mathrm{T}} f^2(x)\mathrm{d}x&=\bigg(\frac{a_0^2}{2}+\sum\limits_{k=1}^{\infty}(a^2_k+b^2_k)\bigg)\\
\end{align*}
Nun gilt für $ T=2\pi, f(x)=x, a_0=2\pi, a_k=0, b_k=-\frac{2}{k} $:\\
\begin{align*}
\frac{2}{2\pi}\int\limits_{0}^{2\pi} x^2\mathrm{d}x=\frac{4\pi^2}{2}+\sum\limits_{k=1}^{\infty}\frac{4}{k^2}
\end{align*}
Integral von $\frac{1}{\pi}\int\limits_{0}^{2\pi} x^2\mathrm{d}x  $ berechnen:\\
\begin{align*}
\frac{1}{\pi}\int\limits_{0}^{2\pi} x^2\mathrm{d}x&=\frac{1}{\pi}\bigg[\frac{x^3}{3}\bigg]^{2\pi}_0=\frac{1}{\pi}\cdot\frac{8\pi^3}{3}=\frac{8\pi^2}{3}
\end{align*}
Dann erhält man:\\
\begin{align*}
&\frac{8\pi^2}{3}=2\pi^2+\sum\limits_{k=1}^{\infty}\frac{4}{k^2}\\
&\Leftrightarrow \frac{8\pi^2}{3}-2\pi^2=\sum\limits_{k=1}^{\infty}\frac{4}{k^2}\\
&\,\;\;\;\;\;\;\;\;\;\;\;\;\;\;\;\frac{2}{3}\pi^2=\sum\limits_{k=1}^{\infty}\frac{4}{k^2}\\
&\Leftrightarrow\\
&\,\;\;\;\;\;\;\;\;\;\;\;\;\;\;\;\frac{1}{3}\pi^2=\sum\limits_{k=1}^{\infty}\frac{2}{k^2}
\end{align*}
\subsubsection{Aufgabe zur Fourierapproximation }
\href{https://youtu.be/USTvFINpj6I}{Video 10}\\
Gegeben ist die $\pi$-periodische Funktion $ f(x)=\big(x-\frac{\pi}{2}\big)^2 $ für $ 0\leq x \leq\pi $.\\
(i) \;\;Skizziere den Graph von $ f $.\\
(ii)\; Berechne das n-te Fourierpolynom von $ f $.\\
(iii) Zeige:\\
\begin{align*}
\sum\limits_{n=1}^{\infty}\frac{(-1)^{n-1}}{n^2}=\frac{\pi^2}{12}
\end{align*}
zu (i): Skizze von der $\pi$ -periodischen Funktion: $ f(x)=\big(x-\frac{\pi}{2}\big)^2$ für $ 0\leq x \leq\pi $\\
	\begin{center}
	
	\begin{tikzpicture}[xscale=1, yscale=1]
	% Gitternetzlinien
	\draw[xstep=1cm, ystep=1cm, very thin, color=lightgray] (-4,4) grid (8,-1);
	% x-Achse
	\draw[->] (-4,0) -- (8,0) node[below,xshift=-1em] {$x$}; 
	%Zahlen auf x-Achse
	%\foreach \x in {-3,...,-1}
	%\draw[shift={(\x,0)},color=black] (0pt,2pt) -- (0pt,-2pt) node[below] {\footnotesize $\x$};
	%\foreach \x in {1,...,3}
	%\draw[shift={(\x,0)},color=black] (0pt,2pt) -- (0pt,-2pt) node[below] {\footnotesize $\x$};
	% y-Achse 
	\draw[->] (0,-1) -- (0,4) node[left,yshift=-1em] {$f(x)$};%node[above left]
	%Zahlen auf y-Achse
	%\foreach \y in {-1,...,-1}
	%\draw[shift={(0,\y)},color=black] (2pt,0pt) -- (-2pt,0pt) node[left] {\footnotesize $\y$};
	%\foreach \y in {1,...,3}
	%\draw[shift={(0,\y)},color=black] (2pt,0pt) -- (-2pt,0pt) node[left] {\footnotesize $\y$};
	%Ursprung
	%\draw[shift={(0,5)},color=black] (2pt,0pt) -- (-2pt,0pt) node[left] {\footnotesize $1$};
	\draw[shift={(0,-5)},color=black] (2pt,0pt) -- (-2pt,0pt) node[left] {\footnotesize $-1$};
	\draw[color=black] (0pt,-10pt) node[right] {\footnotesize $0$};
	\draw[shift={(1.57,0)},color=black] (0pt,2pt) -- (0pt,-2pt) node[below] {\footnotesize $\frac{\pi}{2}$};
	\draw[shift={(3.14,0)},color=black] (0pt,2pt) -- (0pt,-2pt) node[below] {\footnotesize $\pi$};
	
	%\draw[shift={(5*1.57,0)},color=black] (0pt,2pt) -- (0pt,-2pt) node[below] {\footnotesize $5\frac{\pi}{2}$};
	\draw[shift={(6.28,0)},color=black] (0pt,2pt) -- (0pt,-2pt) node[below] {\footnotesize $2\pi$};
	\draw[shift={(-3.14,0)},color=black] (0pt,2pt) -- (0pt,-2pt) node[below] {\footnotesize $-\pi$};
	\draw[shift={(4.71,0)},color=black] (0pt,2pt) -- (0pt,-2pt) node[below] {\footnotesize $\frac{3}{2}\pi$};
		\draw[shift={(-1.57,0)},color=black] (0pt,2pt) -- (0pt,-2pt) node[below] {\footnotesize $-\frac{\pi}{2}$};
	%\draw[shift={(11*1.57,0)},color=black] (0pt,2pt) -- (0pt,-2pt) node[below] {\footnotesize $11\frac{\pi}{2}$};
	%	\draw[shift={(-2*pi,0)},color=black] (0pt,2pt) -- (0pt,-2pt) node[below] {\footnotesize $\mathrm{2\pi}$};
	
	%	\draw[shift={(-4*pi,0)},color=black] (0pt,2pt) -- (0pt,-2pt) node[below] {\footnotesize $\mathrm{-4\pi}$};
	%	\draw[shift={(4*pi,0)},color=black] (0pt,2pt) -- (0pt,-2pt) node[below] {\footnotesize $\mathrm{4\pi}$};
	\draw[thick,color=blue] plot[samples=200, domain=0:3.13] 
	(\x,{(\x-1.57))^2 }); 
		\draw[thick,color=blue] plot[samples=200, domain=3.15:6.27] 
	(\x,{(\x-4.71))^2 }); 
		\draw[thick,color=blue] plot[samples=200, domain=-3.09:0] 
	(\x,{(\x+1.57))^2 }); 
	\end{tikzpicture}\\	
\end{center}
zu (ii): Berechnung von $ \phi_n(x)=\frac{a_0}{2}+\sum\limits_{k=1}^{n}\big(a_k\cos(k\omega x)+b_k\sin(k\omega x)\big) $\\
\begin{quote}
	Berechnen von $ \omega $:\\
	\begin{align*}
	&\mathrm{T}=\pi\\
	&\mathrm{T}=\frac{2\pi}{\omega}\overset{!}{=}\pi\Rightarrow \omega=2
	\end{align*}
	\newpage
	$\bullet$ \;\; Bestimmung von $a_0$:\\
	\begin{align*}
	a_0&=\frac{2}{\mathrm{T}}\int\limits_{0}^{\mathrm{T}} f(x)\mathrm{d}x\\
	&=\frac{2}{\pi}\int\limits_{0}^{\pi} \bigg(x-\frac{\pi}{2}\bigg)^2\mathrm{d}x\\
	&=\frac{2}{\pi}\Bigg[\frac{1}{3}\bigg(x-\frac{\pi}{2}\bigg)^3\Bigg]^\pi_0=\frac{2}{3\pi}\Bigg[\bigg(x-\frac{\pi}{2}\bigg)^3\Bigg]^\pi_0\\
	&=\frac{2}{3\pi}\Bigg[\bigg(0-\frac{\pi}{2}\bigg)^3-\bigg(\pi-\frac{\pi}{2}\bigg)^3\Bigg]\\
		&=\frac{2}{3\pi}\Bigg[\bigg(\frac{\pi}{2}\bigg)^3+\bigg(\frac{\pi}{2}\bigg)^3\Bigg]=\frac{2}{3\pi}\bigg(2\bigg(\frac{\pi}{2}\bigg)^3\bigg)=\frac{\pi^2}{3\cdot 2}\\
		a_0&=\frac{\pi^2}{6}
	\end{align*}
\href{https://youtu.be/BAyYyKjnVSE}{Video 11}\\
$\bullet$ \;\; Bestimmung von $a_k$:\\
\begin{align*}
a_k=\frac{2}{\mathrm{T}}\int\limits_{0}^{\mathrm{T}} f(x)\cos(k\omega x)\mathrm{d}x=\frac{2}{\pi}\int\limits_{0}^{\pi}\bigg(x-\frac{\pi}{2}\bigg)^2\cos(k2x)\mathrm{d}x \\
\end{align*}
NR.: partielle Integration von: $\int\limits_{}^{}\big(x-\frac{\pi}{2}\big)^2\cos(k2x)\mathrm{d}x  $\\
\begin{align*}
&u(x)=\bigg(x-\frac{\pi}{2}\bigg)^2\Rightarrow u'(x)=2\bigg(x-\frac{\pi}{2}\bigg)\\
&v'(x)=\cos(k2x)\;\;\Rightarrow v(x)=\frac{1}{2k}\sin(k2x)
\end{align*}
\begin{align*}
\int\limits_{}^{}\bigg(x-\frac{\pi}{2}\bigg)^2\cos(k2x)\mathrm{d}x& =\bigg(x-\frac{\pi}{2}\bigg)^2\frac{1}{2k}\sin(k2x)-\int\limits_{}^{} 2\bigg(x-\frac{\pi}{2}\bigg)\frac{1}{2k}\sin(k2x)\mathrm{d}x\\
&=\bigg(x-\frac{\pi}{2}\bigg)^2\frac{1}{2k}\sin(k2x)-\frac{1}{k}\int\limits_{}^{} \bigg(x-\frac{\pi}{2}\bigg)\sin(k2x)\mathrm{d}x
\end{align*}
Für das bestimmte Integral von $ \int\limits_{0}^{\pi}  $ gilt:\\
\begin{align*}
\int\limits_{0}^{\pi}\bigg(x-\frac{\pi}{2}\bigg)^2\cos(k2x)\mathrm{d}x&=\bigg[\bigg(x-\frac{\pi}{2}\bigg)^2\frac{1}{2k}\sin(k2x)\bigg]^\pi_0-\int\limits_{0}^{\pi} \bigg(x-\frac{\pi}{2}\bigg)\sin(k2x)\mathrm{d}x \\
&=0-\int\limits_{0}^{\pi} \bigg(x-\frac{\pi}{2}\bigg)\sin(k2x)\mathrm{d}x \\
\end{align*}
Nochmalige partielle Integration von $\int\limits_{}^{} \big(x-\frac{\pi}{2}\big)\sin(k2x)\mathrm{d}x  $\\
\begin{align*}
&u(x)=x-\frac{\pi}{2}\;\;\;\;\Rightarrow u'(x)=1\\
&v'(x)=\sin(2kx)\Rightarrow v(x)=-\frac{1}{2k}\cos(2kx)
\end{align*}
\begin{align*}
\int\limits_{}^{} \bigg(x-\frac{\pi}{2}\bigg)\sin(k2x)\mathrm{d}x&=\bigg(x-\frac{\pi}{2}\bigg)\bigg(-\frac{1}{2k}\cos(2kx)\bigg)-\int\limits_{}^{}-\frac{1}{2k}\cos(2kx)\mathrm{d} x\\
&=\bigg(x-\frac{\pi}{2}\bigg)\bigg(-\frac{1}{2k}\cos(2kx)\bigg)+\frac{1}{2k}\int\limits_{}^{}\cos(2kx)\mathrm{d} x\\
\end{align*}
Für das bestimmte Integral von $ \int\limits_{o}^{\pi}  $ gilt, weil das Integral: $ \int\limits_{0}^{\pi}\cos(2kx) =0 $:\\
\begin{align*}
\int\limits_{0}^{\pi} \bigg(x-\frac{\pi}{2}\bigg)\sin(k2x)\mathrm{d}x=\bigg(x-\frac{\pi}{2}\bigg)\bigg(-\frac{1}{2k}\cos(2kx)\bigg)+0
\end{align*}
Nun ergibt sich das Integral von:$ \int\limits_{}^{}\big(x-\frac{\pi}{2}\big)^2\cos(k2x)\mathrm{d}x $:\\
\begin{align*}
\int\limits_{}^{}\bigg(x-\frac{\pi}{2}\bigg)^2\cos(k2x)\mathrm{d}x=-\frac{1}{k} \bigg(x-\frac{\pi}{2}\bigg)\bigg(-\frac{1}{2k}\cos(2kx)\bigg)=\frac{1}{k} \bigg(x-\frac{\pi}{2}\bigg)\bigg(\frac{1}{2k}\cos(2kx)\bigg)
\end{align*}
Nun ergibt sich für $ a_k $:\\
\begin{align*}
a_k&=\frac{2}{\pi}\int\limits_{0}^{\pi}\frac{1}{k} \bigg(x-\frac{\pi}{2}\bigg)^2\cos(k2x)\mathrm{d}x\\
&=\frac{2}{k\pi}\bigg[\bigg(x-\frac{\pi}{2}\bigg)\bigg(\frac{1}{2k}\cos(2kx)\bigg)\bigg]^\pi_0\\
&=\frac{1}{\pi k^2}\bigg[\bigg(x-\frac{\pi}{2}\bigg)\cos(2kx)\bigg]^\pi_0\\&=\frac{1}{\pi k^2}\bigg[\bigg(\pi-\frac{\pi}{2}\bigg)\cos(2k\pi)-\bigg(0-\frac{\pi}{2}\bigg)\cos(2k 0)\bigg]\\
&=\frac{1}{\pi k^2}\bigg(\frac{\pi}{2}+\frac{\pi}{2}\bigg)=\frac{1}{\pi k^2}\pi\\
a_k&=\frac{1}{k^2}
\end{align*}
\href{https://youtu.be/3sByNrjJVc0}{Video 12}\\
$\bullet$ \;\; Bestimmung von $b_k$:\\
\begin{align*}
b_k=\frac{1}{\mathrm{T}}\int\limits_{0}^{\mathrm{T}}f(x) \sin(k\omega x)\mathrm{d}x=\frac{2}{\pi}\int\limits_{0}^{\pi} \bigg(x-\frac{\pi}{2}\bigg)^2\sin(k2x)\mathrm{d}x
\end{align*}
NR.: partielle Integration von: $\int\limits_{}^{} \big(x-\frac{\pi}{2}\big)^2\sin(k2x)\mathrm{d}x  $\\
\begin{align*}
&u(x)=\bigg(x-\frac{\pi}{2}\bigg)^2\Rightarrow u'(x)=2\bigg(x-\frac{\pi}{2}\bigg)\\
&v'(x)=\sin(2kx)\Rightarrow v(x)\;=-\frac{1}{2k}\cos(2kx)
\end{align*}
\begin{align*}
\int\limits_{}^{} \bigg(x-\frac{\pi}{2}\bigg)^2\sin(k2x)\mathrm{d}x&=\bigg(x-\frac{\pi}{2}\bigg)^2\bigg(-\frac{1}{2k}\cos(2kx)\bigg)-\int\limits_{}^{} 2\bigg(x-\frac{\pi}{2}\bigg)\bigg(-\frac{1}{2k}\cos(2kx)\bigg)\mathrm{d}x\\
&=-\frac{1}{2k}\bigg(x-\frac{\pi}{2}\bigg)^2\bigg(\cos(2kx)\bigg)+\frac{1}{k}\int\limits_{}^{} \bigg(x-\frac{\pi}{2}\bigg)\bigg(\cos(2kx)\bigg)\mathrm{d}x
\end{align*}
Für das bestimmte Integral von $ \int\limits_{0}^{\pi}  $ gilt:\\
\begin{align*}
\int\limits_{0}^{\pi}\bigg(x-\frac{\pi}{2}\bigg)^2\sin(k2x)\mathrm{d}x&=-\frac{1}{2k}\bigg[\bigg(x-\frac{\pi}{2}\bigg)^2\bigg(\cos(2kx)\bigg)\bigg]^\pi_0+\frac{1}{k}\int\limits_{0}^{\pi} \bigg(x-\frac{\pi}{2}\bigg)\bigg(\cos(2kx)\bigg)\mathrm{d}x\\
&=\frac{1}{k}\bigg(\bigg(\frac{\pi}{2}\bigg)^2-\bigg(\frac{\pi}{2}\bigg)^2\bigg)+\frac{1}{k}\int\limits_{0}^{\pi} \bigg(x-\frac{\pi}{2}\bigg)\bigg(\cos(2kx)\bigg)\mathrm{d}x\\
&=0+\frac{1}{k}\int\limits_{0}^{\pi} \bigg(x-\frac{\pi}{2}\bigg)\bigg(\cos(2kx)\bigg)\mathrm{d}x\\
\end{align*}
Nochmalige partielle Integration von: $ \int\limits_{}^{} \big(x-\frac{\pi}{2}\big)\big(\cos(2kx)\big)\mathrm{d}x $
\begin{align*}
&u(x)=x-\frac{\pi}{2}\;\;\;\;\Rightarrow u'(x)=1\\
&v'(x)=\cos(2kx)\Rightarrow v(x)=\frac{1}{2k}\sin(2kx)
\end{align*}
\begin{align*}
\int\limits_{}^{} \bigg(x-\frac{\pi}{2}\bigg)\bigg(-\frac{1}{2k}\cos(2kx)\bigg)\mathrm{d}x=\bigg(x-\frac{\pi}{2}\bigg)\bigg(\frac{1}{2k}\sin(2kx)\bigg)-\int\limits_{}^{} \frac{1}{2k}\sin(2kx)\mathrm{d}x\\
\end{align*}
Für das bestimmte Integral: $ \int\limits_{0}^{\pi}  $ gilt:\\
\begin{align*}
\int\limits_{0}^{\pi} \bigg(x-\frac{\pi}{2}\bigg)\bigg(-\frac{1}{2k}\cos(2kx)\bigg)\mathrm{d}x
&=\frac{1}{2k}\bigg[\bigg(x-\frac{\pi}{2}\bigg)\bigg(\sin(2kx)\bigg)\bigg]^\pi_0-\frac{1}{2k}\int\limits_{0}^{\pi} \sin(2kx)\mathrm{d}x\\
&=\bigg(\pi-\frac{\pi}{2}\bigg) \sin(2k\pi)-\bigg(0-\frac{\pi}{2}\bigg)\sin(2k 0)\bigg)-\frac{1}{2k}\cdot 0\\
&=0-0=0\\
&\Rightarrow \int\limits_{0}^{\pi} \bigg(x-\frac{\pi}{2}\bigg)\sin(k2x)\mathrm{d}x=0
\end{align*}
Nun ergibt sich für $ b_k $:\\
\begin{align*}
b_k=\frac{1}{\mathrm{T}}\int\limits_{0}^{\mathrm{T}}f(x) \sin(k\omega x)\mathrm{d}x&=\frac{2}{\pi}\int\limits_{0}^{\pi} \bigg(x-\frac{\pi}{2}\bigg)^2\sin(k2x)\mathrm{d}x\\
&=\frac{2}{\pi}\cdot 0
\end{align*}
\begin{align*}
b_k=\frac{1}{\pi}\cdot 0 =0
\end{align*}
Für das Fourierplynom mit $ \omega =2,\; a_0=\frac{\pi^2}{6}\pi, \;a_k=\frac{1}{k^2} $ und $ b_k=0 $ gilt dann:\\
\begin{align*}
&\phi(x)=\frac{a_0}{2}+\sum\limits_{k=1}^{n}(a_k\cos(k\omega x)+b_k\sin(k\omega x))\\
&\phi(x)=\frac{\pi^2}{12}+\sum\limits_{k=1}^{n}\frac{1}{k^2}\cos(k2x)
\end{align*}
\end{quote}
\href{https://youtu.be/IcA3p-nQn7Y}{Video 13}\\
zu (iii): Beweise:\\
\begin{align*}
\sum\limits_{n=1}^{\infty}\frac{(-1)^{n-1}}{n^2}=\frac{\pi^2}{12}
\end{align*}
\begin{quote}
 Da $ f(x) $ in $ \frac{\pi}{2} $ stetig ist gilt: $ f(\frac{\pi}{2})=\lim\limits_ {n\to\infty}\phi_n(\frac{\pi}{2})=\frac{\pi^2}{12}+\sum\limits_{n=1}^{\infty}\frac{1}{n}\cos(n2\cdot\frac{\pi}{2}) $.\\
Dann gilt:\\
\begin{align*}
 f(\genfrac{}{}{}{2}{\pi }{2 })=0=\lim\limits_ {n\to\infty}\bigg(\frac{\pi^2}{12}-\sum\limits_{n=1}^{\infty}\frac{1}{n^2}\cos(n\pi)\bigg)
\end{align*}
Für den $ \cos(n\pi) $ gilt:\\
\begin{align*}
&n=1\Rightarrow \cos(1\pi)=1\\
&n=2\Rightarrow \cos(2\pi)=-1\\
&n=3\Rightarrow \cos(3\pi)=1\\
&\colon\\
&\Rightarrow (-1)^{n-1}=\cos(n\pi)\;\;\text{für } n\geq 1
\end{align*}
Dann erhält man:\\
\begin{align*}
&f(\genfrac{}{}{}{2}{\pi }{2 })=0=\lim\limits_ {n\to\infty}\bigg(\frac{\pi}{12}-\sum\limits_{n=1}^{\infty}\frac{1}{n^2}(-1)^{n-1}\bigg)\\
&0=\frac{\pi}{12}-\sum\limits_{n=1}^{\infty}\frac{1}{n^2}(-1)^{n-1}\\
&\Leftrightarrow\\
&\sum\limits_{n=1}^{\infty}\frac{(-1)^{n-1}}{n^2}=\frac{\pi}{12}_{ _\Box}
\end{align*}
\end{quote}
\subsection{Komplexe Fourieranalysis/Komplexes Fourierpolynom}
\href{https://youtu.be/dZAkPwv7v28}{Video 14}\\
Sei $ f:\mathbb{R}\longmapsto \mathbb{R} $ eine $ \mathrm{T}=\frac{2\pi}{\omega} $ periodische monotone Funktion.\\
Wir wissen von  den kompexen Zahlen:\\
\begin{align*}
\mathrm{e}^{\mathrm{i}x}=\cos(x)+\mathrm{i}\sin(x)
\end{align*}
Dann gilt auch:\\
\begin{align*}
\cos(x)=\frac{\mathrm{e}^{\mathrm{ix}}+\mathrm{e}^{-\mathrm{i}x}}{2}\\
\sin(x)=\frac{\mathrm{e}^{\mathrm{ix}}-\mathrm{e}^{-\mathrm{i}x}}{2\mathrm{i}}\\
\end{align*}
Damit bekommt man durch Umformung von:\\
\begin{align*}
\phi_n(x)=\frac{a_0}{2}+\sum\limits_{k=1}^{n}\big(a_k\cos(k\omega x)+b_k\sin(k\omega x)\big)
\end{align*}
das komplexe Fourierpolynom:\\
\begin{align*}
\phi_n(x)&=\sum\limits_{k=-n}^{n}c_k\cdot \mathrm{e}^{\mathrm{i}kwx}\\
c_k&=\frac{1}{\mathrm{T}}\int\limits_{0}^{\mathrm{T}} f(x)\mathrm{e}^{-\mathrm{i}k\omega x}\mathrm{d}x
\end{align*}
\subsubsection{Beispiel zum komplexen Fourierpolynom Teil 1}
\href{https://youtu.be/PLrUWD6DO3g}{Video 15}\\
Betrachte die $ 2\pi $ -periodische Funktion $ f(x)=x $ für $ 0\leq x\leq\pi $ .\\
\\
Da $ \mathrm{T}=\frac{2\pi}{\omega} $ und $ \mathrm{T}=2\pi \Rightarrow \omega=1 $.\\
Zu berechnen:\\
\begin{align*}
\phi_n(x)&=\sum\limits_{k=-n}^{n}c_k\cdot \mathrm{e}^{\mathrm{i}kwx}\\
c_k&=\frac{1}{\mathrm{T}}\int\limits_{0}^{\mathrm{T}} f(x)\mathrm{e}^{-\mathrm{i}k\omega x}\mathrm{d}x
\end{align*}
Mit $f(x)=x,\; \mathrm{T}=2\pi $ und $ \omega=1 $ gilt:\\
\begin{align*}
c_k&=\frac{1}{2\pi}\int\limits_{0}^{\mathrm{2\pi}} x\cdot\mathrm{e}^{-\mathrm{i}k x}\mathrm{d}x
\end{align*}
NR.: partielle Integration von: $ \int\limits_{}^{} x\cdot \mathrm{e}^{-\mathrm{i}kx}\mathrm{d}x $\\
\begin{align*}
&u(x)=x\;\;\;\;\;\;\Rightarrow u'(x)=1\\
&v'(x)=\mathrm{e}^{-\mathrm{i}kx}\Rightarrow v(x)=-\frac{1}{\mathrm{i}k}\mathrm{e}^{-\mathrm{i}kx}
\end{align*}
\begin{align*}
\int\limits_{}^{}x\cdot \mathrm{e}^{-\mathrm{i}kx}\mathrm{d}x&=x\cdot-\frac{1}{\mathrm{i}k} \mathrm{e}^{-\mathrm{i}kx}+\frac{1}{\mathrm{i}k}\int\limits_{}^{}\mathrm{e}^{-\mathrm{i}kx} \mathrm{d}x\\
&=-\frac{1}{\mathrm{i}k}x\cdot \mathrm{e}^{-\mathrm{i}kx}+\frac{1}{\mathrm{i}k}\cdot\frac{1}{-\mathrm{i}k}\mathrm{e}^{-\mathrm{i}kx} \\
\end{align*}
Für das bestimmte Integral: $ \int\limits_{0}^{2\pi}  $ gilt:\\
\begin{align*}
\int\limits_{0}^{2\pi}x\cdot \mathrm{e}^{-\mathrm{i}kx}\mathrm{d}x&=-\frac{1}{\mathrm{i}k}\bigg[x\cdot \mathrm{e}^{-\mathrm{i}kx}\bigg]^{2\pi}_0+\frac{1}{k^2}\bigg[\mathrm{e}^{\mathrm{i}kx}\bigg]^{2\pi}_0\\
&=-\frac{1}{\mathrm{i}k}\big(2\pi\cdot\mathrm{e}^{\mathrm{i}2\pi}-0\cdot\mathrm{e}^{\mathrm{i}k0}\big)+\frac{1}{k^2}\big(\mathrm{e}^{\mathrm{i}k2\pi}-\mathrm{e}^{\mathrm{i}k0}\big)\\
\end{align*}
\begin{quote}
	NR.\\
	\begin{align*}
	\mathrm{e}^{\mathrm{i}k2\pi}=\big(\mathrm{e}^{\mathrm{i}\pi}\big)^{2k}=\big(-1\big)^{2k}=1
	\end{align*}
	oder:\\
	\begin{align*}
	\mathrm{e}^{\mathrm{i}k2\pi}=\cos(k2\pi)+\mathrm{i}\sin(k2\pi)=1-0=1
	\end{align*}
\end{quote}
So ergibt sich das bestimmte Integral:  $ \int\limits_{0}^{2\pi}  $ \\
\begin{align*}
\int\limits_{0}^{2\pi}x\cdot \mathrm{e}^{-\mathrm{i}kx}\mathrm{d}x&=-\frac{1}{\mathrm{i}k}(2\pi\cdot 1-0)+\frac{1}{k^2}(1-1)=-\frac{2\pi}{\mathrm{i}k}
\end{align*}
Dann gilt für $ c_k $:\\
\begin{align*}
c_k&=\frac{1}{2\pi}\int\limits_{0}^{\mathrm{2\pi}} x\cdot\mathrm{e}^{-\mathrm{i}k x}\mathrm{d}x\\
&=\frac{1}{2\pi}\cdot-\frac{2\pi}{\mathrm{i}k}=\frac{1}{-\mathrm{i}k}\\
\end{align*}
\begin{quote}
	Ebenso gilt:\\
	\begin{align*}
	\frac{-1}{\mathrm{i}}=\frac{\mathrm{i}\cdot (-1)}{\mathrm{i}\cdot \mathrm{i}}=\frac{\mathrm{i}\cdot (-1)}{\mathrm{i^2}}=\frac{-\mathrm{i}}{-1}=\mathrm{i}\Rightarrow \frac{-1}{\mathrm{i}k}=\mathrm{i}\frac{1}{k}
	\end{align*}
\end{quote}
Für $ k\neq 0 $ ergibt sich:\\
\begin{align*}
c_k=\mathrm{i}\frac{1}{k} \;\;\text{für } k\neq 0
\end{align*}
Für $ k= 0 $ ergibt sich dann: $ c_0 $\\
\begin{align*}
c_0&=\frac{1}{2\pi}\int\limits_{0}^{2\pi} x\cdot \mathrm{e}^{-\mathrm{i}kx}\mathrm{d}x\\
&=\frac{1}{2\pi}\int\limits_{0}^{2\pi} x\cdot \mathrm{e}^{0}\mathrm{d}x\\
&=\frac{1}{2\pi}\int\limits_{0}^{2\pi} x\mathrm{d}x=\frac{1}{2\pi}\bigg[\frac{x^2}{2}\bigg]^{2\pi}_0=\frac{1}{2\pi}\frac{(2\pi)^2}{2}=\frac{2\pi}{2}=\pi\\
\end{align*}
Fazit:\\
\begin{align*}
c_0&=\pi\\
c_k&=\frac{\mathrm{i}}{k}\;\;\text{für } k\neq 0\\
\phi_n(x)&=\pi+\sum\limits_{\substack{k=-n\\k\neq 0}}^{n}\frac{\mathrm{i}}{k}\mathrm{e}^{\mathrm{i}kx}
\end{align*}
\subsubsection{Beispiel zum komplexen Fourierpolynom Teil 2}
\href{https://youtu.be/Ln1HFYRNBsw}{Video 16}\\
Das komplexe Fourierpolynom in das reelle Fourierpolynom umwandeln.\\
\begin{align*}
\phi_n(x)&=\pi+\sum\limits_{\substack{k=-n\\k\neq 0}}^{n}\frac{\mathrm{i}}{k}\mathrm{e}^{\mathrm{i}kx}\;\;\arrowvert\text{\tiny Summe aufteilen }\\
&=\pi+\sum\limits_{k=-n}^{-1}\frac{\mathrm{i}}{k}\mathrm{e}^{\mathrm{i}kx}+\sum\limits_{k=1}^{n}\frac{\mathrm{i}}{k}\mathrm{e}^{\mathrm{i}kx}\\
&=\pi+\sum\limits_{k=1}^{n}\frac{\mathrm{i}}{-k}\mathrm{e}^{-\mathrm{i}kx}+\sum\limits_{k=1}^{n}\frac{\mathrm{i}}{k}\mathrm{e}^{\mathrm{i}kx}\\
&=\pi+\sum\limits_{k=1}^{n}\frac{\mathrm{i}}{k}\big(\mathrm{e}^{\mathrm{i}kx}-\mathrm{e}^{-\mathrm{i}kx}\big)\;\;\arrowvert\text{\tiny mit 2 erweitern }\\
&=\pi+\sum\limits_{k=1}^{n}\frac{2\mathrm{i}}{k}\bigg(\frac{\mathrm{e}^{\mathrm{i}kx}-\mathrm{e}^{-\mathrm{i}kx}}{2}\bigg)\;\;\arrowvert\text{\tiny $-\mathrm{i}=\frac{1}{\mathrm{i}}$ }\\
&=\pi+\sum\limits_{k=1}^{n}\frac{-2}{\mathrm{i}k}\bigg(\frac{\mathrm{e}^{\mathrm{i}kx}-\mathrm{e}^{-\mathrm{i}kx}}{2}\bigg)\;\;\lvert\text{\tiny mit $ \mathrm{i} $ erweitern}\\
&=\pi+\sum\limits_{k=1}^{n}\frac{-2\mathrm{i}}{\mathrm{i}k}\bigg(\frac{\mathrm{e}^{\mathrm{i}kx}-\mathrm{e}^{-\mathrm{i}kx}}{2\mathrm{i}}\bigg)\\
&=\pi+\sum\limits_{k=1}^{n}\frac{-2}{k}\bigg(\underbrace {\frac{\mathrm{e}^{\mathrm{i}kx}-\mathrm{e}^{-\mathrm{i}kx}}{2\mathrm{i}}}_{=\sin(kx)}\bigg)\\
&=\pi+\sum\limits_{k=1}^{n}\frac{-2}{k}\big(\sin(kx)\big)\\
\phi_n(x)&=\pi-\sum\limits_{k=1}^{n}\frac{2}{k}\sin(kx)\\
\end{align*}
Bei der komplexen Berechnung des Fourierpolynoms muß nur $ c_k $ und vielleicht noch $ c_o $ bestimmt werden. Beim reellen Fourierpolynom muß man $ a_0,\;a_k $ und $ b_k $ bestimmen. Welche Methode besser ist, ist von Fall zu Fall unterschiedlich.\\
%	
\chapter{Reihen}
\setcounter{section}{13}
\setcounter{subsection}{0}
\setcounter{subsubsection}{-1}
\setcounter{secnumdepth}{4}
\subsection{Reihen}
\href{https://youtu.be/-YXCVGfMCcQ}{Video 01}\\
Eine (unendliche) Reihe hat die Form:\\
\begin{align*}
a_0+a_1+\ldots=\sum\limits_{k=0}^{\infty}a_k\;\;\text{mit }\;\;a_k\in\mathbb{R}\;\;\text{ für alle}\;\;k\in\mathbb{N}
\end{align*}
Dazu definieren wie die $ n-\text{te } $ Partialsumme:\\
\begin{align*}
\mathrm{S}_n:=\sum\limits_{k=0}^{n}a_k
\end{align*}
Die unendliche Reihe $ \sum\limits_{k=0}^{\infty}a_k $ heißt konvergent, wenn der Grenzwert von:\\
\begin{align*}
\lim\limits_ {n\to\infty}\mathrm{S}_n=\sum\limits_{k=0}^{\infty}a_k
\end{align*}
existiert.\\
\\
Bevor man eine Reihe bilden kann benötigt man eine Folge $ (a_k)_{\mathbb{N}} $. Nun summiert man die Folgeglieder auf und bildet die Partialsumme.\\
\begin{align*}
&\mathrm{S}_0=a_0=\sum\limits_{k=0}^{0}a_k\\
&\mathrm{S}_1=a_0+a_1=\sum\limits_{k=1}^{1}a_k\\
&\mathrm{S}_2=a_0+a_1+a_2=\sum\limits_{k=1}^{2}a_k\\
&\mathrm{S}_n=a_0+a_1+a_2+\ldots+a_n=\sum\limits_{k=0}^{n}a_k
\end{align*}
Konvergiert nun die Reihe für $ n\rightarrow\infty $, so existiert ein  Grenzwert $ \sum\limits_{k=0}^{\infty} a_k$. Die Reihe konvergiert dann, wenn die Folge konvergiert.\\
Durch $ \lim\limits_ {n\to\infty}\mathrm{S}_n=\sum\limits_{k=0}^{\infty} a_k$ übersetzen wir Begriffe wie Konvergenz, Divergenz und bestimmmte Divergenz.\\
\\
Beispiel:\\
\begin{align*}
\sum\limits_{k=0}^{\infty}a_k\;\;\text{mit }\;\;a_k=k\;\;\text{bzw. }\;\;\sum\limits_{k=0}^{\infty}k
\end{align*}
Es gilt (kann mittels Induktion bewiesen werden - Gaußsche Summenformel ):\\
\begin{align*}
\mathrm{S}_n=\sum\limits_{k=0}^{n}k=\frac{n(n+1)}{2}
\end{align*}
Für den Grenzwert: $ \lim\limits_ {n\to\infty} $ gilt:\\
\begin{align*}
\lim\limits_ {n\to\infty}\mathrm{S}_n=\sum\limits_{k=0}^{\infty}k=\lim\limits_ {k\to\infty}\frac{n(n+1)}{2}=\lim\limits_ {n\to\infty}\frac{n^2+n}{2}=\infty
\end{align*}
Die Partialsumme divergiert gegen das Unendliche.\\
\newpage
\noindent
\subsubsection{Beispiel zum Thema Reihen}
\href{https://youtu.be/hroOOlHucwM}{Video 02}\\
Gegeben die Reihe:\\
\begin{align*}
\sum\limits_{k=2}^{\infty}\bigg(\frac{1}{k-1}-\frac{1}{k}\bigg)
\end{align*}
Bestimme den Grenzwert $\lim\limits_ {n\to\infty} \mathrm{S}_n $:
\begin{align*}
\mathrm{S}_n=\sum\limits_{k=2}^{n}\bigg(\frac{1}{k-1}-\frac{1}{k}\bigg)&=\bigg(1-\frac{1}{2}\bigg)+\bigg(\frac{1}{2}-\frac{1}{3}\bigg)+\bigg(\frac{1}{3}-\frac{1}{4}\bigg)+\ldots+\bigg(\frac{1}{n-1}-\frac{1}{n}\bigg)\\
&=1\underbrace {-\frac{1}{2}+\frac{1}{2}}_{=0}\underbrace {-\frac{1}{3}+\frac{1}{3}}_{=0}\underbrace {-\frac{1}{4}+.}_{=0}\ldots\underbrace {.\bigg(\frac{1}{n-1}}_{=0}-\frac{1}{n}\bigg)\\
&=1-\frac{1}{n}\\
\lim\limits_ {n\to\infty}\mathrm{S}_n&=\lim\limits_ {n\to\infty}\bigg(1-\frac{1}{n}\bigg)=1-0=1\\
\end{align*}
\subsection{Gometrische Reihe}
\href{https://youtu.be/nmtUHgdMWXs}{Video 03}\\
\begin{center}
	\fbox{\parbox  {45mm} {Sei $ q\in\mathbb{R} $. Dann heißt\\$ \sum\limits_{k=0}^{\infty}q^k=1+q+q^2+\ldots $\\die geometrische Reihe}}	
\end{center}
Für die $ n- $te Partialsumme gilt:\\
\begin{align*}
\mathrm{S}_n =   \left\{ \begin{array}{rcl}
n+1 & \mbox{für} & q=1\\
\frac{1-q^2}{1-q}&\mbox{für}& q\neq1
\end{array}\right.
\end{align*}
Untersuchung der Konvergenz der geometrischen Reihe in Abängigkeit von $ q $.\\
\underline{Für $ q=1 $ gilt:}\\
\begin{quote}
	\begin{align*}
	\lim\limits_ {n\to\infty}\mathrm{S}_n=\lim\limits_ {n\to\infty}(n+1)=\infty
	\end{align*}
\end{quote}
\newpage
\noindent
\underline{Für $ q\neq 1 $ gilt:}\\
\begin{quote}
	1. Fall: $ q>1 $\\
	\begin{quote}
		\begin{align*}
		&\lim\limits_ {n\to\infty}\mathrm{S}_n=\lim\limits_ {n\to\infty}\bigg(\frac{1-q^{n+1}}{1-q}\bigg)\\
		\end{align*}
		
		\begin{align*}
		&\text{Es gilt für den Zähler: }\\
		&\lim\limits_ {n\to\infty}1-q^{n+1}=-\infty\\
		&\text{Es gilt für den Nenner: }\\
		&1-q<0\\
		&\Rightarrow \text{Der Bruch geht gegen}+\infty
		\end{align*}
		
	\end{quote}
	2. Fall: $ -1<q<1 $\\
	\begin{quote}
		\begin{align*}
		&\lim\limits_ {n\to\infty}\mathrm{S}_n=\lim\limits_ {n\to\infty}\bigg(\frac{1-q^{n+1}}{1-q}\bigg)\\
		\end{align*}
		\begin{align*}
		&\text{Es gilt für den Zähler: }\\
		&\lim\limits_ {n\to\infty}1-q^{n+1}=1\\
		&\text{Es gilt für den Nenner: }\\
		&1-q>0\\
		&\Rightarrow \lim\limits_ {n\to\infty}\biggl(\frac{1-q^{n+1}}{1-q}\biggr)=\frac{1}{1-q}
		\end{align*}
	\end{quote}
	3. Fall: $ q\leq -1 $\\
	\begin{quote}
		\begin{align*}
		&\lim\limits_ {n\to\infty}\mathrm{S}_n=\lim\limits_ {n\to\infty}\bigg(\frac{1-q^{n+1}}{1-q}\bigg)\\
		\end{align*}
		Es existiert kein Grenzwert, weil $ q^n $ für $ q\leq -1 $ alterniert.\\
		Diese Folge ist nicht konvergent.\\
	\end{quote}
\end{quote}
Fazit:\\
Für $ -1<q<1 $ oder $ |q| <1$ ist die geometrische Reihe konvergent.\\
\begin{center}
	\boxed{ \sum\limits_{k=0}^{\infty}q^k=\lim\limits_ {n\to\infty}\mathrm{S}_n=\frac{1}{1-q} }	
\end{center}
\subsubsection{Aufgabe zur geometrischen Reihe}
\href{https://youtu.be/FlwGJnrlD4c}{Video 04}\\
Zeige, dass die Summe: $ \sum\limits_{k=2}^{\infty}\frac{2^{k+1}}{5\cdot 3^k} $  konvergent ist und berechne den Grenzwert.\\
\begin{align*}
\sum\limits_{k=2}^{\infty}\underbrace {\frac{2^{k+1}}{5\cdot 3^k}}_{=a_k}
\end{align*}
\begin{align*}
a_k=\frac{2^{k+1}}{5\cdot 3^k}=\frac{2}{5}\cdot \frac{2^k}{3^k}=\frac{2}{5}\cdot \bigg(\frac{2}{3}\bigg)^k
\end{align*}
Dann gilt:\\
\begin{align*}
\sum\limits_{k=2}^{\infty}\frac{2^{k+1}}{5\cdot 3^k}=\sum\limits_{k=2}^{\infty}\frac{2}{5}\cdot \bigg(\frac{2}{3}\bigg)^k=\frac{2}{5}\sum\limits_{k=2}^{\infty}\bigg(\frac{2}{3}\bigg)^k
\end{align*}
Damit man die geometrische Reihe anwenden kann muss die Summe mit $ k=0 $ beginnen.\\
\begin{align*}
\frac{2}{5}\sum\limits_{k=2}^{\infty}\bigg(\frac{2}{3}\bigg)^k=\frac{2}{5}\sum\limits_{k=0}^{\infty}\bigg(\frac{2}{3}\bigg)^k-\frac{2}{5}\bigg(\frac{2}{3}\bigg)^0-\frac{2}{5}\bigg(\frac{2}{3}\bigg)^1=\frac{2}{5}\sum\limits_{k=0}^{\infty}\bigg(\frac{2}{3}\bigg)^k-\frac{2}{3}
\end{align*}
Mit der geometrischen Reihe gilt für $ |q|<1:\;\;\sum\limits_{k=0}^{\infty}q^k=\frac{1}{1-q} $:\\
\begin{align*}
\sum\limits_{k=2}^{\infty}\frac{2^{k+1}}{5\cdot 3^k}&=\frac{2}{5}\underbrace {\sum\limits_{k=0}^{\infty}\bigg(\frac{2}{3}\bigg)^k}_{=\frac{1}{1-q}}-\frac{2}{3}\\
&=\frac{2}{5}\bigg(\frac{1}{1-\frac{2}{3}}\bigg)-\frac{2}{3}\\
&=\frac{6}{5}-\frac{2}{3}=\frac{8}{15}
\end{align*}
Fazit:
\begin{align*}
\sum\limits_{k=2}^{\infty}\frac{2^{k+1}}{5\cdot 3^k}=\frac{8}{15} 
\end{align*}
\subsection{Rechenregeln für konvergente Reihen}
\href{https://youtu.be/Hs5s1w53hmI}{Video 05}\\
Seien $ \sum\limits_{k=0}^{\infty} a_k=\mathrm{A}$ und $ \sum\limits_{j=0}^{\infty}b_j=\mathrm{B} $ konvergierende Reihen.\\
Dann gilt:\\
\begin{align*}
\begin{array}{lll}
\text{(i) }&\sum\limits_{k=0}^{\infty}(a_k+b_k)=\sum\limits_{k=0}^{\infty}a_k+\sum\limits_{j=0}^{\infty}b_j=\mathrm{A}+\mathrm{B}\\[15pt]%Zeilenvorschub
\text{(ii) }&\sum\limits_{k=0}^{\infty}(\mathrm{C}\cdot a_k)=\mathrm{C}\sum\limits_{k=0}^{\infty}a_k=\mathrm{C}\cdot\mathrm{A}\;\;\text{für }\mathrm{C}\in\mathbb{R}\;\;\text{beliebig }\\[15pt]
\text{(iii) }&\sum\limits_{k=0}^{\infty}a_k\cdot\sum\limits_{k=0}^{\infty}b_k=\sum\limits_{k=0}^{\infty}c_k\;\;\text{, wobei } c_k=a_0\cdot b_k+a_1\cdot b_{k-1}+\ldots+a_k\cdot b_0\\
&\text{ gilt nur, wenn}\sum\limits_{k=0}^{\infty}c_k \;\;\text{konvergent ist. }
\end{array}
\end{align*}
Hinweis-Cauchy-Produkt: \\
Für $ n\in \mathbb{N} $ setzt man $ c_n:=\sum\limits_{k=0}^{n}a_k\cdot b_{n-k} $, so kann mann alle Summanten bestimmen.\\
\begin{align*}
&c_n=\sum\limits_{k=0}^{n}a_k\cdot b_{n-k}\\
&c_0=a_0b_0\\
&c_1=a_0b_1+a_1b_0\\
&c_2=a_0b_2+a_1b_1+a_2b_0\\
&c_3=a_0b_3+a_1b_2+a_2b_1+a_3b_0\\
&\colon
\end{align*}
\subsubsection{Notwendiges Konvergenzkriterium für Reihen}
\href{https://youtu.be/l_l4p0D7qro}{Video 06}\\
Sei $ \sum\limits_{k=0}^{\infty}a_k $ konvergent gegen $ \mathrm{A} $.\\
Dann gilt:\\
\begin{align*}
\lim\limits_ {n\to\infty}\mathrm{S}_n=\mathrm{A}=\lim\limits_ {n\to\infty}\mathrm{S}_{n-1}
\end{align*}
Bemerkung:\\
Der Grenzwert z.B. von $ \lim\limits_ {n\to\infty}\frac{1}{n}= 0 $ und $ \lim\limits_ {n\to\infty}\frac{1}{n+1}=0 $.\\
\begin{align*}
&\mathrm{S}_n=\sum\limits_{k=0}^{n}a_k\\
&\lim\limits_ {n\to\infty}\mathrm{S}_n=\mathrm{A}=\lim\limits_ {n\to\infty}\mathrm{S}_{n-1}\\
&\Leftrightarrow \lim\limits_ {n\to\infty}(\mathrm{S}_n-\mathrm{S}_{n-1})=\mathrm{A}-\mathrm{A}=0\\
\end{align*}
\begin{align*}
\mathrm{S}_n-\mathrm{S}_{n-1}=\sum\limits_{k=0}^{n}a_k-\sum\limits_{k=0}^{n-1}a_k=a_n
\end{align*}
Bei der Differenz von $ \sum\limits_{k=0}^{n}a_k-\sum\limits_{k=0}^{n-1}a_k $ bleibt nur noch das Glied $ a_n $ über.\\
Also ist der Grenzwert von:\\
\begin{align*}
\lim\limits_ {n\to\infty}(\mathrm{S}_n-\mathrm{S}_{n-1})=\lim\limits_ {n\to\infty}a_n=0
\end{align*}
Fazit:\\
Wenn $ \sum\limits_{k=0}^{\infty}a_k $ ist konvergent $ \Rightarrow \lim\limits_ {n\to\infty}a_n=0$.\\
Wenn $ \lim\limits_ {n\to\infty}a_n\neq 0\Rightarrow\sum\limits_{k=0}^{\infty}a_k $ ist nicht konvergent.\\
\newpage
\noindent
Beispiel:\\
\begin{align*}
\sum\limits_{k=0}^{\infty}\frac{k+1}{k+2}
\end{align*}
Diese Reihe ist nicht konvergent wegen dem notwendigen Konvergenzkriterium:\\
\begin{align*}
\lim\limits_ {k\to\infty}a_k=\lim\limits_ {k\to\infty} \bigg(\frac{k+1}{k+2}\bigg)=\frac{k(1+\frac{1}{k})}{k(1+\frac{2}{k})}=1\neq 0
\end{align*}
Daraus folgt, dass die Reihe:$ \sum\limits_{k=0}^{\infty}\frac{k+1}{k+2}=1 $ divergiert.\\
Achtung:\\
Der Umkehrschluss gilt nicht. Ist die Reihe divergent dann muss die Folge auch nicht divergent sein.\\
Beispiel:\\
$ \sum\limits_{k=1}^{\infty}\frac{1}{k} $ ist divergent, aber der Grenzert von: $ \lim\limits_ {k\to\infty}\frac{1}{k}=0 $.\\
\subsubsection{Majorantenkriterium für Reihen}
\href{https://youtu.be/qShFeav_pYk}{Video 07}\\
Sei $ \sum\limits_{k=0}^{\infty}a_k $ und $ \sum\limits_{k=0}^{\infty}b_k $ gegeben.\\
Existiert ein $ k_o\in \mathbb{N} $, so dass gilt:\\
\begin{align*}
|a_k|\leq b_k\;\;\text{ für alle}\; k\geq k_o
\end{align*}
und ist $ \sum\limits_{k=0}^{\infty}b_k $ konvergent, dann ist auch $ \sum\limits_{k=0}^{\infty}a_k $ konvergent.\\
\\
Beispiel:\\
Ist die Reihe: $ \sum\limits_{k=0}^{\infty}\frac{1}{r^2}=a_k $ konvergent ?\\
Es gilt: Die Reihe: $ \sum\limits_{k=2}^{\infty}\frac{1}{k-1}-\frac{1}{k}=b_k $ ist konvergent.\\
Wir betrachten:\\
\begin{align*}
|a_k|=\frac{1}{k^2}=\frac{1}{k}\cdot\frac{1}{k}&<\frac{1}{(k-1)}\cdot\frac{1}{k}\;\;\text{für }\;k\geq 2\\
&<\frac{k-k+1}{k(k-1)}=\frac{k}{k(k-1)}+\frac{-k+1}{k(k-1)}=\frac{1}{k-1}+\frac{-1(k-1)}{k(k-1)}=\frac{1}{k-1}-\frac{1}{k}=b_k\\
|a_k|&<b_k
\end{align*}
Fazit:\\
Da $ |a_k|=\frac{1}{k^2}<\frac{1}{k-1}-\frac{1}{k}=b_k $ für $ k\geq2 $ und $ \sum\limits_{k=2}^{\infty}\frac{1}{k-1}-\frac{1}{k} $ konvergent ist folgt mit dem Majorantenkriterium:\\
\begin{align*}
\sum\limits_{k=0}^{\infty}\frac{1}{k^2}\;\;\text{ist konvergent. }\\
\end{align*}
\subsubsection{Quotientenkriterium für Reihen}
\href{https://youtu.be/6xdKQ09P7c4}{Video 08}\\
Sei $ \sum\limits_{k=0}^{\infty} a_k $ gegeben. Dann gilt:\\
\begin{align*}
\begin{array}{lll}
\text{(i) }&\sum\limits_{k=0}^{\infty}\big|\frac{a_{k+1}}{a_k}\big|<1\Rightarrow \sum\limits_{k=0}^{\infty}a_k \;\;\text{ist konvergnet }\\[15pt]%Zeilenvorschub
\text{(ii) }&\sum\limits_{k=0}^{\infty}\big|\frac{a_{k+1}}{a_k}\big|>1\Rightarrow \sum\limits_{k=0}^{\infty}a_k \;\;\text{ist divergent }\\[15pt]%Zeilenvorschub
\text{(iii) }&\sum\limits_{k=0}^{\infty}\big|\frac{a_{k+1}}{a_k}\big|=1\;\text{oder nicht existent, dann kann man keine Aussage machen! }\\[15pt]%Zeilenvorschub
\end{array}
\end{align*}
Beispiel: \\
\begin{align*}
\sum\limits_{k=0}^{\infty}\frac{k}{2^k}
\end{align*}
Betrachte:\\
\begin{align*}
\bigg|\frac{a_{k+1}}{a_k}\bigg|
\end{align*}
Denn Betrag kann man weglassen, weil Zähler und Nenner immer  positv werden.\\
\begin{align*}
\frac{a_{k+1}}{a_k}=\frac{\frac{k+1}{2^{k+1}}}{\frac{k}{2^k}}=\frac{(k+1)(2^k)}{(2^{k+1})(k)}
\end{align*}
Für $ k\neq 0 $ gilt dann:\\
\begin{align*}
\frac{a_{k+1}}{a_k}=\frac{(k+1)(2^k)}{(2^{k+1})(k)}=\frac{1}{2}\cdot\frac{k+1}{k}
\end{align*}
Für $ k\mapsto\infty $ gilt:\\
\begin{align*}
\lim\limits_ {k\to\infty}\frac{1}{2}\cdot\frac{k+1}{k}=\frac{1}{2}\cdot\frac{k(1+\frac{1}{k})}{k}=\frac{1}{2}
\end{align*}
Daraus folgt mit dem Quotientenkriterium:\\
\begin{align*}
\lim\limits_ {k\to\infty}\frac{a_{k+1}}{a_k}=\frac{(k+1)(2^k)}{(2^{k+1})(k)}=\frac{1}{2}<1\Rightarrow \;\;\text{Die Reihe ist konvergent. }
\end{align*}
\newpage
\noindent
\paragraph{Aufgabe zum Quotientenkriterium für Reihen}$  $\vspace*{3mm}\\
\href{https://youtu.be/Zp1q1k58Tqk}{Video 09}\\
Ist die Reihe $ \sum\limits_{k=1}^{\infty}\frac{(-1)^k(k+1)!}{2\cdot 4\cdot 6 \cdot\ldots \cdot 2k} $ konvergent?\\
Daraus folgt mit dem Quotientenkriterium mit $ a_k=\frac{(-1)^k(k+1)!}{2\cdot 4\cdot 6 \cdot\ldots \cdot 2k} $ :\\
\begin{align*}
\lim\limits_ {k\to\infty}\bigg|\frac{a_{k+1}}{a_k}\bigg|&=\Bigg|\frac{\frac{(-1)^{k+1}(k+2)!}{2\cdot 4\cdot 6\cdot\ldots\cdot 2(k+1)}}{\frac{(-1)^{k}(k+1)!}{2\cdot 4\cdot 6\cdot\ldots\cdot 2k}}\Bigg|\\
&=\frac{2\cdot 4\cdot 6\cdot\ldots\cdot2k(k+2)!}{2\cdot 4\cdot 6 \cdot\ldots \cdot 2(k+1)(k+1)!}\\
&=\frac{2\cdot 4\cdot 6\cdot\ldots\cdot2k(k+2)(k+1)!}{2\cdot 4\cdot 6 \cdot\ldots \cdot 2(k+1)(k+1)!}\\
&=\frac{2\cdot 4\cdot 6\cdot\ldots\cdot2k(k+2)}{2\cdot 4\cdot 6\cdot\ldots\cdot2(k+1)}\\
&=\frac{2\cdot 4\cdot 6\cdot\ldots\cdot2k(k+2)}{2\cdot 4\cdot 6\cdot\ldots\cdot (2k+2)}\text{\tiny vorletzter Faktor im Nenner ist $ 2k $ und der Letzte $ 2k+2 $}\\
&=\frac{2\cdot 4\cdot 6\cdot\ldots\cdot2k(k+2)}{2\cdot 4\cdot 6\cdot\ldots\cdot 2k\cdot(2k+2)}\\
&=\frac{(k+2)}{(2k+2)}\\
&=\frac{(k+2)}{2(k+1)}=\frac{1}{2}\cdot\frac{k(1+\frac{2}{k})}{k(1+\frac{1}{k})}=\frac{1}{2}<1
\end{align*}
Daraus folgt nach dem Quotientenkriterium, dass die Folge:\\
\begin{align*}
\sum\limits_{k=1}^{\infty}\frac{(-1)^k(k+1)!}{2\cdot 4\cdot 6 \cdot\ldots \cdot 2k}
\end{align*}
konvergent ist.\\
\newpage
\noindent
\subsubsection{Wurzelkriterium für Reihen}
\href{https://youtu.be/axTFmjV9WSA}{Video 10}\\
Sei $ \sum\limits_{k=0}^{\infty} a_k $ gegeben. Dann gilt:\\
\begin{align*}
\begin{array}{lll}
\text{(i) }&\lim\limits_ {k\to\infty}|\sqrt[k]{a_k}|=\lim\limits_ {k\to\infty}|a_k|^{\frac{1}{k}}<1\Rightarrow \sum\limits_{k=0}^{\infty}a_k\;\;\text{ist konvergent }\\[15pt]%Zeilenvorschub
\text{(ii) }&\lim\limits_ {k\to\infty}|\sqrt[k]{a_k}|=\lim\limits_ {k\to\infty}|a_k|^{\frac{1}{k}}\geq 1\Rightarrow \sum\limits_{k=0}^{\infty}a_k\;\;\text{ist für unendlich viele $ k $ divergent}\\[15pt]%Zeilenvorschub
\end{array}
\end{align*}
Beispiel:\\
Ist die Reihe: $ \sum\limits_{k=0}^{\infty}\big(-\frac{1}{2^k}\big)^k $ konvergent ?\\
Daraus folgt mit dem Wurzelkriterium mit $ a_k=\big(-\frac{1}{2^k}\big)^k $ :\\
\begin{align*}
\lim\limits_ {k=0\to\infty}|\sqrt[k]{a_k}|=\bigg|\sqrt[k]{\bigg(-\frac{1}{2^k}\bigg)^k }\bigg|=\bigg|\bigg(\frac{(-1)^k}{2^{k\cdot k}}\bigg)\bigg|^{\frac{1}{k}}=\bigg(\frac{1}{2^{k^2}}\bigg)^{\frac{1}{k}}=\frac{1}{2^k}=0<1
\end{align*}
Daraus folgt nach dem Wurzelkriterium, dass die Folge:\\
\begin{align*}
\sum\limits_{k=0}^{\infty}\bigg(-\frac{1}{2^k}\bigg)^k
\end{align*}
konvergent ist.\\
\paragraph{Aufgabe zum Wurzelkriterium für Reihen}$  $\vspace*{3mm}\\
\href{https://youtu.be/Z95Uw9OgMck}{Video 11}\\
Zeige, dass diese Folge konvergent ist:\\
\begin{align*}
\sum\limits_{k=0}^{\infty}\frac{k}{3^k}
\end{align*}
Def.:\\
\begin{align*}
a_k=\frac{k}{3^k}
\end{align*}
Berechne:\\
\begin{align*}
\sqrt[k]{|a_k|}&=\sqrt[k]{\frac{k}{3^k}}=\frac{\sqrt[k]{k}}{\sqrt[k]{3^k}}=\frac{\sqrt[k]{k}}{3}
\end{align*}
NR.:\\
\begin{align*}
\sqrt[k]{3^k}=\big(3^k\big)^{\frac{1}{k}}=3^{k\cdot\frac{1}{k}}=3^1=1
\end{align*}
Berechne den Grenzwert:\\
\begin{align*}
\lim\limits_ {k\to\infty}\frac{1}{3}\sqrt[k]{k}=\lim\limits_ {k\to\infty}\frac{1}{3}k^{\frac{1}{k}}=\frac{1}{3}\lim\limits_ {k\to\infty}\mathrm{e}^{\ln(k)\cdot\frac{1}{k}}
\end{align*}
Da die Funktion $ \mathrm{e}^x $ stetig ist, kann man den Grenzwert in die \glqq Potenz\grqq  ziehen.\\
\begin{align*}
\lim\limits_ {k\to\infty}\frac{1}{3}\sqrt[k]{k}=\frac{1}{3}\mathrm{e}^{\lim\limits_ {k\to\infty}\ln(k)\cdot\frac{1}{k}}
\end{align*}
Für den Grenzwert von $\lim\limits_ {k\to\infty}\ln(k)\cdot\frac{1}{k} $ gilt:\\
\begin{align*}
\lim\limits_ {k\to\infty}\frac{\ln(k)}{k}=0
\end{align*}
Bemerkung:\\
Der Logarithmus $ \ln(k) $wächst langsamer als der Bruch $ \frac{1}{k} $, deswegen geht der Grenzwert gegen Null.\\
Nun gilt:\\
\begin{align*}
&\lim\limits_ {k\to\infty}\frac{1}{3}\sqrt[k]{k}=\frac{1}{3}\mathrm{e}^{\lim\limits_ {k\to\infty}\ln(k)\frac{1}{k}}=\frac{1}{3}\mathrm{e}^0=\frac{1}{3}<1\\
&\Rightarrow \sum\limits_{k=0}^{\infty}\frac{k}{3^k}\;\;\text{konvergiert}
\end{align*}
\subsubsection{Das Leibnizkriterium}
\href{https://youtu.be/BM329RZgnwg}{Video 12}\\
Eine Reihe $ \sum\limits_{k=0}^{\infty} a_k$ heißt alternierent, wenn gilt: 
\begin{align*}
a_k\cdot a_{k+1}<0
\end{align*}
Die Summanten haben hintereinander verschiedene Vorzeichen.\\
\begin{align*}
\overset{-}{a_0}+\overset{+}{a_0}+\overset{-}{a_0}+\overset{+}{a_0}+\\
\overset{+}{a_0}+\overset{-}{a_0}+\overset{+}{a_0}+\overset{-}{a_0}+\\
\end{align*}
Beispiel:\\
\begin{align*}
\sum\limits_{k=0}^{\infty}(-1)^k=1-1+1-1\pm\ldots=0
\end{align*}
Die Summe wird $ 0 $, wenn man folgende Summanten addiert:\\
\begin{align*}
\sum\limits_{k=0}^{\infty}(-1)^k=\underbrace {1-1}_{=0}+\underbrace {1-1}_{=0}+\underbrace {1-1}_{=0}+\underbrace {1-1}_{=0}\pm\ldots=0
\end{align*}
Die Summe wird $ 1 $, wenn man folgende Summanten addiert:\\
\begin{align*}
\sum\limits_{k=0}^{\infty}(-1)^k=1\underbrace {-1+1}_{=0}\underbrace {-1+1}_{=0}\underbrace {-1+1}_{=0}\underbrace {-1+1}_{=0}\pm\ldots=1
\end{align*}
Die Summe wird $ 20 $, wenn man folgende Summanten addiert:\\
\begin{align*}
\sum\limits_{k=0}^{\infty}(-1)^k=\underbrace {1+\ldots +1}_{20\text{\tiny mal }}\underbrace {+1-1+1-1+1}_{=0}\pm\ldots=20\
\end{align*}
\\
\underline{Leibnizkriterium:}\\
Sei $ \sum\limits_{k=0}^{\infty}a_k $ eine alternierende Reihe\\ und ist $ \big(|a_k|\big)_k $ eine monoton fallende Nullfolge, dann konvergiert $ \sum\limits_{k=0}^{\infty}a_k $.
\\
Bemerkung:\\
Da kein Grenzwert bei der Reihe $\sum\limits_{k=0}^{\infty} (-1)^k  $ existiert und dies ist ja ein notwendiges Kriterium für eine Konvergenz ist divergiert die Reihe $ \sum\limits_{k=0}^{\infty} (-1)^k$.\\
\\
Betrachte die Reihe:\\
\begin{align*}
\sum\limits_{k=1}^{\infty}(-1)^k\frac{1}{k}
\end{align*}
\begin{align*}
a_k=(-1)^k\frac{1}{k}
\end{align*}
Wegen:
\begin{align*}
\big(|a_k|\big)_{k\in \mathbb{N}}=\bigg(\frac{1}{k}\bigg)_{k\in\mathbb{N}\setminus\{0\}}
\end{align*}
Bemerkung:\\
\begin{align*}
\bigg|(-1)^k\cdot\frac{1}{k}\bigg|=\frac{1}{k}
\end{align*}
Da $\big(\frac{1}{k}\big)_{k\in\mathbb{N}\setminus\{0\}}  $ eine streng monoton fallende Nullfolge ist, folgt mit dem Leinbnizkriterium, dass die Folge $ \sum\limits_{k=1}^{\infty}(-1)^k\frac{1}{k} $ konvergent ist.\\
Bemerkung:
Die Reihe $ \sum\limits_{k=1}^{\infty}\frac{1}{k} $ ist divergent.\\
\newpage
\noindent
\paragraph{Aufgabe zum Leibnizkriterium}$  $\vspace*{3mm}\\
\href{https://youtu.be/95O_7qUchvg}{Video 13}\\
Zeige, dass diese Folge konvergent ist.\\
\begin{align*}
\sum\limits_{k=2}^{\infty}(-1)^k\cdot \bigg(\frac{1}{k-1}-\frac{1}{k+2}\bigg)
\end{align*}
Mit dem Leibnizkriterium ist zu zeigen, dass $ \big(|a_k|\big)_{k\in \mathbb{N}_{k\geq 2}} $ eine monoton fallende Nullfolge ist.\\
$ |q_k| $ ist bei dieser Reihe:\\
\begin{align*}
|a_k|&=\bigg|(-1)^k\cdot \bigg(\frac{1}{k-1}-\frac{1}{k+2}\bigg) \bigg|\\
&=\bigg|(-1)^k\cdot\frac{k+2-(k-1)}{(k-1)(k+2)}\bigg|\\
|a_k|&=\bigg|(-1)^k\cdot\frac{3}{(k-1)(k+2)}\bigg|\\
\end{align*}
\underline{I. Monotonie}\\
Zz.: $ |a_k|\geq |a_{k+1}| $ für $ k\in \mathbb{N}_{k\geq 2} $\\
\begin{align*}
\Rightarrow &|a_k|-|a_{k+1}|\geq 0\\
&\bigg|(-1)^k\cdot\frac{3}{(k-1)(k+2)}\bigg|-\bigg|(-1)^k\cdot\frac{1}{(k+1-1)(k+1+2)}\bigg|\geq 0\\
&\frac{3}{(k-1)(k+2)}-\frac{1}{k(k+3)}\geq 0\\
&\frac{3k(k+3)-3(k-1)(k+2)}{(k-1)(k+2)k(k+3)}\geq 0\\
&\frac{3k^2+9k-3(k^2+k-2)}{(k-1)(k+2)k(k+3)}\geq 0\\
&\frac{3k^2+9k-3k^2-3k+6}{(k-1)(k+2)k(k+3)}\geq 0\\
&\frac{6k+6}{(k-1)(k+2)k(k+3)}\geq 0\\
\end{align*}
Da der Zähler immer größer Null ist:\\
\begin{align*}
6k+6>0
\end{align*}
Und der Nenner immer größer Null ist:\\
\begin{align*}
(k-1)(k+2)k(k+3)>0
\end{align*}
Ist der Bruch immer größer Null.
Daraus folgt: $ \big(|a_k| \big)_{k\in \mathbb{N}_{k\geq 2}}>0\Rightarrow |a_k|$ ist monoton fallend.\\
\\
\underline{II. Nullfolge:}\\
\begin{align*}
\lim\limits_ {k\to\infty}|a_k|=\lim\limits_ {k\to\infty}\frac{3}{(k-1)(k+2)}=\lim\limits_ {k\to\infty}\frac{1}{k^2+k+2}=0
\end{align*}
Fazit: Wegen I.und II. ist $ \big(|a_k|\big)_k{k\in\mathbb{N}_{k\geq2}} $ eine monoton fallende Nullfolge.\\

\begin{align*}
\overset{(Leibniz)}{\Rightarrow} \sum\limits_{k=2}^{\infty} (-1)^k\cdot\bigg(\frac{1}{k-1}-\frac{1}{k+2}\bigg)\;\;\text{ist konvergent. }
\end{align*}
\subsubsection{Das Reihen-Integral-Kriterium}
\href{https://youtu.be/x1s0KUSIwaQ}{Video 14}\\
Sei $ f:[n,\infty[ \;\mapsto \mathbb{R}$ eine monoton fallende Funktion mit $ f(x) \geq 0$ für $ x\in [n,\infty[ \;$.\\
\begin{align*}
\int\limits_{n}^{\infty} f(x)dx\;\;\text{existiert }\Leftrightarrow \sum\limits_{k=n}^{\infty}f(k)\;\;\text{konvergiert }
\end{align*}
\begin{center}
	
	Wenn dieses Integral $ f(x) $ existiert,\\	
	\begin{tikzpicture}[xscale=2, yscale=2]
	
	% Gitternetzlinien
	\draw[xstep=1cm, ystep=1cm, very thin, color=lightgray] (0,2) grid (4,0);
	% x-Achse
	\draw[->] (0,0) -- (4,0) node[below,xshift=-1em] {$x$}; 
	%Zahlen auf x-Achse
	%\foreach \x in {-1,...,-1}
	%\draw[shift={(\x,0)},color=black] (0pt,2pt) -- (0pt,-2pt) node[below] {\footnotesize $\x$};
	\foreach \x in {1,...,3}
	\draw[shift={(\x,0)},color=black] (0pt,2pt) -- (0pt,-2pt) node[below] {\footnotesize $\x$};
	% y-Achse 
	\draw[->] (0,0) -- (0,2) node[left,yshift=-1em] {$f(x)$};%node[above left]
	%Zahlen auf y-Achse
	%\foreach \y in {0,...,-1}
	%\draw[shift={(0,\y)},color=black] (2pt,0pt) -- (-2pt,0pt) node[left] {\footnotesize $\y$};
	\foreach \y in {1,...,1}
	\draw[shift={(0,\y)},color=black] (2pt,0pt) -- (-2pt,0pt) node[left] {\footnotesize $\y$};
	%Ursprung
	%\draw[color=black] (0pt,-10pt) node[right] {\footnotesize $0$};
	
	\draw[thick,color=black] plot[samples=200, domain=0:4] 
	(\x,{e^(-\x)}) node[below, xshift=-15em,yshift=53] {$f(x)$};
	%\draw[color=black] (.5,0.2) node[right] {\footnotesize $\mathrm{A}=1$};
	\fill[color =red,fill opacity=0.3] plot[domain=0:4] (\x,{e^(-\x) })% 
	-- plot[domain=4:0] (\x,{0});% 
	
	
	
	\end{tikzpicture}\\
\end{center}
\begin{center}
	dann existiert auch die Folge $ \sum\limits_{k=n}^{\infty}f(k) $.
	
	\begin{tikzpicture}[xscale=2, yscale=2]
	
	% Gitternetzlinien
	\draw[xstep=1cm, ystep=1cm, very thin, color=lightgray] (0,2) grid (4,0);
	% x-Achse
	\draw[->] (0,0) -- (4,0) node[below,xshift=-1em] {$x$}; 
	%Zahlen auf x-Achse
	%\foreach \x in {-1,...,-1}
	%\draw[shift={(\x,0)},color=black] (0pt,2pt) -- (0pt,-2pt) node[below] {\footnotesize $\x$};
	\foreach \x in {1,...,3}
	\draw[shift={(\x,0)},color=black] (0pt,2pt) -- (0pt,-2pt) node[below] {\footnotesize $\x$};
	% y-Achse 
	\draw[->] (0,0) -- (0,2) node[left,yshift=-1em] {$f(x)$};%node[above left]
	%Zahlen auf y-Achse
	%\foreach \y in {0,...,-1}
	%\draw[shift={(0,\y)},color=black] (2pt,0pt) -- (-2pt,0pt) node[left] {\footnotesize $\y$};
	\foreach \y in {1,...,1}
	\draw[shift={(0,\y)},color=black] (2pt,0pt) -- (-2pt,0pt) node[left] {\footnotesize $\y$};
	%Ursprung
	%\draw[color=black] (0pt,-10pt) node[right] {\footnotesize $0$};
	
	\draw[thick,color=white] plot[samples=200, domain=0:4] 
	(\x,{e^(-\x)}) node[below, xshift=-15em,yshift=53] {$f(x)$};
	\draw[color=black] (1,0.45) node[right] {\footnotesize $f(1)$};
	\draw[color=black] (0,1.2) node[right] {\footnotesize $f(0)$};
	\draw[color=black] (2,0.25) node[right] {\footnotesize $f(2)$};
	\draw[color=black] (3,0.17) node[right] {\footnotesize $f(3)$};
	\draw[color=blue] (1,0.36) node {\textbullet};
	\draw[color=blue] (2,0.125) node {\textbullet};
	\draw[color=blue] (3,0.049) node {\textbullet};
	\draw[color=blue] (0,1) node {\textbullet};
	\fill[color =red,fill opacity=0.3] plot[domain=0:4] (\x,{e^(-\x) })% 
	-- plot[domain=4:0] (\x,{0});% 
	
	
	
	\end{tikzpicture}\\
\end{center}
\paragraph{Aufgabe zum Reihen-Integral-Kriterium}$  $\vspace*{3mm}\\
\href{https://youtu.be/x0s61IXjBKI}{Video 15}\\
Für welches $ \alpha\in\mathbb{R} $ konvergiert die Reihe $ \sum\limits_{k=1}^{\infty}\frac{1}{k^{\alpha}} $ ?\\
\renewcommand{\baselinestretch}{1.0}\normalsize
Es gilt:\\
\begin{align*}
\sum\limits_{k=1}^{\infty}\frac{1}{k^{\alpha}}\;\;\text{konvergiert }\Leftrightarrow\int\limits_{1}^{\infty}\frac{1}{x^{\alpha}} dx\;\;\text{existiert }
\end{align*}
\underline{1.Fall: $ \alpha=1 $}\\
\begin{align*}
&\int\limits_{1}^{\infty} \frac{1}{x}=\lim\limits_ {t\to\infty}\big[\ln(x)\big]_1^t=\lim\limits_ {t\to\infty}\ln(t)=\infty\\
&\Rightarrow \sum\limits_{k=1}^{\infty}\frac{1}{k}\;\;\text{divergiert }
\end{align*}
Bemerkung:\\
\begin{align*}
\sum\limits_{k=1}^{\infty}(-1)^k\frac{1}{k}\;\;\text{konvergiert }
\end{align*}
\underline{2.Fall:$ \alpha\neq 1 $}\\
\begin{align*}
\int\limits_{1}^{\infty} \frac{1}{x^{\alpha}}&=\lim\limits_ {t\to\infty}\int\limits_{1}^{t} x^{-\alpha}\\
&=\lim\limits_ {t\to\infty}\bigg|\frac{1}{1-\alpha}x^{1-\alpha}\bigg|_1^t\\
&=\lim\limits_ {t\to\infty}\bigg(\frac{1}{1-\alpha}t^{1-\alpha}-\frac{1}{1-\alpha}\bigg)
\end{align*}
\begin{quote}
	(i)\; wenn $ 1-\alpha>0 \Leftrightarrow\;\;\text{divergent }$\\
	(ii) wenn $ 1-\alpha<0 $
	\begin{quote}
		Es gilt für $ t^{1-\alpha} $ :\\
		Für $ \tilde{\alpha}=1-\alpha< 0\rightarrow \lim\limits_ {t\to\infty}\frac{1}{t^{\tilde{\alpha}}}=0$
	\end{quote}
	
	\begin{align*}
	&\frac{1}{1-\alpha}t^{1-\alpha}\;\;
	&\text{konvergiert gegen Null und der Grenzwert geht gegen } -\frac{1}{1-\alpha}=\frac{1}{\alpha-1}
	\end{align*}
\end{quote}
Fazit: Für $ \alpha>1 $ konvergiert $ \sum\limits_{k=1}^{\infty}\frac{1}{k^{\alpha}} $

%	\chapter{Potenzreihen}
\setcounter{section}{14}
\setcounter{subsection}{0}
\setcounter{subsubsection}{1}
\setcounter{secnumdepth}{4}
\subsection{Potenzreihen}
\href{https://youtu.be/f8ZfnOCjNYM}{Video 01}\\
Eine Potenzreihe hat die Form von:\\
\begin{align*}
	\sum\limits_{k=0}^{\infty}a_k(z-z_0)^k\;\;\text{ für }\;\;a_k, z_0, z\in \mathbb{C}
\end{align*}
Für $ a_k, z_0, z\in \mathbb{R} $ heißt die Potenzreihe reell.\\
Man nennt $ a_k $ die Koeffizienten und $ z_0 $ den Entwicklungspunkt.\\
Für jedes $ a_k, z_0, z $, die ja konstant sind bekommt man dann eine Reihe.\\
\begin{align*}
	\sum\limits_{k=0}^\infty{}\underbrace {a_k(z-z_0)^k}_{b_k}
\end{align*}
Sei $ \mathbb{D} \subset  \mathbb{C}$ die Menge aller $ z $ für die Reihe: $ \sum\limits_{k=0}^{\infty}a_k(z-z_0)^k $ konvergiert, dann ist: $f:\mathbb{D}\mapsto \mathbb{C}$ mit $ f(k)=\sum\limits_{k=0}^{\infty}a_k(z-z_0) ^k$ eine Funktion.\\
Für einen Punkt gilt:\\
\begin{align*}
	z_1\mapsto\sum\limits_{k=0}^{\infty}a_k(z_1-z_0)^k
\end{align*}
Eine Potenzreihe ist eine Reihe mit varablen Wert und für die Punkte für die Reihe konvergiert erhält man eine Funktion.\\
Z.B.:\\
\begin{align*}
	\sum\limits_{k=0}^{\infty}z^k
\end{align*}
Es gilt:\\
\begin{align*}
	&\sum\limits_{k=0}^{\infty}z^k=\sum\limits_{k=0}^{\infty}1\cdot(z-0)^k\\
	&a_k=1\;\;\text{ist der Koeffizient }\\
	&z_0=0\;\;\text{ist der Entwicklungspunkt }
\end{align*}
Diese Reihe konvergiert für $ |z|<1 $ (siehe geometrische Reihe)\\
$\bullet$ Für $  z\in\mathbb{C} $ gilt:\\
\begin{quote}
	\begin{align*}
		&	f:\{z\in \mathbb{C}\;|\;|z|<1\}\longmapsto \mathbb{C}\\
		&z\longmapsto\sum\limits_{k=0}^{\infty}z^k=\frac{1}{1-z}
	\end{align*}
	Welche Zahlen sind dies?
	\begin{center}
		
		\begin{tikzpicture}[scale=2.5]
			\filldraw [->,draw=black,fill=red!20] (0,0)--(10mm,0mm) arc (0:360:10mm);
			%	\filldraw[->,draw=black,fill=green!20] (0,0) -- (7mm,0mm) arc (0:45:7mm); 
			%	\filldraw[->,fill=orange!20,draw=black,rotate=-45] (0,0) -- (8mm,0mm) arc (0:45:8mm); 
			
			% Gitternetzlinien
			\draw[xstep=1cm, ystep=1cm, very thin, color=lightgray] (-2,-2) grid (2,2);
			% x-Achse
			\draw[->] (-2,0) -- (2,0) node[below,xshift=-1em] {$\mathrm{\mathrm{Re}}$}; 
			%Zahlen auf x-Achse
			\foreach \x in {-1,...,-1}
			\draw[shift={(\x,0)},color=black] (0pt,2pt) -- (0pt,-2pt) node[below] {\footnotesize $\x$};
			\foreach \x in {1,...,1}
			\draw[shift={(\x,0)},color=black] (0pt,2pt) -- (0pt,-2pt) node[below] {\footnotesize $\x$};
			% y-Achse 
			\draw[->] (0,-2) -- (0,2) node[left,yshift=-1em] {$\mathrm{\mathrm{Im}}$};%node[above left]
			%Zahlen auf y-Achse
			\foreach \y in {-1,...,-1}
			\draw[shift={(0,\y)},color=black] (2pt,0pt) -- (-2pt,0pt) node[left] {\footnotesize $\y$};
			\foreach \y in {1,...,1}
			\draw[shift={(0,\y)},color=black] (2pt,0pt) -- (-2pt,0pt) node[left] {\footnotesize $\y$};
			%Ursprung
			\draw[color=black] (-4pt,-4pt) node[right] {\footnotesize $0$};
			
			
			%	\draw [blue,->] (0,0)-.(1,0)node[below,xshift= 0.5em,yshift=1em] {\footnotesize $1+\mathrm{i}$};
			
			%	\draw [blue,->] (0,0)-.(1,-1)node[below,xshift= 0.5em,yshift=1em] {\footnotesize $1-\mathrm{i}$};
			
			\draw [blue] (0.2,0.2) node[xshift= 0.8em,yshift=-0.6em,rotate=0] {\footnotesize $|z|<1$};
			%	\draw [blue] (-0.2,-0.2) node[xshift= 3.8em,yshift=0.1em,rotate=0] {\footnotesize $\overset{\AC}{\alpha}=\frac{\pi}{4}$};
			%	\draw [blue] (-0.6,0.6) node[xshift= 1.3em,yshift=-0.9em,rotate=0] {\footnotesize $\beta=2\pi-\overset{\AC}{\alpha}$};
			
		\end{tikzpicture}\\
		Für alle Punkte innerhalb des Kreises konvergiert diese Reihe.
	\end{center}
\end{quote}
$\bullet$ Für $  z\in\mathbb{R} $ gilt:\\
\begin{quote}
	\begin{align*}
		&f:\{z\in\mathbb{R}\;|\;|z|<1=\;]-1,1[ \;\}\longmapsto\mathbb{R}\\
		&z\longmapsto\sum\limits_{k=0}^{\infty}z^k=\frac{1}{1-z}
	\end{align*}
	Für alle Zahlen im Intervall: $ ]-1,1[ $ konvergiert diese Reihe.\\
	\renewcommand{\baselinestretch}{2.0}\normalsize
	
	\begin{tikzpicture}
		\node at (2,1) [rotate=0,color=black] {$]$};
		\node at (2,0.5) [rotate=0, color=black] {$-1$};
		\draw[->] (0,1) -- (6,1) node [below]{$\mathbb{R}$};
		\draw  (2,1) -- (4,1)[ color=red] ;
		\node at (4,1) [rotate=0,color=black] {$[$};
		\node at (4,0.5) [rotate=0, color=black] {$1$};
	\end{tikzpicture}\ \\
\end{quote}
\renewcommand{\baselinestretch}{1.0}\normalsize
Bemerkung:\\
Die Potenzreihe ist eine Reihe mit konstanten Wert. Konvergiert diese Reihe dann ist sie auch eine Funktion.\\
Im komplexen Bereich konvergiert die Reihe innerhalb eines Kreises (2-Dimensional) und im reellen Bereich in einem 1-Dimensonalen Kreis.\\

\subsection{Konvergenz von Potenzreihen}
\href{https://youtu.be/dGwysWKpoE0}{Video 02}\\
Betrachte folgende Reihe:\\
\begin{align*}
	\sum\limits_{k=0}^{\infty}a_k(z-z_0)^k
\end{align*}
Für welche $ z\in \mathbb{C} $($ z\in\mathbb{R} $) ist die Potenzreihe konvergent?\\
Def.:\\
\begin{align*}
	\sum\limits_{k=0}^{\infty}\underbrace {a_k(z-z_0)^k}_{b_k}
\end{align*}
Wann ist $\sum\limits_{k=0}^{\infty}b_k  $ konvergent?\\
Nach dem Quotientenkriterium gilt:\\
\begin{align*}
	\lim\limits_ {k\to\infty}\bigg|\frac{b_{k+1}}{b_k}\bigg|<1\Rightarrow\;\;\text{Konvergenz }
\end{align*}
Für welche $ z $ gilt dies?\\
\begin{align*}
	\lim\limits_{k\to\infty}\bigg|\frac{b_{k+1}}{b_k}\bigg|=\lim\limits_{k\to\infty}\bigg|\frac{a_{k+1}(z-z_0)^{k+1}}{a_k(z-z_0)^k}\bigg|=\bigg|\frac{a_{k+1}}{a_k}\bigg|\big|z-z_0\big|<1
\end{align*}
Es gilt:\\
Für $ <1 $ konvergent.\\
Für $ >1 $ divergent.\\
Wenn kein Grenzwert existiert und wenn der Grenzwert gegen 1 geht, dann kann man keine Aussage treffen!\\
Nennt man nun $ \lim\limits_ {k\to\infty}\bigg|\frac{b_{k+1}}{b_k}\bigg|=\mathrm{A} $
Dann gilt:\\
\begin{align*}
	\lim\limits_ {k\to\infty}\bigg|\frac{b_{k+1}}{b_k}\bigg|=\mathrm{A}(z-z_0)<1
\end{align*}
\underline{1.Fall:}\\
\begin{quote}
	\begin{align*}
		\mathrm{A}|z-z_0|<1\Leftrightarrow |z-z_0|<\frac{1}{\mathrm{A}}
	\end{align*}
	Die Potenzreihe konvergiert.\\
\end{quote}
\underline{2.Fall:}\\
\begin{quote}
	\begin{align*}
		\mathrm{A}|z-z_0|>1\Leftrightarrow |z-z_0|>\frac{1}{\mathrm{A}}
	\end{align*}
	Die Potenzreihe konvergiert nicht.\\
\end{quote}
\underline{Spezielle Fälle:}\\
\begin{quote}
	$\bullet  \;\; \mathrm{A}=1 $ oder kein Grenzwert.\\
	\begin{quote}
		Kann keine Aussage getroffen werden.
	\end{quote}
	$\bullet  \;\; \mathrm{A}=\infty $ wird zugelassen.\\
	\begin{quote}
		\begin{align*}
			|z-z_0|<\frac{1}{\underbrace {\mathrm{A}}_{\mapsto 0}}
		\end{align*}
		Dies gilt nur, wenn $ z-z_0=0 $ also der Grenzwert wird Null (trivial). Die Reihe konvergiert.\\
	\end{quote}
	$\bullet  \;\; \mathrm{A}=0 $
	\begin{quote}
		\begin{align*}
			\mathrm{A}|z-z_0|=0|z-z_0|<1
		\end{align*}
		Potenzreihe konvergiert für alle $ z $.\\
	\end{quote}
\end{quote}
Die Potenzreihe:\\
\begin{align*}
	\sum\limits_{k=0}^{\infty}a_k(z-z_0)^k
\end{align*}
konvergiert für alle $ |z-z_0| $, wenn gilt:\\
\begin{align*}
	|z-z_0|<R
\end{align*}
wobei für $ R $ gilt:\\
\begin{align*}
	R:=\frac{1}{\lim\limits_ {k\to\infty}\bigg|\frac{a_{k+1}}{a_k}\bigg|}
\end{align*}
\newpage
\noindent
Bemerkung:\\
$ |z-z_0|<R\;\;z\in \mathbb{C} $ ist ein Kreis von Radius $ R $ um $ z_0 $\\
\begin{center}
	
	\begin{tikzpicture}[scale=2.5]
		%\filldraw [->,draw=black,fill=red!20] (0,0)--(12mm,0mm) arc (0:360:12mm);
		%	\filldraw[->,draw=black,fill=green!20] (0,0) -- (7mm,0mm) arc (0:45:7mm); 
		%	\filldraw[->,fill=orange!20,draw=black,rotate=-45] (0,0) -- (8mm,0mm) arc (0:45:8mm); 
		\filldraw [->,draw=black,fill=green!20](0.5,0.5) circle (11mm);
		% Gitternetzlinien
		\draw[xstep=1cm, ystep=1cm, very thin, color=lightgray] (-2,-2) grid (2,2);
		% x-Achse
		\draw[->] (-2,0) -- (2,0) node[below,xshift=-1em] {$\mathrm{\mathrm{Re}}$}; 
		%Zahlen auf x-Achse
		\foreach \x in {-1,...,-1}
		\draw[shift={(\x,0)},color=black] (0pt,2pt) -- (0pt,-2pt) node[below] {\footnotesize $\x$};
		\foreach \x in {1,...,1}
		\draw[shift={(\x,0)},color=black] (0pt,2pt) -- (0pt,-2pt) node[below] {\footnotesize $\x$};
		% y-Achse 
		\draw[->] (0,-2) -- (0,2) node[left,yshift=-1em] {$\mathrm{\mathrm{Im}}$};%node[above left]
		%Zahlen auf y-Achse
		\foreach \y in {-1,...,-1}
		\draw[shift={(0,\y)},color=black] (2pt,0pt) -- (-2pt,0pt) node[left] {\footnotesize $\y$};
		\foreach \y in {1,...,1}
		\draw[shift={(0,\y)},color=black] (2pt,0pt) -- (-2pt,0pt) node[left] {\footnotesize $\y$};
		%Ursprung
		\draw[color=black] (-4pt,-4pt) node[right] {\footnotesize $0$};
		
		
		%	\draw [blue,->] (0,0)-.(1,0)node[below,xshift= 0.5em,yshift=1em] {\footnotesize $1+\mathrm{i}$};
		
		%	\draw [blue,->] (0,0)-.(1,-1)node[below,xshift= 0.5em,yshift=1em] {\footnotesize $1-\mathrm{i}$};
		\draw[blue,->](5mm,5mm) node[xshift= 0,yshift=-3mm,rotate=0] {\footnotesize $z_0$};
		\draw[blue,->] (0.5,0.5)--(12mm,13.3mm)node[below,xshift= 0.5em,yshift=1em] {\footnotesize $R$};
		
		%	\draw [blue] (-0.2,-0.2) node[xshift= 3.8em,yshift=0.1em,rotate=0] {\footnotesize $\overset{\AC}{\alpha}=\frac{\pi}{4}$};
		%	\draw [blue] (-0.6,0.6) node[xshift= 1.3em,yshift=-0.9em,rotate=0] {\footnotesize $\beta=2\pi-\overset{\AC}{\alpha}$};
		
	\end{tikzpicture}\\
\end{center}
$ |z-z_0|<\mathrm{R}\;\;z\in \mathbb{R} $ ist eine Gerade mit Abstand $ R $ um $ z_0 $\\
\begin{center}
	Für $ z_0=2$ gilt:\\ 
	
	\begin{tikzpicture}
		\node at (5,1) [rotate=0,color=black] {$|$};
		\node at (2,1) [rotate=0,color=black] {$]$};
		\node at (2,0.5) [rotate=0, color=black] {$2-\mathrm{R}$};
		\draw[->] (0,1) -- (10,1) node [below]{$\mathbb{R}$};
		\draw  (2,1) -- (8,1)[ color=red] ;
		\node at (8,1) [rotate=0,color=black] {$[$};
		\node at (8,0.5) [rotate=0, color=black] {$2+\mathrm{R}$};
		\node at (5,0.5) [rotate=0, color=black] {$2$};
	\end{tikzpicture}\ \\
\end{center}
\subsubsection{ Aufgabe 1 zur Konvergenz von Potenzreihen}
\href{https://youtu.be/PZkcdZNx7I4}{Video 03}\\
Bestimme den Konvergenzradius von:\\
\begin{align*}
	\sum\limits_{k=0}^{\infty}\frac{\mathrm{e}^{2k}\cdot x^k}{k+1}
\end{align*}
Zur Berechnung des Konvergenzradius formt man entsprechend um:\\
\begin{align*}
	\sum\limits_{k=0}^{\infty}\frac{\mathrm{e}^{2k}\cdot x^k}{k+1}=\sum\limits_{k=0}^{\infty}\underbrace {\frac{\mathrm{e}^{2k}}{k+1}}_{a_k}(x-0)^k
\end{align*}
Es gilt mit $ a_k=\frac{\mathrm{e}^{2k}}{k+1} $:\\
\begin{align*}
	\mathrm{R}=\frac{1}{\lim\limits_ {k\to\infty}\big|\frac{a_{k+1}}{a_k}\big|}
\end{align*}
Nun den Grenzwert berechen:\\
\begin{align*}
	\lim\limits_ {k\to\infty}\bigg|\frac{a_{k+1}}{a_k}\bigg|=\lim\limits_ {k\to\infty}\bigg|\frac{\mathrm{e}^{2(k+1)}(k+1)}{(k+2)\mathrm{e}^{2k}}\bigg|=\lim\limits_ {k\to\infty}\mathrm{e}^2\;\;\frac{k+1}{k+2}=\mathrm{e}^2\lim\limits_ {k\to\infty}\frac{k\big(\frac{1}{k}+1\big)}{k\big(\frac{2}{k}+1\big)}=\mathrm{e}^2
\end{align*}
Für den Konvergenzradius gilt:\\
\begin{align*}
	\mathrm{R}=\frac{1}{\lim\limits_ {k\to\infty}\big|\frac{a_{k+1}}{a_k}\big|}=\frac{1}{\mathrm{e}^2}
\end{align*}
Fazit:\\
\begin{align*}
	\mathrm{R}=\frac{1}{\mathrm{e}^2}
\end{align*}
Für $ |x-\underbrace {x_0}_{x_0=0}|=|x|<\frac{1}{\mathrm{e}^2} $ ist die Reihe:$\sum\limits_{k=0}^{\infty}\frac{\mathrm{e}^{2k}\cdot x^k}{k+1}  $ konvergent und definiert dort eine Funktion.\\
\\
Für die Randpunkte muss man noch die Konvergenz untersuchen.\\
Es gilt:\\
\begin{align*}
	|x|<\frac{1}{\mathrm{e}^2}\Rightarrow x_1=\frac{1}{\mathrm{e}^2},\; x_2=-\frac{1}{\mathrm{e}^2}
\end{align*}
\textbullet\;\; Für$\;\;x_1=\frac{1}{\mathrm{e}^2} $ gilt:\\
\begin{align*}
	\sum\limits_{k=0}^{\infty}\frac{\mathrm{e}^{2k}\cdot \big(\frac{1}{\mathrm{e}^2}\big)^k}{k+1}=\sum\limits_{k=0}^{\infty}\frac{\mathrm{e}^{2k}}{k+1}\bigg(\frac{1}{\mathrm{e}^2}\bigg)^k=\sum\limits_{k=0}^{\infty}\frac{\mathrm{e}^{2k}}{(k+1)\mathrm{e}^{2k}}=\sum\limits_{k=0}^{\infty}\frac{1}{k+1}\Rightarrow\text{divergent (Harmonische Reihe) }
\end{align*}
\textbullet\;\; Für$\;\;x_1=-\frac{1}{\mathrm{e}^2} $ gilt:\\
\begin{align*}
	\sum\limits_{k=0}^{\infty}\frac{\mathrm{e}^{2k}}{k+1}\bigg(\frac{-1}{\mathrm{e}^2}\bigg)^k&=\sum\limits_{k=0}^{\infty}\frac{\mathrm{e}^{2k}}{k+1}(-1)^k\bigg(\frac{1}{\mathrm{e}^2}\bigg)^k=\sum\limits_{k=0}^{\infty}\frac{\mathrm{e}^{2k}}{k+1}(-1)^k\bigg(\frac{1}{\mathrm{e}^{2k}}\bigg)\\&=\sum\limits_{k=0}^{\infty}(-1)^k\frac{1}{k+1}\Rightarrow\text{konvergent (Leibnizkritrium) }
\end{align*}
Fazit:\\
Die Reihe $ \sum\limits_{k=0}^{\infty}\frac{\mathrm{e}^{2k}\cdot x^k}{k+1} $ ist konvergent im Bereich von: $ x\in\big[-\frac{1}{\mathrm{e}^2},\frac{1}{\mathrm{e}^2}\big[\; $.\\
\newpage
\noindent
\subsubsection{ Aufgabe 2 zur Konvergenz von Potenzreihen}
\href{https://youtu.be/8TeaUxqFeuk}{Video 04}\\
Bestimme den Konvergenzradius von:\\
\begin{align*}
	\sum\limits_{k=1}^{\infty}\frac{(4z-2)^k}{k^4}\;\;\text{für}\;\;z\in \mathbb{C}
\end{align*}
Zum Konvergenzradius:\\
\begin{align*}
	\sum\limits_{k=1}^{\infty}\frac{(4z-2)^k}{k^4}=\sum\limits_{k=1}^{\infty}\frac{(4(z-\frac{1}{2}))^k}{k^4}=\sum\limits_{k=0}^{\infty}\underbrace {\frac{4^k}{k^4}}_{=a_k}\bigg(z-\frac{1}{2}\bigg)^k
\end{align*}
Es gilt mit $ a_k=\frac{4^k}{k^4} $:
\begin{align*}
	\mathrm{R}=\frac{1}{\lim\limits_ {k\to\infty}\big|\frac{a_{k+1}}{a_k}\big|}
\end{align*}
Nun den Grenzwert berechnen:\\
\begin{align*}
	\lim\limits_ {k\to\infty}\bigg|\frac{a_{k+1}}{a_k}\bigg|=\lim\limits_ {k\to\infty}\frac{4^{k+1}}{(k+1)^4}\frac{k^4}{4^k}=\lim\limits_ {k\to\infty}4\bigg(\frac{k}{k+1}\bigg)^4=\lim\limits_ {k\to\infty}4\bigg(\underbrace {\frac{k}{k(1+\frac{1}{k})}}_{=1}\bigg)=4
\end{align*}
Folglich ist der Konvergenzradius von: $\sum\limits_{k=1}^{\infty}\frac{(4z-2)^k}{k^4}  $:\\
\begin{align*}
	\mathrm{R}&=\frac{1}{\lim\limits_ {k\to\infty}\big|\frac{a_{k+1}}{a_k}\big|}=\frac{1}{4}\\
	\mathrm{R}&=\frac{1}{4}
\end{align*}
\newpage
\noindent
\subsection{Differentiation von Potenzreihen }
\href{https://youtu.be/kpaMK4KSYDg}{Video 05}\\
\begin{center}
	
	
	\fbox{\parbox  {155mm} {
			Sei $\sum\limits_{k=0}^{\infty}a_k(x-x_0)^k  $ eine reelle Potenzreihe mit Konvergenzradius: $ \mathrm{R}>0$.\\ Dann ist: 
			\begin{align*}
				f:\; ]x_o-\mathrm{R} , x_0+\mathrm{R}[\; \longmapsto \mathbb{R}\;\; \text{mit } f(x)=\sum\limits_{k=0}^{\infty}a_k(x-x_0)^k
			\end{align*}\\
			eine differenzierbare Funktion mit: 
			\begin{align*}
				f'(x)=\sum\limits_{k=0}^{\infty}\big(a_k(x-x_0)^k\big)'
	\end{align*}}}
\end{center}
Bemerkung:\\
Wie schon bekannt ist eine Potenzreihe mit $ \mathrm{R}>0 $ zwischen $ ]x_o-\mathrm{R} , x_0+\mathrm{R}[  $ eine Funktion und deshalb auch differenzierbar.\\
\\
Für einen Term gilt:\\
\begin{align*}
	\big(a_k(x-x_0)^k\big)'=a_k\cdot k(x-x_0)^{k-1}\\
\end{align*}
Dann gilt für die Reihe:\\
\begin{align*}
	\sum\limits_{k=0}^{\infty}\big(a_k(x-x_0)^k\big)'=\sum\limits_{k=0}^{\infty}a_k\cdot k(x-x_0)^{k-1}
\end{align*}
Für $ k=0 $ gilt:\\
\begin{align*}
	a_0\cdot 0(x-x_0)^{0-1}=0
\end{align*}
Dann gilt:\\
\begin{align*}
	\sum\limits_{k=1}^{\infty}a_k\cdot k(x-x_0)^{k-1}
\end{align*}
Mit Indexverschiebung ergibt sich:\\
\begin{align*}
	\sum\limits_{k=0}^{\infty}a_{k+1}(k+1)(x-x_0)^{k}
\end{align*}
Für die Potenzreihe mit dem Konvergenzradius: $ \mathrm{R}>0 $ gilt dann:\\
\begin{align*}
	\sum\limits_{k=0}^{\infty}\big(a_k(x-x_0)^k\big)'=\sum\limits_{k=0}^{\infty}a_{k+1}(k+1)(x-x_0)^{k}
\end{align*}
Beispiel:\\
Wir betrachten:\\
\begin{align*}
	\sum\limits_{k=0}^{\infty} x^k\;\;\text{mit Konvergenzradius }\mathrm{R}=1
\end{align*}
Es gilt für:\\
\begin{align*}
	|x|<1 : f(x)=\sum\limits_{k=0}^{\infty}x^k=\frac{1}{1-x}
\end{align*}
Die Funktion $ f(x) $ ist differenzierbar dann gilt:\\
\begin{align*}
	f'(x)&=\sum\limits_{k=0}^{\infty}\big(a_k(x-x_0)^k\big)'=\sum\limits_{k=0}^{\infty}a_{k+1}(k+1)(x-x_0)^{k}\\
	a_k&=1\\
	x_0&=0\\
&=\sum\limits_{k=0}^{\infty}\big(1\cdot(x-0)^k\big)'=\sum\limits_{k=0}^{\infty}1\cdot(k+1)(x-0)^k\\
	f'(x)	&=\sum\limits_{k=0}^{\infty}(k+1)(x)^k\\
\end{align*}
Für $ f'(x)=\big(\frac{1}{1-x}\big)' $ gilt:\\
\begin{align*}
	f'(x)=\bigg(\frac{1}{1-x}\bigg)'=\bigg((1-x)^{-1}\bigg)'=-1\cdot(1-x)^{-2}(-1)=\frac{1}{(1-x)^2}
\end{align*}
Fazit:\\
\begin{align*}
	\sum\limits_{k=0}^{\infty}(k+1)x^k=\frac{1}{(1-x)^2}\;\;\text{und hat Konvergenzradius: }\mathrm{R}=1
\end{align*}
\newpage
\noindent
\subsection{Integration von Potenzreihen}
\href{https://youtu.be/0rKwoivWboE}{Video 06}\\
\begin{center}
	
	
	\fbox{\parbox  {155mm} {
			Sei $\sum\limits_{k=0}^{\infty}a_k(x-x_0)^k  $ eine reelle Potenzreihe mit Konvergenzradius: $ \mathrm{R}>0$.\\ Dann ist: 
			\begin{align*}
				f:\; ]x_o-\mathrm{R} , x_0+\mathrm{R}[\; \longmapsto \mathbb{R}\;\; \text{mit } f(x)=\sum\limits_{k=0}^{\infty}a_k(x-x_0)^k
			\end{align*}\\
			eine integrierbare Funktion mit:
			\begin{align*}
				\int\limits_{}^{}f(x) \mathrm{d}x=\sum\limits_{k=0}^{\infty}\int\limits_{}^{} a_k(x-x_0)^k dx
			\end{align*}	
	}}
\end{center}
Es gilt für die Potenzreihe:\\
\begin{align*}
	\int\limits_{}^{} a_x(x-x_0)^k \mathrm{d}x=a_k\cdot \frac{(x-x_0)^{k+1}}{k+1}+\alpha_k\;\;\text{mit }\alpha_k\in\mathbb{R}
\end{align*}
Dies ergibt eine neue Potenzreihe mit dem gleichen Konvergenzradius: $ \mathrm{R}=1 $\\
\\
Beispiel:\\
Betrachte:\\
\begin{align*}
	f(x)=\sum\limits_{k=0}^{\infty}x^k\;\;\text{mit }\;\;|x|<1\;\;\mathrm{R}=1
\end{align*}
Für das Integral gilt:\\
\begin{align*}
	\int\limits_{}^{}f(x) \mathrm{d}x=\sum\limits_{k=0}^{\infty}\int\limits_{}^{} x^k dx=\sum\limits_{}^{}\frac{x^{k+1}}{k+1}+\alpha\;\;\text{mit }\alpha\in\mathbb{R}
\end{align*}
Nun gilt auch:\\
\begin{align*}
	&\sum\limits_{k=0}^{\infty}x^k=\frac{1}{1-x}\\
	&\Rightarrow \int\limits_{}^{} \frac{1}{1-x} \mathrm{d}x=-\ln(1-x)+\alpha,\;\;\text{mit } \alpha\in\mathbb{R}
\end{align*}
Fazit:\\
\begin{align*}
	\sum\limits_{k=0}^{\infty}\frac{x^{k+1}}{k+1}=-\ln(1-x)=\ln\bigg(\frac{1}{1-x}\bigg)
\end{align*}
Bemerkung:\\
\begin{align*}
	-\ln(a)=\ln(a^{-1})=\ln\bigg(\frac{1}{a}\bigg)
\end{align*}
Durch die Integation kommt man zu einer neuen Potenzreihe mit dem gleichen Konvergenzradius $ \mathrm{R}=1 $ wie bei der Funktion: $ f(x)=\sum\limits_{k=0}^{\infty}x^k $.
\subsection{Die Taylorreihe}
\href{https://youtu.be/lnDbJwbBlqY}{Video 07}\\
\begin{center}
	\underline{Taylorreihe}	
\end{center}
Sei $ f:\mathbb{R}\supset \mathbb{I}\rightarrow\mathbb{R}\;\; n$-mal differenzierbar (im Intervall $  \mathbb{I} $) dann \\
gilt mit $ x_0\in \mathbb{I} $:\\
\begin{align*}
	f(x)=\sum\limits_{k=0}^{n}\frac{f^k (x_0)}{k!}(x-x_0)^k+\mathrm{R}_n(x)
\end{align*}
Wobei der Grenzwert vom Rest gegen Null geht für:\\
\begin{align*}
	\lim\limits_ {x\to x_0}\frac{\mathrm{R}_n (x)}{(x-x_0)^k}=0
\end{align*}
Wenn die Funktion beliebig oft differenzierbar ist, dann gilt für den Rest:\\
\begin{align*}
	\lim\limits_ {n\to\infty}\mathrm{R}_n (x)=0
\end{align*}
Daraus folgt dann, dass man die Funktion auch als Potenzreihe schreiben kann, weil der Rest Null wird.\\
\begin{align*}
	f(x)=\sum\limits_{k=0}^{\infty}a_k(x-x_0)^k=\sum\limits_{k=0}^{\infty}\frac{f^k (x_0)}{k!}(x-x_0)^k
\end{align*}
\subsubsection{Beispiel zur Taylorreihe}
\href{https://youtu.be/AJxupcdeQJc}{Video 08}\\
Beispiel:\\
Betrachten wir eine beliebig oft differnzierbare Funktion:\\
\begin{align*}
	f(x)=\mathrm{e}^x\;\;\text{und  } x_0=0
\end{align*}
Es gilt:\\
\begin{align*}
	&f^k (x)=\mathrm{e}^x\\
	&f^k (x_0)=\mathrm{e}^0=1
\end{align*}
Die Taylorreihe lautet dann:\\
\begin{align*}
	f(x)=\sum\limits_{k=0}^{n}\frac{1}{k!}(x-x_0)^k+\mathrm{R}_n(x)=\frac{x^k}{k!}+\mathrm{R}_n(x)
\end{align*}
\underline{Lagrange'sche Restglieddarstellung:}
\begin{align*}
	\mathrm{R}_n(x)&=\frac{f^{k+1}(a)}{(n+1)!}\cdot x^{n+1}\;\;\text{mit $ a $ zwischen $ x $ und $ x_0 $}\\
	&=\frac{\mathrm{e}^{a}}{(n+1)!}\cdot x^{n+1}
\end{align*}
Betrachte:\\
\begin{align*}
	\lim\limits_ {n\to\infty}f(x)&=\lim\limits_ {n\to\infty}\mathrm{e}^{a}\frac{x^{n+1}}{(n+1)!}=\mathrm{e}^{a}\lim\limits_ {n\to\infty}\frac{x^{n+1}}{(n+1)!}\\
	&=\mathrm{e}^{a}\lim\limits_ {n\to\infty}\frac{x\cdot x\cdot x\cdot\dots}{1\cdot 2 \cdot 3\dots  n (n+1)}\\
	&=\mathrm{e}^{a}\bigg(\lim\limits_ {n\to\infty}\bigg(\frac{x}{1}\bigg)\cdot \lim\limits_ {n\to\infty}\bigg(\frac{x}{2}\bigg)\cdot \lim\limits_ {n\to\infty}\bigg(\frac{x}{3}\bigg)\dots \underbrace {\lim\limits_ {n\to\infty}\bigg(\frac{x}{n}\bigg)}_{=0}\cdot \underbrace {\lim\limits_ {n\to\infty}\bigg(\frac{x}{n+1}\bigg)}_{=0}\bigg)=0
\end{align*}
Fazit:\\
Mit $ f(x)=\mathrm{e}^x $ und $ x_0=0 $ ergibt sich als Taylorreihe:\\
\begin{align*}
	\mathrm{e}^x=f(x)=\sum\limits_{k=0}^{\infty}\frac{1}{k!}(x-x_0)^k=\sum\limits_{k=0}^{\infty}\frac{1}{k!}x^k
\end{align*}
\underline{Bestimmung des Konvergenzradius:}\\
\begin{align*}
	\mathrm{R}=\frac{1}{\lim\limits_ {k\to\infty}\big|\frac{a_{k+1}}{a_k}\big|}
\end{align*}
Es gilt:\\
\begin{align*}
	a_k=\frac{1}{k!}
\end{align*}
\begin{align*}
	\bigg|\frac{a_{k+1}}{a_k}\bigg|=\frac{\frac{1}{(k+1)!}}{\frac{1}{k!}}=\frac{k!}{(k+1)!}
\end{align*}
Für den Grenzwert gilt dann:\\
\begin{align*}
	\lim\limits_ {k\to\infty}\frac{k!}{(k+1)!}=\lim\limits_ {k\to\infty}\frac{1\cdot 2\cdot 3\dots k}{1\cdot 2\cdot 3\dots k(k+1)}=\lim\limits_ {k\to\infty}\frac{1}{k+1}=0\\
\end{align*}
Der Konvergenzradius ist dann:\\
\begin{align*}
	\mathrm{R}=\frac{1}{\lim\limits_ {k\to\infty}\big|\frac{a_{k+1}}{a_k}\big|}=\frac{1}{\lim\limits_ {k\to\infty}\frac{1}{k+1}}=\infty
\end{align*}
Fazit:\\
Die Taylorreihe von $ \mathrm{e}^x $ konvergiert für alle $ x\in\mathbb{R} $\\


































































%	\chapter{Klausurtraining}

\setcounter{section}{15}
\setcounter{subsection}{0}
\setcounter{subsubsection}{1}
\setcounter{secnumdepth}{4}
\subsection{Klausurtipps}
\href{https://youtu.be/uMfjDH8dFl4}{Video 01}\\
Klausurvorbereitungzeit  (Zeitaufwand ca. 7-10 Tage, a 5-8 Stunden):\\

\begin{tabular}[t]{ll}% l=linksbündig r=rechtsbündig
	\textbullet & Alle wichtigen Sätze in einer Zusammenfassung aufschreiben.  \\
	&\\
	\textbullet&Alle Aufgaben aller Aufgabenblätter alleine durchgehen und bearbeiten und danach mit\\
	& Musterlösung vergleichen. \\
	& Liste von problematischen Aufgaben erstellen.\\
	&\\
	\textbullet&Bearbeite alle relevanten Aufgaben aus der Vorlesung.\\
	&Liste von problematischen Aufgaben erstellen.\\
	&\\
	\textbullet&Bearbeite Altklausuren\\
	&\\
	\textbullet&Listen durchgehen\\
	&\\
	\textbullet&Simuliere Klausursituationen\\
	&\\
\end{tabular}\\
\\
In der Klausur:\\

\begin{tabular}[t]{ll}% l=linksbündig r=rechtsbündig
	\textbullet & Überfliege Klausur und schätze Schwierigkeitsgrad der Aufgaben ein.  \\
	&\\
	\textbullet&Bearbeite Klausur von \glqq leicht nach schwer\grqq durch.\\
	&\\
	\textbullet&Lese Aufgabenstellung gründlich und überlege eine Strategie\\
	&zur Bearbeitung und fange dann an!\\
	&\\
	\textbullet&Ruhig bleiben: Bei Fehlern mit falschen Ergebnissen weiter rechnen.\\
	&\\
	\textbullet&Schreibe keinen \glqq Quatsch\grqq: Z.B. keine allgemeine Formel.\\
	&\\
\end{tabular}\\
\\
Nach der Klausur: Gehe in die Einsicht!
\newpage
\noindent
\subsection{Aufgabe 1 - Komplexe Zahlenmenge}
\href{https://youtu.be/mBNzaQIcrM8}{Video 02}\\
Skizziere folgende Menge in der Gauß'schen Zahlenebene:\\
\begin{align*}
\{ z\in\mathbb{C} \;|\; z^4 =8\sqrt{2} (1+\mathrm{i})  \}\cap\{ z\in \mathbb{C}\;|\;  |z| \leq 2 \}
\end{align*}
\underline {Bestimmung der Menge: $  \mathbb{A} $}\\
\begin{align*}
\mathbb{A}:=\{ z\in\mathbb{C} \;|\; z^4 =8\sqrt{2} (1+\mathrm{i})  \}\
\end{align*}
\begin{quote}
	\textbullet\;\;Lösen von $z^4 =8\sqrt{2} (1+\mathrm{i})  $\\
	\begin{quote}
		Es gilt:\\\\
		\begin{align*}
		z^n&=|a|\cdot\mathrm{e}^{\mathrm{i}\alpha}\\
		z_k&=\sqrt[n]{|a|}\cdot\mathrm{e}^{\mathrm{i}\frac{\alpha+2k\pi}{n}}\;\;\text{für }\;\;k=0,1,2,\ldots,n-1
		\end{align*}
		\underline{Darstellung der komplexen Zahl $ a=8\sqrt{2}(1+\mathrm{i}) $ in Polarkoordinaten.}\vspace*{4mm}	\\
		Bestimmung des Betrage von $ |a| $.\\
		\begin{align*}
		|a|=|8\sqrt{2}(1+\mathrm{i})|=8\sqrt{2}|(1+\mathrm{i})|=8\sqrt{2}\cdot\sqrt{|1^2|+|\mathrm{i}^2|}=8\sqrt{2}\sqrt{1^2+1^2}=8\sqrt{2}\sqrt{2}=8\cdot 2=16
		\end{align*}
		
		Bestimmung des Winkels $ \alpha $:\\
		\begin{center}
			
			
			\begin{tikzpicture}[xscale=1, yscale=1]
			
			% Gitternetzlinien
			%\draw[xstep=1cm, ystep=1cm, very thin, color=lightgray] (-4,4) grid (4,-2);
			% x-Achse
			\draw[->] (-4,0) -- (4,0) node[below,xshift=-1em] {Re}; 
			
			% y-Achse 
			\draw[->] (0,-2) -- (0,4) node[left,yshift=-1em] {Im};%node[above left]
			\draw[color=blue] (1.5,0.8)  node[below] {$ \alpha =\frac{\pi}{4} $}; 
			\draw[color=green, ->] (0.6,0) arc (0:180:0.6);
			\draw[color=blue, ->] (1,0) arc (0:45:1);
			\draw[color=green] (-1.2,0.8)  node[below] {$ \beta =\pi $}; 
			\draw[thick,->,color=blue] plot[samples=200, domain=0:2] 
			(\x,{\x}) node[below, xshift=1em]{$\overrightarrow{a} $};
			\end{tikzpicture}
		\end{center}
		Nach Skizze folgt: $ \alpha=\frac{\pi}{4} $
		
	\end{quote}
	Nun kann man die Gleichung lösen mit:\\
	\begin{align*}
	z_k&=\sqrt[n]{|a|}\cdot\mathrm{e}^{\mathrm{i}\frac{\alpha+2k\pi}{n}}\;\;\text{für }\;\;k=0,1,2,\ldots,n-1\\
	z^k&=\sqrt[4]{16}\cdot\mathrm{e}^{\mathrm{i}\cdot \frac{\frac{\pi}{4}+2k\pi}{4}}\\
	z^0&=2\cdot \mathrm{e}^{\mathrm{i}\cdot\frac{\pi}{16}}\\
	z^1&=2\cdot \mathrm{e}^{\mathrm{i}\cdot\frac{\pi}{16}+\frac{\pi}{2}}\\
	z^2&=2\cdot \mathrm{e}^{\mathrm{i}\cdot\frac{\pi}{16}+\pi}\\
	z^3&=2\cdot \mathrm{e}^{\mathrm{i}\cdot\frac{\pi}{16}+\frac{3\pi}{2}}\\
	\end{align*}
	Daraus folgt:\\
	\begin{align*}
	\mathbb{A}=\{z_0,z_1,z_2,z_3\}
	\end{align*}
	Skizze von der Menge $ \mathbb{A} $:\\
	\begin{center}
		
		
		\begin{tikzpicture}[xscale=1, yscale=1]
		
		% Gitternetzlinien
		\draw[xstep=1cm, ystep=1cm, very thin, color=lightgray] (-3,3) grid (3,-3);
		% x-Achse
		\draw[->] (-3,0) -- (3,0) node[below,xshift=-1em] {Re}; 
		%Zahlen auf x-Achse
		\foreach \x in {-1,...,-1}
		\draw[shift={(\x,0)},color=black] (0pt,2pt) -- (0pt,-2pt) node[below] {\footnotesize $\x$};
		\foreach \x in {1,...,2}
		\draw[shift={(\x,0)},color=black] (0pt,2pt) -- (0pt,-2pt) node[below] {\footnotesize $\x$};
		% y-Achse 
		\draw[->] (0,-3) -- (0,3) node[left,yshift=-1em] {Im};%node[above left]
		%Zahlen auf y-Achse
		\foreach \y in {-1,...,-1}
		\draw[shift={(0,\y)},color=black] (2pt,0pt) -- (-2pt,0pt) node[left] {\footnotesize $\y$};
		\foreach \y in {1,...,1}
		\draw[shift={(0,\y)},color=black] (2pt,0pt) -- (-2pt,0pt) node[left] {\footnotesize $\y$};
		%Ursprung
		\draw[color=blue] (2,0) arc (0:360:2);
		\draw[color=black] (0pt,-10pt) node[right] {\footnotesize $0$};
		\draw[color=red] (0,0)--(1.96,0.39);
		\filldraw[black] (1.96,0.39) circle (2pt) node[anchor=west] {$ z_0 $};\filldraw[black] (1.96,0.39) circle (2pt) node[anchor=west] {$ z_0 $};
		\filldraw[black] (-1.96,-0.39) circle (2pt) node[anchor=east] {$ z_2 $};
		\filldraw[black] (-0.39,1.96) circle (2pt) node[anchor=east] {$ z_1 $};
		\filldraw[black] (0.39,-1.96) circle (2pt) node[anchor=east] {$ z_3 $};
		\draw[red] (1.39,0.7)  node[anchor=east] {$\alpha=\frac{\pi}{16} $};
		\end{tikzpicture}
	\end{center}
	Die Menge: $ \mathbb{A}=\{z_0,z_1,z_2,z_3\} $ ist auf dem Kreis mit Radius: $ r=2 $.\\
\end{quote}
\underline{Bestimmung der Menge: $ \mathbb{B} $}\\
\begin{quote}
	\begin{align*}
	\mathbb{B}=\{ z\in \mathbb{C}\;|\;  |z| \leq 2 \}
	\end{align*}
	Es gilt:\\
	\begin{align*}
	\mathbb{B}=\{x+\mathrm{i}y\in \mathbb{C}\;|\;\sqrt{x^2+y^2}\leq 2\}
	\end{align*}
	$ \sqrt{x^2+y^2} $ ist eine Kreisscheibe mit Radius: $ r=2 $ um $ (0,0) $.\\
	\begin{center}
		
		\begin{tikzpicture}[xscale=1, yscale=1]
		
		% Gitternetzlinien
		\draw[xstep=1cm, ystep=1cm, very thin, color=lightgray] (-3,3) grid (3,-3);
		\filldraw[color=green!10] (2,0) arc (0:360:2);
		\draw[color=green](-1,1)node[below]{\footnotesize $\mathbb{B}$};
		% x-Achse
		\draw[->] (-3,0) -- (3,0) node[below,xshift=-1em] {Re}; 
		%Zahlen auf x-Achse
		\foreach \x in {-1,...,-1}
		\draw[shift={(\x,0)},color=black] (0pt,2pt) -- (0pt,-2pt) ;
		\foreach \x in {1,...,2}
		\draw[shift={(\x,0)},color=black] (0pt,2pt) -- (0pt,-2pt) node[below] {\footnotesize $\x$};
		% y-Achse 
		\draw[->] (0,-3) -- (0,3) node[left,yshift=-1em] {Im};%node[above left]
		%Zahlen auf y-Achse
		\foreach \y in {-1,...,-1}
		\draw[shift={(0,\y)},color=black] (2pt,0pt) -- (-2pt,0pt) node[left] {\footnotesize $\y$};
		\foreach \y in {1,...,1}
		\draw[shift={(0,\y)},color=black] (2pt,0pt) -- (-2pt,0pt) node[left] {\footnotesize $\y$};
		%Ursprung
		%\filldraw[color=green] (2,0) arc (0:360:2);
		
		\draw[color=black] (0pt,-10pt) node[right] {\footnotesize $0$};
		%\draw[color=red] (0,0)--(1.96,0.39);
		\filldraw[black] (1.96,0.39) circle (2pt) node[anchor=west] {$ z_0 $};\filldraw[black] (1.96,0.39) circle (2pt) node[anchor=west] {$ z_0 $};
		\filldraw[black] (-1.96,-0.39) circle (2pt) node[anchor=east] {$ z_2 $};
		\filldraw[black] (-0.39,1.96) circle (2pt) node[anchor=east] {$ z_1 $};
		\filldraw[black] (0.39,-1.96) circle (2pt) node[anchor=east] {$ z_3 $};
		%\draw[red] (1.39,0.7)  node[anchor=east] {$\alpha=\frac{\pi}{16} $};
		
		\end{tikzpicture}
	\end{center}
\end{quote}
\underline{Bestimmung von $ \mathbb{A}\cap\mathbb{B} $}\\
\begin{quote}
	Es gilt:\\
	\begin{align*}
	\mathbb{A}\subset\mathbb{B}\Rightarrow\mathbb{A}\cap\mathbb{B}=\mathbb{A}
	\end{align*}
\end{quote}
\subsection{Aufgabe 2 - Grenzwert einer Folge}
\href{https://youtu.be/360OQSMzKcc}{Video 03}\\
Bestimme den Grenzwert folgender Folge:\\
\begin{align*}
x_n:=\frac{\pi\cdot n^3+\cos(n^2)}{4\cdot n^3-\sin(\frac{1}{y})}
\end{align*}
Betrachte:\\
\begin{align*}
x_n&=\frac{\pi\cdot n^3+\cos(n^2)}{4\cdot n^3-\sin(\frac{1}{y})}\\
&=\frac{n^3(\pi+\frac{\cos(n^2)}{n^3})}{n^3(4-\frac{\sin(\frac{1}{n+1})}{n^3})}\\
\lim\limits_ {n\to\infty} x_n&=\frac{\pi+\overbrace {\frac{\cos(n^2)}{n^3}}^{\rightarrow 0}}{4-\underbrace{\frac{\sin(\frac{1}{n+1})}{n^3}}_{\rightarrow 0}}=\frac{\pi}{4}
\end{align*}
\newpage
\noindent
\subsection{Aufgabe 3 - Grenzwert einer Folge}
\href{https://youtu.be/ICQcFrINxAQ}{Video 04}\\
Bestimme den Grenzwert folgender Folge:\\
\begin{align*}
x_n:=\sqrt{2n}\big(\sqrt{n+1}-\sqrt{n}\big)
\end{align*}
Betrachte:\\
\begin{align*}
x_n&=\sqrt{2n}\big(\sqrt{n+1}-\sqrt{n}\big)\\
&=\sqrt{2n}\frac{\big(\sqrt{n+1}-\sqrt{n}\big)\big(\sqrt{n+1}+\sqrt{n}\big)}{\big(\sqrt{n+1}+\sqrt{n}\big)}\\
&=\sqrt{2n}\frac{n+1-n}{\big(\sqrt{n+1}+\sqrt{n}\big)}\\
&=\sqrt{2}\sqrt{n}\frac{1}{\big(\sqrt{n+1}+\sqrt{n}\big)}\\
&=\sqrt{2}\frac{\sqrt{n}}{\big(\sqrt{n+1}+\sqrt{n}\big)}\\
&=\sqrt{2}\frac{\sqrt{n}}{\sqrt{n}\bigg(\frac{\sqrt{n+1}}{\sqrt{n}}+1\bigg)}\\
&=\sqrt{2}\frac{\sqrt{n}}{\sqrt{n}\bigg(\sqrt{\frac{n+1}{n}}+1\bigg)}\\
&=\sqrt{2}\frac{\sqrt{n}}{\sqrt{n}\bigg(\sqrt{\frac{n(\frac{1}{n}+1)}{n}}+1\bigg)}\\
&=\sqrt{2}\frac{1}{\bigg(\sqrt{\frac{1}{n}+1)}+1\bigg)}\\
&=\sqrt{2}\frac{1}{2}=\frac{\sqrt{2}}{2}
\end{align*}
\subsection{Aufgabe 4 - Induktion}
\href{https://youtu.be/nJx5VsSVHaQ?t=5}{Video 05}\\
Beweise mit vollständiger Induktion:\\
\begin{align*}
\sum\limits_{k=1}^{n}k(k+1)=\frac{n(n+1)(n+2)}{3}\;\;\text{ für}\;\;n\in\mathbb{N}
\end{align*}
$ \bullet $\;\;(I.A.):\\
\begin{quote}
	Für $ n=0 $ gilt:\\
	\begin{align*}
	\sum\limits_{k=0}^{0}k(k+1)=0=\frac{0(0+1)(0+2)}{3}\checkmark
	\end{align*}
\end{quote}
$ \bullet $\;\;(I.V.):\\
\begin{align*}
\sum\limits_{k=1}^{n}k(k+1)=\frac{n(n+1)(n+2)}{3}\;\;\text{gilt für \underline{ein} }\;\;n\in\mathbb{N}
\end{align*}
$ \bullet $\;\;(I.S.):\\
\begin{quote}
	Z.z.:
	\begin{align*}
	\sum\limits_{k=1}^{n+1}k(k+1)=\frac{(n+1)(n+2)(n+3)}{3}
	\end{align*}
	Es gilt:\\
	\begin{align*}
	\sum\limits_{k=1}^{n+1}k(k+1)=	\sum\limits_{k=1}^{n}k(k+1)+(n+1)(n+2)&\overset{(I.V.)}{=}\frac{n(n+1)(n+2)}{3}+(n+1)(n+2)\\
	&\;\;=\frac{n(n+1)(n+2)+3(n+1)(n+2)}{3}\\\sum\limits_{k=1}^{n+1}k(k+1)&\;\;=\frac{(n+1)(n+2)(n+3)}{3}\boxed{}\\
	\end{align*}
\end{quote}

\subsection{Aufgabe 5 - Differenzierbarkeit}
\href{https://youtu.be/8EBDpt2f5S0}{Video 06}\\
Für welches $ \alpha=\in\mathbb{R} $ ist $ f:\mathbb{R}\mapsto \mathbb{R}$ mit\\

\begin{align*}
f(x) =     \left\{ \begin{array}{rcl}
x\cdot\mathrm{e}^{\sinh(x)} +\alpha(x+1)& \mbox{für} & x\geq0\\
-1&\mbox{für} & x<0
\end{array}\right.
\end{align*}
in $ x_0=0 $ differenzierbar.\\
\\
Der Grenzwert des Differntialquotienten muss links von $ x_0=0 $ und rechts von $ x_0=0 $ gleich sein.\\
Es muss gelten:\\
\begin{align*}
\lim\limits_ {\substack{x\to x_0\\x>0}}\frac{f(x-f(x_0))}{x-x_0}\overset{!}{=}\lim\limits_ {\substack{x\to x_0\\x<0}}\frac{f(x-f(x_0))}{x-x_0}
\end{align*}
$\bullet$ \;\;Für $ x<0 $ gilt:\\
\begin{align*}
\lim\limits_ {\substack{x\to x_0\\x<0}}\frac{f(x)-f(x_0)}{x-x_0}=\lim\limits_ {\substack{x\to 0\\x<0}}\frac{-1-\big(0\cdot \mathrm{e}^{\sinh(0)})+\alpha(0+1)\big)}{x}=\frac{-1-\alpha}{x}
\end{align*}
Für $ x\to x_0=0 $ geht der Bruch: $ \frac{1-\alpha}{x} $ gegen unendlich. Damit ein Grenzwert entsteht muß der Zähler Null werden.\\ Mit $ \alpha=-1 $ wird der Grenzert Null.\\
Daraus folgt dann mit $ \alpha=-1 $:\\
\begin{align*}
\lim\limits_ {\substack{x\to x0\\x<0}}\frac{f(x)-f(x_0)}{x-x_0}=0
\end{align*}
$\bullet$ \;\;Für $ x>0 $ gilt:\\
\begin{align*}
\lim\limits_ {\substack{x\to x_0\\x>0}}\frac{f(x)-f(x_0)}{x-x_0}&=\lim\limits_ {\substack{x\to 0\\x>0}}\frac{x\cdot\mathrm{e}^{\sinh(x)}+\alpha(x+1)-(0\cdot\mathrm{e}^{\sinh(0)}+\alpha(0+1))}{x}\\
&=\lim\limits_ {\substack{x\to 0\\x>0}}\frac{x\cdot\mathrm{e}^{\sinh(x)}+\alpha(x+1)-\alpha}{x}\\
&=\lim\limits_ {\substack{x\to 0\\x>0}}\frac{x\cdot\mathrm{e}^{\sinh(x)}+\alpha x+\alpha-\alpha}{x}\\
&=\lim\limits_ {\substack{x\to 0\\x>0}}\mathrm{e}^{\sinh(x)}+\alpha\\
&=\lim\limits_ {\substack{x\to 0\\x>0}}=1+\alpha
\end{align*}
Für $ \alpha=-1 $ gilt dann:\\
\begin{align*}
\lim\limits_ {\substack{x\to 0\\x>0}}\frac{f(x)-f(x_0)}{x-x_0}=0
\end{align*}
Da nun gilt:\\
\begin{align*}
\lim\limits_ {\substack{x\to x_0\\x>0}}\frac{f(x-f(x_0))}{x-x_0}&=\lim\limits_ {\substack{x\to x_0\\x<0}}\frac{f(x-f(x_0))}{x-x_0}\\
0&=0
\end{align*}
Da die Grenzwerte des Differentialquotienten, der ja die Ableitung von $ f(x) $ entspricht, gleich sind ist die Funktion an der Stelle $ x_0=0 $ stetig und differnzierbar.\\
Ebenso muss an der Stelle $ x=x_0=0 $ gelten:\\
\begin{align*}
f(0)=-1=0\mathrm{e}^{\sinh(0)}+(-1)(0+1)=-1\\
\Rightarrow\text{keine Sprungfunktion! }
\end{align*}
\newpage
\noindent
\subsection{Aufgabe 6 - Induktion}
\href{https://youtu.be/LnNUw7oT4O4}{Video 07}\\
Betrachte die Funktion: \\
\begin{align*}
f(x) =\cosh(\sqrt{2}x)-\sinh(\sqrt{2}x)
\end{align*}
Zeige mittels vollständiger Induktion:\\
\begin{align*}
f^{(n)}=\big(-\sqrt{2}\big)^n\mathrm{e}^{-\sqrt{2}x}\;\;\text{für }\;\;n\in\mathbb{N}
\end{align*}
Betrachte:\\
\begin{align*}
f(x) =\cosh(\sqrt{2}x)-\sinh(\sqrt{2}x)
\end{align*}
\begin{quote}
	Es gilt:\\
	\begin{align*}
	\cosh(x)&=\frac{\mathrm{e}^x+\mathrm{e}^{-x}}{2}\\
	\sinh(x)&=\frac{\mathrm{e}^x-\mathrm{e}^{-x}}{2}\\
	\end{align*}
	In $ f(x) $ eingesetzt ergibt dann:\\
	\begin{align*}
	f(x)&=\frac{\mathrm{e}^{\sqrt{2}x}+\mathrm{e}^{-\sqrt{2}x}}{2}-\frac{\mathrm{e}^{\sqrt{2}x}-\mathrm{e}^{-\sqrt{2}x}}{2}=\frac{2\mathrm{e}^{-\sqrt{2}x}}{2}	=\mathrm{e}^{-\sqrt{2}x}
	\end{align*}
	
\end{quote}
Beweis  mittels vollständiger Induktion:\\
$ \bullet $\;\;(I.A):\\
\begin{quote}
	Für $ n=0 $ gilt:\\
	\begin{align*}
	f^{(n)}&=\big(-\sqrt{2}\big)^n\mathrm{e}^{-\sqrt{2}x}\\
	f^{(0)}&=\big(-\sqrt{2}\big)^0\mathrm{e}^{-\sqrt{2}x}=\mathrm{e}^{-\sqrt{2}x}=f(x)\checkmark
	\end{align*}
\end{quote}
$ \bullet $\;\;(I.V):\\
\begin{quote}
	Für ein $ n\in\mathbb{N} $ gilt:\\
	\begin{align*}
	f^{(n)}(x)=\big(-\sqrt{2}\big)^n\mathrm{e}^{-\sqrt{2}x}
	\end{align*}
\end{quote}
\newpage
\noindent
$ \bullet $\;\;(I.S):\\
\begin{quote}
	Z.z.:\\
	\begin{align*}
	f^{n+1}(x)=\big(-\sqrt{2}\big)^{n+1}\mathrm{e}^{-\sqrt{2}x}
	\end{align*}
	Es gilt:\\
	\begin{align*}
	f^{n+1}(x)=\big(f^{n}(x)\big)'&\overset{(I.V.)}{=}\bigg(\big(-\sqrt{2}\big)^n\mathrm{e}^{-\sqrt{2}x}\bigg)'\\
	&\;\;=\big(-\sqrt{2}\big)^n\mathrm{e}^{-\sqrt{2}x}\cdot \big(-\sqrt{2}\big)\\
	&\;\;=\big(-\sqrt{2}\big)^{n+1}\mathrm{e}^{-\sqrt{2}x}\boxed{}
	\end{align*}
\end{quote}
\subsection{Aufgabe 7 - Folgen }
\underline{Teil I}\\
\href{https://youtu.be/WVkBwJisfwI}{Video 08 Teil 1}\\
Entscheide über Korrektheit der folgender Aussage und begründe deine Antwort.\\
Sei $ (x_n)_{n\in\mathbb{N}} $ eine Nullfolge und $ (y_n)_{n\in\mathbb{N}} $ beschränkt, dann konvergiert die Folge: $ (x_n\cdot y_n)_{n\in\mathbb{N}} $.\\
\\
Gegeben:\\
\begin{align*}
&(x_n)_{n\in\mathbb{N}}\;\;\text{mit }\;\;\lim\limits_ {x\to\infty}x_n=0\\
&(y_n)_{n\in\mathbb{N}}\;\;\text{beschränkt }
\end{align*}
Da $ (y_n)_{n\in\mathbb{N}} $ beschränkt ist, existiert eine Schranke:\\
\begin{align*}
\mathrm{S}\in \mathbb{R}\;\;\text{mit }\;\;y_n\leq \mathrm{S}\;\;\text{für alle }\;\;n\in\mathbb{N}
\end{align*}
Genauso gilt:\\
\begin{align*}
|y_n|\leq \mathrm{S}\;\;\text{für alle }\;\;n\in\mathbb{N}
\end{align*}
$ y_n $ bewegt sich zwischen zwei Schranken.\\
Betrachte nun:\\
\begin{align*}
&\lim\limits_ {n\to\infty}|x_n\cdot y_n|=\lim\limits_ {n\to\infty}|x_n\cdot \mathrm{S}|=|\lim\limits_ {n\to\infty} x_n\cdot \mathrm{S}|=|0\cdot \mathrm{S}|=0\\
&\Rightarrow \lim\limits_ {n\to\infty}x_n\cdot y_n=0
\end{align*}
Also ist die Aussage korrekt.
\newpage
\noindent
\underline{Teil II}\\
\href{https://youtu.be/gXJcu4IlP0k}{Video 09 Teil 2}\\
Entscheide über Korrektheit der folgender Aussage und begründe deine Antwort.\\
Sei $ (x_n)_{n\in\mathbb{N}} $ eine konvergente  und $ (y_n)_{n\in\mathbb{N}} $ beschränkt, dann konvergiert die Folge: $ (x_n\cdot y_n)_{n\in\mathbb{N}} $.\\
\\
Wenn die Aussage nicht korrekt ist, dann sucht man sich ein Gegenbeispiel.\\
Betrachte:\\
\begin{align*}
&(x_n)_{n\in\mathbb{N}}=(1)_{n\in\mathbb{N}}\;\;\text{mit } \lim\limits_ {n\to\infty}x_n=1 \Rightarrow\;\;\text{ Folge ist konvergent}\\
&(y_n)_{n\in\mathbb{N}}=(-1)^n_{n\in\mathbb{N}}\;\;\text{ beschränkt, wegen:}-1\leq(-1)^n\leq 1\;\;\text{ für alle}\;\;n\in\mathbb{N}
\end{align*}
Aber:\\
\begin{align*}
(x_n\cdot y_n)_{n\in\mathbb{N}}=((-1))^n_{n\in\mathbb{N}}\;\;\text{Folge ist divergent! }\Rightarrow \text{Aussage falsch! }
\end{align*}
\subsection{Aufgabe 8 - Integral}
\href{https://youtu.be/VEksC5lXo9o}{Video 10}\\
Gerechne das folgende Integral:\\
\begin{align*}
\int\limits_{1}^{2} \frac{4\ln(x^3)}{3}\mathrm{d}x
\end{align*}
Das Integral vereinfachen:\\
\begin{align*}
\int\limits_{1}^{2} \frac{4\ln(x^3)}{3}\mathrm{d}x=\int\limits_{1}^{2} \frac{4}{3}\ln(x^3)\mathrm{d}x=\int\limits_{1}^{2} \frac{4*3}{3}\ln(x)\mathrm{d}x=4\int\limits_{1}^{2} \ln(x)\mathrm{d}x
\end{align*}
Es gilt mit Hilfe der partiellen Integration.\\
\begin{align*}
\int\limits_{a}^{b} u(x)\cdot v'(x)=\big [u(x)\cdot v(x)\big]-\int\limits_{a}^{b} u'(x)\cdot v(x)
\end{align*}
Für das Integral: $ \int\limits_{a}^{b}1\cdot \ln(x) $ gilt dann:\\
\begin{align*}
&u(x)\;=\ln(x)\Rightarrow u'(x)=\frac{1}{x}\\
&v'(x)=1\;\;\;\;\;\;\Rightarrow v(x)\;=x
\end{align*}
Dies ergibt dann durch Einsetzen:\\
\begin{align*}
\int\limits_{a}^{b}1\cdot \ln \mathrm{d}x&=  \big[\ln(x)\cdot x\big]^b_a-\int\limits_{a}^{b} \frac{1}{x}\cdot x \mathrm{d}x\\
&=\big[\ln(x)\cdot x\big]^b_a-\int\limits_{a}^{b}1\cdot \mathrm{d}x
\end{align*}
Für das Integral $4\int\limits_{1}^{2} \ln(x) \mathrm{d} x$ gilt:\\
\begin{align*}
\int\limits_{1}^{2} \frac{4\ln(x^3)}{3}\mathrm{d}x&=4\bigg [  \big[\ln(x)\cdot x\big]^2_0-\int\limits_{1}^{2} 1\cdot  \mathrm{d}x\bigg]\\
&=4\bigg [ \big[\ln(x)\cdot x\big]^2_0-\big[x\big]^2_1\bigg]\\
&=4(2\ln(2)-1)\\
&=8\ln(2)-4
\end{align*}
\subsection{Aufgabe 9 - Grenzwert}
\href{https://youtu.be/z0jsH4rjKZM}{Video 11}\\
Berechne den Grenzwert von:\\
\begin{align*}
\lim\limits_ {x\to 0}\frac{\cosh(x)-1}{\sinh(x)}
\end{align*}
Für $ x \to 0 $ gilt:
\begin{align*}
\lim\limits_ {x\to 0}\frac{\cosh(x)-1}{\sinh(x)}=\frac{0}{0}
\end{align*}
Mit Regel von L'Hospital folgt:\\
\begin{align*}
\lim\limits_ {x\to 0}\frac{\cosh(x)-1}{\sinh(x)}=\lim\limits_ {x\to 0}\frac{\big(\cosh(x)-1\big)'}{\big(\sinh(x)\big)'}=\lim\limits_ {x\to 0}\frac{\sinh(x)}{\cosh(x)}
\end{align*}
Es gilt:\\
\begin{align*}
\sinh(x)=\frac{e^x-e^{-x}}{2}\\
\cosh(x)=\frac{e^x+e^{-x}}{2}
\end{align*}
Für den Grenzwert gilt dann:\\
\begin{align*}
\lim\limits_ {x\to 0}\frac{\sinh(x)}{\cosh(x)}=\lim\limits_ {x\to 0}\bigg(\frac{\mathrm{e}^x-\mathrm{e}^{-x}}{\mathrm{e}^x+\mathrm{e}^{-x}}\bigg)=\frac{1-1}{1+1}=\frac{0}{2}=0
\end{align*}
\newpage
\noindent
\subsection{Aufgabe 10 - Integral mt Substitution}
\href{https://youtu.be/lU3LyEOpWQE}{Video 12}\\
Berechne das folgende Integral:\\
\begin{align*}
\int\limits_{\ln(5)}^{\ln(6)} \frac{2\mathrm{e}^x+\frac{1}{2}}{\mathrm{e}^x-4}
\end{align*}
Substitution mit $ x=\ln(t+4) \Rightarrow $Nenner wird zu $ t $:\\
\begin{align*}
x&=\ln(t+4)\\
\frac{\mathrm{d}x}{\mathrm{d}t}&=\frac{\ln(t+4)}{\mathrm{d}t}=\frac{1}{t+4}\\
\Leftrightarrow\\
\mathrm{d}x&=\frac{1}{t+4}\mathrm{d}t
\end{align*}
Für die Grenzen gilt dann:\\
\begin{align*}
\ln(5)&=\ln(t+4)\Rightarrow t=1\\
\ln(6)&=\ln(t+4)\Rightarrow t=2
\end{align*}
Nun gilt:\\
\begin{align*}
\int\limits_{1}^{2} \frac{2\mathrm{e}^{\ln(t+4)}+\frac{1}{2}}{\mathrm{e}^{\ln(t+4)}-4}\frac{1}{t+4}\mathrm{d}t&=\int\limits_{1}^{2}\frac{2(t+4)+\frac{1}{2}}{t} \frac{1}{t+4}\mathrm{d}t=\int\limits_{1}^{2} \frac{2(t+4)+\frac{1}{2}}{t(t+4)}\mathrm{d}t=\int\limits_{1}^{2} \frac{2(t+4)}{t(t+4)}+\frac{\frac{1}{2}}{t(t+4)}\mathrm{d}t\\&=2\int\limits_{1}^{2} \frac{1}{t}\mathrm{d}t+\int\limits_{1}^{2} \frac{\frac{1}{2}}{t(t+4)}\mathrm{d}t
\end{align*}
Mit Partialbruchzerlegung und Zuhaltemethode erhält man für den Bruch: $  \frac{\frac{1}{2}}{t(t+4)} $\\
\begin{align*}
\frac{\frac{1}{2}}{t(t+4)} =\frac{\mathrm{A}}{t}+\frac{\mathrm{B}}{t+4}
\end{align*}
Mit den Nullstellen  0 und -4 erhält man A und B durch die Zuhaltemethode.\\
\begin{align*}
\mathrm{A}&=\frac{1}{8}\\
\mathrm{B}&=-\frac{1}{8}\\
\end{align*}
\begin{align*}
\frac{\frac{1}{2}}{t(t+4)} =\frac{\mathrm{\frac{1}{8}}}{t}+\frac{\mathrm{-\frac{1}{8}}}{t+4}
\end{align*}
Nun gilt:\\
\begin{align*}
\int\limits_{1}^{2} \frac{2\mathrm{e}^{\ln(t+4)}+\frac{1}{2}}{\mathrm{e}^{\ln(t+4)}-4}\frac{1}{t+4}\mathrm{d}t&=2\int\limits_{1}^{2} \frac{1}{t}\mathrm{d}t+\int\limits_{1}^{2} \frac{\frac{1}{2}}{t(t+4)}\mathrm{d}t\\
&=2\int\limits_{1}^{2}\frac{1}{t} \mathrm{d}t+\frac{1}{8}\int\limits_{1}^{2} \frac{1}{t}\mathrm{d}t-\frac{1}{8}\int\limits_{1}^{2} \frac{1}{t+4}\mathrm{d}t\\
&=2\bigg[\ln(t)\bigg]^2_1+\frac{1}{8}\bigg[\ln(t)\bigg]^2_1-\frac{1}{8}\bigg[\ln(t+4)\bigg]^2_1\\
&=2\ln(2)+\frac{1}{8}\ln(2)-\frac{1}{8}\big(\ln(6)-\ln(5)\big)\\
&=\frac{17}{8}\ln(2)-\frac{\ln(6)-\ln(5)}{8}
\end{align*}
\subsection{Aufgabe 11 - Nullstellen bestimmen}
\href{https://youtu.be/LFx7ZRhGa_o}{Video 13}\\
Gegeben sei:\\
\begin{align*}
f:\;\;]0,2\pi]\longmapsto\mathbb{R}\;\;\text{ mit}\;\;f(x)=\sin(\cos(x))
\end{align*}
Bestimme alle Nullstellen von $ f(x) $.\\
\\
\\Es gilt:\\
\begin{align*}
\sin(x)=0\;\; \text{für }\;\; k\cdot \pi\;\;\text{mit }\;\;k\in \mathbb{Z}\\
\end{align*}
Ebenso gilt :\\
\begin{align*}
\sin(\underbrace {\cos(x)}_{|\cos(x)|<1})
\end{align*}
Daraus folgt , wenn der $ \cos(x) =0$ hat der $ \sin(\cos(x)) $ eine Nullstelle.\\
Für den $ \cos(x) $ gilt ja:\\
\begin{align*}
\cos(x)=0 \;\;\text{ für}\;\;\frac{\pi}{2}+k\cdot\pi\;\;\text{ }\;\;k\in \mathbb{Z}
\end{align*}
Wegen $f:\;\;]0,2\pi]\longmapsto\mathbb{R}  $ hat $ f(x) $ zwei Nullstellen bei:\\
\begin{align*}
&x_0=\frac{\pi}{2}<2\pi\\
&x_1=\frac{\pi}{2}+\pi<2\pi
\end{align*}
\begin{center}
	
	
	\begin{tikzpicture}[xscale=1, yscale=1]
	
	% Gitternetzlinien
	\draw[xstep=1cm, ystep=1cm, very thin, color=lightgray] (-1,3) grid (8,-2);
	% x-Achse
	\draw[->] (-1,0) -- (8,0) node[below,xshift=-1em] {$x$}; 
	%Zahlen auf x-Achse
	\foreach \x in {-1,...,-1}
	\draw[shift={(\x,0)},color=black] (0pt,2pt) -- (0pt,-2pt) node[below] {\footnotesize $\x$};
	\foreach \x in {1,...,7}
	\draw[shift={(\x,0)},color=black] (0pt,2pt) -- (0pt,-2pt) node[below] {\footnotesize $\x$};
	% y-Achse 
	\draw[->] (0,-2) -- (0,3) node[left,yshift=-1em] {$f(x)$};%node[above left]
	%Zahlen auf y-Achse
	\foreach \y in {-1,...,-1}
	\draw[shift={(0,\y)},color=black] (2pt,0pt) -- (-2pt,0pt) node[left] {\footnotesize $\y$};
	\foreach \y in {1,...,2}
	\draw[shift={(0,\y)},color=black] (2pt,0pt) -- (-2pt,0pt) node[left] {\footnotesize $\y$};
	%Ursprung
	\draw[color=black] (0pt,-10pt) node[right] {\footnotesize $0$};
	
	\draw[thick,color=red] plot[samples=200, domain=0:6.28] 
	(\x,{sin(cos(\x r)r}) node[ xshift=-8em]{$f(x)=\sin(\cos(x))$};
	
	
	\end{tikzpicture}
\end{center}
\subsection{Aufgabe 12 - Extrema bestimmen}
\href{https://youtu.be/cN-hCDhSnWw}{Video 14}\\
Gegeben sei:\\
\begin{align*}
f:\;\;]0,2\pi]\longmapsto\mathbb{R}\;\;\text{ mit}\;\;f(x)=\sin(\cos(x))
\end{align*}
Bestimme alle Extrema von $ f(x) $.\\
\\
$ \bullet $\;\;Extrema: \glqq innerer Punkte \grqq\\
\begin{quote}
	
	Bestimmung von $ f'(x) $:\\
	\begin{align*}
	f(x)&=\sin(\cos(x))\\
	f'(x)&=-\cos(\cos(x))\cdot \sin(x)
	\end{align*}
	Nullstellen von $ f'(x) $ bestimmen:\\
	$ f'(x) $ ist Null, wenn $ -\cos(\cos(x)) $ oder $ \sin(x) $ Null wird.\\
	Für $ -\cos(\cos(x)) $ gilt:
	\begin{align*}
	-\cos(\underbrace {\cos(x)}_{|\cos(x)|<1})
	\end{align*}
	Für $ \cos(x)=0 $ gilt:\\
	\begin{align*}
	&	\cos(x)=0\;\; \text{für }\;\;\frac{\pi}{2}+ n\cdot \pi\;\;\text{mit }\;\;n\in \mathbb{Z}\\
	&	\Rightarrow |\cos(x)|<\big|\frac{\pi}{2}+ n\cdot \pi\big|\;\;\text{mit }\;\;n\in \mathbb{Z}\\
	&\Rightarrow-\cos(\cos(x))\neq 0\;\;\text{für alle }\;\;x\in\mathbb{R}\\
	\end{align*}
	Für  $ \sin(x) =0$ gilt:\\
	\begin{align*}
	\sin(x)=0 \;\;\text{für }\;\;n\cdot\pi\;\;\text{mit }\;\;n\in \mathbb{Z}
	\end{align*}
	Wegen: $f:\;\;]0,2\pi]\longmapsto\mathbb{R}  $ existiert eine innere potenzielle Extremstelle: $ x_1=\pi $
\end{quote}
Bestimmung von $ f''(x) $:\\
\begin{align*}
f'(x)&=-\cos(\cos(x))\cdot \sin(x)\\
f''(x)&=-\sin(\cos(x))\cdot\sin(x)\cdot\sin(x)-\cos(x)\cdot\cos(\cos(x))\\
&=-\sin(\cos(x))\cdot\sin^2(x)-\cos(x)\cdot\cos(\cos(x))\\
\end{align*}
Bei der Extremstelle $ x_1=\pi $ gilt :\\
\begin{align*}
&f''(x)=-\sin(\cos(\pi))\cdot\underbrace {\sin^2(\pi)}_{=0}-\underbrace {\cos(\pi)}_{=-1}\cdot\underbrace {\cos(\cos(\pi))}_{=\cos(-1)>0}\\
&	f''(x)=0+1\cdot\cos(-1)>0\Rightarrow x_1=\pi\;\; \text{ist ein Minimum }
\end{align*}
Randpunkt und Grenzwertbetrachtung:\\
$ \bullet\;\;f(2\pi) =\sin(\cos(2\pi))=\sin(1)\Rightarrow \text{globales Maximum }$\\
$ \bullet\;\;f(\pi)\;\;=\sin(\cos(\pi)) =\sin(-1)=-\sin(1)\Rightarrow \text{ globales Minimum}$\\
$ \bullet\;\;\lim\limits_ {\substack{x\to 0\\x>0}} f(x)=\lim\limits_ {\substack{x\to 0\\x>0}}\sin(\cos(x))=\sin(1)$

\begin{center}
	
	
	\begin{tikzpicture}[xscale=1, yscale=1]
	
	% Gitternetzlinien
	\draw[xstep=1cm, ystep=1cm, very thin, color=lightgray] (-1,3) grid (8,-2);
	% x-Achse
	\draw[->] (-1,0) -- (8,0) node[below,xshift=-1em] {$x$}; 
	%Zahlen auf x-Achse
	\foreach \x in {-1,...,-1}
	\draw[shift={(\x,0)},color=black] (0pt,2pt) -- (0pt,-2pt) node[below] {\footnotesize $\x$};
	\foreach \x in {1,...,7}
	\draw[shift={(\x,0)},color=black] (0pt,2pt) -- (0pt,-2pt) node[below] {\footnotesize $\x$};
	% y-Achse 
	\draw[->] (0,-2) -- (0,3) node[left,yshift=-1em] {$f(x)$};%node[above left]
	%Zahlen auf y-Achse
	\foreach \y in {-1,...,-1}
	\draw[shift={(0,\y)},color=black] (2pt,0pt) -- (-2pt,0pt) node[left] {\footnotesize $\y$};
	\foreach \y in {1,...,2}
	\draw[shift={(0,\y)},color=black] (2pt,0pt) -- (-2pt,0pt) node[left] {\footnotesize $\y$};
	%Ursprung
	\draw[color=black] (0pt,-10pt) node[right] {\footnotesize $0$};
	
	\draw[thick,color=red] plot[samples=200, domain=0:6.28] 
	(\x,{sin(cos(\x r)r}) node[ xshift=-8em]{$f(x)=\sin(\cos(x))$};
	
	
	\end{tikzpicture}
\end{center}
\subsection{Aufgabe 13 - Fourierpolynom}
\underline{Teil I}\\
\href{https://youtu.be/eVy5SC3mWBM}{Video 15 Teil 1}\\
Bestimme das $ n-k $ Fourierpolynom von:\\
\begin{align*}
f:\mathbb{R}\rightarrow\mathbb{R} \;\;\text{mit  }\;\; f(x)=\sin^2(x)\bigg(\cot^2(x)+\frac{1}{1-\cos^2(x)}\bigg)
\end{align*}
\\
Vereinfachung von $ f(x) $:\\
Es gilt:\\
\begin{align*}
\cot^2(x)&=\frac{\cos^2(x)}{\sin^2(x)}\\
1&=\cos^2(x)+\sin^2(x)\Leftrightarrow\sin^2(x)=1-\cos^2(x)\\
\cos(2x)&=\cos^2(x)-\sin^2(x)
\end{align*}
Eingesetzt ergibt sich dann:\\
\begin{align}
f(x)=\sin^2(x)\bigg(\frac{\cos^2(x)}{\sin^2(x)}+\frac{1}{\sin^2(x)}\bigg)=\cos^2(x)+1 \label{x1}
\end{align}
Ebenso gilt:\\
\begin{align}
\cos(2x)&=\cos^2(x)-\sin^2(x) \nonumber\\
\cos(2x)+1&=\cos^2(x)-\sin^2(x)+1\;\;|\scriptscriptstyle 1=\cos^2(x)+\sin^2(x)\nonumber\\
&=\cos^2(x)-\sin^2(x)+\sin^2(x)+\cos^2(x)\nonumber\\
&=2\cos^2(2x)\nonumber\\
\Leftrightarrow\nonumber\\
\cos^2(x)&=\frac{\cos(x)+1}{2}
\end{align}
(XV.2) in (XV.1) einsetzen:\\
\begin{align*}
f(x)&=\frac{\cos(2x)+1}{2}+1\\&=\frac{1}{2}\cos(2x)+\frac{1}{2}+1\\
f(x)&=\frac{3}{2}+\frac{1}{2}\cos(2x)
\end{align*}
Da $ f(x)=\frac{3}{2}+\frac{1}{2}\cos(2x) $ ein trigonometrisches Polynom ist, folgt:\\
\begin{align*}
\mathrm{Q}_n(x)=f(x)&=\frac{3}{2}+\frac{1}{2}\cos(2x)
\end{align*}
Bemerkung:\\
\begin{align*}
\mathrm{Q}_n(x)=\frac{a_0}{2}+\sum\limits_{k=1}^{n}a_k\cdot\cos(k\omega x)+b_k\cdot\sin(k\omega x)
\end{align*}
Bestimmung von $ \omega, a_k $ und $ b_k $ von $\mathrm{Q}_n(x)$ .\\
Für $ \omega $ gilt:\\
\begin{align*}
\omega =2\pi\cdot \mathrm{f}=\frac{2\pi}{\mathrm{T}}
\end{align*}
Der $ \cos(2x) $ ist $ \pi $-periodisch, daraus folgt:\\
\begin{align*}
&\Rightarrow \mathrm{T}=\pi\\
&\omega=\frac{2\pi}{\mathrm{T}}=\frac{2\pi}{\pi}=2
\end{align*}
Für $ a_0 $ gilt:\\
\begin{align*}
&\frac{a_0}{2}=\frac{3}{2}\\
&\Leftrightarrow\\
&a_0=3
\end{align*}
Für $ a_k $ gilt:\\ 
\begin{align*}
a_1=\frac{1}{2}\\
a_k=0\;\;\text{für }\;\; k>1
\end{align*}
Für $ b_k $ gilt:\\ 
\begin{align*}
b_k=0\;\;\text{für }\;\;k\geq 1
\end{align*}
Für $\mathrm{Q}_n(x)  $ gilt nun:\\
\begin{align*}
\mathrm{Q}_n(x)&=\frac{a_0}{2}+\sum\limits_{k=1}^{n}a_k\cdot\cos(k\omega x)+b_k\cdot\sin(k\omega x)\\
&=\frac{a_0}{2}+a_1\cos(k\omega x)+a_2\cos(k\omega x)\dots+b_1\sin(k\omega x)+b_2\cos(k\omega x)\dots\\
&=\frac{3}{2}+\frac{1}{2}\cos(1\cdot 2x)\\
\mathrm{Q}_n(x)=f(x)&=\frac{3}{2}+\frac{1}{2}\cos( 2x)\\
\end{align*}
\underline{Teil II}\\
\href{https://youtu.be/j2W_UAccyW8}{Video 16 Teil 2}\\
Bestimme  mit dem Fourierpolynom von:\\
\begin{align*}
f:\mathbb{R}\rightarrow\mathbb{R} \;\;\text{mit  }\;\; f(x)=\frac{3}{2}+\frac{1}{2}\cos(2x)
\end{align*}
das Integral:\\
\begin{align*}
\int\limits_{0}^{\pi} f(x)\cdot\sin(8x)\mathrm{d}x
\end{align*}
Es gilt für $ f(x) $:\\
\begin{align*}
b_k&=\frac{2}{\pi}\int\limits_{0}^{\pi} f(x)\cdot\sin(k\omega x)\mathrm{d}x\\
b_k&=0\;\;\text{ für}\;\;k\geq 0\\
\omega&=2\\
\Rightarrow&\\
b_k&=\frac{2}{\pi}\int\limits_{0}^{\pi} f(x)\cdot\sin(k\cdot2x)\mathrm{d}x\\
\end{align*}
Für $ b_k=4 $ gilt:\\
\begin{align*}
b_4&=\frac{2}{\pi}\int\limits_{0}^{\pi} f(x)\cdot\sin(4\cdot2x)\mathrm{d}x\;\;|\cdot\frac{\pi}{2}\\
&=\int\limits_{0}^{\pi} f(x)\cdot\sin(8x)\mathrm{d}x\
\end{align*}
Da $ b_4=0 $:\\
\begin{align*}
&b_4=\int\limits_{0}^{\pi} f(x)\cdot\sin(8x)\mathrm{d}x=0\\
&\Leftrightarrow\\
&\int\limits_{0}^{\pi} f(x)\cdot\sin(8x)\mathrm{d}x=0\\
\end{align*}
\subsection{Aufgabe 14 - Extrema bestimmen}
\href{https://youtu.be/9nOEDykwLso}{Video 17}\\
Bestimme alle Extrema von:\\
\begin{align*}
\mathrm{F}:\;\;]\genfrac{}{}{}{2}{1 }{2 },\infty[\;\;\longrightarrow \mathbb{R}\;\;\text{mit }\;\;\mathrm{F}(x)=\int\limits_{\frac{1}{2}}^{x} \ln(t)\mathrm{d}t
\end{align*}
$ \bullet $\;\;Extrema: \glqq innerer Punkte\grqq\\
\begin{quote}
	Bestimmung von $ \mathrm{F}'(x) $:
	\begin{align*}
	\mathrm{F}'(x)=\bigg(\int\limits_{\frac{1}{2}}^{x}\ln(t)\mathrm{d}t\bigg)'=\ln(x)
	\end{align*}
	Für die Nullstellen gilt :\\
	\begin{align*}
	\mathrm{F}'(x)=\ln(x)=0\;\;\text{für  }\;\;x_0=1
	\end{align*}
	Bestimmung von $ \mathrm{F}''(x) $:\\
	\begin{align*}
	\mathrm{F}''(x)=\bigg(\ln(x)\bigg)'=\frac{1}{x}
	\end{align*}
	Für  $ \mathrm{F}''(x)=\frac{1}{x}$ gilt:\\
	\begin{align*}
	&\frac{1}{x}\;\;\text{für }x\in [\genfrac{}{}{}{2}{1 }{2 },\infty]\;\;\text{gilt }\;\;\frac{1}{x}>0\\
	&\Rightarrow \mathrm{F}\;\;\text{ hat ein lokales Minimum bei:}\;\;x_0=1\\
	\end{align*}
\end{quote}
$ \bullet $\;\;Grenzwertbetrachtung:\\
\begin{quote}
	
	\begin{align*}
	\mathrm{F}(x)&=\int\limits_{\frac{1}{2}}^{x} \ln(t)\mathrm{d}t=\bigg[t(\ln(t)-1)\bigg]^x_{\frac{1}{2}}\;\;|\text{\tiny patielle Integration von $ \ln(t) $ }\\
	&=x(\ln(x)-1)-\frac{1}{2}(\ln(\genfrac{}{}{}{2}{1}{2})-1)
	\end{align*}
	$ \bullet \;\;\lim\limits_ {x\to\infty}\mathrm{F}(x)=\lim\limits_ {x\to\infty}x(\ln(x)-1)-\frac{1}{2}(\ln(\genfrac{}{}{}{2}{1 }{2 })-1)=\infty\Rightarrow\text{kein globales Maximum }$\\
	$ \bullet\;\;\lim\limits_ {\substack{x\to \frac{1}{2}\\x>\frac{1}{2}}}\mathrm{F}(x)= \lim\limits_ {\substack{x\to \frac{1}{2}\\x>\frac{1}{2}}}x(\ln(x)-1)-\frac{1}{2}(\ln(\genfrac{}{}{}{2}{1}{2 })-1)=0$
\end{quote}
Wenn der Funktionswert vom lokalen Minimum bei $ x_0=1 $  kleiner als Null ist, dann wird es wegen des  Grenzwertes bei $\lim\limits_ {\substack{x\to \frac{1}{2}\\x>\frac{1}{2}}}\mathrm{F}(x)= 0  $ ein globales Minimum.\\
Es gilt bei $ x_0=1 $:\\
\begin{align*}
\mathrm{F}(x)&=x(\ln(x)-1)-\frac{1}{2}(\ln(\genfrac{}{}{}{2}{1}{2})-1)\\
\mathrm{F}(1)&=1(\underbrace {\ln(1)}_{=0}-1)-\frac{1}{2}(\ln(\genfrac{}{}{}{2}{1}{2})-1)\\
&=-1-\frac{1}{2}\ln(\genfrac{}{}{}{2}{1 }{2 })+\frac{1}{2}=-\frac{1}{2}-\frac{1}{2}\ln(\genfrac{}{}{}{2}{1}{2 })=\frac{-1-\ln(\genfrac{}{}{}{2}{1 }{2 })}{2}=\frac{-1((\ln(1)-\ln(2)))}{2}\\
&=\frac{-1+\ln(2)}{2}
\end{align*}
Damit der Bruch kleiner Null wird, muß der $ \ln(2) <1$ sein:\\
Es gilt:\\
\begin{align*}
&\ln(2)<\ln(\mathrm{e})=1\\
&\Rightarrow\\
&\mathrm{F}(1)=\frac{-1+\ln(2)}{2}<0\\
&\Rightarrow\text{globales Minimum bei }\;\;x_0=1
\end{align*}
\subsection{Aufgabe 15 - Partialbruchzerlegung}
\href{https://youtu.be/H3AxdhKnQPc}{Video 18}\\
Bestimme alle Stammfunktionen von:\\
\begin{align*}
f(x)=\frac{x^4+4x^3+4}{x^4-1}
\end{align*}
Durch Erweiterung des Nenner mit $ +1-1 $ erreicht man,dass man im Zähler ein $ x^4-1 $ erhält und dies sich dann beim Zerlegenvon $ f(x) $ in einzelnen Summanten zu 1 kürzt.\\
\begin{align*}
f(x)=\frac{x^4-1+1+4x^3+4}{x^4-1}=\frac{x^4-1+4x^3+5}{x^4-1}
\end{align*}
Nun den Bruch in einzelne Summanten zerlegen:\\
\begin{align*}
f(x)=\frac{x^4-1+4x^3+5}{x^4-1}=\frac{x^4-1}{x^4-1}+\frac{4x^3}{x^4-1}+\frac{5}{x^4-1}=1+\frac{4x^3}{x^4-1}+\frac{5}{x^4-1}
\end{align*}
Nun die einzelnen Stammfunktion bilden:\\
$ \bullet $\;\;Stammfunktion von : $f_1(x) =1 $:\\
\begin{align*}
\mathrm{F}_1(x)=x+\mathrm{C}
\end{align*}
$ \bullet $\;\;Stammfunktion von : $f_2(x)=\frac{4x^3}{x^4-1}  $:\\
\begin{quote}
	Wenn im Zähler die Ableitung vom Nenner steht dann gilt:\\
	\begin{align*}
	\mathrm{F}_2(x)=\ln(x^4-1)+\mathrm{c}
	\end{align*}
\end{quote}
$ \bullet $\;\;Stammfunktion ohne Faktor 5 von: $f_3(x)=\frac{1}{x^4-1}   $\\

\begin{quote}
	Es gilt mit Partialbruchzerlegung für : $ f_3(x)=\frac{1}{x^4-1} $:\\
	\begin{align*}
	\frac{1}{x^4-1}=\frac{1}{(x^2-1)(x^2+1)}=\frac{1}{(x-1)(x+1)(x^2+1)}
	\end{align*}
	Es gilt:\\
	\begin{align*}
	\frac{1}{(x-1)(x+1)(x^2+1)}=\frac{\mathrm{A}}{x-1}+\frac{\mathrm{B}}{(x+1)}+\frac{\mathrm{C}x+\mathrm{D}}{(x^2+1)}
	\end{align*}
	Mit der \glqq Zuhaltemethode \grqq bekommt man:\\
	\begin{align*}
	\frac{1}{(1+1)(1^2+1)}=\frac{1}{4}=\mathrm{A}\\
	\frac{1}{(-1-1)(-1^2+1)}=-\frac{1}{4}=\mathrm{B}\\
	\end{align*}
	Durch einsetzen von A und B:\\
	\begin{align*}
	\frac{1}{(x-1)(x+1)(x^2+1)}&=\frac{\frac{1}{4}}{x-1}-\frac{\frac{1}{4}}{(x+1)}+\frac{\mathrm{C}x+\mathrm{D}}{(x^2+1)}\;\;|\cdot \scriptscriptstyle (x-1)(x+1)(x^2+1)\\
	1&=\frac{1}{4}(x+1)(x^2+1)-\frac{1}{4}(x-1)(x^2+1)+(\mathrm{C}x+\mathrm{D})(x-1)(x+1)\\
	&=\frac{1}{4}(x^3+x^2+x+1)-\frac{1}{4}(x^3-x^2+x-1)+(\mathrm{Cx}+\mathrm{D})(x^2-1)\\
		&=\frac{1}{4}(x^3+x^2+x+1)-\frac{1}{4}(x^3-x^2+x-1)+(\mathrm{C}x^3+\mathrm{D}x^2-\mathrm{C}-\mathrm{D})\\
	0\cdot x^3+	0\cdot x^2+	0\cdot x+1&=x^3\bigg(\frac{1}{4}-\frac{1}{4}+\mathrm{C}\bigg)+x^2\bigg(\frac{1}{4}+\frac{1}{4}+\mathrm{D}\bigg)+x\bigg(\frac{1}{4}-\frac{1}{4}-\mathrm{C}\bigg)+\frac{1}{4}+\frac{1}{4}-\mathrm{D}
	\end{align*}
	Mit Koeffizientenvergleich bekommt man:\\
	\begin{align*}
	0&=\frac{1}{4}-\frac{1}{4}+\mathrm{C}\Rightarrow\mathrm{C}=0\\
	1&=\frac{1}{4}+\frac{1}{4}-\mathrm{D}\Rightarrow\mathrm{D}=-\frac{1}{2}
	\end{align*}
	Für $ f_3(x) $ gilt nun:\\
	\begin{align*}
	f_3(x)&=\frac{1}{4}\bigg(\frac{1}{(x-1)}\bigg)-\frac{1}{4}\bigg(\frac{1}{(x+1)}\bigg)-\frac{1}{2}\bigg(\frac{1}{(x^2+1)}\bigg)\\
	\mathrm{F}_3(x)&=\frac{1}{4}\ln(x-1)-\frac{1}{4}\ln(x+1)-\frac{1}{2}\arctan(x^2+1)+\mathrm{c}
	\end{align*}
\end{quote}
Die Stammfunktionen von $ f(x) $ lauten nun:\\
\begin{align*}
\mathrm{F}(x)&=\mathrm{F_1}(x)+\mathrm{F_2}(x)+5(\mathrm{F_3}(x))\\
&=x+\ln(x^4-1)+5\bigg[\frac{1}{4}\ln(x-1)-\frac{1}{4}\ln(x+1)-\frac{1}{2}\arctan(x+1)\bigg]
\end{align*}





\end{document}
